\documentclass[11pt,oneside]{amsbook}

\usepackage[T1]{fontenc}
\usepackage[utf8]{inputenc}
\usepackage[spanish,english]{babel}
\usepackage{lmodern}
\usepackage{microtype}
\usepackage{amsmath,amssymb,mathtools}
\usepackage{ifthen}
\usepackage[all]{xy}
\usepackage[
  a4paper,
  inner=31mm,
  outer=31mm,
  top=28mm,
  bottom=30mm,
  marginparwidth=22mm,
  marginparsep=5mm
]{geometry}
\usepackage[hidelinks,pdfencoding=auto,hypertexnames=false]{hyperref}
\usepackage{bookmark}
\makeindex

\newtheorem{proposition}{Proposición}[section]
\newtheorem{theorem}[proposition]{Teorema}
\newtheorem{theoreme}[proposition]{Teorema}
\newtheorem{corollary}[proposition]{Corolario}
\newtheorem{corollaire}[proposition]{Corolario}
\newtheorem{corollaries}[proposition]{Corolarios}
\newtheorem{lemma}[proposition]{Lema}
\newtheorem{lemme}[proposition]{Lema}
\newtheorem{definition}[proposition]{Definición}
\newtheorem{sectiondefinition}{Definición}[section]
\newtheorem{theoremedepurete}[proposition]{Teorema de pureza}
\newtheorem{theoremedefinition}[proposition]{Teorema y definición}
\newtheorem{subproposition}[subsubsection]{Proposición}
\newtheorem{subdefinition}[subsubsection]{Definición}
\newtheorem{sublemme}[subsubsection]{Sublema}
\newtheorem{subremarque}[subsubsection]{Observación}
\newtheorem{example}[proposition]{Ejemplo}
\newtheorem{exemples}[proposition]{Ejemplos}
\newtheorem{exemple}[proposition]{Ejemplo}
\newenvironment{definitions}
  {\refstepcounter{proposition}\par\medskip\noindent
   \textbf{Definiciones \theproposition.}\ }
  {\par\medskip}
\newcommand{\sganum}{\refstepcounter{proposition}\par\medskip\noindent
  \textbf{\theproposition.}\ }
\newtheorem*{serrecriterion}{Criterio de Serre}
\newtheorem*{theoremstar}{Teorema}
\newtheorem*{corollairestar}{Corolario}
\newtheorem*{propositionstar}{Proposición}
\newtheorem*{definitionstar}{Definición}
\newenvironment{enonce*}[1]
  {\par\medskip\noindent\textbf{#1.}\ \itshape}
  {\par\medskip}
\theoremstyle{remark}
\newtheorem{remarques}[proposition]{Observaciones}
\newtheorem{remarque}[proposition]{Observación}
\newtheorem*{remarquestar}{Observación}
\newtheorem*{remarquesstar}{Observaciones}
\newtheorem*{scholiestar}{Escolio}
\theoremstyle{plain}
\newenvironment{remarqueMR}
  {\par\medskip\noindent\lcrochetbf\ \textit{Añadido en 2003 (MR):}\ }
  {\rcrochetbf\par\medskip}
\newenvironment{enumerateb}{\begin{enumerate}}{\end{enumerate}}

\newboolean{orig}
\ifdefined\SGAOrig
  \setboolean{orig}{true}
\else
  \setboolean{orig}{false}
\fi
\providecommand{\Ref}[1]{}
\renewcommand{\Ref}[1]{\ref{#1}}
\newcommand{\Spec}{\operatorname{Spec}}
\newcommand{\goth}[1]{\mathfrak{#1}}
\newcommand{\cal}[1]{\mathcal{#1}}
\newcommand{\ZZ}{\mathbb Z}
\newcommand{\QQ}{\mathbb Q}
\newcommand{\CC}{\mathbb C}
\newcommand{\RR}{\mathbb R}
\newcommand{\PP}{\mathbb P}
\renewcommand{\AA}{\mathbf A}
\newcommand{\VV}{\mathbf V}
\DeclareMathOperator{\gr}{gr}
\newcommand{\cf}{cf.\ }
\newcommand{\ie}{i.e.\ }
\newcommand{\resp}{respectivamente\ }
\newcommand{\og}{«}
\newcommand{\fg}{»}
\providecommand{\oldindexnot}[1]{}
\DeclareMathAlphabet{\othercal}{U}{eus}{m}{n}
\newcommand{\kres}{\kappa}
\DeclareMathOperator{\Hom}{Hom}
\DeclareMathOperator{\hop}{hop}
\DeclareMathOperator{\Gal}{Gal}
\renewcommand{\H}{\mathrm H}
\DeclareMathOperator{\Aut}{Aut}
\DeclareMathOperator{\Isom}{Isom}
\DeclareMathOperator{\Ker}{Ker}
\renewcommand{\Im}{\operatorname{Im}}
\DeclareMathOperator{\Ind}{Ind}
\DeclareMathOperator{\Pro}{Pro}
\DeclareMathOperator{\Prodash}{Pro\text{-}}
\DeclareMathOperator{\Proj}{Proj}
\DeclareMathOperator{\prof}{prof}
\DeclareMathOperator{\coprof}{coprof}
\DeclareMathOperator{\Coker}{Coker}
\DeclareMathOperator{\Norm}{N}
\DeclareMathOperator{\Ext}{Ext}
\DeclareMathOperator{\Tor}{Tor}
\DeclareMathOperator{\Gl}{\mathbf{GL}}
\DeclareMathOperator{\GL}{\mathbf{GL}}
\DeclareMathOperator{\Pic}{Pic}
\DeclareMathOperator{\Quot}{Quot}
\DeclareMathOperator{\Symp}{\mathbf{Sp}}
\DeclareMathOperator{\Supp}{Supp}
\DeclareMathOperator{\codim}{codim}
\DeclareMathOperator{\car}{car}
\DeclareMathOperator{\card}{card}
\DeclareMathOperator{\trace}{Tr}
\DeclareMathOperator{\divisor}{div}
\DeclareMathOperator{\diag}{diag}
\DeclareMathOperator{\Symm}{Sym}
\newcommand{\res}{\mathrm{res}}
\def\sp{\mathrm{sp}}
\newcommand{\red}{\mathrm{red}}
\newcommand{\R}{\mathrm{R}}
\newcommand{\et}{\mathrm{\acute{e}t}}
\newcommand{\an}{\mathrm{an}}
\newcommand{\alg}{\mathrm{alg}}
\newcommand{\GG}{\mathbf{G}}
\newcommand{\WW}{\mathbf{W}}
\newcommand{\bbalpha}{\boldsymbol{\alpha}}
\newcommand{\bbmu}{\boldsymbol{\mu}}
\newcommand{\LL}{\mathbb L}
\newcommand{\tame}{\mathrm{t}}
\newcommand{\leftexp}[2]{{\vphantom{#2}}^{#1}{#2}}
\newcommand{\skpt}{}
\let\Subsection\subsection
\let\Subsubsection\subsubsection
\newcommand{\Ens}{\mathbf{Ens}}
\newcommand{\Univ}{\mathfrak U}
\newcommand{\Cat}{\mathbf{Cat}}
\newcommand{\Sch}{\mathbf{Sch}}
\DeclareMathOperator{\CatFl}{\mathbf{Fl}}
\DeclareMathOperator{\Ob}{Ob}
\DeclareMathOperator{\Fl}{Fl}
\DeclareMathOperator{\SheafHom}{\mathbf{Hom}}
\DeclareMathOperator{\SheafAut}{\mathbf{Aut}}
\DeclareMathOperator{\SheafIsom}{\mathbf{Isom}}
\newcommand{\quoi}{}
\providecommand{\rd}{\mathrm{d}}
\newcommand{\cart}{\mathrm{cart}}
\newcommand{\Fib}{\mathrm{Fib}}
\newcommand{\Ouv}{\operatorname{Ouv}}
\newcommand{\Fer}{\operatorname{Fer}}
\newcommand{\varprojLim}{\varprojlim}
\newcommand{\id}{\mathrm{id}}
\newcommand{\from}{\longleftarrow}
\newcommand{\mto}{\longmapsto}
\newcommand{\To}{\Longrightarrow}
\newcommand{\From}{\Longleftarrow}
\newcommand{\lto}[1]{\xrightarrow{#1}}
\newcommand{\lfrom}[1]{\xleftarrow{#1}}
\let\dpl\displaystyle
\newcommand{\chapterspace}[1]{}
\newcommand{\isomto}{\xrightarrow{\sim}}
\newcommand{\isomfrom}{\xleftarrow{\sim}}
\newcommand{\Cf}{Cf.\ }
\newcommand{\loccit}{loc.\ cit.\ }
\newcommand{\ptbl}{.\ }
\newcommand{\No}{n.º~}
\newcommand{\no}{n.º~}
\newcommand{\numero}{n.º~}
\newcommand{\nieme}{$n$-ésimo}
\newcommand{\niemes}{$n$-ésimas}
\newcommand{\Nieme}{Enésima}
\newcommand{\sss}{si y solo si}
\newcommand{\tH}{{}^t\!H}
\newcommand{\tbigcap}{\mathop{\textstyle\bigcap}\limits}
\newcommand{\tbigwedge}{\mathop{\textstyle\bigwedge}\limits}
\newcommand{\bbullet}{{\scriptscriptstyle\bullet}}
\newcommand{\pr}{\mathrm{pr}}
\newcommand{\lcrochetbf}{{\fontfamily{ptm}\bfseries\upshape[}}
\newcommand{\rcrochetbf}{{\fontfamily{ptm}\bfseries\upshape]}}

\hypersetup{
  pdftitle={Recubrimientos étales y grupo fundamental (SGA 1)},
  pdfauthor={Seminario dirigido por A. Grothendieck; traducción española},
  pdfsubject={Traducción española de SGA 1 producida por IA},
  pdfkeywords={morfismo étale, esquema, grupo fundamental, recubrimiento, descenso}
}

\renewcommand{\thechapter}{\Roman{chapter}}
\renewcommand{\thesection}{\thechapter.\arabic{section}}
\numberwithin{equation}{section}
\setcounter{secnumdepth}{1}
\setcounter{tocdepth}{1}

\makeatletter
\newenvironment{sgabibliography}[2]{%
  \section{Bibliografía}\label{#2}%
  \normalfont\bibliofont\labelsep .5em\relax
  \renewcommand\theenumiv{\arabic{enumiv}}\let\p@enumiv\@empty
  \list{\@biblabel{\theenumiv}}{\settowidth\labelwidth{\@biblabel{#1}}%
    \leftmargin\labelwidth \advance\leftmargin\labelsep
    \usecounter{enumiv}}%
  \sloppy \clubpenalty\@M \widowpenalty\clubpenalty
  \sfcode`\.=\@m
}{%
  \def\@noitemerr{\@latex@warning{Empty `sgabibliography' environment}}%
  \endlist
}
\makeatother

\begin{document}
\begingroup\small
SGA 1: seminario dirigido por A. Grothendieck, ampliado con dos exposiciones de M. Raynaud. La fuente francesa es la edición \href{https://arxiv.org/abs/math/0206203v2}{arXiv math/0206203v2}. Se conservan la recomposición LaTeX de los voluntarios del proyecto dirigido por Bas Edixhoven, las notas de Michel Raynaud y la contribución de la Société mathématique de France, con los créditos de los prefacios originales. La traducción al español y las posteriores correcciones y composición son producidas por IA; no se atribuyen a los colaboradores de la edición francesa.

La traducción al español está producida por IA. Las correcciones, la edición y la composición también están producidas por IA. El responsable de la publicación no tradujo el texto. Se conservan la autoría y las atribuciones de las obras originales y de las contribuciones heredadas; no se atribuyen a la IA. Los modelos documentados se identifican en cada edición.
\par\endgroup\clearpage
\begingroup\small
Los derechos sobre las obras originales y sobre las contribuciones protegidas a estas traducciones y transcripciones derivadas corresponden a sus respectivos titulares. Todo aquello sobre lo que el editor tenga derechos en estas ediciones queda dedicado al dominio público bajo CC0 1.0 (https://creativecommons.org/publicdomain/zero/1.0/deed.es). Las menciones de derechos de terceros y las autorizaciones ya concedidas siguen siendo aplicables.

Estas ediciones se ofrecen gratuitamente, con fines no comerciales y exclusivamente educativos, bajo la presunción de que los respectivos titulares de derechos no se opondrían a que estas traducciones y transcripciones estuvieran disponibles gratuitamente junto a los originales franceses en \href{https://www.numdam.org/}{NUMDAM}. Esta presunción no significa que se haya obtenido autorización; el acceso gratuito a los originales no autoriza por sí solo su traducción, adaptación o republicación.

Si usted es titular de derechos sobre algún material presentado aquí y se opone a su disponibilidad, puede \href{https://github.com/KokunoYumeto/sga-es/issues}{abrir una incidencia en GitHub}. El responsable de la publicación retirará el material afectado a petición del titular de los derechos.

\par\endgroup\clearpage

\selectlanguage{spanish}
\frontmatter
\begin{titlepage}
  \centering
  \vspace*{18mm}
  {\Large Seminario de Geometría Algebraica del Bois Marie\par}
  \vspace{3mm}
  {\Large 1960--61\par}
  \vspace{16mm}
  {\Huge\bfseries Recubrimientos étales y grupo fundamental\par}
  \vspace{4mm}
  {\LARGE (SGA~1)\par}
  \vspace{20mm}
  {\large Seminario dirigido por A.\ Grothendieck\par}
  \vspace{2mm}
  {\normalsize Ampliado con dos exposiciones de la Sra.\ M.\ Raynaud\par}
  \vfill
  {\small\itshape Versión española modificada, basada en la edición francesa corregida arXiv math/0206203v2\par}
\end{titlepage}

\chapter*{Resumen}
Este volumen es una edición recompuesta y anotada del libro «Recubrimientos étales y grupo fundamental», Lecture Notes in Mathematics, 224, Springer-Verlag, Berlín-Heidelberg-Nueva York, 1971, de Alexander Grothendieck y otros.

El texto expone los fundamentos de una teoría del grupo fundamental en geometría algebraica desde el punto de vista «kroneckeriano», que permite tratar en pie de igualdad tanto el caso de una variedad algebraica en el sentido usual como, por ejemplo, el del anillo de enteros de un cuerpo de números.

\begin{otherlanguage}{english}
\chapter*{Abstract}
This volume is a new updated edition of the book ``Rev\^etements \'Etales et Groupe Fondamental'', Lecture Notes in Mathematics, 224, Springer-Verlag, Berlin-Heidelberg-New York, 1971, by Alexander Grothendieck et al.

The text presents the foundations of a theory of the fundamental group in Algebraic Geometry from the Kronecker point of view, allowing one to treat on an equal footing the case of an algebraic variety in the usual sense, and that of the ring of integers in a number field, for instance.
\end{otherlanguage}

\medskip
\noindent\textbf{Clasificación 2020 de materias de matemáticas:} 14-02, 14A15, 14B25, 14D15, 14E20, 14F35, 14H30.

\medskip
\noindent\textbf{Palabras clave:} morfismo étale, morfismo liso, morfismo plano, esquema, grupo fundamental, recubrimiento, teoría del descenso, especialización.

\begin{otherlanguage}{english}
\medskip
\noindent\textbf{Keywords:} étale morphism, smooth morphism, flat morphism, scheme, fundamental group, covering, theory of descent, specialization.
\end{otherlanguage}

\cleardoublepage

\renewcommand{\baselinestretch}{1.05}\normalfont
\chapter*{Prefacio}
\thispagestyle{empty}
\vfill

El presente texto es una edición recompuesta
\ifthenelse{\boolean{orig}}{\ignorespaces}{y anotada}
del libro «Recubrimientos étales y grupo fundamental», Lecture Notes in Mathematics, 224, Springer-Verlag, Berlín-Heidelberg-Nueva York, 1971, de Alexander Grothendieck y otros.

La composición en \LaTeX 2e fue realizada por voluntarios que participaron en un proyecto dirigido por Bas Edixhoven, sobre el cual pueden encontrarse más detalles en \url{http://www.math.leidenuniv.nl/~edix/}. La maquetación fue terminada por la Société mathématique de France.

\ifthenelse{\boolean{orig}}
{Se supone que esta versión reproduce el \emph{texto original}. Existe otra versión en la que se han corregido algunos errores tipográficos y Michel Raynaud ha realizado algunas actualizaciones. Esta última es publicada por la Société mathématique de France (\url{http://smf.emath.fr/}) en la serie «Documents Mathématiques».}
{Esta es una \emph{versión ligeramente corregida} del texto original. Es publicada por la Société mathématique de France (\url{http://smf.emath.fr/}) en la serie «Documents Mathématiques». Michel Raynaud ha añadido algunas observaciones de actualización. Estas están delimitadas por corchetes \lcrochetbf\,\rcrochetbf\ e indicadas con el símbolo (MR) (observaciones en las páginas \pageref{X.2.14}, \pageref{XI.1.4}, \pageref{XII.5.6}, \pageref{XIII.2.13} y nota \Ref{III.6.6.p24} en la página \pageref{III.6.6.p24}). Para uniformar la notación, el cuerpo residual de un punto $x$ se denota por $\kres(x)$, y el de un anillo local $A$, por $\kres(A)$.\par
Existe asimismo una versión electrónica que pretende reproducir el texto original.}

Las dos versiones se producen a partir de un único archivo fuente, alojado en el servidor de prepublicaciones de arXiv.org, \url{http://arxiv.org/}. Todas las diferencias entre ambas versiones están documentadas en el archivo fuente.

La antigua numeración de páginas se incorpora en el margen (el número~$n$ indica el comienzo de la página~$n$).

\vfill
\newpage

\begin{otherlanguage}{english}
The present text is a new
\ifthenelse{\boolean{orig}}{\ignorespaces}{updated}
edition of the book ``Rev\^etements \'Etales et Groupe Fondamental'', Lecture Notes in Mathematics, 224, Springer-Verlag, Berlin-Heidelberg-New York, 1971, by Alexander Grothendieck et al.

Typesetting in \LaTeX 2e was done by volunteers, participating in a project directed by Bas Edixhoven, on which one can find more details at \url{http://www.math.leidenuniv.nl/~edix/}. It has been completed by the Société mathématique de France.

\ifthenelse{\boolean{orig}}
{This version is meant to reproduce the \emph{original text}. There is another version, in which some typograpical mistakes have been corrected and some notes have been added by Michel Raynaud. The latter has been published by the Société mathématique de France (\url{http://smf.emath.fr/}) as a volume of the series ``Documents Mathématiques''.}
{This is a slightly \emph{corrected version} of the original text. It is published by the Société mathématique de France (\url{http://smf.emath.fr/}) as a volume of the series ``Documents Mathématiques''. Some updating remarks have been added by Michel Raynaud. These are bounded by brackets \lcrochetbf\,\rcrochetbf\ and indicated by the symbol (MR) (Remarks on pages \pageref{X.2.14}, \pageref{XI.1.4}, \pageref{XII.5.6}, \pageref{XIII.2.13} and footnote \Ref{III.6.6.p24} on page \pageref{III.6.6.p24}). In order to unify notation, the residue field of a point $x$ is denoted by $\kres(x)$ and the residue field of a local ring $A$ is denoted by $\kres(A)$.\par
There also exists an electronic version that is meant to reproduce the original text.}

Both versions are produced from the same source file, that can be downloaded from the arXiv.org e-Print server at \url{http://arxiv.org/}. All differences between the two versions are documented in the source file.

The old page numbering is incorporated in the margin (the number~$n$ marking the beginning of page~$n$).
\end{otherlanguage}

\vfill
\cleardoublepage

\renewcommand{\baselinestretch}{1.1}\normalfont
\chapter*{Introducción}
\thispagestyle{empty}
\label{I.0}

En la primera parte de esta introducción damos algunas precisiones acerca del contenido del presente volumen; en la segunda, acerca del conjunto del «\emph{Seminario de Geometría Algebraica del Bois Marie}», del cual el presente volumen constituye el primer tomo.

\let\oldthesubsection\thesubsection
\renewcommand{\thesubsection}{\arabic{subsection}}
\setcounter{secnumdepth}{2}
\subsection{}
\label{I.introduction.1}

El presente volumen expone los fundamentos de una teoría del grupo fundamental en geometría algebraica desde el punto de vista «kroneckeriano», que permite tratar en pie de igualdad tanto el caso de una variedad algebraica en el sentido usual como, por ejemplo, el de un anillo de enteros de un cuerpo de números. Este punto de vista solo puede expresarse satisfactoriamente en el lenguaje de los esquemas, y utilizaremos libremente tanto este lenguaje como los principales resultados expuestos en los tres primeros capítulos de los \emph{Elementos de geometría algebraica} de J.\ \textsc{Dieudonné} y A.\ \textsc{Grothendieck} (citados en lo sucesivo como EGA). El estudio del presente volumen del «\emph{Seminario de Geometría Algebraica del Bois Marie}» no requiere otros conocimientos de geometría algebraica y puede, por tanto, servir de introducción a sus técnicas actuales para quien desee familiarizarse con ellas.

Las exposiciones \Ref{I} a \Ref{XI} de este libro reproducen textualmente, casi sin cambios, las notas mimeografiadas del seminario oral, que fueron distribuidas por el \emph{Institut des Hautes Études Scientifiques\footnote{Así como las notas de los seminarios que siguieron a este. Como este modo de distribución resultó a la larga impracticable e insuficiente, todos los «\emph{Seminarios de Geometría Algebraica del Bois Marie}» aparecerán en lo sucesivo en forma de libro, como el presente volumen.}}. Nos hemos limitado a añadir algunas notas al pie al texto primitivo, corregir algunas erratas y efectuar un ajuste terminológico; en particular, el término «morfismo simple» fue sustituido entretanto por «morfismo liso», que no da lugar a las mismas confusiones.

Las exposiciones \Ref{I} a \Ref{IV} presentan las nociones locales de morfismo \emph{étale} y morfismo \emph{liso}; apenas utilizan el lenguaje de los esquemas, expuesto en el capítulo~I de los \emph{Elementos}\footnote{Ahora se dispone de un estudio más completo en los \emph{Elementos}, cap.~IV, \S\S~17 y 18.}. La exposición~\Ref{V} presenta la descripción axiomática del grupo fundamental de un esquema, útil incluso en el caso clásico en que dicho esquema se reduce al espectro de un cuerpo, donde se obtiene una reformulación muy cómoda de la teoría de Galois usual. Las exposiciones \Ref{VI} y \Ref{VIII} presentan la \emph{teoría del descenso}, que ha adquirido una importancia creciente en geometría algebraica durante los últimos años y que podría prestar servicios análogos en geometría analítica y en topología. Debe señalarse que la exposición VII no llegó a redactarse y que su sustancia se halla
\ifthenelse{\boolean{orig}}
{incorporado}
{incorporada}
en un trabajo de J.\ Giraud (\emph{Méthode de la Descente}, Bull.\ Soc.\ Math.\ France, Mémoire~2, 1964, viii + 150~p.). En la exposición \Ref{IX} se estudia más específicamente el descenso de los morfismos étales, con lo que se obtiene un enfoque sistemático de los teoremas de tipo \textsc{Van Kampen} para el grupo fundamental, que aparecen aquí como simples traducciones de teoremas de descenso. Se trata esencialmente de un procedimiento para calcular el grupo fundamental de un esquema conexo~$X$, provisto de un morfismo propio y sobreyectivo, digamos $X'\to X$, en términos de los grupos fundamentales de las componentes conexas de~$X'$ y de los productos fibrados~$X'\times_X X'$, $X'\times_X X'\times_X X'$, así como de los homomorfismos inducidos entre esos grupos por los morfismos simpliciales canónicos entre los esquemas anteriores. La exposición~\Ref{X} ofrece la teoría de la \emph{especialización del grupo fundamental} para un morfismo propio y liso; su resultado más notable consiste en determinar, salvo pequeños matices, el grupo fundamental de una curva algebraica lisa en característica $p>0$ a partir del resultado conocido por vía trascendente en característica cero. La exposición~\Ref{XI} da algunos ejemplos y complementos; en particular, explicita en términos cohomológicos la teoría de los recubrimientos de \emph{\textsc{Kummer}} y la de \emph{\textsc{Artin}-\textsc{Schreier}}. Para otros comentarios sobre el texto, véase la \emph{Advertencia} de la versión mimeografiada, que sigue a la presente Introducción.

Desde la redacción, en 1961, del presente Seminario, M.\ \textsc{Artin} y yo hemos desarrollado, en colaboración, el lenguaje de la \emph{topología étale} y una teoría cohomológica correspondiente, expuesta en la parte SGA~4, «Cohomología étale de los esquemas», del \emph{Seminario de Geometría Algebraica}, que aparecerá en la misma serie que el presente volumen. Este lenguaje y los resultados ya disponibles proporcionan una herramienta particularmente flexible para estudiar el grupo fundamental, que permite comprender mejor y superar algunos de los resultados aquí expuestos. Convendría, pues, rehacer por completo la teoría del grupo fundamental desde este punto de vista (todos los resultados fundamentales figuran, de hecho, ya en \emph{loc. cit.}). Eso era lo proyectado para el capítulo de los \emph{Elementos} dedicado al grupo fundamental, que debía contener además varios desarrollos que no pudieron encontrar cabida aquí, basados en la técnica de resolución de singularidades: el cálculo del «grupo fundamental local» de un anillo local completo en términos de una resolución de singularidades adecuada de dicho anillo; fórmulas de Künneth locales y globales para el grupo fundamental sin suponer que los morfismos sean propios (cf.\ exposición \Ref{XIII}); y los resultados de M.\ \textsc{Artin} sobre la comparación de los grupos fundamentales locales de un anillo local henseliano excelente y de su completado (SGA~4~XIX). Señalemos asimismo la necesidad de desarrollar una teoría del grupo fundamental de un topos, que abarque a la vez la teoría topológica ordinaria, su versión semisimplicial, la variante «profinita»
\ifthenelse{\boolean{orig}}
{desarrollado}
{desarrollada}
en la exposición~\Ref{V} del presente volumen y la variante prodiscreta, algo más general, de SGA~3~X~7 (adaptada al caso de esquemas no normales y no unirramos). En espera de una refundición completa de la teoría desde esta perspectiva, la exposición \Ref{XIII} de la Sra.\ \textsc{Raynaud}, que utiliza el lenguaje y los resultados de SGA~4, pretende mostrar las posibilidades que estos ofrecen en algunas cuestiones típicas, generalizando en particular ciertos resultados de la exposición~\Ref{X} a esquemas relativos no propios. Allí se da, en particular, la estructura del grupo fundamental «primo con $p$» de una curva algebraica no completa en
\ifthenelse{\boolean{orig}}
{car.}
{característica}
arbitraria (que yo había
\ifthenelse{\boolean{orig}}
{anunciado}
{anunciado}
en 1959, pero cuya demostración no se había publicado hasta ahora).

Pese a estas numerosas lagunas e imperfecciones (otros dirán: a causa de ellas), pienso que el presente volumen podrá ser útil tanto a quien desee familiarizarse con la teoría del grupo fundamental como en calidad de obra de referencia, mientras se redacta y publica un texto que no incurra en las críticas que acabo de enumerar.

\subsection{}
\label{I.introduction.2}

El presente volumen constituye el tomo~1 del «\emph{Seminario de Geometría Algebraica del Bois Marie}», cuyos volúmenes siguientes está previsto que aparezcan en la misma serie que este. El objetivo que se propone el \emph{Seminario}, paralelamente al tratado «\emph{Elementos de geometría algebraica}» de J\ptbl\ \textsc{Dieudonné} y A\ptbl\ \textsc{Grothendieck}, es sentar los fundamentos de la geometría algebraica desde los puntos de vista expuestos en esta última obra. La referencia estándar para todos los volúmenes del \emph{Seminario} está constituida por los capítulos~I, II y III de los «\emph{Elementos de geometría algebraica}» (citados como EGA~I, II y III), y se supone que el lector posee el bagaje de álgebra conmutativa y álgebra homológica que estos capítulos requieren\footnote{Véase la Introducción de EGA~I para más precisiones al respecto.}. Además, en cada volumen del \emph{Seminario} se remitirá libremente, en la medida de lo necesario, a volúmenes anteriores del mismo \emph{Seminario} o a otros capítulos publicados o a punto de aparecer de los «Elementos».

Cada \emph{parte} del \emph{Seminario} se centra en un tema principal, indicado en el título del volumen o de los volúmenes correspondientes; el seminario oral abarca generalmente un año académico, a veces más. Las exposiciones comprendidas en cada parte del \emph{Seminario} mantienen generalmente entre sí una estrecha dependencia lógica; en cambio, las diferentes partes del \emph{Seminario} son, en gran medida, lógicamente independientes unas de otras. Así, la parte «Esquemas en grupos» es casi enteramente independiente de las dos partes del \emph{Seminario} que la preceden cronológicamente; en cambio, recurre con frecuencia a los resultados de EGA~IV. He aquí la lista de las partes del \emph{Seminario} que han de aparecer próximamente (citadas en lo sucesivo como SGA~1 a SGA~7):
\begin{enumerateb}
\item[SGA 1.] Recubrimientos étales y grupo fundamental, 1960 y 1961.
\item[SGA 2.] Cohomología local de los haces coherentes y teoremas de Lefschetz locales y globales, 1961/62.
\item[SGA 3.] Esquemas en grupos, 1963 y 1964 (3~volúmenes), en colaboración con M\ptbl\ \textsc{Demazure}.
\item[SGA 4.] Teoría de los topoi y cohomología étale de los esquemas, 1963/64 (3~volúmenes), en colaboración con M\ptbl\ \textsc{Artin} y J.L\ptbl\ \textsc{Verdier}.
\item[SGA 5.] Cohomología $l$-ádica y funciones~$L$, 1964 y 1965 (2~volúmenes).
\item[SGA 6.] Teoría de intersecciones y teorema de Riemann--Roch, 1966/67 (2~volúmenes), en colaboración con P\ptbl\ \textsc{Berthelot} y L\ptbl\ \textsc{Illusie}.
\item[SGA 7.] Grupos de monodromía local en geometría algebraica.
\end{enumerateb}

Tres de estos \emph{Seminarios} parciales fueron dirigidos en \emph{colaboración} con otros matemáticos, que figurarán como coautores en la cubierta de los volúmenes correspondientes. En cuanto a los demás participantes activos del \emph{Seminario}, cuyo papel (tanto desde el punto de vista de la redacción como desde el de la puesta a punto matemática) ha ido creciendo de año en año, el nombre de cada participante figura al comienzo de las exposiciones de las que es responsable como conferenciante o como redactor, y la lista de quienes figuran en un volumen determinado se encuentra
\ifthenelse{\boolean{orig}}
{indicado}
{indicada}
en la portada de dicho volumen.

Conviene dar algunas precisiones sobre las relaciones entre el \emph{Seminario} y los \emph{Elementos}. Estos últimos estaban destinados en principio a ofrecer una exposición de conjunto de las nociones y técnicas consideradas más fundamentales en geometría algebraica, a medida que estas mismas nociones y técnicas se fueran desprendiendo mediante el juego natural de exigencias de coherencia lógica y estética. Desde esta perspectiva, era natural considerar el \emph{Seminario} como una versión preliminar de los \emph{Elementos}, destinada a quedar incorporada casi por entero, tarde o temprano, en estos últimos. Este proceso ya había comenzado en cierta medida hace algunos años, puesto que las exposiciones I a IV del presente volumen SGA~1 están enteramente incorporadas en EGA~IV, y las exposiciones VI a VIII debían estarlo dentro de algunos años en EGA~VI. Sin embargo, a medida que se desarrolla la labor de edificación emprendida en los \emph{Elementos} y en el \emph{Seminario}, y se precisan las proporciones del conjunto, el principio inicial (según el cual el \emph{Seminario} no constituiría más que una versión preliminar y provisional) parece cada vez menos realista, debido, entre otras cosas, a los límites que la previsora naturaleza impone a la duración de la vida humana. Habida cuenta del cuidado prestado generalmente a la redacción de las distintas partes del \emph{Seminario}, sin duda solo habrá motivo para retomar una parte semejante en los \emph{Elementos} (o en los tratados que los sucedan) cuando los progresos posteriores a su redacción permitan introducir en ella mejoras muy sustanciales, al precio de modificaciones bastante profundas. Este es ya el caso del presente seminario SGA~1, como se ha dicho más arriba, y también de SGA~2 (gracias a los resultados recientes de la Sra.~\textsc{Raynaud}). En cambio, nada indica actualmente que vaya a ser así en un futuro próximo para ninguna de las partes SGA~3 a SGA~7 citadas más arriba.

Las referencias dentro del «Seminario de Geometría Algebraica del
\ifthenelse{\boolean{orig}}
{Bois-Marie}
{Bois Marie}» se dan del modo siguiente. Una \emph{referencia interna} a una de las partes SGA~1 a SGA~7 del Seminario se da en la forma III~9.7, donde la cifra III designa el número de la exposición, que figura en la parte superior de cada página de la exposición en cuestión, y~9.7, el número del enunciado (o de la definición, observación, etc.) dentro de la exposición. Cuando sea necesario, pueden utilizarse números decimales más largos, por ejemplo 9.7.1 y 9.7.2, para designar, por ejemplo, las distintas etapas de la demostración de una proposición~9.7. La referencia III~9 designa el parágrafo~9 de la exposición~III. El número de la exposición se omite en las referencias internas a una exposición. Para una \emph{referencia a otra} de las \emph{partes} del \emph{Seminario}, se utilizan las mismas siglas, pero precedidas de la mención de la parte de los SGA de que se trate: SGA~1~III~9.7. Del mismo modo, la referencia EGA~IV~11.5.7 significa: \emph{Elementos de geometría algebraica}, cap\ptbl\ IV, enunciado (o definición
\ifthenelse{\boolean{orig}}{etc...}{etc.})~11.5.7; aquí, la primera cifra arábiga sigue designando el número del parágrafo. Aparte de estas convenciones vigentes en el conjunto de los SGA, la bibliografía relativa a una exposición estará reunida generalmente al final de esta, y dentro de la exposición se remitirá a ella mediante números entre corchetes, como
\ifthenelse{\boolean{orig}}{[3]}{[\textbf{3}]}, según el uso habitual.

Por último, para comodidad del lector, siempre que parezca necesario añadiremos al final de los volúmenes de los SGA un índice de notaciones y un índice terminológico que contenga, cuando proceda, una traducción inglesa de los términos franceses utilizados.

Deseo añadir a esta introducción un comentario extramatemático. En noviembre de 1969 tuve conocimiento de que el Institut des Hautes Études Scientifiques, donde he sido profesor casi desde su fundación, recibía desde hacía tres años subvenciones del Ministerio de las Fuerzas Armadas. Ya como investigador principiante me parecía sumamente lamentable la falta de escrúpulos con que la mayoría de los científicos aceptan colaborar, de una forma u otra, con los aparatos militares. Mis motivaciones de entonces eran esencialmente de naturaleza moral y, por tanto, poco susceptibles de ser tomadas en serio. Hoy adquieren una fuerza y una dimensión nuevas, en vista del peligro de destrucción de la especie humana con que nos amenaza la proliferación de los aparatos militares y de los medios de destrucción masiva de que estos disponen. En otros lugares me he explicado con mayor detalle sobre estas cuestiones, mucho más importantes que el progreso de cualquier ciencia (incluida la matemática); puede consultarse, por ejemplo, a este respecto el artículo de G\ptbl\ Edwards en el n.º~1 de la revista \emph{Survivre} (agosto de 1970), que resume una exposición más detallada de estas cuestiones que había hecho en otro lugar. Así pues, durante tres años me encontré trabajando en una institución mientras esta participaba, sin que yo lo supiera, en una modalidad de financiación que considero inmoral y peligrosa\footnote{Huelga decir que la opinión que acabo de expresar es de mi exclusiva responsabilidad y no de la editorial Springer, que publica el presente volumen.}. Al ser actualmente el único de mis colegas del IHES que mantiene esta opinión, lo que ha condenado al fracaso mis esfuerzos por lograr que se supriman del presupuesto del IHES las subvenciones militares, he tomado la decisión que se imponía: dejo el IHES el 30 de septiembre de 1970 y suspendo asimismo toda colaboración científica con esta institución mientras siga aceptando tales subvenciones. He pedido al Sr\ptbl\ Motchane, director del IHES, que, a partir del 1 de octubre de 1970, el IHES se abstenga de difundir textos matemáticos de los que soy autor o que formen parte del Seminario de Geometría Algebraica del Bois Marie. Como se ha dicho más arriba, la difusión de este seminario será asumida por la editorial Julius Springer, en la serie Lecture Notes. Me complace agradecer aquí a la editorial Springer y al Sr\ptbl\ K\ptbl\ Peters la ayuda eficaz y cortés que me han prestado para hacer posible esta publicación, haciéndose cargo, en particular, de la mecanografía para la reproducción fotoóffset de las nuevas exposiciones añadidas a los antiguos seminarios y de las exposiciones que faltaban en los seminarios incompletos.

Agradezco asimismo al Sr\ptbl\ J.P\ptbl\ Delale, que se encargó de la ingrata labor de compilar el índice de notaciones y el índice terminológico.

\bigskip\hfill Massy, agosto de 1970.


\cleardoublepage
\chapter*{Advertencia}
\thispagestyle{empty}
\label{I.avertissement}

Cada una de las exposiciones redactadas recoge el contenido de varias exposiciones orales consecutivas. No ha parecido útil precisar sus fechas.

La exposición~VII, a la que se remite en diversas ocasiones a lo largo de la exposición~\Ref{VIII}, no fue redactada por el conferenciante, quien, en las conferencias orales, se había limitado a esbozar el lenguaje del descenso en las categorías generales, situándose desde un punto de vista estrictamente utilitario y sin entrar en las dificultades lógicas que plantea este lenguaje. Resultó que una exposición correcta de este lenguaje rebasaría los límites de las presentes notas, aunque solo fuera por su extensión. Para una exposición formal de la teoría del descenso, remito a un artículo en preparación de Jean \textsc{Giraud}. Mientras aparece\footnote{Publicado entretanto: J\ptbl\ \textsc{Giraud}, \emph{\ifthenelse{\boolean{orig}}{Méthodes}{Méthode} de la descente}, Mémoire n.º~2 de la Société mathématique de France, 1964.\label{footnotegiraud}}, pienso que un lector atento no tendrá dificultad en suplir por sus propios medios las referencias fantasma de la exposición~VIII.

Tampoco fueron redactadas otras exposiciones orales, posteriores a la exposición~\Ref{XI} y a las que se alude en algunos lugares del texto; estaban destinadas a constituir el contenido de una exposición~\Ref{XII} y de una exposición~\Ref{XIII}. Las primeras de estas exposiciones orales retomaban, en el marco de los esquemas y de los espacios analíticos con elementos nilpotentes (tal como se introducen en el Seminario Cartan 1960/61), la construcción del espacio analítico asociado a un preesquema localmente de tipo finito sobre un cuerpo valuado completo~$k$, los teoremas de tipo GAGA en el caso en que $k$ es el cuerpo de los números complejos y la aplicación a la comparación entre el grupo fundamental definido por vía trascendente y el grupo fundamental estudiado en estas notas (compárese A\ptbl\ Grothendieck, \emph{Fondements de la Géométrie Algébrique}, Seminario Bourbaki n.º~190, página~10, diciembre de 1959). Las últimas exposiciones orales esbozaban la generalización de los métodos desarrollados en el texto para estudiar los recubrimientos que admiten ramificación moderada y la estructura del grupo fundamental de una curva completa menos un número finito de puntos (compárese \loccit, n.º~182, página~27, teorema~14). Estas exposiciones no introducen ninguna idea esencialmente nueva, razón por la cual no ha parecido indispensable darles una redacción formal antes de la aparición de los capítulos correspondientes de los \emph{Elementos de geometría algebraica}\footnote{Se incluyen en el presente volumen, en la exp\ptbl\ \Ref{XII} de la Sra.~Raynaud, con una demostración distinta de la demostración original expuesta en el Seminario oral (\cf introducción).}.

En cambio, los teoremas de tipo Lefschetz para el grupo fundamental y el grupo de Picard, tanto desde el punto de vista local como desde el global, fueron objeto de un Seminario separado en 1962, que fue redactado por completo y está a disposición de los usuarios\footnote{\emph{Cohomologie locale des faisceaux cohérents et théorèmes de Lefschetz locaux et globaux} (SGA~2), publicado por North-Holland Publishing Company.}. Señalemos que los resultados desarrollados tanto en el presente Seminario como en el de 1962 se utilizarán de manera esencial para la publicación de varios resultados clave de la cohomología étale de los preesquemas, que serán objeto de un Seminario (dirigido por M\ptbl\ Artin y por mí) en 1963/64, actualmente en preparación\footnote{\emph{Cohomologie étale des schémas} (citado como SGA~4), de próxima aparición en esta misma serie.}.

Las exposiciones \Ref{I} a~\Ref{IV}, de naturaleza esencialmente local y muy elementales, quedarán enteramente incorporadas en el capítulo~IV de los \emph{Elementos de geometría algebraica}, cuya primera parte está en prensa y que se publicará sin duda hacia finales de 1964. No obstante, podrán resultar útiles al lector que desee ponerse al corriente de las propiedades esenciales de los morfismos lisos, étales o planos antes de adentrarse en los arcanos de un tratado sistemático. En cuanto a las demás exposiciones, quedarán incorporadas en el capítulo~VIII\footnote{De hecho, a raíz de una modificación del plan inicialmente previsto para los \emph{Elementos}, el estudio del grupo fundamental queda aplazado allí hasta un capítulo posterior al que acaba de indicarse. Compárese la introducción que precede a la presente «Advertencia».} de los «Elementos», cuya publicación difícilmente podrá contemplarse antes de varios años.

\bigskip\hfill Bures, junio de 1963.

{\def\footnotemark{}
\chapterspace{-4}
\tableofcontents}

\mainmatter
\renewcommand{\chaptername}{Exposición}
\let\thesubsection\oldthesubsection

\chapterspace{-2}

\chapter{Morfismos étales}
\label{I}

\marginpar{1}

Para simplificar la exposición, se supone que todos los preesquemas considerados son localmente noetherianos, al menos después del número~\Ref{I.2}.

\section{Nociones de cálculo diferencial}
\label{I.1}
Sea $X$ un preesquema sobre~$Y$, y sea $\Delta_{X/Y}$, o $\Delta$,
\label{indnot:ab}\oldindexnot{$\Delta_{X/Y}$ o simplemente $\Delta$|hyperpage}%
el morfismo diagonal $X\to X\times_Y X$. Es un morfismo de inmersión y, por tanto, un morfismo de inmersión cerrada de~$X$ en un abierto~$V$ de~$X\times_Y X$. Sea $\cal{I}_X$ el ideal del subpreesquema cerrado correspondiente a la diagonal en~$V$ (N.B.: si se quieren hacer las cosas de manera intrínseca, sin suponer que $X$ es separado sobre~$Y$ ---hipótesis que sería una tomadura de pelo---, habría que considerar la imagen inversa como haz de conjuntos de~$\cal{O}_{X\times X}$ sobre~$X$, y designar por $\cal{I}_X$ el ideal de aumentación de este últimoldots). El haz $\cal{I}_X/\cal{I}_X^2$ puede considerarse como un haz cuasicoherente sobre~$X$; se denota por~$\mathit{\Omega}^1_{X/Y}$.
\label{indnot:ac}\oldindexnot{$\mathit{\Omega}^1_{X/Y}$|hyperpage}%
Es de tipo finito si $X\to Y$ es de tipo finito. Se comporta bien respecto del cambio de base $Y'\to Y$. Se introducen asimismo los haces $\cal{O}_{X\times_Y X}/\cal{I}_X^{n+1}=\cal{P}^n_{X/Y}$,
\label{indnot:ad}\oldindexnot{$\cal{P}_{X/Y}^n$|hyperpage}%
que son haces de \emph{anillos} sobre~$X$ y dotan a~$X$ de una estructura de preesquema que puede denotarse por $\Delta_{X/Y}^n$
\label{indnot:ae}\oldindexnot{$\Delta_{X/Y}^n$|hyperpage}%
y llamarse el \emph{$n$-ésimo entorno infinitesimal de}~$X/Y$.
\index{entorno infinitesimal de~$X/Y$|hyperpage}%
El sorites correspondiente es de una trivialidad absoluta, aunque bastante largo\footnote{\cf EGA~IV~16.3.}; sería prudente no hablar de él hasta el momento en que pueda decirse algo útil al respecto, con los morfismos lisos.


\section{Morfismos cuasifinitos}
\label{I.2}
\begin{proposition}
\label{I.2.1}
Sea $A\to B$ un homomorfismo local (N.B.: los anillos son ahora noetherianos), y sea \ifthenelse{\boolean{orig}}{$m$}{$\goth{m}$} el ideal maximal de~$A$. Son equivalentes las condiciones siguientes:
\begin{enumerate}
\item[(i)] \ifthenelse{\boolean{orig}}
{$B/mB$ es de dimensión finita sobre~$k=A/m$}
{$B/\goth{m}B$ es de dimensión finita sobre~$k=A/\goth{m}$}.
\item[(ii)] \ifthenelse{\boolean{orig}}
{$mB$ es un ideal de definición, y $B/r(B)=\kappa(B)$ es una extensión de~$k=\kappa(A)$}
{$\goth{m}B$ es un ideal de definición, y $B/\goth{r}(B)=\kres(B)$ es una extensión de~$k=\kres(A)$}.
\item[(iii)] El completado $\widehat{B}$ es finito sobre el completado $\widehat{A}$ de~$A$.
\end{enumerate}
\end{proposition}
Se dice entonces
\marginpar{2}
que $B$ \emph{es cuasifinito}
\index{cuasifinito (morfismo, álgebra)|hyperpage}%
sobre~$A$. Se dice que un morfismo $f\colon X\to Y$ es cuasifinito en~$x$ (o que el $Y$-preesquema $f$ es cuasifinito en~$x$) si $\cal{O}_x$ es cuasifinito sobre~$\cal{O}_{f(x)}$. Esto equivale asimismo a decir que $x$ está \emph{aislado en su fibra}~$f^{-1}(f(x))$.% SGA1-CORR-0003: diplomatic French and original print have f^{-1}(x), with x in X.
Se dice que un morfismo es cuasifinito si lo es en todo punto\footnote{En EGA~II~6.2.3 se supone además que $f$ es de tipo finito.}.
\begin{corollaire}
\label{I.2.2}
Si $A$ es completo, ser cuasifinito equivale a ser finito.
\ifthenelse{\boolean{orig}}
{Podría darse el sorites habitual \textup{(i) (ii) (iii) (iv) (v)} para los morfismos cuasifinitos, pero no parece indispensable aquí.}
{}
\end{corollaire}

\ifthenelse{\boolean{orig}}
{}
{Podría darse el sorites habitual \textup{(i), (ii), (iii), (iv), (v),} para los morfismos cuasifinitos, pero no parece indispensable aquí.}


\section{Morfismos no ramificados o netos}
\label{I.3}
\begin{proposition}
\label{I.3.1}
Sea $f\colon X\to Y$ un morfismo de tipo finito, $x\in X$ y $y=f(x)$. Son equivalentes las condiciones siguientes:
\begin{enumerate}
\item[(i)] \ifthenelse{\boolean{orig}}
{$\cal{O}_x/m_y\cal{O}_x$ es una extensión finita separable de~$\kappa(y)$}
{$\cal{O}_x/\goth{m}_y\cal{O}_x$ es una extensión finita separable de~$\kres(y)$}.
\item[(ii)] $\mathit{\Omega}^1_{X/Y}$ es nulo en~$x$.
\item[(iii)] El morfismo diagonal $\Delta_{X/Y}$ es una inmersión abierta en un entorno de~$x$.
\end{enumerate}
\end{proposition}
Para la implicación (i)$\To$(ii), el lema de Nakayama reduce inmediatamente al caso $Y=\Spec(k)$, $X=\Spec(k')$, donde el resultado es bien conocido y, por lo demás, trivial a partir de la definición de separable; (ii)$\To$(iii) se sigue de una caracterización agradable y fácil de las inmersiones abiertas que utiliza a Krull; (iii)$\To$(i), porque se reduce de nuevo al caso $Y=\Spec(k)$ en el que el morfismo diagonal es una inmersión abierta en todas partes. Hay que demostrar entonces que $X$ es finito y de anillo separable sobre~$k$; para ello se reduce al caso en que $k$ es algebraicamente cerrado. Pero entonces todo punto cerrado de~$X$ es aislado (pues coincide con la imagen inversa de la diagonal por el morfismo $X\to X\times_k X$ definido por~$x$), de donde se sigue que $X$ es finito. Puede suponerse entonces que $X$ está reducido a un punto, de anillo~$A$; por tanto, $A\otimes_k A\to A$ es un isomorfismo, de donde $A=k$, c.q.d.

\begin{definition}
\label{I.3.2}
\begin{enumerate}
\item[a)] Se dice entonces que $f$ es \emph{neto},
\index{neto (morfismo, álgebra)|hyperpage}%
o también \emph{no ramificado},
\index{no ramificado (morfismo, álgebra)|hyperpage}%
en~$x$, o que $X$ es neto, o también no ramificado, en~$x$ sobre~$Y$.
\item[b)] Sea $A\to B$ un homomorfismo local. Se dice que es \emph{neto}, o \emph{no ramificado}, o que $B$ es un álgebra local \emph{neta}, o \emph{no ramificada}, sobre~$A$, si
\ifthenelse{\boolean{orig}}
{$B/r(A)B$ es una extensión finita separable de~$A/r(A)$, \ie si $r(A)B=r(B)$}
{$B/\goth{r}(A)B$ es una extensión finita separable de~$A/\goth{r}(A)$, \ie si $\goth{r}(A)B=\goth{r}(B)$},
y $\kres(B)$ es una extensión separable de~$\kres(A)$\kern1pt\footnote{\Cf los remordimientos en \Ref{III.1.2}.}.
\end{enumerate}
\end{definition}

\begin{remarquesstar}
El hecho
\marginpar{3}
de que $B$ sea neto sobre~$A$ se reconoce ya en los completados de~$A$ y de~$B$. Neto implica cuasifinito.
\end{remarquesstar}

\begin{corollaire}
\label{I.3.3}
El conjunto de los puntos en los que $f$ es neto es abierto.
\end{corollaire}

\begin{corollaire}
\label{I.3.4}
Sean $X'$ y $X$ dos preesquemas de tipo finito sobre~$Y$, y sea $g\colon X'\to X$ un $Y$-morfismo. Si $X$ es neto sobre~$Y$, el morfismo grafo $\Gamma_g\colon X'\to X'\times_Y X$ es una inmersión abierta.% SGA1-CORR-0004: source has codomain X x_Y X.
\end{corollaire}

En efecto, es la imagen inversa del morfismo diagonal $X\to X\times_Y X$ por
$$
g\times_Y\id_X\colon X'\times_Y X\to X\times_Y X.
$$

También puede introducirse el ideal anulador $\goth{d}_{X/Y}$
\label{indnot:af}\oldindexnot{$\goth{d}_{X/Y}$|hyperpage}%
de~$\mathit{\Omega}^1_{X/Y}$, llamado \emph{ideal diferente}
\index{diferente (ideal)|hyperpage}%
\index{ideal diferente|hyperpage}%
de~$X/Y$; define un subpreesquema cerrado de~$X$ que, como conjunto, es el conjunto de los puntos en los que $X/Y$ está ramificado, \ie no es neto.

\begin{proposition}
\label{I.3.5}
\begin{enumerate}
\item[(i)] Toda inmersión es neta.
\item[(ii)] El compuesto de dos morfismos netos es neto.
\item[(iii)] Todo cambio de base de un morfismo neto es neto.
\end{enumerate}
\end{proposition}

Esto se ve indistintamente a partir de (ii) o de (iii) (la segunda me parece más divertida). Por supuesto, también puede precisarse dando enunciados puntuales; esto no es más general que en apariencia (salvo en el marco de la definición~b)) y resulta aburrido. Se obtienen, como de costumbre, los corolarios siguientes:

\begin{corollaries}
\label{I.3.6}
\begin{enumerate}
\item[(iv)] El producto cartesiano de dos morfismos netos es neto.
\item[(v)] Si $gf$ es neto, $f$ es neto.
\item[(vi)] Si $f$ es neto, $f_{\text{\textup{red}}}$ es neto.
\end{enumerate}
\end{corollaries}

\begin{proposition}
\label{I.3.7}
Sea $A\to B$ un homomorfismo local; se supone que la extensión residual $\kres(B)/\kres(A)$ es trivial o que $\kres(A)$ es algebraicamente cerrado. Para que $B/A$ sea neto es necesario y suficiente que $\widehat{B}$ sea, como $\widehat{A}$-álgebra, un cociente de~$\widehat{A}$.
\end{proposition}

\begin{remarquesstar}
\begin{enumerate}
\item[--] Cuando no se supone trivial la extensión residual, puede reducirse a este caso mediante una extensión finita plana adecuada de~$A$ que elimine dicha extensión.
\item[--] Considérese el ejemplo en que $A$ es el anillo local de un punto doble ordinario de una curva y $B$ el de un punto de la normalización: entonces $A\subset B$, $B$ es neto sobre
\marginpar{4}
$A$, con extensión residual trivial, y $\widehat{A}\to\widehat{B}$ es sobreyectivo pero \emph{no inyectivo}. Vamos, pues, a reforzar la noción de netitud.
\end{enumerate}
\end{remarquesstar}


\section{Morfismos étales. Recubrimientos étales}
\label{I.4}
Admitiremos todo cuanto nos sea necesario acerca de los morfismos planos; estos hechos se demostrarán más adelante cuando corresponda\footnote{\Cf exp\ptbl\ \Ref{IV}.}.

\begin{definition}
\label{I.4.1}
\begin{enumerate}
\item[a)] Sea $f\colon X\to Y$ un morfismo de tipo finito. Se dice que $f$ es \emph{étale}
\index{étale (morfismo, álgebra)|hyperpage}%
en~$x$ si $f$ es plano en~$x$ y neto en~$x$. Se dice que $f$ es étale si lo es en todos los puntos. Se dice que $X$ es étale en~$x$ sobre~$Y$, o que es un $Y$-preesquema étale en~$x$, \ifthenelse{\boolean{orig}}{etc...}{etc.}
\item[b)] Sea $f\colon A\to B$ un homomorfismo local. Se dice que $f$ es étale, o que $B$ es étale sobre~$A$, si $B$ es plano y no ramificado sobre~$A$.\footnote{\Cf los remordimientos en \Ref{III.1.2}.}
\end{enumerate}
\end{definition}

\begin{proposition}
\label{I.4.2}
Para que $B/A$ sea étale es necesario y suficiente que $\widehat{B}/\widehat{A}$ lo sea.
\end{proposition}

\marginpar{5}
En efecto, esto es cierto por separado para «neto» y para «plano».

\begin{corollaire}
\label{I.4.3}
Sea $f\colon X\to Y$ de tipo finito y sea $x\in X$. El hecho de que $f$ sea étale en~$x$ solo depende del homomorfismo local $\cal{O}_{f(x)}\to\cal{O}_x$, e incluso únicamente del homomorfismo correspondiente entre los completados.
\end{corollaire}

\begin{corollaire}
\label{I.4.4}
Supongamos que la extensión residual $\kres(A)\to\kres(B)$ es trivial o que $\kres(A)$ es algebraicamente cerrado. Entonces $B/A$ es étale \sss $\widehat{A}\to\widehat{B}$ es un isomorfismo.
\end{corollaire}

Se combinan la planitud y~\Ref{I.3.7}.
\begin{proposition}
\label{I.4.5}
Sea $f\colon X\to Y$ un morfismo de tipo finito. Entonces el conjunto de los puntos en los que es étale es abierto.
\end{proposition}

En efecto, esto es cierto por separado para «neto» y para «plano».

Esta proposición muestra que pueden omitirse los enunciados «puntuales» en el estudio de los morfismos de tipo finito que son étales en alguna parte.
\begin{proposition}
\label{I.4.6}
\begin{enumerate}
\item[(i)] Toda inmersión abierta es étale.
\item[(ii)] El compuesto de dos morfismos étales es étale.
\item[(iii)] Cambio de base.
\end{enumerate}
\end{proposition}

En efecto, (i) es trivial, y para (ii) y (iii) basta observar que esto es cierto para «neto» y para «plano». A decir verdad, también hay enunciados correspondientes para los homomorfismos locales (sin condición de finitud), que en todo caso deberán figurar en el \emph{multiplodoque} (comenzando por el caso neto).

\begin{corollaire}
\label{I.4.7}
El producto cartesiano de dos morfismos étales también lo es.
\end{corollaire}

\begin{corollaire}
\label{I.4.8}
Sean $X$ y $X'$ de tipo finito sobre~$Y$, y sea $g\colon X\to X'$ un $Y$-morfismo. Si $X'$ es no ramificado sobre~$Y$ y $X$ es étale sobre~$Y$, entonces $g$ es étale.
\end{corollaire}

En efecto, $g$ es el compuesto del morfismo grafo $\Gamma_g\colon X\to X\times_Y X'$, que es una inmersión abierta por~\Ref{I.3.4}, y del morfismo de proyección, que es étale por obtenerse del morfismo étale $X\to Y$ mediante el «cambio de base» $X'\to Y$.

\begin{definition}
\label{I.4.9}
Se llama \emph{recubrimiento étale} (\resp neto)
\index{étale (recubrimiento)|hyperpage}%
\index{neto (recubrimiento)|hyperpage}%
\index{recubrimiento étale (respectivamente neto)|hyperpage}%
de~$Y$ a todo $Y$-esquema $X$ que es finito sobre~$Y$ y étale (\resp neto) sobre~$Y$.
\end{definition}

\ifthenelse{\boolean{orig}}{}
{\enlargethispage{.5cm}}%
La primera condición significa que $X$ está definido por un haz coherente de álgebras $\cal{B}$ sobre~$Y$. La segunda significa entonces que $\cal{B}$ es localmente libre sobre~$Y$ (\resp no impone nada), \emph{y} que, además, para todo $y\in Y$, la fibra $\cal{B}(y)=\cal{B}_y\otimes_{\cal{O}_y}\kres(y)$ es un álgebra separable (= producto finito de extensiones finitas separables) sobre~$\kres(y)$.

\begin{proposition}
\label{I.4.10}
Sea $X$ un recubrimiento plano de~$Y$ de grado~$n$ (la definición de este término habría merecido figurar en~\Ref{I.4.9}), definido por un haz coherente localmente libre $\cal{B}$ de álgebras. Se define de la manera bien conocida el homomorfismo traza $\cal{B}\to\cal{A}$ (que es un homomorfismo de $\cal{A}$-módulos, donde $\cal{A}=\cal{O}_Y$). Para que $X$ sea étale es necesario y suficiente que la forma bilineal correspondiente $\trace_{\cal{B}/\cal{A}}xy$ defina un isomorfismo de~$\cal{B}$ sobre~$\check{\cal{B}}$ o, lo que es equivalente, que la \emph{sección discriminante}
$$
d_{X/Y}=d_{\cal{B}/\cal{A}}\in
\Gamma(Y,\tbigwedge^n\check{\cal{B}}\otimes_{\cal{A}}
\tbigwedge^n\check{\cal{B}})
$$
sea invertible, o bien, por último, que el ideal discriminante definido por esta sección sea el ideal unidad.
\end{proposition}
En efecto, se reduce al caso $Y=\Spec(k)$, y entonces se trata de un criterio de separabilidad bien conocido, trivial al pasar a la clausura algebraica de~$k$.

\marginpar{6}
\begin{remarquestar}
Más adelante se obtendrá un enunciado menos trivial, cuando no se suponga \emph{a priori} que $X$ es plano sobre~$Y$, sino que se formule una hipótesis de normalidad.
\end{remarquestar}


\section{La propiedad fundamental de los morfismos étales}
\label{I.5}

\begin{theoreme}
\label{I.5.1}
Sea $f\colon X\to Y$ un morfismo de tipo finito. Para que $f$ sea una inmersión abierta es necesario y suficiente que sea un morfismo \emph{étale y radicial}.
\end{theoreme}

Recordemos que radicial
\index{radicial (morfismo)|hyperpage}%
significa: inyectivo, con extensiones residuales radiciales (y también \ifthenelse{\boolean{orig}}{se}{se}\ % La autoridad corrige aquí «en» por «on».
puede recordar que esto significa que el morfismo sigue siendo inyectivo después de todo cambio de base). La necesidad es trivial; queda la suficiencia. Daremos dos demostraciones distintas, la primera más breve y la segunda más elemental.

1) Todo morfismo plano es abierto; por tanto, puede suponerse (sustituyendo $Y$ por~$f(X)$) que $f$ es un \emph{homeomorfismo} \ifthenelse{\boolean{orig}}{sobre}{sobre}. Después de todo cambio de base seguirá siendo cierto que \ifthenelse{\boolean{orig}}{$f$}{$f$}\ % La autoridad corrige aquí $f'$ por $f$.
es plano, radicial y sobreyectivo, y por tanto un homeomorfismo, \emph{a fortiori} \ifthenelse{\boolean{orig}}{cerrado}{cerrado}.\ % La autoridad corrige aquí el acuerdo de «fermé».
Así pues, $f$ es \emph{propio}. Por consiguiente, $f$ es \emph{finito} (referencia: teorema de Chevalley) y está definido por un haz coherente $\mathcal{B}$ de álgebras. El haz $\mathcal{B}$ es localmente libre y, además, en virtud de la hipótesis, tiene rango~$1$ en todas partes; por tanto, $X=Y$, c.q.d.

2) Puede suponerse que $Y$ y $X$ son \emph{afines}. Se reduce además fácilmente a demostrar lo siguiente: si $Y=\Spec(A)$, con $A$ local, y si $f^{-1}(y)$ no es vacío (siendo $y$ el punto cerrado de~$Y$), entonces $X=Y$ (en efecto, esto implicará que todo $y\in f(X)$ posee un entorno abierto $U$ tal que $X|U=U$). Se tendrá $X=\Spec(B)$ y queremos demostrar que $A=B$. Para ello se reduce a demostrar la afirmación análoga sustituyendo $A$ por~$\hat A$ y $B$ por $B\otimes_A\hat A$ (teniendo en cuenta que $\hat A$ es fielmente plano sobre~$A$). Puede suponerse, pues, que $A$ es \emph{completo}. Sea $x$ el punto situado sobre~$y$; por el corolario~\Ref{I.2.2}, $\mathcal{O}_x$ es finito sobre~$A$ y, por tanto (al ser plano y radicial sobre~$A$), es idéntico a~$A$. Así pues, $X=Y\amalg X'$ (suma disjunta). Como $X$ es radicial sobre~$Y$, $X'$ es vacío. Hemos terminado.
\ifthenelse{\boolean{orig}}{}
{\enlargethispage{.5cm}}%
\begin{corollaire}
\label{I.5.2}
Sea $f\colon X\to Y$ un morfismo que es una \emph{inmersión cerrada} y es \emph{étale}. Si $X$ es conexo, $f$ es un isomorfismo de~$X$ sobre una componente conexa de~$Y$.
\end{corollaire}

En efecto, $f$ es también una inmersión abierta. Se deduce:

\begin{corollaire}
\label{I.5.3}
Sea
\marginpar{7}
$X$ un $Y$-esquema neto, con $Y$ conexo. Entonces toda sección de~$X$ sobre~$Y$ es un isomorfismo de~$Y$ sobre una componente conexa de~$X$. Hay, por tanto, una correspondencia biunívoca entre el conjunto de estas secciones y el conjunto de las componentes conexas $X_i$ de~$X$ tales que la proyección $X_i\to Y$ sea un isomorfismo (o, lo que es equivalente por~\Ref{I.5.1}, sobreyectiva y radicial). En particular, una sección queda determinada por su valor en un punto.
\end{corollaire}

Solo la primera afirmación requiere demostración; por~\Ref{I.5.2}, basta observar que una sección es una inmersión cerrada (pues $X$ es separado sobre~$Y$) y étale en virtud de~\Ref{I.4.8}.

\begin{corollaire}
\label{I.5.4}
Sean $X$ e $Y$ dos preesquemas sobre~$S$, con $X$ neto y separado sobre~$S$ e $Y$ conexo. Sean $f$ y $g$ dos $S$-morfismos de~$Y$ en~$X$, y sea $y$ un punto de~$Y$. Supongamos que $f(y)=g(y)=x$ y que son idénticos los homomorfismos residuales $\kres(x)\to\kres(y)$ definidos por $f$ y~$g$ («$f$ y~$g$ coinciden geométricamente en~$y$»). Entonces $f$ y~$g$ son idénticos.
\end{corollaire}

Se sigue de~\Ref{I.5.3}, reduciendo al caso $Y=S$ mediante la sustitución de $X$ por $X\times_S Y$.

He aquí una variante particularmente importante de~\Ref{I.5.3}.

\begin{theoreme}
\label{I.5.5}
Sean $S$ un preesquema, $X$ e $Y$ dos $S$-preesquemas y $S_0$ un subpreesquema cerrado de~$S$ que tenga el mismo espacio subyacente que~$S$; sean $X_0=X\times_S S_0$ e $Y_0=Y\times_S S_0$ las «restricciones» de~$X$ e~$Y$ a~$S_0$. Supongamos que $X$ es étale sobre~$S$. Entonces la aplicación natural
$$
\Hom_S(Y,X)\to\Hom_{S_0}(Y_0,X_0)
$$
es \emph{biyectiva}.
\end{theoreme}

Se reduce de nuevo al caso $Y=S$, y entonces el resultado se sigue de la descripción «topológica» de las secciones de~$X/Y$ dada en~\Ref{I.5.3}.

\begin{scholiestar}
Este resultado comprende una afirmación de \emph{unicidad} y otra de \emph{existencia} de \emph{morfismos}. También puede expresarse (cuando $X$ e~$Y$ se toman ambos étales sobre~$S$) diciendo que el funtor $X\mto X_0$, de la categoría de los $S$-esquemas étales en la categoría de los $S_0$-esquemas étales, es \emph{plenamente fiel}, \ie establece una \emph{equivalencia} de la primera con una \emph{subcategoría plena} de la segunda. Veremos más adelante que incluso es una equivalencia entre la primera y la segunda (lo que constituirá un teorema de \emph{existencia de $S$-esquemas étales}).
\end{scholiestar}

La
\marginpar{8}
forma siguiente, aparentemente más general, de~\Ref{I.5.5} resulta a menudo cómoda:

\begin{corollaire}[«Teorema de prolongación de levantamientos»]
\label{I.5.6}
\index{prolongación de levantamientos (teorema de)|hyperpage}%
\index{teorema de prolongación de levantamientos|hyperpage}%
Consideremos un diagrama conmutativo
$$
\xymatrix{
X\ar[d]&\ar[l]Y_0\ar[d]\\
S&\ar[l]Y
}
$$
de morfismos, donde $X\to S$ es étale e $Y_0\to Y$ es una inmersión cerrada biyectiva. Entonces puede encontrarse un único morfismo $Y\to X$ que haga conmutativos los dos triángulos correspondientes.
\end{corollaire}

En efecto, sustituyendo $S$ por~$Y$ y $X$ por $X\times_S Y$, se reduce al caso $Y=S$, y entonces se trata del caso particular de~\Ref{I.5.5} con $Y=S$.

Señalemos asimismo la siguiente consecuencia inmediata de~\Ref{I.5.1} (que no hemos \ifthenelse{\boolean{orig}}{dado}{dado}\ % La autoridad corrige «donné» por «donnée».
como corolario~1 para no interrumpir la línea de ideas desarrollada después de~\Ref{I.5.1}):

\begin{proposition}
\label{I.5.7}
Sean $X$ y $X'$ dos preesquemas de tipo finito y planos sobre~$Y$, y sea $g\colon X\to X'$ un $Y$-morfismo. Para que $g$ sea una inmersión abierta (\resp un isomorfismo) es necesario y suficiente que, para todo $y\in Y$, lo sea el morfismo inducido en las fibras
$$
g\otimes_Y\kres(y)\colon X\otimes_Y\kres(y)\to X'\otimes_Y\kres(y).
$$
\end{proposition}

Basta demostrar la suficiencia; como esta es cierta para la noción de sobreyección, se reduce al caso de una inmersión abierta. Por~\Ref{I.5.1}, hay que comprobar que $g$ es \emph{radicial}, lo cual es trivial, y que es \emph{étale}, lo cual se sigue del corolario~\Ref{I.5.9} que figura más abajo.

\begin{corollaire}
\label{I.5.8}
(Debería pasar al \No\Ref{I.3}.) Sean $X$ y $X'$ dos $Y$-preesquemas, $g\colon X\to X'$ un $Y$-morfismo, $x$ un punto de~$X$ e $y$ su proyección sobre~$Y$. Para que $g$ sea cuasifinito (\resp neto) en~$x$ es necesario y suficiente que lo mismo ocurra con $g\otimes_Y\kres(y)$.
\end{corollaire}

En efecto, las dos álgebras sobre~$k\bigl(g(x)\bigr)$ que hay que examinar para cerciorarse de que se tiene un morfismo cuasifinito, \resp neto, en~$x$ son las mismas para~$g$ y para $g\otimes_Y\kres(y)$.

\begin{corollaire}
\label{I.5.9}
Con
\marginpar{9}
las notaciones de~\Ref{I.5.8}, supongamos que $X$ y~$X'$ son planos y de tipo finito sobre~$Y$. Para que $g$ sea plano (\resp étale) en~$x$ es necesario y suficiente que $g\otimes_Y\kres(y)$ lo sea.
\end{corollaire}

Para «plano», el enunciado solo se incluye como recordatorio: es uno de los criterios fundamentales de planitud\footnote{\Cf \Ref{IV.5.9}.}. Para «étale», se deduce de este, teniendo en cuenta~\Ref{I.5.8}.

\section{Aplicación a las extensiones étales de los anillos locales completos}
\label{I.6}

Este número es un caso particular de resultados sobre los preesquemas formales que deberán figurar en el \emph{multiplodoque}. No obstante, aquí se sale del paso de manera más económica, \ie sin la determinación local explícita de los morfismos étales del \No\Ref{I.7} (que utiliza el \emph{Main Theorem}). Esta es quizá razón suficiente para mantener el presente número (incluso en el \emph{multiplodoque}) en este lugar.

\begin{theoreme}
\label{I.6.1}
Sea $A$ un anillo local completo (noetheriano, por supuesto), de cuerpo residual~$k$. Para cada $A$-álgebra~$B$, sea $R(B)=B\otimes_A k$, considerado como $k$-álgebra; depende, pues, funtorialmente de~$B$. Entonces $R$ define una \emph{equivalencia} entre la categoría de las $A$-álgebras \emph{finitas y étales sobre~$A$} y la categoría de las álgebras \emph{separables de rango finito} sobre~$k$.
\end{theoreme}

Ante todo, el funtor en cuestión es plenamente fiel, como se sigue del hecho más general siguiente:

\begin{corollaire}
\label{I.6.2}
Sean $B$ y $B'$ dos $A$-álgebras finitas sobre~$A$. Si $B$ es étale sobre~$A$, entonces la aplicación canónica
$$
\Hom_{\textup{$A$-alg}}(B,B')\to
\Hom_{\textup{$k$-alg}}\bigl(R(B),R(B')\bigr)
$$
es biyectiva.
\end{corollaire}

Se reduce al caso en que $A$ es artiniano (sustituyendo $A$ por~$A/\goth{m}^n$), y entonces se trata de un caso particular de~\Ref{I.5.5}.

Queda por demostrar que, para toda $k$-álgebra finita y separable (por qué no decir étale, que es más corto) $L$, existe una $B$ étale sobre~$A$ tal que $R(B)$ sea isomorfa a~$L$. Puede suponerse que $L$ es una extensión separable de~$k$; como tal, admite un generador~$x$, \ie es isomorfa a un álgebra $k[t]/Fk[t]$, donde $F\in k[t]$ es un polinomio mónico. Se levanta~$F$ a un polinomio mónico $F_1$ de~$A[t]$ y se toma $B=A[t]/F_1A[t]$.


\marginpar{10}
\section{Construcción local de los morfismos no ramificados y étales}
\label{I.7}

\begin{proposition}
\label{I.7.1}
Sean $A$ un anillo noetheriano, $B$ un álgebra finita sobre~$A$, $u$ un generador de~$B$ sobre~$A$, $F\in A[t]$ tal que $F(u)=0$ (no se supone que $F$ sea mónico), $u'=F'(u)$ (donde $F'$ es el polinomio derivado), $\goth{q}$ un ideal primo de~$B$ que no contiene a~$u'$, y $\goth{p}$ su contracción a~$A$. Entonces $B_\goth{q}$ es neto sobre~$A_\goth{p}$.
\end{proposition}

En otros términos, poniendo $Y=\Spec(A)$, $X=\Spec(B)$ y $X_{u'}=\Spec(B_{u'})$, el preesquema $X_{u'}$ es no ramificado sobre~$Y$. El enunciado se sigue del siguiente, más preciso:

\begin{corollaire}
\label{I.7.2}
El ideal diferente de~$B/A$ contiene a $u'B$, y es igual a él si el homomorfismo natural $A[t]/FA[t]\to B$ (que envía $t$ en~$u$) es un isomorfismo.
\end{corollaire}

Sea $J$ el núcleo del homomorfismo $C=A[t]\to B$; este núcleo contiene a $FA[t]$ y es igual a él en el segundo caso considerado en~\Ref{I.7.2}. Como el homomorfismo es sobreyectivo, $\Omega^1_{B/A}$ se identifica con el cociente de~$\Omega^1_{C/A}$ por el submódulo generado por $J\Omega^1_{C/A}$ y $d(J)$ (habría sido preciso explicitar en el \No\Ref{I.1} la definición del homomorfismo~$d$ y el cálculo de~$\Omega^1$ para un álgebra de polinomios). Identificando $\Omega^1_{C/A}$ con~$C$ mediante la base~$dt$, se obtiene $B/B\cdot J'$; por tanto, la diferente está generada por el conjunto~$J'$ de las imágenes en~$B$ de las derivadas de los $G\in J$ (y basta tomar unos $G$ que generen~$J$). Como $F\in J$, \resp $F$ es un generador de~$J$, hemos terminado. (N.B.: debería ponerse~\Ref{I.7.2} como proposición y~\Ref{I.7.1} como corolario.) Se obtiene:

\begin{corollaire}
\label{I.7.3}
Bajo las condiciones de~\Ref{I.7.1}, suponiendo que $F$ es mónico y que $A[t]/FA[t]\to B$ es un isomorfismo, para que $B_\goth{q}$ sea étale sobre~$A_\goth{p}$ \ifthenelse{\boolean{orig}}{es n. y s.}{es necesario y suficiente} que~$\goth{q}$ no contenga a~$u'$.
\end{corollaire}

En efecto, como $B$ es plano sobre~$A$, étale equivale a neto y puede aplicarse~\Ref{I.7.2}.

\begin{corollaire}
\label{I.7.4}
Bajo las condiciones de~\Ref{I.7.3}, para que $B$ sea étale sobre~$A$ \ifthenelse{\boolean{orig}}{es n. y s.}{es necesario y suficiente}\ % La autoridad corrige la abreviatura francesa mal puntuada.
que~$u'$ sea invertible, o también que el ideal generado por $F$ y $F'$ en~$A[t]$ sea el ideal unidad.
\end{corollaire}

El último criterio se sigue del primero y del lema de Nakayama (en~$B$).

Se dice que un polinomio mónico $F\in A[t]$ que tiene la propiedad
\marginpar{11}
enunciada en el corolario~\Ref{I.7.4}
es un \emph{polinomio separable}
\index{separable (polinomio)|hyperpage}%
(si $F$ no es mónico, habría que exigir al menos que el coeficiente de su término de mayor grado fuera invertible; cuando $A$ es un cuerpo se recupera la definición usual).

\begin{corollaire}
\label{I.7.5}
Sea $B$ un álgebra \emph{finita} sobre el anillo \emph{local}~$A$. Supongamos que $K(A)$ es infinito o que $B$ es local. Sea~$n$ el rango de $L=B\otimes_AK(A)$ sobre~$K(A)=k$. Para que $B$ sea neta (\resp étale) sobre~$A$ es necesario y suficiente que $B$ sea isomorfa a un cociente de (\resp isomorfa a) $A[t]/FA[t]$, donde $F$ es un polinomio mónico separable que puede suponerse (\resp que es necesariamente) de grado~$n$.
\end{corollaire}

Solo hay que demostrar la necesidad. Supongamos que $B$ es neta sobre~$A$; entonces $L$ es separable sobre~$k$, y de la hipótesis formulada se sigue que $L/k$ admite un generador~$\xi$, por lo que los $\xi^i$ ($0\leq i<n$) forman una base de~$L$ sobre~$k$. Sea $u\in B$ un levantamiento de~$\xi$; por Nakayama, los $u^i$ ($0\leq i<n$) generan el $A$-módulo~$B$ (\resp forman una base de este). En particular, puede encontrarse un polinomio mónico $F\in A[t]$ tal que $F(u)=0$, y $B$ será isomorfa a un cociente de (\resp isomorfa a) $A[t]/FA[t]$. Por último, en virtud de~\Ref{I.7.4} aplicado a $L/k$, $F$ y~$F'$ generan $A[t]$ \ifthenelse{\boolean{orig}}{módulo}{módulo}\ % La autoridad corrige «module» por «modulo».
$\goth{m}A[t]$; por tanto (por Nakayama en $A[t]/FA[t]$), $F$ y~$F'$ generan~$A[t]$. Hemos terminado.

\begin{theoreme}
\label{I.7.6}
Sean $A$ un anillo local y $A\to\cal{O}$ un homomorfismo local tal que~$\cal{O}$ sea isomorfa a la localización de un álgebra de tipo finito sobre~$A$. Supongamos que \ifthenelse{\boolean{orig}}{$\cal{O}$}{$\cal{O}$}\ % La autoridad corrige $O$ por $\cal O$.
es \emph{neta} sobre~$A$. Entonces pueden encontrarse un $A$-álgebra $B$ integral sobre~$A$, un ideal maximal $\goth{n}$ de~$B$, un generador $u$ de~$B$ sobre~$A$ y un polinomio mónico $F\in A[t]$ tales que $\goth{n}\not\ni F'(u)$ y que $\cal{O}$ sea isomorfa, como $A$-álgebra, a~$B_{\goth{n}}$. Si $\cal{O}$ es étale sobre~$A$, puede tomarse $B=A[t]/FA[t]$.
\end{theoreme}

(Por supuesto, estas condiciones son asimismo suficientes\ldots)

Señalemos primero los agradables corolarios siguientes:

\begin{corollaire}
\label{I.7.7}
Para que $\cal{O}$ sea neta sobre~$A$ es necesario y suficiente que $\cal{O}$ sea isomorfa al cociente de un álgebra de forma análoga y \emph{étale} sobre~$A$.
\end{corollaire}

En efecto, se tomará $\cal{O}'=B'_{\goth{n}'}$, donde $B'=A[t]/FA[t]$ y $\goth{n}'$ es la imagen inversa de~$\goth{n}$ en~$B'$.

\marginpar{12}
\begin{corollaire}
\label{I.7.8}
Sea
$f\colon X\to Y$ un morfismo de tipo finito y sea $x\in X$. Para que $f$ sea neto en~$x$ es necesario y suficiente que exista un entorno abierto~$U$ de~$x$ tal que $f|U$ se factorice como $U\to X'\to Y$, donde la primera flecha es una inmersión cerrada y la segunda un morfismo étale.
\end{corollaire}

Esta es una simple traducción de~\Ref{I.7.7}.

Veamos cómo la formulación de~\Ref{I.7.6} se sigue del enunciado principal: en efecto, por~\Ref{I.7.7} existe un epimorfismo $\cal{O}'\to\cal{O}$, donde $\cal{O}'$ tiene las propiedades requeridas;% SGA1-CORR-0005: la autoridad imprime aquí O sin prima.
si además $\cal{O}$ es étale sobre~$A$, tanto $\cal{O}'$ como~$\cal{O}$ son étales sobre~$A$, y el morfismo $\cal{O}'\to\cal{O}$ es étale por~\Ref{I.4.8}, luego es un isomorfismo.

\subsubsection*{Demostración de~\Ref{I.7.6}}
Retoma una demostración del seminario Chevalley. Por el \emph{Main Theorem} se tendrá $\cal{O}=B_{\goth{n}}$, donde $B$ es un álgebra finita sobre~$A$ y $\goth{n}$ uno de sus ideales maximales. Entonces $B/\goth{n}=K(\cal{O})$ es una extensión separable y, por tanto, monógena de~$k$; si $\goth{n}_i$ ($1\leq i\leq r$) son los ideales maximales de~$B$ distintos de~$\goth{n}$, existe, pues, un elemento $u$ de~$B$ que pertenece a todos los~$\goth{n}_i$ y cuya imagen en~$B/\goth{n}$ es un generador. Ahora bien,
$B/\goth{n}=B_\goth{n}/\goth{n}B_\goth{n}=B_\goth{n}/\goth{m}B_\goth{n}$
(donde $\goth{m}$ es el ideal maximal de~$A$). Admitamos por un momento el lema siguiente.

\begin{lemme}
\label{I.7.9}
Sean $A$ un anillo local, $B$ un álgebra finita sobre~$A$, $\goth{n}$ un ideal maximal de~$B$ y $u$ un elemento de~$B$ cuya imagen en $B_\goth{n}/\goth{m}B_\goth{n}$ lo genera como álgebra sobre~$k=A/\goth{m}$ y que pertenece a todos los ideales maximales de~$B$ distintos de~$\goth{n}$. Sean
$B'=A[u]$ y $\goth{n}'=\goth{n}\cap B'$.% SGA1-CORR-0006: la autoridad imprime B'=B[u], n'=nB'.
Entonces el homomorfismo canónico $B'_{\goth{n}'}\to B_\goth{n}$ es un isomorfismo.
\end{lemme}

\begin{lemme}
\label{I.7.10}
(Habría debido figurar como corolario de~\Ref{I.7.1}, antes de~\Ref{I.7.5}, al que implica.) Sea $B$ un álgebra finita sobre~$A$ generada por un elemento~$u$, y sea $\goth{n}$ un ideal maximal de~$B$ tal que $B_\goth{n}$ sea no ramificado sobre~$A$. Entonces existe un polinomio mónico $F\in A[t]$ tal que $F(u)=0$ y $F'(u)\notin\goth{n}$.
\end{lemme}

Sea, en efecto, $n$ el rango de la $k$-álgebra $L=B\otimes_Ak$; por Nakayama existe un polinomio mónico de grado~$n$ en~$A[t]$ tal que $F(u)=0$. Sea $f$ el polinomio deducido de~$F$ por reducción módulo~$\goth{m}$; entonces $L$ es $k$-isomorfa a $k[t]/fk[t]$, y por~\Ref{I.7.3}, $f'(\xi)$ no pertenece al ideal maximal de~$L$ que corresponde a~$\goth{n}$ (donde $\xi$ designa la imagen de~$t$ en~$L$, \ie la imagen de~$u$ en~$L$). Como $f'(\xi)$ es la imagen de~$F'(u)$, hemos terminado.

\marginpar{13}
El teorema~\Ref{I.7.6} se sigue ahora de la conjunción de~\Ref{I.7.9} y~\Ref{I.7.10}. Queda demostrar~\Ref{I.7.9}. Pongamos $S'=B'-\goth{n}'$, de modo que $B'{S'}^{-1}=B'_{\goth{n}'}$.
$$
\begin{array}{ccccc}
B&\longrightarrow&B{S'}^{-1}&\longrightarrow&BS^{-1}=B_\goth{n}\\ \Big\uparrow&&\Big\uparrow&&\\
B'&\longrightarrow&B'{S'}^{-1}
\rlap{\hbox{$\displaystyle\ =B'_{\goth{n}'}$}}&&
\end{array}
$$
Sea análogamente $S=B-\goth{n}$, de modo que $BS^{-1}=B_\goth{n}$. Se tiene, pues, un homomorfismo natural $B{S'}^{-1}\to BS^{-1}=B_\goth{n}$. Demostremos que es un isomorfismo, \ie que los elementos de~$S$ son invertibles en $B{S'}^{-1}$, \ie que todo ideal maximal $\goth{p}$ de este último no interseca a~$S$, \ie induce~$\goth{n}$ sobre~$B$. En efecto, como $B{S'}^{-1}$ es finito sobre~$B'{S'}^{-1}=B'_{\goth{n}'}$, $\goth{p}$ induce el único ideal maximal $\goth{n}'B'_{\goth{n}'}$ de~$B'_{\goth{n}'}$,% SGA1-CORR-0007: la autoridad omite la prima en B'_{n'}.
y, por tanto, induce el ideal maximal $\goth{n}'$ de~$B'$; como $B$ es finito sobre~$B'$, el ideal $\goth{q}$ de~$B$ inducido por~$\goth{p}$, al estar situado sobre~$\goth{n}'$, es necesariamente maximal y no contiene a~$u$, por lo que es idéntico a~$\goth{n}$. (Acabamos de utilizar que $u$ pertenece a todo ideal maximal de~$B$ distinto de~$\goth{n}$.) Demostremos ahora que $B{S'}^{-1}=B'{S'}^{-1}$: como el primero es finito sobre el segundo, Nakayama reduce a demostrar la igualdad módulo $\goth{n}'B{S'}^{-1}$ y, \emph{a fortiori}, basta demostrar la igualdad módulo $\goth{m}B{S'}^{-1}$. Ahora bien,
$B{S'}^{-1}/\goth{m}B{S'}^{-1}=B_\goth{n}/\goth{m}B_\goth{n}$
está generado sobre~$k$ por~$u$ (aquí se utiliza la otra propiedad de~$u$); por tanto, la imagen de~$B'$ (y, \emph{a fortiori}, la de $B'{S'}^{-1}$) en él es todo el anillo, pues es un subanillo que contiene a~$k$ y a la imagen de~$u$.

\begin{remarquestar}
Debería poder enunciarse el teorema~\Ref{I.7.6} para un anillo $\cal{O}$ que solo sea semilocal, de modo que abarque también~\Ref{I.7.5}: se formulará la hipótesis de que $\cal{O}/\goth{m}\cal{O}$ es una $k$-álgebra \emph{monógena}; podrá encontrarse entonces un $u\in B$ cuya imagen en \ifthenelse{\boolean{orig}}{$B/mB$}{$B/\goth{m}B$}\ % La autoridad corrige m por \goth m.
sea un generador y que pertenezca a todos los ideales maximales de~$B$ que no provengan de~$\cal{O}$. Los lemas~\Ref{I.7.9} y~\Ref{I.7.10} deben adaptarse sin dificultad. Más generalmente, \ldots
\end{remarquestar}

\section[Levantamiento infinitesimal de los esquemas étales]{Levantamiento infinitesimal de los esquemas étales. Aplicación a los esquemas formales}
\label{I.8}

\begin{proposition}
\label{I.8.1}
Sean $Y$ un preesquema, $Y_0$ un subpreesquema, $X_0$ un $Y_0$-esquema étale y $x$ un punto de~$X_0$. Entonces existen un $Y$-esquema étale~$X$, un entorno~$U_0$ de~$x$ en~$X_0$ y un $Y_0$-isomorfismo $U_0\isomto X\times_Y Y_0$.
\end{proposition}

Sea, en efecto, $y$ la proyección de~$x$ en~$Y_0$. Aplicando~\Ref{I.7.6} al homomorfismo local étale $A_0\to B_0$ de los anillos locales de~$y$ y~$x$ en $Y_0$ y~$X_0$, se obtiene un isomorfismo
\marginpar{14}%
$$
B_0={C_0}_{\goth{n}_0}\qquad C_0=A_0[t]/F_0A_0[t],
$$
donde $F_0$ es un polinomio mónico y $\goth{n}_0$ un ideal maximal de~$C_0$ que no contiene la clase de $F'_0(t)$ en~$C_0$. Sea $A$ el anillo local de~$y$ en~$Y$, sea $F$ un polinomio mónico de~$A[t]$ que da~$F_0$ mediante el homomorfismo sobreyectivo $A\to A_0$ (se levantan los coeficientes de~$F_0$), y sea, por último,
\ifthenelse{\boolean{orig}}{$\cal{C}=A[t]/FA[t]$}{$C=A[t]/FA[t]$}\ % La autoridad corrige \cal C por C.
y $\goth{n}$ el ideal maximal \ifthenelse{\boolean{orig}}{de~$\cal{C}$}{de~$C$}\ % La autoridad corrige \cal C por C.
que es imagen inversa de \ifthenelse{\boolean{orig}}{$\goth{n}_0x$}{$\goth{n}_0$}\ % La autoridad elimina la x espuria.
por el epimorfismo natural $C\to C\otimes_AA_0=C_0$. Pongamos
$$
\ifthenelse{\boolean{orig}}{B=C_\goth{n}}{B=C_\goth{n}.}% La autoridad añade el punto final.
$$
Por la construcción y~\Ref{I.7.1}, es inmediato que $B$ es étale sobre~$A$ y que se tiene un isomorfismo $B\otimes_AA_0\simeq B_0$.% SGA1-CORR-0008: la autoridad imprime A_0.
Se sabe (cap\ptbl\ I) que existen un $Y$-esquema de tipo finito~$X$ y un punto~$z$ de~$X$ sobre~$y$ tales que $\cal{O}_z$ es $A$-isomorfo a~$B$;% SGA1-CORR-0008: la autoridad imprime C, pero solo B=C_n es étale.
como este último es étale sobre~$A=\cal{O}_y$, puede suponerse (tomando $X$ suficientemente pequeño) que $X$ es étale sobre~$Y$. Sea $X'_0=X\times_YY_0$; entonces el anillo local de~$z$ en~$X'_0$ se identifica con $\cal{O}_z\otimes_AA_0=B\otimes_AA_0$ y es, por tanto, isomorfo a~$B_0$. Este isomorfismo está definido por un isomorfismo de un entorno~$U_0$ de~$x$ en~$X_0$ sobre un entorno de~$z$ en~$X'_0$ (\emph{loc. cit.}), que puede suponerse igual a~$X'_0$ tomando $X$ suficientemente pequeño. Hemos terminado.

\begin{corollaire}
\label{I.8.2}
Enunciado análogo para \emph{recubrimientos} étales, suponiendo infinito el cuerpo residual $\kres(y)$.
\end{corollaire}

La demostración es la misma, sustituyendo~\Ref{I.7.6} por~\Ref{I.7.5}.

\begin{theoreme}
\label{I.8.3}
El funtor considerado en~\Ref{I.5.5} es una \emph{equivalencia de categorías}.
\end{theoreme}

En virtud del teorema~\Ref{I.5.5}, queda demostrar que todo $S_0$-esquema étale~$X_0$ es isomorfo a un $S_0$-esquema $X\times_SS_0$, donde $X$ es un $S$-esquema étale. El espacio topológico subyacente a~$X$ deberá ser necesariamente idéntico al de~$X_0$, identificándose además $X_0$ con un subpreesquema cerrado de~$X$. El problema equivale, pues, al siguiente: encontrar sobre el espacio topológico subyacente $|X_0|$ de~$X_0$ un haz de álgebras~$\cal{O}_X$ sobre~$f_0^*(\cal{O}_S)$ (donde $f_0$ es la proyección $X_0\to S_0$, considerada aquí como aplicación continua de los espacios subyacentes) que haga de~$|X_0|$ un $S$-preesquema étale~$X$, y un homomorfismo de álgebras $\cal{O}_X\to\cal{O}_{X_0}$ compatible, en los haces de escalares, con el homomorfismo $f_0^*(\cal{O}_S)\to f_0^*(\cal{O}_{S_0})$, y que induzca un isomorfismo
$\cal{O}_X\otimes_{f_0^*(\cal{O}_S)}f_0^*(\cal{O}_{S_0})\isomto\cal{O}_{X_0}$.
(Entonces $X$ será un $S$-preesquema étale cuya restricción es~$X_0$,
\marginpar{15}
y será, por tanto, separado sobre~$S$, puesto que $X_0$ lo es sobre~$S_0$; así, $X$ resuelve el problema.) Si $(U_i)$ es un recubrimiento de~$X_0$ por abiertos y se ha encontrado una solución del problema en cada~$U_i$, del teorema de unicidad~\Ref{I.5.5} se sigue que estas soluciones se pegan (\ie se pegan los haces de álgebras que las definen, provistos de sus homomorfismos de aumentación), y se comprueba de inmediato que el espacio anillado así construido sobre~$S$ es un $S$-preesquema étale~$X$ provisto de un isomorfismo $X\times_SS_0\isomfrom X_0$. Basta, pues, encontrar una solución localmente, lo cual está garantizado por~\Ref{I.8.1}.

\begin{corollaire}
\label{I.8.4}
Sea $S$ un preesquema formal localmente noetheriano, provisto de un ideal de definición~$J$, y sea $S_0=\bigl(|S|,\cal{O}_S/\cal{J}\bigr)$ el preesquema ordinario correspondiente. Entonces el funtor $\goth{X}\mto\goth{X}\times_SS_0$, de la categoría de los recubrimientos étales de~$S$ en la categoría de los recubrimientos étales de~$S_0$, es una equivalencia de categorías.
\end{corollaire}

Por supuesto, se llamará recubrimiento étale de un preesquema \emph{formal}~$S$ a un recubrimiento de~$S$, \ie un preesquema formal sobre~$S$ definido mediante un haz coherente de álgebras~$\cal{B}$, tal que $\cal{B}$ sea \emph{localmente libre} y las fibras residuales $\cal{B}_s\otimes_{\cal{O}_s}\kres(s)$ de~$\cal{B}$ sean álgebras \emph{separables} sobre~$\kres(s)$. Si se denota por~$S_n$ el preesquema ordinario
\ifthenelse{\boolean{orig}}{$\bigl(|S|\ \cal{O}_S/J^{n+1}\bigr)$,}
{$\bigl(|S|,\cal{O}_S/J^{n+1}\bigr)$,}% La autoridad corrige el espacio por una coma.
dar un haz coherente de álgebras $\cal{B}$ sobre~$S$ equivale a dar una sucesión de haces coherentes de álgebras
\ifthenelse{\boolean{orig}}{$B_n$}{$\cal{B}_n$}
sobre
\ifthenelse{\boolean{orig}}{los~$X_n$,}{los~$S_n$,}
provistos de un sistema transitivo de homomorfismos
\ifthenelse{\boolean{orig}}{$B_m\to B_n$}{$\cal{B}_m\to\cal{B}_n$}
($m\geq n$) que definen isomorfismos
\ifthenelse{\boolean{orig}}
{$B_m\otimes_{\cal{O}_{S_m}}\cal{O}_{S_n}\isomto B_n$.}
{$\cal{B}_m\otimes_{\cal{O}_{S_m}}\cal{O}_{S_n}\isomto\cal{B}_n$.}
Es inmediato que $\cal{B}$ es localmente libre si y solo si lo son los $\cal{B}_n$ sobre los~$S_n$, y que la condición de separabilidad se verifica si y solo si se verifica para~$\cal{B}_0$, o también para todos
\ifthenelse{\boolean{orig}}{los~$B_n$.}{los~$\cal{B}_n$.}
Así, $\cal{B}$ es étale sobre~$S$ si y solo si lo son los
\ifthenelse{\boolean{orig}}{$B_n$}{$\cal{B}_n$}
sobre los~$S_n$. Teniendo esto en cuenta,~\Ref{I.8.4} se sigue inmediatamente de~\Ref{I.8.3}.

\begin{remarquestar}
En~\Ref{I.8.4} no era necesario limitarse al caso de los \emph{recubrimientos}. Sin embargo, es el único que se utiliza por el momento.
\end{remarquestar}

\section{Propiedades de permanencia}
\label{I.9}
Sea $A\to B$ un homomorfismo local y étale. Examinamos aquí algunos casos en los que cierta propiedad de~$A$ implica la misma propiedad para~$B$, o recíprocamente.
\marginpar{16}
Varias proposiciones de esta índole son ya consecuencias del mero hecho de que $B$ es \emph{cuasifinito} y \emph{plano} sobre~$A$, y nos limitaremos a «recordar» algunas de ellas: $A$ \emph{y} $B$ \emph{tienen la misma dimensión de Krull y la misma profundidad} («codimensión cohomológica» de Serre, según la terminología aún corriente). Se sigue, por ejemplo, que $A$ \emph{es Cohen--Macaulay si y solo si lo es $B$}. Además, para todo ideal primo $\goth{q}$ de~$B$ que induzca $\goth{p}$ sobre~$A$, $B_{\goth{q}}$ seguirá siendo cuasifinito y plano sobre~$A_{\goth{p}}$, siempre que se suponga que $B$ es la localización de un álgebra de tipo finito sobre~$A$ (esto se sigue de que el conjunto de los puntos en los que un morfismo de tipo finito es cuasifinito, \resp plano, es abierto); por otra parte, \emph{todo} ideal primo $\goth{p}$ de~$A$ es inducido por un ideal primo $\goth{q}$ de~$B$ (pues $B$ es \emph{fielmente} plano sobre~$A$). Se sigue, por ejemplo, que \emph{$\goth{p}$ y~$\goth{q}$ tienen la misma altura}; y también que \emph{$A$ no tiene ideales primos embebidos si y solo si $B$ tampoco los tiene}.

Nos limitaremos, pues, a las proposiciones más específicas del caso de los morfismos étales.

\begin{proposition}
\label{I.9.1}
Sea $A\to B$ un homomorfismo local étale. Para que $A$ sea regular es necesario y suficiente que $B$ lo sea.
\end{proposition}
En efecto, sea $k$ el cuerpo residual de~$A$ y $L$ el de~$B$. Como $B$ es plano sobre~$A$ y
\ifthenelse{\boolean{orig}}
{$L=B\otimes_Ak$}
{$L=B\otimes_Ak$,}
\ie $\goth{n}=\goth{m}B$ (donde $\goth{m}$ y~$\goth{n}$ son los ideales maximales de~$A$ y~$B$), la filtración $\goth{m}$-ádica de~$B$ es idéntica a su filtración $\goth{n}$-ádica, y se tendrá
$$
\gr^*(B)=\gr^*(A)\otimes_kL.
$$
Se sigue que $\gr^*(B)$ es un álgebra de polinomios sobre~$L$ si y solo si $\gr^*(A)$ es un álgebra de polinomios sobre~$k$, c.q.d. (N.B.: no se ha utilizado que $L/k$ sea separable.)

\begin{corollaire}
\label{I.9.2}
Sea $f\colon X\to Y$ un morfismo étale. Si $Y$ es regular, $X$ también lo es; el recíproco es cierto si $f$ es \ifthenelse{\boolean{orig}}{sobreyectivo.}{sobreyectivo.}% La autoridad corrige el acuerdo francés.
\end{corollaire}

\setcounter{subsection}{1}
\setcounter{proposition}{1}

\begin{proposition}
\label{prop:I.9.2}
Sea $f\colon X\to Y$ un morfismo étale. Si $Y$ es reducido, $X$ también lo es; el recíproco es cierto si $f$ es sobreyectivo.
\end{proposition}

Esto equivale al corolario siguiente.
\begin{corollaire}
\label{I.9.3}
Sea $f\colon A\to B$ un homomorfismo local étale, siendo $B$ isomorfa a la localización de un $A$-álgebra de tipo finito.
\marginpar{17}
Para que $A$ sea reducido es necesario y suficiente que $B$ lo sea.
\end{corollaire}

La necesidad es trivial, puesto que $A\to B$ es inyectivo (al ser $B$ fielmente plano sobre~$A$). Para la suficiencia, sean
\ifthenelse{\boolean{orig}}
{$p_i$ los ideales primos minimales de~$A$; por hipótesis, la aplicación natural $A\to\prod_iA/p_i$ es inyectiva, de modo que, tensorizando con el $A$-módulo plano $B$, se obtiene que $B\to\prod_iB/p_iB$ es inyectiva, y se reduce a demostrar que los $B/p_iB$ son reducidos. Como $B/p_iB$ es étale sobre~$A/p_i$}
{$\goth{p}_i$ los ideales primos minimales de~$A$; por hipótesis, la aplicación natural $A\to\prod_iA/\goth{p}_i$ es inyectiva, de modo que, tensorizando con el $A$-módulo plano $B$, se obtiene que $B\to\prod_iB/\goth{p}_iB$ es inyectiva, y se reduce a demostrar que los $B/\goth{p}_iB$ son reducidos. Como $B/\goth{p}_iB$ es étale sobre~$A/\goth{p}_i$},
se reduce al caso en que $A$ es íntegro. Sea~$K$ su cuerpo de fracciones; como $A\to K$ es inyectivo, también lo es (por ser $B$ plano sobre~$A$) el homomorfismo $B\to B\otimes_AK$, y se reduce a demostrar que este último anillo es reducido. Ahora bien, como $B$ es la localización de un $A$-álgebra de tipo finito, es el anillo local de un punto~$x$ de un esquema de tipo finito y \emph{étale} $X=\Spec(C)$ sobre~$Y=\Spec(A)$; por tanto, $B\otimes_AK$ es una localización (respecto de un conjunto multiplicativamente estable adecuado) del anillo $C\otimes_AK$ de~$X\otimes_AK$. Como $X\otimes_AK$ es étale sobre~$K$, su anillo es un producto finito de cuerpos (extensiones separables de~$K$); por tanto, lo mismo ocurre con~$B\otimes_AK$, c.q.d.

\begin{corollaire}
\label{I.9.4}
Sea $f\colon A\to B$ un homomorfismo local étale; supongamos que $A$ es analíticamente reducido
\index{analíticamente reducido|hyperpage}%
(\ie que el completado $\hat{A}$ de~$A$ carece de elementos nilpotentes). Entonces $B$ es analíticamente reducido y, \emph{a fortiori}, reducido.
\end{corollaire}

En efecto, \ifthenelse{\boolean{orig}}{$\hat{B}$}{$\hat{B}$}\ % La autoridad corrige B por \hat B.
es \emph{finito} y étale sobre~$\hat{A}$, y se aplica~\Ref{I.9.3}.

\begin{theoreme}
\label{I.9.5}
Sea $f\colon A\to B$ un homomorfismo local, siendo $B$ isomorfa a la localización de un $A$-álgebra de tipo finito. Entonces:
\begin{enumerate}
\item[(i)] Si $f$ es étale, $A$ es normal si y solo si $B$ lo es.
\item[(ii)] Si $A$ es normal, $f$ es étale si y solo si $f$ es inyectivo y neto (y entonces $B$ es normal por~\textup{(i)}).
\end{enumerate}
\end{theoreme}

Daremos dos demostraciones distintas de~(i). La primera utiliza algunas propiedades de los morfismos planos cuasifinitos (recordadas al comienzo de este número), sin utilizar~\Ref{I.7.6} (y, por tanto, el \emph{Main Theorem}); la segunda demostración hace lo contrario. Por último, para~(ii) parece que se necesita en todo caso el \emph{Main Theorem}.

\subsubsection*{Primera demostración}
Se utiliza la siguiente condición necesaria y suficiente
\marginpar{18}
para la normalidad de un anillo local noetheriano $A$ de dimensión distinta de cero.

\begin{enonce*}{Criterio de Serre}
\index{Serre (criterio de)|hyperpage}%
\textup{(i)} Para todo ideal primo $\goth{p}$ de~$A$ de altura~$1$, $A_{\goth{p}}$ es normal (o, lo que es equivalente, regular); \textup{(ii)} para todo ideal primo $\goth{p}$ de~$A$ de altura~$\geq2$, se tiene $\operatorname{profundidad}A_{\goth{p}}\geq2$.\footnote{\Cf EGA~IV~5.8.6.}
\end{enonce*}

Admitiremos aquí este criterio, que se supone debe figurar en el \ifthenelse{\boolean{orig}}{párr.}{parágrafo} de los morfismos planos. Su principal ventaja es que no supone \emph{a priori} que $A$ sea reducido ni, \emph{a fortiori}, íntegro. Aquí puede suponerse ya que $\dim A=\dim B\neq0$.

Según los recordatorios del comienzo de este número, los ideales primos $\goth{p}$ de~$A$ que tienen altura~$1$ (\resp altura~$\geq2$) son exactamente las contracciones a~$A$ de los ideales primos~$\goth{q}$ de~$B$ que tienen altura~$1$ (\resp altura~$\geq2$). Por último, si $\goth{p}$ y $\goth{q}$ se corresponden, $B_{\goth{q}}$ es étale sobre~$A_{\goth{p}}$ y, por tanto, tiene la misma profundidad que $A_{\goth{p}}$, y es regular si y solo si $A_{\goth{p}}$ lo es (~\Ref{I.9.1}). Aplicando el criterio de Serre, se obtiene que $A$ es normal si y solo si $B$ lo es.

\subsubsection*{Segunda demostración}
Supongamos que $B$ es normal, y sean $L$ su cuerpo de fracciones y $K$ el de~$A$ ($A$ es íntegro puesto que $B$ lo es). En la demostración de~\Ref{I.9.3} se vio que $B\otimes_AK$ es un producto finito de cuerpos; como está contenido en~$L$, es un cuerpo, y como contiene a~$B$, es~$L$. Todo elemento de~$K$ entero sobre~$A$ es entero sobre~$B$ y, por tanto, pertenece a~$B$, pues \ifthenelse{\boolean{orig}}{$B$}{$B$}\ % La autoridad corrige b por B.
es normal; pertenece entonces a~$A$, ya que $B\cap K=A$ (como se sigue de que $B$ es fielmente plano sobre~$A$).

Supongamos ahora que $A$ es normal y demostremos que $B$ lo es. En virtud de~\Ref{I.7.6}, se tendrá $B=B'_{\goth{n}}$, donde $B'=A[t]/FA[t]$, siendo $F$ y $\goth{n}$ como en~\Ref{I.7.6}. Entonces $L=B\otimes_AK$ será una localización de~$B'\otimes_AK=K[t]/FK[t]$, y un producto de cuerpos que son extensiones finitas separables de~$K$; este producto (como ocurre siempre que se localiza un anillo artiniano, aquí $B'_K$, respecto de un conjunto multiplicativamente estable) es un factor directo de~$B'_K$. Corresponde, por tanto, a una descomposición $F=F_1F_2$ en~$K[t]$, y el generador de~$L$ correspondiente a~$t$ está ya anulado por~$F_1$. Ahora bien, \emph{como $A$ es normal, los $F_i$ pertenecen a~$A[t]$} (suponiéndolos mónicos). Observando que $B\to L=B\otimes_AK$ es inyectivo (pues $A\to K$ lo es y $B$ es plano sobre~$A$), se sigue que ya se tendrá
\ifthenelse{\boolean{orig}}{$F_1(u)=0$.}{$F_1(u)=0$, donde $u$ es la clase de~$t$ en~$L$.}
Suponiendo que $F$ se haya elegido de grado mínimo, se seguirá que $F_2=1$ (N.B.: se tendrá $F'(u)=F'_1(u)F_2(u)+F_1(u)F'_2(u)=F'_1(u)F_2(u)$, puesto que $F_1(u)=0$; de aquí $F'_1(u)\neq0$, puesto que \ifthenelse{\boolean{orig}}{$F'(u)\neq0$.}{$F'(u)\neq0$.)}% La autoridad añade el paréntesis de cierre.

Por
\marginpar{19}
tanto, se tiene
\begin{equation*}
\label{eq:I.9.5.*}
\tag{$*$}L=B\otimes_AK=K[t]/FK[t]
\end{equation*}
y $F$ es, en consecuencia, un polinomio separable de~$K[t]$ (pero, evidentemente, no es necesariamente separable en~$A[t]$). (N.B.: por el momento solo se ha demostrado, en esencia, que en~\Ref{I.7.6} pueden elegirse $F$ y \ifthenelse{\boolean{orig}}{\ }{$\goth{n}$}\ % La autoridad repone \goth n.
de modo que ---con las notaciones adoptadas aquí--- \ifthenelse{\boolean{orig}}
{$B'\to B'_{\goth{n}}=\cal{B}$}
{$B'\to B'_{\goth{n}}=B$}% La autoridad corrige \cal B por B.
sea \emph{inyectivo}; para ello se ha utilizado la normalidad de~$A$; no sé si esto sigue siendo cierto sin la hipótesis de normalidad.)

Recordemos ahora el conocido lema, extraído del curso de Serre del año pasado:

\begin{lemme}
\label{I.9.6}
Sean $K$ un anillo, $F\in K[t]$ un polinomio mónico separable, $L=K[t]/FK[t]$ y $u$ la clase de~$t$ en~$L$ (de modo que $F'(u)$ es un elemento invertible de~$L$). Entonces se tienen las fórmulas siguientes (donde $n=\deg F$):
\begin{enumerate}
\item[] $\trace_{L/K}u^i/F'(u)=0$ si \ifthenelse{\boolean{orig}}{$0\leq i<N-1$,}{$0\leq i<n-1$,}% La autoridad corrige N por n.
\item[] $\trace_{L/K}u^{n-1}/F'(u)=1$.
\end{enumerate}
\end{lemme}

\begin{corollaire}
\label{I.9.7}
El determinante de la matriz $(\trace_{L/K}(u^j\cdot u^i/F'(u)))_{0\leq i,j\leq n-1}$ es igual a
\ifthenelse{\boolean{orig}}{$(-1)^n$}{$(-1)^{n(n-1)/2}$},
por lo que es invertible en todo subanillo $A$ de~$K$.
\end{corollaire}

\begin{corollaire}
\label{I.9.8}
\ifthenelse{\boolean{orig}}{Sean}{Sean}\ % La autoridad corrige el acuerdo francés.
$A$ un subanillo de~$K$, $V$ el $A$-módulo engendrado por los $u^i$ ($0\leq i\leq n-1$) en~$L$, y $V'$ el sub-$A$-módulo de~$L$ formado por los $x\in L$ tales que $\trace_{L/K}(xy)\in A$ para todo $y\in V$ (\ie para $y$ de la forma $u^i$, $0\leq i\leq n-1$). Entonces $V'$ es el $A$-módulo que tiene por base los $u^i/F'(u)$ ($0\leq i\leq n-1$).
\end{corollaire}

\begin{corollaire}
\label{I.9.9}
Supongamos que $K$ es el cuerpo de fracciones de un anillo íntegro normal~$A$ y que $F$ tiene sus coeficientes en~$A$. Entonces, con las notaciones de~\Ref{I.9.8}, $V'$ contiene la clausura normal $A'$ de~$A$ en~$L$, que está, por tanto, contenida en $A[u]/F'(u)$ y, \emph{a fortiori}, en $A[u][F'(u)^{-1}]$.
\end{corollaire}

Apliquemos este último corolario a la situación obtenida en la demostración: como $F'(u)$ es invertible en~$B$, que contiene a~$A[u]$, $B$ contiene a~$A'$. Por el \emph{Main Theorem} (o a partir del hecho de que $B=A[u]_{\goth{n}}$), $B$ es una localización de~$A'$. Como $A'$ es normal, también lo es~$B$.

\subsubsection*{Demostración de \textup{(ii)}}
Se
\marginpar{20}
procede como en la demostración precedente para demostrar que, en~\Ref{I.7.6}, puede elegirse $F$ de modo que siga cumpliéndose~($*$). El único obstáculo \emph{a priori} es que, al no suponerse ya que $B$ es plano sobre~$A$, no puede afirmarse que $B\to L$ sea inyectivo, por lo que el razonamiento solo se aplicará \emph{a priori} a la imagen~$B_1$ de~$B$ por dicho homomorfismo. Se sigue inmediatamente que $B_1$ es plano sobre~$A$ (por ser una localización de un álgebra libre sobre~$A$). En virtud de~\Ref{I.4.8}, el morfismo $B\to B_1$ es étale y, por tanto, un isomorfismo, lo que concluye la demostración.

(Desde el punto de vista de la redacción, habría que invertir el orden de las dos últimas demostraciones y poner en un número aparte los cálculos formales del lema y de sus corolarios.)

\begin{corollaire}
\label{I.9.10}
Sea $f\colon X\to Y$ un morfismo étale. Si $Y$ es normal, $X$ lo es; el recíproco es cierto si $f$ es \ifthenelse{\boolean{orig}}{sobreyectivo.}{sobreyectivo.}% La autoridad corrige el acuerdo francés.
\end{corollaire}

\begin{corollaire}
\label{I.9.11}
Sea $f\colon X\to Y$ un morfismo dominante, siendo $Y$ normal y $X$ conexo. Si $f$ es neto, $f$ es étale; por tanto, $X$ es normal y, en consecuencia (por ser conexo), irreducible.
\end{corollaire}

Sea $U$ el conjunto de los puntos en los que $f$ es étale; es abierto, y basta demostrar que también es cerrado y no vacío. $U$ contiene la imagen inversa del punto genérico de~$Y$ (pues, para un álgebra sobre un cuerpo, no ramificado $=$ étale) y, por tanto (como $X$ domina a~$Y$), no es vacío. Si $x$ pertenece a la adherencia de~$U$, pertenece a la adherencia de una componente irreducible $U_i$ de~$U$ y, por tanto, a una componente irreducible $X_i\overset{\mathrm{def}}{=}\overline{U_i}$ de~$X$ que corta a~$U$; en consecuencia, $X_i$ domina a~$Y$ (pues toda componente de~$U$, plana sobre~$Y$, domina a~$Y$). Por consiguiente, si $y$ es la proyección de~$x$ sobre~$Y$, $\cal{O}_y\to\cal{O}_x$ es \emph{inyectivo} (teniendo en cuenta que $\cal{O}_y$ es íntegro). Como $\cal{O}_y$ es normal y $\cal{O}_y\to\cal{O}_x$ es neto, se concluye mediante~\Ref{I.9.5}(ii).

\begin{corollaire}
\label{I.9.12}
Sea $f\colon X\to Y$ un morfismo de tipo finito dominante, con $Y$ normal y $X$ irreducible. Entonces el conjunto de los puntos en los que $f$ es étale es el complemento del soporte de $\Omega^1_{X/Y}$, \ie el complemento del subpreesquema de~$X$ definido por el ideal diferente $\goth{d}_{X/Y}$.
\end{corollaire}

(Este
\marginpar{21}
es el enunciado «menos trivial» al que se aludía en la observación del \No\Ref{I.4}.)

\begin{remarquestar}
No debe creerse que un recubrimiento étale conexo de un esquema irreducible sea él mismo irreducible cuando no se supone normal la base. Esta cuestión se estudiará en el \No\Ref{I.11}.
\end{remarquestar}

\section{Recubrimientos étales de un esquema normal}
\label{I.10}

\begin{proposition}
\label{I.10.1}
Sea $X$ un preesquema étale separado sobre~$Y$, normal y conexo, de cuerpo de funciones~$K$. Entonces las componentes conexas $X_i$ de~$X$ son íntegras, sus cuerpos de funciones $K_i$ son extensiones finitas separables de~$K$, y $X_i$ se identifica con una parte abierta no vacía de la normalización de
\ifthenelse{\boolean{orig}}{$X$}{$Y$}\ % SGA1-CORR-0009: la base de normalización es Y.
en~$K_i$ (por tanto, $X$ se identifica con una parte abierta densa de la normalización de~$Y$ en $R(X)=L=\prod K_i$).
\end{proposition}

Según~\Ref{I.9.10}, $X$ es normal; \,\emph{a fortiori}, sus anillos locales son íntegros, de modo que las componentes conexas de~$X$ son irreducibles. Como $X_i$ es normal,
\ifthenelse{\boolean{orig}}{finito}{cuasifinito}\ % SGA1-CORR-0009: la hipótesis étale solo da cuasifinitud.
y dominante sobre~$Y$, se sigue de un caso particular (por lo demás casi trivial) del \emph{Main Theorem} que $X_i$ es un abierto de la normalización de
\ifthenelse{\boolean{orig}}{$X$}{$Y$}\ % SGA1-CORR-0009.
en el cuerpo $K_i$ de
\ifthenelse{\boolean{orig}}{$X$}{$X_i$}.% SGA1-CORR-0009.

\begin{corollaire}
\label{I.10.2}
En las condiciones de~\Ref{I.10.1}, $X$ es finito sobre~$Y$ (\ie un recubrimiento étale de~$Y$) si y solo si $X$ es isomorfo a la normalización $Y'$ de~$Y$ en $L=R(X)$ (anillo de las funciones racionales sobre~$X$).
\end{corollaire}

En efecto, se sabe que esta normalización es finita sobre~$Y$ (por ser $Y$ normal y
\ifthenelse{\boolean{orig}}{$R/K$}{$L/K$}\ % SGA1-CORR-0010: L=R(X); R no está definido.
separable); recíprocamente, si $X$ es finito sobre~$Y$, lo es sobre~$Y'$, por lo que su imagen en~$Y'$ es cerrada y, por otra parte, densa.

Se dirá que un álgebra $L$ de rango finito sobre~$K$ es \emph{no ramificada sobre
\ifthenelse{\boolean{orig}}{$X$}{$Y$}}\ % SGA1-CORR-0011: la base definida es Y.
(o simplemente no ramificada sobre~$K$, si
\ifthenelse{\boolean{orig}}{$X$}{$Y$}\ % SGA1-CORR-0011.
se sobreentiende) si $L$ es un álgebra separable sobre~$K$, \ie un producto directo de extensiones separables $K_i$, y si la normalización $Y'$ de~$Y$ en~$L$ (suma disjunta de las normalizaciones de~$Y$ en los $K_i$) es no ramificada ($=$ étale por~\Ref{I.9.11}) sobre~$Y$. Por tanto:

\begin{corollaire}
\label{I.10.3}
Para todo $X$ finito sobre~$Y$ tal que cada una de sus componentes irreducibles domina a~$Y$, sea $R(X)$ el anillo de las funciones racionales sobre~$X$ (producto de los anillos
\marginpar{22}
locales de los puntos genéricos de las componentes irreducibles de~$X$), de modo que $X\mto R(X)$ es un funtor con valores en las álgebras de rango finito sobre~$K=R(Y)$. Este funtor establece una equivalencia entre la categoría de los recubrimientos étales conexos de~$Y$ y la categoría de las extensiones $L$ de~$K$ no ramificadas sobre~$Y$.
\end{corollaire}

El funtor inverso es el funtor de normalización.

Supongamos que $Y$ es afín, definido, por tanto, por un anillo normal~$A$ de cuerpo de fracciones~$K$. Sea $L$ una extensión finita de~$K$ que es producto directo de cuerpos. Por definición, la normalización $Y'$ de~$Y$ en~$L$ es isomorfa a~$\Spec(A')$, donde $A'$ es la normalización de~$A$ en~$L$. Decir que $L$ es no ramificada sobre~$Y$ significa que $A'$ es no ramificada (o, equivalentemente, étale) sobre~$A$. Si $A$ es local, esto equivale a decir que los anillos locales $A'_{\goth{n}}$ (donde $\goth{n}$ recorre el conjunto finito de los ideales maximales de~$A'$, \ie de sus ideales primos que inducen el ideal maximal $\goth{m}$ de~$A$) son no ramificados ($=$ étales) sobre
\ifthenelse{\boolean{orig}}{$X$ }{}\ % Corrección ya incorporada en la autoridad: se suprime X.
el anillo local~$A$. Observemos, por último, que el criterio del discriminante~(\Ref{I.4.10}) también puede aplicarse en esta situación (más generalmente, una variante de dicho criterio debería enunciarse así, sin condición preliminar de planitud cuando $X$ domina a~$Y$, aunque se suponga que $Y$ es localmente íntegro: $A\to B$ y $B\to B\otimes_AK=L$ son inyectivos ---entonces está definida $\trace_{L/K}$--- y $\trace_{L/K}(xy)$ induce una \emph{forma bilineal fundamental} $B\times B\to A$, \ie existen $x_i\in B$ ($1\leq i\leq n$, donde $n$ es el rango de~$L$ sobre~$K$) tales que $\trace(x_ix_j)\in A$ para todo $i,j$ y
\ifthenelse{\boolean{orig}}{la matriz $(\trace(x_ix_j))$}{$\det(\trace(x_ix_j))$}\ % La autoridad corrige la pérdida de det.
es invertible en~$A$).

El sorites~(4.6) implica de inmediato el sorites de la no ramificación en el marco clásico:

\begin{proposition}
\label{I.10.4}
Sea $Y$ un preesquema normal íntegro, de cuerpo de funciones~$K$.
\textup{(i)} $K$ es no ramificado sobre~$Y$.
\textup{(ii)} Si $L$ es una extensión de~$K$ no ramificada sobre~$Y$, si $Y'$ es un preesquema normal de cuerpo de funciones~$L$ que domina a~$Y$ (por ejemplo, la normalización de~$Y$ en~$L$), y si $M$ es una extensión de~$L$ no ramificada sobre~$Y'$, entonces $M/K$ es no ramificada sobre
\ifthenelse{\boolean{orig}}{$X$}{$Y$}\ % SGA1-CORR-0012: X no se introduce; la base transitiva es Y.
(\emph{transitividad} de la no ramificación).
\textup{(iii)} Sea $Y'$ un preesquema normal íntegro que domina a~$Y$, de cuerpo de funciones $K'/K$; si $L$ es una extensión de~$K$ no ramificada sobre~$Y$, entonces $L\otimes_KK'$ es una extensión de~$K'$ no ramificada sobre~$Y'$ (propiedad de \emph{traslación}).
\end{proposition}

Además:
\marginpar{23}

\begin{corollaire}
\label{I.10.5}
En las condiciones de~\textup{(iii)}, si $Y=\Spec(A)$ y $Y'=\Spec(A')$, entonces la normalización $\bar{A'}$ de~$Y'$ en $L'=L\otimes_KK'$ se identifica con $\bar{A}\otimes_AA'$, donde $\bar{A}$ es la normalización de~$A$ en~$L$.
\end{corollaire}

Habitualmente, quienes rehúyen considerar anillos no íntegros, aunque sean productos directos de cuerpos, enuncian la propiedad de traslación en la forma siguiente, más débil:

\begin{corollaire}
\label{I.10.6}
En las condiciones de~\textup{(iii)}, sea $L_1$ una \emph{extensión compuesta} de~$L/K$ (no ramificada sobre~$Y$) y de~$K'/K$. Entonces $L_1/K'$ es no ramificada sobre~$Y'$. En el caso en que $Y=\Spec(A)$ y $Y'=\Spec(A')$, se tendrá además
\[
\overline{A'}=A[\overline{A},A'].
\]
Es decir, el anillo $\overline{A'}$, normalización de~$A'$ en~$L_1$, es el $A$-álgebra engendrada por~$A'$ y por la normalización $\overline{A}$ de~$A$ en~$L$.
\end{corollaire}

Este último hecho es, por lo demás, falso sin la hipótesis de no ramificación, incluso para extensiones compuestas de cuerpos de números...

Para terminar este número, daremos la interpretación de la noción de recubrimiento étale que corresponde a su imagen intuitiva: debe haber el «número máximo» de puntos por encima del punto considerado $y\in Y$ y, en particular, no debe haber «varios puntos confundidos» por encima de~$y$. Para demostrar los resultados en este sentido con toda la generalidad deseable, admitiremos aquí la proposición~\Ref{I.10.7} que figura más abajo (cuya demostración aparecerá en el multiplodoque, cap.~IV, párr.~15, y utiliza la técnica de los conjuntos constructibles de Chevalley y un poco de teoría del descenso...).

Un morfismo de tipo finito $f\colon X\to Y$ se dice \emph{universalmente abierto} si, para todo cambio de base $Y'\to Y$ (con $Y'$ localmente noetheriano), el morfismo $f'\colon X'=X\times_YY'\to Y'$ es abierto, \ie transforma abiertos en abiertos. Puede uno limitarse, por lo demás, al caso en que $Y'$ es de tipo finito sobre~$Y$ (e incluso al caso en que $Y'$ es de la forma $Y[t_1,\dots,t_r]$, donde los $t_i$ son indeterminadas). Un morfismo universalmente abierto es, \,\emph{a fortiori}, abierto (el recíproco es falso); por otra parte, si $f$ es abierto, siendo $X$ e~$Y$ irreducibles, todas las componentes de todas las fibras de~$f$ tienen la misma dimensión (a saber, la dimensión de la fibra genérica
\marginpar{24}
$f^{-1}(z)$, donde $z$ es el punto genérico de~$Y$). Finalmente, si $Y$ es normal, esta última condición implica ya que $f$ es \emph{universalmente} abierto (teorema de Chevalley). Se sigue, por ejemplo, que si $f\colon X\to Y$ es un morfismo \emph{cuasifinito}, con $Y$ normal e irreducible, entonces $f$ es universalmente abierto (o, equivalentemente, abierto) si y solo si toda componente irreducible de~$X$ domina a~$Y$. Recordemos también que, como un morfismo plano de tipo finito es abierto, también es universalmente abierto. Sentados estos preliminares, «recordemos» la siguiente

\begin{proposition}
\label{I.10.7}
Sea $f\colon X\to Y$ un morfismo cuasifinito, separado y universalmente abierto. Para todo $y\in Y$, sea $n(y)$ el «número geométrico de puntos de la fibra $f^{-1}(y)$», igual a la suma de los grados separables de las extensiones residuales $\kres(x)/\kres(y)$, para los puntos $x\in f^{-1}(y)$. Entonces la función $y\mto n(y)$ sobre~$Y$ es semicontinua superiormente. Para que sea constante en un entorno del punto~$y$ (\ie para que se tenga $n(y)=n(z_i)$, donde los $z_i$ son los puntos genéricos de las componentes irreducibles de~$Y$ que contienen a~$y$), es necesario que exista un entorno $U$ de~$y$ tal que $X|U$ sea \emph{finito} sobre~$U$.\footnote{\Cf EGA~IV~15.5.1.}
\end{proposition}

\begin{corollaire}
\label{I.10.8}
Si $y\mto n(y)$ es constante e~$Y$ es geométricamente unirramoso,\footnote{Para la definición, \cf más abajo, n.º~\Ref{I.11}, p.~\pageref{I.11}.} las componentes irreducibles de~$X$ son disjuntas.
\end{corollaire}

\begin{proposition}
\label{I.10.9}
Sea $f\colon X\to Y$ un morfismo \emph{étale} separado. Con las notaciones de~\Ref{I.10.7}, la función $y\mto n(y)$ es semicontinua superiormente. Para que sea constante en un entorno del punto~$y$ (\ie para que se tenga $n(y)=n(z_i)$, donde los $z_i$ son los puntos genéricos de las componentes irreducibles de~$Y$ que contienen a~$y$), es necesario y suficiente que exista un entorno abierto $U$ de~$y$ tal que $X|U$ sea finito sobre~$U$, \ie que sea un \emph{recubrimiento étale} de~$U$.
\end{proposition}

\begin{corollaire}
\label{I.10.10}
Para que un morfismo étale separado $f\colon X\to Y$, con $Y$ conexo, sea finito (\ie haga de~$X$ un \emph{recubrimiento étale} de~$Y$), es necesario y suficiente que todas las fibras de~$f$ tengan el mismo número geométrico de puntos.
\end{corollaire}

En~\ifthenelse{\boolean{orig}}{\Ref{I.10.8}}{\Ref{I.10.7}}\ % La autoridad corrige la referencia 10.8 por 10.7.
y su corolario no había hipótesis de normalidad sobre~$Y$. Si se añade tal hipótesis, se obtiene el enunciado más fuerte siguiente (tomado con frecuencia como definición de la no ramificación de un recubrimiento):

\begin{theoreme}
\label{I.10.11}
Sea
\marginpar{25}
$f\colon X\to Y$ un morfismo \emph{cuasifinito} separado. Se supone que~$Y$ es \emph{irreducible}, que toda componente de~$X$ domina a~$Y$ y que $X$ es reducido (\ie que $\cal{O}_X$ no tiene elementos nilpotentes). Sea $n$ el grado de~$X$ sobre~$Y$ (suma de los grados, sobre el cuerpo~$K$ de~$Y$, de los cuerpos $K_i$ de las componentes irreducibles $X_i$ de~$X$). Sea $y$ un punto normal de~$Y$. Entonces el número geométrico $n(y)$ de puntos de~$X$ por encima de~$y$ es $\leq n$, y la igualdad se da si y solo si existe un entorno abierto $U$ de~$y$ tal que $X|U$ sea un \emph{recubrimiento étale} de~$U$.
\end{theoreme}

\ifthenelse{\boolean{orig}}
{Siendo trivial el «solo si», demostremos el «si».}
{La existencia de tal entorno étale implica trivialmente que $n(y)=n$; demostremos el recíproco.}% SGA1-CORR-0013: la autoridad intercambia los nombres de las dos implicaciones.
Sea $z$ el punto genérico de~$Y$; se tiene $n(z)=$ (suma de los grados separables de los $K_i/K$) $\leq n$ y, por~\Ref{I.10.7}, se tiene $n(y)\leq n(z)\leq n$; la igualdad implica que $X|U$ es \emph{finito} sobre~$U$ para un entorno adecuado~$U$ de~$y$. Puede suponerse, por tanto, que $X$ es finito sobre~$Y$ y que la función $n(y')$ sobre~$Y$ es constante. Por último, según~\Ref{I.10.8}, $X$ es entonces la unión disjunta de sus componentes irreducibles y, para demostrar que es no ramificado en~$y$, se reduce uno al caso en que $X$ es irreducible y, por tanto, íntegro. Finalmente, puede suponerse que $Y=\Spec(\cal{O}_y)$. El teorema se reduce entonces al enunciado clásico siguiente:

\begin{corollaire}
\label{I.10.12}
Sean $A$ un anillo local normal (noetheriano, como siempre) de cuerpo de fracciones~$K$; $L$ una extensión finita de~$K$ de grado~$n$ y grado separable~$n_s$; $B$ un subanillo de~$L$, finito sobre~$A$ y de cuerpo de fracciones~$L$; $\goth{m}$ el ideal maximal de~$A$; y $n'$ el grado separable de~$B/\goth{m}B$ sobre~$A/\goth{m}A=k$ ($=$ la suma de los grados separables de las extensiones residuales de este anillo). Se tiene $n'\leq n_s$ y, \,\emph{a fortiori}, $n'\leq n$. Esta última desigualdad es una igualdad si y solo si $B$ es no ramificado ($=$ étale) sobre~$A$.
\end{corollaire}

Solo queda demostrar que $n'=n$ implica que $B$ es étale sobre~$A$. Recordemos la demostración cuando $k$ es infinito: basta demostrar que $R=B/\goth{m}B$ es separable sobre~$k$; de no ser así, se seguiría (por un lema conocido) que existe un elemento $a$ de~$R$ cuyo polinomio minimal sobre~$k$ tiene grado $>n'$. Este elemento procede de un elemento $x$ de~$B$, cuyo polinomio minimal sobre~$K$ (como elemento de~$L$) tiene grado $\leq n$; por otra parte, este último tiene sus coeficientes en~$A$, pues $A$ es normal, y da, por reducción módulo~$\goth{m}$, un polinomio mónico $F\in k[t]$, de grado $\leq n=n'$, tal que $F(a)=0$, lo cual es absurdo.

En
\marginpar{26}
el caso general (pudiendo ser finito~$k$), retomemos el lenguaje geométrico y consideremos $Y'=\Spec(A[t])$, que es fielmente plano sobre~$Y$, y el punto genérico $y'$ de la fibra $\Spec(k[t])$ de~$Y'$ sobre~$y$. Entonces $X$ es neto sobre~$Y$ en~$y$ si y solo si $X'=X\times_YY'=\Spec(B[t])$ es neto en~$y'$ sobre~$Y'$, como se comprueba de inmediato. Por otra parte, por la elección de~$y'$, su cuerpo residual es $k(t)$ y, por tanto, infinito. Como $y'$ es un punto normal de~$Y'$, se reduce uno al caso precedente.


\section{Algunos complementos}
\label{I.11}

Ya hemos dicho que un recubrimiento étale conexo de un esquema íntegro no es necesariamente íntegro. He aquí dos ejemplos de este hecho.

\noindent
a) Sea $C$ una curva algebraica con un punto doble ordinario~$x$, sea $C'$ su normalización, y sean $a$ y~$b$ los dos puntos de~$C'$ situados por encima de~$x$. Sean $C'_i$ ($i=1,2$) dos copias de~$C'$, y sean $a_i$ y~$b_i$ los puntos de~$C'_i$ que corresponden a~$a$ y~$b$, respectivamente. En la curva suma $C'_1\amalg C'_2$, identifiquemos
\ifthenelse{\boolean{orig}}{$a_1^x$}{$a_1$}\ % La autoridad corrige a_1^x por a_1.
con~$b_2$, por una parte, y $a_2$ con~$b_1$, por otra (se deja al lector precisar el proceso de identificación; se explicará en el cap.~VI del multiplodoque, pero, para curvas sobre un cuerpo algebraicamente cerrado, se trata en el libro de Serre sobre curvas algebraicas). Se obtiene una curva $C''$ \emph{conexa} y \emph{reducible}, que es un recubrimiento étale de grado~$2$ de~$C$. El lector comprobará que, en general, los recubrimientos étales conexos «galoisianos» $C''$ de~$C$ cuya imagen inversa $C''\times_CC'$ es un recubrimiento \emph{trivial} de~$C'$ (\ie isomorfo a la suma disjunta de cierto número de copias de~$C'$) son «cíclicos» de grado~$n$, y que, para todo entero $n>0$, puede construirse un recubrimiento étale conexo cíclico de grado~$n$. En el lenguaje del grupo fundamental, que se desarrollará más adelante, esto significa que el cociente de~$\pi_1(C)$ por el subgrupo normal cerrado engendrado por la imagen de~$\pi_1(C')\to\pi_1(C)$ (homomorfismo inducido por la proyección) es isomorfo a la completación profinita de~$\ZZ$. Más precisamente, debería poder demostrarse que el grupo fundamental de~$C$ es isomorfo al producto libre topológico del grupo fundamental
\ifthenelse{\boolean{orig}}{de~$C$}{de~$C'$}\ % La autoridad corrige C por C'.
con la completación profinita de~$\ZZ$. Obsérvese que cuestiones de este tipo dieron origen a la «teoría del descenso» para los esquemas.

\noindent
b) Sea $A$ un anillo local completo íntegro. Se sabe que su normalización $A'$ es finita
\marginpar{27}
sobre~$A$ (Nagata); es, por tanto, un anillo semilocal completo y, por ser íntegro, es local. Supongamos que la extensión residual $L/k$ que define no es radicial (en el caso contrario se dirá que $A$ es geométricamente unirramoso; \cf más abajo). Así ocurre, por ejemplo, para el anillo
\ifthenelse{\boolean{orig}}
{$\RR[s,t]/(s^2+t^2)\RR[s,t]$}
{$\RR[[s,t]]/(s^2+t^2)\RR[[s,t]]$},% La autoridad repone el anillo local completo de series formales.
donde $\RR$ es el cuerpo de los números reales. Sea entonces $k'$ una extensión galoisiana finita de~$k$ tal que $L\otimes_kk'$ se descomponga, y sea~$B$ un álgebra finita étale sobre~$A$ correspondiente a la extensión residual~$k'$ (recordemos que $B$ es esencialmente única). Entonces $B'=A'\otimes_AB$, como $B$-álgebra, tiene el álgebra residual $L\otimes_kk'$, que no es local; por tanto, $B'$ no es un anillo local y, como es completo, tiene divisores de cero. Puesto que está contenido en el anillo total de fracciones de~$B$ (pues es libre sobre~$A'$, luego sin torsión sobre~$A'$, luego sin torsión sobre~$A$, y está, por consiguiente, contenido en $B'\otimes_AK=B'_{(K)}=A'_{(K)}\otimes_KB_{(K)}=B_{(K)}$, ya que $A'_{(K)}=K$), se sigue que $B$ no es íntegro. En el caso del anillo $\RR[[s,t]]/(s^2+t^2)\RR[[s,t]]$, tomando $k'/k=\CC/\RR$, se obtiene para~$B$ el anillo local de dos rectas secantes del plano en su punto de intersección.

Obsérvese, además, que, si existe un recubrimiento étale conexo $X$ de un esquema íntegro~$Y$ que no sea irreducible, entonces toda componente irreducible de~$X$ da un ejemplo de un recubrimiento no ramificado $X'$ de~$Y$, dominante sobre~$Y$, que no es étale sobre~$Y$. En el ejemplo~a) se obtiene así que $C'$ es no ramificado sobre~$C$ sin ser étale en los dos puntos $a$ y~$b$ (como también se comprueba directamente inspeccionando los completados de los anillos locales de~$x$ y~$a$: desde el punto de vista «formal», $C'$ en el punto~$a$ se identifica con un subesquema cerrado de~$C$ en el punto~$x$, a saber, una de las dos «ramas» de~$C$ que pasan por~$x$).

En a) y~b), se ve que la falta de validez de las conclusiones de~\Ref{I.9.5}\ptbl\textup{(i)} y~\textup{(ii)} está ligada directamente al hecho de que un punto de~$Y$ «se desdobla» en puntos \emph{distintos} de la normalización (\ifthenelse{\boolean{orig}}{en~b,}{en~b),}\ % La autoridad repone el paréntesis.
el hecho de que la extensión residual no sea radicial debe interpretarse geométricamente de este modo). Más precisamente, diremos que un anillo local íntegro~$A$ es \emph{geométricamente unirramoso}
\index{geométricamente unirramoso|hyperpage}%
\index{unirramoso (geométricamente)|hyperpage}%
si su normalización tiene un único ideal maximal y la extensión residual correspondiente es radicial; se dice que un punto~$y$ de un preesquema íntegro es geométricamente unirramoso si lo es su anillo local. Ejemplos: un punto normal, un punto cuspidal ordinario de una curva, etc. Parece que, si $Y$ admite un punto que no es unirramoso, existe siempre
\marginpar{28}
\ifthenelse{\boolean{orig}}{algún}{un}\ % La autoridad corrige el artículo francés.
recubrimiento étale conexo no irreducible de~$Y$; esto es, al menos, lo que hemos demostrado en el caso~b), cuando $Y$ es el espectro de un anillo local completo. En cambio, puede demostrarse que \emph{si todos los puntos de~$Y$ son geométricamente unirramosos, entonces todo $Y$-preesquema no ramificado, conexo y dominante sobre~$Y$ es étale} e irreducible. La demostración retoma la de~\Ref{I.9.5}, utilizando la generalización siguiente del teorema~\Ref{I.8.3}, que se demostrará más adelante mediante la técnica del descenso:\footnote{\Cf \Ref{IX.4.10}. Para una demostración más directa, \cf EGA~IV~18.10.3, donde se utiliza una variante de~\Ref{I.9.5} para anillos locales geométricamente unirramosos.}

\emph{Sea $Y'\to Y$ un morfismo finito, radicial y sobreyectivo (\ie lo que podría llamarse un «homeomorfismo universal»). Consideremos el funtor $X\mto X\times_YY'=X'$ de los $Y$-preesquemas en los $Y'$-preesquemas. Este funtor induce una equivalencia entre la categoría de los $Y$-esquemas étales y la categoría de los $Y'$-esquemas étales.} Este resultado podrá aplicarse, por ejemplo, en el caso en que $Y'$ es la normalización de~$Y$, suponiendo que $Y$ sea unirramoso (y que $Y'$ sea finito sobre~$Y$, como sucede en todos los casos que se encuentran en la práctica), o al caso de un $Y''$ «intercalado» entre~$Y$ y su normalización (que ya no necesita ser finita sobre~$Y$).

\chapter{Morfismos lisos: generalidades,~propiedades~diferenciales}
\label{II}
\marginpar{29}

Las referencias a la exposición~\Ref{I} se indican mediante~I. Recordemos que los anillos son noetherianos y los preesquemas localmente noetherianos.

\section{Generalidades}
\label{II.1}

Sea $Y$ un preesquema y sean $t_1,\dots,t_n$ indeterminadas. Se pone
\begin{equation}
\label{eq:II.1.1}
Y[t_1,\dots,t_n]=Y\otimes_\ZZ\ZZ[t_1,\dots,t_n]\quoi.
\end{equation}
Así, $Y[t_1,\dots,t_n]$ es un $Y$-esquema, afín sobre~$Y$, definido por el haz cuasicoherente de álgebras $\cal{O}_Y[t_1,\dots,t_n]$. Dar una sección de este preesquema sobre~$Y$ equivale, por tanto, a dar $n$ secciones de~$\cal{O}_Y$ (correspondientes a las imágenes de los~$t_i$ por el homomorfismo correspondiente). Si $Y'$ está sobre~$Y$, se tiene
\begin{equation}
\label{eq:II.1.2}
Y[t_1,\dots,t_n]\times_YY'=Y'[t_1,\dots,t_n]\quoi,
\end{equation}
(lo que implica que dar un $Y$-morfismo de~$Y'$ en $Y[t_1,\dots,t_n]$ equivale a dar $n$ secciones de~$\cal{O}_{Y'}$); por otra parte, se tiene
\begin{equation}
\label{eq:II.1.3}
\bigl(Y[t_1,\dots,t_n]\bigr)[t_{n+1},\dots,t_m]=Y[t_1,\dots,t_m]\quoi,
\end{equation}
en virtud de la fórmula análoga para los anillos de polinomios sobre~$\ZZ$. La fórmula~\eqref{eq:II.1.2} implica que $Y[t_1,\dots,t_n]$ varía funtorialmente con~$Y$.

$Y[t_1,\dots,t_n]$ es de tipo finito y plano sobre~$Y$.

\begin{definition}
\label{II.1.1}
Sea $f\colon X\to Y$ un morfismo que hace de~$X$ un $Y$-preesquema. Se dice que $f$ es \emph{liso}\footnote{Terminología antigua: $f$ es \emph{simple}
\index{simple: sinónimo anticuado de liso|hyperpage}%
en~$x$, o $x$ es un punto \emph{simple} para~$f$. Esta terminología se prestaba a confusión en diversos contextos (álgebras simples, grupos simples) y hubo que abandonarla.}
\index{liso (morfismo, álgebra)|hyperpage}%
en $x\in X$, o que $X$ es \emph{liso sobre~$Y$ en~$x$}, si existen un entero $n\geq0$, un entorno abierto $U$ de~$x$ y un $Y$-morfismo étale $U\to Y[t_1,\dots,t_n]$. Se dice que~$f$ (\resp $X$) es \emph{liso} si es liso en todos los puntos de~$X$. Se dice que un álgebra~$B$ sobre un anillo~$A$ es lisa en un ideal primo $\goth{p}$ de~$B$ si $\Spec(B)$ es liso sobre~$\Spec(A)$ en el punto~$\goth{p}$;
\marginpar{30}
se dice que $B$ es lisa sobre~$A$ si lo es en todo ideal primo~$\goth{p}$ de~$B$. Por último, se dice que un homomorfismo local $A\to B$ de anillos locales es liso (o que $B$ es lisa sobre~$A$)\footnote{Es preferible decir entonces, como en EGA~IV~18.6.1, que $B$ es «\emph{esencialmente lisa}» sobre~$A$.}
\index{esencialmente liso|hyperpage}%
si $B$ es la localización de un álgebra de tipo finito $B_1$ lisa sobre~$A$.
\end{definition}

Obsérvese que la noción de lisitud de~$X$ sobre~$Y$ es local sobre~$X$ y sobre~$Y$; si $X$ es liso sobre~$Y$, es localmente de tipo finito sobre~$Y$.

\setcounter{subsection}{0}
\setcounter{proposition}{0}
\begin{proposition}
\label{prop:II.1.1}
El conjunto de los puntos $x$ de~$X$ en los que $f$ es liso es abierto.
\end{proposition}

Es trivial a partir de la definición.

\begin{corollaire}
\label{II.1.2}
Si $B$ es lisa sobre~$A$ en~$\goth{p}$, entonces es lisa sobre~$A$ en~$\goth{q}$ para todo ideal primo~$\goth{q}$ de~$B$ contenido en~$\goth{p}$.
\end{corollaire}

\Ref{prop:II.1.1} implica también que las dos últimas definiciones de~\Ref{II.1.1} coinciden en su dominio común de existencia.

\begin{proposition}
\label{II.1.3}
\textup{(i)} Un morfismo étale ---en particular, una inmersión abierta o un morfismo identidad--- es liso. \textup{(ii)} Un cambio de base en un morfismo liso da un morfismo liso. \textup{(iii)} El compuesto de dos morfismos lisos es liso.
\end{proposition}

\textup{(i)} es trivial a partir de la definición; más precisamente:

\begin{corollaire}
\label{II.1.4}
étale $=$ cuasifinito $+$ liso.
\end{corollaire}

\textup{(ii)} se sigue inmediatamente del hecho análogo para los morfismos étales (\Ref{I.4.6}) y para las proyecciones $Y[t_1,\dots,t_n]\to Y$ (\cf \eqref{eq:II.1.2}). Para~\textup{(iii)}, se sigue formalmente de que esto es cierto por separado para «étale» (\Ref{I.4.6}) y para las proyecciones del tipo \ifthenelse{\boolean{orig}}{$Y[t_1,\dots,t_n]$}{$Y[t_1,\dots,t_n]\to Y$}\ % La autoridad repone el morfismo de proyección.
(\cf \eqref{eq:II.1.3}), junto con los dos hechos citados para~\textup{(ii)}. Supongamos que $Y$ es liso sobre~$Z$ y que $X$ es liso sobre~$Y$; demostremos que $X$ es liso sobre~$Z$. Puede suponerse que $Y$ es étale sobre~$Z[t_1,\dots,t_n]$ y que $X$ es étale sobre~$Y[s_1,\dots,s_m]$; la primera hipótesis implica entonces que $Y[s_1,\dots,s_m]$ es étale sobre
\[
Z[t_1,\dots,t_n][s_1,\dots,s_m]
=\ifthenelse{\boolean{orig}}{Z[t_1,\dots,s_m]}{Z[t_1,\dots,t_n,s_1,\dots,s_m]},
\]% SGA1-CORR-0014: se conservan los dos bloques completos de variables.
de modo que $X$ es étale sobre~\ifthenelse{\boolean{orig}}{$Z[t_1,\dots,s_m]$}{$Z[t_1,\dots,t_n,s_1,\dots,s_m]$},\ % SGA1-CORR-0014.
c.q.d.

\begin{remarque}
\label{II.1.5}
El entero $n$ que figura en la def.~\Ref{II.1.1} está bien determinado, pues se comprueba
\marginpar{31}
inmediatamente que es la dimensión del anillo local de~$x$ en su fibra $f^{-1}\bigl(f(x)\bigr)$. Se llama «dimensión relativa» de~$X$ sobre~$Y$. Se comporta aditivamente respecto de la composición de morfismos.
\end{remarque}

\section{Algunos criterios de lisitud de un morfismo}
\label{II.2}

\begin{theoreme}
\label{II.2.1}
Sea $f\colon X\to Y$ un morfismo localmente de tipo finito, sea $x\in X$ y sea $y=f(x)$. Para que $f$ sea liso en~$x$, es necesario y suficiente que \textup{(a)} $f$ sea plano en~$x$, y \textup{(b)} $f^{-1}(y)$ sea liso sobre~$\kres(y)$ en~$x$.
\end{theoreme}

Como el compuesto de dos morfismos planos es plano y $Y[t_1,\dots,t_n]\to Y$ es un morfismo plano, se ve que liso implica plano; teniendo en cuenta~\Ref{II.1.3}~\textup{(ii)}, esto demuestra la necesidad. Supongamos verificadas \textup{(a)} y~\textup{(b)}; sean $V$ un entorno afín de~$y$, de anillo~$A$, y $U$ un entorno afín de~$x$ situado sobre~$V$, de anillo~$B$. Tomando $U$ suficientemente pequeño, puede suponerse por~\textup{(b)} que existe un $\kres(y)$-morfismo \emph{étale}
\[
g\colon U|_{f^{-1}(y)}\longrightarrow\Spec k[t_1,\dots,t_n]\quad (k=\kres(y))
\]
definido por $n$ secciones $g_i$ del haz estructural de $U|_{f^{-1}(y)}$. Se comprueba fácilmente que puede suponerse que las~$g_i$ (que a priori son elementos de $B\otimes_Ak=BS^{-1}$, donde $S=A-\goth{p}$ y $\goth{p}$ es el ideal primo de~$A$ correspondiente a~$y$) provienen de secciones del haz estructural de
\ifthenelse{\boolean{orig}}{$V$}{$U$},\ % rama corregida ya incorporada por la autoridad: U, no V.
de modo que $g$ está inducido por un morfismo, denotado asimismo por~$g$,
\[
\ifthenelse{\boolean{orig}}{g\colon V\to Y[t_1,\dots,t_n]}{g\colon U\to Y[t_1,\dots,t_n]}% rama corregida ya incorporada por la autoridad.
\]
(tras multiplicar, si es preciso, las~$g_i$ por un mismo elemento no nulo de~$k$). Ahora bien,
\ifthenelse{\boolean{orig}}{$V$}{$U$}\ % rama corregida ya incorporada por la autoridad: U.
es plano sobre~$Y$ por~\textup{(a)}, y lo mismo ocurre con $Y[t_1,\dots,t_n]$; por otra parte, $g$ induce un morfismo étale entre las fibras sobre~$y$, luego $g$ es étale en~$x$ por~(\Ref{I.5.8}), c.q.d.

\begin{corollaire}
\label{II.2.2}
Sean $S$ un preesquema, $f\colon X\to Y$ un $S$-morfismo de tipo finito, siendo $Y$ de tipo finito y plano sobre~$S$, $x\in X$, y sea $s$ la proyección de~$x$ sobre~$S$. Para que $f$ sea liso en~$x$, es necesario y suficiente que $X$ sea plano\ifthenelse{\boolean{orig}}{ (o también: liso)}{}\ % SGA1-CORR-0017: la equivalencia impresa es falsa sin suponer Y liso sobre S.
sobre~$S$ en~$x$, y que el morfismo $f_s\colon X_s\to Y_s$ inducido entre las fibras sobre~$s$ sea liso en~$x$.
\end{corollaire}

Solo la suficiencia requiere demostración, y resulta del criterio~\Ref{II.2.1}, junto con el criterio de planitud~(\Ref{I.5.9}).

Para
\marginpar{32}
enunciar el resultado siguiente, «recordemos» que un morfismo $f\colon X\to Y$ localmente de tipo finito se dice \emph{equidimensional}
\index{equidimensional (morfismo)|hyperpage}%
en el punto $x\in X$ si (poniendo $y=f(x)$) puede encontrarse un entorno abierto~$U$ de~$x$, cada una de cuyas componentes domina una componente de~$Y$, tal que, para todo $y'\in Y$, las componentes irreducibles de $f^{-1}(y')\cap U$ tengan todas una misma dimensión, independiente de~$y'$. Basta, por lo demás, tomar en esta condición como~$y'$ los puntos genéricos de las componentes irreducibles de~$Y$ que pasan por~$y$, \emph{y} el propio punto~$y$. Si, por ejemplo, $X$ e~$Y$ son íntegros y $f$ es dominante, la condición significa que las componentes de $f^{-1}(y)$ que pasan por~$x$ tienen la dimensión «correcta», \ie la dimensión de la fibra genérica (recordemos que su dimensión es siempre $\geq$ la de la fibra genérica). Si $f$ es equidimensional en~$x$, la dimensión de su fibra en~$x$ es~$n$, y $g\colon U\to Y'=Y[t_1,\dots,t_n]$ es un $Y$-morfismo definido en un entorno~$U$ de~$x$ que induce sobre las fibras de~$y$ un morfismo cuasifinito en~$x$ (o, lo que equivale a lo mismo, si $g$ es cuasifinito en~$x$), entonces se demuestra que toda componente irreducible de~$U$ que pasa por~$x$ domina una componente irreducible de~$Y'$. Además, en virtud del «lema de normalización», siempre existe un~$g$ tal (y, recíprocamente, si existe un $Y$-morfismo cuasifinito~$g$ de un entorno abierto~$U$ de~$x$ a un $Y$-esquema de la forma $Y'=Y[t_1,\dots,t_n]$, tal que toda componente de~$U$ que pasa por~$x$ domina una componente de~$Y'$, entonces $f$ es equidimensional en~$x$). Sentado esto:

\begin{proposition}
\label{II.2.3}
Sean $f\colon X\to Y$ un morfismo localmente de tipo finito, $x$ un punto de~$X$ y $y=f(x)$; supóngase que $\cal{O}_y$ es normal. Para que $f$ sea liso en~$x$, es necesario y suficiente que $f$ sea equidimensional en~$x$ y que $f^{-1}(y)$ sea liso sobre~$\kres(y)$ en~$x$.
\end{proposition}

Se ve inmediatamente a partir de la definición que un morfismo liso es equidimensional (N.B.: un morfismo plano de tipo finito no es necesariamente equidimensional en~$x$, aunque su fibra en~$x$ sea irreducible). Demostremos el recíproco. Como $f^{-1}(y)$ es liso sobre~$\kres(y)$ en~$x$, puede suponerse (reemplazando si es preciso $X$ por un entorno adecuado de~$x$) que existe un $Y$-morfismo
\[
g\colon X\to Y[t_1,\dots,t_n]=Y'
\]
que induce un morfismo étale sobre las fibras de~$y$ y, a fortiori, es cuasifinito en~$x$. Por tanto,
\marginpar{33}
$g$ es no ramificado y (como $f$ es equidimensional en~$x$) las componentes irreducibles de~$X$ que pasan por~$x$ dominan cada una una componente de~$Y'$; a fortiori, el homomorfismo $\cal{O}_{y'}\to\cal{O}_x$ deducido de~$g$ (donde $y'=g(x)$) es \emph{inyectivo}. Este homomorfismo es además no ramificado, y $\cal{O}_{y'}$ es normal por ser una localización del anillo $\cal{O}_y[t_1,\dots,t_n]$, que es normal porque $\cal{O}_y$ lo es. Por consiguiente, el homomorfismo $\cal{O}_{y'}\to\cal{O}_x$ es étale (\Ref{I.9.5}~\textup{(ii)}).

\begin{remarques}
\label{II.2.4}
El enunciado anterior sigue siendo válido si se reemplaza la hipótesis de que $\cal{O}_y$ es normal por la hipótesis, más débil, de que $Y$ es \emph{geométricamente unirramoso} en~$y$ (\cf \Ref{I.11}), pues (\Ref{I.9.5}) vale bajo esta hipótesis. Aprovechemos la ocasión para señalar asimismo que, si el cuerpo residual de un anillo local íntegro~$A$ es algebraicamente cerrado, ser analíticamente irreducible
\index{analíticamente irreducible|hyperpage}%
(\ie que $\widehat{A}$ sea íntegro) implica ser geométricamente unirramoso; el recíproco es además cierto en toda categoría de «buenos anillos», más precisamente en una categoría de anillos estable bajo las operaciones usuales y en la que la completación de un anillo local normal sea normal (condición satisfecha, en virtud del «teorema de normalidad analítica» de Zariski, en la categoría de las álgebras afines y sus localizaciones)\footnote{\Cf EGA~IV~7.8.}.

«Recordemos», por último, en el contexto actual, el resultado siguiente, debido a Hironaka\footnote{\Cf EGA~IV~5.12.10.}, que permite a veces asegurarse de que $f^{-1}(y)$ sea un esquema reducido, \ie que sea también lo que numerosos geómetras algebraicos considerarían abusivamente como la fibra (sin multiplicidad) de~$f$ sobre~\ifthenelse{\boolean{orig}}{$x$}{$y$}\ % SGA1-CORR-0015: la fibra f^{-1}(y) está sobre y en Y.
(a saber, $f^{-1}(y)_{\red}$):
\end{remarques}

\begin{proposition}
\label{II.2.5}
Sean $f\colon X\to Y$ un morfismo dominante de tipo finito de preesquemas reducidos, y $y$ un punto de~$Y$ tal que $\cal{O}_y$ sea regular. Supóngase que todas las componentes de $f^{-1}(y)$ son de multiplicidad~$1$
\index{componente de multiplicidad~$1$|hyperpage}%
(\cf la definición más abajo), y que $f^{-1}(y)_{\red}$ es normal. Entonces $f^{-1}(y)$ es reducido y, por tanto, normal; $X$ es normal en todos los puntos de $f^{-1}(y)$; por último, $X$ es plano sobre~$Y$ en todos los puntos de $f^{-1}(y)$.
\end{proposition}
Se
\marginpar{34}
dice que una componente~$Z$ de $f^{-1}(y)$ es de \emph{multiplicidad}~$1$
\index{multiplicidad de una componente|hyperpage}%
si, denotando por~$x$ el punto genérico de~$Z$, se tiene \textup{(i)} $\dim\cal{O}_x=\dim\cal{O}_y$ (\ie que $Z$ no es una «componente excedentaria», es decir, no es «de dimensión demasiado grande»\ifthenelse{\boolean{orig}}{}{)};\ % SGA1-CORR-0016: se cierra la explicación parentética antes de (ii).
\textup{(ii)} el ideal maximal de~$\cal{O}_x$ está engendrado por el ideal maximal de~$\cal{O}_y$ (que a priori, en virtud de la elección de~$x$, engendra un ideal de definición de~$\cal{O}_x$).

Teniendo en cuenta~\Ref{II.2.3} o~\Ref{II.2.1}, se obtiene, pues:

\begin{corollaire}
\label{II.2.6}
Sean $f\colon X\to Y$ un morfismo dominante de tipo finito de preesquemas reducidos, y $y$ un punto de~$Y$ tal que $\cal{O}_y$ sea regular. Para que $f$ sea liso en los puntos de~$X$ situados sobre~$y$, es necesario y suficiente que las componentes de $f^{-1}(y)$ sean de multiplicidad~$1$ y que $f^{-1}(y)_{\red}$ sea liso sobre~$\kres(y)$.
\end{corollaire}

Esta situación se consideraba sobre todo antiguamente cuando $Y$ era el espectro de un anillo de valoración discreta~$A$, y se designaba comúnmente mediante expresiones tales como: «si la reducción de~$X$ respecto de la valoración dada es bonita»\dots. Además, $X$ designaba entonces un subesquema cerrado (si cabe decirlo) de algún $\PP_K^n$ (siendo $K$ el cuerpo de fracciones de~$A$) y, a falta de un
\ifthenelse{\boolean{orig}}{\emph{language}}{lenguaje}\ % corrección léxica ya incorporada por la autoridad: langage.
adecuado, apenas se percibía el papel más intrínseco de un objeto «definido sobre~$A$» (y no solamente sobre~$K$).

\section{Propiedades de permanencia}
\label{II.3}

\begin{proposition}
\label{II.3.1}
Sea $f\colon X\to Y$ un morfismo, sea $x\in X$ y sea $y=f(x)$. Supóngase que $f$ es liso en~$x$. Para que $\cal{O}_x$ sea reducido (\resp regular, \resp normal), es necesario y suficiente que $\cal{O}_y$ lo sea.
\end{proposition}

En efecto, este enunciado es conocido cuando $X$ es de la forma $Y[t_1,\dots,t_n]$, y se ha demostrado en~(I, \no \Ref{I.9}) para un morfismo étale; el caso general se deduce inmediatamente gracias a la definición~\Ref{II.1.1}.

No detallaremos aquí las demás propiedades de permanencia, que resultan ya de la sola planitud, o del hecho de que $X$ sea localmente cuasifinito y plano sobre un $Y$-preesquema de la forma $Y[t_1,\dots,t_n]$ (o, como
\marginpar{35}
diremos, de que $X$ sea
\ifthenelse{\boolean{orig}}{Cohen--Macauley}{Cohen--Macaulay}\ % corrección ortográfica ya incorporada por la autoridad.
\index{Cohen--Macaulay (esquema de)|hyperpage}%
sobre~$Y$). Señalemos solamente que de este último hecho se sigue que
\begin{equation}
\label{eq:II.3.1}
{\dim\cal{O}_x=\dim\cal{O}_y+n-d,\quad
\prof\cal{O}_x=\prof\cal{O}_y+n-d},
\end{equation}
donde $n$ es la dimensión de la fibra de~$f$ en~$x$, y $d$ el grado de trascendencia de~$\kres(x)$ sobre~$\kres(y)$; de donde (poniendo $\coprof=\dim-\prof$)\footnote{Para estas fórmulas, \cf EGA~IV~6.1 y~6.3.}
\begin{equation}
\label{eq:II.3.2}
{\coprof\cal{O}_x=\coprof\cal{O}_y}.
\end{equation}
De ello se sigue, por ejemplo, que $\cal{O}_x$ es Cohen--Macaulay (\resp no tiene componentes embebidas) si y solo si lo mismo ocurre con~$\cal{O}_y$.


\section{Propiedades diferenciales de los morfismos lisos}
\label{II.4}
Para simplificar, nos restringiremos esencialmente al cálculo diferencial de orden~$1$ y nos limitaremos a dar indicaciones rápidas para los órdenes superiores (donde los resultados son igual de sencillos).

Para la definición del haz de $1$-diferenciales $\it{\Omega}^1_{X/Y}$ de un $Y$-preesquema $X$, \cf (I~\No\Ref{I.1}). Supongamos que $X$ e $Y$ son $S$-preesquemas y que el morfismo estructural $f\colon X\to Y$ es un $S$-morfismo. Entonces $f$ define un homomorfismo de módulos (compatible con $f$)
\begin{equation}
\label{eq:II.4.1}
{f^*\colon \it{\Omega}^1_{Y/S}\to\it{\Omega}^1_{X/S}}
\end{equation}
es decir, $\it{\Omega}^1_{X/S}$ es \emph{contravariante} en el $S$-preesquema $X$. Por otra parte,~\eqref{eq:II.4.1} equivale a un homomorfismo de módulos sobre~$X$
\begin{equation*}
\label{eq:II.4.1bis}
\tag{4.1\textrm{bis}}
{f^*\big(\it{\Omega}^1_{Y/S}\big)\to\it{\Omega}^1_{X/S}}
\end{equation*}
que, a falta de una notación mejor, se denota también por $f^*$ y se inserta en una sucesión exacta canónica de homomorfismos de módulos
\begin{equation}
\label{eq:II.4.2}
{f^*\big(\it{\Omega}^1_{Y/S}\big)\to
\it{\Omega}^1_{X/S}\to\it{\Omega}^1_{X/Y}\to 0}.
\end{equation}
Todos estos homomorfismos se definen por la condición de ser de naturaleza local (lo que reduce al caso afín) y de conmutar con los operadores $\rd$. La exactitud de~\eqref{eq:II.4.2} es clásica e inmediata y, en el caso afín, se traduce en la sucesión exacta siguiente (correspondiente a un homomorfismo $B\to C$ de $A$-álgebras):
\begin{equation*}
\label{eq:II.4.2bis}
\tag{4.2\textrm{bis}}
{\Omega^1_{B/A}\otimes_B C\to\Omega^1_{C/A}\to\Omega^1_{C/B}\to 0}.
\end{equation*}

\begin{lemme}
\label{II.4.1}
Sea
\marginpar{36}
$f\colon X\to Y$ un morfismo de $S$-preesquemas. Si $f$ es no ramificado (\resp étale), entonces
$f^*\big(\it{\Omega}^1_{Y/S}\big)\to\it{\Omega}^1_{X/S}$
es sobreyectivo (\resp un isomorfismo). El recíproco es cierto en el caso «no ramificado» si se supone que $f$ es localmente de tipo finito.
\end{lemme}

El caso no ramificado resulta de la sucesión exacta~\eqref{eq:II.4.2} y de (\Ref{I.3.1}), pero el caso étale también puede verse directamente. Consideremos el diagrama
\begin{equation*}
\xymatrix@C=1.5cm{
X \ar[r]^-{\Delta_{X/Y}} & X\times_Y X \ar[d] \ar[r] &
X\times_S X \ar[d] \\ & Y \ar[r]^-{\Delta_{Y/S}} & Y\times_S Y
}
\end{equation*}
en el que $X\times_Y X$ se identifica con el producto fibrado de $Y$ y $X\times_S X$ sobre $Y\times_S Y$. Como $f$ es no ramificado, $X\to X\times_Y X$ es una inmersión abierta; por tanto, el haz conormal de la inmersión compuesta $\Delta_{X/S}$ de esta última con $X\times_Y X\to X\times_S X$ es isomorfo a la imagen inversa sobre~$X$ del haz conormal de la inmersión $X\times_Y X\to X\times_S X$. Por otra parte, en el caso étale, $X\to Y$ es plano, luego $X\times_S X\to Y\times_S Y$ es plano; por consiguiente, el haz conormal de la inmersión $X\times_Y X\to X\times_S X$ es isomorfo a la imagen inversa del haz conormal de la inmersión $Y\to Y\times_S Y$, es decir, a la imagen inversa de~$\it{\Omega}^1_{Y/S}$. De aquí se sigue la conclusión.

\begin{lemme}
\label{II.4.2}
Sea $X=Y[t_1,\dots,t_n]$, donde $Y$ es un $S$-preesquema. Entonces la sucesión de homomorfismos canónicos
\[
0\to f^*\big(\it{\Omega}^1_{Y/S}\big)\to
\it{\Omega}^1_{X/S}\to\it{\Omega}^1_{X/Y}\to 0
\]
es exacta y $\it{\Omega}^1_{X/Y}$ es libre con base $\rd_{X/Y}t_i$.
\end{lemme}

La verificación (puramente afín) es inmediata. (Nótese que ya conocemos la exactitud de~\eqref{eq:II.4.2}).

Combinando estos dos enunciados con la definición~\Ref{II.1.1}, se obtiene:

\begin{theoreme}
\label{II.4.3}
Sea $f\colon X\to Y$ un morfismo liso de $S$-preesquemas. Entonces:
\begin{enumerate}
\item[(i)] La sucesión de homomorfismos canónicos
\[
0\to f^*\big(\it{\Omega}^1_{Y/S}\big)\to
\it{\Omega}^1_{X/S}\to\it{\Omega}^1_{X/Y}\to 0
\]
es exacta.
\item[(ii)] $\it{\Omega}^1_{X/Y}$ es localmente libre y su rango $n$ en $x$ es igual a la dimensión relativa de~$f$ en~$x$.
\end{enumerate}
\end{theoreme}

\begin{corollaire}
\label{II.4.4}
El homomorfismo
\marginpar{37}
$f^*\big(\it{\Omega}^1_{Y/S}\big)\to\it{\Omega}^1_{X/S}$ es inyectivo, y su imagen en $\it{\Omega}^1_{X/S}$ es localmente un sumando directo.
\end{corollaire}

Sea $u\colon F\to G$ un homomorfismo de módulos sobre el preesquema $X$. Se dice que es \emph{universalmente inyectivo}
\index{universalmente inyectivo (morfismo de módulos)|hyperpage}%
en $x\in X$ si el homomorfismo $F_x\to G_x$ de $\cal{O}_x$-módulos es inyectivo y sigue siéndolo después de tensorizar con cualquier $\cal{O}_x$-álgebra (o, lo que por lo demás es equivalente, con cualquier $\cal{O}_x$-módulo). Basta, por ejemplo, con que exista un entorno abierto $U$ de~$x$ tal que $u$ induzca un isomorfismo de $F|U$ sobre un sumando directo de $G|U$; esta condición es también necesaria cuando $F$ y $G$ son libres (y de tipo finito) en un entorno de~$x$. Más precisamente, en este caso son equivalentes las condiciones siguientes:
\begin{enumerate}
\item[(i)] $u$ es inyectivo en $x$ y $\Coker u$ es libre en $x$;
\item[(ii)] existe un entorno abierto $U$ de~$x$ tal que $u$ induzca un isomorfismo de $F|U$ sobre un sumando directo de $G|U$;
\item[(iii)] $u$ es universalmente inyectivo en $x$;
\item[(iv)] el homomorfismo
\ifthenelse{\boolean{orig}}
{$\cal{F}_x\otimes\kres(x)\to\cal{G}_x\otimes\kres(x)$}
{$F_x\otimes\kres(x)\to G_x\otimes\kres(x)$}
sobre las «fibras restringidas» inducido por $u$ es inyectivo;
\item[(v)] el homomorfismo transpuesto $\check{G}\to\check{F}$ es sobreyectivo en el punto $x$ (o, lo que es equivalente, en un entorno de~$x$).
\end{enumerate}
(Demostración circular: (iv)$\To$(v) resulta de Nakayama; por otra parte, (v)$\To$(i),
\ifthenelse{\boolean{orig}}{puesto que un}{pues un}
haz cociente localmente libre es necesariamente un sumando directo). Geométricamente, la situación considerada significa que $u$ corresponde a un isomorfismo del fibrado vectorial cuyo haz de secciones
\ifthenelse{\boolean{orig}}{es~$\cal{F}$,}{es~$F$,}
sobre un subfibrado del fibrado vectorial análogo definido por $G$. Naturalmente, para ello no basta con que
\ifthenelse{\boolean{orig}}{$\cal{F}\to\cal{G}$}{$F\to G$}
sea inyectivo.

\begin{corollaire}
\label{II.4.5}
Sea $f\colon X\to Y$ un morfismo de $S$-preesquemas, localmente de tipo finito; sean $x\in X$, $y=f(x)$ y $s$ la proyección de $x$ e $y$ sobre~$S$. Supongamos que $Y$ es liso sobre~$S$ en $y$. Son equivalentes las condiciones:
\begin{enumerate}
\item[(i)] $f$ es liso en $x$.
\item[(ii)] $X$ es liso sobre
\ifthenelse{\boolean{orig}}{$Y$}{$S$}
en~$x$, y $f^*\big(\it{\Omega}^1_{Y/S}\big)\to\it{\Omega}^1_{X/S}$ es universalmente inyectivo en~$x$, es decir, es un homomorfismo inyectivo en~$x$ y su conúcleo $\it{\Omega}^1_{X/Y}$ es libre
\ifthenelse{\boolean{orig}}{en~$x$).}{en~$x$.}
\end{enumerate}
\end{corollaire}

La necesidad resulta de~\Ref{II.1.3}~(iii) y de~\Ref{II.4.3}~(i),\,(ii). Probemos la suficiencia. Como los $\rd g$ ($g\in\cal{O}_x$) engendran el módulo $\it{\Omega}^1_{X/Y}$ en $x$, se pueden encontrar
\marginpar{38}
$g_i$ ($1\le i\le n$) tales que las imágenes de los $\rd g_i$ en $\big(\it{\Omega}^1_{X/Y}\big)_x$ formen una base de este módulo. Tomando $X$ suficientemente pequeño, podemos suponer que los $g_i$ provienen de secciones de~$\cal{O}_X$ y definen, por tanto, un $Y$-morfismo $g\colon X\to Y'=Y[t_1,\dots,t_n]$. Utilizando la hipótesis y el lema~\Ref{II.4.2}, se ve fácilmente que el homomorfismo correspondiente $g^*\big(\it{\Omega}^1_{Y'/S}\big)\to\it{\Omega}^1_{X/S}$ es biyectivo en $x$, lo que nos reduce a demostrar el

\begin{corollaire}
\label{II.4.6}
Sea $f\colon X\to Y$ un morfismo de $S$-preesquemas lisos. Para que $f$ sea étale en $x\in X$, es necesario y suficiente que $f^*\big(\it{\Omega}^1_{Y/S}\big)\to\it{\Omega}^1_{X/S}$ sea un isomorfismo en $x$.
\end{corollaire}

Por~\Ref{II.4.1} sabemos que es necesario, y la misma condición implica, por ese lema, que $f$ es no ramificado en $x$. En virtud de~\Ref{II.2.2}, nos reducimos al caso $S=\Spec(k)$. Como $Y$ es liso sobre~$k$, es regular y, con mayor razón, normal; en virtud de (\Ref{I.9.5}~(ii)), nos reducimos a probar que $\cal{O}_y\to\cal{O}_x$ es inyectivo, o bien que $\cal{O}_y$ y $\cal{O}_x$ tienen la misma dimensión. Ahora bien, estas dimensiones son, respectivamente, los rangos de $\it{\Omega}^1_{Y/k}$ y $\it{\Omega}^1_{X/k}$ en $y$ \resp $x$, y por la hipótesis son iguales.

\begin{remarques}
\label{II.4.7}
Suponiendo que $X$ e $Y$ son lisos sobre~$S$, el criterio~\Ref{II.4.5}~(ii) para la lisura de $f\colon X\to Y$ puede enunciarse también diciendo que, para todo $x\in X$, la aplicación \emph{tangente} (relativa a la base $S$) de~$f$ en $x$, es decir, la transpuesta del homomorfismo de espacios vectoriales de dimensión finita sobre $\kres(x)$ que son las fibras restringidas en $x$ de $f^*\big(\it{\Omega}^1_{Y/S}\big)$ y $\it{\Omega}^1_{X/S}$, es \emph{sobreyectiva}. Esta es una hipótesis muy familiar, en particular entre quienes trabajan con espacios analíticos. La hipótesis de nonsingularidad que suelen imponer (que significa que $X$ e $Y$ son «lisos sobre~$\CC$», \cf \No\Ref{II.5}) parece deberse únicamente al temor que los puntos singulares de las variedades algebraicas o de los espacios analíticos todavía inspiran a muchos geómetras.
\end{remarques}

Señalemos el siguiente caso particular de~\Ref{II.4.6}:
\begin{corollaire}
\label{II.4.8}
Sean $X$ un $S$-preesquema, $g\colon X\to S[t_1,\dots,t_n]$ un $S$-morfismo definido por las secciones $g_i$ ($1\le i\le n$) de~$\cal{O}_X$, y $x$ un punto de~$X$ tal que $X$ sea liso sobre~$S$ en $x$. Para que $g$ sea étale en $x$, es necesario y suficiente que los $\rd g_i$
\marginpar{39}
($1\le i\le n$) formen una base de~$\it{\Omega}^1_{X/S}$ en $x$ (o, lo que es equivalente, que sus imágenes en
$\it{\Omega}^1_{X/S}(x)=\big(\it{\Omega}^1_{X/S}\big)_x\otimes_{\cal{O}_x}\kres(x)$ formen una base de este espacio vectorial sobre~$\kres(x)$).
\end{corollaire}

Sean $X$ un preesquema e $Y$ un subpreesquema cerrado de~$X$ definido por un haz coherente $\cal{J}$ de ideales. Entonces $\cal{J}/\cal{J}^2$ puede considerarse como un haz coherente sobre~$Y$ (el \emph{haz conormal}
\index{conormal (haz)|hyperpage}%
\index{haz conormal|hyperpage}%
de~$Y$ en~$X$). Si ahora $X$ es un $S$-preesquema, se tiene una sucesión exacta canónica de haces cuasicoherentes sobre~$Y$
\begin{equation}
\label{eq:II.4.3}
{\cal{J}/\cal{J}^2\lto{\rd}
\it{\Omega}^1_{X/S}\otimes_{\cal{O}_X}\cal{O}_Y\to
\it{\Omega}^1_{Y/S}\to 0}
\end{equation}
cuya parte derecha no es sino~\eqref{eq:II.4.2} (con los papeles de $X$ e $Y$ intercambiados y teniendo en cuenta que $\it{\Omega}^1_{Y/X}=0$), mientras que el homomorfismo $\cal{J}/\cal{J}^2\to\it{\Omega}^1_{X/S}\otimes_{\cal{O}_X}\cal{O}_Y$ se deduce del homomorfismo (en general no lineal) $g\mapsto\rd g$ pasando a los cocientes. La exactitud de~\eqref{eq:II.4.3} es clásica e inmediata y, en el caso afín, se interpreta mediante la sucesión exacta siguiente (correspondiente a un homomorfismo sobreyectivo $B\to C$ de $A$-álgebras, de núcleo $J$):
\begin{equation*}
\label{eq:II.4.3bis}
\tag{4.3\textrm{bis}}
{J/J^2\to\Omega^1_{B/A}\otimes_B C\to\Omega^1_{C/A}\to 0\qquad
(C=B/J)}
\end{equation*}
(sucesión exacta que ya se había utilizado implícitamente en la demostración de (\Ref{I.7.2})).


\begin{proposition}
\label{II.4.9}
Sean $X$ un $S$-preesquema, $Y$ un subpreesquema cerrado de~$X$ definido por un haz coherente $\cal{J}$ de ideales sobre~$X$, $x$ un punto de~$X$, $g_i$ ($1\leq i\leq n$) secciones de~$\cal{O}_X$ que definen un $S$-morfismo
\[
g\colon X\to S[t_1,\dots,t_n]=X',
\]
y, por último, $p$ un entero, $0\leq p\leq n$. Se supone que $X$ es \emph{liso sobre~$S$ en $x$}. Las condiciones siguientes son equivalentes:
\begin{enumerate}
\item[(i)] Existe un entorno abierto $X_1$ de~$x$ en~$X$ tal que $g|X_1$ es \emph{étale} y $Y_1=Y\cap X_1$ (la traza de~$Y$ sobre~$X_1$) es la \emph{imagen inversa} del subpreesquema cerrado $Y'=S[t_{p+1},\dots,t_n]$ de $X'=S[t_1,\dots,t_n]$ (es decir, los $g_i$ ($1\leq i\leq p$) engendran $\cal{J}|X_1$):
\[
\xymatrix{Y_1 \ar[d] \ar[r] & X_1 \ar[d]^{\hbox{étale}} \\
Y'=S[t_{p+1},\dots,t_n] \ar[r] & X'=S[t_1,\dots,t_n].}
\]
\item[(ii)]
$Y$
\marginpar{40}
es \emph{liso sobre~$S$ en $x$}; los $g_i$ ($1\leq i\leq p$) definen \emph{elementos de} $\cal{J}_x$; los $\rd g_i(x)$ ($1\leq i\leq n$) forman una \emph{base de} $\it{\Omega}^1_{X/S}(x)$ sobre~$\kres(x)$; y los $\rd g'_i(x)$ ($p+1\leq i\leq n$) forman una \emph{base de} $\it{\Omega}^1_{Y/S}(x)$ sobre~$\kres(x)$ (donde $g'_i$ denota la restricción de $g_i$ a~$Y$; las diferenciales se toman con respecto a~$S$).
\item[(iii)] Los $g_i$ ($1\leq i\leq p$) definen un \emph{sistema de generadores} de~$\cal{J}_x$, y los $\rd g_i(x)$ ($1\leq i\leq n$) forman una \emph{base de} $\it{\Omega}^1_{X/S}(x)$ sobre~$\kres(x)$.
\item[(iv)] $Y$ es \emph{liso sobre~$S$} en $x$; los $g_i$
\ifthenelse{\boolean{orig}}{\ignorespaces}{($1\leq i\leq p$)}
forman un \emph{sistema minimal de generadores de} $\cal{J}_x$; y los $\rd g'_i(x)$ ($p+1\leq i\leq n$) forman una \emph{base de} $\it{\Omega}^1_{Y/S}(x)$ sobre~$\kres(x)$.
\end{enumerate}

Además, bajo estas condiciones, $\cal{J}/\cal{J}^2$ es un módulo libre sobre~$Y$ en $x$, que admite como \emph{base en $x$} las clases de los $g_i$ ($1\leq i\leq p$), \emph{y el homomorfismo} canónico $\cal{J}/\cal{J}^2\to\it{\Omega}^1_{X/S}\otimes\cal{O}_Y$ es \emph{universalmente inyectivo en $x$}.
\end{proposition}

\begin{remarquestar}
Esto implica que $p$ queda bien determinado por las demás condiciones: como el \emph{rango} del módulo libre $\cal{J}/\cal{J}^2$ sobre~$Y$ en $x$; como el \emph{número mínimo de generadores} de~$\cal{J}_x$ sobre~$X$; o, finalmente, por el hecho de que la dimensión relativa de~$Y$ respecto de~$S$ en $x$ es $n-p$.
\end{remarquestar}

\subsubsection*{Demostración}
Supongamos primero que se cumple (i). Entonces, por (\Ref{I.4.6}~(iii)), $Y_1$ es étale sobre~$Y'$; por definición, es por tanto liso sobre~$S$ en $x$ (de dimensión relativa $n-p$), y lo mismo ocurre con~$Y$. De~\Ref{II.4.8} se sigue que los $\rd g_i$ ($1\leq i\leq n$) forman una base de~$\it{\Omega}^1_{X/S}$ en $x$, y que los $\rd g'_i$ ($p+1\leq i\leq n$) forman una base de~$\it{\Omega}^1_{Y/S}$ en $x$. La sucesión exacta~\eqref{eq:II.4.3} implica entonces que los $g_i$ ($1\leq i\leq p$) son linealmente independientes en $\cal{J}/\cal{J}^2$ (considerado como módulo sobre~$Y$) en $x$; como los $g_i$ ($1\leq i\leq p$) engendran $\cal{J}_x$, se sigue que los $g_i\bmod\cal{J}_x^2$ forman una \emph{base de} $\cal{J}/\cal{J}^2$ en $x$. Esto implica, por una parte, que los $g_i$ ($1\leq i\leq p$) forman un sistema \emph{minimal} de generadores de~$\cal{J}_x$ y, por otra, que el homomorfismo $\cal{J}/\cal{J}^2\to\it{\Omega}^1_{X/S}\otimes\cal{O}_Y$ de~\eqref{eq:II.4.3} es universalmente inyectivo en $x$ (pues lleva una base de un módulo libre en $x$ a una parte de una base de un módulo libre en $x$; nótese que se trata de $Y$-módulos). Esto prueba que (i) implica (ii), (iii), (iv), así como las últimas afirmaciones de la proposición~
\ifthenelse{\boolean{orig}}{\Ref{I.4.9}}{\Ref{II.4.9}}.

(iii) implica (i) en virtud del corolario~
\ifthenelse{\boolean{orig}}{\Ref{I.4.8}}{\Ref{II.4.8}}.

(ii)
\marginpar{41}
implica (i). En efecto, la primera hipótesis de (ii) significa que (reemplazando, si es preciso, $X$ por un entorno abierto de $x$ en $X$) $g$ induce un morfismo $h\colon Y\to Y'$. Según~\Ref{II.4.8}, las otras dos hipótesis de (ii) significan que $g$ es étale en $x$ y que $h$ es étale en $x$. Sea entonces $Y''$ la imagen inversa de~$Y'$ por $g$. Así, $Y$ es un subpreesquema cerrado de~$Y''$, el cual es étale sobre~$Y'$ en $x$ por (\Ref{I.4.6}~(iii)), pues $g$ es étale en $x$. Por tanto, el morfismo de inmersión $Y\to Y''$ es él mismo étale (\Ref{I.4.8}) y, en consecuencia, una inmersión abierta (\Ref{I.5.8} o \Ref{I.5.2}); reemplazando de nuevo $X$ por un entorno abierto conveniente $X_1$ de~$x$, se obtiene (i).

Lo anterior establece la equivalencia de las condiciones
\ifthenelse{\boolean{orig}}{(i) (ii) (iii)}{(i), (ii), (iii)},
y el hecho de que implican (iv). Falta probar que (iv)$\To$(ii), lo cual es inmediato (teniendo en cuenta que $\it{\Omega}^1_{X/S}$ es libre sobre~$X$ en $x$) una vez que se sabe que el hecho de que $Y$ sea liso sobre~$S$ en $x$ implica que $\cal{J}/\cal{J}^2$ es libre sobre~$Y$ en $x$ y que el homomorfismo $\cal{J}/\cal{J}^2\to\it{\Omega}^1_{X/S}\otimes\cal{O}_Y$ es universalmente inyectivo en $x$. Este último punto está incluido en el

\begin{theoreme}
\label{II.4.10}
Sean $X$ un $S$-preesquema liso, $Y$ un subpreesquema cerrado de~$X$ definido por un haz coherente $\cal{J}$ de ideales sobre~$X$, y $x$ un punto de~$X$. Las condiciones siguientes son equivalentes:
\begin{enumerate}
\item[(i)] $Y$ es \emph{liso sobre~$S$ en $x$}.
\item[(ii)] Existen un entorno abierto $X_1$ de~$x$ en~$X$ y un $S$-morfismo \emph{étale}
\[
g\colon X_1\to X'=S[t_1,\dots,t_n]
\]
tal que $Y_1=Y\cap X_1$ (la traza de~$Y$ sobre~$X_1$) sea el subpreesquema de~$X_1$ que es \emph{imagen inversa} por $g$ del subpreesquema cerrado $Y'=S[t_{p+1},\dots,t_n]$ de $X'=S[t_1,\dots,t_n]$, para un $p$ conveniente.
\item[(iii)] Existen \emph{generadores $g_i$ ($1\leq i\leq p$) de~$\cal{J}_x$} tales que los $\rd g_i$ formen parte de una base de~$\it{\Omega}^1_{X/S}$ en $x$ (o, lo que es equivalente, tales que los $\rd g_i(x)$ en $\it{\Omega}^1_{X/S}(x)$ sean linealmente independientes sobre~$\kres(x)$).
\item[(iv)] El haz $\cal{J}/\cal{J}^2$ es libre sobre~$Y$ en $x$, y el homomorfismo canónico
\[
\rd\colon\cal{J}/\cal{J}^2\to\it{\Omega}^1_{X/S}\otimes\cal{O}_Y
\]
es universalmente inyectivo en~$x$; o, lo que es equivalente, la sucesión de homomorfismos canónicos
\[
0\to\cal{J}/\cal{J}^2\to\it{\Omega}^1_{X/S}\otimes\cal{O}_Y\to
\it{\Omega}^1_{Y/S}\to 0
\]
es exacta en $x$, y $\it{\Omega}^1_{Y/S}$ es localmente libre en~$x$.
\end{enumerate}
\end{theoreme}

\subsubsection*{Demostración}
Ya
\marginpar{42}
sabemos, por lo anterior, que (ii) implica (i), (iii), (iv). Probemos que (i)$\To$(ii) (lo que concluirá al mismo tiempo la demostración
\ifthenelse{\boolean{orig}}{de~\Ref{I.4.9})}{de~\Ref{II.4.9}).}
En virtud del teorema~\Ref{II.4.3}~(ii), los dos últimos términos de la sucesión exacta~\eqref{eq:II.4.3} son módulos libres sobre~$Y$. Como las imágenes en $\it{\Omega}^1_{X/S}\otimes_{\cal{O}_X}\cal{O}_Y$ de los $\rd g$ ($g\in\cal{O}_X$) engendran este módulo en $x$, y sus imágenes en $\it{\Omega}^1_{Y/S}$ engendran por tanto este último en~$x$, se pueden encontrar $g_i$ ($p+1\leq i\leq n$) en $\cal{O}_X$ tales que los $\rd g'_i$ formen una base de~$\it{\Omega}^1_{Y/S}$; luego, por la exactitud de~\eqref{eq:II.4.3}, se puede completar el sistema de los $\rd g_i$ ($p+1\leq i\leq n$) hasta una base del término intermedio con elementos de la forma $\rd g_i$ ($1\leq i\leq p$), donde los $g_i$ ($1\leq i\leq p$) \emph{pertenecen a $\cal{J}_x$}. Los $g_i$ provienen de secciones de~$\cal{O}_X$ sobre un entorno de~$x$ en $X$, que podemos suponer igual a~$X$. Estamos entonces bajo la condición
\ifthenelse{\boolean{orig}}{\Ref{II.4.8}~(ii)}{\Ref{II.4.9}~(ii)},
y hemos establecido que esta implica la condición
\ifthenelse{\boolean{orig}}{\Ref{II.4.8}~(i)}{\Ref{II.4.9}~(i)},
de donde se sigue~\Ref{II.4.10}~(ii).

La implicación (iii)$\To$(ii) en~\Ref{II.4.10} se sigue inmediatamente de la implicación (iii)$\To$(i) en
\ifthenelse{\boolean{orig}}{\Ref{II.4.8}}{\Ref{II.4.9}}.
Así,
\ifthenelse{\boolean{orig}}{(i) (ii) (iii)}{(i), (ii), (iii)}
son equivalentes e implican (iv). Finalmente, es trivial que (iv) implica (iii), teniendo en cuenta que unos $g_i\in\cal{J}_x$ que forman una base de~$\cal{J}_x\bmod\cal{J}_x^2$ engendran $\cal{J}_x$ (Nakayama).

Además, la demostración precedente muestra lo siguiente:
\begin{corollaire}
\label{II.4.11}
Sean $X$ un $S$-preesquema, $Y$ un subpreesquema cerrado definido por un haz coherente $\cal{J}$ de ideales sobre~$X$, y $x$ un punto de~$Y$. Supongamos que \emph{$X$ e $Y$ son lisos sobre~$S$ en $x$}. Sean $g_i$ secciones de~$\cal{J}$ ($1\leq i\leq p$). Las condiciones siguientes son equivalentes:
\begin{enumerate}
\item[(i)] Los $g_i$ \emph{engendran} $\cal{J}_x$ y los $\rd g_i(x)$ son \emph{linealmente independientes} en $\it{\Omega}^1_{X/S}(x)$
\ifthenelse{\boolean{orig}}{sobre~$(x)$.}{sobre~$\kres(x)$.}
\item[(ii)] Los $g_i\bmod\cal{J}^2$ forman una base de $\cal{J}/\cal{J}^2$ en $x$.
\item[(iii)] Los $g_i$ forman un sistema minimal de generadores de~$\cal{J}_x$.
\item[(iv)] Se pueden encontrar otras secciones $g_i$ ($p+1\leq i\leq n$) de~$\cal{O}_X$ sobre un entorno $X_1$
\ifthenelse{\boolean{orig}}{de~$X$}{de~$x$ en~$X$},
que, junto con las anteriores, definan un morfismo \emph{étale} $X_1\to X'=S[t_1,\dots,t_n]$ tal que $Y_1=Y\cap X_1$ sea la \emph{imagen inversa} por $g$ del subpreesquema cerrado $Y'=S[t_{p+1},\dots,t_n]$ de~$X'=S[t_1,\dots,t_n]$.
\end{enumerate}
\end{corollaire}

En
\marginpar{43}
particular:

\begin{corollaire}
\label{II.4.12}
Sean $X$ un $S$-preesquema, $F$ una sección de~$\cal{O}_X$, $Y$ el subpreesquema de los ceros de~$F$ (definido por el ideal coherente $F\cdot\cal{O}_X$), y $x$ un punto de~$Y$. Supongamos que $X$ es simple sobre~$S$ en $x$. Para que $Y$ sea liso sobre~$S$ en $x$, es necesario y suficiente que, o bien $F$ sea nulo en un entorno de $x$, o bien $\rd F(x)\neq0$ (donde $\rd F(x)$ designa la imagen de $\rd F$ en el espacio vectorial $\it{\Omega}^1_{X/S}(x)$ sobre~$\kres(x)$).
\end{corollaire}

Es suficiente por~\Ref{II.4.10}, criterio~(iii). Es necesario porque, como $\cal{J}$ está engendrado por un elemento, hace falta primero que $\cal{J}/\cal{J}^2$ en el punto $x$ sea libre de rango $\leq1$. Si este rango es $0$, es decir, si $\cal{J}/\cal{J}^2=0$ en $x$, se sigue que $\cal{J}=0$ en $x$ por Nakayama, es decir, $F$ es nulo en un entorno de~$x$. Si este rango es $1$, entonces $F$ forma un sistema minimal de generadores de~$\cal{J}$ en $x$, y se concluye por~\Ref{II.4.11}, equivalencia de (i) y~(iii).

\begin{corollaire}
\label{II.4.13}
Sean $Y$ un $S$-preesquema localmente de tipo finito, $S'$ un $S$-preesquema \emph{plano}, $Y'=Y\times_S S'$, $x'$ un punto de~$Y'$ y $x$ su imagen canónica en~$Y$. Para que $Y$ sea liso sobre~$S$
\ifthenelse{\boolean{orig}}{en~$X$}{en~$x$},
es necesario y suficiente que $Y'$ sea liso sobre~$S'$ en~$x'$. En particular, si $S'\to S$ es plano y sobreyectivo, $Y$ es liso sobre~$S$
\ifthenelse{\boolean{orig}}{\emph{sss}}{si y solo si\ }
$Y'$ es liso sobre~$S'$.
\end{corollaire}

Solo hay que demostrar la suficiencia (la necesidad se vio en~\Ref{II.1.3}~(ii)). Podemos suponer (reemplazando $Y$ por un entorno conveniente de~$x$ y $Y'$ por su imagen inversa) que $Y$ es afín de tipo finito sobre~$S$ afín; por tanto, $Y$ es isomorfo a un subpreesquema cerrado de un esquema $X=S[t_1,\dots,t_n]$. En consecuencia, $Y'$ se identifica con un subpreesquema cerrado de~$X'=X\times_S S'$. Como $X$ es liso sobre~$S$, también $X'$ es liso sobre~$S'$, y se pueden aplicar los criterios de lisura~\Ref{II.4.10}. Aquí, el criterio~(iv) da inmediatamente el resultado.


\begin{remarques}
\label{II.4.14}
El criterio~(iii) del teorema~\Ref{II.4.10} merece llamarse \emph{criterio jacobiano de lisura}.
\index{criterio jacobiano de lisura|hyperpage}%
\index{jacobiano (criterio de lisura)|hyperpage}%
\index{lisura (criterio jacobiano de)|hyperpage}%
Este permite reconocer, teóricamente, si un $S$-preesquema dado $Y$ es liso sobre~$S$ en un punto $x$ de~$Y$, pues siempre existe un entorno de~$x$ en~$Y$ isomorfo a un subpreesquema de un $S$-preesquema liso $X$, por ejemplo $X=S[t_1,\dots,t_n]$. Es, por otra parte, para $X=S[t_1,\dots,t_n]$, $S=\Spec(A)$, como suele enunciarse el criterio jacobiano (naturalmente, en el caso clásico considerado por Zariski, $A$ era
\marginpar{44}
un cuerpo). Se deja al lector formular el enunciado correspondiente a un ideal $J$ de~$A[t_1,\dots,t_n]$ y a un ideal primo que lo contenga, al que así se llega. Señalemos que, al parecer, en la actualidad (y sobre todo desde que Nagata consiguió generalizar por métodos no diferenciales el teorema de Zariski según el cual el conjunto de puntos regulares de un esquema algebraico es abierto), el criterio jacobiano apenas presenta interés salvo en la forma en que lo damos aquí (es decir, utilizando exclusivamente diferenciales \emph{relativas}, y no diferenciales \emph{absolutas}, esto es, relativas al anillo de constantes absolutas $\ZZ$). Como ocurre con frecuencia, la consideración de las diferenciales es aquí más cómoda que la de las derivaciones. Señalemos finalmente que, si $Y$ es liso sobre~$S$ en $x$, de dimensión relativa
\ifthenelse{\boolean{orig}}{$n$}{$m$},
entonces existe un entorno abierto de~$x$ en~$Y$ isomorfo a un subpreesquema de $X=S[t_1,\dots,t_n]$, con $n=m+1$, como resulta de la definición y
\ifthenelse{\boolean{orig}}{de~\Ref{I.7.6}).}{de~(\Ref{I.7.6}).}

Sean $A$ un anillo noetheriano, $x_i$ ($1\leq i\leq n$) elementos de~$A$, y $J$ el ideal engendrado por los $x_i$. Se dice que los $x_i$ forman un \emph{sistema regular de generadores}
\index{regular (sistema de generadores)|hyperpage}%
\index{sistema regular de generadores|hyperpage}%
de~$J$ si el homomorfismo sobreyectivo canónico
\[
(A/J)[t_1,\dots,t_n]\to\gr^J(A)
\]
definido por los $x_i$ (donde el segundo miembro designa el anillo graduado asociado a $A$, filtrado por las potencias de~$J$) es un \emph{isomorfismo}. Esta condición significa también que
\begin{enumerate}
\item[(i)] el homomorfismo sobreyectivo canónico
\[
S_{A/J}(J/J^2)\to\gr^J(A)
\]
(donde el primer miembro designa el álgebra simétrica del $A/J$-módulo $J/J^2$) es un isomorfismo, y
\item[(ii)] $J/J^2$ es libre y admite como base las clases de los $x_i\bmod J^2$.
\end{enumerate}
En esta forma se ve que, si $J\neq A$, los $x_i$ forman un \emph{sistema minimal de generadores} de~$J$, y que \emph{cualquier otro sistema minimal de generadores} de~$J$ \emph{es un sistema regular de generadores} (nótese que «minimal» se toma en sentido estricto: número mínimo de elementos, sentido que solo equivale a «minimal para la inclusión» cuando $A$ es local); por otra parte, si $J=A$, todo sistema de generadores de~$J$ es regular.

La
\marginpar{45}
condición de regularidad de un sistema de generadores de un ideal es estable por localización respecto de cualquier conjunto multiplicativamente cerrado. Por otra parte, se ve enseguida que, para que $(x_i)$ sea un sistema minimal de generadores de~$J$, basta ya con que, para todo ideal \emph{maximal $\goth m$ que contenga a $J$}, los $x_i$ definan un sistema regular de generadores de~$JA_{\goth m}$ en $A_{\goth m}$. Esto nos reduce al caso en que $A$ es un anillo local de ideal maximal $\goth m$ y los $x_i$ pertenecen a $\goth m$. \emph{Entonces los $x_i$ forman un sistema regular de generadores de~$J$ si y solo si forman una $A$-sucesión en el sentido de Serre}\footnote{Ahora diremos más bien «sucesión $A$-regular», \cf EGA $0_{\textup{IV}}$~15.1.7 y 15.1.11.}, es decir, si para todo $i$ tal que $1\leq i\leq n$, $x_i$ no es divisor de cero en
\ifthenelse{\boolean{orig}}{$A/(x_1,\dots,x_i)A$.}{$A/(x_1,\dots,x_{i-1})A$.}

Finalmente, en el caso en que $A$ sea un álgebra sobre un anillo $B$ y $A/J$ sea isomorfo como $B$-álgebra a~$B$ (de modo que $J$ sea el núcleo de un homomorfismo de $B$-álgebras $A\to B$), los $x_i$ forman un sistema regular de generadores de~$J$ si y solo si el homomorfismo canónico
\[
B[[t_1,\dots,t_n]]\to\hat{A}
\]
definido por los $x_i$ (donde el segundo miembro designa la completación separada $\varprojlim A/J^{n+1}$ de~$A$ para la topología definida por las potencias de~$J$) es un \emph{isomorfismo} (en todo caso, es \emph{sobreyectivo}).

Todos estos hechos son bien conocidos y, salvo pequeñas diferencias, figuran sin duda en el curso de álgebra conmutativa de Serre redactado por Gabriel. (Allí se encuentran otras $N$ caracterizaciones de las $A$-sucesiones cuando $A$ es un anillo local).

Sea $J$ un ideal de un anillo noetheriano $A$. Diremos que $J$ es un \emph{ideal regular}
\index{ideal regular|hyperpage}%
\index{regular (ideal)|hyperpage}%
si, para todo ideal primo $\goth p$ de~$A$, $JA_{\goth p}$ admite un sistema regular de generadores. Evidentemente, basta comprobarlo para $\goth p\supset J$ y, además, se puede restringir a $\goth p$ maximal. Más generalmente, sea $\cal{J}$ un ideal sobre un preesquema localmente noetheriano $X$; se dice que $\cal{J}$ es un \emph{ideal regular} si, para todo $x\in X$, $\cal{J}_x$ es un ideal de $\cal{O}_x$ que admite un sistema regular de generadores. Esto equivale a la conjunción de las dos condiciones siguientes:
\begin{enumerate}
\item[(a)] El homomorfismo sobreyectivo canónico
\[
\ifthenelse{\boolean{orig}}
{S_{\cal{O}/\cal{J}}(\cal{J}/\cal{J}^2)\to\gr^{\cal{J}}(\cal{O}_X)}
{S_{\cal{O}_X/\cal{J}}(\cal{J}/\cal{J}^2)\to\gr^{\cal{J}}(\cal{O}_X)}
\]
es un isomorfismo, y
\item[(b)] el haz de $\cal{O}_X/\cal{J}$-módulos $\cal{J}/\cal{J}^2$ es localmente libre.
\end{enumerate}

Se
\marginpar{46}
dice entonces también que el subpreesquema $Y$ de~$X$ definido por $\cal{J}$ (por tanto, tal que $\cal{O}_Y$, prolongado por cero, sea isomorfo a $\cal{O}_X/\cal{J}$) está \emph{inmerso regularmente}
\index{inmerso regularmente (preesquema)|hyperpage}%
en~$X$, y se define del mismo modo (de manera evidente) la noción de morfismo que es una \emph{inmersión regular},
\index{inmersión regular|hyperpage}%
\index{regular (inmersión)|hyperpage}%
(\resp \emph{regular en un punto $x$}), es decir, un morfismo de inmersión $Y\to X$ que identifica $Y$ (\resp un entorno conveniente de $x$) con un subpreesquema cerrado inmerso regularmente en un abierto de~$X$. (No se debe decir «subpreesquema regular», pues esto significaría que los anillos locales de~$Y$ son regulares). Finalmente, unas secciones $x_i$ de~$\cal{J}$ se llaman \emph{sistema regular de generadores} si, para todo $x\in X$, los elementos correspondientes de~$\cal{O}_x$ forman un sistema regular de generadores de~$\cal{J}_x$ (terminología compatible con la introducida para los generadores de un ideal de un anillo). Esto significa también que el homomorfismo sobreyectivo canónico
\[
\cal{O}_Y[t_1,\dots,t_n]\to\gr^{\cal{J}}(\cal{O}_X)
\]
definido por los $x_i$ es un isomorfismo. Si se sabe de antemano que el ideal $\cal{J}$ es regular, esto significa también, sencillamente, que en todo punto $x$ \emph{de~$Y$}, los $x_i$ definen una \emph{base} de~$\cal{J}/\cal{J}^2$ sobre~$\cal{O}_{Y,x}$. (Nótese que esta condición es vacía si $Y$ es vacío). Así, para que $\cal{J}$ admita un sistema regular de generadores, es necesario y suficiente que $\cal{J}$ sea regular y que el $\cal{O}_Y$-módulo $\cal{J}/\cal{J}^2$ sea globalmente libre (y no solo localmente libre), es decir, que el homomorfismo canónico $S_{\cal{O}_Y}(\cal{J}/\cal{J}^2)\to\gr^{\cal{J}}(\cal{O}_X)$ sea sobreyectivo y que el $\cal{O}_Y$-módulo $\cal{J}/\cal{J}^2$ sea globalmente libre.

Un \emph{anillo aumentado se dice regular}
\index{anillo aumentado regular|hyperpage}%
\index{regular (anillo aumentado)|hyperpage}%
si el ideal de aumento es regular. Así, si $A$ es un anillo local considerado como aumentado sobre su cuerpo residual $k$, entonces $A$ es un anillo local regular si y solo si es un anillo aumentado regular.

(A decir verdad, parece que era innecesario comenzar con el sorites preliminar para los anillos; conviene empezar directamente con los haces. Si se quiere algo en el caso noetheriano, es la definición adoptada aquí ---a priori menos restrictiva que la dada mediante las $A$-sucesiones de Serre--- la que parece preferible para las necesidades del cálculo diferencial. Naturalmente, para hacer bien las cosas, habría que desarrollar también al menos una parte de la teoría de los morfismos lisos en el marco no noetheriano\footnote{Como se dice en el prefacio, esto ya se ha hecho, \cf EGA IV~17, 18}, probablemente
\marginpar{47}
partiendo del criterio jacobiano, de modo que se obtengan, si es posible, todas las propiedades formales esenciales de los morfismos lisos y de los morfismos étale, es decir, lisos y cuasifinitos; solo los recíprocos recurrirían a hipótesis noetherianas).
\end{remarques}

Después de estos largos preliminares terminológicos, un pequeño teorema:


\begin{theoreme}
\label{II.4.15}
Sean $X$ un $S$-preesquema localmente de tipo finito, $Y$ un subpreesquema cerrado de~$X$ definido por un haz coherente $\cal{J}$ de ideales sobre~$X$, y $x$ un punto de~$X$. Supongamos ahora que \emph{$Y$ es liso sobre~$S$ en $x$} (sin suponer nada sobre~$X$). Entonces son equivalentes las condiciones siguientes:
\begin{enumerate}
\item[(i)] $X$ es liso sobre~$S$ en $x$;
\item[(ii)] la inmersión $i\colon Y\to X$ es regular en $x$, es decir, $\cal{J}_x$ es un ideal regular de~$\cal{O}_x$.
\end{enumerate}
\end{theoreme}

\begin{corollaire}
\label{II.4.16}
Supongamos que $Y$ es \emph{liso} sobre~$S$. Para que $X$ sea liso sobre~$S$ en un entorno de~$Y$ (es decir, en los puntos de~$Y$), es necesario y suficiente que $Y$ esté inmerso regularmente en~$X$, es decir, que la inmersión $i\colon Y\to X$ sea regular.
\end{corollaire}

\subsubsection*{Demostración}
(i) implica (ii). Aplicamos el criterio~(ii) de~\Ref{II.4.10}; como
\ifthenelse{\boolean{orig}}{$g\colon X_1\to X$}{$g\colon X_1\to X'$}
es \emph{plano}, para demostrar que la imagen inversa por $g$ del subpreesquema $Y'$ de~$X'$ está inmersa regularmente, nos reducimos a probar que $Y'=S[t_{p+1},\dots,t_n]$ está inmerso regularmente en $S[t_1,\dots,t_n]$, lo cual es trivial (los $t_i$ ($1\leq i\leq p$) forman un sistema regular de generadores del ideal que define $Y'$ en $X'$).

(ii) implica (i). Sea $g_i$ ($1\leq i\leq p$) un sistema regular de generadores de~$\cal{J}_x$, y sean $g_i$ ($p+1\leq i\leq n$) elementos de~$\cal{O}_{X,x}$ tales que sus imágenes $g'_i$ en $\cal{O}_{Y,x}$ definan un morfismo \emph{étale}
\[
Y_1\to Y'=S[t_{p+1},\dots,t_n]
\]
\ifthenelse{\boolean{orig}}
{de un entorno $Y_1$ de~$Y$ en~$Y'$}
{de un entorno $Y_1$ de~$x$ en~$Y$}.
Los $g_i$ ($1\leq i\leq n$) provienen de secciones (denotadas del mismo modo) de~$\cal{O}_X$ sobre un entorno $X_1$ de~$x$, y podemos suponer que $X_1=X$, $Y_1=Y$. Se obtiene así un morfismo
\[
g\colon X\to X'=S[t_1,\dots,t_n],
\]
y todo se reduce a probar que este morfismo es \emph{étale} en~$x$. Tomando $X_1$ suficientemente pequeño,
\marginpar{48}
podemos suponer que los $g_i$ ($1\leq i\leq p$) forman un sistema regular de generadores de~$\cal{J}$ en todo~$X$. En particular, engendran $\cal{J}$, de modo que el subpreesquema $Y$ de~$X$ se identifica con la imagen inversa por $g$ del subpreesquema $Y'$ de~$X'$. Sea $x'=g(x)$. Entonces
\ifthenelse{\boolean{orig}}
{la fibra de~$X'\to X$ en~$x'$}
{la fibra de~$X\to X'$ sobre~$x'$}
es
\ifthenelse{\boolean{orig}}{por tanto}{\ignorespaces}
idéntica a la fibra de~$Y\to Y'$ sobre~$x'$ y, por consiguiente, es étale sobre~$\kres(x')$; por tanto, $g$ es \emph{no ramificado} en $x$. Falta probar que $g$ es \emph{plano} en $x$. Ahora bien, el graduado asociado a~$\cal{O}_{X',x'}$, filtrado por las potencias de
\ifthenelse{\boolean{orig}}{$\cal{J}'_x$}{$\cal{J}'_{x'}$},
es \emph{libre} sobre~$\cal{O}_{Y',x'}$ en todos los grados. Por otra parte, el graduado asociado a~$\cal{O}_{X,x}$, filtrado por las potencias de
\ifthenelse{\boolean{orig}}
{$\cal{J}_x=\cal{J}'_x\cal{O}_{X,x}$}
{$\cal{J}_x=\cal{J}'_{x'}\cal{O}_{X,x}$},
es isomorfo (por el homomorfismo canónico) al producto tensorial del anterior por $\cal{O}_{Y,x}$, pues uno y otro anillo son anillos de polinomios en
\ifthenelse{\boolean{orig}}{$n-p$}{$p$}
indeterminadas, con anillo de constantes $\cal{O}_{Y',x'}$, \resp $\cal{O}_{Y,x}$. Finalmente, sobre $\cal{O}_{X',x'}/\cal{J}'_{x'}=\cal{O}_{Y',x'}$, $\cal{O}_{X,x}/\cal{J}_x=\cal{O}_{Y,x}$ es plano.

Según un criterio general de planitud (válido para un homomorfismo local de anillos locales noetherianos $A'\to A$, estando $A'$ dotado de un ideal $J'\neq A'$ tal que el graduado asociado sea libre sobre~$A'/J'$ en todo grado), se sigue que $X$ es plano sobre~$X'$ en $x$, como se quería demostrar.

\begin{corollaire}
\label{II.4.17}
Sean $X$ un preesquema localmente de tipo finito sobre~$Y$, $i$ una sección de~$X$ sobre~$Y$, $y$ un punto de~$Y$, $x=i(y)$, y $\cal{J}$ el haz de ideales sobre~$X$ definido por el subpreesquema $i(Y)$ (que suponemos cerrado para simplificar el enunciado, condición que se cumple si $X$ es un esquema).

Las condiciones siguientes son equivalentes:
\begin{enumerate}
\item[(i)] $X$ es liso sobre~$Y$ en $x$;
\item[(ii)] $i$ es una inmersión regular en $y$;
\item[(iii)] la $\cal{O}_y$-álgebra obtenida completando $\cal{O}_x$ para la topología definida por las potencias de~$\cal{J}_x$ es isomorfa a un álgebra de series formales $\cal{O}_y[[t_1,\dots,t_n]]$;
\item[(iii~bis)] existe un entorno abierto $U$ de~$y$ tal que el haz de álgebras $\varprojlim i^*(\cal{O}_X/\cal{J}^{n+1})$ sobre~$\cal{O}_Y$ sea isomorfo, sobre~$U$, a un haz de la forma $\cal{O}_Y[[t_1,\dots,t_n]]$;
\item[(iv)] existen un entorno abierto $U$ de~$y$, un entorno abierto $V$ de~$x$ y, finalmente, un $Y$-morfismo $g\colon V\to U[t_1,\dots,t_n]$ tal que $g$ sea étale, que
\ifthenelse{\boolean{orig}}{$s$}{$i$}
induzca una sección de~$V$ sobre~$U$, transformada por $g$ en la sección cero de~$U[t_1,\dots,t_n]$ sobre~$U$.
\end{enumerate}
\end{corollaire}

La equivalencia
\marginpar{49}
de (i) y (ii) es un caso particular del teorema~\Ref{II.4.15}, tomando $Y=S$; la equivalencia de (ii) y (iii) (y, moralmente, de (ii) y (iii~bis)) se señaló junto con los «recordatorios». En cuanto a la equivalencia de (i) y (iv), se deduce fácilmente del teorema~\Ref{II.4.10} (equivalencia de las condiciones (i) y (ii) de dicho teorema).

\begin{corollaire}
\label{II.4.18}
Sea $X$ un preesquema liso sobre~$S$. Entonces el morfismo diagonal
\[
\Delta_{X/S}\colon X\to X\times_S X
\]
es una
\ifthenelse{\boolean{orig}}
{\emph{inmersión regular}.}
{\emph{inmersión regular},}
o, como se dice también, $X$ es «\emph{diferencialmente liso}» sobre~$S$.
\index{diferencialmente liso|hyperpage}%
\end{corollaire}

En efecto, este es un caso particular del corolario~\Ref{II.4.16}, pues $X$ y $X\times_S X$ son ambos lisos sobre~$S$.

\setcounter{subsection}{17}
\setcounter{proposition}{17}

\begin{remarques}
\label{rem:II.4.18}
Recordemos (\Ref{I.1}) que, si $X$ es un preesquema sobre~$S$, se introducen los haces cuasicoherentes
\ifthenelse{\boolean{orig}}{de Álgebras}{de álgebras}
$\cal{P}_{X/S}^n=\cal{O}_{X\times_S X}/\cal{I}_X^{n+1}$ sobre~$X$ (donde $\cal{I}_X$ designa el haz
\ifthenelse{\boolean{orig}}{de Ideales}{de ideales}
que define la diagonal en $X\times_S X$), considerado como haz de $\cal{O}_X$-álgebras mediante la primera proyección $\pr_1\colon X\times_S X\to X$. Los $\cal{P}_{X/S}^n$ forman un sistema proyectivo de álgebras sobre~$X$, cuyo límite proyectivo se denota por $\cal{P}_{X/S}^{\infty}$ y no es sino el haz estructural de la completación formal de~$X\times_S X$ a lo largo de la diagonal (suponiendo ahora que $X$ es localmente de tipo finito sobre~$S$, de modo que los $\cal{P}_{X/S}^n$ sean coherentes). Decir que $X$ es diferencialmente liso sobre~$S$ (es decir, que el morfismo diagonal $\Delta_{X/S}$ es una inmersión regular) significa también que $\cal{P}_{X/S}^{\infty}$ es regular como haz de álgebras aumentado hacia $\cal{O}_X$, es decir, que $\it{\Omega}_{X/S}^1$ es localmente libre y que el homomorfismo sobreyectivo canónico
\[
S_{\cal{O}_X}(\it{\Omega}_{X/S}^1)\to
\mathrm{gr}_*(\cal{P}_{X/S}^{\infty})
\]
es un isomorfismo; o, finalmente, que todo punto de~$X$ tiene
un entorno abierto sobre el cual el haz
\ifthenelse{\boolean{orig}}{de Álgebras}{de álgebras}
aumentadas $\cal{P}_{X/S}^{\infty}$ es isomorfo a un haz $\cal{O}_X[[t_1,\dots,t_n]]$.

Sea $s$ una sección de~$X$ sobre~$S$, y sea $\cal{J}$ el haz de ideales sobre~$X$ que ella define (suponiendo,
\marginpar{50}
para simplificar, que $s(S)$ es cerrado). Se tienen entonces isomorfismos canónicos de $\cal{O}_X$-álgebras aumentadas:
\begin{equation}
\label{eq:II.4.4}
s^*(\cal{P}_{X/S}^n)=\cal{O}_X/\cal{J}^{n+1}\quoi,\quad
s^*(\cal{P}_{X/S}^{\infty})=\varprojlim_n
\cal{O}_X/\cal{J}^{n+1}.
\end{equation}
Estos isomorfismos son funtoriales, en un sentido evidente, por cambio de base y (teniendo en cuenta este hecho) vuelven a dar una caracterización de los haces de álgebras $\cal{P}_{X/S}^n$ sobre~$S$. Si, por ejemplo, $S=\Spec(k)$, con $k$ un cuerpo, dar una sección $s$ de~$X$ sobre~$S$ equivale a dar un punto $x$ de~$X$ racional sobre~$k$, y las fórmulas precedentes significan que se tiene un isomorfismo de $k$-álgebras
\begin{equation}
\label{eq:II.4.5}
\cal{P}_{X/S}^n(x)=\cal{O}_x/\goth{m}_x^{n+1},
\end{equation}
lo que justifica el nombre «\emph{haz de partes principales de orden~$n$ sobre~$X$ relativo a~$S$}»
\index{partes principales de orden~$n$ (haz de, sistema de)|hyperpage}%
\label{indnot:ad1}\oldindexnot{$\cal{P}_{X/Y}^n$|hyperpage}%
dado a~$\cal{P}_{X/S}^n$. Se ve además en~\eqref{eq:II.4.4} que \emph{si $X$ es diferencialmente liso sobre~$S$ en todo punto de~$s(S)$, entonces $X$ es liso sobre~$S$ en todo punto de~$s(S)$} (corolario~\Ref{II.4.17}), \emph{siendo también cierto el recíproco} (corolario~\Ref{II.4.18}). Teniendo en cuenta~\Ref{II.4.13}, se concluye fácilmente que, si $X$ es un $S$-preesquema localmente de tipo finito, \emph{$X$ es liso sobre~$S$ si y solo si es plano sobre~$S$ y diferencialmente liso sobre~$S$}. (Nótese que la hipótesis de planitud es esencial, como se ve tomando por $X$ un subpreesquema cerrado de~$S$).

Recordemos también que se obtiene una \emph{segunda estructura de álgebra} sobre~$\cal{P}_{X/S}^n$ mediante la proyección $\pr_2\colon X\times_S X\to X$, la cual se deduce, por otra parte, de la anterior mediante la \emph{involución canónica} del haz de anillos $\cal{P}_{X/S}^n$, inducida por el automorfismo de simetría de~$X\times_S X$. Se denota por $d_{X/S}^n$
\label{indnot:bb}\oldindexnot{$d_{X/S}^n$|hyperpage}%
o simplemente por $d^n$ el homomorfismo de haces de anillos
\begin{equation}
\label{eq:II.4.6}
d_{X/S}^n\colon\cal{O}_X\to\cal{P}_{X/S}^n
\end{equation}
que corresponde a esta segunda estructura de álgebra. Teniendo en cuenta el isomorfismo~\eqref{eq:II.4.4}, este homomorfismo transforma una sección $f$ de~$\cal{O}_X$ en una sección $d^n(f)$ de~$\cal{P}_{X/S}^n$ cuya imagen inversa por una sección $s$ de~$X$ sobre~$S$ se identifica con la imagen canónica de~$f$ en $\Gamma(X,\cal{O}_X/\cal{J}^{n+1})$. Esto justifica el nombre de «\emph{sistema de partes principales de orden~$n$ de~$f$}»
\index{sistema de partes principales de orden n de una sección@sistema de partes principales de orden~$n$ de una sección|hyperpage}%
dado a~$d^nf$,
\marginpar{51}
especialmente en el caso $S=\Spec(k)$ considerado en la fórmula~\eqref{eq:II.4.5}.

Para terminar, señalemos que el homomorfismo~\eqref{eq:II.4.6} puede considerarse como el \emph{operador diferencial universal de orden}~$\leq n$\kern1pt\footnote{Para todo lo relativo al presente párrafo, puede consultarse EGA IV~16.8 a 16.12.} (relativo al preesquema de constantes $S$) sobre~$\cal{O}_X$, conviniendo en llamar operador diferencial de orden~$\leq n$
\index{operador diferencial universal de orden~$\leq n$|hyperpage}%
\index{operador diferencial de orden~$\leq n$|hyperpage}%
de~$\cal{O}_X$ en un módulo $F$ a un homomorfismo de haces $D$ que se factoriza como
\[
D\colon\cal{O}_X\lto{d^n}\cal{P}_{X/S}^n\lto{u}F,
\]
donde $u$ es un homomorfismo \emph{de $\cal{O}_X$-módulos}, determinado de manera única por~$D$. Esta definición concuerda con la definición recursiva intuitiva ($D$ es un operador diferencial de orden $\leq n$ si, para toda sección $g$ de~$\cal{O}_X$ sobre un abierto $U$ de~$X$,
\ifthenelse{\boolean{orig}}
{$f\mto D(fg)-D(f)$}
{$f\mto D(fg)-gD(f)$}
es un operador diferencial de orden $\leq n-1$ sobre~$U$). Se sigue que \emph{si $X$ es diferencialmente liso sobre~$S$, el haz de anillos de operadores diferenciales de todos los órdenes tiene la estructura sencilla bien conocida} en el cálculo diferencial sobre variedades diferenciables y, en particular, admite localmente una base de $\cal{O}_X$-módulo
\ifthenelse{\boolean{orig}}{cerrada}{formada}
por las \emph{potencias divididas} de unos operadores que conmutan $\delta/\delta x_i$ ($1\leq i\leq n$). Si
\ifthenelse{\boolean{orig}}{$S$}{$\cal{O}_S$}
es un haz de $\QQ$-álgebras ($\QQ=$ el cuerpo de los números racionales), basta tomar los polinomios ordinarios en los $\delta/\delta x_i$. En este caso, además, y de modo muy excepcional, para que $X$ sea diferencialmente liso sobre~$S$, basta ya con que $\it{\Omega}_{X/S}^1$ sea localmente libre.
\end{remarques}

\begin{remarque}
\label{II.4.19}
\index{Cohen--Macaulay (esquema de)|hyperpage}%
La terminología «inmersión regular», «ideal regular», etc., introducida en este número, encontró una oposición bastante viva y general (Chevalley, Serre). Se propusieron «ideal de Cohen--Macaulay», «ideal de Macaulay» o «ideal macaulayano», lo cual obligaría moralmente a adoptar también «inmersión de Cohen--Macaulay» o «inmersión de Macaulay». Sin embargo, esta terminología entra en conflicto con otra ya empleada en futuras redacciones del \emph{multiplodoque}, donde un morfismo (de tipo finito) se dice «Cohen--Macaulay» en un punto si es plano en dicho punto y la fibra que pasa por él tiene allí un anillo local que es un anillo de Cohen--Macaulay. Mientras se encuentra una solución satisfactoria, mantendremos, con toda reserva, la terminología introducida en este número\footnote{Es la adoptada en EGA~$\mathrm{0}_{\mathrm{IV}}$~15.1.7.}.
\end{remarque}


\section{Caso de un cuerpo base}
\label{II.5}
\marginpar{52}

\begin{proposition}
\label{II.5.1}
Sean $k$ un cuerpo, $X$ un preesquema de tipo finito sobre~$k$, $x$ un punto de~$X$ y $n$ la dimensión de~$X$ en~$x$, y sea
\ifthenelse{\boolean{orig}}
{$f\colon X\to\Spec k[t_1,\dots,tn]=Y$}
{$f\colon X\to\Spec k[t_1,\dots,t_n]=Y$}
un morfismo definido por elementos
\ifthenelse{\boolean{orig}}
{$f_i\in\Gamma(X:\cal{O}_X)$.}
{$f_i\in\Gamma(X,\cal{O}_X)$.}
Las condiciones siguientes son equivalentes (e implican que $X$ es liso sobre~$k$ en $x$ y, con mayor razón, regular en $x$ por~\Ref{II.3.1}):
\begin{enumerate}
\item[(i)] $f$ es étale en $x$.
\item[(ii)] Los $df_i$ forman una base de~$\it{\Omega}_{X/k}^1$ en $x$.
\item[(iii)] Los $df_i$ engendran $\it{\Omega}_{X/k}^1$ en $x$.
\end{enumerate}
\end{proposition}

Como (i) implica que $X$ es liso sobre~$k$ en $x$, la implicación (i)$\To$(ii) es un caso particular de~\Ref{II.4.8}; (ii)$\To$(iii) es trivial, y queda probar (iii)$\To$(i). Ahora bien, si se cumple (iii), $f$ es neto en $x$ por el lema~\Ref{II.4.1}; por tanto, sustituyendo $x$ por un entorno abierto, es cuasifinito y, por razones de dimensión, dominante. Como $Y$ es regular, se sigue que $f$ es étale por (\Ref{I.9.5}~(ii)) o (\Ref{I.9.11}).

\begin{corollaire}
\label{II.5.2}
Bajo las hipótesis preliminares de~\Ref{II.5.1}, supongamos que
\ifthenelse{\boolean{orig}}{$\kappa(x)$}{$\kres(x)$}
es una extensión \emph{finita separable} de~$k$ y que los $f_i$ ($1\leq i\leq n$) definen elementos de~$\goth{m}_x$. Entonces las condiciones anteriores equivalen a:
\begin{enumerate}
\item[(iv)] Los $f_i$ forman un sistema de generadores de~$\goth{m}_x$ (o, equivalentemente, los $f_i$ módulo $\goth{m}_x^2$ forman una base de~$\goth{m}_x/\goth{m}_x^2$ sobre
\ifthenelse{\boolean{orig}}{$\kappa(x)$}{$\kres(x)$}).
\end{enumerate}
\end{corollaire}

En efecto, (iv)$\To$(iii) en virtud de la sucesión exacta
\begin{equation}
\label{eq:II.5.1}
\ifthenelse{\boolean{orig}}
{\goth{m}_x/\goth{m}_x^2\to\it{\Omega}_{\cal{O}_x/k}^1\to
\it{\Omega}_{\kappa(x)/k}^1\to 0}
{\goth{m}_x/\goth{m}_x^2\to\Omega_{\cal{O}_x/k}^1\to
\Omega_{\kres(x)/k}^1\to 0}.
\end{equation}
Además,
\ifthenelse{\boolean{orig}}
{$\Omega_{\kappa(x)/k}^1=0$, pues $\kappa(x)$ es étale sobre~$k$. Por otra parte, (ii) implica (iv), ya que, como $X$ y $\Spec(\kappa(x))$ son lisos sobre~$k$ en $x$, en la sucesión exacta anterior puede ponerse un $0$ a la izquierda por~\Ref{II.4.10}~(iv).}
{$\Omega_{\kres(x)/k}^1=0$, pues $\kres(x)$ es étale sobre~$k$. Por otra parte, (ii) implica (iv), ya que, como $X$ y $\Spec(\kres(x))$ son lisos sobre~$k$ en $x$, en la sucesión exacta anterior puede ponerse un $0$ a la izquierda por~\Ref{II.4.10}~(iv).}

\begin{corollaire}
\label{II.5.3}
Sea $x$ un punto de~$X$, de tipo finito sobre~$k$. Si $X$ es liso sobre~$k$ en $x$, entonces $\cal{O}_x$ es regular; el recíproco es cierto si
\ifthenelse{\boolean{orig}}{$\kappa(x)$}{$\kres(x)$}
es una extensión finita separable de~$k$.
\end{corollaire}

En
\marginpar{53}
efecto, el recíproco se deduce de~\Ref{II.5.2}, tomando un sistema regular $(f_i)$ de generadores de~$\goth{m}_x$. (N.B. En lugar de~\Ref{II.5.2}, también puede invocarse el teorema~\Ref{II.4.15}). Se concluye:

\begin{proposition}
\label{II.5.4}
Sea $X$ un preesquema de tipo finito sobre~$k$. Si $X$ es liso sobre~$k$, entonces es regular; el recíproco es cierto si $k$ es perfecto.
\end{proposition}

Para el recíproco, obsérvese que, por~\Ref{II.5.3}, $X$ es liso sobre~$k$ en todo punto cerrado y, por tanto, en todas partes, pues el conjunto de puntos en los que es liso es abierto.

\begin{theoreme}
\label{II.5.5}
Sean $X$ un preesquema de tipo finito sobre~$k$, $x$ un punto de~$X$, $n$ la dimensión de~$X$ en $x$ y $k'$ una extensión perfecta de~$k$. Las condiciones siguientes son equivalentes:
\begin{enumerate}
\item[(i)] $X$ es liso sobre~$k$ en~$x$.
\ifthenelse{\boolean{orig}}
{\item[(ii)] $\Omega_{X/k}^1$ es libre de rango $n$ en~$x$.
\item[(ii~bis)] $\Omega_{X/k}^1$ está engendrado por $n$}
{\item[(ii)] $\it{\Omega}_{X/k}^1$ es libre de rango $n$ en~$x$.
\item[(ii~bis)] $\it{\Omega}_{X/k}^1$ está engendrado por $n$}
elementos en~$x$.
\item[(iii)] $X$ es diferencialmente liso sobre~$k$ en $x$.
\item[(iv)] Existe un entorno abierto $U$ de~$x$ tal que $U\otimes_k k'$ es regular (es decir, los anillos locales de sus puntos son regulares).
\end{enumerate}
\end{theoreme}

Se tiene (i)$\To$(ii) por~\Ref{II.4.3}~(ii), (ii)$\To$(ii~bis) trivialmente y (ii~bis)$\To$(i) gracias a~\Ref{II.5.1}. Como $X$ es plano sobre~$k$, se tiene (i)$\Leftrightarrow$(iii) por~\Ref{II.4.18}. Se tiene (i)$\To$(iv), pues la lisura es invariante por extensión de la base e implica regularidad; y (iv)$\To$(i), porque la proposición~\Ref{II.5.4} muestra que $U\otimes_k k'$ es liso sobre~$k'$ y, por tanto, $U$ es liso sobre~$k$ por~\Ref{II.4.13}.

Tomando como $x$ el punto genérico de~$X$, supuesto irreducible, se obtiene:

\begin{corollaire}
\label{II.5.6}
Sea $K$ un anillo local artiniano obtenido por localización de un álgebra de tipo finito sobre el cuerpo~$k$ (por ejemplo, $K$ puede ser una extensión de tipo finito de~$k$), y sea $n$ el grado de trascendencia
\ifthenelse{\boolean{orig}}{sobre~$K$}{sobre~$k$}
de su cuerpo residual. Son equivalentes las condiciones siguientes:
\begin{enumerate}
\item[(i)] $K$ es una extensión finita separable de una extensión puramente trascendente $k(t_1,\dots,t_n)$ de~$k$.
\item[(ii)]
\marginpar{54}
\ifthenelse{\boolean{orig}}{$\it{\Omega}^1_{X/k}$}{$\Omega^1_{K/k}$}
es un $K$-módulo libre de rango~$n$.
\item[(ii~bis)]
\ifthenelse{\boolean{orig}}{$\it{\Omega}^1_{X/k}$}{$\Omega^1_{K/k}$}
es un $K$-módulo que admite $n$ generadores.
\item[(iii)] La completación $O'$ de~$K\otimes_k K$ para la topología definida por las potencias del ideal de aumentación $K\otimes_k K\to K$ es un álgebra aumentada «regular» sobre~$K$, es decir, isomorfa a un álgebra de series formales sobre~$K$ (N.B. Si $K$ es un cuerpo, esto equivale a decir que $O'$ es un anillo local regular).
\item[(iv)] $K$ es una extensión separable de~$k$.
\end{enumerate}
\end{corollaire}

En efecto, siempre puede considerarse $K$ como el anillo local del punto genérico de un esquema irreducible de tipo finito~$X$ sobre~$k$, y las condiciones consideradas son las condiciones del mismo nombre de~\Ref{II.5.5}, tomando en~(iv) como~$k'$ una extensión algebraicamente cerrada de~$k$. Solo la implicación «$K$ separable sobre~$k$» $\To$ «$X$ liso sobre~$k$ en~$x$» requiere una demostración. Por~\Ref{II.4.13}, uno se reduce inmediatamente al caso en que el cuerpo base es~$k'$, por tanto algebraicamente cerrado, y existe entonces un punto $a$ de~$X$ racional sobre~$k$. Pero $X$ es liso sobre~$k$ en~$a$ por~\Ref{II.5.4}; con mayor razón, es liso sobre~$k$ en el punto genérico~$x$, c.q.d.\footnote{\Cf los Errata al final del presente Exposé II (p\ptbl \pageref{II.fin.errata}).}


Obsérvese que, cuando $K$ es una extensión de tipo finito de~$k$, la equivalencia entre
\ifthenelse{\boolean{orig}}
{(i) (ii) (ii~bis) (iv)}
{(i), (ii), (ii~bis), (iv)}
es bien conocida, pero no hemos utilizado ninguna de
\ifthenelse{\boolean{orig}}{sus}{estas}
\ifthenelse{\boolean{orig}}{equivalencia}{equivalencias}
ya conocidas. Naturalmente, la proposición~\Ref{II.5.1} contiene como caso particular el hecho de que una sucesión de elementos~$x_i$ ($1\leq i\leq n$) es
\ifthenelse{\boolean{orig}}{un}{una}
«base de trascendencia separante» de~$K$ sobre~$k$ si y solo si los~$dx_i$ forman una base del $K$-módulo $\Omega^1_{K/k}$, hecho también bien conocido.

\begin{corollaire}
\label{II.5.7}
Sea $X$ un preesquema de tipo finito sobre un cuerpo~$k$. Para que $X$ sea liso sobre~$k$, es necesario y suficiente que
\ifthenelse{\boolean{orig}}{$\Omega^1_{X/k}$}{$\it{\Omega}^1_{X/k}$}
sea localmente libre y que los anillos locales de los puntos genéricos de las componentes irreducibles de~$X$ sean extensiones separables de~$k$ (esta última condición se cumple automáticamente si $k$ es perfecto y $X$ es reducido).
\end{corollaire}

Puede suponerse que $X$ es conexo; sea $n$ el rango de
\ifthenelse{\boolean{orig}}{$\Omega^1_{X/k}$}{$\it{\Omega}^1_{X/k}$},
supuesto localmente libre. Por la hipótesis y~\Ref{II.5.6}, este es también el grado de trascendencia de las extensiones
\marginpar{55}
de~$k$ definidas por los anillos locales de los puntos genéricos de~$X$; por tanto, todas las componentes irreducibles de~$X$ tienen dimensión~$n$. Se concluye entonces por~\Ref{II.5.5}.

Debe advertirse que, si $K$ es una extensión finita (no necesariamente separable) de~$k$, entonces $\Omega^1_{K/k}$ es un
\ifthenelse{\boolean{orig}}{$k$-módulo}{$K$-módulo}
libre; por tanto, poniendo $X=\Spec(K)$,
\ifthenelse{\boolean{orig}}{$\Omega^1_{X/k}$}{$\it{\Omega}^1_{X/k}$}
es un haz localmente libre y $X$ es reducido, sin que $X$ sea necesariamente liso sobre~$k$. Extendiendo entonces los escalares a la clausura algebraica de~$k$, se obtiene un ejemplo análogo en el que $k$ es algebraicamente cerrado, pero $X$, en cambio, no es reducido.

\begin{corollaire}
\label{II.5.8}
Sean $X$ un preesquema de tipo finito sobre el cuerpo~$k$, $x$ un punto de~$X$, $n$ la dimensión de~$X$ en~$x$ y $p$ la dimensión
\ifthenelse{\boolean{orig}}{de~$\cal{O}_x$}{de~$\cal{O}_x$,}
es decir, la codimensión en $X$ de la adherencia $Y$ de~$x$ en~$X$; así, $n-p$ es el grado de trascendencia de
\ifthenelse{\boolean{orig}}{$\kappa(x)$}{$\kres(x)$}
sobre~$k$. Sean $f_i$ ($1\leq i\leq n$) elementos de~$\cal{O}_x$ tales que $f_i\in\goth{m}_x$ para $1\leq i\leq p$. Las condiciones siguientes son equivalentes:
\begin{enumerate}
\item[(i)] El germen de morfismo en~$x$
\[
X\to\Spec\bigl(k[t_1,\dots,t_n]\bigr)
\]
definido por los $f_i$ es étale en~$x$.
\item[(ii)] Los $f_i$ ($1\leq i\leq p$)
\ifthenelse{\boolean{orig}}{engendran~$\goth{m}_x$}{engendran~$\goth{m}_x$,}
es decir, forman un sistema regular de parámetros de~$\cal{O}_x$, y las clases en
\ifthenelse{\boolean{orig}}
{$\kappa(x)$ de los~$f_j$ ($p+1\leq j\leq n$) forman una base de trascendencia separante (es decir, los~$d\overline{f}_j$ ($p+1\leq j\leq n$) forman una base de~$\Omega^1_{\kappa(x)/k}$, o bien engendran~$\Omega^1_{\kappa(x)/k}$).}
{$\kres(x)$ de los~$f_j$ ($p+1\leq j\leq n$) forman una base de trascendencia separante (es decir, los~$d\overline{f}_j$ ($p+1\leq j\leq n$) forman una base de~$\Omega^1_{\kres(x)/k}$, o bien engendran~$\Omega^1_{\kres(x)/k}$).}
\end{enumerate}
\end{corollaire}

Supongamos que se cumple~(i). Se sigue que los $df_i(x)$ forman una base de~$\it{\Omega}^1_{X/k}(x)$
\ifthenelse{\boolean{orig}}{\eqref{II.4.8}}{por~\Ref{II.4.8}};
por tanto, sus imágenes $d\overline{f}_i(x)$ en
\ifthenelse{\boolean{orig}}{$\Omega^1_{\kappa(x)/k}$}{$\Omega^1_{\kres(x)/k}$}
engendran este espacio vectorial
\ifthenelse{\boolean{orig}}{sobre~$k$}{sobre~$\kres(x)$}.
Como los $\overline{f}_i$ son nulos para $1\leq i\leq p$, basta tomar los $d\overline{f}_i(x)$ con $p+1\leq i\leq n$. Puesto que el grado de trascendencia de
\ifthenelse{\boolean{orig}}{$\kappa(x)$}{$\kres(x)$}
sobre~$k$ es $n-p$, se deduce entonces del corolario~\Ref{II.5.6}, criterio
\ifthenelse{\boolean{orig}}{(iii)}{(ii~bis)},
aplicado a
\ifthenelse{\boolean{orig}}{$K=\kappa(x)$}{$K=\kres(x)$},
que $Y$ es liso sobre~$k$ en su punto genérico~$x$ y que los $d\overline{f}_i(x)$ ($p+1\leq i\leq n$) forman una \emph{base} de
\ifthenelse{\boolean{orig}}{$\Omega^1_{\kappa(x)/k}$ sobre~$\kappa(x)$}{$\Omega^1_{\kres(x)/k}$ sobre~$\kres(x)$}.
Por consiguiente, se cumple la condición~(ii) de~\Ref{II.4.9}, luego también la condición~(iii) y, en particular, los $f_i$ ($1\leq i\leq p$) forman un sistema de generadores de~$\goth{m}_x$. Como
\marginpar{56}
$\cal{O}_x$ tiene dimensión~$p$, forman un sistema regular de parámetros en~$x$. Esto prueba~(ii).

Supongamos que se cumple~(ii). En virtud de la sucesión exacta
\ifthenelse{\boolean{orig}}{\eqref{II.5.1}}{\eqref{eq:II.5.1}},
se sigue que los $df_i(x)$
\ifthenelse{\boolean{orig}}{engendran~$\Omega^1_{X/k}$}{engendran~$\it{\Omega}^1_{X/k}(x)$},
de donde se obtiene~(i) por la prop\ptbl \Ref{II.5.1}.

\begin{corollaire}
\label{II.5.9}
Sean $X$ un preesquema de tipo finito sobre el cuerpo~$k$, $x$ un punto de~$X$, $n$ la dimensión de~$X$ en~$x$ y $p$ la dimensión
\ifthenelse{\boolean{orig}}{de~$\cal{O}_x$}{de~$\cal{O}_x$,}
es decir, la codimensión de la adherencia~$Y$ de~$x$ en~$X$; así, $n-p$ es el grado de trascendencia de
\ifthenelse{\boolean{orig}}{$\kappa(x)$}{$\kres(x)$}
sobre~$k$. Son equivalentes las condiciones siguientes:
\begin{enumerate}
\item[(i)] $\cal{O}_x$ es regular y
\ifthenelse{\boolean{orig}}{$\kappa(x)$}{$\kres(x)$}
es una extensión separable de~$k$.
\item[(ii)] $X$ es liso sobre~$k$ en~$x$, y el homomorfismo canónico
\[
\ifthenelse{\boolean{orig}}
{\goth{m}_x/\goth{m}_x^2\to
\Omega^1_{\cal{O}_x/k}\otimes_{\cal{O}_x}
\kappa(x)=\it{\Omega}^1_{X/k}(x)}
{\goth{m}_x/\goth{m}_x^2\to
\Omega^1_{\cal{O}_x/k}\otimes_{\cal{O}_x}
\kres(x)=\it{\Omega}^1_{X/k}(x)}
\]
es inyectivo.
\item[(iii)] Existen $f_i\in\cal{O}_x$ ($1\leq i\leq n$), con $f_i\in\goth{m}_x$ para $1\leq i\leq p$, tales que el germen en~$x$ del morfismo de~$X$ a~$\Spec\bigl(k[t_1,\dots,t_n]\bigr)$ definido por los~$f_i$ es étale en~$x$ (es decir, por~\Ref{II.5.1}, tales que los~$df_i(x)$ engendran~$\it{\Omega}^1_{X/k}(x)$).
\item[(iv)] Existen $f_i\in\cal{O}_x$ ($1\leq i\leq n$) tales que los $f_i$ ($1\leq i\leq p$) engendran~$\goth{m}_x$ y los $df_j(x)$ ($p+1\leq j\leq n$) engendran
\ifthenelse{\boolean{orig}}{$\Omega^1_{\kappa(x)/k}$ sobre~$\kappa(x)$}{$\Omega^1_{\kres(x)/k}$ sobre~$\kres(x)$}.
\end{enumerate}
\end{corollaire}

La equivalencia entre~(iii) y~(iv) se deduce del corolario~\Ref{II.5.8}; estas condiciones equivalen asimismo, por~\Ref{II.4.9}, a que $X$ sea liso sobre~$k$ en~$x$ y a que la condición~(ii)
\ifthenelse{\boolean{orig}}{de~\Ref{II.4.10},}{de~\Ref{II.4.10}}
se cumpla. Por tanto, equivalen a que $X$ sea liso sobre~$k$ en~$x$ y a que se cumpla la condición~(iv) de~\Ref{II.4.10}, es decir, a~\Ref{II.5.9}~(ii). O bien, a que $X$ sea liso sobre~$k$ en~$x$ y a que se cumpla la condición~(i) de~\Ref{II.4.10}, que aquí significa simplemente que
\ifthenelse{\boolean{orig}}{$\kappa(x)$}{$\kres(x)$}
es separable sobre~$k$. Esto implica~\Ref{II.5.9}~(i); queda probar que~\Ref{II.5.9}~(i) lo implica, es decir, probar el

\begin{corollaire}
\label{II.5.10}
Sea~$x$ un punto
\ifthenelse{\boolean{orig}}{de una preesquema}{de un preesquema~$X$}
de tipo finito sobre el cuerpo~$k$, tal que
\ifthenelse{\boolean{orig}}{$\kappa(x)$}{$\kres(x)$}
sea separable sobre~$k$. Para que $X$ sea liso sobre~$k$ en~$x$, es necesario y suficiente que sea regular en~$x$ (es decir, que el anillo local~$\cal{O}_x$ sea regular).
\end{corollaire}

En efecto, si así ocurre, es evidente que pueden encontrarse $f_i\in\cal{O}_x$ ($1\leq i\leq n$) que satisfagan la condición~\Ref{II.5.9}~(iv).

\subsection*{Errata}
\label{II.fin.errata}
En
\marginpar{57}
este número, en la demostración de~\Ref{II.5.6}, se ha utilizado el hecho de que todo esquema reducido no vacío de tipo finito sobre un cuerpo algebraicamente cerrado admite al menos un punto regular (por tanto liso), hecho que suele demostrarse por métodos diferenciales (mediante el teorema de Zariski según el cual el conjunto de puntos regulares de~$X$ es abierto). Si se quiere evitar un círculo vicioso, hay que demostrar que, si $K/k$ es una extensión separable de tipo finito y los $f_i\in K$ son tales que los $d_{K/k}f_i$ forman una base de~$\Omega^1_{K/k}$ ($1\leq i\leq n$), entonces $n$ es el grado de trascendencia de~$K$
\ifthenelse{\boolean{orig}}{sobre~$k$}{sobre~$k$,}
es decir, los $f_i$ son algebraicamente independientes. La demostración de este hecho mediante el criterio de
\ifthenelse{\boolean{orig}}{Mac-Lane}{Mac\kern2pt Lane}
es bien conocida; \cf Bourbaki, Álgebra, cap\ptbl V, §~9, teorema~2: se toma un polinomio~$g\in k[t_1,\dots,t_n]$ de grado mínimo tal que~$g(f_1,\dots,f_n)=0$; se tiene entonces
\[
\sum\frac{dg}{d t_i}(f_1,\dots,f_n)df_i=0,
\]
de donde se deduce (puesto que los $df_i$ forman una base de~$\Omega^1_{K/k}$) que los $dg/dt_i$ se anulan en~$(f_1,\dots,f_n)$ y, por la minimalidad de~$g$, son nulos. Por tanto, si $k$ es de característica~$0$, se tiene $g=0$; y si $k$ es de característica~$p\neq0$, se tiene $g=h(t_1^p,\dots,t_n^p)$. Utilizando el criterio de
\ifthenelse{\boolean{orig}}{Mac-Lane}{Mac\kern2pt Lane},
se ve que el polinomio~$h\in k[t_1,\dots,t_n]$ también se anula en~$(f_1,\dots,f_n)$, de donde nuevamente $g=0$ por la minimalidad de~$g$.


\chapter{Morfismos lisos: propiedades de prolongación}
\label{III}
\marginpar{58}

\section{Homomorfismos formalmente lisos}
\label{III.1}

En~\Ref{II} nos limitamos a considerar homomorfismos de tipo finito y, por consiguiente, en el caso de homomorfismos locales~$A\to B$ de anillos locales, al caso en que~$B$ es isomorfa a una localización de una $A$-álgebra de tipo finito. Este caso es insuficiente para diversas aplicaciones, en particular en geometría formal o en geometría analítica. Por ejemplo, el anillo de series formales~$B=A[[t_1,\dots,t_n]]$ tiene (desde el punto de vista de la geometría formal) las propiedades de un álgebra lisa sobre~$A$. En geometría analítica, lo mismo ocurre con el anillo local de un punto~$(y,z)$ de un producto $Y\times\CC^n$, considerado como álgebra sobre el anillo local de~$y$; además, la completación de esta álgebra es isomorfa al álgebra de series formales en~$n$ indeterminadas sobre la completación del anillo base
\ifthenelse{\boolean{orig}}{$\cal{O}_x$}{$\cal{O}_y$}.
Esto conduce a formular la definición siguiente.

\begin{definition}
\label{III.1.1}
Sea~$u\colon A\to B$ un homomorfismo local de anillos locales (noetherianos, como recordamos). Se supone que
\ifthenelse{\boolean{orig}}
{$\kappa(B)$ es finito sobre~$\kappa(A)$}
{$\kres(B)$ es finito sobre~$\kres(A)$}.
Se dice que~$u$ es un \emph{homomorfismo formalmente liso},
\index{formalmente liso (homomorfismo, álgebra)|hyperpage}%
\index{liso (formalmente)|hyperpage}%
o que el álgebra~$B$ es \emph{formalmente lisa sobre}~$A$, si existe una $\overline{A}$-álgebra local
\ifthenelse{\boolean{orig}}{finita,}{finita~$A'$,}
libre sobre~$\overline{A}$, tal que las componentes locales del anillo semilocal~$\overline{B}\otimes_{\overline{A}}A'=B'$ sean isomorfas, como $A'$-álgebras, a álgebras de series formales sobre~$A'$\kern1pt\footnote{Para una definición más general y conceptual, motivada por~\Ref{III.2.1} más adelante, \cf EGA~$\mathrm{0}_{\mathrm{IV}}$~19.3.1.}.
\end{definition}
(Se denotan por~$\overline{A}$ y $\overline{B}$ las completaciones de~$A$ y~$B$). Como~$B'$ es finito y libre sobre~$\overline{B}$, es efectivamente un anillo semilocal, producto directo de anillos locales completos, cada uno de los cuales sigue siendo un módulo libre sobre~$\overline{B}$ y tiene, por tanto, la misma dimensión que~$\overline{B}$, luego que~$B$. Se sigue que el número de variables~$t_i$ en los anillos de series formales considerados en~\Ref{III.1.1} es igual a $\dim\overline{B}-\dim\overline{A}=\dim B-\dim A$ y, en particular, independiente de la componente local
\marginpar{59}
elegida. Se ve de inmediato que también es la dimensión del anillo
\ifthenelse{\boolean{orig}}
{$B\otimes k=B/\goth{m}B$}
{$B\otimes_A k=B/\goth{m}B$},
donde~$k=A/\goth{m}$ es el cuerpo residual de~$A$; se la llamará \emph{dimensión relativa de}~$B$
\index{dimensión relativa|hyperpage}%
\emph{respecto de}~$A$.

\begin{remarques}
\label{III.1.2}
Es evidente que la definición~\Ref{III.1.1} solo depende del homomorfismo entre las completaciones $\overline{A}\to\overline{B}$ deducido de~$A\to B$, lo que justifica en cierta medida la terminología. Nos retractamos aquí de la definición~\Ref{I.3.2}~b) y~\Ref{I.4.1}~b), que puede inducir a error, y preferimos decir «formalmente no ramificado»
\index{formalmente no ramificado (respectivamente neto)|hyperpage}%
y «formalmente étale»
\index{formalmente étale|hyperpage}%
en los casos considerados en esas definiciones, reservando la terminología «no ramificado» y «étale» para el caso en que $B$ es una localización de una $A$-álgebra de tipo finito\footnote{O, mejor, «esencialmente no ramificado», respectivamente «esencialmente étale»; compárese EGA IV 18.6.1.}.
\index{esencialmente étale|hyperpage}%
\index{esencialmente no ramificado (respectivamente neto)|hyperpage}%
El lector comprobará enseguida que «formalmente étale» equivale a «formalmente liso y cuasifinito». Señalemos, por último, que existe una definición razonable de «formalmente liso» sin ninguna hipótesis previa sobre la extensión residual
\ifthenelse{\boolean{orig}}
{$\kappa(B)/\kappa(A)$}
{$\kres(B)/\kres(A)$}
(que aquí se supone finita), la cual abarca, entre otros, los homomorfismos locales $A\to B$ tales que $B$ sea \emph{plano} sobre~$A$ y $B/\goth{m}B$ sea una \emph{extensión separable} de~$A/\goth{m}=k$ (no necesariamente de tipo finito); por ejemplo, un $p$-anillo de Cohen es formalmente
\ifthenelse{\boolean{orig}}{simple}{liso}
sobre el anillo de los enteros $p$-ádicos. En este caso general, es la propiedad de levantamiento de los homomorfismos (compárese~\Ref{III.2.1}) la que debe convertirse en definición. Para las aplicaciones que tenemos en mente, bastará el caso tratado en la definición~\Ref{III.1.1}; en adelante, en «formalmente liso» se sobrentenderá «con extensión residual finita».
\end{remarques}

\begin{lemme}
\label{III.1.3}
Si $B$ es formalmente lisa sobre~$A$, entonces $B$ es plana sobre~$A$.
\end{lemme}

\ifthenelse{\boolean{orig}}
{\itshape}{}
Como la planitud es invariante por completación, puede suponerse que $A$ y $B$ son completos. Como la planitud es invariante por extensión local plana (y, por tanto, fielmente plana) del anillo base, la definición~\Ref{III.1.1} reduce el problema al caso en que $B$ es un álgebra de series formales sobre~$A$. Pero entonces, como $A$-módulo, $B$ es isomorfo a un producto de $A$-módulos isomorfos a~$A$ y, por tanto (al ser noetheriano el anillo base~$A$), es $A$-plano por ser un producto de $A$-módulos planos.
\ifthenelse{\boolean{orig}}
{\upshape}{}

\medskip
Situémonos en las condiciones de~\Ref{III.1.1}. Como las extensiones
\marginpar{60}
residuales de las componentes locales de~$B'$ sobre~$A'$ son triviales, se sigue que $L\otimes_k k'$ es una $k'$-álgebra artiniana cuyas componentes locales tienen extensiones residuales triviales (donde $L,\,k,\,k'$ son los cuerpos residuales
\ifthenelse{\boolean{orig}}
{de~$A,\,B,\,A'$}
{de~$B,\,A,\,A'$}).
Esta condición, necesaria para que la extensión finita libre~$A'$ satisfaga la condición enunciada en~\Ref{III.1.1}, es también suficiente, como resulta inmediatamente de~\Ref{III.1.4}~(i) y~\Ref{III.1.5} más adelante.

\begin{proposition}
\label{III.1.4}
Sea $A\to B$ un homomorfismo local de anillos locales, con extensión residual finita
\ifthenelse{\boolean{orig}}{\ignorespaces}{y}
sea $A'$ una $A$-álgebra local finita sobre~$A$, de modo que $B'=B\otimes_A A'$ es finita sobre~$B$ y, por tanto, un anillo semilocal asimismo noetheriano. \textup{(i)} Si $B$ es formalmente lisa sobre~$A$, las localizaciones de
\ifthenelse{\boolean{orig}}{$B$}{$B'$}
en sus ideales maximales son formalmente lisas sobre~$A'$. \textup{(ii)} El recíproco es cierto si $A'$ es libre sobre~$A$.
\end{proposition}

Se reduce de inmediato al caso en que $A$ y $B$ son completos.

(i) Sea $A''$ una extensión local finita libre de~$A$ tal que las componentes locales de~$B''=B\otimes_A A''$ sean álgebras de series formales sobre~$A''$. Mediante la extensión de escalares $A''\to A''\otimes_A A'\to A'''$, donde $A'''$ es una componente local de~$A''\otimes_A A'$, se ve que las componentes locales de $B''\otimes_{A''}A'''=B\otimes_A A'''$ son álgebras de series formales sobre~$A'''$. Por otra parte, también se tiene
\ifthenelse{\boolean{orig}}
{$B\otimes_AA'''=(B\otimes_A A')\otimes_{A'}A'''=B\otimes_{A'}A'''$}
{$B\otimes_AA'''=(B\otimes_A A')\otimes_{A'}A'''=B'\otimes_{A'}A'''$};
además, como $A''$ es libre sobre~$A$, $A''\otimes_A A'$ es libre sobre~$A'$ y, por consiguiente, lo mismo ocurre con~$A'''$, que es un factor directo suyo. Esto prueba que $B'$ es formalmente lisa sobre~$A'$.

(ii) Sea $A''$ una $A'$-álgebra local finita libre tal que las componentes locales de~$B'\otimes_{A'}A''=B\otimes_A A''$ sean álgebras de series formales sobre~$A''$. Como $A'$ es libre sobre~$A$, también lo es~$A''$; por tanto, $B$ es formalmente \emph{lisa} sobre~$A$.


\begin{proposition}
\label{III.1.5}
Sea $A\to B$ un homomorfismo local de anillos locales, con extensión residual \emph{trivial}. Para que $B$ sea formalmente lisa sobre~$A$, es necesario y suficiente que $\overline{B}$ sea isomorfa a un álgebra de series formales sobre~$\overline{A}$.
\end{proposition}

Solo queda demostrar la necesidad, y puede suponerse que $A$ y $B$ son completos. Sea $\goth{m}$ (respectivamente $\goth{n}$) el ideal maximal de~$A$ (respectivamente de~$B$), y sean $t_1,\dots,t_n$ elementos de~$\goth{n}$ que determinan una base del espacio vectorial
$$
(\goth{n}/\goth{n}^2)/\Im(\goth{m}/\goth{m}^2)=
\goth{n}/(\goth{n}^2+\goth{m}B)
$$
Estos elementos determinan, por tanto, un homomorfismo de $A$-álgebras locales
$$
B_1=A[[t_1,\dots,t_n]]\to B
$$
Probemos
\marginpar{61}
que es un isomorfismo. Basta probar que, para toda potencia $\goth{m}^q$ de~$\goth{m}$, se obtiene un isomorfismo al reducir módulo~$\goth{m}^q$ (pues $B_1$ y $B$ son los límites proyectivos de los correspondientes anillos reducidos módulo~$\goth{m}^q$, al variar~$q$). Como $B$ y $B_1$ son $A$-módulos planos, los graduados asociados a la filtración $\goth{m}$-ádica se obtienen tensorizando con $\gr(A)$, sobre~$k=A/\goth{m}$, los anillos $B_1/\goth{m}B_1$ y, respectivamente, $B/\goth{m}B$; por tanto, se reduce a demostrar que $B_1/\goth{m}B_1\to B/\goth{m}B$ es un isomorfismo. Teniendo en cuenta~\Ref{III.1.3}, el problema queda así reducido al caso en que $A$ es un \emph{cuerpo}~$k$. Por otra parte, si $A'$ es una $A$-álgebra local finita libre tal que $B\otimes_A A'$ sea un álgebra de series formales sobre~$A'$ (N.B. Esta álgebra es local, pues la extensión residual de~$B$ sobre~$A$ es trivial), para probar que $B_1\to B$ es un isomorfismo basta probar que lo es $B_1\otimes_A A'\to B\otimes_A A'$. Esto reduce el problema al caso en que $B$ ya es un álgebra de series formales (habría sido preciso comenzar por esta reducción, antes de pasar al caso de un cuerpo base). Pero entonces $B$ es un anillo local regular con cuerpo de coeficientes~$k$, y es bien sabido (e inmediato al considerar los graduados asociados a las filtraciones $\goth{n}_1$-ádica y $\goth{n}$-ádica sobre~$B_1$ y~$B$) que $B_1\to B$ es un isomorfismo, lo que concluye la demostración.

\begin{corollaire}
\label{III.1.6}
Si $B$ es formalmente lisa sobre~$A$, existe una $A$-álgebra local finita $A'$ tal que las componentes locales de
$\overline{B}\otimes_{\overline{A}}\overline{A'}=\overline{(B\otimes_A A')}$
sean isomorfas a álgebras de series formales sobre~$\overline{A'}$.
\end{corollaire}
En efecto, si $L/k$ es la extensión residual de~$B/A$, se considera una extensión $k'/k$ tal que las extensiones residuales de la $k'$-álgebra $L\otimes_k k'$ sean triviales. Sea $A'$ una $A$-álgebra finita libre tal que $A'/\goth{m}A'=k'$ (se sabe que existe una; por ejemplo, se reduce paso a paso al caso en que $k'/k$ es monógena, y entonces se levantan a~$A$ los coeficientes del polinomio minimal de un generador de~$k'$ sobre~$k$). Es local. Entonces las extensiones residuales de~$B\otimes_A A'$ en sus ideales maximales sobre el cuerpo residual~$k'$ de~$A'$ son triviales, y se concluye mediante~\Ref{III.1.5}.

\begin{corollaire}
\label{III.1.7}
Sea $A\to B$ un homomorfismo local de anillos locales. Para que $B$ sea formalmente lisa sobre~$A$, es necesario y suficiente que $B$ sea plana sobre~$A$ y que $B/\goth{m}B$ sea formalmente lisa sobre~$k=A/\goth{m}$.
\end{corollaire}
Efectuando una extensión local finita libre adecuada $A'$ de~$A$ y usando~\Ref{III.1.4}~(ii), se reduce al caso en que la extensión residual de~$B/A$ es trivial.
\marginpar{62}
Por otra parte, se sabe por~\Ref{III.1.4}~(i) y~\Ref{III.1.3} que las condiciones enunciadas son necesarias. Para la suficiencia, basta observar que la demostración de~\Ref{III.1.5} prueba, bajo las hipótesis formuladas aquí, que $B$ es un álgebra de series formales sobre~$A$ (suponiendo, como es lícito, que $A$ y $B$ son completos).

\begin{remarque}
\label{III.1.8}
No sería difícil desarrollar, para los homomorfismos formalmente lisos, el análogo de todas las propiedades de los morfismos lisos estudiadas en~\Ref{II}. Sin embargo, para las propiedades diferenciales esto exige modificar la definición habitual de las diferenciales de Kähler (\cf \Ref{I.1}), sustituyendo los productos tensoriales ordinarios por productos tensoriales completados. Nos limitamos aquí a evocar estos abismos, pues lo precedente basta para nuestro propósito.

Queda establecer el vínculo entre la lisura formal y la noción de lisura desarrollada en~\Ref{II} (y de la que aún no hemos hecho uso alguno):
\end{remarque}

\begin{proposition}
\label{III.1.9}
Sea $A\to B$ un homomorfismo local, siendo $B$ una localización de una $A$-álgebra de tipo finito. Para que $B$ sea lisa sobre~$A$, es necesario y suficiente que sea formalmente lisa sobre~$A$.
\end{proposition}

Mediante~\Ref{III.1.7} y~\Ref{II.2.1}, se reduce al caso en que $A$ es un cuerpo.

Mediante~\Ref{III.1.4}~(ii) y~\Ref{II.4.13}, una extensión adecuada $k'$ de~$k$ reduce el problema al caso en que la extensión residual de~$B/k$ es trivial. En virtud de
\ifthenelse{\boolean{orig}}
{\Ref{III.1.5} (respectivamente \Ref{II.5.2})}
{\Ref{II.5.2} (respectivamente \Ref{III.1.5})},
$B$ es entonces lisa sobre~$k$ (respectivamente, formalmente \emph{lisa sobre}~$k$) si y solo si $B$ es un anillo local regular (respectivamente, si su completación es un álgebra de series formales sobre~$k$). Ahora bien, es bien sabido que estas dos condiciones son equivalentes (pues la extensión residual es trivial).


\section[Propiedad de levantamiento]{Propiedad de levantamiento característica de los homomorfismos formalmente lisos}
\label{III.2}

\begin{theoreme}
\label{III.2.1}
Sea $A\to B$ un homomorfismo local de anillos locales que define una extensión residual finita. Las condiciones siguientes son equivalentes:
\begin{enumerate}
\item[(i)]
\marginpar{63}
$B$ es formalmente lisa sobre~$A$.
\item[(ii)] Para todo homomorfismo local $A\to C$, donde $C$ es un anillo local \emph{completo}, todo ideal~$J$ de~$C$ contenido en
\ifthenelse{\boolean{orig}}
{$r(C)$}
{$\goth{r}(C)$},
y todo $A$-homomorfismo local $B\to C/J$, existe un $A$-homomorfismo (necesariamente local) $B\to C$ que lo levanta.
\item[(iii)] Para toda $A$-álgebra~$C$ (no necesariamente un anillo noetheriano), todo ideal nilpotente~$J$ de~$C$ y
\ifthenelse{\boolean{orig}}
{todo $A$-homomorfismo continuo $B\to C$}
{todo $A$-homomorfismo continuo $B\to C/J$}
(es decir, que se anula en una potencia
\ifthenelse{\boolean{orig}}
{de~$r(B)$}
{de~$\goth{r}(B)$}),
existe un $A$-homomorfismo $B\to C$ (necesariamente también continuo) que lo levanta.
\item[(iv)] El mismo enunciado que en \textup{(ii)} y \textup{(iii)}, pero siendo~$C$ un anillo local artiniano finito sobre~$A$.
\item[(iv~bis)] Como \textup{(iv)}, pero suponiendo además que~$J$ es de cuadrado nulo.
\end{enumerate}
\end{theoreme}

\begin{remarquestar}
En lo que resta de este Exposé solo utilizaremos la implicación (iv)$\To$(i), o bien (iv~bis)$\To$(i); la implicación directa (i)$\To$(ii) se demostrará por otro método en el número siguiente cuando~$B$ sea una localización de una $A$-álgebra de tipo finito. Recordemos que, en la «buena» teoría de los teoremas de Cohen\footnote{\Cf EGA $\mathrm{0}_{\textup{IV}}$ 19.3, 19.8}, la propiedad~(ii) o~(iii) pasa a ser la definición de los homomorfismos formalmente lisos, mientras que~\Ref{III.1.1} pasa a ser una propiedad característica válida únicamente cuando la extensión residual es finita. Conviene advertir que ninguna de las propiedades~(ii) y~(iii) es más general que la otra. Podría darse una propiedad equivalente que englobase a ambas introduciendo un anillo linealmente topologizado~$C$, \emph{separado} y \emph{completo}, un ideal \emph{cerrado topológicamente nilpotente} de~$C$ y un homomorfismo continuo $A\to C$ (que hace, por tanto, de~$C$ una $A$-álgebra topológica); dejamos esta modificación al lector.
\end{remarquestar}

\subsubsection*{Demostración de~\Ref{III.2.1}}
Demostraremos (i)$\To$(iii)$\To$(ii), y después (iv)$\To$(i). Como (ii)$\To$(iv) es trivial y la equivalencia entre~(iv) y~(iv~bis) se obtiene por una inducción inmediata sobre el entero~$n$ tal que $J^n=0$, esto concluirá la demostración.

(i)$\To$(iii). Una inducción inmediata nos reduce al caso en que $J^2=0$.
\ifthenelse{\boolean{orig}}
{Como $C$ es finita sobre~$A$, existe una potencia $\goth{m}^q$ del ideal maximal de~$A$ que anula~$C$.}
{Como $u$ es continua y $J^2=0$, existe una potencia $\goth{m}^q$ del ideal maximal de~$A$ cuya imagen en~$C$ es nula.}
Dividiendo por~$\goth{m}^q$ y observando que $B/\goth{m}^qB$ sigue siendo formalmente lisa
\marginpar{64}
sobre $A/\goth{m}^q$ por~\Ref{III.1.4}~(i), puede suponerse que~$A$ es artiniano.
Como $B$ es plana sobre~$A$ por~\Ref{III.1.3}, $B$ \emph{es libre sobre}~$A$, pues~$A$ es artiniano. Existe, por tanto, un \emph{homomorfismo de} $A$-\emph{módulos}
$$
w\colon B\to C
$$
que levanta el homomorfismo dado $u\colon B\to C/J$. Pongamos
$$
f(x,y)=w(xy)-w(x)w(y)\qquad(x,y\in B).
$$
Entonces $f(x,y)\in J$ y, por tanto, $f$ es una aplicación $A$-bilineal de~$B\times B$ en~$J$. Para que exista un levantamiento $v\colon B\to C$ de~$u$ que sea un homomorfismo de álgebras, es necesario y suficiente que exista una aplicación $A$-lineal $g\colon B\to J$ tal que $v=w+g$ sea un homomorfismo de álgebras, lo que se escribe
\begin{align*}
g(1)&=1-w(1),\\
\ifthenelse{\boolean{orig}}
{g(x,y)}
{g(xy)}-u(x)g(y)-u(y)g(x)&=-f(x,y)\qquad(\text{para }x,y\in B).
\end{align*}
Este es un sistema de ecuaciones \emph{lineales} en $\Hom_A(B,J)$, con segundos miembros en~$J$; por tanto, tiene solución si y solo si la tiene el sistema correspondiente en $\Hom_A(B,J)\otimes_A A'$, con segundos miembros en $J'=J\otimes_A A'$, donde~$A'$ designa un álgebra fielmente plana sobre~$A$. Sea ahora~$A'$ una $A$-álgebra local finita y libre tal que $B'=B\otimes_A A'$ sea un álgebra de series formales sobre~$A'$ (N.B. en nuestra demostración puede suponerse que~$A$ y~$B$ son completos, como se comprueba de inmediato). Como $A'$ es libre de tipo finito sobre~$A$, se tiene
$$
\Hom_A(B,J)\otimes_A A'=\Hom_{A'}(B',J'),
$$
y se comprueba que el sistema de ecuaciones obtenido en $\Hom_{A'}(B',J')$ es el que determina los homomorfismos de $A'$-álgebras $B'\to C'=C\otimes_A A'$ que levantan el homomorfismo $u'\colon B'\to C'/J'$ deducido de~$u$ por extensión de escalares, «corrigiendo» mediante un homomorfismo de $A'$-módulos $g'\colon B'\to J'$ el homomorfismo de $A'$-módulos $w'\colon B'\to C'$ deducido de~$w$ por extensión de escalares. (Obsérvese que~$B$ genera~$B'$ como $A'$-módulo). Quedamos así reducidos a demostrar~(iii) cuando~$B$ es un \emph{álgebra de series formales} sobre~$A$, $B=A[[t_1,\dots,t_n]]$. Levantemos de manera arbitraria las imágenes de los~$t_i$ en~$C/J$ a elementos~$z_i$ de~$C$. Como los $z_i$ módulo~$J$ son nilpotentes (pues $u\colon B\to C/J$ es continua), también lo son los~$z_i$ (ya que~$J$ es nilpotente); por tanto, los~$z_i$ definen un homomorfismo continuo de $A$-álgebras topológicas de~$B$ en~$C$ provisto de la topología discreta, que evidentemente levanta~$u$, como se quería demostrar.


(iii)$\To$(ii).
\marginpar{65}
Sea~$\goth{n}$ el ideal maximal de~$C$ y, para todo entero $q>0$, sea
$$
C_q=C/\goth{n}^q,\qquad J_q=(J+\goth{n}^q)/\goth{n}^q.
$$
Así, $C_q/J_q$ se identifica con un álgebra cociente de~$C/J$; por otra parte, el homomorfismo compuesto $u_q\colon B\to C/J\to C_q/J_q$ es continuo de~$B$ en el anillo discreto~$C_q/J_q$, y~$J_q$ es un ideal nilpotente de~$C_q$. Se construyen sucesivamente $A$-homomorfismos
$$
v_q\colon B\to C_q
$$
tales que (a)~$v_q$ levanta~$u_q$ y (b)~$v_q$ levanta~$v_{q-1}$. La posibilidad de la inducción se comprueba fácilmente, pues, como
$$
u_q\colon B\to C/(J+\goth{n}^q)\qquad\text{y}\qquad
v_{q-1}\colon B\to C/\goth{n}^{q-1}
$$
definen el mismo homomorfismo
$$
B\to C/((J+\goth{n}^q)+\goth{n}^{q-1})
=C/(J+\goth{n}^{q-1})=C_{q-1}/J_{q-1},
$$
a saber, $u_{q-1}$, determinan un homomorfismo
$$
B\to C/J_q'\qquad\text{donde}\qquad
J_q'=(J+\goth{n}^q)\cap\goth{n}^{q-1}\supset\goth{n}^q
$$
(del cual ambos proceden por reducción). Se reduce, pues, el problema a levantar un homomorfismo $B\to C/J_q'$ de~$B$ en un cociente de~$C_q$ por un ideal $J_q'/\goth{n}^q$ contenido en~$J_q$ y, por tanto, nilpotente; ello es posible por la hipótesis~(iii).

Hecho esto, los~$v_q$ definen un homomorfismo de~$B$ en el límite proyectivo~$C$ de los~$C_q$. Como~$J$ es cerrado, $J$ es el límite proyectivo de los~$J_q$; por consiguiente, $v$ levanta~$u$, como se quería demostrar.

(iv)$\To$(i). Se comprueba primero de inmediato que, si se verifica~(iv), sigue verificándose para los componentes locales de $B\otimes_A A'$ sobre~$A'$, cuando~$A'$ es una $A$-álgebra local finita. Tomando~$A'$ libre sobre~$A$ y tal que las extensiones residuales de~$B'$ sobre~$A'$ sean triviales, se reduce el problema, gracias a~\Ref{III.1.4}~(ii), al caso en que la extensión residual de~$B$ sobre~$A$ es trivial. Demostraremos entonces el resultado algo más preciso siguiente.

\begin{corollaire}
\label{III.2.2}
Bajo las condiciones de~\Ref{III.2.1}, supongamos además que la extensión residual de~$B$ sobre~$A$ es trivial. Entonces las condiciones equivalentes
\marginpar{66}
de~\Ref{III.2.1} son también equivalentes a las dos condiciones siguientes (suponiendo en \textup{(v)} que~$A$ y~$B$ son completos):
\begin{enumerate}
\item[(iv~ter)] Como \textup{(iv)}, pero restringiendo el anillo local artiniano~$C$, finito sobre~$A$, a tener una extensión residual trivial (y, además, si se desea, suponiendo que el ideal~$J$ es de cuadrado nulo).
\item[(v)] Existe un $A$-homomorfismo local (donde $n=\dim \goth{n}/(\goth{n}^2+\goth{m}B)$)
$$
u\colon B\to B_1=A[[t_1,\dots,t_n]]
$$
que induce un \emph{isomorfismo}
$$
\goth{n}/(\goth{n}^2+\goth{m}B)\isomto
\goth{n}_1/(\goth{n}_1^2+\goth{m}B_1),
$$
donde~$\goth{n}$ (respectivamente $\goth{n}_1$) es el ideal maximal de~$B$ (respectivamente de~$B_1$), y~$\goth{m}$ el de~$A$.
\end{enumerate}
\end{corollaire}

\subsubsection*{Demostración}
Como (iv~bis) implica evidentemente (iv~ter) ---dejando de lado la broma del ideal de cuadrado nulo---, bastará demostrar (iv~ter)$\To$(v)$\To$(i).

(iv~ter)$\To$(v). Elijamos una base $a_1,\dots,a_n$ de $\goth{n}/(\goth{n}^2+\goth{m}B)$, lo que define un homomorfismo local de $A$-álgebras
$$
B\to B_1/(\goth{n}_1^2+\goth{m}B_1)
=k[t_1,\dots,t_n]/(t_1,\dots,t_n)^2,
$$
que, en virtud de~(iv~ter), puede levantarse sucesivamente a homomorfismos de $A$-álgebras de~$B$ en $B_1/\goth{n}_1^2$, $B_1/\goth{n}_1^3$, etc.; pasando al límite proyectivo se obtiene el homomorfismo $B\to B_1$ con la propiedad deseada.

(v)$\To$(i). Como, en el diagrama conmutativo
$$
\ifthenelse{\boolean{orig}}
{\xymatrix{ \goth{m}/\goth{m}^2 \ar[r] & \goth{n}/\goth{n}^2 \ar[r]
\ar[d] & \goth{n}/(\goth{n}^2+\goth{m}B) \ar[r] \ar[d] & 0 \\
\goth{m}/\goth{m}^2 \ar[r] & \goth{n}_1/\goth{n}_1^2 \ar[r] &
\goth{n}_1/(\goth{n}_1^2+\goth{m}B_1) \ar[r] & 0}}
{\xymatrix{ \goth{m}/\goth{m}^2 \ar[r] \ar[d] & \goth{n}/\goth{n}^2
\ar[r] \ar[d] & \goth{n}/(\goth{n}^2+\goth{m}B) \ar[r] \ar[d] & 0 \\
\goth{m}/\goth{m}^2 \ar[r] & \goth{n}_1/\goth{n}_1^2 \ar[r] &
\goth{n}_1/(\goth{n}_1^2+\goth{m}B_1) \ar[r] & 0}}
$$
las dos filas son exactas y las flechas verticales extremas son sobreyectivas, la flecha intermedia es sobreyectiva; de ello se sigue (al ser~$B$ completo) que $B\to B_1$ es sobreyectivo. Sean~$x_i$ ($1\leq i\leq n$) elementos de~$B$ que levantan los~$t_i$. Definen un homomorfismo de $A$-álgebras $B_1\to B$, que será sobreyectivo por la misma razón que~$u$ y cuya composición con~$u$ es la identidad por construcción. Por tanto, $B_1\to B$ es también inyectivo y, en consecuencia, un isomorfismo. Se obtiene así el corolario siguiente.

\begin{corollaire}
\label{III.2.3}
Bajo
\marginpar{67}
las condiciones de~\Ref{III.2.2}~\textup{(v)}, $u$ es necesariamente un isomorfismo.
\end{corollaire}

Esto concluye la demostración de que~$B$ es formalmente lisa sobre~$A$. Al mismo tiempo hemos recuperado~\Ref{III.1.5} (aunque no hay gran mérito en ello).


\section[Prolongación infinitesimal local de los morfismos]{Prolongación infinitesimal local de los morfismos en un $S$-esquema liso}
\label{III.3}

\begin{theoreme}
\label{III.3.1}
Sea $f\colon X\to Y$ un morfismo localmente de tipo finito. Las condiciones siguientes son equivalentes:
\begin{enumerate}
\item[(i)] $f$ es liso.
\item[(ii)] Para todo preesquema~$Y'$ sobre~$Y$, todo subpreesquema cerrado~$Y'_0$ de~$Y'$ con el mismo espacio subyacente que~$Y'$, todo $Y$-morfismo $g_0\colon Y'_0\to X$ y todo $z\in Y'_0$, existen un entorno abierto~$U$ de~$z$ en~$Y'$ y una prolongación~$g$ de $g_0|_{Y'_0\cap U}$ a un $Y$-morfismo $U\to X$.
\item[(ii~bis)] Para $Y'$, $Y'_0$ y~$z$ como en \textup{(ii)}, poniendo $X'=X\times_Y Y'$ y $X'_0=X\times_Y Y'_0$, toda sección de~$X'_0$ sobre~$Y'_0$ se prolonga a una sección de~$X'$ sobre un entorno abierto~$U$ de~$z$.
\item[(iii)] Para todo $Y$-esquema~$Y'$ que sea el espectro de un anillo local artiniano finito sobre algún~$\cal{O}_y$ ($y\in Y$), todo subpreesquema cerrado no vacío~$Y'_0$ de~$Y'$ y todo $Y$-morfismo $g_0\colon Y'_0\to X$, existe un $Y$-morfismo $g\colon Y'\to X$ que prolonga~$g_0$.
\item[(iii~bis)] Para $Y'$, $Y'_0$ como en \textup{(iii)}, poniendo $X'=X\times_Y Y'$ y $X'_0=X\times_Y Y'_0$, toda sección de~$X'_0$ sobre~$Y'_0$ se prolonga a una sección de~$X'$ sobre~$Y'$.
\end{enumerate}
\end{theoreme}

\subsubsection*{Demostración}
La equivalencia entre~(ii) y~(ii~bis), por una parte, y entre~(iii) y~(iii~bis), por otra, es trivial, al igual que la implicación (ii)$\To$(iii). Queda, pues, demostrar (i)$\To$(ii) y (iii)$\To$(i).

(i)$\To$(ii). Sea $x=g_0(z)$. Reemplazando~$X$ por un entorno abierto conveniente de~$x$ y~$Y'$ por el preesquema inducido en la imagen inversa de ese abierto por~$g_0$, puede suponerse que~$X$ es étale sobre~$Y[t_1,\dots,t_n]$. Consideremos el $Y$-morfismo compuesto $Y'_0\to X\to Y[t_1,\dots,t_n]$; está definido por~$n$ secciones del haz~$\cal{O}_{Y'_0}$, que pueden prolongarse, en un entorno de~$z$, a secciones de~$\cal{O}_{Y'}$. Puede suponerse, por tanto, que
\marginpar{68}
el morfismo considerado se ha prolongado a un
\ifthenelse{\boolean{orig}}
{$Y$-morfismo $Y'\to Y'_0$}
{$Y$-morfismo $Y'\to Y[t_1,\dots,t_n]$}.
En virtud de~(\Ref{I.5.6}), existe entonces un único $Y$-morfismo $g\colon Y'\to X$ que levanta el anterior y prolonga al mismo tiempo~$g_0$, como se quería demostrar.

(iii)$\To$(i). Como el conjunto de los puntos en los que~$f$ es liso es abierto, basta demostrar que contiene todo $x\in X$ que sea \emph{cerrado} en su fibra. Sea $y=f(x)$; entonces~$\cal{O}_x$ es una $\cal{O}_y$-álgebra, localización de un álgebra de tipo finito, con extensión residual finita. Por otra parte, la hipótesis~(iii) implica que todo homomorfismo de~$\cal{O}_x$ en un álgebra~$A/J$, donde~$A$ es un álgebra local artiniana finita sobre~$\cal{O}_y$ y~$J$ un ideal contenido en su radical, se levanta a un homomorfismo de~$\cal{O}_x$ en el álgebra~$A$ (teniendo en cuenta que
\ifthenelse{\boolean{orig}}
{un morfismo de un $\Spec(B)$, con~$B$ local, en~$X$ está determinado biunívocamente por un homomorfismo local de un~$\cal{O}_x$ ($x\in X$) en~$B$}
{un morfismo $\Spec(B)\to X$, con~$B$ local y cuyo punto cerrado se envía a~$x$, está determinado biunívocamente por un homomorfismo local $\cal{O}_x\to B$}).
Por~\Ref{III.2.1} se sigue que~$\cal{O}_x$ es formalmente lisa sobre~$\cal{O}_y$ y, por tanto, lisa sobre~$\cal{O}_y$ en virtud de~\Ref{III.1.9}.

\begin{corollaire}
\label{III.3.2}
Sea $f\colon X\to Y$ como en~\Ref{III.3.1}. Las condiciones siguientes son equivalentes:
\begin{enumerate}
\item[(i)] $f$ es étale.
\item[(ii)] Se cumple la condición~\textup{(ii)} de~\Ref{III.3.1}, con \emph{unicidad} de la prolongación~$g$ de~$g_0$ a~$U$.
\item[(iii)] Se cumple la condición~\textup{(iii)} de~\Ref{III.3.1}, con \emph{unicidad} de~$g$.
\end{enumerate}
\end{corollaire}

Basta observar, en la demostración anterior de (i)$\To$(ii), que solo puede haber unicidad (cuando~$Y'_0$ no coincide con~$Y'$ en un entorno de~$z$) si~$n=0$, condición que se sabe suficiente.

\begin{corollaire}
\label{III.3.3}
Sea~$X$ un preesquema localmente de tipo finito sobre un anillo local \emph{completo}~$A$, sea~$y$ el punto cerrado de~$Y=\Spec(A)$
\ifthenelse{\boolean{orig}}
{y sea~$x$ un punto de~$f^{-1}(y)$}
{, denótese por $f\colon X\to Y$ el morfismo estructural, y sea~$x$ un punto de~$f^{-1}(y)$}
\emph{racional}
\ifthenelse{\boolean{orig}}
{sobre~$\kappa(y)$}
{sobre~$\kres(y)$}.
Si~$X$ es \emph{liso sobre~$A$ en~$x$}, existe una sección~$s$ de~$X$ sobre~$Y$ «que pasa por~$x$», es decir, tal que $s(y)=x$.
\end{corollaire}

En particular, si~$X$ es liso sobre~$A$, la aplicación natural
$$
\Gamma(X/Y)\to\Gamma(X\otimes_A k/k)
$$
de las secciones de~$X$ sobre~$Y$ en el conjunto de los puntos de la fibra $f^{-1}(y)=X\otimes_A k$ racionales sobre~$k$ es sobreyectiva. Este hecho era especialmente conocido y utilizado cuando~$A$ es un anillo de valoración discreta y~$X$
\marginpar{69}
es propio sobre~$A$ (de hecho, proyectivo sobre~$A$); en tal caso, las secciones de~$X$ sobre~$Y$ (es decir, los «puntos de~$X$ con valores en~$A$») se identifican también con las secciones racionales, es decir, con los puntos de $X\otimes_A K=X_K$ (que es un esquema propio y
\ifthenelse{\boolean{orig}}{simple}{liso}
sobre~$K$) con valores en $K=$ cuerpo de fracciones
\ifthenelse{\boolean{orig}}
{de~$A$}
{de~$A$,}
es decir, con los puntos de~$X$ racionales sobre~$K$.


\section{Prolongación infinitesimal local de los $S$-esquemas lisos}
\label{III.4}

\begin{theoreme}
\label{III.4.1}
Sean~$Y$ un preesquema localmente noetheriano, $Y_0$ un subpreesquema cerrado con el mismo espacio subyacente, $X_0$ un $Y_0$-preesquema liso y~$x$ un punto de~$X_0$. Existen entonces un entorno abierto~$U_0$ de~$x$, un preesquema~$X$ liso sobre~$Y$ y un $Y_0$-isomorfismo
$$
h\colon U_0\isomto X\times_Y Y_0.
$$
Además, si $(U'_0,X',h')$ es otra solución de este problema, «es isomorfa a la primera en un entorno de~$x$».
\end{theoreme}

Se deja al lector precisar qué se entiende con ello. Puede observarse que, dado~$U_0$, una solución del problema planteado equivale a dar, sobre~$U_0$, un haz de álgebras~$\cal{B}$ sobre~$f_0^{-1}(\cal{O}_Y)$ (donde~$f_0$ es la aplicación continua subyacente al morfismo estructural $U_0\to Y_0$), provisto de un homomorfismo de anillos $\cal{B}\to\cal{O}_{U_0}$ compatible con el homomorfismo
\ifthenelse{\boolean{orig}}
{$f^{-1}(\cal{O}_Y)\to f^{-1}(\cal{O}_{Y_0})$}
{$f_0^{-1}(\cal{O}_Y)\to f_0^{-1}(\cal{O}_{Y_0})$},
y que satisfaga:

(a) este homomorfismo induce un isomorfismo
$$
\ifthenelse{\boolean{orig}}
{\cal{B}\otimes_{f^{-1}(\cal{O}_Y)} f^{-1}(\cal{O}_{Y_0})\isomto
\cal{O}_{Y_0}}
{\cal{B}\otimes_{f_0^{-1}(\cal{O}_Y)} f_0^{-1}(\cal{O}_{Y_0})\isomto
\cal{O}_{U_0}};
$$

(b) $U_0$, provisto de~$\cal{B}$, se convierte en un $Y$-preesquema liso.

De este modo, el sentido preciso de la afirmación de unicidad local se hace particularmente evidente.

\subsubsection*{Demostración}
Puede suponerse ya que~$X_0$ es étale sobre un $Y_0[t_1,\dots,t_n]=Y'_0$. Este último puede considerarse como un subpreesquema cerrado de $Y'=Y[t_1,\dots,t_n]$ con el mismo espacio subyacente. Por~(\Ref{I.8.3}), existe un~$X$ étale sobre~$Y'$ y un $Y'_0$-isomorfismo
\ifthenelse{\boolean{orig}}
{$X\times_{Y'}Y'_0\isomto X'$}
{$X\times_{Y'}Y'_0\isomto X_0$}.
Con esto queda probada la existencia. Para la unicidad se utiliza la propiedad~\Ref{III.3.1}~\textup{(ii)} de los morfismos lisos, teniendo en cuenta el siguiente lema.
\marginpar{70}

\begin{lemme}
\label{III.4.2}
Sean~$Y$ un preesquema, $Y_0$ un subpreesquema cerrado definido por un haz de ideales~$\cal{J}$ localmente nilpotente, $X$ y~$X'$ dos $Y$-preesquemas y $u\colon X\to X'$ un $Y$-morfismo. Supóngase que~$X$ es plano sobre~$Y$. Para que~$u$ sea un isomorfismo es necesario y suficiente que lo sea $u_0\colon X\times_Y Y_0\to X'\times_Y Y_0$.
\end{lemme}

La demostración es sencilla: basta pasar al caso afín y considerar los graduados asociados. Obsérvese, además, que el enunciado análogo obtenido sustituyendo «isomorfismo» por «inmersión cerrada» sigue siendo válido, y sin hipótesis de planitud.

\begin{remarque}
\label{III.4.3}
Es esencial observar que la prolongación local~$X$ obtenida en~\Ref{III.4.1} \emph{no es canónica}; en otros términos, el isomorfismo local entre dos soluciones no es único, es decir, existen en general $Y$-automorfismos no triviales de~$X$ que inducen la identidad sobre el subpreesquema cerrado $X_0=X\times_Y Y_0$. Por ello, al construir prolongaciones infinitesimales \emph{globales} de preesquemas
\ifthenelse{\boolean{orig}}
{simples}
{lisos},
cabe esperar la existencia de una obstrucción de naturaleza cohomológica, que se precisará más adelante (\no \Ref{III.6}).
\end{remarque}


\section{Prolongación infinitesimal global de los morfismos}
\label{III.5}

Sean~$T$ un espacio topológico, $\cal{G}$ un haz de grupos
\ifthenelse{\boolean{orig}}
{sobre~$X$}
{sobre~$T$},
\ifthenelse{\boolean{orig}}{$P$}{$\cal{P}$}
un haz de conjuntos sobre~$T$ sobre el que opera~$\cal{G}$ (por la derecha, para fijar ideas). Se dice que~$\cal{P}$ es \emph{formalmente homogéneo principal}
\index{formalmente homogéneo principal (haz)|hyperpage}%
bajo~$\cal{G}$ si el conocido homomorfismo
$$
\cal{G}\times\cal{P}\to\cal{P}\times\cal{P}
$$
de haces de conjuntos, deducido de la operación de~$\cal{G}$ sobre~$\cal{P}$, es un \emph{isomorfismo}. Equivale a decir que, para todo $x\in T$, $\cal{P}_x$ es \emph{vacío o un espacio homogéneo principal} bajo el grupo ordinario~$\cal{G}_x$, o también que, para todo abierto~$U$ de~$T$, $\cal{P}(U)$ es vacío o un espacio homogéneo principal bajo el grupo ordinario~$\cal{G}(U)$. Se dice que
\ifthenelse{\boolean{orig}}{$P$}{$\cal{P}$}
es un \emph{haz homogéneo principal}
\index{homogéneo principal (haz)|hyperpage}%
\ifthenelse{\boolean{orig}}{bajo~$G$}{bajo~$\cal{G}$}
si lo es formalmente y si, además, los~$\cal{P}_x$ no son vacíos (en otros términos, si \emph{todos} los~$\cal{P}_x$ son espacios homogéneos principales, y por tanto no vacíos, bajo los~$\cal{G}_x$).
\marginpar{71}
\footnote{Parece preferible adoptar el término más breve y expresivo
\ifthenelse{\boolean{orig}}{«\emph{torsor bajo}~$G$»}{«\emph{torsor bajo}~$\cal{G}$»},
introducido en la tesis de J\ptbl\ \textsc{Giraud}.}
Recuérdese que el conjunto de las clases (salvo isomorfismo) de haces homogéneos principales bajo~$\cal{G}$ se identifica con el conjunto de cohomología~$\H^1(T,\cal{G})$, que es también el grupo de cohomología usual de~$T$ con coeficientes en~$\cal{G}$ cuando~$\cal{G}$ es conmutativo. Se tiene así, para todo~$\cal{P}$ homogéneo principal, una clase característica $c(\cal{P})\in\H^1(T,\cal{G})$, cuya trivialidad es necesaria y suficiente para que~$\cal{P}$ sea trivial (es decir, isomorfo a~$\cal{G}$, sobre el que~$\cal{G}$ opera por traslaciones a la derecha), o también para que~$\cal{P}$ tenga una sección.

\begin{proposition}
\label{III.5.1}
Sean~$S$ un preesquema, $X$ e~$Y$ preesquemas sobre~$S$, $Y_0$ un subpreesquema cerrado de~$Y$ definido por un haz de ideales~$\cal{J}$ sobre~$Y$ \emph{de cuadrado nulo}. Sea~$g_0$ un $S$-morfismo de~$Y_0$ en~$X$, y sea~$\cal{P}(g_0)$ el haz sobre~$Y$ cuyas secciones sobre un abierto~$U$ son las prolongaciones $g\colon U\to X$ de~$g_0|_{U\cap Y_0}$ a un $S$-morfismo~$g$. Entonces~$\cal{P}(g_0)$ es, de manera natural, un haz \emph{formalmente homogéneo principal} bajo el haz de grupos conmutativo
$$
\cal{G}=\SheafHom_{\cal{O}_{Y_0}}(g_0^*(\it{\Omega}^1_{X/S}),\cal{J}).
$$
\end{proposition}

Pongamos~$\cal{P}=\cal{P}(g_0)$. Debemos definir, para todo abierto~$U$ de~$Y$, una aplicación
$$
\cal{P}(U)\times\cal{G}(U)\to\cal{P}(U)
$$
de modo que: (a) para $g\in\cal{P}(U)$ fijo, la aplicación $s\mto gs$ de~$\cal{G}(U)$ en~$\cal{P}(U)$ sea biyectiva; (b) $\cal{P}(U)$ se convierta en un conjunto con grupo de operadores~$\cal{G}(U)$; (c) las aplicaciones precedentes sean compatibles con los operadores de restricción para un abierto $V\subset U$. La verificación de~(c) será trivial; para simplificar, puede suponerse, pues, que~$U=Y$. La verificación de~(b) (que, si se quiere, es de naturaleza local) se dejará al lector; nos limitaremos, por tanto, para un $g\in\cal{P}(Y)$ dado, a definir una biyección natural de~$\cal{G}(Y)$ sobre~$\cal{P}(Y)$. Se supone ya dado, pues, un $S$-morfismo $g\colon Y\to X$, y se busca una biyección canónica
\begin{equation*}
\label{eq:III.5.1.*}
\tag{$*$}
\Hom_{\cal{O}_{Y_0}}(g_0^*(\it{\Omega}^1_{X/S}),\cal{J})\isomto
\cal{P}(Y),
\end{equation*}
donde~$\cal{P}(Y)$ es el conjunto de los $S$-morfismos $g'$ de~$Y$ en~$X$ que inducen el mismo morfismo $g_0\colon Y_0\to X$ que~$g$. Dar un tal~$g'$ equivale a dar un $S$-morfismo $h\colon Y\to X\times_S X$ tal que ${\pr}_1\circ h=g$ y $h\circ i=(g_0,g_0)$, donde
\marginpar{72}
\ifthenelse{\boolean{orig}}
{${\pr}_1\colon X\times_S X$}
{${\pr}_1\colon X\times_S X\to X$}
es la primera proyección, $i\colon Y_0\to Y$ la inmersión canónica, y $(g_0,g_0)\colon Y_0\to X\times_S X$ es el morfismo $\Delta_{X/S}g_0$ de
\ifthenelse{\boolean{orig}}{componente}{componentes}
$g_0,g_0$:
\begingroup
\shorthandoff{>}% Spanish babel otherwise treats the Xy-pic arrow tip as a shorthand.
$$
\xymatrix@C=2.2cm{ X\times_S X \ar[d]_{{\pr}_1} & Y_0
\ar[l]_-{h_0=(g_0,g_0)} \ar[d]^i \\ X & Y \ar[l]_-g
\ar@{-->}[lu]_-h }
$$
\endgroup
Como~$h_0$ se factoriza por la inmersión diagonal~$\Delta_{X/S}$ y~$Y$ está en el entorno infinitesimal de primer orden de~$Y_0$ (es decir, $\cal{J}^2=0$), los~$h$ buscados se factorizan necesariamente, y de manera única, por el entorno infinitesimal de primer orden de la diagonal. Este se identifica (como $X$-preesquema gracias a~${\pr}_1$) con el espectro~$X'$ del haz de álgebras $\cal{O}_X+\it{\Omega}^1_{X/S}$ (donde el segundo sumando se considera un ideal de cuadrado nulo), y el morfismo diagonal $X\to X'$ corresponde a la aumentación canónica de este haz de álgebras. Pongamos $Y'=X'\times_XY$ y $Y'_0=Y'\times_YY_0=X'\times_XY_0$, de modo que los~$h$ buscados están en correspondencia biunívoca con las secciones~$u$ de~$Y'$ sobre~$Y$ que prolongan una sección dada~$u_0$ de~$Y'_0$ sobre~$Y_0$. Puede, además, identificarse~$Y'$ con el espectro del haz de álgebras sobre~$Y$ $\cal{A}=g^*(\cal{O}_X+\it{\Omega}^1_{X/S})=\cal{O}_Y+g^*(\it{\Omega}^1_{X/S})$, y~$Y'_0$
\ifthenelse{\boolean{orig}}
{con el haz de álgebras}
{con el espectro del haz de álgebras}
$\cal{A}_0=\cal{A}\otimes_{\cal{O}_Y}\cal{O}_{Y_0}=\cal{O}_{Y_0}+g_0^*(\it{\Omega}^1_{X/S})$.
Entonces la sección~$u_0$ es la definida por la aumentación canónica
\ifthenelse{\boolean{orig}}
{de~$A_0$ en~$\cal{O}_{Y_0}$}
{de~$\cal{A}_0$ en~$\cal{O}_{Y_0}$}.
Por tanto, $\cal{P}(Y)$ se identifica con el conjunto de los homomorfismos de álgebras $\cal{A}\to\cal{O}_Y$ que inducen la aumentación canónica $\cal{A}_0\to\cal{O}_{Y_0}$. Ahora bien, los homomorfismos de álgebras $\cal{A}\to\cal{O}_Y$
\ifthenelse{\boolean{orig}}{se corresponde}{se corresponden}
biunívocamente con los homomorfismos de módulos
\ifthenelse{\boolean{orig}}
{$\cal{M}\to\cal{A}$}
{$\cal{M}\to\cal{O}_Y$}
(poniendo, para simplificar, $\cal{M}=g^*(\it{\Omega}^1_{X/S})$), y nos interesan los que inducen el homomorfismo \emph{nulo} $\cal{M}_0\to\cal{O}_{Y_0}$ (donde $\cal{M}_0=\cal{M}\otimes_{\cal{O}_Y}\cal{O}_{Y_0}$), es decir, los que envían~$\cal{M}$ al ideal~$\cal{J}$ de la aumentación. Se obtiene, pues, el conjunto $\Hom_{\cal{O}_Y}(\cal{M},\cal{J})=\Hom_{\cal{O}_{Y_0}}(\cal{M}_0,\cal{J})$
(puesto que~$\cal{J}$ es anulado por~$\cal{J}$). Esta es la biyección canónica buscada~\eqref{eq:III.5.1.*}.

Teniendo en cuenta la implicación (i)$\To$(iii) de~\Ref{III.3.1}, se obtiene:

\begin{corollaire}
\label{III.5.2}
\marginpar{73}
Si~$X$ es liso sobre~$S$ (al menos en los puntos de~$g_0(Y_0)$), entonces~$\cal{P}$ es incluso un \emph{haz homogéneo principal} bajo el haz de grupos conmutativo~$\cal{G}$, que en este caso puede escribirse también
\begin{equation*}
\cal{G}=g_0^*(\goth{g}_{X/S})\otimes_{\cal{O}_{Y_0}}\cal{J},
\end{equation*}
donde~$\goth{g}_{X/S}$
\label{indnot:cb}\oldindexnot{$\goth{g}_{X/S}$|hyperpage}%
es el haz sobre~$X$ dual de~$\it{\Omega}_{X/S}^1$, es decir, el \emph{haz tangente}
\index{haz tangente|hyperpage}%
\index{tangente (haz)|hyperpage}%
(o \emph{haz de derivaciones})
\index{derivaciones (haz de)|hyperpage}%
\index{haz de derivaciones|hyperpage}%
de~$X$ respecto de~$S$. (Esta última fórmula proviene del hecho de que~$\it{\Omega}_{X/S}^1$ es entonces
\ifthenelse{\boolean{orig}}
{libre de tipo finito}
{localmente libre de tipo finito}).
\end{corollaire}

En particular, a este haz homogéneo principal le corresponde una clase de cohomología en~$\H^1(Y_0,\cal{G})$, cuya anulación es necesaria y suficiente para la existencia de un $S$-morfismo~$g$ que prolongue~$g_0$. Y si existe tal prolongación, el conjunto de todas las prolongaciones posibles es un espacio homogéneo bajo el grupo~$\H^0(Y_0,\cal{G})$.


En la aplicación de los métodos de la geometría formal, la situación es casi siempre la siguiente: se dan dos $S$-preesquemas $X$ e~$Y$ y un haz de ideales coherente~$\cal{I}$ sobre~$S$; se denota por~$S_n$ el subpreesquema cerrado de~$S$ definido por~$\cal{I}^{n+1}$, y se pone
\begin{equation*}
X_n=X\times_S S_n\quoi,\quad Y_n=Y\times_S S_n.
\end{equation*}
Se supone dado un $S_n$-morfismo
\begin{equation*}
g_n\colon Y_n\to X_n
\end{equation*}
(o, lo que equivale a lo mismo, un $S$-morfismo $Y_n\to X$, o también un $S_{n+1}$-morfismo $Y_n\to X_{n+1}$, puesto que tal morfismo induce necesariamente $Y_n\to X_n$), y se busca prolongarlo a un $S_{n+1}$-morfismo
\begin{equation*}
g_{n+1}\colon Y_{n+1}\to X_{n+1}.
\end{equation*}
(Si puede continuarse indefinidamente, se obtiene, pues, un morfismo $\widehat{Y}\to\widehat{X}$ entre los preesquemas formales completados de~$Y$ y~$X$ respecto de los ideales
\ifthenelse{\boolean{orig}}
{$\cal{IO}_X$ y~$\cal{IO}_Y$.)}
{$\cal{IO}_Y$ y~$\cal{IO}_X$.)}
Puede aplicarse~\Ref{III.5.1} sustituyendo allí $(S,X,Y,Y_0,g_0)$ por $(S_{n+1},X_{n+1},Y_{n+1},Y_n,g_n)$. El haz~$\cal{G}$ se convierte aquí en el haz de homomorfismos de módulos de~$g_n^*(\it{\Omega}_{X_{n+1}/S_{n+1}}^1)$ en~$\cal{J}=\cal{I}^{n+1}\cal{O}_Y/\cal{I}^{n+2}\cal{O}_Y$. Como~$\cal{J}$ es anulado por~$\cal{IO}_Y$, puede sustituirse entonces~$g_n^*(\it{\Omega}_{X_{n+1}/S_{n+1}}^1)$
\marginpar{74}
por el haz que induce sobre~$Y_0$, a saber, $h_0^*(\it{\Omega}_{X/S}^1)$, donde~$h_0$ es el compuesto $Y_0\to Y_n\to X_{n+1}$, o también el compuesto $Y_0\to X_0\to X_{n+1}$, donde $g_0\colon Y_0\to X_0$ es inducido por~$g_n$. Como la imagen inversa de~$\it{\Omega}_{X_{n+1}/S_{n+1}}^1$ sobre~$X_0=X_{n+1}\times_{S_{n+1}}S_0$ es~$\it{\Omega}_{X_0/S_0}^1$, se ve que también se tiene
\begin{equation*}
\cal{G}=\SheafHom_{\cal{O}_{Y_0}}
(g_0^*(\it{\Omega}_{X_0/S_0}^1),
\cal{I}^{n+1}\cal{O}_Y/\cal{I}^{n+2}\cal{O}_Y).
\end{equation*}
Se obtiene, pues, el

\begin{corollaire}
\label{III.5.3}
Sean~$S,X,Y,\cal{I},g_n$ como antes, y sea~$\cal{P}(g_n)$ el haz sobre~$Y$ cuyas secciones sobre un abierto~$U$ son
\ifthenelse{\boolean{orig}}
{las prolongaciones~$g_{n+1}$ de~$g_n$ a un $S_{n+1}$-morfismo $Y_{n+1}\to X_{n+1}$}
{las prolongaciones $g_{n+1}\colon U_{n+1}\to X_{n+1}$ de $g_n|_{U_n}$, donde $U_m=U\times_S S_m$}.
Entonces
\ifthenelse{\boolean{orig}}{$\cal{P}$}{$\cal{P}(g_n)$}
es un haz formalmente homogéneo principal bajo el haz de grupos
\begin{equation*}
\cal{G}=\SheafHom_{\cal{O}_{Y_0}}
(g_0^*(\it{\Omega}_{X_0/S_0}^1),
\gr_{\cal{IO}_Y}^{n+1}(\cal{O}_Y)).
\end{equation*}
\end{corollaire}

En particular:

\begin{corollaire}
\label{III.5.4}
Si, además, $X$ es liso sobre~$S$ (al menos en los puntos de~$g_0(Y_0)$), entonces
\ifthenelse{\boolean{orig}}{$\cal{P}$}{$\cal{P}(g_n)$}
es incluso un haz homogéneo principal.
\ifthenelse{\boolean{orig}}{}
{\enlargethispage{.5cm}}%
En particular, define una clase de obstrucción en~$\H^1(Y_0,\cal{G})$, cuya anulación es necesaria y suficiente para la existencia de una prolongación global~$g_{n+1}$ de~$g_n$. Y si existe tal prolongación, el conjunto de todas las prolongaciones globales es un espacio homogéneo principal bajo~$\H^0(Y_0,\cal{G})$. Por último, en el caso considerado, el haz~$\cal{G}$ puede escribirse también
\begin{equation*}
\cal{G}=g_0^*(\goth{g}_{X_0/S_0})\otimes_{\cal{O}_{Y_0}}
\gr_{\cal{IO}_Y}^{n+1}(\cal{O}_Y).
\end{equation*}
\end{corollaire}


Procediendo paso a paso, se ve, pues, que si todos los~$\H^1(Y_0,\cal{G}_n)$
son nulos (donde $\cal{G}_n=g_0^*(\goth{g}_{X_0/S_0})\otimes\gr_{\cal{IO}_Y}^{n}(\cal{O}_Y)$), entonces, partiendo de un~$g_k$ cualquiera, puede prolongarse sucesivamente a~$g_{k+1},\dots$. En particular, si~$\cal{I}$ es nilpotente, podrá hallarse una prolongación~$g$ de~$g_k$ a~$Y$. La condición de anulación de los~$\H^1$ se cumple, en particular, si~$Y_0$ es afín. Se obtiene, por tanto:

\begin{corollaire}
\label{III.5.5}
En
\marginpar{75}
el enunciado del teorema~\Ref{III.3.1}, se obtiene una condición necesaria y suficiente equivalente a las demás suponiendo que el~$Y'$ que interviene en~\textup{(ii)} (o~\textup{(ii~bis)}) sea afín, y exigiendo la existencia de una \emph{prolongación global}~$g$ de~$g_0$ a todo~$Y'$.
\end{corollaire}

Obsérvese que la demostración de~\Ref{III.3.1} no podía darnos directamente este resultado.

Un caso importante es aquel en que~$Y$ es \emph{plano} sobre~$S$; se tiene entonces
\begin{equation*}
\gr^n(\cal{O}_Y)=\gr^n(\cal{O}_S)\otimes_{\cal{O}_{S_0}}\cal{O}_{Y_0},
\end{equation*}
y cuando, además, los~$\gr^n(\cal{O}_S)$ son \emph{localmente libres sobre}
\ifthenelse{\boolean{orig}}
{~$S$}
{~$S_0$},
se obtiene
\begin{equation*}
\cal{G}_n=
\SheafHom_{\cal{O}_{Y_0}}(g_0^*(\it{\Omega}_{X_0/S_0}^1),\cal{O}_{Y_0})
\otimes_{\cal{O}_{S_0}}\gr^n(\cal{O}_S),
\end{equation*}
o también, si~$\it{\Omega}_{X_0/S_0}^1$ es asimismo localmente libre (por ejemplo, si~$X$ es liso sobre~$S$),
\begin{equation*}
\cal{G}_n=g_0^*(\goth{g}_{X_0/S_0})
\otimes_{\cal{O}_{S_0}}\gr^n(\cal{O}_S).
\end{equation*}
Si, por ejemplo, $S$ es afín,
de anillo~$A$, y~$\cal{I}$ está definido por un ideal~$I$ de~$A$, se obtiene
\begin{equation*}
\H^i(Y_0,\cal{G}_n)=\H^i(Y_0,\cal{G}_0)\otimes_A\gr_I^n(A)
\end{equation*}
para todo~$i$ (en efecto, la cuestión es local sobre~$S_0$, y se reduce al caso en que se tensoriza por un módulo libre). \emph{En este caso, la anulación de~$\H^1(Y_0,\cal{G}_0)$ implica que todas las obstrucciones a las prolongaciones sucesivas de~$g_n$ son nulas}. Se obtiene, pues:

\begin{corollaire}
\label{III.5.6}
Sean $(S,X,Y,\cal{I},g_n)$ como antes; supóngase además que~$X$ es liso sobre~$S$, que~$Y$ es plano sobre~$S$, que~$S$ es afín y que los~$\gr^n(\cal{O}_S)=\cal{I}^n/\cal{I}^{n+1}$ son localmente libres. Entonces la obstrucción para construir~$g_{n+1}$ se encuentra en~$\H^1(Y_0,\cal{G}_0)\otimes_A\gr_I^{n+1}(A)$ (donde~$A$ es el anillo de~$S$ e~$I$ el ideal de~$A$ que define~$\cal{I}$), poniendo
\begin{equation*}
\cal{G}_0=g_0^*(\goth{g}_{X_0/S_0}).
\end{equation*}
Si~$\H^1(Y_0,\cal{G}_0)=0$, entonces~$g_n$ puede prolongarse a un $\widehat{S}$-morfismo $\widehat{g}\colon\widehat{Y}\to\widehat{X}$.
\end{corollaire}


Por supuesto, este resultado seguiría siendo válido tal cual si, en lugar de partir de $S$-preesquemas ordinarios~$X$ e~$Y$, se partiese de preesquemas formales $\widehat{\cal{I}}$-ádicos
\marginpar{76}
$\goth{X}$ y~$\goth{Y}$ sobre~$\widehat{S}$. Permite demostrar, por ejemplo, que ciertos esquemas formales propios sobre un anillo local completo son, de hecho, algebraicos. En efecto, procediendo como en el lema~\Ref{III.4.2}, se obtiene:

\begin{corollaire}
\label{III.5.7}
Bajo las condiciones de~\Ref{III.5.6}, si~$g_0$ es un isomorfismo, lo mismo ocurre con~$\widehat{g}$.
\end{corollaire}

\noindent
(N.B.: el mismo resultado vale para las inmersiones cerradas).

Se obtiene así:

\begin{proposition}
\label{III.5.8}
Sean~$A$ un anillo local completo de ideal maximal~$\goth{m}$,
\ifthenelse{\boolean{orig}}{}{de}
cuerpo residual~$k$, y sean~$\goth{X}$ y~$\goth{Y}$ dos preesquemas formales $\goth{m}$-ádicos sobre~$A$, planos sobre~$A$ (es decir, para todo~$n$, $X_n$ e~$Y_n$ son planos sobre~$A_n=A/\goth{m}^{n+1}$). Supóngase que~$X_0=\goth{X}\otimes_Ak$ es liso sobre~$k$ y que~$\H^1(X_0,\goth{g}_{X_0/k})=0$. Entonces todo $k$-isomorfismo de~$Y_0$ sobre~$X_0$ se prolonga a un $A$-isomorfismo de~$\goth{Y}$ sobre~$\goth{X}$; esta prolongación es única si, además, $\H^0(X_0,\goth{g}_{X_0/k})=0$.
\end{proposition}

Esto da, en particular, un resultado de \emph{unicidad} para un preesquema formal liso sobre~$A$ que se reduce a un preesquema~$X_0$ dado (siempre que $\H^1(X_0,\goth{g}_{X_0/k})=0$). Además, si~$\goth{X}$ y~$\goth{Y}$ provienen de esquemas ordinarios propios sobre~$A$, sean~$X$ e~$Y$, se sabe, por el teorema de existencia de haces en geometría formal (cf. la exposición del Seminario Bourbaki, \No 182\footnote{Cf. EGA III 5.4.1 para la demostración.}), que hay una correspondencia biunívoca entre los $A$-isomorfismos $Y\isomto X$ y
los $A$-isomorfismos de los completados formales; por tanto:

\begin{corollaire}
\label{III.5.9}
El enunciado precedente~\Ref{III.5.8} sigue siendo válido sustituyendo~$\goth{X}$ y~$\goth{Y}$ por esquemas ordinarios~$X$ e~$Y$ sobre~$A$ que sean \emph{propios}.
\end{corollaire}

Por último, cuando~$\goth{X}$ es un esquema formal propio sobre~$A$ y~$\goth{Y}$ es de la forma~$\widehat{Y}$, donde~$Y$ es un esquema ordinario propio sobre~$A$, la proposición~\Ref{III.5.8} da condiciones suficientes para poder hallar un isomorfismo de~$\goth{X}$ con~$\widehat{Y}$, y por tanto para que el esquema formal~$\goth{X}$ sea, de hecho, «algebraico» (es decir, isomorfo a un~$\widehat{X}$, con~$X$ un esquema ordinario propio sobre~$A$, que quedará entonces determinado canónicamente). Esto ocurre, en particular, si~$X_0=\PP_k^r$ (o, más generalmente, si~$X_0$ es un esquema de
\ifthenelse{\boolean{orig}}{Sévéri--Brauer}{Severi--Brauer},
es decir, si se vuelve
\marginpar{77}
isomorfo al espacio proyectivo tipo sobre la clausura algebraica de~$k$): todo esquema formal propio y plano sobre~$A$, de fibra~$\PP_k^r$, es algebraizable y, más precisamente, es isomorfo al completado formal $\goth{m}$-ádico de~$\PP_A^r$. En particular (gracias al «teorema de existencia»), todo esquema ordinario propio sobre~$A$, de fibra~$\PP_k^r$, es isomorfo a~$\PP_A^r$ (siendo~$A$ un anillo local completo). Utilizando la teoría del descenso, puede demostrarse que, si~$A$ no es completo, $X$ se vuelve isomorfo a~$\PP^r$ tras una extensión finita y étale de la base $A\to A'$ (y, bajo esta forma, el resultado sigue siendo válido para una fibra que sea un esquema de
\ifthenelse{\boolean{orig}}{Sévéri--Brauer}{Severi--Brauer}).


\section{Prolongación infinitesimal global de los $S$-esquemas lisos}
\label{III.6}

Bajo las condiciones del teorema~\Ref{III.4.1}, se trata de averiguar si existe un preesquema~$X$ liso sobre~$Y$ tal que $X\times_YY_0$ sea $Y_0$-isomorfo a~$X_0$, sabiendo que tal esquema «existe localmente sobre~$X_0$». Retomando el método de construcción paso a paso, se llega a sustituir la letra~$Y$ por~$S$, a suponer dado un subpreesquema cerrado~$S_0$ de~$S$ definido por un haz de ideales~$\cal{I}$ (que ya no es necesario suponer localmente nilpotente), a introducir los subpreesquemas cerrados~$S_n$ de~$S$ definidos por los~$\cal{I}^{n+1}$ y a suponer dado un preesquema~$X_n$ liso sobre~$S_n$. Se busca un $S_{n+1}$-preesquema~$X_{n+1}$ «que se reduzca a~$X_n$», es decir, provisto de un isomorfismo
\begin{equation*}
X_{n+1}\times_{S_{n+1}}S_n\isomfrom X_n,
\end{equation*}
que sea \emph{liso} sobre~$S_{n+1}$ (o, lo que equivale a lo mismo por~\Ref{II.2.1}, \emph{plano} sobre~$S_{n+1}$). Como se señaló en el \No\Ref{III.4}, dar un objeto de este tipo equivale a dar un haz de álgebras~$\cal{B}$ sobre~$f^{-1}(\cal{O}_{S_{n+1}})$ (donde~$f$ es la aplicación continua subyacente al morfismo estructural $X_n\to S_n$), provisto de una aumentación $\cal{B}\to\cal{O}_{X_n}$ compatible con la aumentación $f^{-1}(\cal{O}_{S_{n+1}})\to f^{-1}(\cal{O}_{S_n})$ y que satisfaga dos condiciones (a) y (b) que no volveremos a escribir; nos limitaremos a observar que son de \emph{naturaleza local} sobre el espacio topológico subyacente a~$X_n$. Por~\Ref{III.4.1} se sabe que existe una solución localmente. Además, es única salvo isomorfismo (no único), al menos localmente. Comencemos por precisar este punto:

\begin{proposition}
\label{III.6.1}
Sea
\marginpar{78}
$X_{n+1}$ sobre~$S_{n+1}$ que se reduzca a~$X_n$ sobre~$S_n$. Entonces el haz (sobre el espacio topológico subyacente a~$X_n$, o también a~$X_0$) de los $S_{n+1}$-automorfismos de~$X_{n+1}$ que inducen la identidad sobre~$X_n$ es canónicamente isomorfo a
\begin{equation*}
\cal{G}=\goth{g}_{X_0/S_0}\otimes_{\cal{O}_{S_0}}
\gr_{\cal{I}}^{n+1}(\cal{O}_S)
\end{equation*}
(como haz de grupos).
\end{proposition}

En efecto, por~\Ref{III.5.4} y~\Ref{III.4.2}, este haz es un haz homogéneo principal bajo~$\cal{G}$. Como está provisto de una sección privilegiada (el automorfismo identidad de~$X_{n+1}$), se identifica, pues, como haz de conjuntos con~$\cal{G}$. Hay que comprobar que esta identificación es compatible con las estructuras de grupo. Esto es fácil y, por lo demás, un caso particular de un resultado más general sobre la compatibilidad de las estructuras de fibrados principales de~\Ref{III.5.1} y~\Ref{III.5.3} con la composición de morfismos (resultado que no enunciamos aquí, pero que deberá figurar en el \emph{hyperplodoque}).

En particular, el haz sobre~$X_0$ de los gérmenes de automorfismos de~$X_{n+1}$, con las estructuras explicitadas, es \emph{conmutativo}. Se sigue que, si~$X_{n+1}'$ es otra solución del problema, isomorfa a~$X_{n+1}$ sobre el abierto~$U$ de~$X_0$, el isomorfismo de~$\SheafAut(X_{n+1})|U$ sobre~$\SheafAut(X_{n+1}')|U$ deducido, por transporte de estructura, de un isomorfismo $X_{n+1}|U\isomto X_{n+1}'|U$ \emph{no depende} de la elección de este último. (No es, por lo demás, sino el isomorfismo identidad de~$\cal{G}$ cuando se identifican uno y otro haz de automorfismos con~$\cal{G}$ gracias a~\Ref{III.6.1}).

De~\Ref{III.6.1} se deduce:

\begin{corollaire}
\label{III.6.2}
Sean~$X_{n+1}$ y~$X_{n+1}'$ lisos sobre~$S_{n+1}$ y «que se reduzcan a~$X_n$». Entonces el haz (sobre el espacio subyacente a~$X_0$) de los $S_{n+1}$-isomorfismos de~$X_{n+1}$ sobre~$X_{n+1}'$ que inducen la identidad sobre~$X_n$ es, de manera natural, un haz homogéneo principal bajo~$\cal{G}$.
\end{corollaire}

Esto expresa, en efecto, que~$X_{n+1}$ y~$X_{n+1}'$ son localmente isomorfos y que el haz de gérmenes de automorfismos del primero es~$\cal{G}$.

Obsérvese ahora que, en virtud de~\Ref{III.4.1}, siempre puede hallarse un recubrimiento~$(U_i)$
\marginpar{79}
de~$X_n$ por abiertos, que pueden suponerse afines, y para todo~$i$ un esquema liso~$X^i$ sobre~$S_{n+1}$ que se reduzca
\ifthenelse{\boolean{orig}}{a~$U_i=U_{in}$.}{a~$U_i$.}
Supongamos, para simplificar, que~$X_n$ es \emph{separado}; entonces los~$U_{ij}=U_i\cap U_j$ siguen siendo abiertos \emph{afines} de~$X_n$. Como el~$\H^1$ de un abierto de este tipo, con valores en el haz cuasicoherente~$\cal{G}$, es nulo, se deduce por
\ifthenelse{\boolean{orig}}{}{el}
corolario~\Ref{III.6.2} que~$X^i|U_{ij}$ es isomorfo a~$X^j|U_{ij}$; sea
$$
f_{ji}\colon X^i|U_{ij}\isomto X^j|U_{ij}
$$
un isomorfismo de este tipo. Está determinado salvo una sección de~$\cal{G}$ sobre~$U_{ij}$.
Pongamos, para toda terna de índices:
$$
f_{ji}^{(k)}=f_{ji}|U_{ijk}\quad\text{donde }U_{ijk}=U_i\cap U_j\cap U_k.
$$
Si se tuviera
\begin{equation*}
\label{eq:III.6.1}
\tag{1} {f_{kj}^{(i)}f_{ji}^{(k)}=f_{ki}^{(j)},}
\end{equation*}
se seguiría que los~$X^i$ «se pegan» mediante los~$f_{ji}$ y, por tanto, que definen una solución~$X=X_{n+1}$ del problema planteado. Más generalmente, existe una solución de este tipo si pueden modificarse los~$f_{ji}$ en unos~$f_{ji}'$:
\begin{equation*}
\label{eq:III.6.2}
\tag{2} {f_{ji}'=f_{ji}g_{ji}\quad (g_{ji}\in\Gamma(U_{ij},\cal{G}))}
\end{equation*}
de modo que los~$f_{ji}'$ satisfagan la condición de transitividad anterior. Esta condición suficiente para la existencia de una solución es también necesaria, como se ve recordando que una solución~$X$ de este tipo debe ser isomorfa a~$X^i$ sobre cada~$U_i$. Esto permite elegir isomorfismos
$$
f_i\colon X|U_i\isomto X^i
$$
y definir
$$
f_{ji}'=(f_j|U_{ij})(f_i|U_{ij})^{-1}\colon X^i|U_{ij}\isomto X^j|U_{ij}
$$
que satisfacen la condición de pegado.

Ahora bien, pongamos
\begin{equation*}
\label{eq:III.6.3}
\tag{3}
\ifthenelse{\boolean{orig}}
{{f_{ijk}=(f_{ki}^{(j)})^{-1}f_{kj}^if_{ji}^k,}}
{{f_{ijk}=(f_{ki}^{(j)})^{-1}f_{kj}^{(i)}f_{ji}^{(k)},}}
\end{equation*}
que es un automorfismo de~$X^i|U_{ijk}$, que identificaremos con una sección de~$\cal{G}$ gracias a~\Ref{III.6.1}. Se comprueba, mediante un pequeño cálculo formal que se deja al lector, que es un $2$-\emph{cociclo}~$f$ del recubrimiento abierto
\marginpar{80}
$\cal{U}=(U_i)$ con coeficientes en~$\cal{G}$. El mismo cálculo muestra que, bajo~\eqref{eq:III.6.2}, la condición de pegado~\eqref{eq:III.6.1} \emph{para los}~$f_{ij}'$ equivale a la fórmula
\begin{equation*}
\label{eq:III.6.4}
\tag{4} {f=dg,}
\end{equation*}
donde~$g=(g_{ij})$ se considera una $1$-cocadena de~$\cal{U}$ con coeficientes en~$\cal{G}$. Por tanto, \emph{la condición necesaria y suficiente para la existencia de una solución del problema es que sea nula la clase de cohomología en~$\H^2(\cal{U},\cal{G})$ definida por el cociclo~\eqref{eq:III.6.3}}. Obsérvese, además, que, como~$\cal{U}=(U_i)$ es un recubrimiento afín de~$X_0$, que es un \emph{esquema}, $\H^2(\cal{U},\cal{G})$ se identifica con~$\H^2(X_0,\cal{G})$. Es inmediato, por lo demás, que la clase de cohomología así obtenida en~$\H^2(X_0,\cal{G})$ no depende del recubrimiento afín considerado. La llamaremos \emph{clase de obstrucción a la prolongación de~$X_n$ a un esquema~$X_{n+1}$ liso sobre~$S_{n+1}$}.

Supongamos nula esta obstrucción. Entonces un razonamiento esbozado antes muestra que toda solución~$X=X_{n+1}$ es isomorfa a una solución obtenida por pegado a partir de isomorfismos~$f_{ji}'$, que pueden suponerse de la forma~\eqref{eq:III.6.2}; la condición de pegado no es otra que
\ifthenelse{\boolean{orig}}{\eqref{eq:III.6.3}}{\eqref{eq:III.6.4}}.
El conjunto de los~$g$ admisibles es, pues, un espacio homogéneo principal bajo el grupo~$\mathrm{Z}^1(\cal{U},\cal{G})$ de los $1$-cociclos de~$\cal{U}$ con coeficientes en~$\cal{G}$. Además, se comprueba inmediatamente que \emph{dos cocadenas~$g$ y~$g'$} (tales que $dg=dg'=f$) \emph{definen soluciones isomorfas si y solo si el cociclo~$g-g'$ es de la forma~$dh$}, donde $h=(h_i)\in\mathrm{C}^0(\cal{U},\cal{G})$. Se obtiene, por tanto:


\begin{theoreme}
\label{III.6.3}
\index{obstrucción a la prolongación de un esquema a un esquema liso|hyperpage}%
Sean $(S,\cal{I},X_n)$ como antes, y supóngase que~$X_n$ es separado\footnote{En realidad, esta condición es innecesaria, y pueden evitarse los cálculos de cociclos precedentes. Cf. el libro de J.\ \textsc{Giraud}, \emph{Cohomologie Non Abélienne} (de próxima aparición en Springer Verlag, 1971). Compárense las observaciones~\Ref{III.6.4}.}. Entonces puede definirse canónicamente una clase de obstrucción en~$\H^2(X_0,\cal{G})$ (donde~$\cal{G}$ se define en~\Ref{III.6.1}), cuya anulación es necesaria y suficiente para la existencia de un esquema~$X_{n+1}$ liso sobre~$S_{n+1}$ que se reduzca a~$X_n$. Si esta obstrucción es nula, el conjunto de las clases de $S_{n+1}$-preesquemas~$X_{n+1}$ que se reduzcan a~$X_n$, salvo isomorfismo que induzca la identidad sobre~$X_n$, es, de manera natural, un espacio homogéneo principal bajo~$\H^1(X_0,\cal{G})$.
\end{theoreme}

\begin{remarques}
\label{III.6.4}
A partir de~\Ref{III.6.1}, los razonamientos hechos aquí son puramente formales y se transcriben ventajosamente en el marco de las categorías locales,
\marginpar{81}
o incluso de las categorías fibradas generales. La clase de obstrucción a la existencia de un objeto «global» de una categoría ---en la que puede encontrarse un objeto «localmente», dos objetos cualesquiera son siempre «localmente isomorfos» y el grupo de automorfismos de todo objeto es conmutativo--- obtenida así en el contexto general contiene, como caso particular, el «segundo morfismo de borde» de una sucesión exacta de haces de grupos no necesariamente conmutativos, estudiado, por ejemplo, por Grothendieck en Kansas o
\ifthenelse{\boolean{orig}}
{Tohoku}
{T\^ohoku}.
El cálculo simplón mediante cociclos efectuado aquí debe considerarse, pues, un mal menor, debido a la inexistencia de un texto de referencia satisfactorio.
\end{remarques}

\refstepcounter{proposition}% III.6.5
\setcounter{subsection}{4}%
\subsection{}
\label{III.6.5}
Obsérvese que en~\Ref{III.6.3} no hay, en general, ningún elemento privilegiado en el espacio homogéneo principal considerado bajo~$\H^1(X_0,\cal{G})$. Esto se traduce, en particular, en que se obtiene ---localizando sobre~$S$--- un haz homogéneo principal sobre~$S_0$, de grupo estructural~$\R^1f_*(\cal{G})$, que no es necesariamente trivial; es decir, define una clase de cohomología en~$\H^1(S_0,\R^1f_*(\cal{G}))$ que no es necesariamente nula. (Esto ocurre cuando se supone que la clase $d\in\H^2(X_0,\cal{G})$ no es nula, pero es nula «localmente sobre~$S$»; es decir, define la sección nula de~$\R^2f_*(\cal{G})$, o sea, el elemento nulo de~$\H^0(S_0,\R^2f_*(\cal{G}))$).

\refstepcounter{proposition}% III.6.6
\setcounter{subsection}{5}%
\subsection{}
\label{III.6.6}
Por ahora no se sabe casi nada del mecanismo algebraico general de las clases de cohomología introducidas en este número, de sus relaciones con el número precedente, ni puede decirse nada preciso de ellas en los casos particulares más sencillos, como el de los esquemas abelianos sobre anillos artinianos\footnote{\label{III.6.6.p24}%
\ifthenelse{\boolean{orig}}
{Ahora se sabe que esta obstrucción es siempre nula en este caso.}
{Ahora se sabe que esta obstrucción es siempre nula en este caso \lcrochetbf añadido en 2003 (MR): cf. F\ptbl Oort, \emph{Finite group schemes, local moduli for abelian varieties and lifting problems}, Algebraic Geometry Oslo 1970, Wolters-Noordhoff, 1972, pp\ptbl223--254\rcrochetbf.}}. Cabe esperar que haya quien estudie a fondo la cuestión, que parece particularmente interesante. Está íntimamente ligada, en particular, a la «teoría de los espacios de módulos» de las estructuras algebraicas.

\begin{corollaire}
\label{III.6.7}
Supongamos que~$\H^2(X_0,\cal{G})=0$. Entonces existe un~$X_{n+1}$, y es único salvo isomorfismo si, además, $\H^1(X_0,\cal{G})=0$.
\end{corollaire}

En particular, procediendo paso a paso y observando que un esquema afín es acíclico para todo haz cuasicoherente, se concluye:

\begin{corollaire}
\label{III.6.8}
Bajo
\marginpar{82}
las condiciones del teorema~\Ref{III.4.1}, si~$X_0$ es afín, existe un~$X$ liso sobre~$Y$ que se reduce a~$X_0$, y este~$X$ es único salvo un isomorfismo, que no es único.
\end{corollaire}

Obsérvese que la demostración directa del teorema~\Ref{III.4.1} no podía darnos este resultado.

\begin{corollaire}
\label{III.6.9}
Bajo las condiciones de~\Ref{III.6.3}, supóngase que~$S$ es afín de anillo~$A$, que~$\cal{I}$ está definido por un ideal~$I$ de~$A$ y, por último, que los~$\gr^n_{\cal{I}}(\cal{O}_S)=\cal{I}^n/\cal{I}^{n+1}$ son localmente libres. Entonces~$\H^i(X_0,\cal{G})$ se identifica con~$\H^i(X_0,\cal{G}_0)\otimes_A\gr_I^{n+1}(A)$, donde
$$
\cal{G}_0=\goth{g}_{X_0/S_0},
$$
de modo que la clase de obstrucción a la prolongación de~$X_n$ se encuentra en~$\H^2(X_0,\cal{G}_0)\otimes_A\gr_I^{n+1}(A)$ y, si es nula, el conjunto de las clases de soluciones, salvo isomorfismo, es un espacio homogéneo principal bajo~$\H^1(X_0,\cal{G}_0)\otimes_A\gr_I^{n+1}(A)$.
\end{corollaire}

En particular:

\begin{corollaire}
\label{III.6.10}
Bajo las condiciones de~\Ref{III.6.9}, supóngase que
$$
\H^2(X_0,\goth{g}_{X_0/S_0})=0.
$$
Entonces existe un esquema formal
\ifthenelse{\boolean{orig}}{}{~$\goth{X}$}
que es $\hat I$-ádico sobre el completado formal $\cal{I}$-ádico~$\hat{S}$ de~$S$, que es «liso sobre~$S$» (es decir, los~$\goth{X}_p$ son lisos sobre los~$S_p$) y que se reduce a~$X_n$, es decir, provisto de un isomorfismo
$$
\goth{X}\times_S S_n\isomfrom X_n.
$$
Si, además, $\H^1(X_0,\goth{g}_{X_0/S_0})=0$, tal~$\goth{X}$ es único salvo isomorfismo.
\end{corollaire}

En efecto, se construyen paso a paso~$X_{n+1},X_{n+2},\dots$, de donde se obtiene~$\goth{X}$ pasando al límite inductivo de los~$X_i$. La afirmación de unicidad figura ya en el número precedente.


\section[Aplicación a la construcción de esquemas formales...]{Aplicación a la construcción de esquemas formales y de esquemas ordinarios lisos sobre un anillo local completo~$A$}
\label{III.7}

Los resultados del número precedente permiten a veces demostrar la existencia
\marginpar{83}
de un esquema formal $\goth{m}$-ádico sobre un anillo de este tipo, que se reduzca a un esquema liso~$X_0$ dado sobre~$k$. Distinguiremos dos casos:

\medskip
\noindent a) \emph{$A$ es equicaracterístico} (es el caso, en particular, si~$k$ es de característica~$0$). Se sabe entonces que existe un \emph{cuerpo de coeficientes de}~$A$, es decir, un subcuerpo~$k'$ tal que~$A\to k$ induce un isomorfismo $k'\isomto k$. \emph{Existe entonces incluso un esquema ordinario liso sobre~$A$ que se reduce a~$X_0$}, a saber, $X=X_0\otimes_kA$, donde~$A$ se considera como álgebra sobre~$k$ mediante el homomorfismo $k\to k'\to A$ definido por~$k'$. Conviene observar, sin embargo, que esta construcción no es «natural»; es fácil convencerse (ya en el caso en que $A=k[t]/(t^2)$, álgebra de los números duales) de que otro homomorfismo de levantamiento~$k\to A$ (definido, en este caso, por una derivación absoluta de~$k$ en sí mismo) define un~$X'$ sobre~$A$ que, en general, \emph{no es isomorfo a~$X$} (si~$\H^1(X_0,\goth{g}_{X_0/k})\neq0$). Sería interesante, por lo demás, estudiar (para~$k$ de característica~$0$, o no perfecto y de
\ifthenelse{\boolean{orig}}
{car.}
{característica}~$p>0$) cuáles son los~$X$ lisos sobre~$A$ que se obtienen así y bajo qué condición dos homomorfismos~$k\to A$ definen $A$-esquemas isomorfos. No obstante, la existencia de~$k'$ basta para implicar que la primera obstrucción al levantamiento de~$X_0$, que se encuentra en~$\H^2(X_0,\goth{g}_{X_0/k})\otimes_k\goth{m}/\goth{m}^2$, es necesariamente nula. Naturalmente, una vez levantado~$X_0$ a un~$X_1$ liso sobre~$A/\goth{m}^2$, la nueva obstrucción a la construcción de~$X_2$ no será nula en general: dependerá de un elemento variable de cierto espacio homogéneo principal bajo~$\H^1(X_0,\cal{G}_0)\otimes\goth{m}/\goth{m}^2$ y se encontrará en~$\H^2(X_0,\cal{G}_0)\otimes\goth{m}^2/\goth{m}^3$; convendría estudiar detenidamente la situación\footnote{Sin duda está descrita por la operación corchete de Kodaira--Spencer (cf. Seminario Cartan, 1960/61, Exp\ptbl 4).}.

\medskip
\noindent b) \emph{$A$ es de característica mixta.} En este caso no se sabe nada (salvo si, por fortuna, $\H^2(X_0,\goth{g}_{X_0/k})=0$, en cuyo caso puede construirse un esquema formal $\goth{m}$-ádico
\ifthenelse{\boolean{orig}}{simple}{liso}
sobre~$A$ que se reduzca
\ifthenelse{\boolean{orig}}{a~$k$}{a~$X_0$}). Incluso si $A=\ZZ/p^2\ZZ$ y si~$X_0$ es un esquema «abeliano» de dimensión~$2$, no se sabe si puede levantarse a un~$X=X_1$ liso sobre~$A$\kern1pt\footnote{Esto ya está demostrado; cf. la nota~\Ref{III.6.6.p24} de la página~\pageref{III.6.6.p24}.}; por otra parte, no se tiene ningún ejemplo de un~$X_0$ del que se haya demostrado que no procede de un esquema ordinario~$X$ liso sobre~$A$. (Tengo la impresión de que tal ejemplo debe existir, con~$X_0$ una superficie proyectiva)\footnote{\label{nIII.3}Desde entonces, J.-P\ptbl Serre ha construido un ejemplo de este tipo
\ifthenelse{\boolean{orig}}
{(Proc.\ Nat.\ Acad.\ Sc.\ USA, vol\ptbl 47, \No 1, pp~108--109, 1961)}
{(Proc.\ Nat.\ Acad.\ Sc.\ USA \textbf{47} (1961), \No 1, pp\ptbl108--109)}, al menos para ciertas dimensiones. D\ptbl Mumford ha dado un ejemplo (no publicado) con una \emph{superficie} algebraica.}. Señalemos simplemente que, por el teorema de Cohen, existe un $p$-anillo de Cohen~$B$ de cuerpo residual~$k$ y un homomorfismo~$B\to A$ que induce el isomorfismo identidad sobre los cuerpos residuales; por consiguiente, el resultado «más fuerte» de levantamiento se obtendría tomando por~$A$
\marginpar{84}
un $p$-anillo de Cohen: si existe una solución (ordinaria o formal) sobre un anillo de este tipo, existe una sobre todo anillo local completo de cuerpo residual~$k$. En particular, como para un $p$-anillo de Cohen~$\goth{m}/\goth{m}^2$ se identifica canónicamente con~$k$, se ve que \emph{para todo esquema liso~$X_0$ sobre un cuerpo~$k$ de característica~$p>0$, existe una clase de cohomología en~$\H^2(X_0,\goth{g}_{X_0/k})$}, primera obstrucción al levantamiento de~$X_0$ a un esquema liso sobre un $p$-anillo de Cohen; se ignora si puede ser no nula\footnote{Puede ser no nula, como se señala en la nota~\Ref{nIII.3} de la página~\pageref{nIII.3}.}.

Incluso si se consigue construir paso a paso los~$X_n$ que se reducen a~$X_0$, esto no da, en general, más que un esquema \emph{formal}~$\goth{X}$ liso sobre~$A$, que se reduce a~$X_0$. Cuando~$X_0$ es propio sobre
\ifthenelse{\boolean{orig}}{~$A$}{~$k$}, queda por determinar si~$\goth{X}$ es de hecho algebraizable, para poder obtener un esquema \emph{ordinario} propio sobre~$A$ y
\ifthenelse{\boolean{orig}}{simple}{liso}
sobre~$A$, que se reduzca a~$X_0$. El único criterio conocido (señalado en el Seminario Bourbaki y que figura en los \emph{Éléments}, cap\ptbl III, 4.7.1) es el siguiente: si~$\goth{X}$ es propio sobre~$A$, y si~$\cal{L}$ es un haz invertible sobre~$\goth{X}$ tal que el haz inducido~$\cal{L}_0$ sobre~$X_0$ sea amplio (es decir, una potencia tensorial conveniente~$\cal{L}_0^{\otimes n}$, $n>0$, procede de una inmersión proyectiva de~$X_0$), entonces existen un esquema~$X$ proyectivo sobre~$A$ y un haz invertible amplio sobre~$X$, tales que el par~$(\goth{X},\cal{L})$ se obtenga de ellos por completación $\goth{m}$-ádica. Esto nos lleva, dado un haz localmente libre~$\cal{E}_0$ sobre~$X_0$ (que elegiremos invertible y amplio para nuestro propósito), a prolongarlo a un haz localmente libre~$\cal{E}$ sobre~$\goth{X}$. Para ello, se reduce la cuestión a construir paso a paso haces localmente libres~$\cal{E}_n$ sobre los~$X_n$. La discusión es completamente análoga a la del número~\Ref{III.6} (cf. la observación~\Ref{III.6.4}), y el papel esencial lo desempeña el \emph{haz de automorfismos} de un~$\cal{E}_{n+1}$ que inducen la identidad sobre~$\cal{E}_n$: se demuestra inmediatamente que este haz se identifica con
\begin{equation*}
\label{eq:III.7.*}\tag{$*$}{}
\cal{G}=\SheafHom_{\cal{O}_{X_0}}(\cal{E}_0,
\cal{E}_0\otimes\gr_{\cal{I}}^{n+1}(\cal{O}_X))=
\SheafHom_{\cal{O}_{X_0}}(\cal{E}_0,
\cal{E}_0)\otimes\gr_{\cal{I}}^{n+1}(\cal{O}_X)
\end{equation*}
que sigue siendo un haz de grupos conmutativos. Se obtiene:


\begin{proposition}
\label{III.7.1}
Sea~$S$ un preesquema provisto de un haz cuasicoherente de ideales~$\cal{I}$, sea~$X$ un preesquema sobre~$S$, sea~$S_n$ el subpreesquema de~$S$ definido por~$\cal{I}^{n+1}$, y sea~$X_n=X\times_SS_n$ (para todo entero~$n$). Sea~$\cal{E}_n$ un haz localmente libre
\marginpar{85}
sobre~$X_n$; se pretende prolongarlo a un haz localmente libre~$\cal{E}_{n+1}$ sobre~$X_{n+1}$. Entonces~$\cal{E}_n$ define una clase de obstrucción canónica en~$\H^2(X_0,\cal{G})$, donde~$\cal{G}$ es el haz cuasicoherente dado por la fórmula anterior, clase cuya anulación es necesaria y suficiente para la existencia de un~$\cal{E}_{n+1}$ que prolongue~$\cal{E}_n$. Si esta clase es nula, el conjunto de las clases de soluciones~$\cal{E}_{n+1}$, salvo isomorfismo que induzca la identidad sobre~$\cal{E}_n$, es un espacio homogéneo principal bajo~$\H^1(X_0,\cal{G})$.
\end{proposition}

Esta proposición da lugar a los corolarios habituales. Señalemos únicamente que, si~$X$ es \emph{plano} sobre~$S$, puede escribirse
$$
\cal{G}=\SheafHom_{\cal{O}_{X_0}}(\cal{E}_0,\cal{E}_0)
\otimes_{\cal{O}_{S_0}}\gr_{\cal{I}}^{n+1}(\cal{O}_S),
$$
de donde, si~$S$ es afín de anillo~$A$ y si los~$\cal{I}^n/\cal{I}^{n+1}$ son localmente libres, se obtiene la condición suficiente
$$
\H^2(X_0,\cal{G}_0)=0\qquad\text{con}\qquad
\cal{G}_0=\SheafHom_{\cal{O}_{X_0}}(\cal{E}_0,\cal{E}_0)
$$
para la existencia de un~$\cal{E}_{n+1}$ y, por tanto, paso a paso, para la existencia de prolongaciones sucesivas~$\cal{E}_m$ ($m=n,n+1,$
\ifthenelse{\boolean{orig}}{etc...}{etc.}).

Volviendo a la situación inicial, se obtiene:

\begin{proposition}
\label{III.7.2}
Sean~$A$ un anillo local completo y~$\goth{X}$ un esquema formal propio y plano sobre~$A$, tal que~$X_0$ sea proyectivo y que~$\H^2(X_0,\cal{O}_{X_0})=0$. Entonces existe un esquema~$X$ proyectivo sobre~$A$ cuyo completado formal $\goth{m}$-ádico es isomorfo a~$\goth{X}$.
\end{proposition}

Combinando este resultado con~\Ref{III.6.10}, se obtiene:

\begin{theoreme}
\label{III.7.3}
Sean~$A$ un anillo local completo de cuerpo residual~$k$ y~$X_0$ un esquema proyectivo y liso sobre~$k$, tal que
$$
\H^2(X_0,\goth{g}_{X_0/k})=\H^2(X_0,\cal{O}_{X_0})=0.
$$
Entonces existe un esquema liso y proyectivo~$X$ sobre~$A$ que se reduce a~$X_0$.
\end{theoreme}

Más generalmente, si se da un~$X_n$ liso sobre~$A_n=A/\goth{m}^{n+1}$ que se reduce a~$X_0$, entonces existe un~$X$ liso y propio sobre~$A$ y un isomorfismo~$X\otimes_AA_n=X_n$.

\begin{corollaire}
\label{III.7.4}
Toda curva lisa y propia sobre~$k$ se obtiene por reducción de una curva lisa y propia sobre~$A$.
\end{corollaire}

Este resultado será el instrumento esencial (junto con el teorema de existencia para haces en geometría formal) para estudiar el grupo fundamental de~$X_0$ por métodos trascendentes.
\marginpar{86}


\chapter{Morfismos planos}\label{IV}
\marginpar{87}

Exponemos aquí principalmente las propiedades de planitud que hemos utilizado en los exposés precedentes. Se encontrará un estudio más detallado en el capítulo~IV de los «Éléments de Géométrie Algébrique» en preparación\footnote{\Cf EGA IV 11 y 12.}, donde se estudia sistemáticamente la situación siguiente: dado~$X$ localmente de tipo finito sobre~$Y$ localmente noetheriano, y~$\cal{F}$ coherente sobre~$X$ y plano sobre~$Y$, establecer relaciones entre las propiedades de~$Y$, las de~$\cal{F}$ y las de los haces coherentes inducidos por~$\cal{F}$ sobre las fibras de~$X\to Y$ (en particular desde el punto de vista de la dimensión, de la dimensión cohomológica, de la profundidad, \ifthenelse{\boolean{orig}}{etc...}{etc.}).
Se dispone, en particular, de un modo sistemático de obtener teoremas de tipo \emph{Seidenberg} o \emph{Bertini} (para las secciones hiperplanas). El resultado esencial para aplicar los métodos de planitud en este contexto es el siguiente (que se demostrará más adelante): si~$Y$ es íntegro, $X$ de tipo finito sobre~$Y$ y~$\cal{F}$ coherente sobre~$X$, existe un abierto no vacío~$U$ de~$Y$ tal que~$\cal{F}$ sea plano sobre~$Y$ en los puntos de~$X$ situados sobre~$U$. Una segunda manera, sin duda aún más importante, en que la planitud interviene en geometría algebraica es la \emph{teoría del descenso}: véanse, por ejemplo, los dos exposés de Grothendieck sobre el tema en el Séminaire Bourbaki\footnote{Y, para una exposición más detallada, los exposés~\Ref{VIII} y~\Ref{IX} más adelante.}. La planitud aparece así como una de las nociones técnicas centrales de la geometría algebraica.

Recordemos que Serre introdujo en GAGA las nociones de planitud y planitud fiel. También se encontrará una exposición de los números~\Ref{IV.1} y~\Ref{IV.2} siguientes en el Álgebra Conmutativa de Bourbaki (que, por supuesto, como indica el título del libro, no se limita al caso de anillos de base conmutativos)\footnote{N. Bourbaki, Algèbre Commutative, cap.~I (Modules Plats), Act. Sc. Ind. 1290, París, Hermann (1961).}.

A diferencia de los exposés precedentes, no suponemos que los anillos considerados sean necesariamente noetherianos.

\section{Sorites sobre los módulos planos}
\label{IV.1}

Se dice que un módulo~$M$ sobre el anillo~$A$ es \emph{plano}
\index{plano, fielmente plano (módulo)|hyperpage}%
(o plano sobre~$A$, si se quiere precisar~$A$)
\marginpar{88}
si el funtor
$$
T_M \colon N \mto M \otimes_{A}N
$$
(que es, en cualquier caso, exacto por la derecha) es
\ifthenelse{\boolean{orig}}
{\emph{exacto}}
{\emph{exacto},}
\ie transforma monomorfismos en monomorfismos. Esto equivale a decir que
\ifthenelse{\boolean{orig}}
{el primer funtor derivado por la derecha, o todos los funtores derivados por la derecha,}
{el primer funtor derivado por la izquierda, o todos los funtores derivados por la izquierda,}
son nulos, \ie que se tiene%
\ifthenelse{\boolean{orig}}
{}
{\enlargethispage{.5cm}}
$$
\Tor^{A}_{1}(M,N)=0 \qquad \text{para todo $N$,}
$$
\resp que se tiene
$$
\Tor^{A}_{i}(M,N)=0 \qquad \text{para $i>0$ y todo $N$.}
$$
Como los~$\Tor_i$ conmutan con los límites inductivos, basta además verificar estas condiciones para~$N$ de tipo finito, e incluso (tomando entonces una sucesión de composición de~$N$ con cocientes cíclicos) que se tenga
$$
\Tor^{A}_{1}(M,N)=0
$$
si~$N$ es cíclico, \ie de la forma~$A/I$, siendo~$I$ un ideal de~$A$. Obsérvese además que
$$
\Tor^{A}_{1}(M,A/I)=0 \quad\Longleftrightarrow\quad I\otimes_{A}M \rightarrow
M=A\otimes_{A}M\text{ es \emph{inyectivo}},
$$
como se ve en la sucesión exacta de los~$\Tor$, teniendo en cuenta que~$\Tor^{A}_{1}(M,A)=0$. Por tanto, que~$M$ sea plano equivale a que, para todo ideal~$I$, el homomorfismo natural
$$
I\otimes_{A}M \to IM
$$
sea un isomorfismo. Basta, además, verificarlo para~$I$ de tipo finito; a fortiori, basta verificar que el funtor~$M\otimes$ sea exacto sobre los módulos \emph{de tipo finito}.

Como ocurre siempre que se tiene un funtor exacto~$T$, si para un subobjeto~$N'$ de~$N$ se identifica~$T(N')$ con un subobjeto de~$T(N)$, se tiene
$$
T(N' \cap N'')=T(N') \cap T(N'')
$$
$$
T(N' + N'')=T(N') + T(N'')
$$
para dos subobjetos~$N'$, $N''$ de~$N$.

Una suma directa de módulos planos y un sumando directo de un módulo plano son planos. En particular, como~$A$ es plano, todo módulo \emph{libre}, y por tanto también todo módulo \emph{proyectivo}, es plano. El producto tensorial de dos módulos planos es plano y, si~$M$ es plano sobre~$A$, entonces~$M\otimes_AB$ es plano sobre~$B$ para todo cambio de base
\ifthenelse{\boolean{orig}}{$A \ B$}{$A \to B$}
(por la asociatividad del producto tensorial y porque un compuesto de funtores
\marginpar{89}
exactos es exacto). Si~$M$ es plano sobre~$B$ y~$B$ plano sobre~$A$, entonces~$M$ es plano sobre~$A$ (por la misma razón).

La sucesión exacta de los~$\Tor$, junto con la «conmutatividad» del~$\Tor$, da:

\begin{proposition}
\label{IV.1.1}
Sea~$0 \to M' \to M \to M'' \to 0$ una sucesión exacta de módulos sobre~$A$, siendo~$M''$ plano. Entonces
\begin{enumerate}
\item[(i)] esta sucesión sigue siendo exacta al tensorizar por cualquier módulo~$N$ sobre~$A$;
\item[(ii)] para que~$M$ sea plano es necesario y suficiente que~$M'$ lo sea.
\end{enumerate}
\end{proposition}

Puede decirse, por tanto, que desde el punto de vista del comportamiento respecto de los productos tensoriales, los módulos planos son «tan buenos» como los módulos libres o proyectivos (y que, en particular, la sucesión exacta de~\Ref{IV.1.1} es «tan buena» como si se escindiese).

Sea~$S$ un conjunto multiplicativamente cerrado de~$A$; entonces~$S^{-1}A$ es plano sobre~$A$, pues~$S^{-1}A\otimes N=S^{-1}N$ es un funtor exacto en~$N$. Si~$M$ es plano sobre~$A$, entonces~$S^{-1}M=S^{-1}A\otimes M$ es plano sobre~$S^{-1}A$; el recíproco es cierto si~$M\to S^{-1}M$ es un
\ifthenelse{\boolean{orig}}
{isomorfismo}
{isomorfismo,}
\ie \ifthenelse{\boolean{orig}}{Si}{si}
los~$s\in S$ actúan biyectivamente sobre~$M$ (por la transitividad de la planitud, ya que~$S^{-1}A$ es plano sobre~$A$). Más generalmente, el caso de un morfismo de preesquemas~$X\to Y$ y de un haz cuasicoherente~$\cal{F}$ sobre~$X$ cuya planitud respecto de~$Y$ se quiere estudiar conduce a la situación con dos anillos:

\begin{proposition}
\label{IV.1.2}
Sean~$A\to B$ un homomorfismo de anillos, $M$ un módulo sobre~$B$ y~$T$ un conjunto multiplicativamente cerrado de~$B$.
\begin{enumerate}
\item[(i)] Si~$M$ es plano sobre~$A$, entonces~$T^{-1}M$ es plano sobre~$A$ (y, por tanto, también sobre~$S^{-1}A$ para todo conjunto multiplicativamente cerrado~$S$ de~$A$ cuya imagen esté contenida en~$T$).
\item[(ii)] \ifthenelse{\boolean{orig}}
{Recíprocamente, si~$M_{\goth{n}}$ es plano sobre~$A_{\goth{n}}$ para todo ideal maximal~$\goth{n}$ de~$B$, $M_{\goth{n}}$ es plano sobre~$A$ (o, lo que equivale a lo mismo, sobre~$A_{\goth{m}}$, donde~$\goth{m}$ es el ideal primo de~$A$ imagen inversa de~$\goth{n}$), entonces~$M$ es plano sobre~$A$.}
{Recíprocamente, si para todo ideal maximal~$\goth{n}$ de~$B$ el módulo~$M_{\goth{n}}$ es plano sobre~$A$ (o, lo que equivale a lo mismo, sobre~$A_{\goth{m}}$, donde~$\goth{m}$ es el ideal primo de~$A$ imagen inversa de~$\goth{n}$), entonces~$M$ es plano sobre~$A$.}
\end{enumerate}
\end{proposition}

En efecto, se tiene la fórmula, funtorial respecto del módulo~$N$ sobre~$A$:
$$
T^{-1}M \otimes_{A} N=T^{-1}(M \otimes_{A} N)
$$
pues ambos miembros son funtorialmente isomorfos a~$T^{-1}B\otimes_BM\otimes_BN_{(B)}$ (con~$N_{(B)}=N\otimes_AB$), en virtud de las fórmulas de asociatividad de~$\otimes$. Se sigue inmediatamente que, si~$M\otimes_AN$ es exacto en~$N$, ocurre lo mismo con~$T^{-1}M\otimes_AN$ (como compuesto de dos funtores exactos), de donde~(i). Y se deduce asimismo~(ii), pues para verificar la exactitud de una sucesión de módulos sobre~$B$ basta verificar la exactitud de los localizados en todos los ideales maximales de~$B$.

\begin{proposition}
\label{IV.1.3}
\begin{enumerate}
\item[(i)] Sea
\marginpar{90}
$M$ un módulo plano sobre~$A$. Si~$x\in A$ no es divisor de~$0$ en~$A$, tampoco lo es en~$M$. En particular, si~$A$ es íntegro, $M$ es libre de torsión.
\item[(ii)] Supongamos que~$A$ es íntegro y que, para todo ideal maximal~$\goth{m}$ de~$A$, $A_{\goth{m}}$ es principal (por ejemplo, si~$A$ es un anillo de Dedekind, o incluso principal). Para que el módulo~$M$ sobre~$A$ sea plano, \ifthenelse{\boolean{orig}}{es necesario y suficiente}{es necesario y suficiente} que sea libre de torsión.
\end{enumerate}
\end{proposition}

Se obtiene~(i) observando que la multiplicación por~$x$ en~$M$ se obtiene tensorizando por~$M$ la multiplicación por~$x$ en~$A$. Para~(ii), por~\Ref{IV.1.2}(ii) puede suponerse ya que~$A$ es principal; hay que demostrar que, si~$M$ es libre de torsión, entonces para todo ideal~$I$ de~$A$ la inyección~$I\to A$, tensorizada por~$M$, es una inyección, lo que significa que el generador~$x$ de~$I$ no es divisor de~$0$ en~$M$. O.K.


\section{Módulos fielmente planos}
\label{IV.2}

Un funtor~$F$ de una categoría en otra se dice 
\emph{fiel} si, para cualesquiera~$X$, $Y$, la aplicación~$\Hom(X,Y)\to
\Hom(F(X),F(Y))$ es inyectiva. Si se trata de un funtor aditivo entre
categorías aditivas, equivale a decir que~$F(u)=0$ implica~$u=0$, y esto
implica que~$F(X)=0$ implica~$X=0$. Para que~$F$ sea \emph{fiel y exacto},
es necesario y suficiente que se cumpla la condición siguiente: para toda
sucesión~$M'\to M\to M''$ de morfismos en~$\cal{C}$, la sucesión transformada
\ifthenelse{\boolean{orig}}{$F(M)\to F(M')\to F(M'')$}
{$F(M')\to F(M)\to F(M'')$}
es exacta \emph{si y solo si} lo es la anterior. O, lo que es lo mismo:
$F$ es exacto, y~$F(X)=0$ implica~$X=0$. (N.B.: para poder hablar de
exactitud hay que suponer que las categorías consideradas son
\emph{abelianas}). Supongamos dada una familia~$(M_i)$ de objetos no nulos
de~$\cal{C}$ tales que todo objeto no nulo de~$\cal{C}$ tenga un subobjeto
que admita un cociente isomorfo a algún~$M_i$. Entonces, que~$F$ sea fiel y
exacto equivale a que~$F$ sea exacto y~$F(M_i)\neq0$ para todo~$i$. Si
$\cal{C}$ es la categoría de los módulos sobre un anillo~$A$, puede tomarse,
por ejemplo, como~$(M_i)$ la familia de los~$A/\goth{m}$, donde~$\goth{m}$
recorre los ideales maximales de~$A$. (En efecto, todo módulo no nulo admite
un submódulo cíclico no nulo, por tanto isomorfo a un~$A/I$, con~$I$ un
ideal distinto de~$A$, el cual, por Krull, admite un cociente~$A/\goth{m}$).
De estos sorites se deduce, en particular:

\begin{proposition}
\label{IV.2.1}
Sea~$M$ un $A$-módulo. Son equivalentes las condiciones siguientes:
\begin{enumerate}
\item[(i)] El funtor~$M\otimes_A$ es fiel y exacto.
\item[(i~bis)]
\marginpar{91}
$M$ es plano, y~$M\otimes_A N=0$ implica~$N=0$.
\item[(i~ter)] $M$ es plano, y~$M\otimes_A A/\goth{m}\neq0$ para todo ideal
maximal~$\goth{m}$ de~$A$.
\item[(ii)] Para toda sucesión de homomorfismos~$N'\to N\to N''$, la
sucesión tensorizada por~$M$ es exacta si y solo si lo es la sucesión
inicial.
\end{enumerate}
\end{proposition}

Se dice entonces que~$M$ es un $A$-módulo \emph{fielmente plano}.
\index{fielmente plano (módulo)|hyperpage}%
\index{plano, fielmente plano (módulo)|hyperpage}%
En particular, si~$M$ es fielmente plano, entonces~$N\to N'$ es un
monomorfismo (epimorfismo, isomorfismo) si y solo si lo es el homomorfismo
tensorizado por~$M$. Un módulo fielmente plano es \emph{fiel}, puesto que la
multiplicación por~$f$ en~$M$ se obtiene tensorizando por~$M$ la
multiplicación por~$f$ en~$A$.

\ifthenelse{\boolean{orig}}
{Como en~\Ref{IV.1}, se obtienen}
{Como en~\Ref{IV.1}, se obtienen}
las propiedades habituales de transitividad: el producto tensorial de dos
módulos fielmente planos es fielmente plano; si~$M$ es fielmente plano
sobre~$A$, $M\otimes_A B$ es fielmente plano sobre~$B$ para toda extensión
de la base~$A\to B$; si~$B$ es una $A$-álgebra fielmente plana sobre~$A$ y
$M$ es un $B$-módulo fielmente plano, entonces~$M$ es un $A$-módulo
fielmente plano.

\begin{corollaire}
\label{IV.2.2}
Sean~$A\to B$ un homomorfismo local de anillos locales y~$M$ un $B$-módulo
de tipo finito. Para que~$M$ sea fielmente plano sobre~$A$, es necesario y
suficiente que sea plano sobre~$A$ y no nulo.
\end{corollaire}

Esto resulta del criterio~(i~ter) y de Nakayama. En particular,
\emph{para que~$B$ sea plano sobre~$A$, es necesario y suficiente que sea
fielmente plano sobre~$A$}.

\begin{proposition}
\label{IV.2.3}
Sean~$A\to B$ un homomorfismo de anillos y~$M$ un $B$-módulo fielmente plano
sobre~$A$. Para todo ideal primo~$\goth{p}$ de~$A$ existe un ideal
primo~$\goth{q}$ de~$B$ cuya contracción a~$A$ es~$\goth{p}$.
\end{proposition}

Al pasar a cocientes por~$\goth{p}$, se reduce al caso~$\goth{p}=0$.
Localizando en el ideal primo~$0$, se reduce al caso en que~$A$ es un
cuerpo. Pero, como~$M$ es fielmente plano sobre~$A$, no es nulo; con mayor
razón, $B\neq0$, de modo que~$B$ tiene un ideal primo, cuya contracción solo
puede ser el único ideal primo de~$A$. Geométricamente, puede decirse que la
existencia de un haz cuasicoherente~$\cal{F}$ sobre~$X=\Spec(B)$ que sea
«fielmente plano» sobre~$A$ implica que~$X\to Y=\Spec(A)$ es
\emph{sobreyectivo}.

\begin{corollaire}
\label{IV.2.4}
Supongamos que~$M$ es plano sobre~$A$, de tipo finito sobre~$B$ y
\ifthenelse{\boolean{orig}}{sup.~$M=\Spec(B)$}{$\Supp M=\Spec(B)$}
(\ie $M_{\goth{q}}\neq0$ para todo ideal primo~$\goth{q}$ de~$B$). Entonces,
los ideales primos~$\goth{q}$ de~$B$ que contienen~$\goth{p}B$ y son
minimales entre ellos se contraen a~$\goth{p}$.
\end{corollaire}

\marginpar{92}
Se reduce de nuevo al caso~$\goth{p}=0$ (pues todas las hipótesis se
conservan al pasar a cocientes), y por tanto~$A$ es íntegro. Se reduce así al
enunciado siguiente:

\begin{corollaire}
\label{IV.2.5}
Si~$M$ es como antes, todo ideal primo minimal~$\goth{q}$ de~$B$ se contrae
a un ideal primo~$\goth{p}$ de~$A$ que es minimal.
\end{corollaire}

En efecto, localizando en~$\goth{p}$ y~$\goth{q}$, se reduce a demostrar que,
si~$A$ y~$B$ son locales, el homomorfismo~$A\to B$ es local, $M$ es un
$B$-módulo no nulo plano sobre~$A$ y~$B$ tiene dimensión~$0$, entonces~$A$
tiene dimensión~$0$. Por~\Ref{IV.2.2} y~\Ref{IV.2.3}, se concluye que todo
ideal primo de~$A$ es la contracción de un ideal primo de~$B$, por tanto del
ideal maximal
\ifthenelse{\boolean{orig}}{de~$A$,}{de~$B$,}
y por consiguiente es el ideal maximal, como se quería. Geométricamente,
\Ref{IV.2.5} significa que toda componente irreducible de
\ifthenelse{\boolean{orig}}{$\overline{X}=\Spec(B)$}{$X=\Spec(B)$}
domina alguna componente irreducible de~$Y=\Spec(A)$ (si existe un haz
cuasicoherente de tipo finito sobre~$X$, con soporte~$X$, y plano sobre~$Y$).

Obsérvese que en~\Ref{IV.2.4} no fue necesario suponer que~$M$ fuera
fielmente plano sobre~$A$, pero entonces nada garantiza la existencia de un
ideal primo que contenga~$\goth{p}B$ (y, por tanto, de uno minimal entre
ellos).

\begin{proposition}
\label{IV.2.6}
Sea~$i\colon A\to B$ un homomorfismo de anillos. Son equivalentes las
condiciones siguientes:
\begin{enumerate}
\item[(i)] $B$ es un $A$-módulo fielmente plano.
\item[(ii)] $B$ es plano sobre~$A$, y~$\Spec(B)\to\Spec(A)$ es sobreyectivo.
\item[(ii~bis)] $B$ es plano sobre~$A$, y
\ifthenelse{\boolean{orig}}
{todo ideal maximal es la contracción de un ideal de~$B$.}
{todo ideal maximal de~$A$ es la contracción de un ideal de~$B$.}
\item[(iii)] $i$ es inyectivo y~$\Coker i$ es un $A$-módulo plano.
\item[(iv)] El funtor~$M\mapsto M_{(B)}=M\otimes_A B$ sobre los $A$-módulos
es exacto, y el homomorfismo funtorial canónico~$M\to M_{(B)}$ es inyectivo.
\item[(iv~bis)] Para todo ideal~$I$ de~$A$, $I\otimes_A B\to IB$ es un
isomorfismo, y la imagen inversa de~$IB$ en~$A$ es igual a~$I$.
\end{enumerate}
\end{proposition}

Se tiene (i)$\To$(ii) por~\Ref{IV.2.3}; (ii)$\To$(ii~bis) es trivial, y
(ii~bis)$\To$(i) por el criterio~(i~ter) de~\Ref{IV.2.1}. Se tiene
(iii)$\To$(iv) por~\Ref{IV.1.1}; (iv)$\To$(iv~bis) trivialmente (tomando
$M=A/I$ en la segunda condición
\ifthenelse{\boolean{orig}}{(iv~bis)}{(iv)}),
y (iv~bis)$\To$(i) en virtud del criterio de planitud por ideales visto al
comienzo de~\Ref{IV.1} y del criterio~\Ref{IV.2.1}~(i~ter). Por último,
(iv)$\To$(iii) por un recíproco fácil de~\Ref{IV.1.1}, y (i)$\To$(iv) porque
\marginpar{93}
si~$N$ es el núcleo de~$M\to M\otimes_A B=T(M)$, entonces (como~$T$ es
exacto) $N\to T(N)$ es nulo, de donde~$T(N)=N\otimes_A B=0$, y por tanto
$N=0$, como se quería.


\section{Relaciones con la completación}
\label{IV.3}

Sean~$A$ un anillo noetheriano, $I$ un ideal de~$A$,
$\widehat{A}$ la completación separada de~$A$ para la topología
$I$-preádica y, para todo $A$-módulo~$M$, sea~$\widehat{M}$ su
completación para la topología $I$-preádica. Este es un
$\widehat{A}$-módulo, de donde se obtiene un homomorfismo canónico
$$
M\otimes_A\widehat{A}\to\widehat{M}.
$$
Cuando~$M$ recorre los módulos \emph{de tipo finito}, el funtor
$M\mto\widehat{M}$ es exacto, como resulta fácilmente del
\emph{teorema de Krull: si~$N\subset M$, la topología de~$N$ es la
inducida por la topología de~$M$}. Como el funtor
$M\mto M\otimes_A\widehat{A}$ es exacto por la derecha, se concluye
fácilmente (resolviendo~$M$ mediante~$L\to L'\to M$, con~$L$ y~$L'$
libres de tipo finito) que el homomorfismo funtorial anterior es un
\emph{isomorfismo} (pues el funtor de completación también es exacto por
la derecha) y, por consiguiente, que~$M\mto M\otimes_A\widehat{A}$ es
también un funtor \emph{exacto}. Por tanto:

\begin{proposition}
\label{IV.3.1}
Sean~$A$ un anillo noetheriano e~$I$ un ideal de~$A$. Entonces la
completación separada~$\widehat{A}$ de~$A$ para la topología
$I$-preádica es \emph{plana} sobre~$A$.
\end{proposition}

\begin{corollaire}
\label{IV.3.2}
Para que~$\widehat{A}$ sea fielmente plano sobre~$A$, es necesario y
suficiente que~$I$ esté contenido en el radical de Jacobson de~$A$.
\end{corollaire}

En efecto, basta aplicar el criterio~\Ref{IV.2.1}~(i~ter).

Estos resultados resumen todo cuanto puede decirse, desde el punto de vista
del álgebra lineal, sobre las relaciones entre~$A$ y~$\widehat{A}$. El
corolario~\Ref{IV.3.2} se utiliza sobre todo cuando~$A$ es un anillo local
noetheriano e~$I$ está contenido en el ideal maximal~$\goth{m}$ (y, casi
siempre, es igual a él).



\section{Relaciones con los módulos libres}
\label{IV.4}

\begin{proposition}
\label{IV.4.1}
Sean~$A$ un anillo, $I$ un ideal de~$A$ y~$M$ un $A$-módulo.
Supongamos que se cumple una de las dos hipótesis siguientes:
\begin{enumerate}
\item[(a)] $I$ es nilpotente.
\item[(b)] $A$ es noetheriano, $I$ está contenido en el radical de
Jacobson de~$A$ y~$M$ es de tipo finito.
\end{enumerate}
Para que~$M$ sea libre sobre~$A$, es necesario y suficiente que
$M\otimes_A A/I$ sea libre sobre~$A/I$ y que
$\Tor_1^A(M,A/I)=0$.
\end{proposition}

La necesidad es clara;
\marginpar{94}
demostremos la suficiencia. Sea~$(e_i)$ una familia
de elementos de~$M$ cuya imagen en~$M\otimes_A A/I=M/IM$ define una base
sobre~$A/I$ (la familia es finita en el caso~(b)). Sea~$L$ el $A$-módulo
libre construido sobre el mismo conjunto de índices. Se tiene entonces un
homomorfismo~$L\to M$ tal que el funtor~$T$ de tensorización por~$A/I$
induce un isomorfismo~$T(L)\isomto T(M)$. Si~$Q$ es el conúcleo de
$L\to M$, se tiene~$T(Q)=0$, de donde~$Q=0$ por Nakayama (válido bajo
cualquiera de las condiciones~(a) y~(b)). Por tanto, $L\to M$ es
sobreyectivo. Sea~$R$ su núcleo; se tiene una sucesión exacta
$$
0\to R\to L\to M\to0.
$$
Como~$\Tor_1^A(M,A/I)=0$, se obtiene una sucesión exacta
$0\to T(R)\to T(L)\to T(M)\to0$, de donde~$T(R)=0$ y, de nuevo,
\ifthenelse{\boolean{orig}}{$T(R)=0$}{$R=0$}
por Nakayama (teniendo en cuenta que, en el caso~(b), $R$ es de tipo
finito porque~$A$ se supuso noetheriano).

\ifthenelse{\boolean{orig}}
{\begin{corollairestar}}
{\begin{corollaire}\label{IV.4.2}}
La condición~$\Tor_1^A(M,A/I)=0$ puede sustituirse por la siguiente: el
homomorfismo canónico sobreyectivo
\begin{equation*}
\ifthenelse{\boolean{orig}}
{\gdef\IVFourStarLabel{eq:IV.2.*}}
{\gdef\IVFourStarLabel{eq:IV.4.*}}
\label{\IVFourStarLabel}
\tag{$*$} {\gr_I^0(M)\otimes_{A/I}\gr_I(A)\to\gr_I(M)}
\end{equation*}
es un isomorfismo.
\ifthenelse{\boolean{orig}}
{\end{corollairestar}}
{\end{corollaire}}

En efecto, si~$M$ es libre, esto se cumple sin duda. Hay que demostrar,
pues, que si~$M\otimes_A A/I$ es libre sobre~$A/I$ y se cumple la condición
sobre los graduados asociados, entonces~$M$ es libre. Se retoma la
demostración anterior construyendo~$L\to M$. La hipótesis implica que este
homomorfismo induce un isomorfismo entre los graduados asociados; por
tanto, su núcleo está contenido en la intersección de los~$I^nL$ y es nulo
(como es trivial en~(a) y bien conocido en~(b)). c.q.d.

\ifthenelse{\boolean{orig}}{\refstepcounter{subsection}}{}

\begin{corollaire}
\label{IV.4.3}
Supongamos que~$A/I$ es un cuerpo. Entonces son equivalentes las condiciones
siguientes sobre~$M$:
\begin{enumerate}
\item[(i)] $M$ es libre.
\item[(ii)] $M$ es proyectivo.
\item[(iii)] $M$ es plano.
\item[(iv)] $\Tor_1^A(M,A/I)=0$.
\item[(v)] El homomorfismo canónico~\eqref{\IVFourStarLabel}
es biyectivo.
\end{enumerate}
\end{corollaire}

En efecto, en el caso considerado, $M\otimes_A A/I$ es automáticamente
libre.

El resultado anterior es válido en los dos casos siguientes:
\begin{enumerate}
\item[(a)] $M$ es un módulo \emph{arbitrario} sobre un anillo local~$A$
cuyo ideal maximal~$I$ es \emph{nilpotente} (por ejemplo, un anillo local
artiniano).
\item[(b)] $M$ es un módulo \emph{de tipo finito} sobre un anillo
\emph{local noetheriano}.
\end{enumerate}

Recordemos, a modo de referencia:
\marginpar{95}
\begin{corollaire}
\label{IV.4.4}
Supongamos que~$A$ es un dominio \emph{local noetheriano}, de ideal
maximal~$\goth{m}=I$, cuerpo residual~$k=A/I$ y cuerpo de
fracciones~$K$. Sea~$M$ un módulo de tipo finito sobre~$A$. Entonces las
condiciones equivalentes anteriores~\textup{(i)} a~\textup{(v)} son
también equivalentes a
\ifthenelse{\boolean{orig}}{}
{\enlargethispage{.5cm}}%
\begin{enumerate}
\item[(vi)] $M\otimes_A K$ y~$M\otimes_A k$ son espacios vectoriales de
la misma dimensión (es decir, el rango de~$M$ sobre~$A$ es igual al número
mínimo de generadores del $A$-módulo~$M$).
\end{enumerate}
\end{corollaire}

La demostración es inmediata; se deja al lector la tarea de generalizar al
caso en que~$A$ solo se supone reducido: hay que exigir entonces que los
rangos de~$M$ en los ideales primos minimales de~$A$ sean iguales a la
dimensión del espacio vectorial~$M\otimes_A k$.


\section{Criterios locales de planitud}
\label{IV.5}

\begin{proposition}
\label{IV.5.1}
Sean~$A$ un anillo provisto de un ideal~$I$ y~$M$ un $A$-módulo.
Supongamos que
$$
\Tor_1^A(M,A/I^n)=0 \qquad \text{para $n>0$};
$$
entonces el homomorfismo canónico sobreyectivo
\begin{equation*}
\label{eq:IV.5.1.*}
\tag{$*$} {\gr_I^0(M)\otimes_{A/I}\gr_I(A)\to\gr_I(M)}
\end{equation*}
es un isomorfismo. El recíproco es cierto si~$I$ es nilpotente.
\end{proposition}

La hipótesis significa que los homomorfismos
$$
I^n\otimes_A M\to I^nM
$$
son isomorfismos, de donde se sigue inmediatamente que los homomorfismos
$$
I^n/I^{n+1}\otimes_A M\to I^nM/I^{n+1}M
$$
son isomorfismos. Recíprocamente, supongamos que se cumple esta condición
y que~$I$ es nilpotente; demostremos que~$\Tor_1^A(M,A/I^n)=0$ para todo
$n$. Esto es cierto para~$n$ grande; procedamos por inducción descendente
sobre~$n$, suponiéndolo demostrado para~$n+1$. Se tiene un diagrama
conmutativo
$$
\xymatrix{ & M\otimes I^{n+1}\ar[r]\ar[d] & M\otimes I^n\ar[r]
\ar[d] & M\otimes(I^n/I^{n+1})\ar[r]\ar[d] & 0 \\ 0\ar[r] &
MI^{n+1}\ar[r] & MI^n\ar[r] & MI^n/MI^{n+1}\ar[r] & 0 }
$$
cuyas filas son exactas.
\marginpar{96}
Por hipótesis, la última flecha vertical es un isomorfismo, y la hipótesis
de inducción significa asimismo que lo es la primera flecha vertical. Lo
mismo vale, por tanto, para la flecha vertical intermedia, lo que termina
la demostración.

Serre aisló la proposición siguiente durante el Seminario; permite
simplificaciones sustanciales en esta sección.
\begin{proposition}
\label{IV.5.2}
Sean~$A\to B$ un homomorfismo de anillos y~$M$ un $A$-módulo. Son
equivalentes las condiciones siguientes:
\begin{enumerate}
\item[(i)] Para todo $B$-módulo~$N$, se tiene $\Tor_1^A(M,N)=0$.
\item[(ii)] $\Tor_1^A(M,B)=0$ y $M_{(B)}=M\otimes_A B$ es plano
sobre~$B$.
\end{enumerate}
\end{proposition}

Se tiene un isomorfismo funtorial
$$
M\otimes_A N=(M\otimes_A B)\otimes_B N
$$
que expresa el primer miembro, considerado como funtor en~$M$, como
compuesto de los dos funtores $M\mto M\otimes_A B$ y
$P\mto P\otimes_B N$. Como el primero transforma módulos libres sobre~$A$
en módulos libres sobre~$B$, y por tanto proyectivos en proyectivos, se
tiene la sucesión espectral de los funtores compuestos
$$
\Tor_n^A(M,N)\From\Tor_p^B(\Tor_q^A(M,B),N),
$$
de donde resulta una sucesión exacta para los términos de grado bajo
$$
0\from\Tor_1^B(M\otimes_A B,N)\from\Tor_1^A(M,N)
\from\Tor_1^A(M,B)\ifthenelse{\boolean{orig}}{\otimes_A}{\otimes_B}N.
$$
Si se cumple~(i), entonces se concluye de esta sucesión exacta que
$\Tor_1^B(M\otimes_A B,N)=0$ para todo~$N$, es decir,
$M\otimes_A B$ es plano sobre~$B$, de donde~(ii). Si, recíprocamente, se
cumple~(ii), los términos que flanquean a~$\Tor_1^A(M,N)$ en la sucesión
exacta son nulos, y se obtiene~(i).

\begin{corollaire}
\label{IV.5.3}
Supongamos que~$B=A/I$; entonces las condiciones precedentes equivalen a
la siguiente:
\begin{enumerate}
\item[(iii)] $\Tor_1^A(M,N)=0$ para todo $A$-módulo~$N$ aniquilado por
una potencia de~$I$.
\end{enumerate}
\end{corollaire}

En efecto, (i) dice que esto vale cuando~$N$ está aniquilado por~$I$. Se
deduce~(iii) aplicando la hipótesis a los~$I^nN/I^{n+1}N$.

\begin{corollaire}
\label{IV.5.4}
Bajo las condiciones de~\Ref{IV.5.3}, las condiciones consideradas
implican que
\marginpar{97}
el homomorfismo funtorial
\begin{equation*}
\label{eq:IV.5.*}
\tag{$*$} {\gr_I^0(M)\otimes_{A/I}\gr_I(A)\to\gr_I(M)}
\end{equation*}
es un isomorfismo y que~$M\otimes_A A/I$ es plano sobre~$A/I$.
\end{corollaire}

Basta aplicar~(iii) y
\ifthenelse{\boolean{orig}}{\eqref{IV.5.1}}{\Ref{IV.5.1}}.
Utilizando el recíproco de~\Ref{IV.5.1} cuando~$I$ es nilpotente, se
obtiene:
\begin{corollaire}
\label{IV.5.5}
Sean~$A$ un anillo provisto de un ideal nilpotente~$I$ y~$M$ un
$A$-módulo. Son equivalentes las condiciones siguientes:
\begin{enumerate}
\item[(i)] $M$ es plano sobre~$A$.
\item[(ii)] $M\otimes_A A/I$ es plano sobre~$A/I$ y
$\Tor_1^A(M,A/I)=0$.
\item[(iii)] $M\otimes_A A/I$ es plano sobre~$A/I$ y el homomorfismo
canónico~\eqref{eq:IV.5.*} sobre los graduados asociados es un
isomorfismo.
\end{enumerate}
\end{corollaire}

Estas son, respectivamente, las condiciones (iii) y (ii) precedentes y
las condiciones del corolario~\Ref{IV.5.4}.

No supongamos ya que~$I$ sea nilpotente. Entonces, a priori, solo se
tienen en~\Ref{IV.5.5} las implicaciones
(i)$\To$(ii)$\To$(iii). Por otra parte, como la condición~(iii) permanece
estable al pasar al cociente por una potencia de~$I$, se ve, en virtud
de~\Ref{IV.5.5}, que implica
\begin{enumerate}
\item[(iv)] \emph{Para todo entero
\ifthenelse{\boolean{orig}}{$n$}{$n>0$},
$M\ifthenelse{\boolean{orig}}{\otimes}{\otimes_A}A/I^n$ es plano
sobre~$A/I^n$.}
\end{enumerate}
Nos proponemos dar condiciones bajo las cuales puede deducirse~(i), es
decir, que~$M$ es plano sobre~$A$. Afirmo que basta para ello que~$A$ sea
noetheriano y que~$M$ satisfaga la condición de finitud siguiente:
\emph{para todo módulo de tipo finito~$N$ sobre~$A$, $M\otimes_A N$ es
separado para la topología $I$-preádica}. (Bastaría verificarlo cuando~$N$
es un ideal de tipo finito de~$A$). En efecto, demostremos que, bajo estas
condiciones, si~$N'\to N$ es un monomorfismo de módulos de tipo finito,
$M\otimes_A N'\to M\otimes_A N$ es un monomorfismo. Basta mostrar que el
núcleo está contenido en todos los
$I^n(M\otimes_A N')=\Im(M\otimes_A I^nN'\to M\otimes_A N')$, o, lo que
equivale a lo mismo, en todos los
$\Im(M\otimes_A V'_n\to M\otimes_A N')=
\Ker(M\otimes_A N'\to M\otimes_A(N'/V'_n))$, donde~$V'_n$ recorre un
sistema fundamental numerable de entornos de~$0$ en~$N'$ (provisto de su
topología $I$-ádica). Por el teorema de Krull, la topología $I$-ádica
de~$N'$ es inducida por la de~$N$; se puede tomar, por tanto,
$V'_n=N'\cap I^nN$. Consideremos entonces el diagrama conmutativo
$$
\xymatrix{M\otimes_A N'\ar[r]\ar[d] & M\otimes_A(N'/V'_n)\ar[d]
\\ M\otimes_A N\ar[r] & M\otimes_A(N/I^nN)}.
$$
\marginpar{98}%
Como~$N'/V'_n$ y~$N/I^nN$ están aniquilados por~$I^n$, el segundo
homomorfismo vertical se identifica con el obtenido del homomorfismo
\emph{inyectivo} $N'/V'_n\to N/I^nN$ tensorizando sobre~$A/I^n$ con el
$(A/I^n)$-módulo \emph{plano} $M\otimes_A A/I^n$; es, por tanto,
\emph{inyectivo}. En consecuencia, el núcleo de
$M\otimes_A N'\to M\otimes_A N$ está contenido en el de
$M\otimes_A N'\to M\otimes_A(N'/V'_n)$, como se quería.

La condición de «finitud» considerada sobre~$M$ se cumple, en particular,
si~$M$ es un módulo de tipo finito sobre una $A$-álgebra noetheriana
\ifthenelse{\boolean{orig}}{,}{}~$B$ tal que~$IB$ esté contenido en el
radical de Jacobson de~$B$:
en efecto, entonces~$M\otimes_A N$ es un módulo de tipo finito sobre~$B$
para todo módulo de tipo finito~$N$ sobre~$A$, y por Krull es separado
para la topología $I$-ádica, que es su topología $(IB)$-ádica. Se obtiene
así:

\begin{theoreme}
\label{IV.5.6}
Sean~$A\to B$ un homomorfismo de anillos noetherianos, $I$ un ideal de~$A$
tal que~$IB$ esté contenido en el radical de Jacobson de~$B$ y~$M$ un
$B$-módulo de tipo finito. Son equivalentes las condiciones siguientes:
\begin{enumerate}
\item[(i)] $M$ es plano sobre~$A$.
\item[(ii)] $M\otimes_A A/I$ es plano sobre~$A/I$ y
$\Tor_1^A(M,A/I)=0$.
\item[(iii)] $M\otimes_A A/I$ es plano sobre~$A/I$ y el homomorfismo
canónico
$$
\gr_I^0(M)\otimes_{A/I}\gr_I(A)\to\gr_I(M)
$$
es un isomorfismo.
\item[(iv)] Para todo entero
\ifthenelse{\boolean{orig}}{$n$}{$n>0$}, $M\otimes_A A/I^n$ es plano
sobre~$A/I^n$.
\end{enumerate}
\end{theoreme}

Este resultado se aplica sobre todo cuando~$A$ y~$B$ son anillos
\emph{locales} noetherianos, $A\to B$ es un homomorfismo local e~$I$ es
un ideal de~$A$ contenido en su ideal maximal (y~\Ref{IV.5.6} puede
reducirse de inmediato a este caso). Un caso interesante es aquel en que
$A/I$ es un cuerpo, es decir, $I$ es maximal; entonces la condición de que
$M\otimes_A A/I$ sea plano sobre~$A/I$ es automática. Además, como los
$A/I^n$ son en tal caso anillos locales artinianos, la condición~(iv)
significa que los~$M\otimes_A A/I^n$ son \emph{libres} sobre~$A/I^n$.

\begin{corollaire}
\label{IV.5.7}
Sea
\marginpar{99}
$A\to B$ un homomorfismo local de anillos locales noetherianos y sea
$u\colon M'\to M$ un homomorfismo de $B$-módulos de tipo finito.
Supongamos que~$M$ es plano sobre~$A$. Son equivalentes las condiciones
siguientes:
\begin{enumerate}
\item[(i)] $u$ es inyectivo y~$\Coker u$ es plano sobre~$A$.
\item[(ii)] $u\otimes_A k\colon M'\otimes_A k\to M\otimes_A k$ es
inyectivo.
\end{enumerate}
(Aquí~$k$ denota el cuerpo residual de~$A$).
\end{corollaire}

(i)$\To$(ii) en virtud de~\Ref{IV.1.1}; demostremos el recíproco. En
primer lugar, $u$ es inyectivo, pues basta verificarlo en los graduados
asociados, donde resulta de un cuadrado conmutativo que el lector
escribirá. Sea~$M''$ su conúcleo; se tiene, por tanto, una sucesión exacta
$$
0\to M'\to M\to M''\to0,
$$
de donde la sucesión exacta de los~$\Tor$, teniendo en cuenta la
hipótesis~(ii) y~$\Tor_1^A(M,k)=0$, da la relación
$\Tor_1^A(M'',k)=0$; por el teorema~\Ref{IV.5.6}, $M''$ es plano
sobre~$A$.

\begin{corollaire}
\label{IV.5.8}
Bajo las condiciones de~\Ref{IV.5.6}, sea~$J$ un ideal de~$B$ que
contiene a~$IB$ y está contenido en el radical de Jacobson. Sean
$\hat A$ la completación
\ifthenelse{\boolean{orig}}{$I$-ádicas}{$I$-ádica} de~$A$ y~$\hat B$
y~$\hat M$ las completaciones $J$-ádicas de~$B$ y~$M$. Para que~$M$ sea
plano sobre~$A$, es necesario y suficiente que
\ifthenelse{\boolean{orig}}{$M$}{$\hat M$} sea plano sobre~$\hat A$.
\end{corollaire}

(N.B. la suficiencia resultaría ya fácilmente de~\Ref{IV.3.2}). Se
utiliza el criterio~(iii) de~\Ref{IV.5.6} en la situación~$(A,B,I,M)$ y
en la situación~$(\hat A,\hat B,I\hat A,\hat M)$. Se comprueba que las
condiciones obtenidas en ambos casos son equivalentes gracias
a~\Ref{IV.3.2}.

\begin{corollaire}
\label{IV.5.9}
Sean~$A\to B\to C$ homomorfismos locales de anillos locales
noetherianos y~$M$ un $C$-módulo de tipo finito (N.B. $C$ interviene
únicamente para poder imponer una condición de finitud a~$M$).
Supongamos que~$B$ es plano sobre~$A$. Sea~$k$ el cuerpo residual de~$A$.
Son equivalentes las condiciones siguientes:
\begin{enumerate}
\item[(i)] $M$ es plano sobre~$B$.
\item[(ii)] $M$ es plano sobre~$A$ y~$M\otimes_A k$ es plano
sobre~$B\otimes_A k$.
\end{enumerate}
\end{corollaire}

La implicación (i)$\To$(ii) es trivial; demostremos (ii)$\To$(i). Se
aplica el criterio~(iii) de~\Ref{IV.5.6} a~$(B,C,\goth{m}B=I,M)$. Como
$M\otimes_B(B/I)=M\otimes_B(B\otimes_A k)=M\otimes_A k$, la primera
condición de este criterio significa precisamente que~$M\otimes_A k$ es
plano sobre~$B\otimes_A k$, perfecto.
\marginpar{100}
La segunda condición del criterio se cumple porque~$M$ es plano sobre~$A$
y~$B$ es plano sobre~$A$, por una fórmula de asociatividad del producto
tensorial. --- Naturalmente, remitiéndose a~\Ref{IV.5.5} en lugar de
a~\Ref{IV.5.6}, se obtiene un enunciado análogo sin condición noetheriana
ni de finitud, si se supone, en cambio, que el ideal~$\goth{m}$ de~$A$ es
nilpotente. (El hecho de que se haya tomado~$\goth{m}$ maximal no ha
intervenido; pero, en cierto sentido, el caso «$\goth{m}$ maximal» es «el
mejor posible»).


\section{Morfismos planos y conjuntos abiertos}
\label{IV.6}

Recordemos primero algunos resultados sobre los conjuntos constructibles,
demostrados también en unas notas en circulación del Seminario
Dieudonné--Rosenlicht sobre los esquemas\footnote{\Cf EGA~0$_\mathrm{III}$~9,
EGA~IV~1.8 y~1.10.}.

Sea~$X$ un espacio topológico. Se dice, siguiendo a Chevalley, que un
subconjunto de~$X$ es \emph{constructible}
\index{constructible (subconjunto)|hyperpage}%
si es una unión finita de subconjuntos localmente cerrados.

\begin{lemme}
\label{IV.6.1}
Sea~$X$ un espacio topológico noetheriano y sea~$Z$ un subconjunto de~$X$.
Para que~$Z$ sea constructible es necesario y suficiente que, para todo
subconjunto cerrado irreducible~$Y$ de~$X$, $Z\cap Y$
\ifthenelse{\boolean{orig}}{no es}{no sea} denso en~$Y$ o
\ifthenelse{\boolean{orig}}{contiene}{contenga} un subconjunto abierto no
vacío del espacio~$Y$.
\end{lemme}

De ello se deduce, utilizando un lema bien conocido de álgebra conmutativa:

\begin{lemme}[Chevalley]
\label{IV.6.2}
Sea $f\colon X\to Y$ un morfismo de tipo finito de preesquemas, con~$Y$
noetheriano. Entonces $f(X)$ es constructible.
\end{lemme}

\begin{lemme}
\label{IV.6.3}
Sean~$X$ un espacio topológico noetheriano en el que todo subconjunto cerrado
irreducible admite un punto genérico, $U$ un subconjunto constructible de~$X$ y
$x\in X$. Para que~$U$ sea un entorno de~$x$ es necesario y suficiente que
toda generalización~$y$ de~$x$ (\ie todo $y\in X$ tal que $x\in\bar y$)
pertenezca a~$U$.
\end{lemme}

En particular:

\begin{corollaire}
\label{IV.6.4}
Sea~$X$ un espacio topológico noetheriano en el que todo subconjunto cerrado
irreducible admite un punto genérico, y sea~$U$ un subconjunto de~$X$. Para que
$U$ sea abierto es necesario y suficiente que satisfaga las dos condiciones
siguientes:
\begin{enumerate}
\item[(a)] $U$ contiene toda generalización de cada uno de sus puntos;
\item[(b)] si $x\in U$, entonces
\marginpar{101}
$U\cap\bar x$ contiene un subconjunto abierto no vacío del espacio~$\bar x$.
\end{enumerate}
\end{corollaire}

En efecto, $U$ es necesariamente constructible gracias a~\Ref{IV.6.1}, y se
aplica el criterio~\ifthenelse{\boolean{orig}}{\Ref{IV.6.2}}{\Ref{IV.6.3}},
que demuestra que~$U$ es un entorno de cada uno de sus puntos.

\begin{corollaire}
\label{IV.6.5}
Sea $f\colon X\to Y$ un morfismo de tipo finito de preesquemas, con~$Y$
localmente noetheriano, sea~$x$ un punto de~$X$ y sea~$y=f(x)$. Para que~$f$
transforme \ifthenelse{\boolean{orig}}{toda}{todo} entorno de~$x$ en un
entorno de~$y$ es necesario y suficiente que, para toda generalización~$y'$
de~$y$, exista una generalización~$x'$ de~$x$ tal que $f(x')=y'$.
\end{corollaire}

Puede suponerse evidentemente que~$X$ e~$Y$ son afines y, por tanto,
noetherianos. La condición es suficiente, pues basta demostrar que $f(X)$ es
un entorno de~$y$; ahora bien, $f(X)$ es constructible
por~\ifthenelse{\boolean{orig}}{\Ref{IV.6.1}}{\Ref{IV.6.2}}, y basta aplicar
el criterio~\Ref{IV.6.3}.
La condición es necesaria: sea $Y'=\overline{y'}$ y sea~$F$ la unión de las
componentes irreducibles de~$f^{-1}(Y')$ que no contienen a~$x$. Entonces
$X-F$ es un entorno abierto de~$x$, por lo que su imagen es un entorno de~$y$
y, con mayor razón, contiene a~$y'$; existe, pues, $x'_1\in X-F$ tal que
$f(x'_1)=y'$. Consideremos una componente irreducible de~$f^{-1}(Y')$ que
contenga a~$x'_1$; contiene necesariamente a~$x$ (pues, de lo contrario,
estaría contenida en~$F$). Sea~$x'$ su punto genérico. Es una generalización
de~$x$, y $f(x')$ es una generalización de $f(x'_1)=y'$ contenida en~$Y'$;
por tanto, es igual a~$y'$, c.q.d.

\begin{theoreme}
\label{IV.6.6}
Sean $f\colon X\to Y$ un morfismo localmente de tipo finito, con~$Y$
localmente noetheriano, y~$F$ un haz coherente sobre~$X$, de soporte~$X$ y
plano con respecto a~$Y$. Entonces~$f$ es un morfismo abierto (\ie
transforma abiertos en abiertos).
\end{theoreme}

Basta demostrar el criterio~\Ref{IV.6.5} para todo punto $x\in X$. Ahora
bien, las generalizaciones~$x'$ de~$x$ corresponden a los ideales primos
de~$\cal O_x$, y las generalizaciones~$y'$ de~$y$ corresponden a los ideales
primos de~$\cal O_y$; hay que comprobar, por tanto, que todo ideal primo
de~$\cal O_y$ es la contracción de
\ifthenelse{\boolean{orig}}{una}{un} ideal primo de~$\cal O_x$.
Ahora bien, $F_x$ es un $\cal O_x$-módulo no nulo y plano
sobre~$\cal O_y$, por tanto fielmente plano sobre~$\cal O_y$
por~\Ref{IV.2.2}. Puede aplicarse entonces~\Ref{IV.2.3}, lo que concluye la
demostración.

\begin{remarquesstar}
Como la planitud se conserva por extensión de la base, se ve que, bajo las
condiciones de~\ifthenelse{\boolean{orig}}{\Ref{IV.6.5}}{\Ref{IV.6.6}}, $f$
es incluso \emph{universalmente abierto}.
Ignoro, sin embargo, si, cuando~$Y$ es íntegro y~$X$ es de tipo finito
sobre~$Y$, la restricción de~$f$ a toda componente~$X_i$ de~$X$ es un
morfismo abierto, o siquiera equidimensional\footnote{La respuesta a la segunda pregunta es
afirmativa y la respuesta a la primera es negativa incluso si~$f$ es étale;
\cf EGA~IV~12.1.1.5 y EGA~Err$_{\rm IV}$~33.}, \ie cuyas fibras tienen
todas sus componentes de la misma dimensión (solo se sabe que~$X_i$
\emph{domina} a~$Y$).
\marginpar{102}
La pregunta está relacionada con la siguiente: sea $A\to B$ un homomorfismo
local de anillos locales noetherianos, tal que~$B$ sea plano sobre~$A$ y
$\goth mB$ sea un ideal de definición de~$B$ (lo que implica, además, que
$\dim B=\dim A$). ¿Es cierto que, para todo ideal primo minimal~$\goth p_i$
de~$B$, se tiene $\dim B/\goth p_i=\dim B$? Señalemos únicamente que la
respuesta a la primera pregunta es negativa cuando se sustituye la hipótesis
de planitud de~\ifthenelse{\boolean{orig}}{\Ref{IV.6.5}}{\Ref{IV.6.6}} por
la sola hipótesis de que~$f$ sea universalmente abierto.
\end{remarquesstar}

\begin{lemme}
\label{IV.6.7}
Sean~$A$ un dominio noetheriano, $B$ un $A$-álgebra de tipo finito y~$M$ un
$B$-módulo de tipo finito. Entonces existe un elemento no nulo~$f$ de~$A$
tal que~$M_f$ sea un módulo libre (y, con mayor razón, plano) sobre~$A_f$.
\end{lemme}

Sea~$K$ el cuerpo de fracciones de~$A$; entonces $B\otimes_A K$ es un
álgebra de tipo finito sobre~$K$, y $M\otimes_A K$ un módulo de tipo finito
sobre esta última. Sea~$n$ la dimensión del soporte de este módulo;
razonaremos por inducción sobre~$n$.
Si $n<0$\ifthenelse{\boolean{orig}}{}{,}
\ie si $M\otimes_A K=0$, tomando un número finito de generadores de~$M$
sobre~$B$, se ve que existe un $f\in A$ que anula estos generadores y, por
tanto, a~$M$; de aquí $M_f=0$, y hemos terminado. Supongamos que $n\geq0$.
Se sabe que el $B$-módulo~$M$ admite una serie de composición cuyos
cocientes sucesivos son isomorfos a módulos $B/\goth p_i$, donde los
$\goth p_i$ son ideales primos de~$B$.
Como una extensión de módulos libres es libre, se reduce al caso en que el
propio~$M$ es de la forma~$B/\goth p$, o bien idéntico a~$B$, siendo~$B$ un
$A$-álgebra \emph{íntegra}. Aplicando el lema de normalización de Noether a
la $K$-álgebra $B\otimes_A K$, se ve fácilmente que existe un elemento no
nulo~$f$ de~$A$ tal que~$B_f$ sea entero sobre el subanillo
$A_f[t_1,\dots,t_n]$, donde los~$t_i$ son indeterminadas. Puede suponerse ya
\ifthenelse{\boolean{orig}}{}{que }$B$ es entero sobre
$C=A[t_1,\dots,t_n]$; es, por tanto, un $C$-módulo de tipo finito sin
torsión.
Sea~$m$ su rango; existe entonces una sucesión exacta de $C$-módulos:
$$
0\to C^m\to B\to M'\to0
$$
donde~$M'$ es un $C$-módulo de torsión. Se sigue que la dimensión de Krull
del $C\otimes_A K$-módulo $M'\otimes_A K$ es estrictamente menor que la
dimensión~$n$ de~$C\otimes_A K$. Por la hipótesis de inducción, después de
localizar con respecto a un elemento no nulo conveniente~$f$ de~$A$, puede
suponerse que~$M'$ es un $A$-módulo
\marginpar{103}
libre. Por otra parte, $C^m$ es un $A$-módulo libre. Por tanto, $B$ es
entonces un $A$-módulo libre, y hemos terminado.

\begin{lemme}
\label{IV.6.8}
Sean~$A$ un anillo noetheriano, $B$ un álgebra de tipo finito sobre~$A$,
$M$ un $B$-módulo de tipo finito, $\goth p$ un ideal primo de~$B$ y~$\goth q$
su contracción a~$A$. Supongamos que~$M_{\goth p}$ es plano
sobre~$A_{\goth q}$ (o sobre~$A$, lo que es lo mismo). Entonces existe un
$g\in B-\goth p$ tal que
\begin{enumerate}
\item[(a)] $(M/\goth qM)_g$ es plano sobre~$A/\goth q$.
\item[(b)] $\Tor_1^A(M,A/\goth q)_g=0$.
\end{enumerate}
\end{lemme}

En efecto, aplicando~\Ref{IV.6.7} a
$(A/\ifthenelse{\boolean{orig}}{\goth(q)}{\goth q},B/\goth qB,
M/\goth qM)$, se ve primero que existe un $f\in A-\goth q$ tal que
$(M/\goth qM)_f$ sea plano sobre~$A/\goth q$.
Por otra parte, como $M_{\goth p}$ es plano sobre~$A$, se tiene
$\Tor_1^A(M,A/\goth q)_{\goth p}=
\Tor_1^A(M_{\goth p},A/\goth q)=0$; puesto que
$\Tor_1^A(M,A/\goth q)$ es un $B$-módulo de tipo finito, existe un
$g\in B-\goth p$ tal que se cumple~(b). Puede entonces suponerse que se
cumple a la vez~(a), sustituyendo~$g$ por~$gf$, lo que demuestra el
\ifthenelse{\boolean{orig}}{corolario}{lema}.

\begin{corollaire}
\label{IV.6.9}
Con las notaciones de~\Ref{IV.6.8}, para todo ideal primo~$\goth p'$ de~$B$
que contenga a~$\goth p$ y no contenga a~$g$, $M_{\goth p'}$ es plano
sobre~$A$ (o, lo que es lo mismo, sobre~$A_{\goth q'}$, donde~$\goth q'$ es
la contracción de~$\goth p'$ a~$A$).
\end{corollaire}

Basta aplicar el criterio~\Ref{IV.5.6}~(ii) al sistema
$(A,B_{\ifthenelse{\boolean{orig}}{\goth q'}{\goth p'}},\goth q,
M_{\ifthenelse{\boolean{orig}}{\goth q'}{\goth p'}})$, utilizando la
localización de los~$\Tor$.

\begin{theoreme}
\label{IV.6.10}
Sea $f\colon X\to Y$ un morfismo de tipo finito, con~$Y$ localmente
noetheriano, y sea~$F$ un haz coherente sobre~$X$. Sea~$U$ el conjunto de
los puntos $x\in X$ tales que~$F_x$ sea plano sobre~$\cal O_{f(x)}$.
Entonces~$U$ es un conjunto \emph{abierto}.
\end{theoreme}

\subsubsection*{Demostración}
Puede suponerse que~$X$ e~$Y$ son afines, de anillos~$B$ y~$A$, y, por
tanto, que~$F$ está definido por un $B$-módulo~$M$ de tipo finito. Se aplica
el criterio~\Ref{IV.6.4}. La condición~(a) se satisface trivialmente
por~\Ref{IV.1.2}~(i); queda por comprobar la condición~(b)
de~\Ref{IV.6.4}. Esto es lo que se ha hecho en el lema~\Ref{IV.6.8} y
\ifthenelse{\boolean{orig}}{}{el }corolario~\Ref{IV.6.9}.

En muchas cuestiones basta la forma más débil siguiente del
teorema~\Ref{IV.6.10} (que se sigue ya del lema~\Ref{IV.6.7} y no requiere,
por tanto, ni la técnica
\marginpar{104}
de los constructibles ni el teorema~\Ref{IV.5.6}):

\begin{corollaire}
\label{IV.6.11}
Bajo las condiciones de~\Ref{IV.6.10}, si se supone que~$Y$ es íntegro,
existe un abierto no vacío~$V$ de~$Y$ tal que~$F$ sea plano con respecto
a~$Y$ en todos los puntos de~$f^{-1}(V)$.
\end{corollaire}

En efecto, el conjunto abierto~$U$ contiene la fibra del punto genérico
de~$Y$ (pues el anillo local de este punto es un cuerpo); por tanto, contiene
un abierto de la forma~$f^{-1}(V)$, ya que~$X$ es de tipo finito sobre~$Y$.
De~\Ref{IV.6.11} se deduce también fácilmente el resultado siguiente, donde
se supone que~$Y$ es noetheriano (pero no necesariamente íntegro): existe
una partición de~$Y$ en subconjuntos localmente cerrados~$Y_i$ tales que,
dotando a~$Y_i$ de la estructura reducida inducida, el haz inducido por~$F$
en cada $X_i=X\times_Y Y_i$ sea plano con respecto a~$Y_i$.


\chapter{El grupo fundamental: generalidades}
\label{V}
\marginpar{105}

\setcounter{section}{-1}
\section{Introducción}

El presente Seminario es continuación del Seminario de 1960. Nos referiremos
a este último mediante siglas como \Ref{I.9.7}, que significa: Seminario de
Geometría Algebraica, exposición~I, \No9.7. La numeración de las exposiciones
de 1961 continuará la de 1960. Nos referiremos a los \emph{Elementos de
geometría algebraica} de Dieudonné--Grothendieck mediante siglas como
(EGA~I~8.7.3).

La presente exposición resume (con ligeros complementos) las últimas
exposiciones de 1960, que no habían sido redactadas.

Como en 1961, nos limitaremos por regla general a preesquemas localmente
noetherianos, aunque a menudo esta restricción sea inesencial. En la
exposición~\Ref{VI} daremos por supuesta la teoría del descenso fielmente
plano, resumida en el Seminario Bourbaki \No190. Si fuera necesario,
ofreceremos una exposición más detallada en una exposición
ulterior\footnote{\Cf exp\ptbl\ \Ref{VI} y exp\ptbl\ \Ref{VIII}}, una vez que
el lector haya tenido ocasión de convencerse de la utilidad de esta técnica
para la teoría del grupo fundamental.


\section{Preesquema con grupo finito de operadores; preesquema cociente}
\label{V.1}

Sean $X$ un preesquema y $G$ un grupo finito que actúa sobre~$X$ por
automorfismos, por la derecha para fijar ideas. Si $X$ es afín de anillo $A$,
$G$ actúa, por tanto, por automorfismos a la izquierda sobre~$A$.

Para todo preesquema $Z$, $G$ actúa a la izquierda sobre el conjunto
$\Hom(X,Z)$; por consiguiente, puede considerarse el conjunto
$$
\Hom(X,Z)^G
$$
de los morfismos invariantes por~$G$. Este conjunto depende funtorialmente
de~$Z$, y cabe preguntarse si este funtor es «representable», 
\ie isomorfo a un funtor $Z\mto\Hom(Y,Z)$.
\marginpar{106}

Esto significa que pueden hallarse un preesquema $Y$ y un morfismo invariante
por~$G$
$$
p\colon X\to Y
$$
tales que, para todo~$Z$, la aplicación correspondiente $g\mto gp$
$$
\Hom(Y,Z)\to \Hom(X,Z)^G
$$
sea biyectiva. Se dice entonces que $(Y,p)$ es un
\emph{preesquema cociente}
\index{cociente (preesquema)|hyperpage}%
de~$X$ \ifthenelse{\boolean{orig}}{por~$g$}{por~$G$}
(está determinado salvo isomorfismo único).

\begin{proposition}
\label{V.1.1}
Sean $A$ un anillo sobre el que el grupo finito $G$ actúa a la izquierda,
$B=A^G$ el subanillo de invariantes de~$A$, $X=\Spec(A)$ e $Y=\Spec(B)$, y
sea $p\colon X\to Y$ el morfismo canónico (evidentemente invariante
por~$G$). Entonces:
\begin{enumerate}
\item[(i)] $A$ es entero sobre~$B$, \ie $p$ es un morfismo
\emph{entero}.
\item[(ii)] El morfismo $p$ es sobreyectivo, sus fibras son las órbitas
de~$G$ y la topología de~$Y$ es el cociente de la de~$X$.
\item[(iii)] Sean $x\in X$, $y=p(x)$ y $G_x$ el estabilizador de~$x$.
Entonces
\ifthenelse{\boolean{orig}}
{$\kappa(x)$ es una extensión algebraica cuasigaloisiana de $\kappa(y)$ y
la aplicación canónica de~$G_x$ en el grupo
$\Gal(\kappa(x)/\kappa(y))$ de los $\kappa(y)$-automorfismos de
$\kappa(x)$}
{$\kres(x)$ es una extensión algebraica cuasigaloisiana de $\kres(y)$ y
la aplicación canónica de~$G_x$ en el grupo
$\Gal(\kres(x)/\kres(y))$ de los $\kres(y)$-automorfismos de
$\kres(x)$}
es \ifthenelse{\boolean{orig}}{sobreyectivo.}{sobreyectiva.}
\item[(iv)] $(Y,p)$ es un preesquema cociente de~$X$ por~$G$.
\end{enumerate}
\end{proposition}

Los enunciados
\ifthenelse{\boolean{orig}}
{(i) (ii) (iii)}
{(i), (ii), (iii)}
son bien conocidos en álgebra conmutativa\footnote{\Cf N. Bourbaki,
Alg. Comm. Chap\ptbl~5, \S~1 y \S~2, th\ptbl~2.} y solo se recuerdan aquí,
salvo la afirmación sobre la topología, que procede del siguiente hecho
general, consecuencia fácil del teorema de Cohen--Seidenberg: un morfismo
entero es cerrado (\ie transforma conjuntos cerrados en conjuntos cerrados).
Observemos enseguida lo siguiente:

\begin{corollaire}
\label{V.1.2}
Bajo las condiciones anteriores, el homomorfismo natural
$$\cal{O}_Y{\to} p_*(\cal{O}_X)^G$$
es un isomorfismo.
\end{corollaire}

Esto resulta inmediatamente de la fórmula
$$
(S^{-1}A)^G = S^{-1}(A^G),
$$
válida para toda parte multiplicativamente estable~$S$ de~$B=A^G$
(fórmula que admite una versión con módulos y que se enuncia, más
generalmente, para un cambio de base $A\to A'$ que sea \emph{plano}),
\ifthenelse{\boolean{orig}}{aplicado}{aplicada}
al caso en que~$S$ está engendrada por un elemento~$f$ de~$B$.

La afirmación~(ii) y el cor\ptbl~\Ref{V.1.2} implican fácilmente~(iv); más
generalmente, se tiene lo siguiente:

\begin{proposition}
\label{V.1.3}
Sean
\marginpar{107}
$X$ un preesquema con un grupo finito de automorfismos~$G$ y
$p\colon{}X\to{}Y$ un morfismo afín invariante tal que
$\cal{O}_Y\isomto p_*(\cal{O}_X)^G$. Entonces siguen siendo válidas las
conclusiones
\ifthenelse{\boolean{orig}}
{\textup{(i) (ii) (iii) (iv)}}
{\textup{(i), (ii), (iii), (iv)}}
de~\Ref{V.1.1}.
\end{proposition}

En efecto, para
\ifthenelse{\boolean{orig}}
{(i) (ii) (iii)}
{(i), (ii), (iii),}
puede suponerse que~$Y$, y por tanto~$X$, es afín; si $B,A$ son sus anillos,
la hipótesis implica $B=A^G$, y basta aplicar~\Ref{V.1.1}. Para~(iv), se
utilizan~(ii) y $\cal{O}_Y=p_*(\cal{O}_X)^G$.

\begin{corollaire}
\label{V.1.4}
Bajo las condiciones de~\Ref{V.1.3}, para todo abierto~$U$ de~$Y$, $U$ es
un cociente de $X|U=p^{-1}(U)$ por~$G$.
\end{corollaire}

En efecto, $p^{-1}(U)\to{}U$, inducido por~$p$, satisface las mismas
hipótesis que~$p$.

Si ahora~$X$ es un $Z$-preesquema y las operaciones de~$G$ son
$Z$-automorfismos, entonces, por~(iv), $Y$ es un $Z$-preesquema. Dicho esto:

\begin{corollaire}
\label{V.1.5}
Para que $X$ sea afín, \resp separado sobre~$Z$, es necesario y suficiente
que~$Y$ lo sea. Si $X$ es de tipo finito sobre~$Z$, es \emph{finito}
sobre~$Y$; si además~$Z$ es localmente noetheriano, $Y$ es de tipo finito
sobre~$Z$.
\end{corollaire}

Como $X$ es afín y, con mayor razón, separado sobre~$Y$, si $Y$ es afín,
\resp separado sobre~$Z$, también lo es~$X$. Recíprocamente, supongamos que
$X$ es afín sobre~$Z$ y probemos que~$Y$ lo es: gracias a~\Ref{V.1.4} puede
suponerse que~$Z$ es afín, y la cuestión se reduce a demostrar que si~$X$ es
afín, también lo es~$Y$, lo cual resulta de la determinación explícita de~$Y$
como
\ifthenelse{\boolean{orig}}{$\Spec(A)^G$}{$\Spec(A^G)$}
efectuada en~\Ref{V.1.1}.
Del mismo modo, como $p\colon{}X\to{}Y$ es entero, y por tanto
universalmente cerrado, y es sobreyectivo, se sigue que, si~$X$ es separado
sobre~$Z$, también lo es~$Y$ (¡lema que convendría enunciar por separado!);
en efecto, en el diagrama
$$
\xymatrix@C=1.5cm{ X \times_ZX\ar[r]^-{p\times_Zp} & Y\times_ZY \\
X \ar[u]^{\Delta_{X/Z}}\ar[r]^{p} & Y \ar[u]_{\Delta_{Y/Z}} }
$$
el morfismo $X\times_ZX\to Y\times_ZY$ es cerrado, por lo que transforma la
diagonal (cerrada) de~$X\times_ZX$ en una parte cerrada de~$Y\times_ZY$, que
no es sino la diagonal de este último, pues~$p$ es sobreyectivo.
--- Si $X$ es de tipo finito sobre~$Z$, lo es con mayor razón sobre~$Y$ y,
por consiguiente, es finito sobre~$Y$ (pues el morfismo $X\to Y$ ya es
entero).
Supongamos además que~$Z$ es localmente noetheriano y probemos que~$Y$ es de
tipo finito sobre~$Z$. Gracias a~\Ref{V.1.4} puede suponerse que~$Z$ es afín.
Como el espacio topológico~$X$ es cuasicompacto y
$p\colon X\to Y$ es sobreyectivo, $Y$ es
\ifthenelse{\boolean{orig}}{tambien}{también}
cuasicompacto y, por tanto, unión finita de abiertos
\marginpar{108}
afines; por~\Ref{V.1.4}, la cuestión se reduce al caso en que~$Y$ es afín y,
por tanto, también~$X$. Ahora bien, el anillo~$A$ de~$X$ es un álgebra de
tipo finito sobre el anillo noetheriano~$C$ de~$Z$, y se sabe que
$B=A^G$ también es un álgebra de tipo finito sobre~$C$ (pues $A$ será
entera, y por tanto finita, sobre una subálgebra~$B'$ de~$B$ de tipo finito
sobre~$C$; como~$B'$ es noetheriana, el álgebra~$B$ también es
\ifthenelse{\boolean{orig}}{finito}{finita}
sobre~$B'$, y por consiguiente es de tipo finito sobre~$C$).


\begin{corollaire}
\label{V.1.6}
Para que $X$ sea afín, \resp un esquema,
\ifthenelse{\boolean{orig}}{es n.~y~s.}{es necesario y suficiente}
que~$Y$ lo sea.
\end{corollaire}

\begin{definition}
\label{V.1.7}
Sea
\marginpar{109}
$X$ un preesquema sobre el que un grupo finito~$G$ actúa por la derecha. Se
dice que la acción de~$G$ es \emph{admisible} si existe un morfismo
$p\colon X\to Y$ que tiene las propiedades de~\Ref{V.1.3} (lo cual implica
que $X/G$ existe y es isomorfo a~$Y$).
\end{definition}

\begin{proposition}
\label{V.1.8}
Sea $X$ un preesquema sobre el que el grupo finito~$G$ actúa por la derecha.
Para que la acción de~$G$ sea admisible es necesario y suficiente que~$X$
sea unión de abiertos afines invariantes por~$G$, o, equivalentemente, que
toda órbita de~$G$ en~$X$ esté contenida en un abierto afín.
\end{proposition}

La primera condición implica evidentemente la segunda, y esta, a su vez,
implica la primera. En efecto, sean~$T$ una órbita de~$G$ y~$U$ un abierto
afín que la contiene. La intersección de los transformados de~$U$ por
\ifthenelse{\boolean{orig}}{los~$g$ $G$}{los~$g$ de~$G$}
es entonces un abierto~$U'$ estable por~$G$, que contiene a~$T$ y está
contenido en el abierto afín~$U$. Como, dentro de~$U$, toda parte finita
posee un sistema fundamental de entornos abiertos afines, existe un entorno
abierto afín~$V$ de~$T$ contenido en~$U'$. Sus transformados por
\ifthenelse{\boolean{orig}}{los~$g$ $G$}{los~$g$ de~$G$}
son, pues, afines y están contenidos en~$U'$, que es \emph{separado}; por
tanto, su intersección~$U''$ es un abierto afín, invariante por~$G$, que
contiene a~$T$.
--- Sentado esto, la condición considerada en~\Ref{V.1.8} es
\emph{necesaria}, pues se toman las imágenes inversas~$X_i$ de abiertos
afines~$Y_i$ que recubren~$Y$. Es suficiente porque, por~\Ref{V.1.1}, pueden
construirse entonces los cocientes $Y_i=X_i/G$; en cada~$Y_i$, la imagen de
$X_i\cap X_j$ es un abierto~$Y_{ij}$ que se identifica con~$X_{ij}/G$
por~\Ref{V.1.4}. En particular, se deducen isomorfismos
$Y_{ij}\isomto Y_{ji}$ que permiten pegar los~$Y_i$ para construir~$Y$.
Serre prefiere construir directamente el espacio topológico cociente~$Y$
de~$X$ por~$G$, dotarlo del haz $p_*(\cal{O}_X)^G$ y comprobar que~$Y$ se
convierte en un preesquema y que se está entonces bajo las condiciones
de~\Ref{V.1.3}.

\setcounter{proposition}{6}

\begin{corollaire}
\label{cor:V.1.7}
Si la acción de~$G$ sobre~$X$ es admisible, también lo es la de todo subgrupo
$H$ de~$G$ (por tanto, $X/H$ existe).
\end{corollaire}

Esto puede comprobarse también directamente en la situación de~\Ref{V.1.3},
observando que siempre puede suponerse que~$X$ es afín sobre algún~$Z$ y que
los $s\in G$ actúan por $Z$-automorfismos (se toma, por ejemplo, $Z=Y$). En
efecto:

\begin{corollaire}
\label{cor:V.1.8}
Supongamos que~$X$ es afín sobre~$Z$ y que las operaciones de~$G$ son
$Z$-automorfismos. Entonces la acción de~$G$ sobre~$X$ es admisible. Si~$X$
está definido por un haz cuasicoherente~$\cal{A}$ de álgebras, $Y$ está
definido por el haz~$\cal{A}^G$ de invariantes de~$\cal{A}$ por~$G$.
\end{corollaire}

\begin{proposition}
\label{V.1.9}
Supongamos que la acción de~$G$ sobre~$X$ es admisible y que $X/G=Y$ es un
preesquema sobre~$Z$. Consideremos un morfismo de cambio de base $Z'\to Z$ y
pongamos $X'=X\times_Z Z'$, $Y'=Y\times_Z Z'$, de modo que~$G$ sigue
actuando, por transporte de estructura, sobre~$X'$, y el morfismo
$p'\colon X'\to Y'$ es invariante. Si~$Z'$ es \emph{plano} sobre~$Z$,
entonces~$p'$ sigue satisfaciendo las hipótesis
\ifthenelse{\boolean{orig}}
{de~\Ref{V.1.3}}
{de~\Ref{V.1.3},}
\ie $\cal{O}_{Y'}\to p_*'(\cal{O}_{X'})^G$ es un isomorfismo (pues~$p'$ es
afín en cualquier caso). Por consiguiente, la acción de~$G$ sobre~$X'$ es
admisible y
$(X/G)\times_Z Z'\thickapprox (X\times_Z Z')/G$.
\end{proposition}

Puede suponerse evidentemente que $Z=Y$, y la cuestión se reduce al caso en
que, además, $Y$ e~$Y'$ son afines. Hay que demostrar que, si~$B$ es el
subanillo de invariantes de la acción de~$G$ sobre~$A$ y si~$B'$ es un
$B$-álgebra plana sobre~$B$, entonces~$B'$ es la subálgebra de invariantes de
$A'=A\otimes_B B'$. Esto es inmediato, pues la sucesión exacta
\ifthenelse{\boolean{orig}}{}
{\enlargethispage{.5cm}}%
$$
0 \to B \lto{i} A \lto{j} A^{(G)}
$$
(donde el último término denota una potencia de~$A$ y $j(x)$ es el sistema
de los $s\cdot x-x$, $s\in G$) sigue siendo exacta al tensorizar por
\ifthenelse{\boolean{orig}}
{$A'$}
{$B'$}.

Debe advertirse que la hipótesis de planitud era esencial para la validez
del resultado; en particular, si~$Y'$ es un subpreesquema cerrado de
\ifthenelse{\boolean{orig}}
{$X$}
{$Y$}
(por ejemplo, incluso un punto cerrado de
\ifthenelse{\boolean{orig}}
{$X$}
{$Y$})
y~$X'$ es su imagen inversa en~$X$, entonces, en general, $Y'$ \emph{no se
identifica} con~$X'/G$. Veremos que, no obstante, sí se identifica cuando~$X$
es étale sobre~$Y$.

Para terminar, demos un formalismo tan cómodo como trivial. Sea~$Y$ un
preesquema. Como en la categoría de los preesquemas existen las sumas
disjuntas, para todo conjunto~$E$ puede considerarse el preesquema suma de
una familia $(Y_i)_{i\in E}$ de preesquemas, todos idénticos a~$Y$; este
preesquema se denotará por~$Y\times E$.
\marginpar{110}
Está caracterizado por la fórmula
\begin{equation*}
\label{eq:V.1.*}
\tag{$*$} \Hom(Y\times E,Z) = \Hom(E,\Hom(Y,Z)),
\end{equation*}
donde el segundo~$\Hom$ denota, evidentemente, el conjunto de aplicaciones
del conjunto~$E$ en el conjunto~$\Hom(Y,Z)$. Hay un morfismo canónico
$$
Y\times E\to Y
$$
que hace de~$Y\times E$ un preesquema sobre~$Y$. Como los productos fibrados
conmutan con las sumas disjuntas (en la categoría de los preesquemas), si~$Y$
es un preesquema sobre otro~$Z$, se tiene, para un cambio de base $Z'\to Z$:
$$
(Y\times E)\times_Z Z'=(Y\times_Z Z')\times E
$$
(fórmula especialmente útil si $Z=Y$). Por otra parte, se deduce
trivialmente de la definición que
$$
(Y\times E)\times F=Y\times(E\times F)
=(Y\times E)\times_Y(Y\times F)
$$
(aunque la última fórmula resulta también de la conmutatividad señalada más
arriba).

Para~$Y$ fijo, puede considerarse $Y\times E$ como un funtor de~$E$, con
valores en los preesquemas sobre~$Y$; según la fórmula anterior, este funtor
conmuta con los productos finitos. Esto permite, por ejemplo, asociar a todo
grupo ordinario~$G$ un esquema en grupos $Y\times G$ sobre~$Y$, que será
finito sobre~$Y$ si
\ifthenelse{\boolean{orig}}{$Y$ lo es}{$G$ es finito},
\ifthenelse{\boolean{orig}}{etc...}{etc.}
Más generalmente, este funtor es «exacto por la izquierda», pero no tendremos
que utilizar aquí este hecho. También conmuta trivialmente con las sumas
disjuntas y es «exacto por la derecha», como se ve de inmediato en la fórmula
de definición~\eqref{eq:V.1.*}. En particular, si el grupo finito~$G$ actúa
por la derecha sobre el conjunto~$E$, entonces actúa por la derecha sobre
$Y\times E$, y se tiene
$$
(Y\times E)/G = Y\times (E/G),
$$
donde, de hecho, el cociente del primer miembro satisface las condiciones
de~\Ref{V.1.3} (esto es inmediato).


\section{Grupos de descomposición y de inercia. Caso étale}
\label{V.2}
Sea~$G$
\ifthenelse{\boolean{orig}}{\ignorespaces}{un}
grupo finito que actúa por la derecha sobre el preesquema~$X$. Si $x\in X$,
se llama
\marginpar{111}
\emph{grupo de descomposición de~$x$}
\index{descomposición (grupo de)|hyperpage}%
\index{grupo de descomposición|hyperpage}%
al estabilizador $G_d(x)$ de~$x$. Este grupo actúa canónicamente (por la
izquierda) sobre el cuerpo residual
\ifthenelse{\boolean{orig}}
{$\kappa(x)$}
{$\kres(x)$}, y el conjunto de los elementos de~$G_d(x)$ que actúan
trivialmente se llama 
\emph{grupo de inercia}
\index{grupo de inercia|hyperpage}%
\index{inercia (grupo de)|hyperpage}%
de~$x$, y se denota por~$G_i(x)$.

Supongamos que la acción de~$G$ sobre~$X$ es admisible y que~$Y$ es un
preesquema sobre un preesquema~$Z$. Fijemos $z\in Z$ y una extensión
algebraicamente cerrada~$\Omega$
\ifthenelse{\boolean{orig}}
{de~$\kappa(z)$ cuyo grado de trascendencia sea mayor que el de las
$\kappa(x)/\kappa(z)$}
{de~$\kres(z)$ cuyo grado de trascendencia sea mayor que el de las
$\kres(x)/\kres(z)$}, donde~$x$ es un punto de~$X$ situado sobre~$z$. Puede
considerarse $\Spec(\Omega)$ como un $Z$-esquema, y los puntos de~$X$ con
valores en~$\Omega$ corresponden a los homomorfismos
\ifthenelse{\boolean{orig}}
{de $\kappa(z)$-álgebras $\kappa(x)\to\Omega$}
{de $\kres(z)$-álgebras $\kres(x)\to\Omega$}, donde~$x$ es un punto de~$X$
situado sobre~$z$. Como $\Omega$ se ha tomado suficientemente grande, todo
punto~$x$ de~$X$ situado sobre~$z$ es el punto subyacente a un punto de~$X$
con valores en~$\Omega$. Si $X(\Omega)$ e~$Y(\Omega)$ designan,
respectivamente, los conjuntos de puntos de~$X$ y de~$Y$ con valores
en~$\Omega$, se tiene una aplicación natural
$$
X(\Omega)\to Y(\Omega),
$$
y, por otra parte, $G$ actúa sobre~$X(\Omega)$ y la aplicación precedente es
invariante por~$G$. Sentado esto, las conclusiones (ii) y (iii)
de~\Ref{V.1.3} pueden interpretarse también así: \emph{la aplicación
precedente es sobreyectiva e identifica $Y(\Omega)$ con el cociente
$X(\Omega)/G$. Además, si~$x$ es el punto subyacente a $a\in X(\Omega)$,
entonces el estabilizador de~$a$ en~$G$ no es otro que el grupo de inercia
$G_i(x)$}. Todo esto es cierto, por lo demás, sin suponer que~$\Omega$ sea
«suficientemente grande»; esta última hipótesis sirve únicamente para
asegurar que el grupo de inercia de todo punto de~$X$ situado sobre~$z$ pueda
caracterizarse como estabilizador «geométrico». Se deduce inmediatamente,
por ejemplo:

\begin{proposition}
\label{V.2.1}
Hagamos una extensión de la base $Z'\to Z$, de donde
$X'=X\times_ZZ'$. Sea~$x'$ un punto de~$X'$ y~$x$ su imagen en~$X$.
Entonces $G_i(x)=G_i(x')$.
\end{proposition}

En las consideraciones precedentes basta tomar como~$\Omega$ una extensión
suficientemente grande
\ifthenelse{\boolean{orig}}
{de~$\kappa(z')$}
{de~$\kres(z')$}
(donde $z,z'$ son las imágenes de $x,x'$ en $Z,Z'$, respectivamente).

\begin{proposition}
\label{V.2.2}
Bajo las condiciones de~\Ref{V.1.3}, supongamos que~$Y$ es localmente
noetheriano y que~$X$ es finito sobre~$Y$. Sea~$H$ un subgrupo de~$G$;
consideremos $X'=X/H$ (\cf \Ref{cor:V.1.7}) y sean
\marginpar{112}
$x\in X$, $x'$ su imagen en~$X'$ e~$y$ su imagen en~$Y$.
\begin{enumerate}
\item[(i)] Si $H\supset G_d(x)$, entonces el homomorfismo
$\cal{O}_y\to\cal{O}_{x'}$ induce un isomorfismo entre las completaciones.
\item[(ii)] Si $H\supset G_i(x)$, entonces el homomorfismo
$\cal{O}_y\to\cal{O}_{x'}$ es étale, \ie $X'$ es étale sobre~$Y$ en~$x'$.
\end{enumerate}
\end{proposition}

Sea $Y_1=\Spec(\widehat{\cal{O}_y})$ y hagamos el cambio de base
$Y_1\to Y$. Se obtiene $X_1=X\times_Y Y_1$, finito sobre~$Y_1$, sobre el
que actúa~$G$, siendo el cociente~$Y_1$ por~\Ref{V.1.9}. Sea~$y_1$ el único
punto de~$Y_1$ situado sobre~$y$. Como
\ifthenelse{\boolean{orig}}
{$\kappa(y)=\kappa(y_1)$}
{$\kres(y)=\kres(y_1)$},
la fibra de~$X$ en~$y$ es isomorfa a la de~$X_1$ en~$y_1$; de aquí resulta
un único punto~$x_1$ de~$X_1$ situado sobre~$x$. Además, por~\Ref{V.1.9},
se tiene $X_1/H=X'_1=X'\times_Y Y_1$. Sea~$x'_1$ la imagen de~$x_1$ en
$X'_1$; está situado sobre~$x'$, y se comprueba fácilmente (pues~$X'$ es de
tipo finito sobre~$Y$) que el homomorfismo
$\cal{O}_{x'}\to\cal{O}_{x'_1}$ induce un isomorfismo entre las
completaciones. Por tanto, la cuestión se reduce al caso en que~$Y$ es el
espectro de un anillo local completo~$B$ y, por consiguiente, $X$ es el
espectro de una $B$-álgebra finita~$A$, producto de un número finito de
anillos locales~$A_{x_i}$
\ifthenelse{\boolean{orig}}{correspondientes}{correspondientes}
a los puntos~$x_i$ de~$X$ situados sobre~$Y$. Si~$A_0$ corresponde a
$x=x_0$, entonces~$A$ se identifica con el anillo
$\Hom_{G_d}(G,A_0)$ de las funciones $f\colon G\to A_0$ tales que
$f(st)=sf(t)$ para $s\in G_d$, estando definidas las operaciones de~$G$ sobre
estas funciones por $(uf)(t)=f(tu)$. Se ve, pues, que si~$H$ es un subgrupo
cualquiera de~$G$, entonces~$A^H$ es el anillo de las funciones
$f\colon G\to A_0$ tales que
$$
f(stu)=sf(t)\qquad\text{para}\ s\in G_d,\ u\in H.
$$
Por tanto, es un anillo semilocal cuyos componentes locales corresponden a
las clases dobles $G_daH$ en~$G$; a la clase doble definida por $a\in G$
corresponde (mediante la aplicación $f\mto f(a)$) el subanillo
$A_0^{H(a)}$ de~$A_0$, donde $H(a)=G_d\cap aHa^{-1}$. Además, el componente
local de~$A^H$ correspondiente a la imagen~$x'$ de~$x$ es también el que
corresponde a la clase doble~$G_dH$ del elemento neutro; su componente local
es, por tanto, $A_0^{G_d\cap H}$. Si, pues,
\ifthenelse{\boolean{orig}}{$G_d\in H$,}{$G_d\subset H$,}
se obtiene $A_0^{G_d}=A^G=B$, lo cual demuestra (i). Para demostrar (ii),
puede pasarse a una extensión finita plana conveniente de~$A$ y utilizarse
\Ref{V.2.1}, reduciendo así la cuestión al caso en que la extensión residual
\ifthenelse{\boolean{orig}}
{$\kappa(x)/\kappa(y)$}
{$\kres(x)/\kres(y)$}
es trivial. Pero entonces $G_i(x)=G_d(x)$, y se vuelve al caso precedente.

\begin{corollaire}
\label{V.2.3}
Bajo
\marginpar{113}
las condiciones de~\Ref{V.2.2}, supongamos que $G_i(x)=(e)$. Entonces~$X$
es étale sobre~$Y$ en~$x$. Por consiguiente, si $G_i(x)=(e)$ para todo
$x\in X$, entonces $X\to Y$ es un morfismo étale.
\end{corollaire}

Hay un recíproco parcial:

\begin{corollaire}
\label{V.2.4}
Supongamos que~$X$ es conexo y que la acción de~$G$ sobre~$X$ es fiel. Para
que $p\colon X\to Y=X/G$ sea étale es necesario y suficiente que los grupos
de inercia de los puntos de~$X$ se reduzcan al elemento neutro. Si así
ocurre, $G$ se identifica con el grupo de todos los $Y$-automorfismos del
$Y$-esquema~$X$.
\end{corollaire}

En vista de~\Ref{V.2.3}, puede suponerse que~$X$ es étale sobre~$Y$. Pero si
un $s\in G$ pertenece a algún $G_i(x)$, de~\Ref{I.5.4} se sigue entonces
\ifthenelse{\boolean{orig}}
{que~$s$ actúa trivialmente sobre~$G$}
{que~$s$ actúa trivialmente sobre~$X$},
y, por tanto, es el elemento neutro, pues la acción de~$G$ es fiel; esto
demuestra la primera afirmación. Sea~$u$ un $Y$-automorfismo de~$X$ y sea
$x\in X$. Según la Proposición~\Ref{V.1.3}, existe $s\in G$ tal que
$x'=s(x)=u(x)$ y que induce el mismo homomorfismo residual
\ifthenelse{\boolean{orig}}
{$\kappa(x)\to\kappa(x')$}
{$\kres(x')\to\kres(x)$}
que~$u$. Por~\loccit, se tiene entonces $s=u$, lo cual concluye la
demostración.


\begin{remarque}
\label{V.2.5}
\ifthenelse{\boolean{orig}}
{La hipótesis de que~$G$ actúa fielmente no es evidentemente por superflua}
{La hipótesis de que~$G$ actúa fielmente evidentemente no es superflua}
en el Corolario~\Ref{V.2.4}. Lo mismo ocurre con la hipótesis de que~$X$ es
conexo, como se ve, por ejemplo, tomando $X=Y\times E$, donde~$E$ es un
conjunto finito, y tomando como~$G$ el grupo de permutaciones de~$E$: $G$
actúa con mucha inercia; no obstante,
$(Y\times E)/G=Y\times(E/G)=Y$, y~$X$ es étale sobre~$Y$. Tomando como~$G$
un grupo estrictamente menor que el grupo simétrico de~$E$, pero que actúe
transitivamente sobre~$E$, se ve que habrá también $Y$-automorfismos de~$X$
que no proceden de~$G$.
\end{remarque}

El ejemplo típico de un grupo~$G$ que actúa sin inercia es el de
$Y\times G$, sobre el cual se hace actuar~$G$ mediante sus traslaciones por
la derecha sobre el factor~$G$: un $Y$-preesquema~$X$ con grupo de operadores
por la derecha~$G$ se dice \emph{trivial} si es isomorfo a~$Y\times G$.

Para relacionar
\ifthenelse{\boolean{orig}}{el preesquemas}{los preesquemas}
con grupos finitos de operadores y la noción de fibrado principal en una
categoría (relación que no necesitaremos, por lo demás, en el resto del
seminario, pero que es importante en otros contextos), son útiles las
consideraciones siguientes. Fijamos un preesquema base~$Y$ y nos situamos
\ifthenelse{\boolean{orig}}{en el categoría}{en la categoría}
de los $Y$-preesquemas. Si~$G$ es un grupo
\marginpar{114}
finito, pondremos para abreviar $G_Y=Y\times G$; se trata, pues, de un
esquema en grupos finito sobre~$Y$ (\cf \numero 1), y si~$X$ es un
$Y$-preesquema, se tiene
$$
X\times_YG_Y=X\times G
$$
(misma referencia). Dar un $Y$-morfismo $X\times_YG_Y\to X$ equivale, por
tanto, a dar un $Y$-morfismo $X\times G\to X$, \ie a dar, para todo
$g\in G$, un $Y$-morfismo $T_g\colon X\to X$. Se comprueba inmediatamente
que, para que los~$T_g$ definan sobre~$X$ una estructura de preesquema con
grupo de operadores por la derecha~$G$ (\ie $T_{gg'}=T_{g'}T_g$ y
\ifthenelse{\boolean{orig}}{$T_e=\id_x$}{$T_e=\id_X$}),
es necesario y suficiente que el $Y$-morfismo correspondiente
$X\times_YG_Y\to X$ defina sobre~$X$ una estructura de $Y$-preesquema con
acción del $Y$-esquema en grupos~$G_Y$ (en el sentido general de los objetos
con acción de un $\cal{C}$-grupo en una categoría~$\cal{C}$). Supongamos que
así ocurre. Recordemos que~$X$ se dice \emph{formalmente homogéneo principal}
\index{formalmente homogéneo principal (preesquema)|hyperpage}%
bajo~$G_Y$\kern1pt\footnote{Ahora se dice más bien que~$X$ es un
pseudotorsor bajo~$G_Y$.} si el morfismo canónico
$$
X\times_YG_Y\to X\times_YX
$$
cuyas componentes son, respectivamente, $\pr_1$ y el morfismo de acción
$\pi\colon X\times_YG_Y\to X$, es un isomorfismo. En este caso,
identificando el primer miembro con~$X\times G$, el morfismo considerado es
el que asocia a todo $g\in G$ el morfismo
$$
(\id_X,T_g)=(\id_X\times_YT_g)\Delta_{X/Y}\colon X\to X\times_YX.
$$
Por consiguiente, decir que~$X$ es formalmente homogéneo principal bajo~$G_Y$
significa también que $X\times_YX$ es isomorfo a la suma disjunta de los
transformados de la diagonal por los elementos $(e,g)$ de~$G\times G$ (que
actúan sobre~$X\times_YX$ de manera evidente), donde~$e$ designa el elemento
neutro de~$G$. Si no se quiere distinguir entre izquierda y derecha y se
desea una fórmula que siga siendo aplicable al producto de más de dos
factores idénticos a~$X$, puede formularse la condición diciendo que el
morfismo canónico
$$
X\times_G(G\times G)\to X\times_YX
$$
obtenido asociando al par $(g,g')$ el morfismo
$$
(T_g,T_{g'})=(T_g\times_YT_{g'})\Delta_{X/Y}\colon X\to X\times_YX
$$
y haciendo actuar~$G$ por la izquierda sobre~$G\times G$ mediante el
homomorfismo diagonal
$$
s(g,g')=(sg,sg'),
$$
es un \emph{isomorfismo}.

La
\marginpar{115}
noción de \emph{espacio homogéneo principal}
\index{homogéneo principal (preesquema, fibrado)|hyperpage}%
se obtiene de la de espacio formalmente homogéneo principal añadiendo un
axioma suplementario que asegure que el «cociente» de~$X$ por~$G_Y$ existe y
es precisamente el objeto unidad a la derecha de la categoría, aquí~$Y$.
Este axioma puede variar según el contexto, y a menudo se explicita de la
manera más cómoda (en el yoga del «descenso») exigiendo que el objeto con
operadores se vuelva «trivial», \ie isomorfo al producto
$X\times_YG_Y$ (en este caso, $X\times G$), tras un cambio de base adecuado,
de tipo especificado (de modo que, en la práctica, permita la técnica del
descenso; \cf Grothendieck, \emph{Technique de descente et théorèmes
d'existence en Géométrie Algébrique}, Sém. Bourbaki \No190, pp. 26--28)%
\footnote{\Cf exp\ptbl\ \Ref{VIII} para la teoría del descenso plano.}.
En este orden de ideas, señalemos aquí la caracterización de los fibrados
homogéneos principales de grupo~$G$ (en el sentido de~\loccit):

\begin{proposition}
\label{V.2.6}
Sean~$Y$ un preesquema localmente noetheriano y~$X$ un $Y$-preesquema con un
grupo finito de operadores~$G$ que actúa por la derecha. Las condiciones
siguientes son equivalentes:
\begin{enumerate}
\item[(i)] $X$ es finito sobre~$Y$, $Y=X/G$, y los grupos de inercia de los
puntos de~$X$ se reducen al elemento neutro.
\item[(ii)] Existe un cambio de base fielmente plano y cuasicompacto
$Y_1\to Y$ tal que $X_1=X\times_YY_1$ sea un $Y_1$-preesquema con operadores
trivial, \ie isomorfo a~$Y_1\times G$.
\item[(ii~bis)] Como en (ii), pero siendo $Y_1\to Y$ finito, étale y
sobreyectivo.
\item[(iii)] $X$ es formalmente homogéneo principal bajo~$G_Y$, y es
fielmente plano y cuasicompacto sobre~$Y$.
\end{enumerate}
\end{proposition}
\textit{Demostración.} (i)$\To$(ii~bis). Se toma $Y_1=X$, observando que
$X\to Y$ es, en efecto, finito, étale por~\Ref{V.2.3} y sobreyectivo.
Demostremos que~$X_1$ es entonces trivial sobre~$Y_1$,
\ifthenelse{\boolean{orig}}
{lo que resultará del}
{como se deduce del}

\begin{corollaire}
\label{V.2.7}
Si se verifica \textup{(i)} y~$X$ admite una sección sobre~$Y$, entonces~$X$
es un espacio con operadores trivial.
\end{corollaire}

En efecto, esta sección permite definir un $G$-morfismo $\ifthenelse{\boolean{orig}}{X}{Y}\times G\to X$,
sobreyectivo porque~$G$ es transitivo sobre las fibras de~$X$, inyectivo
porque~$G$ actúa sin inercia y, por último, isomorfismo local en virtud
de~\Ref{I.5.3}, pues~$X$ es étale sobre~$Y$. Por tanto, es un isomorfismo.

(ii~bis)
\marginpar{116}
implica trivialmente (ii), que implica (i), pues los ingredientes de (i) son
«invariantes» por extensión fielmente plana y cuasicompacta de la base
(\cf el Seminario Bourbaki citado anteriormente para «finito»; para los
grupos de inercia se aplica~\Ref{V.2.1}, y para $Y=X/G$, un recíproco
de~\Ref{V.1.9} en el caso de un cambio de base \emph{fielmente plano}, que
habíamos olvidado explicitar).

\ifthenelse{\boolean{orig}}
{Hemos demostrado (i)$\To$(iii) al pasar demostrando (i)$\To$(ii~bis).}
{Al demostrar (i)$\To$(ii~bis), se demostró también
(i)$\To$(iii).}
Finalmente, (iii)$\To$(ii), pues la primera hipótesis de (iii) significa
precisamente que~$X$ se vuelve trivial al hacer el cambio de base $Y_1=X$;
de aquí (ii), ya que~$X$ es fielmente plano y cuasicompacto sobre~$Y$.

\begin{definition}
\label{V.2.8}
Un $Y$-preesquema~$X$ con grupo de operadores por la derecha~$G$ que satisface
las condiciones equivalentes de~\Ref{V.2.6} se llama \emph{recubrimiento
principal de~$Y$, de grupo de Galois~$G$}.
\index{grupo de Galois de un recubrimiento principal|hyperpage}%
\index{principal (recubrimiento)|hyperpage}%
\index{recubrimiento principal|hyperpage}%
\end{definition}


\section{Automorfismos y morfismos de recubrimientos étales}
\label{V.3}

\begin{proposition}
\label{V.3.1}
Sea~$X$ étale, separado y de tipo finito sobre un preesquema~$Y$ localmente
noetheriano, y sea~$G$ un grupo finito que actúa sobre~$X$ por
$Y$-automorfismos. Entonces la acción de~$G$ es admisible y el preesquema
cociente~$X/G$ es étale sobre~$Y$.
\end{proposition}

No se supone que~$X$ sea finito sobre~$Y$; sin embargo, $X$ es
cuasiproyectivo sobre~$Y$, de donde se sigue la existencia de~$X/G$
por~\Ref{V.1.8}. Demostremos primero:

\begin{corollaire}
\label{V.3.2}
El morfismo $X\to X/G$ es étale.
\end{corollaire}

Puede suponerse, evidentemente, que~$G$ actúa transitivamente sobre el
conjunto de las componentes conexas de~$X$ y, después, considerando el
estabilizador de una componente conexa, que~$X$ es incluso conexo. Por último,
puede suponerse que la acción de~$G$ es fiel. Pero entonces se ve, como
en~\Ref{V.2.4}, que~$G$ actúa sin inercia; por~\Ref{V.2.3} se sigue que
$X\to X/G$ es étale. Se concluye gracias al lema siguiente.

\begin{lemme}[remordimiento por el Exposé I]
\label{V.3.3}
Sean $X\to X'\to Y$ morfismos de tipo finito, sea~$x$ un punto de~$X$ y sean
$x'$ e~$y$ sus imágenes. Supongamos que~$Y$ es localmente noetheriano. Si dos
de los tres morfismos considerados son étales en los puntos señalados, también
lo es el tercero.
\end{lemme}

Solo queda considerar el caso en que $X\to X'$ y $X\to Y$ son étales en~$x$,
y demostrar que $X'\to Y$ lo es en~$x'$ (que es el caso que necesitamos
\marginpar{117}
para~\Ref{V.3.1}). Mediante una extensión plana conveniente de la base~$Y$, se
reduce la cuestión al caso en que la extensión residual
\ifthenelse{\boolean{orig}}
{$\kappa(x)/\kappa(y)$}
{$\kres(x)/\kres(y)$}
es trivial. Consideremos los homomorfismos
$\cal{O}_y\to\cal{O}_{x'}\to\cal{O}_x$ y los homomorfismos obtenidos al pasar
a las completaciones. La hipótesis significa que
$\hat{\cal{O}_y}\to\hat{\cal{O}_x}$ y
$\hat{\cal{O}_{x'}}\to\hat{\cal{O}_x}$ son isomorfismos; por tanto,
$\hat{\cal{O}_y}\to\hat{\cal{O}_{x'}}$ también es un isomorfismo, lo que
demuestra el lema.

\begin{corollaire}
\label{V.3.4}
Si~$X$ es finito y étale sobre~$Y$, entonces~$X/G$ es finito y étale
sobre~$Y$.
\end{corollaire}

\begin{proposition}
\label{V.3.5}
Sean~$X$ y~$X'$ dos recubrimientos étales de~$Y$. Todo $Y$-morfismo
$f\colon X\to X'$ se factoriza como la composición de un morfismo étale
sobreyectivo $X\to X''$ y la inmersión canónica $X''\to X'$ de un subconjunto
$X''$ de~$X'$ que es a la vez abierto y cerrado.
\end{proposition}

Se sabe por~\Ref{I.4.8} que~$f$ es étale y, por tanto, un morfismo abierto. Por
otra parte, como~$X$ es finito sobre~$Y$, $f$ es cerrado; así,
$f(X)=X''$ es un subconjunto a la vez abierto y cerrado de~$X'$. Esto concluye
la demostración. (N.B.: habría bastado que~$X'$, en vez de ser un recubrimiento
étale, fuese no ramificado sobre~$Y$.)

\begin{corollaire}
\label{V.3.6}
Con las notaciones precedentes,
\ifthenelse{\boolean{orig}}
{$X\to X'$ es un epimorfismo estricto en la categoría de los preesquemas, y
$X'\to X''$ es un monomorfismo (incluso un monomorfismo estricto) en dicha
categoría}
{$X\to X''$ es un epimorfismo estricto en la categoría de los preesquemas, y
$X''\to X'$ es un monomorfismo (incluso un monomorfismo estricto) en dicha
categoría}.
\end{corollaire}

Por definición, la primera afirmación significa que la sucesión de morfismos
\begingroup
\shorthandoff{<>}% Spanish babel otherwise treats Xy-pic's bend delimiters as shorthands.
$$
\xymatrix@C=.7cm{X \times_{X''} X \ar@<2pt>[r]^-{\pr_1}
\ar@<-2pt>[r]_-{\pr_2} & X \ar[r] & X''}
$$
\endgroup
es exacta. Esto resulta de que $X\to X''$ es finito y fielmente plano, como se
ve fácilmente (\cf Grothendieck, \loccit). La afirmación dual para
$X''\to X'$ es aún más inmediata.

El Corolario~\Ref{V.3.6} nos será útil para la teoría del grupo fundamental en
el número siguiente. Quienes no gusten de la noción de epimorfismo estricto
pueden sustituir el Corolario~\Ref{V.3.6} por la variante que prefieran.
Aprovechemos solamente la ocasión para señalar que una factorización
\marginpar{118}
$f=f'f''$, con~$f''$ epimorfismo estricto y~$f'$ monomorfismo, es
necesariamente única salvo isomorfismo único (en cualquier categoría). Sin
embargo, puede existir al mismo tiempo una factorización $f=f_1f_2$ con las
propiedades duales: $f_2$ es un epimorfismo y~$f_1$ un monomorfismo estricto
(también única salvo isomorfismo único), que no sea isomorfa a la anterior.
Basta tomar, por ejemplo, la categoría de los espacios vectoriales topológicos
(separados, si así se desea) y un morfismo $u\colon X\to X'$ tal que~$u(X)$ no
sea cerrado.

\begin{proposition}
\label{V.3.7}
Sean~$Y$ un preesquema \emph{conexo} localmente noetheriano, $y$ un punto
de~$Y$ y~$\mathit{\Omega}$ una extensión algebraicamente cerrada
\ifthenelse{\boolean{orig}}
{de~$\kappa(y)$}
{de~$\kres(y)$}.
Para todo~$X$ sobre~$Y$, denotemos por~$X(\Omega)$ el conjunto de los puntos
de~$X$ con valores en~$\Omega$. Sean~$X$ y~$X'$ recubrimientos étales de~$Y$,
y sea $u\colon X\to X'$ un $Y$-morfismo tal que la aplicación correspondiente
$X(\Omega)\to X'(\Omega)$ sea una biyección. Entonces~$u$ es un isomorfismo.
\end{proposition}

Se reduce inmediatamente al caso en que~$X'$ es conexo. Como $X\to X'$ es
finito y étale, el número geométrico de puntos de una fibra de $X\to X'$ es
constante y es igual a~$1$ si y solo si el morfismo considerado es un
isomorfismo. Ahora bien, ese número es también el número de elementos de una
fibra de $X(\Omega)\to X'(\Omega)$, de donde se sigue la conclusión.
\ifthenelse{\boolean{orig}}{}
{\enlargethispage{.5cm}}%


\section{Condiciones axiomáticas de una teoría de Galois}
\label{V.4}
Sea~$\cal{C}$ una categoría y~$F$ un funtor covariante de~$\cal{C}$ en la
categoría de los conjuntos finitos. Supongamos que se satisfacen las
condiciones siguientes:

\begin{enumerateb}

\item[(G 1)] $\cal{C}$ tiene un objeto final\footnote{Recordemos que un
objeto~$e$ de~$\cal{C}$ se llama \emph{objeto final} si, para todo~$X$ de
$\cal{C}$, $\Hom(X,e)$ tiene exactamente un elemento. Se define de manera
dual un \emph{objeto inicial} de~$\cal{C}$.}, y existe el producto fibrado de
dos objetos sobre un tercero en~$\cal{C}$ (este axioma también puede
enunciarse diciendo que en~$\cal{C}$ existen los límites proyectivos
finitos).

\item[(G 2)] Existen las sumas finitas en~$\cal{C}$ (y, por tanto, también
un objeto inicial $\emptyset_{\cal{C}}$, que desempeña el papel del conjunto
vacío), así como el cociente de un objeto de~$\cal{C}$ por un grupo finito
de automorfismos.

\item[(G 3)] Sea $u\colon X\to Y$ un morfismo en~$\cal{C}$. Entonces~$u$ se
factoriza como una composición $X\lto{u'}Y'\lto{u''}Y$, donde~$u'$ es un
epimorfismo \emph{estricto} y~$u''$ un monomorfismo que es un isomorfismo
sobre un sumando directo de~$Y$.

\item[(G 4)]
\marginpar{119}
El funtor~$F$ es exacto por la izquierda (\ie transforma un objeto final en
un objeto final y conmuta con los productos
\ifthenelse{\boolean{orig}}{fibrado}{fibrados}).% SGA1-CORR-0137

\item[(G 5)] $F$ conmuta con las sumas directas finitas, transforma los
epimorfismos estrictos en epimorfismos y conmuta con el paso al cociente por
un grupo finito de automorfismos.

\item[(G 6)] Sea $u\colon X\to Y$ un morfismo en~$\cal{C}$ tal que~$F(u)$
sea un isomorfismo. Entonces~$u$ es un isomorfismo.
\end{enumerateb}

Nuestro objetivo es construir un grupo topológico~$\pi$, límite proyectivo
de grupos finitos, y una equivalencia de la categoría~$\cal{C}$ con la
categoría~$\cal{C}(\pi)$
\label{indnot:eb}\oldindexnot{$\cal{C}(\pi)$|hyperpage}%
\emph{de los conjuntos finitos sobre los que} $\pi$ \emph{actúa
continuamente} (\ie de tal modo que el estabilizador de un punto sea un
subgrupo abierto, o, de manera equivalente, que exista un grupo cociente
discreto que ya actúe sobre el conjunto considerado); la equivalencia
construida transforma el funtor dado~$F$ en el funtor de inclusión evidente
de~$\cal{C}(\pi)$ en la categoría de los conjuntos finitos. Señalemos desde
ahora que la categoría~$\cal{C}(\pi)$ construida a partir de un grupo
topológico~$\pi$, junto con el funtor de inclusión precedente, satisface en
efecto las condiciones (G 1) a (G 6).

Procedemos en varias etapas.
\begin{enumerateb}
\item[a)] \emph{Sea} $u\colon X\to Y$ \emph{un morfismo en} $\cal{C}$.
\emph{Para que} $u$ \emph{sea un monomorfismo, es necesario y suficiente que}
$F(u)$ \emph{lo sea}. (Se utilizan (G 1), (G 4) y (G 6).)

En efecto, decir que~$u$ es un monomorfismo significa que la proyección
$\pr_1\colon X\times_YX\to X$ es un isomorfismo.

\item[b)] \emph{Todo objeto} $X$ \emph{de} $\cal{C}$ \emph{es artiniano}.

En efecto, si $X'\to X''\to X$ son monomorfismos tales que~$F(X')$ y
$F(X'')$ tienen la misma imagen en~$F(X)$, entonces, por a),
$F(X')\to F(X'')$ es un isomorfismo; por (G 6), también lo es
$X'\to X''$.

\item[c)] \emph{El funtor} $F$ \emph{es estrictamente prorepresentable}.
(\cf Grothendieck, \emph{Technique de descente et théorèmes d'existence en
Géométrie Algébrique}, II, Séminaire Bourbaki 195, febrero de 1960.)

En efecto, según \loccit prop\ptbl \Ref{V.3.1}, esto resulta de b) y (G 4).
Por tanto, puede encontrarse un sistema proyectivo sobre un conjunto ordenado
filtrante~$I$,
$$
P=(P_i)_{i\in I},
$$
\marginpar{120}%
en~$\cal{C}$, considerado como proobjeto de~$\cal{C}$, y un isomorfismo
funtorial
\begin{equation*}
\label{eq:V.4.*}\tag{$*$}
F(X)=\Hom_{\Pro(\cal{C})}(P,X)=\varinjlim_i\Hom_{\cal{C}}(P_i,X).
\end{equation*}
Este isomorfismo viene realizado por un elemento
$$
\varphi\in\varprojlim_iF(P_i)=F(P).
$$
Puede suponerse además que los morfismos de transición
$\varphi_{ji}\colon P_i\to P_j$ ($i\geq j$) son
\emph{epimorfismos}, y que \emph{todo epimorfismo} $P_i\to P'$ es
equivalente a un epimorfismo $P_i\to P_j$ para algún $j\leq i$ adecuado
(lo que determina el sistema proyectivo~$P$ de manera esencialmente única).

Un objeto $X\in\cal{C}$ se llama \emph{conexo}
\index{conexo (objeto de una categoría)|hyperpage}%
si no es isomorfo a la suma de dos objetos de~$\cal{C}$ ninguno de los cuales
sea isomorfo al objeto inicial~$\emptyset_{\cal{C}}$.

\item[d)] \emph{Los} $P_i$ \emph{son conexos y no son isomorfos a}
$\emptyset_{\cal{C}}$.

Si~$X$ es un objeto inicial, se tiene $F(X)=\emptyset$ por (G 5), aplicado a
la suma directa de una familia vacía, y recíprocamente por (G 6). Por tanto,
si~$X'$ es un objeto de~$\cal{C}$ que no es inicial, \ie tal que
$F(X')\neq\emptyset$, no existe ningún morfismo de~$X'$ a~$X$. Así, si
algún~$P_i$ es inicial, entonces~$i$ es un elemento máximo del conjunto
ordenado filtrante de índices~$I$, y la fórmula~\eqref{eq:V.4.*}
significaría que $F(X)=\Hom(P_i,X)$ es un conjunto reducido a un elemento
para todo~$X$, lo cual es absurdo porque
$F(\emptyset_{\cal{C}})=\emptyset$. Por consiguiente, los~$P_i$ no son
isomorfos a~$\emptyset_{\cal{C}}$.

Supongamos que $P_i=A\amalg B$; por (G 5), se sigue
$F(P_i)=F(A)\amalg F(B)$. En particular, el elemento~$a_i$ de~$F(P_i)$ que
corresponde mediante~\eqref{eq:V.4.*} al morfismo identidad
$P_i\to P_i$ pertenece a $F(A)\amalg F(B)$, por ejemplo a~$F(A)$. Esto
significa que existe $j\geq i$ tal que
$\varphi_{ij}\colon P_j\to P_i$ se factoriza como
$P_j\to A\to P_i=A\amalg B$, donde la segunda flecha es el morfismo
canónico. Por tanto, $F(P_j)\to F(P_i)$ se factoriza como
$F(P_j)\to F(A)\to F(P_i)=F(A)\amalg F(B)$ y, como
$F(P_j)\to F(P_i)$ es sobreyectivo por (G 5), se sigue que
$F(B)=\emptyset$; luego~$B$ es isomorfo a~$\emptyset_{\cal{C}}$.

\item[e)]
\marginpar{121}
\emph{Todo morfismo} $u\colon X\to Y$ \emph{en} $\cal{C}$, \emph{con} $X$
\emph{no isomorfo a} $\emptyset_{\cal{C}}$ \emph{e} $Y$ \emph{conexo, es
un epimorfismo estricto. Todo endomorfismo de un objeto conexo es un
automorfismo}.

Consideremos la factorización de~$u$ dada por (G 3). Como
$X\neq\emptyset_{\cal{C}}$, de (G 6) resulta que
$F(X)\neq\emptyset$, luego $F(Y')\neq\emptyset$, luego
$Y'\neq\emptyset_{\cal{C}}$; como~$Y$ es conexo, $Y'$ se identifica con~$Y$
y, por tanto, $u$ es un epimorfismo estricto. Supongamos que~$u$ es un
endomorfismo del objeto conexo~$X$ y demostremos que es un automorfismo. En
efecto, puede suponerse que~$X$ no es isomorfo a~$\emptyset_{\cal{C}}$; por
lo anterior, $u$ es un epimorfismo estricto y, por (G 5), $F(u)$ es un
epimorfismo. Como~$F(X)$ es un conjunto finito, $F(u)$ es biyectivo. Por
(G 6), $u$ es un automorfismo.

En particular, \emph{todo endomorfismo de un} $P_i$ \emph{es un
automorfismo}.

\item[f)] \emph{Las condiciones siguientes sobre un} $P_i$ \emph{son
equivalentes}: (i) La aplicación inyectiva natural
$\Hom(P_i,P_i)\to\Hom(P,P_i)\simeq F(P_i)$ es también sobreyectiva, \ie para
todo $u\colon P\to P_i$ existe $v\colon P_i\to P_i$ tal que
$u=v\varphi_i$ (donde~$\varphi_i$ es el morfismo canónico
$P\to P_i$). (ii) El grupo~$\Aut(P_i)$ actúa transitivamente sobre
$F(P_i)$. (iii)~El grupo~$\Aut(P_i)$ actúa simplemente transitivamente
sobre~$F(P_i)$.

En efecto, al identificar $\Hom(P,P_i)$ con~$F(P_i)$, la aplicación
considerada en (i) no es sino $v\mto F(v)(\varphi_i)$. La equivalencia de las
tres condiciones resulta entonces de que
$\Hom(P_i,P_i)=\Aut(P_i)$ y de que la aplicación precedente ya es inyectiva.

Un~$P_i$ que satisface las condiciones equivalentes
\ifthenelse{\boolean{orig}}
{(i) (ii) (iii)}
{(i), (ii), (iii)}
de~f) se llama \emph{galoisiano}.
\index{galoisiano (objeto)|hyperpage}%

\item[g)] \emph{Para todo} $X$ \emph{de} $\cal{C}$, \emph{existe un}
$P_i$ \emph{galoisiano tal que todo} $u\in\Hom(P,X)$ \emph{se factoriza
como} $P\lto{\varphi_i}P_i\longrightarrow X$.

Sea $J=\Hom(P,X)=F(X)$; es un conjunto finito, por lo que existe un~$P_j$ tal
que todo $u\colon P\to X$ se factoriza como
\ifthenelse{\boolean{orig}}
{$P\lto{j}P_j\longrightarrow X$}
{$P\lto{\varphi_j}P_j\longrightarrow X$},
o, equivalentemente, tal que el morfismo natural
$$
P\to X^J\qquad(J=\Hom(P,X))
$$
se factoriza como
$$
P\lto{\varphi_j}P_j\to X^J.
$$
En virtud de
\marginpar{122}
(G~3), el morfismo $P_j\to X^J$ se factoriza como la composición de un
epimorfismo estricto y un monomorfismo, pudiéndose tomar el primero de la
forma $\varphi_{ij}\colon P_j\to P_i$. Queda, pues, demostrar que~$P_i$ es
galoisiano. Sea~$k$ un índice $\geq j$ tal que todo morfismo $P\to P_i$ se
factoriza por $P\lto{\varphi_k}P_k\longrightarrow P_i$. Obsérvese que el
morfismo natural $P_k\to X^J$ sigue factorizándose como la composición
$$
P_k\lto{\varphi_{ik}}P_i\lto{U}X^J,
$$
donde la primera \ifthenelse{\boolean{orig}}{flecha}{flecha} es un epimorfismo
estricto por e), y la segunda, un monomorfismo. Queremos demostrar que, dado
un morfismo $\psi\colon P_k\to P_i$, existe un endomorfismo~$v$ de~$P_i$ tal
que $\psi=v\varphi_{ik}$. Para todo $u\in\Hom(P_i,X)$, consideremos
$u\psi\in\Hom(P_k,X)$; es de la forma $u'\varphi_{ik}$, con un
$u'\in\Hom(P_i,X)$ bien determinado. La aplicación $u\mto u'$ de~$J$ en~$J$
así definida por~$\psi$ es, además, inyectiva, porque~$\psi$ es un
epimorfismo en virtud de e); por tanto, es biyectiva, pues~$J$ es finito. La
aplicación biyectiva $u\mto u'$ de~$J$ en~$J$ define entonces un isomorfismo
$\alpha\colon X^J\isomto X^J$ que hace conmutativo el diagrama
$$
\xymatrix{P_k \ar[r]^-{\varphi_{ik}} \ar@{=}[d] & P_i \ar[r]^-U & X^J
\ar[d]^-{\alpha}_-{\simeq} \\ P_k \ar[r]^{\psi} & P_i \ar[r]^-U & X^J}.
$$
De las propiedades de unicidad de la factorización de un morfismo como
composición de un epimorfismo estricto y un monomorfismo se sigue (pues
$\psi$ también es un epimorfismo estricto por e)) que puede encontrarse un
morfismo $v\colon P_i\to P_i$ que haga conmutativo el diagrama, como se
quería.

Se concluye, en particular, que \emph{los} $P_i$ \emph{galoisianos forman un
sistema cofinal en el sistema de los} $(P_j)$. Por consiguiente, como para un
objeto galoisiano~$P_i$ se tiene
$$
\Hom(P,P_i)=\Hom(P_i,P_i)=\Aut(P_i),
$$
al pasar al límite se obtiene
$$
\Hom(P,P)=\varprojlim_i\Hom(P,P_i)=\varprojlim_i\Hom(P_i,P_i)
=\varprojlim_i\Aut(P_i),
$$
donde
\ifthenelse{\boolean{orig}}
{\ignorespaces}
{\enlargethispage{.5cm}}
\marginpar{123}
el límite proyectivo se toma sobre los~$P_i$ galoisianos. Además, mediante
la identificación $\Hom(P,P_i)=F(P_i)$ y teniendo en cuenta que~$F$
transforma epimorfismos en epimorfismos, se ve que los morfismos de transición
del sistema proyectivo precedente son sobreyectivos. De todo ello se concluye:

\item[h)] \emph{Se tiene}
$$
\Hom(P,P)=\Aut(P)=\varprojlim_iF(P_i)=\varprojlim_i\Aut(P_i),
$$
\emph{donde el límite proyectivo se toma sobre los} $P_i$
\emph{galoisianos}.

En particular, $\Aut(P)$ aparece como límite proyectivo de un sistema
proyectivo de grupos finitos (cuyos morfismos de transición son
sobreyectivos); lo dotaremos de la topología de límite proyectivo de las
topologías discretas. \emph{Denotaremos por}~$\pi$
\label{indnot:ec}\oldindexnot{$\pi$|hyperpage}%
\emph{y llamaremos grupo fundamental}
\index{grupo fundamental de una categoría galoisiana|hyperpage}%
(de~$\cal{C}$ provista de~$F$) \emph{al grupo opuesto a}~$\Aut(P)$. Este
grupo actúa, pues, \emph{por la derecha} sobre~$P$; es el límite proyectivo
de grupos finitos~$\pi_i$ que actúan por la derecha sobre los~$P_i$
galoisianos, siendo~$\pi_i$ el grupo opuesto a~$\Aut(P_i)$.

Habida cuenta del isomorfismo funtorial
$$
F(X)=\Hom(P,X)
$$
y de la definición de~$\pi$, se ve que~$\pi$ actúa \emph{por la izquierda}
sobre~$F(X)$ y, además, de manera continua por g) (pues, con las notaciones
de g), es en realidad~$\pi_i$ quien actúa sobre~$F(X)$). Es inmediato que,
para todo morfismo $u\colon X\to Y$ en~$\cal{C}$, el morfismo
$F(u)\colon F(X)\to F(Y)$ es compatible con las acciones de~$\pi$.
\emph{Por tanto, en lo sucesivo puede considerarse} $F$ \emph{como un funtor
covariante}
$$
F\colon\cal{C}\to\cal{C}(\pi),
$$
\emph{donde} $\cal{C}(\pi)$ \emph{es la categoría de los conjuntos finitos
sobre los que} $\pi$ \emph{actúa continuamente por la izquierda}.

Definamos ahora un funtor en sentido inverso:
$$
G\colon\cal{C}\from\cal{C}(\pi)
$$
mediante la fórmula
$$
G(E)=P\times_{\pi}E,
$$
donde $P\times_{\pi}E$ se define como solución del problema universal
resumido por
$$
\Hom_{\cal{C}}(P\times_{\pi}E,X)\isomto
\Hom_{\pi}(E,\Hom(P,X))
$$
(donde,
\marginpar{124}
en el segundo miembro, $\Hom(P,X)=F(X)$ se considera como un conjunto sobre
el que~$\pi$ actúa por la izquierda). Hay que demostrar la existencia del
objeto~$P\times_{\pi}E$.

\item[i)] \emph{Sea} $Q$ \emph{un objeto de} $\cal{C}$ \emph{sobre el que
un grupo finito} $G$ \emph{actúa por la derecha, y sea} $E$ \emph{un conjunto
finito sobre el que} $G$ \emph{actúa por la izquierda. Entonces}
\ifthenelse{\boolean{orig}}
{$G\times_GE$}
{$Q\times_GE$}
\emph{existe, y la aplicación canónica}
$$
F(Q)\times_GE\to F(Q\times_GE)
$$
\emph{es un isomorfismo}.

Como las sumas directas finitas existen en~$\cal{C}$ por (G 2), y~$F$ conmuta
con ellas por (G 5), se reduce inmediatamente al caso en que~$G$ actúa
transitivamente sobre~$E$, pues, si los~$E_j$ son las órbitas de~$G$ en~$E$,
se tendrá
$$
Q\times_GE=\underset{j}{\amalg}\,Q\times_GE_j.
$$
Sea entonces $a\in E$ y sea~$H$ su estabilizador. Se ve de inmediato a partir
de la definición que $Q\times_GE$ se identifica con~$Q/H$. De aquí la
existencia, gracias a (G~2), y la propiedad de conmutación para~$F$, gracias
a (G~5).

\item[j)] \emph{Sea} $E$ \emph{un objeto de} $\cal{C}(\pi)$, \emph{y sea}
$P_i$ \emph{galoisiano tal que} $\pi_i$ \emph{ya actúa sobre} $E$.
\emph{Entonces} $P_i\times_{\pi_i}E$ \emph{existe y se tiene un isomorfismo
canónico}
$$
E\isomto F(P_i\times_{\pi_i}E).
$$
\emph{Si} $j\geq i$ \emph{es tal que} $P_j$ \emph{es galoisiano, entonces
el morfismo canónico} $P_j\times_{\pi_j}E\to\allowbreak
P_i\times_{\pi_i}E$ \emph{es un isomorfismo}.

La primera afirmación es un caso particular de i), teniendo en cuenta que
$\pi_i$ actúa simplemente transitivamente sobre~$F(P_i)$, provisto del punto
marcado~$\varphi_i$; de ahí un isomorfismo
$F(P_i)\times_{\pi_i}E\simeq E$. Para la segunda afirmación se utiliza, por
ejemplo, (G 6).

Sea, \ifthenelse{\boolean{orig}}{para}{para} todo~$j$, $\cal{C}_j$ la
subcategoría plena de~$\cal{C}$ formada por los~$X$ tales que
$\Hom(P_j,X)\to\Hom(P,X)\simeq F(X)$ sea biyectiva. Por g), se sabe que
$\cal{C}$ \emph{es la unión filtrante de las} $\cal{C}_j$. Por tanto, para
$X\in\cal{C}_j$,
\begin{align*}
\Hom_{\pi}(E,\Hom(P,X))&\simeq\Hom_{\pi}(E,\Hom(P_j,X))
\simeq\Hom_{\pi_j}(E,\Hom(P_j,X))\\
&\simeq\Hom(P_j\times_{\pi_j}E,X),
\end{align*}
y,
\marginpar{125}
teniendo en cuenta la última afirmación de j), se obtiene un isomorfismo
funtorial en el objeto~$X$ de~$\cal{C}_j$:
$$
\Hom_{\pi}(E,\Hom(P,X))\simeq\Hom(P_i\times_{\pi_i}E,X).
$$
Como esto es cierto para todo~$j$ y estos isomorfismos funtoriales, al variar
$j$, se inducen mutuamente, se concluye:

\item[k)] \emph{En las condiciones de} j), \emph{la composición de los
morfismos canónicos}
$$
E\to\Hom(P_i,P_i\times_{\pi_i}E)\to\Hom(P,P_i\times_{\pi_i}E)
$$
\emph{hace de} $P_i\times_{\pi_i}E$ \emph{una solución del problema
universal que define} $P\times_{\pi}E$, \ie \emph{este último existe y se
tiene un isomorfismo}
$$
P\times_{\pi}E\isomto P_i\times_{\pi_i}E.
$$

Esto termina la construcción del funtor~$G(E)$. Por otra parte, se tiene un
morfismo funtorial
$$
\alpha\colon\id_{\cal{C}(\pi)}\to FG,
$$
\ie un morfismo funtorial en el objeto~$E$ de~$\cal{C}(\pi)$:
$$
\alpha(E)\colon E\to FG(E)=F(P\times_{\pi}E),
$$
a saber, la composición de los morfismos canónicos
$$
E\to F(P)\times_{\pi}E\to F(P\times_{\pi}E)
$$
(donde el primero procede del punto marcado $\varphi\in F(P)$). Combinando
j) y k), se obtiene:

\item[l)] \emph{El morfismo} $\alpha$ \emph{es un isomorfismo}.

Se define del mismo modo un morfismo funtorial
$$
\beta\colon GF\to\id_{\cal{C}},
$$
\ie un morfismo funtorial en el objeto~$X$ de~$\cal{C}$:
$$
\beta(X)\colon P\times_{\pi}F(X)\to X,
$$
como
\marginpar{126}
el asociado al $\pi$-morfismo
$$
F(X)\to\Hom(P,X)
$$
inverso del isomorfismo canónico $\Hom(P,X)\isomto F(X)$.

\item[m)] \emph{Las composiciones}
\begin{gather*}
F(X)\lto{\alpha(F(X))}FGF(X)\lto{F(\beta(X))}F(X)\\
G(E)\lto{G(\alpha(E))}GFG(E)\lto{\beta(G(E))}G(E)
\end{gather*}
\emph{son los isomorfismos identidad}.

Asno que trota.

Teniendo en cuenta l), se sigue:

\item[n)] \emph{El morfismo} $\beta$ \emph{es un isomorfismo}.

Hemos obtenido así el resultado prometido:
\end{enumerateb}

\begin{theoreme}
\label{V.4.1}
Sea~$\cal{C}$ una categoría que satisface las condiciones \emph{(G 1)},
\emph{(G 2)} y \emph{(G~3)} del comienzo de este número, y sea~$F$ un
funtor covariante de~$\cal{C}$ en la categoría de los conjuntos finitos que
satisface las condiciones \emph{(G 4)}, \emph{(G 5)} y \emph{(G 6)}.
Entonces las construcciones canónicas precedentes definen equivalencias de
categorías $F\colon\cal{C}\to\cal{C}(\pi)$ y
$G\colon\cal{C}(\pi)\to\cal{C}$ cuasiinversas entre sí. Más precisamente,
existe un proobjeto~$P$ de~$\cal{C}$ y un isomorfismo funtorial
$F(X)\isomfrom\Hom(P,X)$; $\pi$ es el grupo opuesto al grupo de
automorfismos de~$P$, dotado de una topología adecuada, de modo que~$\pi$
actúa continuamente sobre los conjuntos $\Hom(P,X)\simeq F(X)$. Por último,
$G$ viene dado por $G(E)\simeq P\times_{\pi}E$.
\end{theoreme}

\begin{remarques}
\label{V.4.2}
El enunciado de las condiciones (G 1) a (G 6) se vuelve más sencillo y
agradable si se sustituyen (G 2) y (G 5), respectivamente, por:
\begin{enumerateb}
\item[(G' 2)] Existen los límites inductivos finitos en
\ifthenelse{\boolean{orig}}{$C$}{$\cal{C}$}.% SGA1-CORR-0144
\item[(G' 5)] El funtor~$F$ es exacto por la derecha (\ie conmuta con los
límites inductivos finitos).
\end{enumerateb}
Estas condiciones son en apariencia más fuertes que (G~2) y (G~5), pero del
teorema de estructura~\Ref{V.4.1} resulta inmediatamente que se deducen de
(G~1) a (G~6). Obsérvese, sin embargo, que en los casos que nos interesarán,
la comprobación de (G~2) y (G~5) parece, en efecto, más sencilla que la de
(G'~2) y (G'~5). Ignoro si, en la condición (G~3), podría omitirse que~$u''$
sea un isomorfismo sobre un sumando directo de~$Y$.
\end{remarques}


\section{Categorías galoisianas}
\label{V.5}
\marginpar{127}

\begin{definition}
\label{V.5.1}
Se llama \emph{categoría galoisiana}
\index{galoisiana (categoría)|hyperpage}%
a una categoría~$\cal{C}$ equivalente a una categoría $\cal{C}(\pi)$, donde
$\pi$ es un grupo compacto, límite proyectivo de grupos finitos (\ie
totalmente desconectado).

Para la definición de~$\cal{C}(\pi)$, \cf el comienzo del \No \Ref{V.4}. En
virtud del teorema~\Ref{V.4.1}, $\cal{C}$ es galoisiana si y solo si
satisface las condiciones (G~1) a (G~3), y si existe un funtor~$F$ de
\ifthenelse{\boolean{orig}}{$C$}{$\cal{C}$}
en la categoría de los conjuntos finitos que satisface las condiciones
(G~4) a (G~6) (\ie que es \emph{exacto} y \emph{conservativo}, en la
terminología general). Tal funtor se llamará \emph{funtor fundamental}
\index{funtor fundamental|hyperpage}%
de la categoría galoisiana~$\cal{C}$\kern1pt\footnote{Parece preferible
adoptar el término más expresivo de «funtor fibra».};
\index{funtor fibra|hyperpage}%
es prorepresentable por un proobjeto que denotaremos por~$P_F$; un
proobjeto~$P$ tal que el funtor asociado~$F$ sea fundamental se llama
\emph{proobjeto fundamental}.
\index{proobjeto fundamental|hyperpage}%
De este modo, la categoría de los funtores fundamentales sobre~$\cal{C}$ es
antiequivalente a la categoría de los proobjetos fundamentales; si~$F$ y~$P$
se corresponden, el grupo $\Aut F$ es, por tanto, isomorfo al grupo opuesto
a~$\Aut P$, de modo que el grupo denotado por~$\pi$ en el número precedente
no es sino
\ifthenelse{\boolean{orig}}{$\Aut P$}{el grupo opuesto a~$\Aut P$}
(o, equivalentemente, $\Aut F$). Recordemos que, en el número precedente,
construimos a partir de un funtor fundamental \emph{dado}~$F$ una
equivalencia de~$\cal{C}$ con $\cal{C}(\pi)$ (donde $\pi=\Aut(F)$) que
transforma el funtor~$F$ en el funtor canónico de
\ifthenelse{\boolean{orig}}{$\cal{C}(\pi)_x$}{$\cal{C}(\pi)$}
en la categoría de los conjuntos finitos. En el caso modelo
$\cal{C}=\cal{C}(\pi)$, con~$F$ el funtor canónico, el proobjeto fundamental
asociado a~$F$ no es sino el sistema proyectivo de los cocientes discretos
$\pi_i$ de~$\pi$.
\end{definition}

Puede ser útil explicitar la categoría de los proobjetos de
$\cal{C}(\pi)$. Se obtiene:

\begin{proposition}
\label{V.5.2}
La categoría $\Prodash\cal{C}(\pi)$
\label{indnot:ed}\oldindexnot{$\Prodash\cal{C}(\pi)$|hyperpage}%
es canónicamente equivalente a la categoría $\cal{C}'(\pi)$ de los
espacios compactos y totalmente desconectados sobre los que actúa
continuamente el grupo topológico~$\pi$.
\end{proposition}

Como esta última contiene a~$\cal{C}(\pi)$ como subcategoría plena
(correspondiente a los espacios discretos compactos con operadores) y en
ella existen los límites proyectivos, se tiene en todo caso un funtor
canónico $g\colon\Prodash\cal{C}(\pi)\to\cal{C}'(\pi)$, que asocia al
sistema proyectivo $Q=(Q_i)$ el objeto
$X=\varprojlim_iQ_i$ de~$\cal{C}'(\pi)$. Para definir un funtor en sentido
inverso, basta definir un funtor contravariante de~$\cal{C}'(\pi)$ en la
categoría de los funtores $\cal{C}\to\Ens$ exactos por la izquierda; para
cada $X\in\cal{C}'(\pi)$ se toma el funtor $h(X)(E)=\Hom(X,E)$
\marginpar{128}
(donde $\Hom$ se toma en~$\cal{C}'(\pi)$). De las definiciones se sigue de
inmediato que los funtores~$h$ y~$g$ son adjuntos entre sí y que~$hg$ es
canónicamente isomorfo al funtor identidad de~$\Prodash\cal{C}(\pi)$. Para
probar que~$g$ y~$h$ son cuasiinversos, resta probar que todo objeto de
$\cal{C}'(\pi)$ es isomorfo a uno de la forma~$g(Q)$, con
$Q\in\Prodash\cal{C}(\pi)$; en otros términos: \emph{todo espacio~$X$
compacto y totalmente desconectado, provisto de una acción continua de
$\pi$, es isomorfo a un límite proyectivo de espacios finitos discretos con
operadores}. Como~$X$ es el límite proyectivo de sus cocientes finitos
discretos (como espacio topológico sin operadores), basta probar que entre
estos cocientes hay un sistema cofinal invariante bajo~$\pi$. Para ello
basta probar que, dado tal cociente~$X'$, el conjunto de sus transformados
por las operaciones de~$\pi$ es finito (se tomará entonces el supremo de
dichos transformados, que será un cociente invariante que domina a~$X'$).
Equivalentemente, basta hallar un subgrupo normal abierto~$\pi'$ de~$\pi$
cuyos elementos dejen invariante~$X'$. Ahora bien, $X'$ corresponde a una
partición finita de~$X$ en abiertos~$X_i$. Por la continuidad y la
compacidad de~$\pi$, existe un entorno~$V$ del elemento neutro de~$\pi$ tal
que $s\in V$ implica $s\cdot X_i\subset X_i$ para todo~$i$, y por tanto~$s$
deja invariante~$X'$. Como los subgrupos normales abiertos de~$\pi$ forman
un sistema fundamental de entornos del elemento neutro, queda concluida la
demostración.

Observemos que resulta aún más sencillo ver que la categoría
$\Ind\cal{C}(\pi)$ es canónicamente equivalente a la categoría de los
conjuntos sobre los que~$\pi$ actúa continuamente. No necesitaremos aquí
este hecho.

\begin{proposition}
\label{V.5.3}
Sean~$\cal{C}$ una categoría galoisiana, $F$ un funtor fundamental
\ifthenelse{\boolean{orig}}{sobre~$C$,}{sobre~$\cal{C}$,}
$P=(P_i)$ el proobjeto asociado, normalizado de la manera habitual. Sea
\ifthenelse{\boolean{orig}}{$X\in C$;}{$X\in\cal{C}$;}
para que~$X$ sea conexo, es necesario y suficiente que~$\pi$ actúe
transitivamente sobre~$E=F(X)$.
\end{proposition}

Se reduce al caso modelo $\cal{C}=\cal{C}(\pi)$, con~$F$ el funtor
canónico, donde la afirmación es trivial.

\begin{corollaire}
\label{V.5.4}
Son equivalentes las condiciones siguientes sobre~$X$: \textup{(i)} $X$ es
conexo y $\not\simeq\emptyset_{\cal{C}}$; \textup{(ii)} el grupo~$\pi$ actúa
transitivamente sobre~$E=F(X)$ y $F(X)\ne\emptyset$; \textup{(iii)} $X$ es
isomorfo a algún~$P_i$.
\end{corollaire}

La equivalencia
\marginpar{129}
entre (i) y (iii) se deduce también fácilmente del apartado e) del
\No~\Ref{V.4}.

\begin{proposition}
\label{V.5.5}
Sea $Q=(Q_i)_{i\in I}$ un proobjeto
\ifthenelse{\boolean{orig}}{de~$C$,}{de~$\cal{C}$,}
normalizado de la manera habitual, y sea~$G$ el funtor correspondiente
$G(X)=\Hom(Q,X)$ de~$\cal{C}$
\ifthenelse{\boolean{orig}}{(en~$\Ens$.}{(en~$\Ens$).}
Las condiciones siguientes son equivalentes:
\begin{enumerate}
\item[(i)] $G$ conmuta con las sumas directas finitas;
\item[(ii)] $G$ conmuta con la suma de dos objetos;
\item[(iii)] los~$Q_i$ son conexos y
\ifthenelse{\boolean{orig}}{$\not\simeq\emptyset_C$}
{$\not\simeq\emptyset_{\cal{C}}$};
\item[(iv)] $Q$ es isomorfo a~$\pi/H$, donde~$H$ es un subgrupo cerrado
de~$\pi$;
\item[(v)] el funtor~$G$ es isomorfo al funtor $E\mto E^H$ (conjunto de
los invariantes bajo~$H$) definido por un subgrupo cerrado~$H$ de~$\pi$.
\end{enumerate}
\end{proposition}

N.B.: en los enunciados de (iv) y (v) se supone elegido un funtor
fundamental que permita identificar~$\cal{C}$ con la categoría
$\cal{C}(\pi)$.

\subsubsection*{Demostración} Podemos suponer
$\cal{C}=\cal{C}(\pi)$. La implicación (i)$\To$(ii) es trivial, y
(ii)$\To$(iii) se demuestra como la propiedad d) del \No~\Ref{V.4}.
Probemos (iii)$\To$(iv). En efecto, se sabe que $\varprojlim_iQ_i$ no es
vacío, por ser un límite proyectivo de conjuntos finitos no vacíos. Sea~$a$
un punto de~$\varprojlim_iQ_i$; este define un homomorfismo de espacios con
operadores
$$
\pi\to Q
$$
que es \emph{sobreyectivo}, pues para todo~$i$ lo es la composición
$\pi\to Q\to Q_i$, ya que~$\pi$ actúa transitivamente sobre~$Q_i$ en virtud
de~\Ref{V.5.3}. Si~$H$ es el subgrupo estabilizador de~$a$, se obtiene así
un isomorfismo $\pi/H\isomto Q$. Las implicaciones (iv)$\To$(v) y
(v)$\To$(i) son nuevamente triviales.

\begin{proposition}
\label{V.5.6}
Sean~$\cal{C}$ una categoría galoisiana, $P$ un proobjeto fundamental
de~$\cal{C}$ y~$F$ el funtor fundamental asociado. Sea
$P'=(P'_i)_{i\in I}$ un proobjeto de~$\cal{C}$ puesto en forma normal, y
sea~$F'$ el funtor asociado $F'(X)=\Hom(P',X)$ de~$\cal{C}$ en~$\Ens$. Son
equivalentes las condiciones siguientes:
\begin{enumerate}
\item[(i)] $P'\simeq P$, o, equivalentemente, $F'\simeq F$;
\item[(ii)] $P'$ es fundamental, o, equivalentemente, $F'$ es fundamental;
\item[(iii)] $F'$ transforma la suma de dos objetos en su suma, y
$X\not\simeq\emptyset_{\cal{C}}$ implica
\ifthenelse{\boolean{orig}}{$F(X)\ne\emptyset$}{$F'(X)\ne\emptyset$};
\item[(iv)]
\marginpar{130}
los objetos de~$\cal{C}$ conexos y distintos de~$\emptyset_{\cal{C}}$ son
exactamente los objetos isomorfos a algún~$P'_i$.
\end{enumerate}
\end{proposition}

Se tienen trivialmente (i)$\To$(iii) y (i)$\To$(ii); además,
(ii)$\To$(iv) en virtud de~\Ref{V.5.4}, aplicado a~$P'$ en lugar de~$P$.
Por otra parte, (iii) o (iv) implica, en virtud de~\Ref{V.5.5}, que~$P'$ es
de la forma~$\pi/H$, donde~$H$ es un subgrupo cerrado de~$\pi$. En el caso
(iii), para cada subgrupo normal abierto~$\pi'$ de~$\pi$ existe un
$\pi$-homomorfismo $P'=\pi/H\to\pi/\pi'$; de aquí $H\subset\pi'$, luego
$H=(0)$ y, por consiguiente, (i), cqfd.

\begin{corollaire}
\label{V.5.7}
Sea~$\cal{C}$ una categoría galoisiana. Los proobjetos fundamentales son
isomorfos entre sí, y los funtores fundamentales son isomorfos entre sí.
\end{corollaire}

En otros términos, \emph{la categoría de los funtores fundamentales es un
grupoide conexo~$\Gamma$},
\label{indnot:ef}\oldindexnot{$\Gamma$|hyperpage}%
que puede llamarse el \emph{grupoide fundamental}
\index{grupoide fundamental|hyperpage}%
de la categoría galoisiana~$\cal{C}$. Si $\cal{C}=\cal{C}(\pi)$, el grupo
de automorfismos de un objeto del grupoide fundamental es isomorfo a~$\pi$,
y este isomorfismo está bien determinado salvo automorfismo interior.
(N.B.: se llama \emph{grupoide} a una categoría en la que todos los
morfismos son isomorfismos, y \emph{grupoide conexo} a un grupoide cuyos
objetos son todos isomorfos entre sí.) Los proobjetos fundamentales de
$\cal{C}$ forman un grupoide conexo equivalente al \emph{opuesto} del
grupoide fundamental. Si~$F,F'$ son dos funtores fundamentales asociados a
proobjetos fundamentales~$P,P'$, entonces
$\Hom(F,F')=\Isom(F,F')$ se denota a veces por~$\pi_{F',F}$ y desempeña el
papel de un «conjunto de \emph{clases de caminos} de~$F$ a~$F'$». En
particular, $\pi_{F,F}=\pi_F$ no es sino el \emph{grupo fundamental de
$\cal{C}$} con base~$F$ construido en el número precedente. En cuanto al
proobjeto~$P$ asociado a~$F$, desempeña el papel de un \emph{recubrimiento
universal con base~$F$} del objeto final~$e_{\cal{C}}$ de~$\cal{C}$.

Puede resultar cómodo describir~$\cal{C}$ (salvo equivalencia) en términos
de su grupoide fundamental~$\Gamma$, sin elegir un objeto particular~$F$
de este último. A cada objeto~$X$ de~$\cal{C}$ se asocia el funtor~$E_X$
sobre el grupoide fundamental definido por
$$
E_X(F)=F(X),
$$
con valores en~$\Ens$. (En topología, un funtor de esta clase se conoce con
el nombre de «sistema local» sobre el grupoide.) $F(X)=E_X(F)$ puede
llamarse la \emph{fibra} de~$X$ en~$F$, y~$E_X$ el funtor fibra asociado
a~$X$. El funtor~$E_X$ tiene la propiedad siguiente:
\marginpar{131}
\ifthenelse{\boolean{orig}}
{\emph{\underline{para todo} $F$, $E_X(F)$
\underline{es un conjunto finito sobre el que el grupo topológico}
$\pi_F=\Aut(F)$ \underline{actúa continuamente}}}
{\emph{para todo~$F$, $E_X(F)$ es un conjunto finito sobre el que el grupo
topológico $\pi_F=\Aut(F)$ actúa continuamente}}.
Para un funtor covariante dado~$\xi$ del grupoide fundamental en~$\Ens$, la
condición precedente equivale, por otra parte, a la misma condición para
\emph{un}~$F$ fijo cualquiera. Sentado esto:


\begin{proposition}
\label{V.5.8}
El funtor $X\mto E_X$ es una equivalencia de la categoría~$\cal{C}$ con la
categoría de los funtores covariantes del grupoide fundamental~$\Gamma$
de~$\cal{C}$ en~$\Ens$ que satisfacen la condición
\ifthenelse{\boolean{orig}}
{subrayada más arriba.}
{destacada más arriba.}
\end{proposition}

En efecto, sea~$F_0$ un objeto del grupoide fundamental y sea
$\pi_0=\Aut(F_0)$. El funtor $\xi\mto\xi(F_0)$ es entonces una equivalencia
de la segunda categoría considerada en~\Ref{V.5.8} con la categoría
$\cal{C}(\pi_0)$, como se comprueba inmediatamente. Por otra parte, la
composición de este último con el funtor $X\mto E_X$ es la equivalencia
natural $\cal{C}\to\cal{C}(\pi_0)$. Se sigue que el propio funtor
$X\mto E_X$ es una equivalencia.

\begin{corollaire}
\label{V.5.9}
La categoría $\Prodash\cal{C}$
es canónicamente equivalente a la categoría de los funtores covariantes
$\xi$ del grupoide fundamental~$\Gamma$ en la categoría de los espacios
topológicos que satisfacen la condición siguiente: para todo objeto~$F$
de~$\Gamma$, $\xi(F)$ es un espacio compacto totalmente desconectado sobre
el que actúa el grupo topológico~$\pi_F$.
\end{corollaire}

También aquí, para verificar esta condición sobre~$\xi$, basta verificarla
para \emph{un}~$F$. La demostración es la misma que la de~\Ref{V.5.8}.

\begin{remarque}
\label{V.5.10}
Sea $(F_s)_{s\in S}$ una familia
\ifthenelse{\boolean{orig}}{de objeto}{de objetos}
del grupoide fundamental~$\Gamma$. Para $s,s'\in S$, pongamos
$$
\Hom(s,s')=\Hom(F_s,F_{s'})
$$
de modo que~$S$ se convierte a su vez en un grupoide conexo y la aplicación
$s\mto F_s$ en un funtor plenamente fiel de~$S$ en~$\Gamma$, que
denotaremos por~$f$. Al considerar entonces el funtor $X\mto E_X\circ f$
de~$\cal{C}$ en la categoría de funtores $\Hom(S,\Ens)$, se obtiene una
variante de~\Ref{V.5.8} (y de~\Ref{V.5.9}) con~$\Gamma$ reemplazado por~$S$.
El enunciado así obtenido se reduce al teorema~\Ref{V.4.1} cuando~$S$ se
reduce a un punto, y no es sino el propio enunciado~\Ref{V.5.8} si~$S$ es
el conjunto de los objetos de~$\Gamma$.
\end{remarque}

Utilizaremos~\Ref{V.5.9} para definir un proobjeto canónico de~$\cal{C}$.
Para ello consideramos el funtor de~$\Gamma$ en la categoría de los espacios
topológicos
\marginpar{132}
(e incluso de los grupos topológicos):
$$
f\colon F\mto\Aut(F)=\pi_F.
$$
Este funtor satisface la condición considerada en~\Ref{V.5.8}: el espacio
con operadores~$f(F)$ bajo~$\pi_F$ no es sino~$\pi_F$, considerado como
espacio con operadores bajo sí mismo mediante los automorfismos interiores.
Por tanto, el funtor~$f$ corresponde a un proobjeto de~$\cal{C}$
determinado salvo isomorfismo único, que es incluso un progrupo de
$\cal{C}$ y se llama el \emph{progrupo fundamental de~$\cal{C}$},
\index{progrupo fundamental de una categoría galoisiana|hyperpage}%
el cual desempeña el papel de un sistema local de grupos fundamentales. Es,
por tanto, un progrupo~$\Pi$ de~$\cal{C}$ definido por la condición de que
exista un isomorfismo funtorial en~$F$
$$
F(\Pi)\simeq\pi_F.
$$

Si~$X$ es un proobjeto cualquiera de~$\cal{C}$, se tiene un morfismo
canónico
$$
\Pi\times X\to X
$$
que hace de~$X$ un objeto con grupo de operadores izquierdos
\ifthenelse{\boolean{orig}}{$G$}{$\Pi$}
en~$\Prodash\cal{C}$. Para verlo basta observar que, para~$F$ variable, se
tiene una aplicación canónica
$$
\Pi(F)\times X(F)\to X(F)
$$
\ie
$$
\Aut(F)\times E_X(F)\to E_X(F),\quad\text{o}\quad
\pi_F\times F(X)\to F(X)
$$
que es funtorial en~$F$. También es funtorial en~$X$; por consiguiente, para
todo morfismo $X\to Y$ de proobjetos, el diagrama correspondiente
$$
\xymatrix{
\Pi\times X\ar[r]\ar[d]&X\ar[d]\\
\Pi\times Y\ar[r]&Y
}
$$
es conmutativo.

\begin{remarque}
\label{V.5.11}
No debe confundirse un proobjeto fundamental~$P$ (que no está provisto de
una estructura de grupo y es conexo) con el progrupo fundamental (que es un
pro\emph{grupo} y, en general, no es conexo). Más precisamente,
\ifthenelse{\boolean{orig}}{$G$}{$\Pi$}
es conexo si
\ifthenelse{\boolean{orig}}{\ignorespaces}{y}
solo si~$\pi_F$, actuando sobre sí mismo mediante automorfismos interiores,
actúa transitivamente, \ie si~$\pi$ se reduce al elemento neutro o,
equivalentemente, si~$\cal{C}$ es equivalente a la categoría de los
conjuntos finitos. Otra diferencia
\marginpar{133}
esencial es que
\ifthenelse{\boolean{orig}}{$G$}{$\Pi$}
está determinado salvo isomorfismo único, mientras que~$P$ solo está
determinado salvo isomorfismo no único.

Sea~$E$ un conjunto finito y consideremos el funtor constante sobre el
grupoide~$\Gamma$, de valor~$E$; en virtud de~\Ref{V.5.8}, este define un
objeto de~$\cal{C}$, denotado por~$E_{\cal{C}}$, que puede interpretarse
también como la suma de~$E$ ejemplares del objeto final~$e_{\cal{C}}$
de~$\cal{C}$. Puede considerarse~$E_{\cal{C}}$ como un funtor en~$E$, de la
categoría de los conjuntos finitos en la categoría~$\cal{C}$, y este funtor
es \emph{exacto}; por tanto, transforma grupos finitos en
$\cal{C}$-grupos,
\ifthenelse{\boolean{orig}}{etc...}{etc.}
Así pues, si~$X$ es un objeto de~$\cal{C}$ sobre el que el grupo finito~$G$
actúa por la derecha, puede considerarse~$X$ como un objeto de~$\cal{C}$
provisto del $\cal{C}$-grupo de operadores derechos~$G_{\cal{C}}$. Por
extensión de la terminología general relativa a objetos con grupos de
operadores que son $\cal{C}$-grupos, se dirá que~$X$ es \emph{formalmente
homogéneo principal}
\index{formalmente homogéneo principal (objeto)|hyperpage}%
bajo~$G$ si es formalmente homogéneo principal bajo~$G_{\cal{C}}$, \ie si
el morfismo canónico
$$
X\times G_{\cal{C}}\to X\times X
$$
deducido de la acción derecha de~$G_{\cal{C}}$ sobre~$X$, es un
isomorfismo. Se dice que~$X$ es \emph{homogéneo principal}
\index{homogéneo principal (objeto)|hyperpage}%
bajo~$G$ si lo es bajo~$G_{\cal{C}}$, \ie si lo es formalmente y, además,
el cociente $X/G=X/G_{\cal{C}}$ es~$e_{\cal{C}}$. Si se fija un funtor
fundamental, y por tanto una equivalencia de~$\cal{C}$ con una categoría
$\cal{C}(\pi)$, el objeto~$X$ corresponde a un conjunto sobre el que~$\pi$
actúa continuamente por la izquierda, sea~$E=F(X)$. Hacer actuar~$G$ por la
derecha sobre~$X$ equivale entonces a hacerlo actuar por la derecha sobre
el conjunto~$E$, de modo que las operaciones de~$G$ conmuten con las
de~$\pi$. Se comprueba de inmediato que~$X$ es homogéneo principal bajo~$G$
si y solo si el conjunto~$E$ es un espacio homogéneo principal
\ifthenelse{\boolean{orig}}
{bajo~$G$}
{bajo~$G$,}
\ie si y solo si~$G$ actúa sobre él de manera simplemente transitiva. (Por
otra parte, $X$ es formalmente homogéneo principal \sss\ 
$E$ es homogéneo principal \emph{o} vacío.) Comparando con~\Ref{V.5.3}, se
ve que, si~$X$ es homogéneo principal bajo~$G$ \emph{y conexo}, entonces el
homomorfismo dado de~$G$ en el grupo opuesto a~$\Aut(X)$ es un
\emph{isomorfismo}; además, para que un objeto~$X$
\ifthenelse{\boolean{orig}}{de~$C$}{de~$\cal{C}$}
sea conexo y homogéneo principal bajo el grupo opuesto a~$\Aut(X)$, es
necesario y suficiente, con las notaciones del \No~\Ref{V.4}, que sea
isomorfo a un~$P_i$ \emph{galoisiano}. En el caso modelo
$\cal{C}=\cal{C}(\pi)$, esto significa que~$X$ es isomorfo a un cociente
de~$\pi$ por un subgrupo normal.

Supongamos aún dado un funtor fundamental~$F$. Entonces, dar un
\marginpar{134}
objeto~$X$ homogéneo principal bajo un grupo finito~$G$ que actúa por la
derecha, junto con un punto $a\in F(X)$, equivale a dar un homomorfismo
de~$\pi$ en el grupo~$G$. En efecto, a tal homomorfismo se le hace
corresponder el conjunto~$E=G$, haciendo actuar~$\pi$ por la izquierda
mediante el homomorfismo dado $\pi\to G$ y las traslaciones izquierdas
de~$G$, y haciendo actuar~$G$ por la derecha mediante traslación derecha;
el punto marcado~$a$ de~$E$ es el elemento neutro de~$G$. Gracias a lo
anterior se obtiene así, de manera esencialmente única, todo triple
$(X,G,a)$ con las propiedades consideradas,
\ifthenelse{\boolean{orig}}
{puesto que un}
{puesto que un}
conjunto homogéneo principal bajo un grupo~$G$ y provisto de un punto
marcado se identifica con dicho grupo. De este modo se obtiene una
interpretación geométrica directa del funtor $G\mto\Hom(\pi,G)$ de la
categoría de los grupos finitos en~$\Ens$, funtor que es prorepresentable
mediante~$\pi$ y cuya consideración proporcionaría, por tanto, otra
construcción del grupo~$\pi$ asociado a~$F$.
\end{remarque}


\section{Funtores exactos entre categorías galoisianas}
\label{V.6}

\begin{proposition}
\label{V.6.1}
Sean $\cal{C}$ y $\cal{C'}$ dos categorías galoisianas,
$H\colon \cal{C}\to \cal{C'}$ un funtor covariante, $F'$ un funtor
fundamental sobre~$\cal{C'}$ y $F=F'\circ H$. Son equivalentes las
condiciones siguientes:
\begin{enumerate}
\item[(i)] $H$ es \emph{exacto}, \ie exacto por la izquierda y exacto
por la derecha.
\item[(ii)] $H$ es exacto por la izquierda, transforma sumas finitas en
sumas finitas y epimorfismos en epimorfismos (o, equivalentemente,
objetos $\not\approx\emptyset_{\cal C}$ en objetos
$\not\approx\emptyset_{\cal C'}$).
\item[(iii)] $F$ es un funtor fundamental sobre~$\cal C$.
\end{enumerate}
\end{proposition}

La implicación (i)$\To$(ii) es un hecho general de la teoría de categorías. Además,
la primera forma dada a (ii) implica la segunda, como se ve observando que,
si~$X$ es un objeto de~$\cal C$, entonces
$X\not\approx\emptyset_{\cal C}$ \sss\ el morfismo
$X\to e_{\cal C}$ es un epimorfismo; se observará que~$F$, al suponerse
exacto por la izquierda, transforma $e_{\cal C}$ en~$e_{\cal C'}$. La
segunda forma dada a (ii) implica (iii), pues~$F$, al ser exacto por la
izquierda y, por tanto, prorepresentable, satisface el criterio (iii)
de~\Ref{V.5.6}. Por último, (iii) implica (i), como se sigue de que~$F$ es
exacto y ``conservativo'' (\ie satisface el axioma (G~6) del
\No~\Ref{V.4}).

Sea entonces~$\Gamma$ el grupoide fundamental de~$\cal C$ y~$\Gamma'$ el
de~$\cal C'$. Así, si
\marginpar{135}
$H$ es exacto, entonces $F'\mto F'\circ H$ es un funtor del
grupoide~$\Gamma'$ en el grupoide~$\Gamma$, que denotaremos por~$\tH$:
$$
\tH(F')(X)=F'(H(X)).
$$
También puede escribirse, con la notación $F(X)=E_X(F)$ introducida en el
\ifthenelse{\boolean{orig}}
{\No~\Ref{V.6}:}
{\No~\Ref{V.5}:}
$$
\ifthenelse{\boolean{orig}}
{E_{H(X)}(F')=E_X(\tH(F'))}
{E_{H(X)}(F')=E_X(\tH(F')).}
$$
Esta última fórmula muestra, teniendo en cuenta~\Ref{V.5.8}
o~\Ref{V.4.1}, que el funtor exacto~$H$ queda determinado (salvo
isomorfismo único) cuando se conoce el funtor correspondiente~$\tH$.
Fijemos un~$F'$ y sea $F=\tH(F')$; entonces~$\tH$ define un homomorfismo
de~$\pi_{F'}=\Aut(F')$ en~$\pi_F=\Aut(F)$:
$$
\tH\colon\pi_{F'}\to\pi_F
\qquad(F=\tH(F')=F'\circ H).
$$
Por otra parte, la fórmula anterior muestra (teniendo en cuenta~\Ref{V.5.8})
que este homomorfismo tiene la propiedad siguiente: para todo conjunto
finito~$E$ sobre el que~$\pi_F$ actúa continuamente, el grupo~$\pi_{F'}$
actúa también \emph{continuamente} por medio del homomorfismo anterior
$\pi_{F'}\to\pi_F$. Aplicándolo a los cocientes de~$\pi_F$ por sus
subgrupos normales abiertos, se ve que la condición anterior significa
también que el homomorfismo considerado es continuo. Recíprocamente,
démonos un objeto~$F$ de~$\Gamma$, un objeto~$F'$ de~$\Gamma'$ y un
homomorfismo continuo
$$
u\colon\pi_{F'}\to\pi_F.
$$
Le corresponde entonces un funtor de~$\cal C(\pi)$ en~$\cal C(\pi')$,
manifiestamente exacto; por tanto, en virtud de~\Ref{V.4.1}, le corresponde
un funtor $H$ de~$\cal C$ en~$\cal C'$ que es exacto y tal que
$\tH\colon\pi_{F'}\to\pi_F$ es precisamente~$u$. En lugar de partir de un
homomorfismo de grupos, también puede partirse de un \emph{funtor}
$$
U\colon\Gamma'\to\Gamma
$$
tal que, para \emph{todo} $F'\in\Gamma'$ (o para \emph{un}
$F'\in\Gamma'$, lo que es lo mismo), el homomorfismo correspondiente
$\pi_{F'}\to\pi_F$ sea continuo: tal funtor es isomorfo a un funtor de la
forma~$\tH$, donde $H\colon\cal C\to\cal C'$ es un funtor exacto
determinado salvo isomorfismo único. Así:

\begin{corollaire}
\label{V.6.2}
Para
\marginpar{136}
que un funtor $H\colon\cal C\to\cal C'$ de categorías galoisianas sea
exacto, es necesario y suficiente que existan equivalencias
$\cal C(\pi)\to\cal C$ y
\ifthenelse{\boolean{orig}}
{$\cal C'\to\cal C(\pi)$}
{$\cal C'\to\cal C(\pi')$}
que transformen el funtor~$H$ en el funtor
$\cal C(\pi)\to\cal C(\pi')$ asociado a un homomorfismo de grupos
topológicos $\pi'\to\pi$.
\end{corollaire}

\begin{corollaire}
\label{V.6.3}
Sean $\cal C$ y $\cal C'$ dos categorías galoisianas y sean~$\Gamma$ y
$\Gamma'$ sus grupoides fundamentales. Entonces la categoría de los
funtores exactos de~$\cal C$ en~$\cal C'$ es equivalente a la categoría de
los funtores $U\colon\Gamma'\to\Gamma$ que tienen la propiedad siguiente:
para todo $F'$ de~$\Gamma'$ (o para \emph{un} $F'$ de~$\Gamma'$, lo que es
lo mismo), poniendo $F=U(F')$, el homomorfismo
$$
\pi_{F'}=\Aut(F')\to\pi_F=\Aut(F)
$$
definido por~$U$ es continuo.
\end{corollaire}

Consideremos el progrupo fundamental~$\Pi$ de~$\cal C$. Entonces un funtor
exacto~$H$ lo transforma en un progrupo $H(\Pi)$ de~$\cal C'$. Vamos a
definir un homomorfismo
$$
\Pi'\to H(\Pi)
$$
(donde~$\Pi'$ es el progrupo fundamental de~$\cal C'$), mediante la
condición de que, para todo objeto~$F'$ de~$\Gamma'$, el homomorfismo
correspondiente
$$
F'(\Pi')=\pi_{F'}\to F'(H(\Pi))=\pi_F\qquad
\ifthenelse{\boolean{orig}}
{\text{(donde $F=F'\circ H=\tH(F')$)};}
{\text{(donde $F=F'\circ H=\tH(F')$)}}
$$
sea el homomorfismo natural
$$
\Aut(F')\to\Aut(F'\circ H).
$$
Como este último es funtorial en~$F'$, define efectivamente, en virtud
de~\Ref{V.5.8}, un homomorfismo de proobjetos, y de hecho de progrupos,
en~$\cal C'$. Este homomorfismo se llama \emph{asociado} al funtor~$H$.

Sea ahora~$H'$ un segundo funtor exacto, de la categoría galoisiana
$\cal C'$ en una categoría galoisiana~$\cal C''$. Es inmediato que
$$
{}^t(H'H)=\tH\,\tH'
$$
(N.B.: se trata de una identidad de funtores, y no solamente de un
isomorfismo canónico). Hay una propiedad de transitividad análoga para los
homomorfismos asociados de los progrupos fundamentales.

Vamos
\marginpar{137}
a interpretar ahora las propiedades del funtor exacto~$H$ en términos del
homomorfismo correspondiente
$$
u\colon\pi_{F'}\to\pi_F\qquad
\ifthenelse{\boolean{orig}}
{\text{(donde $F'=F'\circ H$)}.}
{\text{(donde $F=F'\circ H$)}.}
$$
Conviene introducir la noción de \emph{objeto punteado}
\index{punteado (objeto)|hyperpage}%
de la categoría galoisiana~$\cal C$ (provista de su funtor fundamental~$F$):
es, por definición, un objeto~$X$ de~$\cal C$ provisto de un elemento~$a$
de~$F(X)$. Se interpreta, pues, como un conjunto finito sobre el que
$\pi_F$ actúa continuamente por la izquierda, provisto de un punto~$a$.
Por tanto, los objetos punteados \emph{conexos} de~$\cal C$ se identifican,
en virtud de~\Ref{V.5.3}, con los subgrupos abiertos de~$\pi_F$. Si~$U$ y
$V$ son dos de tales subgrupos, correspondientes a objetos punteados
conexos~$X$ e~$Y$ de~$\cal C$, entonces existe un homomorfismo punteado
de~$X$ en~$Y$ si y solo si $U\subset V$, y dicho homomorfismo es entonces
único. Por supuesto, el funtor~$H$ transforma objetos punteados en objetos
punteados (pues $F=F'\circ H$). Obsérvese, por otra parte, que un subgrupo
cerrado de un grupo como~$\pi_F$ es la intersección de los subgrupos
abiertos que lo contienen; por consiguiente,
\ifthenelse{\boolean{orig}}{\ignorespaces}{si}
$M$ y~$N$ son dos subgrupos cerrados, entonces $M\subset N$ si y solo si
todo subgrupo abierto que contiene a~$N$ contiene también a~$M$. Gracias a
estas observaciones se demuestran fácilmente los resultados siguientes:


\begin{proposition}
\label{V.6.4}
Sea~$X$ un objeto punteado conexo de~$\cal C$, asociado a un subgrupo
abierto~$U$ de~$\pi_F$. Para que
\ifthenelse{\boolean{orig}}{$u$}{$U$}
contenga a~$u(\pi_{F'})$ (\resp al subgrupo normal cerrado engendrado
por~$u(\pi_{F'})$), es necesario y suficiente que~$H(X)$ admita una
sección punteada (\resp esté completamente descompuesto).
\end{proposition}

Se llamará \emph{sección} ---sobrentendido: sobre el objeto final--- de un
objeto~$X$ de una categoría galoisiana~$\cal C$ a un morfismo del objeto
final~$e_{\cal C}$ en~$X$, lo que equivale a dar un elemento~$a$ de~$F(X)$
invariante por~$\pi_F$; si~$X$ es punteado, se dice que se tiene una
\emph{sección punteada} si es compatible con las estructuras punteadas de
$X$ y~$e_{\cal C}$, \ie si~$a$ es precisamente el
\ifthenelse{\boolean{orig}}
{objeto marcado}
{elemento marcado}
de~$F(X)$. Tal sección es, pues, única, y existe si y solo si el
\ifthenelse{\boolean{orig}}
{objeto marcado}
{elemento marcado}
de~$F(X)$ es invariante por~$\pi_F$. Por último, se dice que un objeto de
una categoría galoisiana está \emph{completamente descompuesto}
\index{completamente descompuesto (objeto)|hyperpage}%
\index{descompuesto (objeto completamente)|hyperpage}%
si es isomorfo a una suma de objetos finales, \ie si~$\pi_F$ actúa
trivialmente sobre~$F(X)$, condición evidentemente más fuerte que la
existencia de una sección punteada cuando~$X$ es punteado. La
Proposición~\Ref{V.6.4} resulta inmediatamente de las definiciones y
observaciones anteriores.

\begin{corollaire}
\label{V.6.5}
Para
\marginpar{138}
que~$u$ sea trivial, es necesario y suficiente que, para todo objeto~$X$
de~$\cal C$, $H(X)$ esté completamente descompuesto.
\end{corollaire}

\begin{proposition}
\label{V.6.6}
Sea~$X'$ un objeto punteado conexo de~$\cal C'$, asociado a un subgrupo
abierto~$U'$ de~$\pi_{F'}$. Para que~$U'$ contenga a~$\Ker u$, es necesario
y suficiente que exista un objeto punteado conexo~$X$ de~$\cal C$ y un
homomorfismo punteado de la componente conexa punteada~$X'_0$ de~$H(X)$
en~$X'$ (por tanto, que
\ifthenelse{\boolean{orig}}{$X$}{$X'$}
sea isomorfo, como objeto punteado, a un cociente de la componente conexa
neutra de la imagen inversa de un objeto punteado de~$\cal C$). Si~$u$ es
sobreyectivo, la condición anterior equivale también a la siguiente:
\ifthenelse{\boolean{orig}}{$X$}{$X'$}
es isomorfo a un objeto de la forma~$H(X)$, donde~$X$ es un objeto punteado
de~$\cal C$.
\end{proposition}

(Se llama \emph{componente conexa neutra}
\index{neutra (componente conexa de un objeto punteado)|hyperpage}%
de un objeto punteado~$X$ de una categoría galoisiana~$\cal C$ al único
subobjeto conexo punteado de~$X$; en virtud de~\Ref{V.5.3}, corresponde a
la órbita bajo~$\pi_F$ del punto marcado de~$F(X)$). Como el hecho de que
$U'$ contenga a~$\Ker u$ no depende del punto elegido en~$X'$ (pues elegir
otro punto equivale a reemplazar
\ifthenelse{\boolean{orig}}{$U$}{$U'$}
por un subgrupo conjugado), se obtiene:

\begin{corollaire}
\label{V.6.7}
Para que~$U'$ contenga a~$\Ker u$, es necesario y suficiente que exista un
objeto~$X$ de~$\cal C$ (que puede suponerse conexo) y un morfismo de una
componente conexa de~$H(X)$ en~$X'$. Si~$u$ es sobreyectivo, esto significa
también que~$X'$ es isomorfo a un objeto de la forma~$H(X)$.
\end{corollaire}

\begin{corollaire}
\label{V.6.8}
Para que~$u$ sea inyectivo, es necesario y suficiente que, para todo
objeto~$X'$ de~$\cal C'$, exista un objeto~$X$ de~$\cal C$ y un
homomorfismo de una componente conexa de~$H(X)$ en~$X'$.
\end{corollaire}

\begin{proposition}
\label{V.6.9}
Las condiciones siguientes son equivalentes:
\begin{enumerate}
\item[(i)] El homomorfismo $u\colon\pi_{F'}\to\pi_F$ es sobreyectivo.
\item[(ii)] Para todo objeto conexo~$X$ de~$\cal C$, $H(X)$ es conexo.
\item[(iii)] El funtor~$H$ es plenamente fiel.
\end{enumerate}
\end{proposition}

Esta última condición significa que, para dos objetos~$X$, $Y$ de~$\cal C$,
la aplicación natural
$$
\Hom(X,Y)\to\Hom(H(X),H(Y))
$$
es biyectiva.

\begin{corollaire}
\label{V.6.10}
Para
\marginpar{139}
que~$u$ sea un isomorfismo, es necesario y suficiente que~$H$ sea una
equivalencia de categorías, o, equivalentemente, que se verifiquen las dos
condiciones siguientes:
\begin{enumerate}
\item[a)] para todo objeto conexo~$X$ de~$\cal C$, $H(X)$ es conexo;
\item[b)] todo objeto de~$\cal C'$ es isomorfo a un objeto de la
forma~$H(X)$.
\end{enumerate}
\end{corollaire}

\begin{proposition}
\label{V.6.11}
Sean $H\colon\cal C\to\cal C'$ y
\ifthenelse{\boolean{orig}}
{$H'\colon\cal C\to\cal C''$}
{$H'\colon\cal C'\to\cal C''$}
funtores exactos entre categorías galoisianas; sea~$F''$ un funtor
fundamental sobre~$\cal C''$, pongamos $F'=F''H'$ y $F=F'H$, y consideremos
los homomorfismos asociados
$$
u'\colon\pi_{F''}\to\pi_{F'}\qquad
u\colon\pi_{F'}\to\pi_F.
$$
Para que
\ifthenelse{\boolean{orig}}
{$\Ker u\subset\Im u'$}
{$\Ker u\supset\Im u'$},
\ie para que~$uu'$ sea el homomorfismo trivial, es necesario y suficiente
que, para todo objeto~$X$ de~$\cal C$, $H'(H(X))$ esté completamente
descompuesto. Para que
\ifthenelse{\boolean{orig}}
{$\Ker u\supset\Im u'$}
{$\Ker u\subset\Im u'$},
es necesario y suficiente que, para todo objeto punteado conexo~$X'$ de
$\cal C'$ tal que~$H'(X')$ admita una sección punteada, exista un objeto~$X$
de~$\cal C$ y un homomorfismo de una componente conexa de~$H(X)$ en~$X'$.
\end{proposition}

La primera afirmación resulta de la última afirmación de~\Ref{V.6.4}. La
segunda resulta de la conjunción de~\Ref{V.6.4} y~\Ref{V.6.6}.

\begin{remarque}
\label{V.6.12}
En general, bajo las condiciones de~\Ref{V.6.8}, no es cierto que~$X'$ sea
isomorfo a un objeto de la forma~$H(X)$. Puede demostrarse que, para que
todo objeto conexo (y, por tanto, todo objeto) de~$\cal C'$ sea isomorfo a
un objeto de la forma~$H(X)$, es necesario y suficiente que~$u$ sea un
isomorfismo de~$\pi_{F'}$ sobre un subgrupo \emph{factor directo}
de~$\pi_F$. Sin embargo, en la práctica se construye directamente un
homomorfismo $\pi_F\to\pi_{F'}$
\ifthenelse{\boolean{orig}}
{inverso por la derecha}
{inverso por la izquierda}
de~$u$ mediante un funtor exacto conveniente de~$\cal C'$ en~$\cal C$.
\end{remarque}

\begin{proposition}
\label{V.6.13}
Sean~$\cal C$ una categoría galoisiana provista de un funtor fundamental~$F$,
$S$ un objeto conexo de~$\cal C$ y~$\cal C'$ la categoría de los objetos
de~$\cal C$ sobre~$S$. Entonces~$\cal C'$ es una categoría galoisiana, y el
funtor $X\mto H(X)=X\times S$ de~$\cal C$ en~$\cal C'$ es exacto. Sea
$a\in F(S)$, y sea~$F'$ el funtor de~$\cal C'$ en la categoría de los
conjuntos finitos definido por
$$
F'(X')=\text{preimagen de~$a$ por }F(X')\to F(S).
$$
Entonces se tiene un isomorfismo $F\cong F'\circ H$, y el homomorfismo
correspondiente
$$
u\colon\pi_{F'}\to\pi_F
$$
\marginpar{140}%
es un isomorfismo de~$\pi_{F'}$ sobre el subgrupo abierto~$U$ de~$\pi_F$
estabilizador del elemento marcado~$a$ de
\ifthenelse{\boolean{orig}}{$F(X)$}{$F(S)$}.
\end{proposition}

La demostración se deja al lector.

\section{Caso de los preesquemas}
\label{V.7}
Sea~$S$ un preesquema localmente noetheriano y \emph{conexo}, y sea
$$
a\colon \Spec(\Omega)\to S
$$
un punto geométrico de~$S$, con valores en un cuerpo algebraicamente
cerrado~$\Omega$. Pongamos
$$
\cal{C}=\text{categoría de los recubrimientos étales de~$S$}
$$
y, para un objeto~$X$ de~$\cal{C}$, \ie un recubrimiento étale~$X$
de~$S$, pongamos
$$
F(X)=\text{conjunto de los puntos geométricos de~$X$ situados sobre~$a$.}
$$
Así, $F$ se convierte en un funtor de~$\cal{C}$ en la categoría de los
conjuntos finitos. Se satisfacen las propiedades (G~1)
a (G~6): para (G~1), esto está contenido en los sorites de
\Ref{I.4.6}; (G~2) resulta de~\Ref{V.3.4}; (G~3), de~\Ref{V.3.5};
(G~4) es trivial por definición; (G~5) resulta de~\Ref{V.3.5} y del
comienzo del \No~\Ref{V.2}; por último, (G~6) se demuestra
en~\Ref{V.3.7}. Pueden aplicarse, pues, los resultados de los
\No~\Ref{V.4}, \Ref{V.5}, \Ref{V.6}. Esto permite, en particular,
definir un proobjeto~$P$ de~$\cal{C}$ que representa a~$F$, llamado
\emph{recubrimiento universal de~$S$ en el punto~$a$},
\index{recubrimiento universal de un preesquema en un punto|hyperpage}%
y un grupo topológico~$\pi=\Aut(F)=\Aut^0 P$, llamado
\emph{grupo fundamental de~$S$ en~$a$},
\index{grupo fundamental de un preesquema en un punto|hyperpage}%
y denotado por~$\pi_1(S,a)$.
\label{indnot:eg}\oldindexnot{$\pi_1(S,a)$|hyperpage}%
El funtor~$F$ define entonces una equivalencia entre la categoría~$\cal{C}$
y la categoría de los conjuntos finitos sobre los que
$\pi=\pi_1(S,a)$ actúa continuamente. Esta equivalencia permite interpretar
las operaciones usuales de límites proyectivos y límites inductivos finitos
de recubrimientos (productos, productos fibrados, sumas, paso al cociente,
\ifthenelse{\boolean{orig}}{etc...}{etc.}) en términos de las operaciones
análogas en~$\cal{C}(\pi)$, \ie de las operaciones evidentes sobre conjuntos
finitos en los que actúa~$\pi$. Obsérvese además que, como las componentes
conexas topológicas de un recubrimiento étale son también recubrimientos
étales, \emph{un objeto~$X$ de~$\cal{C}$ es conexo en~$\cal{C}$ si y solo si
es topológicamente conexo}; en virtud de~\Ref{V.5.3}, esto significa que
$\pi_1$ actúa transitivamente sobre~$F(X)$.
\marginpar{141}%
Obsérvese que, para que un objeto~$X$ de~$\cal{C}$ sea fielmente plano y
cuasicompacto sobre~$S$ (puesto que ya es plano y cuasicompacto sobre~$S$),
es necesario y suficiente que~$X\to S$ sea sobreyectivo, \ie que sea un
epimorfismo en~$\cal{C}$, o también que~$X\neq\emptyset$. Del
criterio~\Ref{V.2.6}~(iii) se concluye entonces que~$X$
\emph{es un recubrimiento principal de
\ifthenelse{\boolean{orig}}{$X$}{$S$},
de grupo~$G$, si y solo si es un espacio homogéneo principal bajo~$G$ en
la categoría~$\cal{C}$}, tal como se definió en el \No~\Ref{V.5}.

Si~$a'$ es otro punto geométrico de~$S$ (correspondiente a un cuerpo
algebraicamente cerrado~$\Omega'$, que puede ser distinto de~$\Omega$ e
incluso tener característica distinta), define un funtor fibra
$F'=F_{a'}$ de~$\cal{C}$ en la categoría de los conjuntos finitos, que
sigue siendo exacto y, por tanto, es isomorfo a~$F=F_a$. En consecuencia,
los grupos fundamentales~$\pi_1(S;a)$
\label{indnot:eg1}\oldindexnot{$\pi_1(S,a)$|hyperpage}%
al variar~$a$ son isomorfos entre sí. Si se denota por
$\pi_1(S;a,a')$
\label{indnot:eh}\oldindexnot{$\pi_1(S;a,a')$|hyperpage}%
el conjunto de los isomorfismos (o, lo que equivale a lo mismo,
\ifthenelse{\boolean{orig}}{el conjunto}{el conjunto}
de los homomorfismos) $F_a\to F_{a'}$ de los funtores fibra asociados,
se obtiene así un \emph{grupoide} cuyo conjunto de objetos es el conjunto
de los puntos geométricos de~$S$, siendo los grupos fundamentales los
grupos de automorfismos de los objetos de dicho grupoide. El conjunto
$\pi_1(S;a',a)$ puede llamarse \emph{conjunto de las clases de caminos
de~$a$ a~$a'$}. Estas clases se componen, pues, de manera evidente.
Por último, puede definirse un progrupo~$\Pi_1^S$ de~$\cal{C}$, que podrá
llamarse \emph{progrupo fundamental de~$S$}
\index{progrupo fundamental de un preesquema|hyperpage}%
o \emph{sistema local de los grupos fundamentales sobre~$S$},
\index{sistema local de los grupos fundamentales|hyperpage}%
definido, salvo isomorfismo único, por la condición de que exista un
isomorfismo, funtorial en el punto geométrico~$a$ de~$S$:
$$
F_a(\Pi_1^S)=\pi_1(S;a)
$$
(\cf la Observación~\Ref{V.5.10}). En particular, si~$s$ es un punto
ordinario de~$S$, la fibra de
\ifthenelse{\boolean{orig}}{$G$}{$\Pi_1^S$}
en~$s$ es un progrupo sobre~$\kres(s)$, límite proyectivo de grupos finitos
étales sobre~$\kres(s)$; este progrupo podría llamarse \emph{grupo
fundamental de~$S$ en el punto ordinario~$s$ de~$S$} y denotarse
por~$\pi_1(S,s)$. Por definición, estos puntos
con valores en una extensión algebraicamente cerrada~$\Omega$
de~$\kres(s)$ son los elementos de~$\pi_1(S;a)$, donde~$a$ es el punto
geométrico de~$S$ definido por dicha extensión. En particular (tomando
como~$S$ el espectro de un cuerpo), a todo cuerpo~$k$ se asocia canónica y
funtorialmente un progrupo sobre~$k$, que podría denotarse por~$\pi_1(k)$,
límite proyectivo de grupos finitos étales sobre~$k$, cuyos puntos con
valores en una extensión algebraicamente cerrada~$\Omega$ de~$k$ se identifican
\marginpar{142}%
con los elementos del grupo de Galois topológico de~$\bar{k}/k$, donde
$\bar{k}$ es la
\ifthenelse{\boolean{orig}}{clausura de Galois}{clausura separable}
de~$k$ en~$\Omega$ (\cf~\Ref{V.8.1}). Este grupo~$\pi_1(k)$ no parece haber
atraído todavía la atención de los algebristas.

Sea ahora
$$
f\colon S'\to S
$$
un morfismo de un preesquema conexo localmente noetheriano a otro;
sea~$a'$ un punto geométrico de~$S'$ y sea~$a=f(a')$ su imagen directa
en~$S$. Entonces el funtor de imagen inversa induce un funtor de la
categoría~$\cal{C}(S)$
\label{indnot:ei}\oldindexnot{$\cal{C}(S)$|hyperpage}%
de los recubrimientos étales de~$S$ en la categoría~$\cal{C}(S')$ de los
recubrimientos étales de~$S'$:
$$
f^\bbullet\colon\cal{C}(S)\to\cal{C}(S').
$$
Se tiene además un isomorfismo de funtores
$$
F_a\cong F_{a'}\circ f^\bbullet,
$$
de modo que~$f^\bbullet$ es un funtor \emph{exacto}, al que se aplican los
resultados del \No~\Ref{V.6}. Se tiene, en particular, un homomorfismo
canónico
$$
u=\pi_1(f;a')\colon\pi_1(S',a')\to\pi_1(S;a)\qquad
\ifthenelse{\boolean{orig}}{(a'=f(a))}{(a=f(a'))}.
$$
\label{indnot:ej}\oldindexnot{$\pi_1(f;a')$|hyperpage}%
Este homomorfismo permite reconstruir el funtor de imagen inversa como una operación de
restricción de los grupos de operadores. Las propiedades del
funtor~$f^\bbullet$ se expresan, pues, de manera sencilla mediante las
propiedades del homomorfismo de grupos
\ifthenelse{\boolean{orig}}{asociado}{asociado},
como se explicitó en el \No~\Ref{V.6}. En particular, si~$S'$ es un
recubrimiento étale de~$S$, el homomorfismo~$u$ es un isomorfismo
de~$\pi_1(S',a')$ sobre el subgrupo abierto de~$\pi_1(S,a)$ que define el
recubrimiento étale conexo punteado~$S'$ de~$S$ (\ie el estabilizador~$U$
de~$a'\in F_a(S')$ en
\ifthenelse{\boolean{orig}}{$\pi_1(S;A)$}{$\pi_1(S;a)$}).

Si se quieren interpretar los homomorfismos~$\pi_1(f;a')$ para un punto
geométrico variable~$a'$, hay que considerar, de acuerdo con lo dicho en
el \No~\Ref{V.6}, un homomorfismo
$$
\Pi_1(f)\colon\Pi_1^{S'}\to f^\bbullet(\Pi_1^S)
$$
de progrupos sobre
\ifthenelse{\boolean{orig}}{$S$}{$S'$},
y tomar el homomorfismo correspondiente en las fibras geométricas.

\section{Caso de un preesquema base normal}
\label{V.8}
\marginpar{143}%
\begin{proposition}
\label{V.8.1}
Sea~$S$ el espectro de un cuerpo~$k$, y sea~$\Omega$ una extensión
algebraicamente cerrada de~$k$ que define un punto geométrico~$a$ de~$S$
con valores en~$\Omega$. Sea~$\bar{k}$ la clausura separable de~$k$
en~$\Omega$. Entonces existe un isomorfismo canónico de~$\pi_1(S,a)$
sobre el grupo de Galois topológico de~$\bar{k}/k$.
\end{proposition}

Sea~$k'$ la clausura algebraica de~$k$ en~$\Omega$; corresponde, pues, a un
punto geométrico~$b$ de~$S$, con valores en~$k'$. El homomorfismo natural
de funtores~$F_b\to F_a$ es evidentemente un isomorfismo, pues todo
$k$-homomorfismo de una extensión finita separable de~$k$ en~$\Omega$
toma necesariamente sus valores en~$\bar{k}$ y, con mayor razón, en~$k'$.
Por otra parte, el grupo~$\pi'$ de los $k$-automorfismos de~$k'$ actúa de
manera evidente sobre~$F_b$, de donde se obtiene un homomorfismo
$\pi'\to\Aut(F_b)\isomto\Aut(F_a)=\pi_1(S;a)$.
Por otra parte, es bien sabido que el homomorfismo natural de~$\pi'$ en el
grupo~$\pi$ de los automorfismos de~$\bar{k}/k$ es un isomorfismo. Se
obtiene así un homomorfismo canónico~$\pi\to\pi_1(S;a)$; queda por demostrar
que es un isomorfismo. En efecto, este homomorfismo es inyectivo, pues un
elemento del núcleo es un automorfismo de~$\bar{k}/k$ que induce la identidad
en toda subextensión finita separable y, por tanto, es trivial. Este
homomorfismo es sobreyectivo, pues si~$X$ es un recubrimiento étale
\emph{conexo} de~$S$, definido, por tanto, por una \emph{extensión} finita
separable~$L/k$, entonces~$\pi$ actúa transitivamente sobre el conjunto de
los $k$-homomorfismos de~$L$ en~$k'$,
\ifthenelse{\boolean{orig}}{como bien sabido}{como es bien sabido}.

\begin{proposition}
\label{V.8.2}
Sean~$S$ un preesquema conexo, localmente noetheriano y normal,
$K=\kres(s)$ su cuerpo de funciones $=$ el cuerpo residual en su punto
genérico~$s$, y~$\Omega$ una extensión algebraicamente cerrada de~$K$ que
define un punto geométrico~$a'$ de~$S'=\Spec(K)$ y un punto geométrico~$a$
de~$S$. Entonces el homomorfismo
$\pi_1(S';a')\to\pi_1(S;a)$ es sobreyectivo. Cuando se identifica el primer
grupo con el grupo de Galois de la clausura separable~$\bar{K}$ de~$K$
en~$\Omega$ (\cf~\Ref{V.8.1}), el núcleo del homomorfismo anterior
corresponde, por la teoría de Galois, a la subextensión de~$\bar{K}/K$ que
es el compuesto de las extensiones finitas de~$K$ en~$\Omega$ no ramificadas
sobre~$S$.
\end{proposition}

La primera afirmación significa que la imagen inversa sobre~$S'$ de un
recubrimiento étale conexo~$X$ de~$S$ es conexa, \ie que~$X$ es íntegro;
esto no es sino~\Ref{I.10.1}. El núcleo del homomorfismo anterior se
interpreta entonces como el conjunto de los automorfismos de~$\bar{K}/K$
que inducen la identidad sobre los conjuntos~$F_a(X)$,
\marginpar{144}%
donde puede suponerse conexo el recubrimiento étale~$X$ de~$S$. Pero esto
significa que tal automorfismo induce la identidad en las subextensiones
finitas de~$\bar{K}/K$ no ramificadas sobre~$S$, lo que demuestra la última
afirmación.

\begin{remarquestar}
Gracias a esta interpretación del grupo fundamental del preesquema normal~$S$
en términos de la
\ifthenelse{\boolean{orig}}{teoría de los Galois usual}{teoría de Galois usual},
la definición se conocía desde hacía mucho tiempo en este caso.
\end{remarquestar}

\section{Caso de los preesquemas no conexos: categorías multigaloisianas}
\label{V.9}

Sea~$S$ un preesquema localmente noetheriano, y sean
$(S_i)_{i\in I}$ sus componentes conexas. Entonces la categoría
$\cal{C}(S)$ de los recubrimientos étales de~$S$ es equivalente
al producto de las categorías~$\cal{C}(S_i)$, que se interpretan
en términos de los grupos fundamentales de las~$S_i$, una vez
\ifthenelse{\boolean{orig}}{elegido}{elegido} un punto geométrico en
cada~$S_i$. En la aplicación de la teoría del descenso a los morfismos
étales, resulta a veces incómodo elegir un punto geométrico de~$S_i$
para cada~$S_i$. Es entonces más cómodo recurrir a la generalización
natural de~\Ref{V.5.8} para interpretar~$\cal{C}(S)$ como una categoría
de funtores sobre el grupoide de los puntos geométricos de~$S$,
considerado como suma de los grupoides correspondientes a las componentes
conexas de~$S$; los funtores en cuestión son los funtores con valores en
la categoría de los conjuntos finitos que satisfacen la propiedad de
continuidad análoga a la invocada en~\Ref{V.5.8}. En la práctica, se tendrá
una familia~$(a_t)_{t\in E}$ de puntos geométricos de~$S$, tal que toda
componente conexa~$S_i$ de~$S$ contenga al menos uno, y podremos entonces,
como en~\Ref{V.5.10}, reemplazar el grupoide de todos los puntos geométricos
de~$S$ por el grupoide análogo cuyo conjunto subyacente es~$E$. Por supuesto,
estas consideraciones deberían insertarse en definiciones generales acerca
de las categorías equivalentes a productos de categorías de la
forma~$\cal{C}(\pi)$, que podrán llamarse
\emph{categorías multigaloisianas}.
\index{multigaloisiana (categoría)|hyperpage}%
Dejaremos los detalles al lector.

\chapter{Categorías fibradas y descenso}
\label{VI}
\marginpar{145}%

\setcounter{section}{-1}
\section{Introducción}

Contrariamente a lo anunciado en la introducción de la exposición
anterior, resultó imposible efectuar el descenso en la categoría de los
preesquemas, aun en casos particulares, sin haber desarrollado antes con
suficiente cuidado el lenguaje del descenso en las categorías generales.

La noción de «descenso» proporciona el marco general para todos los
procedimientos de «pegado» de objetos y, por consiguiente, de «pegado» de
categorías. El caso más clásico de pegado corresponde al dato de un espacio
topológico~$X$ y de un recubrimiento de~$X$ por abiertos~$X_i$; si para cada~$i$
se da un espacio fibrado (digamos)~$E_i$ sobre~$X_i$, y para cada par~$(i,j)$
un isomorfismo~$f_{ji}$ de~$E_i|X_{ij}$ sobre~$E_j|X_{ij}$ (donde se pone
$X_{ij}=X_i\cap X_j$), de modo que estos isomorfismos satisfacen una conocida
condición de transitividad
(escrita abreviadamente~$f_{kj}f_{ji}=f_{ki}$), se sabe que existe un espacio
fibrado~$E$ sobre~$X$, determinado salvo isomorfismo por la condición de que
existan isomorfismos~$f_i\colon E|X_i\isomto E_i$ que satisfagan las relaciones
$f_{ji}=f_jf_i^{-1}$ (con el abuso de notación habitual). Sea~$X'$ la suma
topológica de los~$X_i$, que es, por tanto, un espacio fibrado sobre~$X$ (es decir,
dotado de una aplicación continua~$X'\to X$). El dato de los~$E_i$ puede
interpretarse de manera más concisa como un espacio fibrado~$E'$ sobre~$X'$,
y el dato de los~$f_{ji}$ como un isomorfismo entre las dos imágenes inversas
(por las dos proyecciones canónicas)~$E''_1$ y~$E''_2$ de~$E'$ sobre
$X''=X'\times_X X'$; la condición de pegado puede escribirse entonces como
una identidad entre isomorfismos de los espacios fibrados~$E'''_1$ y~$E'''_3$
sobre el producto fibrado triple~$X'''=X'\times_X X'\times_X X'$ (donde
$E'''_i$ designa la imagen inversa de~$E'$ sobre~$X'''$ por la proyección
canónica de índice~$i$). La construcción de~$E$ a partir de~$E'$ y de~$f$ es
un caso típico
\marginpar{146}%
de procedimiento de «descenso». Por otra parte, partiendo de un espacio
fibrado~$E$ sobre~$X$, se dice que
\ifthenelse{\boolean{orig}}{$X$}{$E$} es «localmente trivial», de fibra~$F$,
si existe un recubrimiento abierto~$(X_i)$ de~$X$ tal que los~$E|X_i$ sean
isomorfos a~$F\times X_i$, o, lo que es equivalente, tal que la imagen inversa
$E'$ de~$E$ sobre~$X'=\coprod_i X_i$ sea isomorfa a~$X'\times F$.

Así, tanto la noción de «pegado» de objetos como la de «localización» de una
propiedad están ligadas al estudio de ciertos tipos de «cambios de base»
$X'\to X$. En geometría algebraica, muchos otros tipos de cambio de base, y
en particular los morfismos~$X'\to X$ fielmente planos, deben considerarse
como correspondientes a un procedimiento de «localización» relativo a los
preesquemas, u otros objetos, «sobre»~$X$. Este tipo de localización se utiliza
tanto como la simple localización topológica (que, por otra parte, es un caso
particular de aquel). Lo mismo sucede (en menor medida) en geometría analítica.

La mayoría de las demostraciones, que se reducen a verificaciones, se omiten
o solo se esbozan; cuando procede, precisamos los diagramas menos evidentes
que aparecen en una demostración.

\vfill
\section{Universos, categorías, equivalencia de categorías}
\label{VI.1}

Para evitar ciertas dificultades lógicas, admitiremos aquí la noción de
\emph{universo}, que es un conjunto «suficientemente grande» para no salir de
él mediante las operaciones habituales de la teoría de conjuntos; un
«\emph{axioma de los universos}» garantiza que todo objeto pertenece a algún
universo. Para más detalles, véase un libro en preparación de C\ptbl Chevalley
y el conferenciante\footnote{%
\ifthenelse{\boolean{orig}}
{Los autores definitivos son C\ptbl Chevalley y P\ptbl Gabriel. El libro debe
publicarse en el año 2000. \Cf también SGA~4 VI.7.1.}
{Los autores definitivos son C\ptbl Chevalley y P\ptbl Gabriel. El libro debe
publicarse en el año 3000. Mientras tanto, \cf también SGA~4 I.}}. Así, el
símbolo~$\Ens$
\label{indnot:fb}\oldindexnot{$\Ens$|hyperpage}%
designa, no la categoría de todos los conjuntos (noción que carece de sentido),
sino la categoría de los conjuntos que pertenecen a un universo dado (que aquí
no haremos constar en la notación). Del mismo modo, $\Cat$
\label{indnot:fc}\oldindexnot{$\Cat$|hyperpage}%
designará la categoría de las categorías pertenecientes al universo en
cuestión; los «morfismos» de un objeto~$X$ de~$\Cat$ en otro~$Y$ son, por
definición, los \emph{funtores} de~$X$ en~$Y$.

Si
\marginpar{147}
$\cal{C}$ es una categoría, denotamos por $\Ob(\cal{C})$
\label{indnot:fd}\oldindexnot{$\Ob(\cal{C})$|hyperpage}%
\emph{el conjunto de objetos} de~$\cal{C}$, y por~$\Fl(\cal{C})$
\label{indnot:fe}\oldindexnot{$\Fl(\cal{C})$|hyperpage}%
\emph{el conjunto de flechas} de~$\cal{C}$ (o morfismos de~$\cal{C}$).
Escribiremos, por tanto, $X\in\Ob(\cal{C})$, evitando el abuso de notación
habitual $X\in\cal{C}$. Si $\cal{C}$ y $\cal{C}'$ son dos categorías, un
\emph{funtor} de~$\cal{C}$ en~$\cal{C}'$ será siempre lo que comúnmente se
llama un funtor \emph{covariante} de~$\cal{C}$ en~$\cal{C}'$; su dato incluye
el de la
\ifthenelse{\boolean{orig}}{categoría de llegada y la categoría de partida}
{categoría de partida y la categoría de llegada}, $\cal{C}$ y $\cal{C}'$.
Los funtores de~$\cal{C}$ en~$\cal{C}'$ forman un conjunto, denotado
$\Hom(\cal{C},\cal{C}')$, que es el conjunto de objetos de una categoría
denotada $\SheafHom(\cal{C},\cal{C}')$.
\label{indnot:ff}\oldindexnot{$\SheafHom(\cal{C},\cal{C}')$|hyperpage}%
Por definición, un \emph{funtor contravariante} de~$\cal{C}$ en~$\cal{C}'$ es
un funtor de la \emph{categoría opuesta} $\cal{C}^\circ$
\label{indnot:fg}\oldindexnot{$\cal{C}^\circ$|hyperpage}%
de~$\cal{C}$ en~$\cal{C}'$.

Admitiremos las nociones de \emph{límite proyectivo} y de \emph{límite
inductivo} de un funtor $F\colon \cal{I}\to \cal{C}$ y, en particular, los
casos particulares más habituales de estas nociones: productos cartesianos y
productos fibrados, las nociones duales de sumas directas y sumas amalgamadas,
y las propiedades formales habituales de estas operaciones.

Por ejemplo, en la categoría $\Cat$ introducida más arriba, existen los límites
proyectivos (relativos a categorías~$\cal{I}$ pertenecientes al universo
elegido); el conjunto de objetos (\resp el conjunto de flechas) de la categoría
límite proyectivo~$\cal{C}$ de las~$\cal{C}_i$ se obtiene tomando el límite
proyectivo de los conjuntos de objetos (\resp de los conjuntos de flechas) de
las categorías~$\cal{C}_i$. El caso más conocido es el producto de una familia
de categorías. En lo sucesivo utilizaremos constantemente el producto fibrado
de dos categorías sobre una tercera.

Para todo lo relativo a categorías y funtores, mientras aparece el libro en
preparación ya mencionado, véase \cite{VI.1} (que es necesariamente muy
incompleto, incluso en lo que respecta a las generalidades esbozadas en el
presente número).

Aprovechemos esta ocasión para hacer explícita la noción de equivalencia de
categorías,
\index{equivalencia de categorías|hyperpage}%
que no se expone de manera satisfactoria en~\cite{VI.1}. Se dice que un funtor
$F\colon\cal{C}\to \cal{C}'$ es \emph{fiel}
\index{fiel (funtor)|hyperpage}%
(\resp \emph{plenamente fiel})
\index{plenamente fiel (funtor)|hyperpage}%
si, para todo par de objetos~$S$, $T$ de~$\cal{C}$, la aplicación
$u\mto F(u)$ de~$\Hom(S,T)$ en~$\Hom(F(S),F(T))$ es inyectiva
(\resp biyectiva). Se dice que $F$ es una \emph{equivalencia} de categorías si
\marginpar{148}
$F$ es plenamente fiel y, además, todo objeto~$S'$ de~$\cal{C}'$ es isomorfo a
un objeto de la forma~$F(S)$. Se demuestra que equivale a decir que existe un
funtor~$G$ de~$\cal{C}'$ en~$\cal{C}$ \emph{cuasi-inverso de}~$F$, es decir,
\ifthenelse{\boolean{orig}}
{tal que $GF$ sea isomorfo a $\id_{\cal{C}}$}
{tal que $GF$ sea isomorfo a $\id_{\cal{C}}$ y $FG$ a
$\id_{\cal{C}'}$}. Cuando así ocurre, el dato de un funtor
$G\colon\cal{C}'\to \cal{C}$ y de un isomorfismo
$\varphi\colon \ifthenelse{\boolean{orig}}{GF}{FG}\to \id_{\cal{C}'}$
equivale al dato, para todo $S'\in\Ob(\cal{C}')$, de un par~$(S,u)$ formado por
un objeto~$S$ de~$\cal{C}$ y un isomorfismo $u\colon F(S)\to S'$, a saber,
\ifthenelse{\boolean{orig}}{$(G(S),\varphi(S))$}
{$(G(S'),\varphi(S'))$}. (Con estas notaciones, existe un funtor único
$\cal{C}'\to\cal{C}$ cuya aplicación en los objetos es la aplicación dada
\ifthenelse{\boolean{orig}}{$S\mto G(S)$}{$S'\mto G(S')$}, y tal que la
aplicación \ifthenelse{\boolean{orig}}{$S\mto \varphi(S)$}
{$S'\mto \varphi(S')$} sea una transformación natural de funtores
$FG\to\id_{\cal{C}'}$). Por último, si $G$ es un funtor cuasi-inverso de~$F$
y se eligen isomorfismos $\varphi\colon FG\isomto\id_{\cal{C}'}$ y
$\psi\colon GF\isomto\id_{\cal{C}}$, entonces las dos condiciones de
\ifthenelse{\boolean{orig}}{compatibilidades}{compatibilidad} sobre~$\varphi$,
$\psi$ enunciadas en \cite[I.1.2]{VI.1} son, de hecho, equivalentes entre sí;
y, para todo isomorfismo~$\varphi$ elegido, existe un único
isomorfismo~$\psi$ que satisface dichas condiciones.

\section{Categorías sobre otra categoría}
\label{VI.2}

Sea $\cal{E}$ una categoría en~$\Univ$; es, por tanto, un objeto de~$\Cat$, y
se puede considerar la categoría
$\Cat_{/\cal{E}}$
\label{indnot:fh}\oldindexnot{$\Cat_{/\cal{E}}$|hyperpage}%
de los «objetos de~$\Cat$ sobre~$\cal{E}$». Un objeto de esta categoría es,
pues, un funtor
$$
p\colon \cal{F}\to\cal{E}
$$
También se dice que la categoría~$\cal{F}$, provista de tal funtor, es una
\emph{categoría sobre}~$\cal{E}$, o una $\cal{E}$\emph{-categoría}. Se llamará,
por tanto, $\cal{E}$-funtor de una categoría~$\cal{F}$ sobre~$\cal{E}$ en una
categoría~$\cal{G}$ sobre~$\cal{E}$ a un funtor
$$
f\colon\cal{F}\to\cal{G}
$$
tal que
$$
qf=p
$$
donde $p$ y $q$ son los funtores de proyección de~$\cal{F}$ y~$\cal{G}$,
respectivamente. El conjunto de los $\cal{E}$-funtores~$f$ de~$\cal{F}$ en
$\cal{G}$ está, por tanto, en correspondencia biunívoca con el conjunto de las
flechas de origen~$\cal{F}$ y extremo~$\cal{G}$ en~$\Cat_{/\cal{E}}$, sin
\marginpar{149}
que se trate por ello de una identidad (pues el dato de un~$f$ como el anterior
no determina $\cal{F}$ y $\cal{G}$ como categorías sobre~$\cal{E}$); pero, por
supuesto, como en cualquier otra categoría~$\cal{C}_{/S}$, se cometerá
habitualmente el abuso de lenguaje consistente en identificar los
$\cal{E}$-funtores (en el sentido precisado más arriba) con flechas de una
categoría~$\Cat_{/\cal{E}}$.

Denotaremos por
$$
\Hom_{\cal{E}}(\cal{F},\cal{G})
$$
el conjunto de los $\cal{E}$-funtores de~$\cal{F}$ en~$\cal{G}$. Por supuesto,
un compuesto de $\cal{E}$-funtores es un $\cal{E}$-funtor (la composición en
cuestión corresponde, por definición, a la composición de las flechas
de~$\Cat_{/\cal{E}}$).

Consideremos ahora dos $\cal{E}$-funtores
$$
f,g\colon \cal{F}\to\cal{G}
$$
y una transformación natural de funtores:
$$
u\colon f\to g
$$
Se dice que $u$ es un $\cal{E}$-\emph{homomorfismo}, o un
«\emph{homomorfismo de $\cal{E}$-funtores}», si para todo
$\xi\in\Ob(\cal{F})$ se tiene
$$
q(u(\xi))=\id_{p(\xi)} \quoi\text{,}
$$
en palabras: poniendo
\ifthenelse{\boolean{orig}}
{$S=p(\xi)=qf(\xi)=qg(\xi)\in\Ob\cal{E}$,}
{$S=p(\xi)=qf(\xi)=qg(\xi)\in\Ob(\cal{E})$,}
el morfismo
$$
u(\xi)\colon f(\xi)\to g(\xi)
$$
de~$\cal{G}$ es un $\id_S$-morfismo. (En general, para todo morfismo
$\alpha\colon T \to S$ de~$\cal{E}$ y toda categoría~$\cal{G}$ sobre~$\cal{E}$,
un morfismo~$v$ de~$\cal{G}$ se llama $\alpha$-\emph{morfismo} si
$q(v)=\alpha$, donde $q$ denota el funtor de proyección
$\cal{G}\to \cal{E}$). Si se tiene un tercer $\cal{E}$-funtor
$h\colon \cal{F}\to \cal{G}$ y un $\cal{E}$-homomorfismo $v\colon g\to h$,
entonces $vu$ es también un $\cal{E}$-homomorfismo. Así,
\marginpar{150}
\emph{los $\cal{E}$-funtores de~$\cal{F}$ en~$\cal{G}$ y los
$\cal{E}$-homomorfismos entre ellos forman una subcategoría de la categoría
$\SheafHom(\cal{F},\cal{G})$ de todos los funtores de~$\cal{F}$ en~$\cal{G}$,
que se llamará categoría de los $\cal{E}$-funtores de~$\cal{F}$ en~$\cal{G}$ y
se denotará por}
$$
\SheafHom_{\cal{E}/-}(\cal{F},\cal{G})
$$
\label{indnot:fi}\oldindexnot{$\SheafHom_{\cal{E}/-}(\cal{F},\cal{G})$|hyperpage}%
Es también la subcategoría núcleo del par de funtores
\begingroup
\shorthandoff{<>}% Spanish babel otherwise treats Xy-pic's bend delimiters as shorthands.
$$
\xymatrix@C=.5cm{R,S\colon\SheafHom(\cal{F},\cal{G})
\ar@<2pt>[r]\ar@<-2pt>[r]&\SheafHom(\cal{F},\cal{E})}\quoi\text{,}
$$
\endgroup
donde $R$ es el funtor constante definido por el objeto~$p$ de
$\SheafHom(\cal{F},\cal{E})$, y $S$ es el funtor $f\mto q\circ f$ definido por
$q\colon\cal{G}\to\cal{E}$.

Para terminar estas generalidades, queda definir las operaciones naturales de
composición entre las categorías $\SheafHom_{\cal{E}/-}(\cal{F},\cal{G})$
mediante la composición de $\cal{E}$-funtores. En otras palabras, se quiere
definir un
«funtor de composición»:
\begin{equation*}
\label{eq:VI.2.i}
\tag{i} {\SheafHom_{\cal{E}/-}(\cal{F},\cal{G})\times
\SheafHom_{\cal{E}/-}(\cal{G},\cal{H})\to
\SheafHom_{\cal{E}/-}(\cal{F},\cal{H})}
\end{equation*}
cuando $\cal{F}$, $\cal{G}$ y $\cal{H}$ son tres categorías sobre~$\cal{E}$,
de modo que este funtor induzca, en los objetos, la aplicación de composición
$(f,g)\mto gf$ de los $\cal{E}$-funtores
$f\colon\cal{F}\to\cal{G}$ y $g\colon \cal{G}\to\cal{H}$. Para ello,
recordemos que se define un funtor canónico:
\begin{equation*}
\label{eq:VI.2.ii}
\tag{ii}
{\SheafHom(\cal{F},\cal{G})\times\SheafHom(\cal{G},\cal{H})\to
\SheafHom(\cal{F},\cal{H})}
\end{equation*}
que, en los objetos, no es sino la aplicación de composición $(f,g)\mto gf$ de
funtores y que transforma una flecha~$(u,v)$, donde
$$
u\colon f\to f'\quoi,\quad v\colon g\to g'
$$
son flechas de~$\SheafHom(\cal{F},\cal{G})$ y de
$\SheafHom(\cal{G},\cal{H})$, respectivamente, en la flecha
$$
v*u\colon gf\to g'f'
$$
\label{indnot:fj}\oldindexnot{$v*u$|hyperpage}%
definida por la relación:
\marginpar{151}
$$
v*u(\xi)=v(f'(\xi))\cdot g(u(\xi))=g'(u(\xi))\cdot v(f(\xi))
$$
Es bien sabido que así se obtiene efectivamente una transformación natural
de~$gf$ en~$g'f'$ y que (para $f,g$ y $u,v$ variables) se obtiene de este modo
un funtor~\eqref{eq:VI.2.ii}, es decir, que se tiene
\begin{equation*}
\label{eq:VI.2.I}
\tag{I} {\id_g*\id_f=\id_{gf}\quoi,}
\end{equation*}
\begin{equation*}
\label{eq:VI.2.II}
\tag{II} {(v'*u')\circ(v*u)=(v'\circ v)*(u'\circ u)}
\end{equation*}
Recordemos también que existe una fórmula de asociatividad para las
composiciones canónicas~\eqref{eq:VI.2.ii}, expresada, por una parte, por la
asociatividad $(hg)f=h(gf)$ de la composición de funtores y, por otra, por la
fórmula
\begin{equation*}
\label{eq:VI.2.III}
\tag{III} {(w*v)*u=w*(v*u)}
\end{equation*}
para las composiciones horizontales de transformaciones naturales (donde
$u\colon f\to f'$ y $v\colon g\to g'$ son como antes, y se supone dada además
una transformación natural $w\colon h\to h'$ de los funtores
$h,h'\colon \cal{H}\to\cal{K}$). Afirmamos ahora que
\emph{cuando~$\cal{F}$ y $\cal{G}$ son $\cal{E}$-categorías, el funtor de
composición canónico}~\eqref{eq:VI.2.ii}
\emph{induce un funtor}~\eqref{eq:VI.2.i}. Como ya se sabe que el compuesto de
dos $\cal{E}$-funtores es un $\cal{E}$-funtor, esto equivale a decir que
\emph{cuando $u\colon f\to f'$ y $v\colon g\to g'$ son homomorfismos de
$\cal{E}$-funtores, entonces $v*u\colon gf\to g'f'$ es también un homomorfismo
de $\cal{E}$-funtores.} Esto resulta, en efecto, inmediatamente de las
definiciones. Como las operaciones de composición~\eqref{eq:VI.2.i} son
inducidas por las operaciones de composición~\eqref{eq:VI.2.ii}, satisfacen la
misma propiedad de
asociatividad, expresada también por las fórmulas $(hg)f=h(gf)$ y
$(w*v)*u=w*(v*u)$ para $\cal{E}$-funtores y $\cal{E}$-homomorfismos de
$\cal{E}$-funtores.

Para completar el formulario~\eqref{eq:VI.2.I}, \eqref{eq:VI.2.II},
\eqref{eq:VI.2.III}, recordemos también las fórmulas:
\begin{equation*}
\label{eq:VI.2.IV}
\tag{IV} {v*\id_{\cal{F}}=v\quad\text{y}\quad\id_{\cal{G}}*u=u}\quoi,
\end{equation*}
donde
\marginpar{152}
para abreviar se escribe $v*f$ o
\ifthenelse{\boolean{orig}}{$u*g$}{$g*u$} en lugar de~$v*u$, cuando $u$,
respectivamente $v$, es el automorfismo identidad de~$f$, respectivamente
de~$g$.

De la definición de las operaciones de composición~\eqref{eq:VI.2.i} se sigue
que
$\SheafHom_{\cal{E}/-}(\cal{F},\cal{G})$ \emph{es un funtor en
\ifthenelse{\boolean{orig}}
{$\cal{E}$,~$\cal{G}$,}
{$\cal{F}$,~$\cal{G}$,}
de la categoría producto ${\Cat_{/\cal{E}}}^\circ\times\Cat_{/\cal{E}}$ en la
categoría~$\Cat$}. En efecto, si $g\colon \cal{G}\to\cal{G}_1$ es un
$\cal{E}$-funtor, es decir, un objeto de
$\SheafHom_{\cal{E}/-}(\cal{G},\cal{G}_1)$, entonces, haciendo
$\cal{H}=\cal{G}_1$ en~\eqref{eq:VI.2.i}, le corresponde un funtor
$$
g_*\colon\SheafHom_{\cal{E}/-}(\cal{F},\cal{G})\to
\SheafHom_{\cal{E}/-}(\cal{F},\cal{G}_1)
$$
De manera análoga, para un $\cal{E}$-funtor
$f\colon \cal{F}_1\to \cal{F}$ se define un funtor
$$
f^*:\SheafHom_{\cal{E}/-}(\cal{F},\cal{G})\to
\SheafHom_{\cal{E}/-}(\cal{F}_1,\cal{G})
$$
Para abreviar, estos funtores se denotan también por los símbolos
$f\mto g\circ f$ y $g\mto g\circ f$, respectivamente (que, en realidad,
designan únicamente las aplicaciones correspondientes en los conjuntos de
objetos). De la propiedad de asociatividad indicada más arriba se sigue que,
de este modo, se obtiene efectivamente el funtor anunciado
$\Cat_{/\cal{E}}^\circ\times\Cat_{/\cal{E}}\to\Cat$.

\section{Cambio de base en las categorías sobre~$\mathcal{E}$}
\label{VI.3}

Como en~$\Cat$ existen los límites proyectivos (respecto de categorías~$\cal{I}$
que son elementos de~$\Univ$), lo mismo ocurre en~$\Cat_{/\cal{E}}$; en
particular, existen allí los productos cartesianos, que se interpretan como
productos fibrados en~$\Cat$. De acuerdo con las notaciones generales, si
$\cal{F}$ y $\cal{G}$ son categorías sobre~$\cal{E}$, se denota por
$$
\cal{F}\times_{\cal{E}}\cal{G}
$$
\label{indnot:fk}\oldindexnot{$\cal{F}\times_{\cal{E}}\cal{G}$|hyperpage}%
su producto en~$\Cat_{/\cal{E}}$, es decir, su producto fibrado sobre~$\cal{E}$
en~$\Cat$,
\marginpar{153}
considerado como categoría sobre~$\cal{E}$. Así,
$\cal{F}\times_{\cal{E}}\cal{G}$ está provisto de dos $\cal{E}$-funtores
$\pr_1$ y $\pr_2$, que definen, para toda categoría $\cal{H}$ sobre~$\cal{E}$,
una biyección
$$
\Hom_{\cal{E}}(\cal{H},\cal{F}\times_{\cal{E}}\cal{G})\isomto
\Hom_{\cal{E}}(\cal{H},\cal{F})\times\Hom_{\cal{E}}(\cal{H},\cal{G}).
$$
Esta biyección procede, además, de un isomorfismo de categorías
$$
\SheafHom_{\cal{E}/-}(\cal{H},\cal{F}\times_{\cal{E}}
\cal{G})\isomto\SheafHom_{\cal{E}/-}(\cal{H},\cal{F})\times
\SheafHom_{\cal{E}/-}(\cal{H},\cal{G})
$$
al tomar los conjuntos de objetos de ambos miembros, donde el funtor escrito es
el que tiene por componentes los funtores $h\mto\pr_1\circ h$ y
$h\mto\pr_2\circ h$ del primer miembro en los dos factores del segundo. Se
deja al lector comprobar que se obtiene efectivamente un isomorfismo (pues el
hecho análogo es válido, más generalmente, siempre que se tenga un límite
proyectivo de categorías, y no solo en el caso de un producto fibrado).

Recordemos, además, que (como se dijo en el~\No \Ref{VI.1}):
\begin{eqnarray*}
\Ob(\cal{F}\times_\cal{E}\cal{G})&=&
\Ob(\cal{F})\times_{\Ob(\cal{E})}\Ob(\cal{G})\\
\Fl(\cal{F}\times_\cal{E}\cal{G})&=&
\Fl(\cal{F})\times_{\Fl(\cal{E})}\Fl(\cal{G})\quoi,
\end{eqnarray*}
y la composición de las flechas se realiza componente a componente.

En lo que sigue, consideraremos un funtor
$$
\lambda\colon\cal{E}'\to\cal{E}
$$
y, para toda categoría $\cal{F}$ sobre~$\cal{E}$, consideraremos
$\cal{F}\times_\cal{E}\cal{E}'$ como categoría sobre~$\cal{E}'$ por medio
de~$\pr_2$; en otras palabras, interpretaremos la operación «producto fibrado»
como una operación de \emph{«cambio de base»}, y el funtor
$\lambda\colon\cal{E}'\to\cal{E}$ recibe el nombre de \emph{«funtor de cambio
de base»}.
\index{cambio de base (funtor)|hyperpage}%
\index{funtor de cambio de base|hyperpage}%
De acuerdo con los hechos generales bien conocidos, se obtiene así un funtor,
llamado
\marginpar{154}
\emph{funtor de cambio de base} para~$\lambda$:
$$
\lambda^*\colon\Cat_{/\cal{E}}\to\Cat_{/\cal{E}'}\quoi,
$$
\label{indnot:fl}\oldindexnot{$\lambda^*\colon\Cat_{/\cal{E}}\to\Cat_{/\cal{E}'}$|hyperpage}%
(adjunto al funtor de «restricción de la base» que, a toda categoría~$\cal{F}'$
sobre~$\cal{E}'$ con funtor de proyección~$p'$, asocia~$\cal{F}'$, considerada
como categoría sobre~$\cal{E}$ mediante el funtor~$p=\lambda p'$). Como es bien
sabido en el caso general de un funtor de cambio de base en una categoría, el
funtor de cambio de base «conmuta con los límites proyectivos» y, en particular,
«transforma» productos fibrados sobre~$\cal{E}$ en productos fibrados
sobre~$\cal{E}'$.

Sean $\cal{F}$ y $\cal{G}$ dos categorías sobre~$\cal{E}$; vamos a definir un
\emph{isomorfismo canónico}:
\begin{multline*}
\label{eq:VI.3.i}
\tag{i}
\SheafHom_{\cal{E}'/-}(\cal{F}',\cal{G}')\isomto
\SheafHom_{\cal{E}/-}(\cal{F}\times_{\cal{E}}\cal{E}',\cal{G})\\
\text{donde }\cal{F}'=\cal{F}\times_{\cal{E}}\cal{E}'\text{, }
\cal{G}'=\cal{G}\times_{\cal{E}}\cal{E}'.
\end{multline*}
Para ello, consideremos el funtor
$$
\pr_1\colon\cal{G}'=\cal{G}\times_{\cal{E}}\cal{E}'\to\cal{G}\quoi,
$$
y definamos~\eqref{eq:VI.3.i} por
$$
F\mto\pr_1\circ F\quoi,
$$
que, a priori, designa un funtor
\begin{equation*}
\label{eq:VI.3.ii}
\tag{ii} {\SheafHom(\cal{F}',\cal{G}')\to\SheafHom(\cal{F}',\cal{G})}
\end{equation*}
Solo hay que comprobar, por tanto, que este último induce un funtor en las
subcategorías~\eqref{eq:VI.3.i} y que dicho funtor es un isomorfismo. Que
\eqref{eq:VI.3.ii} induce una biyección
$$
\Hom_{\cal{E}'/-}(\cal{F}',\cal{G}')\isomto
\Hom_{\cal{E}/-}(\cal{F}\times_{\cal{E}}\cal{E}',\cal{G})
$$
es la propiedad característica del funtor de cambio de base. Queda, pues,
\marginpar{155}
probar que, si $F$ y $G$ son $\cal{E}'$-funtores
$\cal{F}'\to\cal{G}'$, entonces \emph{la aplicación}
$$
u\mto\pr_1\circ u
$$
\emph{induce una biyección}
$$
\Hom_{\cal{E}'}(F,G)\isomto\Hom_{\cal{E}}(\pr_1\circ F,\pr_1\circ
G)\quoi.
$$
La comprobación de este hecho es inmediata y se deja al lector.

De este isomorfismo~\eqref{eq:VI.3.i} y del final del número anterior se deduce
que
$$
\SheafHom_{\cal{E}'/-}(\cal{F}\times_{\cal{E}}\cal{E}',\cal{G}
\times_{\cal{E}}\cal{E}')
$$
\emph{puede considerarse como un funtor en}
$\cal{E}',\cal{F},\cal{G}$, \emph{de la categoría}
$\Cat_{/\cal{E}}^\circ\times\Cat_{/\cal{E}}^\circ\allowbreak\times\Cat_{/\cal{E}}$
\emph{en la categoría}~$\Cat$, isomorfo al funtor definido por la expresión
$\SheafHom_{\cal{E}/-}(\cal{F}\times_{\cal{E}}\cal{E}',\cal{G})$. En
particular, para $\cal{F}$ y~$\cal{G}$ fijos, se obtiene un funtor en~$\cal{E}'$,
\ifthenelse{\boolean{orig}}
{$\cal{E}'\to\SheafHom_{\cal{E}''/-}(\cal{F}',\cal{G}')=
\SheafHom_{\cal{E}'/-}(\cal{F}\times_{\cal{E}}
\cal{E}',\cal{G}\times_{\cal{E}}\cal{E}')$}
{$\cal{E}'\mto\SheafHom_{\cal{E}'/-}(\cal{F}',\cal{G}')=
\SheafHom_{\cal{E}'/-}(\cal{F}\times_{\cal{E}}
\cal{E}',\cal{G}\times_{\cal{E}}\cal{E}')$}, y, en particular, el
$\cal{E}$-funtor de proyección $\lambda\colon\cal{E}'\to\cal{E}$ define un
morfismo, es decir, un funtor
$$
\lambda^*_{\cal{F},\cal{G}}\colon
\SheafHom_{\cal{E}/-}(\cal{F},\cal{G})\to
\SheafHom_{\cal{E}'/-}(\cal{F}',\cal{G}')
$$
que vamos a explicitar. En los conjuntos de objetos de ambos miembros, se trata
de la aplicación
$$
f\mto f'=f\times_{\cal{E}}\cal{E}'
$$
que expresa la dependencia funtorial de
$\cal{F}\times_{\cal{E}}\cal{E}'$ respecto del objeto~$\cal{F}$
sobre~$\cal{E}$. Por otra parte, consideremos dos $\cal{E}$-funtores
$$
f,g\colon\cal{F}\to\cal{G}
$$
y un homomorfismo de $\cal{E}$-funtores
$$
u\colon f\to g\quoi,
$$
vamos
\marginpar{156}
a explicitar el homomorfismo correspondiente de $\cal{E}'$-funtores:
$$
u'\colon f'\to g'\quoi.
$$
Para todo
$$
\xi'=(\xi,S')\in\Ob(\cal{F}')
$$
con
$$
\xi\in\Ob(\cal{F}),\quad S'\in\Ob(\cal{E}'),\quad p(\xi)=\lambda(S')=S
$$
el morfismo
$$
u'(\xi')\colon f'(\xi')=(f(\xi),S')\to g'(\xi')=(g(\xi),S')\quad
\text{en }\cal{G}'
$$
se define mediante la fórmula
$$
u'(\xi')=(u(\xi),\id_{S'})
$$
(lo cual es, en efecto, un $S'$-morfismo en~$\cal{G}'$, pues
$q(u(\xi))=\lambda(\id_{S'})=\id_S$).

Consideremos ahora un $\cal{E}$-funtor cualquiera
$$
\lambda'\colon\cal{E}''\to\cal{E}'
$$
y el funtor correspondiente
$$
\SheafHom_{\cal{E}'/-}(\cal{F}\times_{\cal{E}}\cal{E}',
\cal{G}\times_{\cal{E}}\cal{E}')\to
\SheafHom_{\cal{E}''/-}(\cal{F}\times_{\cal{E}}\cal{E}'',
\cal{G}\times_{\cal{E}}\cal{E}'')\quoi,
$$
Afirmamos que este funtor no es sino el obtenido por el procedimiento anterior,
partiendo de~$\cal{F}'$ y $\cal{G}'$ sobre~$\cal{E}'$ y considerando
$\cal{E}''$ como una categoría sobre~$\cal{E}'$, habida cuenta de los
isomorfismos de \emph{«transitividad del cambio de base»}:
$$
\ifthenelse{\boolean{orig}}
{\cal{F}'\times_{\cal{E}'}\cal{E}''\isomto \cal{F}''=
\cal{F}\times_{\cal{E}}\cal{E}''\quad
\text{y}\quad\cal{G}'\times_{\cal{E}}\cal{E}''\isomto\cal{G}''
=\cal{G}\times_{\cal{E}}\cal{E}''\quoi,}
{\cal{F}'\times_{\cal{E}'}\cal{E}''\isomto
\cal{F}''=\cal{F}\times_{\cal{E}}\cal{E}''\quad
\text{y}\quad\cal{G}'\times_{\cal{E}'}\cal{E}''\isomto\cal{G}''
=\cal G\times_{\cal E}\cal E''\quoi,}
$$
que implican un isomorfismo canónico
$$
\SheafHom_{\cal{E}''/-}(\cal{F}'\times_{\cal{E}'}\cal{E}'',
\cal{G}'\times_{\cal{E}'}\cal{E}'')\isomto
\SheafHom_{\cal{E}''/-}(\cal{F}\times_{\cal{E}}\cal{E}'',
\cal{G}\times_{\cal{E}}\cal{E}'')
$$
La
\marginpar{157}
comprobación de esta compatibilidad es inmediata y se deja al lector.

Los funtores que se acaban de definir son compatibles con las operaciones de
composición definidas en el número anterior, en el sentido preciso siguiente:
si $\cal{F}$, $\cal{G}$ y~$\cal{H}$ son categorías sobre~$\cal{E}$ y se pone
$$
\cal{F}'=\cal{F}\times_{\cal{E}}\cal{E}'\quoi,\quad \cal{G}'=
\cal{G}\times_{\cal{E}}\cal{E}'\quoi,\quad \cal{H}'=
\cal{H}\times_{\cal{E}}\cal{E}'\quoi,
$$
el siguiente diagrama de funtores es conmutativo:
\ifthenelse{\boolean{orig}}
{\def\VIthreebase{\cal{E}}}
{\def\VIthreebase{\cal{E}'}}
$$
\xymatrix{ \SheafHom_{\cal{E}/-}(\cal{F},\cal{G})\times
\SheafHom_{\cal{E}/-}(\cal{G},\cal{H}) \ar[r]
\ar[d]_{\lambda^*_{\cal{F},\cal{G}}\times \lambda^*_{\cal{G},\cal{H}}}
& \SheafHom_{\cal{E}/-}(\cal{F},\cal{H})
\ar[d]^{\lambda^*_{\cal{F},\cal{H}}} \\
\SheafHom_{\VIthreebase/-}(\cal{F}',\cal{G}')\times
\SheafHom_{\cal{E}'/-}(\cal{G}',\cal{H}')\ar[r] &
\SheafHom_{\cal{E}'/-}(\cal{F}',\cal{H}')}
$$
donde las flechas horizontales son los funtores de composición definidos en el
número anterior. Esta conmutatividad se expresa mediante las fórmulas
$$
(gf)'=g'f'
$$
para $f\in\Hom_{\cal{E}}(\cal{F},\cal{G})$,
$g\in\Hom_{\cal{E}}(\cal{G},\cal{H})$ (fórmula que expresa simplemente la
funtorialidad del cambio de base), y mediante la fórmula
$$
(v*u)'=v'*u'
$$
cuando $u\colon f\to f_1$ es una flecha de
$\SheafHom_{\cal{E}/-}(\cal{F},\cal{G})$, mientras que
$v\colon g\to g_1$ es una flecha
de~$\SheafHom_{\cal{E}/-}(\cal{G},\cal{H})$. La comprobación de esta fórmula se
deduce fácilmente de las definiciones.

En lo que sigue, nos interesaremos sobre todo por
$\SheafHom_{\cal{E}}(\cal{F},\cal{G})$ (y por ciertas subcategorías notables de
esta) cuando $\cal{F}=\cal{E}$, e introducimos
\ifthenelse{\boolean{orig}}{por estas razones}{por esta razón} una notación
especial:
$$
\bf{\Gamma}(\cal{G}/\cal{E})=
\SheafHom_{\cal{E}}(\cal{E},\cal{G})\quoi,\quad
\Gamma(\cal{G}/\cal{E})=\Ob(\bf{\Gamma}(\cal{G}/\cal{E}))=
\ifthenelse{\boolean{orig}}
{\Hom_{\cal{E}}(\cal{G},\cal{E})\quoi.}
{\Hom_{\cal{E}}(\cal{E},\cal{G})\quoi.}
$$
\label{indnot:fm}\oldindexnot{$\bf{\Gamma}(\cal{G}/\cal{E})$ y $\Gamma(\cal{G}/\cal{E})$|hyperpage}%

\begin{remarquesstar}
Cuando
\marginpar{158}
$\cal{E}$ es una categoría puntual, es decir, cuando $\Ob(\cal{E})$ y
$\Fl(\cal{E})$ se reducen cada uno a un solo elemento, lo que significa también
que $\cal{E}$ es un objeto final de la categoría~$\Cat$, el dato de una
categoría sobre~$\cal{E}$ equivale al dato de una categoría sin más (pues hay
un único funtor de~$\cal{F}$ en~$\cal{E}$). Más precisamente,
$\Cat_{/\cal{E}}$ es entonces isomorfa a~$\Cat$. Además, las categorías
$\SheafHom_{\cal{E}/-}(\cal{F},\cal{G})$ no son sino las categorías
$\SheafHom(\cal{F},\cal{G})$. Recordemos entonces que la fórmula fundamental
$$
\Hom(\cal{H},\SheafHom(\cal{F},\cal{G}))\isomto
\Hom(\cal{F}\times\cal{H},\cal{G})
$$
(isomorfismo funtorial en los tres argumentos que figuran en ella) permite
interpretar $\SheafHom(\cal{F},\cal{G})$ axiomáticamente, en términos internos
a la categoría~$\Cat$, de manera que el formalismo conocido de las categorías
$\SheafHom$ aparece como un caso particular de un formalismo válido en
categorías tales como~$\Cat$, en las que existen «objetos $\SheafHom$ internos»
(definidos por la fórmula anterior). Hay una interpretación análoga de
$\SheafHom_{\cal{E}/-}(\cal{F},\cal{G})$ cuando se supone de nuevo que
$\cal{E}$ es arbitraria, mediante la fórmula
$$
\Hom(\cal{H},\SheafHom_{\cal{E}/-}(\cal{F},\cal{G}))\isomto
\Hom_{\cal{E}}(\cal{F}\times\cal{H},\cal{G})
$$
(isomorfismo funtorial en los tres argumentos). De este modo, las propiedades
formales expuestas en los \No \Ref{VI.2},~\Ref{VI.3} son casos particulares de
resultados más generales, válidos en las categorías en las que existen los
objetos internos $\SheafHom_{\cal{E}/-}(\cal{F},\cal{G})$ (cuando $\cal{F}$ y~$\cal{G}$
son dos objetos de la categoría de objetos sobre un tercero~$\cal{E}$).
\end{remarquesstar}

\section[Categorías fibra]{Categorías fibra; equivalencia de
$\mathcal{E}$-categorías}
\label{VI.4}

Sea~$\cal{F}$ una categoría sobre~$\cal{E}$, y sea $S\in
\Ob({\cal{E}})$. Se llama \emph{categoría fibra de}~$\cal{F}$
en~$S$
\index{categoría fibra|hyperpage}%
a la subcategoría~$\cal{F}_S$
\label{indnot:fn}\oldindexnot{$\cal{F}_S$|hyperpage}%
de~$\cal{F}$ que es imagen inversa de la
\marginpar{159}
subcategoría puntual de~$\cal{E}$ definida por~$S$. Por tanto, los objetos
de~$\cal{F}_S$ son los objetos~$\xi$ de~$\cal{F}$ tales que $p(\xi)=S$, y
sus morfismos son los morfismos~$u$ de~$\cal{F}$ tales que
$p(u)=\id_S$, es decir, los $S$-morfismos en~$\cal{F}$. Naturalmente,
$\cal{F}_S$ es canónicamente isomorfa al producto fibrado
$\cal{F}\times_{\cal{E}}\{S\}$, donde~$\{S\}$ designa la subcategoría
puntual de~$\cal{E}$ definida por~$S$, provista de su funtor de inclusión
en~$\cal{E}$. Se deduce de ello (habida cuenta de la transitividad del cambio
de base) que, si se efectúa el cambio de base
$\lambda\colon\cal{E}'\to\cal{E}$, entonces, para todo
$S'\in\Ob(\cal{E}')$, \emph{la proyección $\pr_1\colon
\cal{F}'=\cal{F}\times_{\cal{E}}\cal{E}'\to\cal{F}$ induce un
isomorfismo}
$$
\cal{F}'_{S'}\to\cal{F}_S\qquad(\text{donde $S=\lambda(S')$})%
\ifthenelse{\boolean{orig}}{}{.}
$$

\begin{proposition}
\label{VI.4.1} Sea $f\colon\cal{F}\to\cal{G}$ un
$\cal{E}$-funtor. Si~$f$ es plenamente fiel, entonces, para todo cambio de
base $\cal{E}'\to\cal{E}$, el funtor correspondiente
$f'\colon\cal{F}'=\cal{F}\times_{\cal{E}}\cal{E}'\to\allowbreak
\cal{G}'=\cal{G}\times_{\cal{E}}\cal{E}'$ es plenamente fiel.
\end{proposition}

La comprobación es inmediata; más generalmente, puede demostrarse que todo
límite proyectivo de
\ifthenelse{\boolean{orig}}{funtor}{funtores}
plenamente
\ifthenelse{\boolean{orig}}{fiel}{fieles}
(aquí,~$f$ y los funtores identidad en~$\cal{E},\cal{E}'$) es un funtor
plenamente fiel.

Obsérvese que la afirmación análoga a~4.1,
\ifthenelse{\boolean{orig}}{o}{donde}
«plenamente fiel» se sustituye por «equivalencia de categorías», es falsa ya
para~$\cal{G}=\cal{E}$. Sin embargo:

\begin{proposition}
\label{VI.4.2} Sea $f\colon\cal{F}\to\cal{G}$ un
$\cal{E}$-funtor. Las condiciones siguientes son equivalentes:
\begin{enumerate}
\item[(i)] Existe un $\cal{E}$-funtor $g\colon\cal{G}\to\cal{F}$ y
$\cal{E}$-isomorfismos
$$
gf\isomto\id_{\cal{F}},\quad fg\isomto\id_{\cal{G}}
\ifthenelse{\boolean{orig}}{}{.}
$$
\item[(ii)] Para toda categoría $\cal{E}'$ sobre~$\cal{E}$, el funtor
$$
f'=f\times_{\cal{E}}\cal{E}'\colon\cal{F}'=\cal{F}
\times_{\cal{E}}\cal{E}'\to\cal{G}'=\cal{G}\times_{\cal{E}}\cal{E}'
$$
es una equivalencia de categorías.
\item[(iii)]
\marginpar{160}
$f$ es una equivalencia de categorías y, para todo $S\in\Ob(\cal{E})$, el
funtor $f_S\colon\cal{F}_S\to\cal{G}_S$ inducido por~$f$ es una equivalencia
de categorías.
\item[(iii~bis)] $f$ es plenamente fiel y, para todo $S\in\Ob(\cal{E})$ y
todo
\ifthenelse{\boolean{orig}}
{$\eta\in\Ob\cal{G}_S$,}
{$\eta\in\Ob(\cal{G}_S)$,}
existen un $\xi\in\Ob(\cal{F}_S)$ y un $S$-isomorfismo
\ifthenelse{\boolean{orig}}
{$u\colon f((\xi)\to\eta$.}
{$u\colon f(\xi)\to\eta$.}
\end{enumerate}
\end{proposition}

\subsubsection*{Demostración} Evidentemente, (i) implica que~$f$ es una
equivalencia de categorías (noción que se define mediante la misma condición,
pero sin exigir que los isomorfismos de funtores sean $\cal{E}$-morfismos).
Por otra parte, de las funtorialidades del número anterior se deduce que la
condición~(i) se conserva después del cambio de base
$\cal{E}'\to\cal{E}$. Por consiguiente, (i)$\To$(ii). Evidentemente,
(ii)$\To$(iii), pues basta tomar $\cal{E}'=\cal{E}$ y
$\cal{E}'=\{S\}$. Es aún más trivial que (iii)$\To$(iii~bis); queda por
probar que (iii~bis)$\To$(i). Para ello, elijamos, para todo
$\eta\in\Ob(\cal{G})$, un $g(\eta)\in\Ob(\cal{F})$ y un isomorfismo
$u(\eta)\colon f(g(\eta))\to\eta$ tal que $q(u(\eta))=\id_S$, donde
$S=q(\eta)$. Esto es posible gracias a la segunda condición de (iii~bis).
Como es bien sabido e inmediato, el hecho de que~$f$ sea plenamente fiel
implica que~$g$ puede considerarse, de manera única, como un funtor de
$\cal{G}$ en~$\cal{F}$, de modo que los~$u(\eta)$ definan una transformación
natural (y, por tanto, un isomorfismo)
$u\colon fg\isomto\id_{\cal{G}}$. Además, por construcción,~$g$ es un
$\cal{E}$-funtor y~$u$ un $\cal{E}$-homomorfismo. A los datos anteriores
corresponde entonces un isomorfismo funtorial
$v\colon gf\to\id_{\cal{F}}$, definido por la condición $f*v=u*f$, y se
comprueba inmediatamente que también es un $\cal{E}$-morfismo, cqfd.

\begin{definition}
\label{VI.4.3} Si se cumplen las condiciones anteriores, se dice que~$f$ es
una equivalencia de categorías sobre~$\cal{E}$, o una
$\cal{E}$-equivalencia.
\end{definition}

\begin{corollaire}
\label{VI.4.4} Supongamos que el funtor de proyección
$p\colon\cal{F}\to\cal{E}$ es un funtor transportable, es decir, que para
todo isomorfismo $\alpha\colon T\to S$ en~$\cal{E}$ y todo objeto~$\xi$
de~$\cal{F}_T$, existe un isomorfismo~$u$ en~$\cal{F}$ de fuente~$\xi$ tal
que $p(u)=\alpha$. Entonces todo $\cal{E}$-funtor
$f\colon\cal{F}\to\cal{G}$ que sea una equivalencia de categorías es una
$\cal{E}$-equivalencia.
\end{corollaire}

Se deduce
\marginpar{161}
del criterio (iii~bis).

\begin{corollaire}
\label{VI.4.5} Sea $f\colon\cal{F}\to\cal{G}$ una
$\cal{E}$-equivalencia. Entonces, para toda categoría~$\cal{H}$
sobre~$\cal{E}$, los funtores correspondientes:
\begin{align*}
\SheafHom_{\cal{E}/-}(\cal{G},\cal{H})&\to
\SheafHom_{\cal{E}/-}(\cal{F},\cal{H})\\
\SheafHom_{\cal{E}/-}(\cal{H},\cal{F})
&\to\SheafHom_{\cal{E}/-}(\cal{H},\cal{G})
\end{align*}
(\cf \No \Ref{VI.2}) son equivalencias de categorías.
\end{corollaire}

Esto se deduce del criterio~(i) mediante el razonamiento habitual.

\section{Morfismos cartesianos, imágenes inversas, funtores
cartesianos}
\label{VI.5}

Sea~$\cal{F}$ una categoría sobre~$\cal{E}$, con funtor de
proyección~$p$.
\ifthenelse{\boolean{orig}}
{}
{\enlargethispage{.5cm}}

\begin{definition}
\label{VI.5.1} Consideremos un morfismo
$$
\alpha\colon\eta\to\xi
$$
en~$\cal{F}$, y sean
$$
S=p(\xi),\quad T=p(\eta),\quad f=p(\alpha)%
\ifthenelse{\boolean{orig}}{}{.}
$$
Se dice que~$\alpha$ es un \emph{morfismo cartesiano}
\index{cartesiano (morfismo)|hyperpage}%
si, para todo $\eta'\in\Ob(\cal{F}_T)$ y todo $f$-morfismo
$u\colon\eta'\to\xi$, existe un único $T$-morfismo
$\overline{u}\colon\eta'\to\eta$ tal que
$u=\alpha\circ\overline{u}$.
\end{definition}

Esto significa, pues, que para todo $\eta'\in\Ob(\cal{F}_T)$, la
aplicación $v\mto\alpha\circ v$:
\begin{equation}
\tag{i} \Hom_T(\eta',\eta)\to\Hom_f(\eta',\xi)
\end{equation}
es biyectiva. Significa también que el par $(\eta,\alpha)$
\emph{representa el funtor en~$\eta'$} $\cal{F}^{\circ}_T\to\Ens$ del
segundo miembro. Si, para un morfismo dado $f\colon T\to S$ en~$\cal{E}$
y un $\xi\in\Ob(\cal{F}_S)$ dado, existe tal par $(\eta,\alpha)$, es decir,
un morfismo
\marginpar{162}
cartesiano~$\alpha$ en~$\cal{F}$ de destino~$\xi$, tal que
$p(\alpha)=f$, entonces~$\eta$ queda determinado en~$\cal{F}_T$ salvo
isomorfismo único. Se dice entonces que \emph{existe la imagen inversa
de~$\xi$ por~$f$}, y un objeto~$\eta$ de~$\cal{F}_T$ provisto de un
$f$-morfismo cartesiano $\alpha\colon\eta\to\xi$ se llama \emph{una imagen
inversa de~$\xi$ por~$f$}.
\index{imagen inversa|hyperpage}%
A menudo se supone
\ifthenelse{\boolean{orig}}{elegido}{elegida}
tal imagen inversa cada vez que existe (con~$\cal{F}$ fijada); se denotará
entonces la imagen inversa mediante símbolos como~$f^*_{\cal{F}}(\xi)$, o
simplemente~$f^*(\xi)$
\label{indnot:fo}\oldindexnot{$f^*_{\cal{F}}(\xi)$ o $f^*(\xi)$|hyperpage}%
\ifthenelse{\boolean{orig}}{donde~$\xi\times_S T$}
{o bien~$\xi\times_S T$}
cuando estas notaciones no ocasionen
\ifthenelse{\boolean{orig}}{unas confusiones;}{confusión;}
el morfismo canónico $\alpha\colon\eta\to\xi$ se denotará, en lo que
sigue, por~$\alpha_f(\xi)$.
\label{indnot:fp}\oldindexnot{$\alpha_f(\xi)$|hyperpage}%
Si para todo $\xi\in\Ob(\cal{F}_S)$ existe la imagen inversa de~$\xi$
por~$f$, se dirá también que \emph{existe el funtor de imagen inversa
por~$f$ en~$\cal{F}$}, y~$f^*(\xi)$ se convierte entonces en un
\emph{funtor covariante en~$\xi$}, de~$\cal{F}_S$ en~$\cal{F}_T$. Esto
procede de que el segundo miembro de~(i) depende covariantemente de~$\xi$,
es decir, más precisamente, designa un funtor de
$\cal{F}^{\circ}_T\times\cal{F}_S$ en~$\Ens$. Esta dependencia funtorial
de~$f^*(\xi)$ se explicita así: consideremos $f$-morfismos cartesianos
$$
\alpha\colon\eta\to\xi,\quad\alpha'\colon\eta'\to\xi'
$$
y un $S$-morfismo $\lambda\colon\xi\to\xi'$; entonces existe un único
$T$-morfismo $\mu\colon\eta\to\eta'$ tal que
$$
\alpha'\mu=\lambda\alpha
$$
(como se deduce de que~$\alpha'$ es cartesiano).

Señalemos también el hecho inmediato siguiente: consideremos un diagrama
conmutativo
$$
\xymatrix@C=1.2cm{
\xi\ar[d]_-{\lambda}&\ar[l]_-{\alpha}\eta\ar[d]^-{\mu}\\
\xi'&\ar[l]_-{\alpha'}\eta'
}
$$
en~$\cal{F}$,
\marginpar{163}
donde~$\alpha$ y~$\alpha'$ son $f$-morfismos,~$\lambda$ un
$S$-isomorfismo y~$\mu$ un $T$-isomorfismo. \emph{Para que~$\alpha$ sea
cartesiano, es necesario y suficiente que~$\alpha'$ lo sea}.

\begin{definition}
\label{VI.5.2} Un $\cal{E}$-funtor $F\colon\cal{F}\to\cal{G}$ se llama
\emph{funtor cartesiano}
\index{cartesiano (funtor)|hyperpage}%
si transforma los morfismos cartesianos en morfismos cartesianos.
Se denota por $\SheafHom_\cart(\cal{F},\cal{G})$
\label{indnot:fq}\oldindexnot{$\SheafHom_\cart(\cal{F},\cal{G})$|hyperpage}%
la subcategoría plena de
$\SheafHom_{\cal{E}/-}(\cal{F},\cal{G})$ formada por los funtores
cartesianos.

Por ejemplo, considerando~$\cal{E}$ como una categoría sobre~$\cal{E}$
mediante el funtor identidad, todo morfismo de~$\cal{E}$ es cartesiano; por
tanto, un funtor cartesiano de~$\cal{E}$ en~$\cal{F}$ es un funtor sección
$F\colon\cal{E}\to\cal{F}$ que transforma todo morfismo de~$\cal{E}$ en
un morfismo cartesiano; tal funtor se llama una \emph{sección cartesiana}
de~$\cal{F}$ sobre~$\cal{E}$.
\end{definition}

\begin{proposition}
\label{VI.5.3} \textup{(i)}~Un funtor $F\colon\cal{F}\to\cal{G}$ que sea
una $\cal{E}$-equivalencia es un funtor cartesiano.
\textup{(ii)}~Sean~$F,G$ dos $\cal{E}$-funtores
\emph{isomorfos}~$\cal{F}\to\cal{G}$. Si uno es cartesiano, el otro también
lo es. \textup{(iii)}~El compuesto de dos funtores cartesianos
$\cal{F}\to\cal{G}$ y $\cal{G}\to\cal{H}$ es un funtor cartesiano.
\end{proposition}

La afirmación~(iii) es trivial por la definición; (ii) se deduce de la
observación que precede a~\Ref{VI.5.2}; (i) se deduce fácilmente de la
definición y del criterio~\Ref{VI.4.2}~(iii). Más precisamente, un
morfismo~$\alpha$ en~$\cal{F}$ es cartesiano si y solo si~$F(\alpha)$ lo
es.

\begin{corollaire}
\label{VI.5.4}
Sea $F\colon\cal{F}\to\cal{G}$ una $\cal{E}$-equivalencia. Entonces, para
toda categoría~$\cal{H}$ sobre~$\cal{E}$, los funtores correspondientes
$G\mto G\circ F$ y $G\mto F\circ G$ inducen equivalencias de categorías:
\begin{eqnarray*}
\SheafHom_\cart(\cal{G},\cal{H})\lto{\approx}
\SheafHom_\cart(\cal{F},\cal{H})\\
\SheafHom_\cart(\cal{H},\cal{F})\lto{\approx}
\SheafHom_\cart(\cal{H},\cal{G})
\end{eqnarray*}
\end{corollaire}

Esto se deduce, de la manera habitual, del criterio~(i) de~\Ref{VI.4.2} y
de~\Ref{VI.5.3}~(i), (ii) y~(iii).
\marginpar{164}
Puede precisarse que \emph{el $\cal{E}$-funtor
$G\colon\cal{G}\to\cal{H}$ es cartesiano si y solo si $G\circ F$ lo es},
y, análogamente, \emph{un $\cal{E}$-funtor
$G\colon\cal{H}\to\cal{F}$ es cartesiano si y solo si $F\circ G$ lo es}.

De
\ifthenelse{\boolean{orig}}{\Ref{VI.5.4}~(iii)}{\Ref{VI.5.3}~(iii)}
se deduce que, si se considera la subcategoría~$\Cat^\cart_{/\cal{E}}$
\label{indnot:fr}\oldindexnot{$\Cat^\cart_{/\cal{E}}$|hyperpage}%
de~$\Cat_{/\cal{E}}$ cuyos objetos son los mismos que los de
$\Cat_{/\cal{E}}$ y cuyos morfismos son los funtores \emph{cartesianos},
se tienen, como en el~\No \Ref{VI.2}, operaciones de composición:
$$
\SheafHom_\cart(\cal{F},\cal{G})\times
\SheafHom_\cart(\cal{G},\cal{H})\to\SheafHom_\cart(\cal{F},\cal{H})
$$
\ifthenelse{\boolean{orig}}{inducida}{inducidas}
por las del~\No \Ref{VI.2}, que permiten considerar
$\SheafHom_\cart(\cal{F},\cal{G})$ como un funtor en~$\cal{F}$,
$\cal{G}$, de la categoría $\left(\Cat^\cart_{/\cal{E}}\right)^\circ
\times\Cat^\cart_{/\cal{E}}$
\ifthenelse{\boolean{orig}}{en~$\Cat$).}{en~$\Cat$.}
Necesitaremos esta observación sobre todo para el caso en que
$\cal{F}=\cal{G}$:

\begin{definition}
\label{VI.5.5}
Sea $\cal{F}$ una categoría sobre~$\cal{E}$. Se denota por
$$
\varprojLim\cal{F}/\cal{E}
$$
\label{indnot:fs}\oldindexnot{$\varprojLim\cal{F}/\cal{E}$|hyperpage}%
la categoría de los $\cal{E}$-funtores cartesianos
$\cal{E}\to\cal{F}$, es decir, de las secciones cartesianas de~$\cal{F}$
sobre~$\cal{E}$.
\end{definition}

De acuerdo con lo que acaba de decirse,
$\varprojLim\cal{F}/\cal{E}$ es un funtor en~$\cal{F}$, de la categoría
$\Cat^\cart_{/\cal{E}}$ en la categoría~$\Cat$.

Veremos más adelante las relaciones entre esta operación~$\varprojLim$ y la
noción de límite proyectivo de categorías, así como numerosos ejemplos.

\section[Categorías fibradas y categorías prefibradas]{Categorías fibradas y
categorías prefibradas. Productos y cambio de base en ellas}
\label{VI.6}

\begin{definition}
\label{VI.6.1}
Una categoría $\cal{F}$ sobre~$\cal{E}$ se llama una
\emph{categoría fibrada}
\index{fibrada (categoría)|hyperpage}%
(y se dice entonces que el funtor $\cal{F}\to\cal{E}$ es un
\emph{fibrado})
\index{fibrado (funtor)|hyperpage}%
si satisface los dos axiomas siguientes:
\begin{itemize}
\item{$\Fib_{\mathrm{I}}$}
\marginpar{165}
Para todo morfismo $f\colon T\to S$ en~$\cal{E}$, existe el funtor de
imagen inversa por~$f$ en~$\cal{F}$.
\item{$\Fib_{\mathrm{II}}$} El compuesto de dos morfismos cartesianos es
cartesiano.
\end{itemize}
Una categoría $\cal{F}$ sobre~$\cal{E}$ que satisface la condición
$\Fib_{\mathrm{I}}$ se llama una \emph{categoría prefibrada sobre~$\cal{E}$}.
\index{prefibrada (categoría)|hyperpage}%
\end{definition}

Si~$\cal{F}$ es una categoría fibrada (\resp prefibrada) sobre~$\cal{E}$,
una subcategoría $\cal{G}$ de~$\cal{F}$ se llama una
\emph{subcategoría fibrada}
\index{subcategoría fibrada (\resp prefibrada)|hyperpage}%
(\resp una \emph{subcategoría prefibrada}) si es una categoría fibrada
(\resp prefibrada) sobre~$\cal{E}$ y, además, el funtor de inclusión es
cartesiano. Si, por ejemplo, $\cal{G}$ es una subcategoría \emph{plena}
de~$\cal{F}$, se ve que esto significa que, para todo morfismo
$f\colon T\to S$ en~$\cal{E}$ y todo $\xi\in\Ob(\cal{G}_S)$,
$f^*_{\cal{F}}(\xi)$ es $T$-isomorfo a un objeto de~$\cal{G}_T$. Otro
caso interesante es el siguiente: siendo $\cal{F}$ una categoría fibrada
sobre~$\cal{E}$, consideremos la subcategoría $\cal{G}$ de~$\cal{F}$ que
tiene los mismos objetos y cuyos morfismos son los morfismos
\emph{cartesianos} de~$\cal{F}$; en particular, los morfismos de
$\cal{G}_S$ son los isomorfismos de~$\cal{F}_S$. Se ve inmediatamente que
es, en efecto, una subcategoría fibrada de~$\cal{F}$, pues en la biyección
$$
\Hom_T(\eta',\eta) \isomto \Hom_f(\eta',\xi)
$$
relativa a un $f$-morfismo cartesiano $\alpha$ en~$\cal{F}$, a los
$T$-isomorfismos del primer miembro corresponden los morfismos cartesianos
del segundo. Por definición, las secciones cartesianas
$\cal{E}\to\cal{F}$
\ifthenelse{\boolean{orig}}{(corresponden}{corresponden}
entonces biunívocamente a los $\cal{E}$-funtores cualesquiera
$\cal{E}\to\cal{G}$ (pero obsérvese que el funtor natural
$$
\SheafHom_{\cal{E}/-}(\cal{E},\cal{G})\to
\SheafHom_\cart(\cal{E},\cal{F})=
\varprojLim(\cal{F}/\cal{E})
$$
es fiel, pero en general no es plenamente fiel, es decir, no es un
isomorfismo).

\begin{remarquesstar}
Sea $\cal{F}$ una categoría sobre~$\cal{E}$. Las condiciones siguientes
son equivalentes:~(i)~Todos los morfismos de~$\cal{F}$ son cartesianos;
(ii)~$\cal{F}$ es una categoría fibrada sobre~$\cal{E}$ y las
$\cal{F}_S$ son grupoides%
\ifthenelse{\boolean{orig}}{,}{}
(es decir, todo morfismo en~$\cal{F}_S$ es un isomorfismo). Se dice
entonces que $\cal{F}$ es una categoría \emph{fibrada en grupoides}
\index{fibrada en grupoides|hyperpage}%
\marginpar{166}
sobre~$\cal{E}$. Son estas las que se encuentran sobre todo en la
\og teoría de los espacios de módulos\fg. Si $\cal{E}$ es un grupoide, se
demuestra que las condiciones~(i) y~(ii) equivalen también a la siguiente:
(iii)~$\cal{F}$ es un grupoide y el funtor de proyección
$p\colon\cal{F}\to\cal{E}$ es transportable (\cf \Ref{VI.4.4}). Por
ejemplo, si $\cal{E}$ y $\cal{F}$ son grupoides tales que
\ifthenelse{\boolean{orig}}
{$\Ob \cal{E}$ y $\Ob \cal{F}$}
{$\Ob(\cal{E})$ y $\Ob(\cal{F})$}
se reducen a un punto, de modo que $\cal{E}$ y $\cal{F}$ quedan definidos,
salvo isomorfismo, por grupos $E$ y $F$, y el funtor
$p\colon\cal{F}\to\cal{E}$ queda definido por un homomorfismo de grupos
$p\colon F\to E$, entonces $\cal{F}$ está
\ifthenelse{\boolean{orig}}{fibrado}{fibrada}
sobre~$\cal{E}$ si y solo si $p$ es sobreyectivo, es decir, si $p$ define
una \emph{extensión} del grupo $E$ por el grupo $G=\Ker p$.
\end{remarquesstar}

\begin{proposition}
\label{VI.6.2}
Sea $F\colon\cal{F}\to\cal{G}$ una $\cal{E}$-equivalencia. Para que
$\cal{F}$ sea una categoría fibrada (\resp prefibrada) sobre~$\cal{E}$,
es necesario y suficiente que $\cal{G}$ lo sea.
\end{proposition}

Se deduce fácilmente de las definiciones y de la observación señalada más
arriba, según la cual un morfismo
\ifthenelse{\boolean{orig}}{$\alpha x$}{$\alpha$}
en~$\cal{F}$ es cartesiano si y solo si $F(\alpha)$ lo es.

\begin{proposition}
\label{VI.6.3}
Sean $\cal{F}_1$, $\cal{F}_2$ dos categorías sobre~$\cal{E}$, y sea
$\alpha=(\alpha_1,\alpha_2)$ un morfismo en
$\cal{F}=\cal{F}_1\times_{\cal{E}}\cal{F}_2$. Para que $\alpha$ sea
cartesiano, es necesario y suficiente que sus componentes lo sean.
\end{proposition}

Sea, en efecto, $\xi_i$ el destino y $\eta_i$ la fuente de~$\alpha_i$, y
sea $f\colon T\to S$ el morfismo de~$\cal{E}$ tal que $\alpha_1$ y
$\alpha_2$ sean $f$-morfismos. Para todo
$\eta'=(\eta'_1,\eta'_2)$ en~$\cal{F}_T$, se tiene un diagrama conmutativo
$$
\xymatrix{ \Hom_T(\eta',\eta) \ar[r] \ar[d] & \Hom_f(\eta',\xi) \ar[d]
\\ \Hom_T(\eta'_1,\eta_1)\times\Hom_T(\eta'_2,\eta_2) \ar[r] &
\Hom_f(\eta'_1,\xi_1)\times \Hom_f(\eta'_2,\xi_2) }
$$
donde las flechas verticales son biyecciones. Por tanto, si una de las
flechas horizontales es una biyección, también lo es la otra. Esto muestra
ya que, si $\alpha_1$, $\alpha_2$ son cartesianos (y, por tanto, la segunda
flecha horizontal es biyectiva), entonces $\alpha$ lo es. El recíproco se
ve tomando en el diagrama anterior $\eta'_i=\eta_i$, de donde
$\Hom_T(\eta'_i,\eta_i)\neq\emptyset$: primero para $i=2$, lo que prueba
que $\alpha_1$ es cartesiano, y luego para $i=1$, lo que prueba que
$\alpha_2$ es cartesiano.

\begin{corollaire}
\label{VI.6.4}
Sea
\marginpar{167}
$\cal{F}=\cal{F}_1\times_{\cal{E}}\cal{F}_2$, y sea $F=(F_1,F_2)$ un
$\cal{E}$-funtor $\cal{G}\to\cal{F}$. Para que
\ifthenelse{\boolean{orig}}{$\cal{F}$}{$F$}
sea cartesiano, es necesario y suficiente que $F_1$ y $F_2$ lo sean. Se
obtiene así un isomorfismo de categorías
$$
\SheafHom_\cart(\cal{G}, \cal{F}_1 \times_{\cal{E}} \cal{F}_2) \isomto
\SheafHom_\cart(\cal{G},\cal{F}_1)
\times\SheafHom_\cart(\cal{G},\cal{F}_2)
$$
y, en particular (tomando $\cal{G}=\cal{E}$), un isomorfismo de categorías
$$
\varprojLim(\cal{F}_1 \times_{\cal{E}}
\cal{F}_2/\cal{E}) \isomto
\varprojLim(\cal{F}_1/\cal{E}) \times
\varprojLim(\cal{F}_2/\cal{E})
$$
\end{corollaire}

\begin{corollaire}
\label{VI.6.5}
Sean $\cal{F}_1$ y $\cal{F}_2$ dos categorías fibradas (\resp prefibradas)
sobre~$\cal{E}$; entonces su producto fibrado
$\cal{F}=\cal{F}_1\times_{\cal{E}}\cal{F}_2$ es una categoría fibrada
(\resp prefibrada) sobre~$\cal{E}$.
\end{corollaire}

Estos resultados se extienden, por lo demás, al producto fibrado de una
familia cualquiera de categorías sobre~$\cal{E}$.

\begin{proposition}
\label{VI.6.6}
Sean $\cal{F}$ una categoría sobre~$\cal{E}$, con funtor de proyección
$p$, y $\lambda\colon\cal{E}'\to\cal{E}$ un funtor; consideremos
$\cal{F}'=\cal{F}\times_{\cal{E}}\cal{E}'$ como una categoría
sobre~$\cal{E}'$ mediante el funtor de proyección
$p'=p\times_{\cal{E}}\id_{\cal{E}'}$. Sea $\alpha'$ un morfismo
de~$\cal{F}'$. Para que $\alpha'$ sea un morfismo cartesiano, es necesario
y suficiente que su imagen $\alpha$ en~$\cal{F}$ lo sea.
\end{proposition}

La demostración es inmediata y se deja al lector.

\begin{corollaire}
\label{VI.6.7}
Para todo funtor cartesiano $F\colon\cal{F}\to\cal{G}$ de categorías
sobre~$\cal{E}$, el funtor $F'=F\times_{\cal{E}}\cal{E}'$ de
$\cal{F}'=\cal{F}\times_{\cal{E}}\cal{E}'$ en
$\cal{G}'=\cal{G}\times_{\cal{E}}\cal{E}'$ es cartesiano.
\end{corollaire}

Por consiguiente, el funtor $\SheafHom_{\cal{E}}(\cal{F},\cal{G})\to
\SheafHom_{\cal{E}'}(\cal{F}',\cal{G}')$ considerado en
\ifthenelse{\boolean{orig}}{\ignorespaces}{el}
\No \Ref{VI.3} induce un funtor
$$
\SheafHom_\cart(\cal{F},\cal{G})\to
\SheafHom_\cart(\cal{F}',\cal{G}') \rm{;}
$$
en otras palabras, para $\cal{F}$, $\cal{G}$ fijas, \emph{puede
considerarse}
$$
\SheafHom_\cart(\cal{F}\times_{\cal{E}}\cal{E}',
\cal{G}\times_{\cal{E}}\cal{E}')
$$
\emph{como
\marginpar{168}
un funtor en~$\cal{E}'$, de la categoría
${\Cat_{/\cal{E}}}^\circ$ en~$\Cat$}. Si se hacen variar asimismo
$\cal{F}$, $\cal{G}$, se obtiene%
\ifthenelse{\boolean{orig}}{,}{}
un funtor de la categoría ${\Cat_{/\cal{E}}}^\circ\times\big(
\Cat^\cart_{/\cal{E}}\big)^\circ\times\Cat^\cart_{/\cal{E}}$
en~$\Cat$. Al tener en cuenta el isomorfismo
$$
\Hom_{\cal{E}'}(\cal{F}',\cal{G}') \isomto
\Hom_\cal{E}(\cal{F}\times_\cal{E}\cal{E}',\cal{G})
$$
considerado en el~\No \Ref{VI.3}, los $\cal{E}'$-funtores cartesianos
de~$\cal{F}'$ en~$\cal{G}'$ corresponden a los $\cal{E}$-funtores
$\cal{F}\times_\cal{E}\cal{E}'\to\cal{G}$ que transforman todo morfismo
cuya primera proyección es un morfismo cartesiano de~$\cal{F}$ en un
morfismo cartesiano de~$\cal{G}$. Tomando $\cal{F}=\cal{E}$, se obtiene
(tras un cambio de notación):

\begin{corollaire}
\label{VI.6.8}
$\varprojLim(\cal{F}'/\cal{E}')$ es isomorfa a la subcategoría plena de
$\SheafHom_{\cal{E}/-}(\cal{E}',\cal{F})$ formada por los
$\cal{E}$-funtores $\cal{E}'\to\cal{F}$ que transforman los morfismos
cualesquiera en morfismos cartesianos.
En particular, si $\cal{F}$ es una categoría fibrada y
$\tilde{\cal{F}}$ es la subcategoría de~$\cal{F}$ cuyos morfismos son
los morfismos cartesianos de~$\cal{F}$, entonces se tiene una biyección
$$
\Ob \varprojLim(\cal{F}'/\cal{E}') \isomto
\Hom_{\cal{E}/-}(\cal{E}',\tilde{\cal{F}}) \rm{.}
$$
\end{corollaire}

Esto precisa la manera en que la expresión
$\varprojLim(\cal{F}\times_{\cal{E}}\cal{E}'/\cal{E}')$ debe considerarse
como un funtor en~$\cal{E}'$ y en~$\cal{F}$, de la categoría
${\Cat_{/\cal{E}}}^\circ\times\Cat^\cart_{/\cal{E}}$ en la categoría
$\Cat$. Más adelante se verá una dependencia funtorial más completa con
respecto a~$\cal{E}'$, cuando se exija que $\cal{F}$ sea una categoría
fibrada.

\begin{corollaire}
\label{VI.6.9}
Sea $\cal{F}$ una categoría fibrada (\resp prefibrada) sobre~$\cal{E}$;
entonces $\cal{F}'=\cal{F}\times_{\cal{E}}\cal{E}'$ es una categoría
fibrada (\resp prefibrada) sobre~$\cal{E}'$.
\end{corollaire}

\begin{proposition}
\label{VI.6.10} Sean $\cal{F}$, $\cal{G}$ categorías prefibradas
sobre~$\cal{E}$, y sea $F$ un $\cal{E}$-funtor cartesiano de~$\cal{F}$
en~$\cal{G}$. Para que $F$ sea fiel,
\marginpar{169}
\ifthenelse{\boolean{orig}}
{\resp plenamente fiel, (\resp una $\cal{E}$-equivalencia)}
{(\resp plenamente fiel, \resp una $\cal{E}$-equivalencia)}
es necesario y suficiente que, para todo
\ifthenelse{\boolean{orig}}
{$S\in\Ob \cal{E}$,}
{$S\in\Ob(\cal{E})$,}
el funtor inducido $F_S\colon\cal{F}_S\to\cal{G}_S$ sea fiel
(\resp plenamente fiel, \resp una equivalencia).
\end{proposition}

Demostración inmediata a partir de las definiciones.

Para terminar este número, damos algunas propiedades de las categorías
fibradas que utilizan el axioma $\Fib_{\mathrm{II}}$.
\begin{proposition}
\label{VI.6.11} Sea $\cal{F}$ una categoría prefibrada sobre~$\cal{E}$.
Para que $\cal{F}$ sea fibrada, es necesario y suficiente que satisfaga la
condición siguiente:

$\Fib_{\mathrm{II}}'$: Sea $\alpha\colon\eta\to\xi$ un morfismo
cartesiano en~$\cal{F}$ sobre el morfismo $f\colon T\to S$ de~$\cal{E}$.
Para todo morfismo $g\colon U\to T$ en~$\cal{E}$ y todo
\ifthenelse{\boolean{orig}}
{$\zeta\in\Ob \cal{F}_U$,}
{$\zeta\in\Ob(\cal{F}_U)$,}
la aplicación $u\mto\alpha\circ u$:
$$
\Hom_g(\zeta,\eta)\to\Hom_{fg}(\zeta,\xi)
$$
es biyectiva.
\end{proposition}

En otras palabras, en una categoría \emph{fibrada} sobre~$\cal{E}$, los
diagramas cartesianos se caracterizan por una propiedad a priori más
fuerte que la de la
\ifthenelse{\boolean{orig}}
{definición, (que se obtiene tomando $g=\id_T$ en}
{definición (que se obtiene tomando $g=\id_T$ en}
el enunciado precedente).
\begin{corollaire}
\label{VI.6.12} Sean $\cal{F}$ una categoría sobre~$\cal{E}$ y $\alpha$
un morfismo en~$\cal{F}$. Para que $\alpha$ sea un isomorfismo, es
necesario que $p(\alpha)=f$ sea un isomorfismo y que $\alpha$ sea
cartesiano; el recíproco es cierto si $\cal{F}$ está fibrada
sobre~$\cal{E}$.
\end{corollaire}

En efecto, si $\alpha$ es un isomorfismo, evidentemente también lo es
$f=p(\alpha)$; para todo
\ifthenelse{\boolean{orig}}
{$\eta'\in\Ob\cal{F}_T$,}
{$\eta'\in\Ob(\cal{F}_T)$,}
la aplicación $u\mto\alpha\circ u$
$$
\Hom(\eta',\eta)\to\Hom(\eta',\xi)
$$
es biyectiva; como $f$ es un isomorfismo, se ve inmediatamente
\ifthenelse{\boolean{orig}}{que un}{que un}
elemento del primer miembro es un $T$-morfismo si y solo si su imagen en el
segundo es un $f$-morfismo; se obtiene, por tanto, una biyección
$$
\Hom_T(\eta',\eta)\to\Hom_f(\eta',\xi)
$$
lo
\marginpar{170}
que prueba la primera afirmación. Recíprocamente, supongamos que $f$ es un
isomorfismo y que $\alpha$ satisface la condición enunciada en
\ifthenelse{\boolean{orig}}{$\Fib_{\mathrm{II}'}$}{$\Fib_{\mathrm{II}}'$}
(lo que significa, cuando $\cal{F}$ está
\ifthenelse{\boolean{orig}}
{fibrado}
{fibrada}
sobre~$\cal{E}$, que $\alpha$ es cartesiano); entonces se ve
inmediatamente que, para todo
\ifthenelse{\boolean{orig}}
{$\zeta\in\Ob\cal{F}$}
{$\zeta\in\Ob(\cal{F})$}%
, la aplicación $u\mto\alpha\circ u$ de~$\Hom(\zeta,\eta)$ en
$\Hom(\zeta,\xi)$ es biyectiva; por tanto, $\alpha$ es un isomorfismo.
\begin{corollaire}
\label{VI.6.13} Sean $\alpha\colon\eta\to\xi$ y
$\beta\colon\zeta\to\eta$ dos morfismos componibles en la categoría
$\cal{F}$ fibrada sobre~$\cal{E}$. Si $\alpha$ es cartesiano, entonces
$\beta$ lo es si y solo si $\alpha\beta$ lo es.
\end{corollaire}

Se utiliza la definición de los morfismos cartesianos en la forma reforzada
de~\Ref{VI.6.11}.

\section{Categorías con clivaje sobre~$\mathcal{E}$}
\label{VI.7}

\begin{definition}
\label{VI.7.1} Sea $\cal{F}$ una categoría sobre~$\cal{E}$. Se llama
\emph{clivaje}
\index{clivaje|hyperpage}%
de~$\cal{F}$ sobre~$\cal{E}$ a una función que asigna a todo
$f\in\Fl(\cal{E})$ un funtor de imagen inversa por~$f$ en~$\cal{F}$,
denotado por~$f^*$. El clivaje se dice \emph{normalizado}
\index{clivaje normalizado|hyperpage}%
\index{normalizado (clivaje)|hyperpage}%
si $f=\id_S$ implica $f^*=\id_{\cal{F}_S}$. Se llama
\emph{categoría con clivaje} (\resp \emph{categoría con clivaje
normalizado})
\index{categoría con clivaje (\resp con clivaje normalizado)|hyperpage}%
a una categoría $\cal{F}$ sobre~$\cal{E}$ dotada de un clivaje
(\resp de un clivaje normalizado).
\end{definition}

Es evidente que $\cal{F}$ admite un clivaje si y solo si $\cal{F}$ es
prefibrada sobre~$\cal{E}$, y en tal caso $\cal{F}$ admite un clivaje
normalizado. El conjunto de los clivajes sobre~$\cal{F}$ está en
correspondencia biunívoca con el conjunto de los subconjuntos $K$ de
$\Fl(\cal{F})$ que satisfacen las condiciones siguientes:
\begin{enumerate}
\item[a)] Los $\alpha\in K$ son morfismos cartesianos.
\item[b)] Para todo morfismo $f\colon T\to S$ en~$\cal{E}$ y todo
$\xi\in\Ob(\cal{F}_S)$, existe un único $f$-morfismo en~$K$ de destino
$\xi$.
\end{enumerate}

Para que el clivaje definido por~$K$ sea normalizado, es necesario y
suficiente que $K$ satisfaga además la condición
\begin{enumerate}
\item[c)] Los morfismos identidad de~$\cal{F}$ pertenecen a~$K$.
\end{enumerate}

Los
\marginpar{171}
morfismos que son elementos de~$K$ podrán llamarse
\emph{\og morfismos de transporte\fg}
\index{transporte (morfismo de)|hyperpage}%
para el clivaje considerado.

La noción de isomorfismo de categorías con clivaje sobre~$\cal{E}$ es
clara. Más generalmente, pueden definirse los morfismos de
$\cal{E}$-categorías con clivaje como los funtores de
$\cal{E}$-categorías $\cal{F}\to\cal{G}$ que aplican los morfismos de
transporte en morfismos de transporte.
(Estos son, en particular, funtores cartesianos). De este modo, las
categorías con clivaje sobre~$\cal{E}$ son los objetos de una categoría,
la \emph{categoría de las categorías con clivaje sobre}~$\cal{E}$. El
lector explicitará la existencia de productos, ligada al hecho siguiente:
si una categoría sobre~$\cal{E}$ es producto de categorías $\cal{F}_i$
sobre~$\cal{E}$, provista cada una de un clivaje, entonces $\cal{F}$ queda
dotada del clivaje natural correspondiente. Se deja asimismo al lector
explicitar la noción de cambio de base en las categorías con clivaje.

Denotaremos por $\alpha_f(\xi)$ el morfismo canónico
$$
\alpha_f(\xi)\colon f^*(\xi)\to\xi.
$$
Como ya se ha dicho, es funtorial en~$\xi$, es decir, se tiene una
transformación natural
$$
\alpha_f\colon i_Tf^*\to i_S,
$$
donde, para todo $S\in\Ob(\cal{E})$, $i_S$ denota el funtor de inclusión
$$
i_S\colon\cal{F}_S\to\cal{F}
$$

Consideremos ahora morfismos
$$
f\colon T\to S\quad\text{ y }\quad g\colon U\to T
$$
en~$\cal{E}$, y sea $\xi\in\Ob(\cal{F}_S)$. Existe entonces un único
$U$-morfismo
$$
c_{f,g}(\xi)\colon g^*f^*(\xi)\to(fg)^*(\xi)
$$
que hace
\marginpar{172}
conmutativo el diagrama
$$
\xymatrix{
f^*(\xi) \ar[d]_{\alpha_f(\xi)} & & \ar[ll]_{\alpha_g(f^*(\xi))}
g^*(f^*(\xi)) \ar[d]^{c_{f,g}(\xi)} \\
\xi & & \ar[ll]^{\alpha_{fg}(\xi)} (fg)^*(\xi) }
$$
(en virtud de la definición de~$(fg)^*(\xi)$). Para $\xi$ variable,
este morfismo es funtorial; es decir, se tiene una transformación natural
$$
c_{f,g}\colon g^*f^*\to(fg)^*
$$
de funtores $\cal{F}_S\to\cal{F}_U$. Observemos inmediatamente:
\begin{proposition}
\label{VI.7.2} Para que la categoría con clivaje $\cal{F}$
sobre~$\cal{E}$ sea fibrada, es necesario y suficiente que los
$c_{f,g}$ sean isomorfismos.
\end{proposition}

Tomando para~$f$ un isomorfismo y para~$g$ su inverso, y considerando los
isomorfismos $c_{f,g}$ y $c_{g,f}$, se concluye:
\begin{corollaire}
\label{VI.7.3} Si $\cal{F}$ es una categoría \emph{fibrada con clivaje}
sobre~$\cal{E}$, entonces, para todo isomorfismo $f\colon T\to S$ en
$\cal{E}$, $f^*$ es una equivalencia de categorías
$\cal{F}_S\to\cal{F}_T$.
\end{corollaire}

\begin{proposition}
\label{VI.7.4} Sea $\cal{F}$ una categoría con clivaje sobre~$\cal{E}$.
Se tiene

$\qquad\qquad A)\qquad
\left\{\begin{array}{ccc}c_{f,\id_T}(\xi)&=&\alpha_{\id_T}(f^*(\xi))\\
c_{\id_S,f}(\xi)&=&f^*(\alpha_{\id_S}(\xi))
\end{array}\right.$

\smallskip
$\qquad\qquad B)\qquad c_{f,gh}(\xi)\cdot
c_{g,h}(f^*(\xi))=c_{fg,h}(\xi)\cdot h^*(c_{f,g}(\xi))$

\noindent(En
\marginpar{173}
estas fórmulas, $f,g,h$ denotan morfismos
$$
V\to U\to T\to S
$$
y $\xi$ un objeto de~$\cal{F}_S$).
\end{proposition}

La primera y la segunda relación, en el caso de un clivaje normalizado,
adoptan la forma más sencilla
\begin{equation*}
\label{eq:VI.7.A'}\quad A')\qquad {c_{f,\id_T}=\id_{f^*},\quad
c_{\id_S,f}=\id_{f^*}.\qquad\qquad\qquad\qquad\qquad\qquad}
\end{equation*}

En cuanto a la tercera, puede visualizarse mediante la conmutatividad del
diagrama
\begin{equation*}
\label{eq:VI.7.D} \tag{$D$}
\begin{split}
\xymatrix@C=2.2cm{
h^*g^*f^*(\xi) \ar[r]^-{c_{g,h}(f^*(\xi))} \ar[d]_{h^*(c_{f,g}(\xi))}
& (gh)^*(f^*(\xi)) \ar[d]^{c_{f,gh}(\xi)} \\
h^*(fg)^*(\xi) \ar[r]_-{c_{fg,h}(\xi)} & (fgh)^*(\xi) }
\end{split}
\end{equation*}
En el caso de las categorías
\ifthenelse{\boolean{orig}}
{fibradas,}
{fibradas}
(donde los $c_{f,g}$ son isomorfismos), esta conmutatividad puede expresarse
intuitivamente por el hecho de que \emph{la utilización sucesiva de los
isomorfismos de la forma $c_{f,g}$ no conduce a \og identificaciones
contradictorias\fg}. También puede escribirse esta fórmula sin el argumento
$\xi$, utilizando el producto de convolución de transformaciones naturales:
$$
c_{fg,h}\circ (h^**c_{f,g})=c_{f,gh}\circ(c_{g,h}*f^*).
$$

\ifthenelse{\boolean{orig}}
{\ignorespaces}
{\enlargethispage{.5cm}}
La demostración de las dos primeras fórmulas de~\Ref{VI.7.4} es trivial;
esbocemos la de la tercera. Para ello, consideremos, además del cuadrado
$(D)$, el cuadrado de transformaciones naturales:
\begin{equation*}
\label{eq:VI.7.D'} \tag{$D'$}
\begin{split}
\xymatrix@C=2cm{
g^*f^*(\xi) \ar[r]^-{\alpha_g(f^*(\xi))}
\ar[d]_{c_{f,g}(\xi)} & f^*(\xi) \ar[d]^{\alpha_f(\xi)} \\
(fg)^*(\xi) \ar[r]_-{\alpha_{fg}(\xi)} & \xi }
\end{split}
\end{equation*}
que
\marginpar{174}
es conmutativo por definición de $c_{f,g}(\xi)$. Consideremos el diagrama
obtenido uniendo los vértices de~$(D)$ a los vértices correspondientes de
$(D')$ mediante las transformaciones naturales de la forma~$\alpha$:
$$
\begin{array}{ll}
\alpha_h(g^*f^*(\xi)),&\alpha_{gh}(f^*(\xi))\\
\alpha_h((fg)^*(\xi)),&\alpha_{fgh}(\xi).
\end{array}
$$
Las cuatro caras laterales del cubo así obtenido son también conmutativas:
para la de la izquierda, esto se debe a que la columna izquierda de~$(D)$
se deduce de la columna izquierda de~$(D')$
\ifthenelse{\boolean{orig}}{por aplicación de~$h$}{por aplicación de~$h^*$}%
, y a que $\alpha_h$ es una transformación natural; para las otras tres,
no es sino la definición de las operaciones~$c$ de los tres lados restantes
de~$(D)$. Así, las cinco caras del cubo distintas de la cara superior son
conmutativas. Se sigue que los dos $(fgh)$-morfismos
$h^*g^*f^*(\xi)\to(fgh)^*(\xi)$ definidos por~$(D)$ tienen el mismo
compuesto con $\alpha_{fgh}(\xi)\colon(fgh)^*(\xi)\to\xi$; por tanto, son
iguales por definición de~$(fgh)^*$.

Para lo que sigue, limitémonos a las categorías con clivaje
\emph{normalizado}. Una categoría tal da lugar a los objetos siguientes:
\begin{enumerate}
\item[a)] Una aplicación $S\mto\cal{F}_S$ de~$\Ob(\cal{E})$ en~$\Cat$.
\item[b)] Una aplicación $f\mto f^*$ que asigna a todo
$f\in\Fl(\cal{E})$, de fuente~$T$ y destino~$S$,
\ifthenelse{\boolean{orig}}
{una}
{un}
funtor $f^*\colon\cal{F}_S\to\cal{F}_T$.
\item[c)] Una aplicación $(f,g)\mto c_{f,g}$ que asigna a todo par de
flechas $(f,g)$ de~$\cal{E}$ una transformación natural
$c_{f,g}\colon g^*f^*\to(fg)^*$.
\end{enumerate}

Además, estos datos satisfacen las condiciones expresadas en las fórmulas
$A')$ y~$B)$ dadas más arriba. (N.B. Si no nos hubiéramos limitado al caso
de un clivaje normalizado, habría sido necesario introducir un dato
suplementario, a saber, una función $S\mto\alpha_S$ que asocia a todo objeto
$S$ de~$\cal{E}$ una transformación natural
$\alpha_S\colon(\id_S)^*\to\id_{\cal{F}_S}$; la condición $A')$ sería
entonces sustituida por la condición~$A)$).

Vamos
\marginpar{175}
a mostrar ahora cómo puede reconstruirse (salvo isomorfismo único) la
categoría con clivaje normalizado $\cal{F}$ sobre~$\cal{E}$ a partir de los
objetos precedentes.

\section[Categoría con clivaje definida por un pseudofuntor]{Categoría con clivaje definida por un pseudofuntor
$\mathcal{E}^\circ\to\Cat$} \label{VI.8}

Llamemos, para abreviar, \emph{pseudofuntor} de~$\cal{E}^\circ$
a~$\Cat$ (habría que decir pseudofuntor \emph{normalizado}) a un
conjunto de datos a), b), c) como los anteriores, que satisfacen las
condiciones $A')$ y~$B)$. En el número precedente hemos asociado a una
categoría con clivaje normalizado sobre~$\cal{E}$ un pseudofuntor
$\cal{E}^\circ\to\Cat$; aquí vamos a indicar la construcción inversa.
Dejaremos al lector la verificación de la mayor parte de los detalles, así
como del hecho de que estas construcciones son efectivamente inversas entre
sí. Con mayor precisión, convendría considerar los pseudofuntores
$\cal{E}^\circ\to\Cat$ como los objetos de una nueva categoría y demostrar
que nuestras construcciones proporcionan equivalencias, cuasi-inversas
entre sí, entre esta última y la categoría de las categorías con clivaje
sobre~$\cal{E}$ definida en el número precedente.

Pongamos
$$
\cal{F}_\circ=\underset{S\in\Ob(\cal{E})}\coprod
\ifthenelse{\boolean{orig}}
{\Ob\cal{F}(S),}
{\Ob(\cal{F}(S)),}
$$
la suma disjunta de los conjuntos
\ifthenelse{\boolean{orig}}
{$\Ob\cal{F}(S)$}
{$\Ob(\cal{F}(S))$}
(N.B. escribiremos aquí $\cal{F}(S)$ y no $\cal{F}_S$ para el valor en el
objeto $S$ de~$\cal{E}$ del pseudofuntor dado, a fin de evitar confusiones
de notación más adelante). Tenemos, pues, una aplicación evidente:
$$
p_\circ\colon\cal{F}_\circ\to\Ob\cal{E}.
$$
Sean
$$
\ifthenelse{\boolean{orig}}
{\xi}
{\overline{\xi}}
=(S,\xi),\quad \overline{\eta}=(T,\eta)\qquad\qquad
\text{(con $\xi\in\Ob\cal{F}(S)$, $\eta\in\Ob\cal{F}(T)$)}
$$
dos elementos de~$\cal{F}_\circ$, y sea $f\in\Hom(T,S)$; pondremos
$$
h_f(\overline{\eta},\overline{\xi})=\Hom_{\cal{F}(T)}(\eta,f^*(\xi)).
$$
\marginpar{176}%
Si además tenemos un morfismo $g\colon U\to T$ en~$\cal{E}$ y un
$\zeta\in\Ob\cal{F}(U)$, definimos una aplicación, denotada por
$(u,v)\mto u\circ v$:
$$
h_f(\overline{\eta},\overline{\xi})\times
h_g(\overline{\zeta},\overline{\eta})\to
h_{fg}(\overline{\zeta},\overline{\xi}),
$$
es decir, una aplicación
$$
\Hom_{\cal{F}(T)}(\eta,f^*(\xi))\times\Hom_{\cal{F}(U)}(\zeta,g^*(\eta))
\to \Hom_{\cal{F}(U)}(\zeta,(fg)^*(\xi)),
$$
mediante la fórmula
$$
u\circ v=c_{f,g}(\xi)\cdot g^*(u)\cdot v,
$$
es decir, $u\circ v$ es el compuesto de la sucesión
$$
\zeta
\ifthenelse{\boolean{orig}}
{\lto{u}}
{\lto{v}}
g^*(\eta)\lto{g^*(u)}
g^*f^*(\xi)\lto{c_{f,g}(\xi)} (fg)^*(\xi).
$$
Pondremos, por otra parte,
$$
h(\overline{\eta},\overline{\xi})=\underset{f\in\Hom(T,S)}\coprod
h_f(\overline{\eta},\overline{\xi}),
$$
y las aplicaciones de composición precedentes definen aplicaciones
$$
h(\overline{\eta},\overline{\xi})\times
h(\overline{\zeta},\overline{\eta})\to
h(\overline{\zeta},\overline{\xi}),
$$
mientras que la definición de los $h(\overline{\eta},\overline{\xi})$
implica una aplicación evidente:
$$
p_{\overline{\eta},\overline{\xi}}\colon
h(\overline{\eta},\overline{\xi})\to \Hom(T,S).
$$
Dicho esto, se verifican los puntos siguientes:
\begin{enumerate}
\item[1)] La composición entre elementos de los
\ifthenelse{\boolean{orig}}
{$h(\eta,\xi)$}
{$h(\overline{\eta},\overline{\xi})$}
es \emph{asociativa}.
\item[2)]
\marginpar{177}
Para todo
\ifthenelse{\boolean{orig}}
{$\overline{\xi}=(\xi,S)$}
{$\overline{\xi}=(S,\xi)$}
en~$\cal{F}_\circ$, consideremos el elemento
\ifthenelse{\boolean{orig}}
{}
{identidad}
de
$$
\ifthenelse{\boolean{orig}}
{h_{\id_S}(\overline{\xi},\overline{\xi})=
\Hom_{\cal{F}_S}(\id_S^*(\xi),\xi)=\Hom_{\cal{F}_S}(\xi,\xi),}
{h_{\id_S}(\overline{\xi},\overline{\xi})=
\Hom_{\cal{F}(S)}(\id_S^*(\xi),\xi)=\Hom_{\cal{F}(S)}(\xi,\xi),}
$$
y su imagen en $h(\overline{\xi},\overline{\xi})$. Este elemento es un
\emph{elemento neutro} por la izquierda y por la derecha para la
composición entre elementos de los
$h(\overline{\eta},\overline{\xi})$.

Esto muestra ya que \emph{se obtiene una categoría} $\cal{F}$ al poner
$$
\Ob\cal{F}=\cal{F}_\circ,\quad \Fl{\cal{F}}=
\ifthenelse{\boolean{orig}}
{\underset{\xi,\overline{\eta}\in\cal{F}_\circ} \coprod
h(\overline{\eta},\overline{\xi}).}
{\underset{\overline{\xi},\overline{\eta}\in\cal{F}_\circ}
\coprod h(\overline{\eta},\overline{\xi}).}
$$
(N.B. no se puede tomar simplemente como $\Fl\cal{F}$ la \emph{unión}
de los conjuntos $h(\overline{\eta},\overline{\xi})$, pues estos no son
necesariamente disjuntos). Además:
\item[3)] Las aplicaciones $p_\circ\colon \Ob\cal{F}\to\Ob\cal{E}$ y
$p_1=(p_{\overline{\eta},\overline{\xi}})\colon\Fl\cal{F}\to\Fl{\cal{E}}$
definen un \emph{funtor} $p\colon\cal{F}\to\cal{E}$. De este modo,
$\cal{F}$ se convierte en una categoría sobre~$\cal{E}$; además, la
aplicación evidente
$h_f(\overline{\eta},\overline{\xi})\to
\Hom(\overline{\eta},\overline{\xi})$ induce una \emph{biyección}
$$
h_f(\overline{\eta},\overline{\xi})\isomto
\Hom_f(\overline{\eta},\overline{\xi}).
$$
\item[4)] Las aplicaciones evidentes
$$ \Ob\cal{F}(S)\to\cal{F}_\circ=\Ob\cal{F},\qquad \Fl\cal{F}(S)\to
\Fl\cal{F},
$$
donde la segunda viene definida por las aplicaciones evidentes
$$
\Hom_{\cal{F}(S)}(\xi,\xi')=
h_{\id_S}(\overline{\xi},\overline{\xi}')\to
\Hom(\overline{\xi},\overline{\xi}')
$$
definen un \emph{isomorfismo}
$$
i_S \colon \cal{F}(S)\isomto\cal{F}_S.
$$
\item[5)] Para todo objeto $\overline{\xi}=(S,\xi)$ de~$\cal{F}$ y todo
morfismo $f\colon T\to S$ de~$\cal{E}$, consideremos
\marginpar{178}
el elemento $\overline{\eta}=(T,\eta)$ de~$\cal{F}_T$, con
$\eta=f^*(\xi)$, y el elemento $\alpha_f(\xi)$ de
$\Hom(\overline{\eta},\overline{\xi})$, imagen de~$\id_{f^*(\xi)}$ por
el morfismo
$\Hom_{\cal{F}(T)}(f^*(\xi),f^*(\xi))=
h_f(\overline{\eta},\overline{\xi})
\to \Hom_f(\overline{\eta},\overline{\xi})$. \emph{Este elemento es
cartesiano y es la identidad de $\overline{\xi}$ si} $f=\id_S$; en otros
términos, el conjunto de los $\alpha_f(\xi)$ define un
\emph{clivaje normalizado de~$\cal{F}$ sobre} $\cal{E}$. Además, por
construcción, es conmutativo el diagrama de funtores
$$
\xymatrix{
\cal{F}(S) \ar[r]^{f^*} \ar[d]_{i_S} & \cal{F}(T) \ar[d]^{i_T} \\
\cal{F}_S \ar[r]^{f^*_{\cal{F}}} & \cal{F}_T }
$$
donde $f^*_{\cal{F}}$ es el funtor imagen inversa por~$f$, relativo al
clivaje considerado sobre~$\cal{F}$. Finalmente:
\item[6)] las transformaciones naturales $c_{f,g}$ dadas con el
pseudofuntor son llevadas, mediante los isomorfismos $i_S$, a las
transformaciones naturales $c_{f,g}$ asociadas al clivaje de~$\cal{F}$.
\end{enumerate}

Nos limitamos a dar la verificación de 1) (que es, si acaso, menos trivial
que las demás). Basta probar la asociatividad de la composición entre los
\ifthenelse{\boolean{orig}}
{objetos}
{elementos}
de conjuntos de la forma $h_f(\overline{\eta},\overline{\xi})$.
Consideremos, pues, en~$\cal{E}$ morfismos
$$
S\lfrom{f} T\lfrom{g} U\lfrom{h} V
$$
y objetos
$$
\xi,\ \eta,\ \zeta,\ \tau
$$
en $\cal{F}(S),\cal{F}(T),\cal{F}(U),\cal{F}(V)$, respectivamente, y,
por último, elementos
$$
\begin{array}{ccc}
u\in
h_f(\overline{\eta},\overline{\xi})&=&\Hom_{\cal{F}(T)}(\eta,f^*(\xi))\\
v\in
h_g(\overline{\zeta},\overline{\eta})&=&\Hom_{\cal{F}(U)}(\zeta,g^*(\eta))\\
w\in
h_h(\overline{\tau},\overline{\zeta})&=&\Hom_{\cal{F}(V)}(\tau,h^*(\zeta)).
\end{array}
$$
Queremos
\marginpar{179}
probar la fórmula
$$
(u\circ v)\circ w = u\circ (v\circ w),
$$
que es una igualdad en
$\Hom_{\cal{F}(V)}(\tau,(fgh)^*(\xi))$. En virtud de las definiciones,
los dos miembros de esta igualdad se obtienen componiendo a lo largo de
los contornos superior e inferior, respectivamente, del diagrama
siguiente:
$$
\xymatrix{
\tau \ar[r]^w \ar[drr]_{v\circ w} & h^*(\zeta)
\ar[r]^{h^*(v)} \ar@/^1.4pc/[rrrrr]^{h^*(u\circ v)} & h^*g^*(\eta)
\ar[rr]^{h^*g^*(u)}
\ar[d]_{c_{g,h}(\eta)} & & h^*g^*f^*(\xi) \ar[rr]^{h^*(c_{f,g}(\xi))}
\ar[d]^{c_{g,h}(f^*(\xi))} & & h^*(fg)^*(\xi) \ar[d]^{c_{fg,h}(\xi)} \\
& & (gh)^*(\eta) \ar[rr]^{(gh)^*(u)} & & (gh)^*f^*(\xi)
\ar[rr]^{c_{f,gh}(\xi)} & & (fgh)^*(\xi) }
$$
Ahora bien, el cuadrado central es conmutativo porque $c_{g,h}$ es una
transformación natural, y el cuadrado de la derecha es conmutativo en
virtud de la condición~$B)$ para un pseudofuntor. De aquí se sigue el
resultado anunciado.

Naturalmente, queda por precisar, cuando el pseudofuntor considerado
proviene ya de una categoría con clivaje normalizado $\cal{F}'$
sobre~$\cal{E}$, cómo se obtiene un isomorfismo natural entre $\cal{F}'$
y~$\cal{F}$. Dejamos el detalle al lector.

Dejamos también al lector la interpretación, en términos de
pseudofuntores, de la noción de imagen inversa de una categoría con
clivaje $\cal{F}$ sobre~$\cal{E}$ por un funtor de cambio de base
$\cal{E}'\to\cal{E}$.

\section[Ejemplo]{Ejemplo: categoría con clivaje definida por un funtor
$\cal{E}^\circ\to\Cat$; categorías escindidas sobre~$\cal{E}$}
\label{VI.9}
Supongamos dado un funtor
$$
\ifthenelse{\boolean{orig}}
{\phi\colon\cal{E}^\circ\to\Cat,}
{\varphi\colon\cal{E}^\circ\to\Cat,}
$$
Este define entonces un pseudofuntor al poner
$$
\ifthenelse{\boolean{orig}}
{\cal{F}(S)=\phi(S),}
{\cal{F}(S)=\varphi(S),}
\quad f^*=\varphi(f),\quad c_{f,g}=\id_{(fg)^*}
$$
\marginpar{180}%
Por tanto, la construcción del número precedente nos da una categoría
$\cal{F}$ con clivaje sobre~$\cal{E}$, llamada asociada al funtor
$\varphi$. Para que una categoría con clivaje sobre~$\cal{E}$ sea
isomorfa a una categoría con clivaje definida por un funtor
$\varphi:\cal{E}^\circ\to\Cat$, es necesario y suficiente,
manifiestamente, que satisfaga las condiciones
$$
(fg)^*=g^*f^*,\quad c_{f,g}=\id_{(fg)^*}\quoi.
$$
En términos del conjunto $K$ de los morfismos de transporte, esto significa
sencillamente que \emph{el compuesto de dos morfismos de transporte es un
morfismo de transporte}. Un clivaje de una categoría $\cal{F}$
sobre~$\cal{E}$ que satisface la condición precedente se llama una
\emph{escisión}
\index{escision@escisión|hyperpage}%
de~$\cal{F}$ sobre~$\cal{E}$, y una categoría $\cal{F}$ sobre~$\cal{E}$
dotada de una escisión se llama una \emph{categoría escindida
sobre~$\cal{E}$}.
\index{escindida (categoria)@escindida (categoría)|hyperpage}%
Se trata, pues, de un caso particular de la noción de categoría con
clivaje. La categoría de las categorías escindidas sobre~$\cal{E}$ es,
por tanto, equivalente a~$\SheafHom(\cal{E}^\circ,\Cat)$. Obsérvese que
una categoría escindida sobre~$\cal{E}$ es, con mayor razón, una
categoría con clivaje sobre~$\cal{E}$.

Si $\cal{F}$ es una categoría fibrada sobre~$\cal{E}$, no siempre existe
una escisión de~$\cal{F}$. Supongamos, por ejemplo, que
$\Ob\cal{E}$ y $\Ob\cal{F}$ se reducen cada uno a un elemento y que los
endomorfismos de estos elementos forman grupos $E$ y $F$,
respectivamente, de modo que el funtor de proyección $p$ viene dado por un
homomorfismo de grupos $p\colon F\to E$, sobreyectivo porque $p$ es un
fibrado. Se comprueba entonces inmediatamente que el conjunto de los
clivajes de~$\cal{F}$ sobre~$\cal{E}$ está en correspondencia biunívoca
con el conjunto de las aplicaciones $s\colon E\to F$ tales que
$ps=\id_E$ (es decir, el conjunto de los «sistemas de representantes» de
las clases módulo el subgrupo $G$, núcleo del homomorfismo sobreyectivo
$p\colon F\to E$). Un clivaje es una escisión si y solo si $s$ es un
homomorfismo de grupos. Decir que existe una escisión significa, pues, que
la extensión de grupos $F$ de~$E$ por~$G$ es escindida, lo cual, cuando
$G$ es conmutativo, se expresa por la anulación de cierta clase de
cohomología en~$\H^2(E,G)$ (donde se considera $G$ como un grupo sobre el
que actúa~$E$).

Supongamos, sin embargo, que $\cal{F}$ es una categoría fibrada
sobre~$\cal{E}$ tal que las
\marginpar{181}
$\cal{F}_S$ son categorías \emph{rígidas}, es decir, el grupo de
automorfismos de todo objeto de~$\cal{F}_S$ se reduce a la identidad. Es
entonces fácil probar que $\cal{F}$ admite una escisión sobre~$\cal{E}$.
En efecto, se observa primero que la cuestión de la existencia de una
escisión no cambia si se sustituye $\cal{F}$ por una categoría
$\cal{E}$-equivalente, lo que nos reduce en este caso a aquel en que las
$\cal{F}_S$ son categorías rígidas \emph{y reducidas} (es decir, dos
objetos isomorfos de~$\cal{F}_S$ son idénticos). Pero si $G$ es una
categoría rígida y reducida, todo isomorfismo entre dos funtores $H\to G$
(donde $H$ es una categoría cualquiera) es una identidad. Se sigue que si
$\cal{F}$ es una categoría fibrada sobre~$\cal{E}$ cuyas categorías fibra
son rígidas y reducidas, entonces existe un clivaje \emph{único} de
$\cal{F}$ sobre~$\cal{E}$, que es necesariamente una escisión. Por tanto,
$\cal{F}$ es isomorfa a la categoría definida por un funtor
$\varphi\colon \cal{E}^\circ\to \Cat$ tal que las $\varphi(S)$ sean
categorías rígidas y
\ifthenelse{\boolean{orig}}
{discretas}
{reducidas}
, y el funtor $\varphi$ queda determinado salvo isomorfismo.

\section{Categorías cofibradas, categorías bifibradas}
\label{VI.10}
Consideremos una categoría $\cal{F}$ sobre~$\cal{E}$, con el funtor de
proyección
$$
p\colon \cal{F}\to \cal{E},
$$
que define una categoría $\cal{F}^\circ$ sobre $\cal{E}^\circ$ mediante el
funtor de proyección
$$
p^\circ\colon \cal{F}^\circ\to \cal{E}^\circ.
$$
Un morfismo $\alpha\colon \eta\to\xi$ de $\cal{F}$ se llama
\emph{cocartesiano} si es un morfismo cartesiano para $\cal{F}^\circ$
sobre~$\cal{E}^\circ$. Explicitando la definición, se ve que esto significa
que, para todo objeto $\xi'$ de~$\cal{F}_S$, la aplicación
$u\mto u\circ \alpha$
$$
\Hom_S(\xi,\xi')\to\Hom_f(\eta,\xi')
$$
es biyectiva. Se dice también entonces que $(\xi,\alpha)$ es una
\emph{imagen directa}
\index{imagen directa|hyperpage}%
de~$\eta$ por~$f$, en la categoría $\cal{F}$ sobre~$\cal{E}$. Si existe
para todo $\eta$ de~$\cal{F}_T$, se dice que existe el funtor de imagen
directa por~$f$, y se denota este funtor
\marginpar{182}
$f_*^{\cal{F}}$ o~$f_*$,
\label{indnot:ft}\oldindexnot{$f_*^{\cal{F}}$ o~$f_*$|hyperpage}%
una vez escogido. Queda, pues, definido por un isomorfismo de bifuntores
sobre~$\cal{F}_T^\circ\times\cal{F}_S$:
$$
\Hom_S(f_*(\eta),\xi)\isomto \Hom_f(\eta,\xi).
$$
Por consiguiente, si $f_*$ existe, para que $f^*$ exista es necesario y
suficiente que
\ifthenelse{\boolean{orig}}
{$f$}
{$f_*$}
admita un funtor adjunto, es decir, que exista un funtor
$f^*\colon \cal{F}_S\to\cal{F}_T$ y un isomorfismo de bifuntores
$$
\ifthenelse{\boolean{orig}}
{\Hom_E(f_*(\eta),\xi)\isomto \Hom_T(\eta,f^*(\xi)).}
{\Hom_S(f_*(\eta),\xi)\isomto \Hom_T(\eta,f^*(\xi)).}
$$
Sea $g\colon U\to T$ otro morfismo de $\cal{E}$, y supongamos que existen
las imágenes inversas y directas por $f,g$ y~$fg$. Consideremos entonces
las transformaciones naturales
$$
\begin{array}{cccc}
c^{f,g}\colon&f_*g_*&\dpl\from&(fg)_*\\
c_{f,g}\colon&g^*f^*&\dpl\to&(fg)^*.
\end{array}
$$
Se comprueba que, si se consideran $f_*g_*$ y $g^*f^*$ como un par de
funtores adjuntos, al igual que $(fg)_*$ y $(fg)^*$, las dos
transformaciones precedentes son adjuntas entre sí. Por tanto, una es un
isomorfismo si y solo si la otra lo es. En particular:
\begin{proposition}
\label{VI.10.1} Supongamos que la categoría $\cal{F}$ sobre~$\cal{E}$ es
prefibrada y coprefibrada. Para que sea fibrada, es necesario y suficiente
que sea cofibrada.
\end{proposition}

Por supuesto, se dice que $\cal{F}$ es coprefibrada, \resp cofibrada,
\index{cofibrada (categoría)|hyperpage}%
sobre~$\cal{E}$ si $\cal{F}^\circ$ es prefibrada, \resp fibrada%
\ifthenelse{\boolean{orig}}
{, sobre~$\cal{E}$.}
{, sobre~$\cal{E}^\circ$.}
Diremos que $\cal{F}$ es bifibrada sobre~$\cal{E}$,
\index{bifibrada (categoría)|hyperpage}%
si es a la vez fibrada y cofibrada sobre~$\cal{E}$.

\section{Ejemplos diversos}
\label{VI.11}

\begin{enumerate}

\item[a)] \textbf{Categorías de flechas de}~$\cal{E}$. Sea
$\cal{E}$ una categoría. Denotemos por $\mathbf{\Delta^1}$ la categoría
asociada al conjunto totalmente ordenado de dos elementos $[0,1]$;
\marginpar{183}
tiene, pues, dos objetos, 0 y 1, y, además de los dos morfismos identidad,
una flecha $(0,1)$ con fuente 0 y destino 1. Sea
$$
\CatFl(\cal{E})=\SheafHom(\mathbf{\Delta^1},\cal{E});
$$
se llama la \emph{categoría de flechas de}~$\cal{E}$. El objeto 1
de~$\mathbf{\Delta^1}$ define un funtor canónico, llamado
\emph{funtor codominio},
$$
\CatFl(\cal{E})\to \cal{E}
$$
(el funtor definido por el objeto 0 de~$\mathbf{\Delta^1}$ se llama
\emph{funtor dominio}). Para todo objeto $S$ de~$\cal{E}$, la categoría
fibra $\CatFl(\cal{E})_S$ es canónicamente isomorfa a la categoría
$\cal{E}_{/S}$ de los objetos de~$\cal{E}$ sobre~$S$.

Consideremos un morfismo $f\colon T\to S$ de $\cal{E}$. Le corresponde
entonces un funtor canónico
$$
f_*\colon \cal{E}_{/T}=\cal{F}_T\to \cal{E}_{/S}=\cal{F}_S
$$
y un isomorfismo natural
$$
\Hom_S(f_*(\eta),\xi)\isomto\Hom_f(\eta,\xi)
$$
que hace de~$f_*$ un funtor de imagen directa para $f$ en~$\cal{F}$.
Además, en este caso se tiene
$$
(\id_S)_*=\id_{\cal{F}_S},\quad (fg)_*=f_*g_*,\quad
\ifthenelse{\boolean{orig}}
{c^{f,g}=\id_{(fg)},}
{c^{f,g}=\id_{(fg)_*},}
$$
es decir, $\cal{F}$ está dotada de una coescisión sobre~$\cal{E}$. Con
mayor razón, $\cal{F}$ es cofibrada sobre~$\cal{E}$. Observemos ahora que
el conjunto de morfismos de $\cal{F}$ está en correspondencia biunívoca
con el conjunto de diagramas cuadrados conmutativos de~$\cal{E}$:
$$
\xymatrix{
X\ar[d]_-u&\ar[l]_-{f'}Y\ar[d]^-v\\
S&\ar[l]_-{f}T
}
$$
Por
\marginpar{184}
definición, el morfismo en cuestión es cartesiano si el cuadrado es
cartesiano en~$\cal{E}$, es decir, si hace de~$Y$ un producto fibrado de
$X$ y $T$ sobre~$S$. Por tanto, el funtor de imagen inversa
\ifthenelse{\boolean{orig}}
{$f$}
{$f^*$}
existe si y solo si para todo objeto $X$ sobre~$S$ existe el producto
fibrado $X\times_ST$. Se deduce de~\Ref{VI.10.1} que, si en~$\cal{E}$
existe siempre el producto de dos objetos sobre un tercero, es decir, si
$\cal{F}$ es prefibrada sobre~$\cal{E}$, entonces $\cal{F}$ es incluso
fibrada sobre~$\cal{E}$.

\medskip
\item[b)] \textbf{Categoría de prehaces o haces sobre espacios variables}

Sea $\cal{E}=\mathbf{Top}$ la categoría de los espacios topológicos. Si
$T$ es un espacio topológico, denotaremos por $\cal{U}(T)$ la categoría de
los abiertos de~$T$, cuyos morfismos son las aplicaciones de inclusión. Si
$\cal{C}$ es una categoría, un funtor
$\cal{U}(T)^\circ\to \cal{C}$ se llama un \emph{prehaz} sobre~$T$ con
valores en~$\cal{C}$, y un \emph{haz} si satisface una condición de
exactitud por la izquierda que no repetiremos aquí. La \emph{categoría
$\cal{P}(T)$ de prehaces sobre~$T$ con valores en~$\cal{C}$} es, por
definición, la categoría $\SheafHom(\cal{U}(T)^\circ,\cal{C})$, y la
categoría $\cal{F}(T)$ de haces sobre~$T$ con valores en~$\cal{C}$ es la
subcategoría plena cuyos objetos son los objetos de
$\SheafHom(\cal{U}(T)^\circ,\cal{C})$ que son haces. Si
$f\colon T\to S$ es un morfismo de~$\cal{E}$, es decir, una aplicación
continua de espacios topológicos, la aplicación creciente
$U\mto f^{-1}(U)$ le hace corresponder un funtor
$\cal{U}(S)\to\cal{U}(T)$ y, por tanto, un funtor
$$
f_*\colon\SheafHom(\cal{U}(T)^\circ,\cal{C})\to
\SheafHom(\cal{U}(S)^\circ,\cal{C}),
$$
llamado \emph{funtor de imagen directa de prehaces por}~$f$. Se ve de
inmediato que la imagen directa de un haz es un haz; por tanto, el funtor
\ifthenelse{\boolean{orig}}
{$f_*\cal{P}(T)\to\cal{P}(S)$}
{$f_*\colon \cal{P}(T)\to\cal{P}(S)$}
induce un funtor, denotado también
\ifthenelse{\boolean{orig}}
{$f_*\colon\cal{P}(T)\to\cal{P}(S)$.}
{$f_*\colon\cal{F}(T)\to\cal{F}(S)$.}
Además, se comprueba trivialmente, por la asociatividad de la composición
de funtores, que para una segunda aplicación continua $g\colon U\to T$ se
tiene la identidad
$$
\ifthenelse{\boolean{orig}}
{(gf)_* =g_*f_*,}
{(fg)_*=f_*g_*,}
\qquad
\text{y asimismo}\quad (\id_S)_*=\id_{\cal{P}(S)}.
$$
De esta forma se ha obtenido un funtor
$$
S\mto \cal{P}(S)
$$
\resp $$
S\mto \cal{F}(S)
$$
de
\marginpar{185}
$\cal{E}$ en~$\Cat$. En realidad, nos interesa el funtor correspondiente
$$
S\mto \cal{P}(S)^\circ,\qquad\resp\qquad S\mto
\cal{F}(S)^\circ.
$$
Define una categoría cofibrada, e incluso coescindida, sobre la categoría
de los espacios topológicos, que se llama la \emph{categoría cofibrada de
los prehaces} (\resp \emph{haces}) \emph{con valores en}~$\cal{C}$
(se sobreentiende: sobre espacios variables). Explicitando la construcción
del \No~\Ref{VI.8}, se ve que un morfismo de un prehaz $B$ sobre~$T$ en un
prehaz $A$ sobre~$S$ es un par $(f,u)$ formado por una aplicación continua
de~$T$ en~$S$ y un morfismo $u\colon A\to f_*(B)$ en la categoría
$\cal{P}(S)$. Esta descripción vale también para los morfismos de haces,
pues $\cal{F}$ es una subcategoría plena de~$\cal{P}$.

En los casos más importantes, la categoría $\cal{P}$ y la categoría
$\cal{F}$ sobre~$\cal{E}$ son también categorías fibradas, es decir, para
toda aplicación continua los funtores de imagen directa
$\cal{P}(T)\to\cal{P}(S)$ y $\cal{F}(T)\to\cal{F}(S)$ tienen un funtor
adjunto, que se denota entonces por $f^*$ y se llama funtor de imagen
inversa de prehaces \resp funtor de imagen inversa de haces por la
aplicación continua~$f$. Este funtor existe, por ejemplo, si
$\cal{C}=\Ens$. Se puede demostrar que el funtor
$f^*\colon\cal{P}(S)\to\cal{P}(T)$ existe siempre que en~$\cal{C}$
existan los límites inductivos relativos a diagramas del universo
considerado. La cuestión es menos fácil para $\cal{F}$; en efecto, incluso
en el caso $\cal{C}=\Ens$, la imagen inversa de un prehaz que es un haz no
es, en general, un haz. En otras palabras, el funtor de imagen inversa de
haces no es isomorfo al funtor inducido por el funtor de imagen inversa de
prehaces, pese a la notación común $f^*$. Así, $\cal{F}$ es una
subcategoría cofibrada de~$\cal{P}$, pero no una subcategoría fibrada, es
decir, \emph{el funtor de inclusión $\cal{F}\to\cal{P}$ no es fibrante}.

La categoría cofibrada $\cal{P}$ puede deducirse de una categoría
cofibrada, o más bien fibrada, más general, obtenida como sigue. Para toda
categoría $\cal{U}$ del universo fijado, pongamos
$$
\cal{P}(\cal{U})=\SheafHom(\cal{U},\cal{C})
$$
y
\marginpar{186}
observemos que $\cal{U}\mto\cal{P}(\cal{U})$ es de manera natural un
funtor contravariante en~$\cal{U}$, de la categoría $\Cat$ en~$\Cat$.
Define, pues, una categoría escindida sobre~$\cal{E}=\Cat$, que denotaremos
por $\Cat_{/\!/\cal{C}}$. Los objetos de esta categoría son los pares
$(\cal{U},p)$ formados por una categoría $\cal{U}$ y un funtor
$p\colon\cal{U}\to\cal{C}$, y un morfismo de $(\cal{U},p)$ en
$(\cal{V},q)$ es esencialmente un par $(f,u)$, donde $f$ es un funtor
$\cal{U}\to\cal{V}$ y $u$ una transformación natural de funtores
$u\colon p\to qf$. Dejamos al lector la tarea de explicitar la composición
de los morfismos en $\Cat_{/\!/\cal{C}}$. El funtor de proyección
$$
\cal{F}=\Cat_{/\!/\cal{C}}\to\cal{E}=\Cat
$$
asocia al par $(\cal{U},p)$ el objeto $\cal{U}$; la categoría fibra en
$\cal{U}$ es la categoría $\SheafHom(\cal{U},\cal{C})$, salvo isomorfismo.
Cuando existen en~$\cal{C}$ los límites inductivos, se demuestra fácilmente
que la categoría fibrada $\Cat_{/\!/\cal{C}}$ sobre~$\Cat$ es también
cofibrada sobre~$\Cat$, es decir, se puede definir la noción de
\emph{imagen directa de un funtor} $p\colon\cal{U}\to\cal{C}$ por un
funtor $f\colon\cal{U}\to\cal{V}$. La categoría de prehaces se deduce de
la categoría fibrada precedente mediante el cambio de base
$$
\mathbf{Top}^\circ\to\Cat
$$
(el funtor $S\mto\cal{U}(S)$ definido más arriba), lo que da una categoría
fibrada sobre~$\mathbf{Top}^\circ$; al pasar a la categoría opuesta se
obtiene la categoría cofibrada $\cal{P}$ de los prehaces sobre
$\mathbf{Top}$. La noción de imagen inversa de un funtor corresponde a la
de imagen directa de un prehaz, y la noción de imagen directa de un funtor,
a la de imagen inversa de un prehaz.

\medskip
\item[c)] \textbf{Objetos con operadores sobre un objeto con operadores}

Sea $\cal{F}$ una categoría sobre~$\cal{E}$, y sea $S$ un objeto de
$\cal{E}$ sobre el que actúa un grupo $G$, por la izquierda para fijar
ideas. Este objeto con operadores puede interpretarse como correspondiente
a un funtor $\lambda\colon\cal{E}'\to\cal{E}$ desde la categoría (con un
único objeto y con $G$ como grupo de endomorfismos) $\cal{E}'$, definida
por~$G$, a la categoría $\cal{E}$. Por cambio de base define,
por tanto, una categoría $\cal{F}'$ sobre~$\cal{E}'$, que es fibrada
\resp cofibrada cuando $\cal{F}$ lo es sobre~$\cal{E}$. Una
\marginpar{187}
\ifthenelse{\boolean{orig}}
{sección de~$\cal{E}'$ sobre~$\cal{F}'$}
{sección de~$\cal{F}'$ sobre~$\cal{E}'$}
(necesariamente cartesiana, pues $\cal{E}'$ es un grupoide y todo
isomorfismo de~$\cal{F}'$ es cartesiano en virtud de~\Ref{VI.6.12}) puede
interpretarse también como un $\cal{E}$-funtor
$\cal{E}'\to\cal{F}$ sobre~$\lambda$, o como un objeto con operadores
$\xi$ en~$\cal{F}$ «sobre» el objeto con operadores~$S$.

\medskip
\item[d)] \textbf{Pares de funtores adjuntos cuasi-inversos;
autodualidades}

Cuando la categoría base $\cal{E}$ se reduce a dos objetos $a,b$ y, además
de las flechas identidad, a dos isomorfismos $f\colon a\to b$ y
$g\colon b\to a$ inversos entre sí, es decir, cuando $\cal{E}$ es un
grupoide conexo rígido con dos objetos, una categoría con clivaje
normalizado sobre~$\cal{E}$ es esencialmente lo mismo que el sistema
formado por dos categorías $\cal{F}_a$ y $\cal{F}_b$ y un \emph{par de
funtores adjuntos} $G\colon\cal{F}_a\to\cal{F}_b$ y
$F\colon\cal{F}_b\to\cal{F}_a$ que sean equivalencias de categorías y,
por tanto, cuasi-inversos entre sí. Se toman como $\cal{F}_a$ y
$\cal{F}_b$ las categorías fibra de~$\cal{F}$, como $F$ y $G$ los funtores
$f^*$ y $g^*$, y los dos isomorfismos
$$
u\colon FG\isomto\id_{\cal{F}_a}\qquad v\colon
GF\isomto\id_{\cal{F}_b}
$$
son $c_{g,f}$ y $c_{f,g}$. Las dos condiciones usuales de compatibilidad
entre $u$ y $v$ no son sino la condición \Ref{VI.7.4}~$B)$ para los
compuestos $fgf$ y $gfg$. Es fácil demostrar que estas condiciones bastan
para implicar que se tiene efectivamente un pseudofuntor
$\cal{E}^\circ\to\Cat$.

Un caso interesante es aquel en que
$$
\cal{F}_b=\cal{F}_a^\circ,\quad G=F^\circ,\quad v=u^\circ.
$$
Se llama \emph{autodualidad}
\index{autodualidad|hyperpage}%
en una categoría $\cal{C}$ al dato de un funtor
$D\colon\cal{C}\to\cal{C}^\circ$ y de un isomorfismo
\ifthenelse{\boolean{orig}}
{$u\colon DD^\circ\isomto\id_{\cal{C}}$,}
{$u\colon D^\circ D\isomto\id_{\cal{C}}$,}
tales que $u$ y el isomorfismo
\ifthenelse{\boolean{orig}}
{$u^\circ\colon D^\circ D\isomto\id_{\cal{C}^\circ}$}
{$u^\circ\colon DD^\circ\isomto\id_{\cal{C}^\circ}$}
hagan de $(D,D^\circ)$ un par de funtores adjuntos, necesariamente
cuasi-inversos entre sí. Esta condición se escribe
$$
D(u(x))=u(D(x))\qquad \text{para todo $x\in\Ob(\cal{C})$.}
$$

\item[e)]
\marginpar{188}
\textbf{Categorías sobre una categoría discreta $\cal{E}$.} Se dice que
$\cal{E}$ es una \emph{categoría discreta} si toda flecha de ella es una
flecha identidad, de modo que $\cal{E}$ queda definida, salvo isomorfismo
único, por el conjunto $I=\Ob(\cal{E})$. El dato de una categoría
$\cal{F}$ sobre~$\cal{E}$ equivale, pues, salvo isomorfismo único, al dato
de una familia de categorías $\cal{F}_i$ ($i\in I$), las categorías fibra.
Toda categoría $\cal{F}$ sobre~$\cal{E}$ es fibrada, todo $\cal{E}$-funtor
$\cal{F}\to\cal{G}$ es cartesiano y se tiene un isomorfismo canónico
$$
\SheafHom_{\cal{E}/-}(\cal{F},\cal{G}) \isomto
\underset{i}\prod\SheafHom(\cal{F}_i,\cal{G}_i).
$$
En particular, se obtiene
$$
\bf{\Gamma}(\cal{F}/\cal{E})=\varprojLim\cal{F}/\cal{E}\isomto
\ifthenelse{\boolean{orig}}
{\underset{i}\prod \cal{E}_i.}
{\underset{i}\prod \cal{F}_i.}
$$

\item[f)] \textbf{Supongamos que $\cal{E}$ tiene exactamente dos objetos
$S$ y $T$,} y, además de los \textbf{morfismos identidad, un morfismo}
$f\colon T\to S$. Entonces una categoría $\cal{F}$ sobre~$\cal{E}$ queda
definida, salvo $\cal{E}$-isomorfismo único, por el dato de dos categorías
$\cal{F}_S$ y $\cal{F}_T$ y de un bifuntor $H(\eta,\xi)$ sobre
$\cal{F}_T^\circ\times\cal{F}_S$ con valores en~$\Ens$. En efecto, si
$\cal{F}$ es una categoría sobre~$\cal{E}$, se le asocian las dos
categorías fibra $\cal{F}_S$ y $\cal{F}_T$ y el bifuntor
$H(\eta,\xi)=\Hom_f(\eta,\xi)$. Dejamos al lector la tarea de explicitar la
construcción en sentido inverso. Para que la categoría considerada sea
fibrada, o prefibrada, lo que equivale a lo mismo, es necesario y
suficiente que el funtor $H$ sea representable respecto del argumento
$\xi$. Para que sea cofibrada, es necesario y suficiente que $H$ sea
representable respecto del argumento~$\eta$.

\item[g)] Sea $\cal{F}=\cal{C}\times\cal{E}$, considerada como categoría
sobre~$\cal{E}$ mediante~$\pr_2$. Entonces $\cal{F}$ es fibrada y
cofibrada sobre~$\cal{E}$, e incluso está dotada de una escisión y una
coescisión canónicas, correspondientes al funtor constante sobre
$\cal{E}$, \resp sobre~$\cal{E}^\circ$, con valores en~$\Cat$ y valor
$\cal{C}$. Se tiene
$$
\bf{\Gamma}(\cal{F}/\cal{E})\simeq\SheafHom(\cal{E},\cal{C})
$$
y
\marginpar{189}
$\varprojLim\cal{F}/\cal{E}$ corresponde a la subcategoría plena formada
por los funtores $F\colon\cal{E\to\cal{C}}$ que transforman morfismos
cualesquiera en isomorfismos.
\end{enumerate}

\section{Funtores sobre una categoría con clivaje}
\label{VI.12}

Sea $\cal{F}$ una categoría con clivaje normalizado sobre~$\cal{E}$.
Para todo objeto $S$ de~$\cal{E}$, denotemos por
$$
i_S\colon \cal{F}_S\to\cal{F}
$$
el funtor de inclusión. Se tiene, pues, una transformación natural, para
todo morfismo $f\colon T\to S$ de~$\cal{E}$:
$$
\alpha_f\colon i_T f^*\to i_S,
$$
donde $f^*$ es el funtor de cambio de base
$\cal{F}_S\to\cal{F}_T$ para $f$ definido por el clivaje. Sea ahora
$$
F\colon\cal{F}\to\cal{C}
$$
un funtor de~$\cal{F}$ en una categoría $\cal{C}$. Pongamos, para todo
$S\in\Ob(\cal{E})$,
$$
F_S=F\circ i_S\colon\cal{F}_S\to\cal{C}
$$
y, para todo $f\colon T\to S$ de~$\cal{E}$,
$$
\varphi_f=F*\alpha_f\colon F_T f^*\to F_S.
$$
Así, a todo funtor $F\colon\cal{F}\to\cal{C}$ se le asocian una familia
$(F_S)$ de funtores $\cal{F}_S\to\cal{C}$ y una familia $(\varphi_f)$ de
transformaciones naturales $F_T f^*\to F_S$. Estas familias satisfacen
las condiciones siguientes:

a) $\varphi_{\id_S}=\id_{F_S}$.

b) Para dos morfismos $f\colon T\to S$ y $g\colon U\to T$ de~$\cal{E}$,
es conmutativo el cuadrado de transformaciones naturales
$$
\xymatrix@C=2cm{ F_U g^* f^* \ar[r]^-{F_U*c_{f,g}} \ar[d]_{\varphi_g * f^*} &
F_U(fg)^* \ar[d]^{\varphi_{fg}} \\ F_T f^*\ar[r]^-{\varphi_f} & F_S.}
$$
\marginpar{190}%
La primera relación es inmediata, y la segunda se obtiene aplicando el
funtor~$F$ al diagrama conmutativo
$$
\xymatrix@C=2cm{ g^* f^* (\xi) \ar[r]^-{c_{f,g}(\xi)}
\ar[d]_{\alpha_g(f^*(\xi))} & (fg)^*(\xi) \ar[d]^{\alpha_{fg}(\xi)}\\
f^*(\xi) \ar[r]^-{\alpha_f(\xi)}& \xi }
$$
para un objeto variable $\xi$ de~$\cal{F}_S$.

Si $G$ es un segundo funtor $\cal{F}\to\cal{C}$, que da lugar a funtores
$G_S\colon\cal{F}_S\to\cal{C}$ y a transformaciones naturales
\ifthenelse{\boolean{orig}}
{$\psi_f\colon G_Tf\to G_S$}
{$\psi_f\colon G_Tf^*\to G_S$},
y si $u\colon F\to G$ es una transformación natural, le corresponden las
transformaciones naturales $u*i_S$:
$$
u_S\colon F_S\to G_S,
$$
y se comprueba de inmediato que, para todo morfismo $f\colon T\to S$ de
$\cal{E}$, son conmutativos los cuadrados
$$
\begin{array}{c}
\xymatrix@C=1.5cm{ F_Tf^* \ar[r]^-
{\ifthenelse{\boolean{orig}}{\phi_f}{\varphi_f}}
\ar[d]_{u_T * f^*} & F_S
\ar[d]^{u_S} \\ G_T f^* \ar[r]^-{\psi_f}& G_S}
\end{array}
\leqno\quad \textup{c)}
$$

\begin{proposition}
\label{VI.12.1}
Sea $\cal{H}(\cal{F},\cal{C})$ la categoría cuyos objetos son los pares
formados por familias $(F_S)$
\ifthenelse{\boolean{orig}}
{($S\in\Ob(\cal{F})$)}
{($S\in\Ob(\cal{E})$)}
de funtores $\cal{F}_S\to\cal{C}$ y familias $(\varphi_f)$
\ifthenelse{\boolean{orig}}
{($f\in\Fl(\cal{F})$)}
{($f\in\Fl(\cal{E})$)}
de transformaciones naturales $F_Tf^*\to F_S$ que satisfacen las
condiciones \emph{a) y~b)}, y cuyos morfismos son las familias $(u_S)$
\ifthenelse{\boolean{orig}}
{($S\in\Ob(\cal{F})$)}
{($S\in\Ob(\cal{E})$)}
de transformaciones naturales
\marginpar{191}
$F_S\to G_S$ que verifican la condición de conmutatividad \emph{c)}
escrita más arriba (la composición de los morfismos se efectúa componiendo
las transformaciones naturales entre funtores $\cal{F}_S\to\cal{C}$).
Entonces las dos reglas explicitadas más arriba definen un
\emph{isomorfismo} $K$ de la categoría
$\SheafHom(\cal{F},\cal{C})$ con la categoría
$\cal{H}(\cal{F},\cal{C})$.
\end{proposition}

Es inmediato que se tiene así un \emph{funtor} de la primera categoría en
la segunda. Este funtor es plenamente fiel, pues, dados $F,G$, la
aplicación $\Hom(F,G)\to\allowbreak\Hom(K(F),K(G))$ es manifiestamente
inyectiva. Para demostrar que es sobreyectiva, basta observar que la
condición de conmutatividad c) expresa la functorialidad de las
aplicaciones
$u(\xi)=u_S(\xi)\colon
\ifthenelse{\boolean{orig}}
{F(S)}
{F(\xi)}
=F_S(\xi)\to G(\xi)=G_S(\xi)$ para los morfismos de la forma
$\alpha_f(\xi)$ en~$\cal{F}$. Por otra parte, se tiene la functorialidad
en cada categoría fibra, es decir, para los morfismos de~$\cal{F}$ que
son $T$-morfismos ($T\in\Ob(\cal{E})$), de donde resulta la functorialidad
para todo morfismo de~$\cal{F}$,
\ifthenelse{\boolean{orig}}
{puesto que un $f$-morfismo (donde $f\colon T\to S$ es un morfismo en
($\cal{E}$) es de manera única) un compuesto de un morfismo
$\alpha_f(\xi)$ y de un $T$-morfismo.}
{puesto que todo $f$-morfismo (donde $f\colon T\to S$ es un morfismo en
$\cal{E}$) es, de manera única, un compuesto de un morfismo
$\alpha_f(\xi)$ y de un $T$-morfismo.}
Queda por demostrar que el funtor~$K$ es biyectivo sobre los objetos. El
argumento precedente muestra ya que $K$ es inyectivo sobre los objetos;
queda probar que es sobreyectivo, es decir, que, si se parte de un sistema
$(F_S),(\varphi_f)$ que satisface a) y b), y se define una aplicación
$\Ob\cal{F}\to\Ob\cal{C}$ por
$$
F(\xi)=F_S(\xi) \quad\text{para }
\xi\in\Ob\cal{F}_S\subset\Ob\cal{F}
$$
y una aplicación $\Fl(\cal{F})\to\Fl(\cal{C})$ por
$$
F(\alpha_f(\xi)u')=\varphi_f(\xi)\,F_T(u'),
$$
para todo morfismo $f\colon T\to S$ de~$\cal{E}$, todo objeto $\xi$
de~$\cal{F}_S$ y todo $T$-morfismo $u'$ con destino $f^*(\xi)$, se
obtiene un \emph{funtor} $F$ de~$\cal{F}$ en~$\cal{C}$. En efecto, la
relación $F(\id_\xi)=\id_{F(\xi)}$ es inmediata; queda demostrar la
multiplicatividad $F(uv)=F(u)F(v)$ cuando se tiene un $f$-morfismo
$u\colon\eta\to\xi$ y un $g$-morfismo
\ifthenelse{\boolean{orig}}
{$v\colon\zeta\to\nu$}
{$v\colon\zeta\to\eta$},
con $f\colon T\to S$ y $g\colon U\to T$ morfismos de~$\cal{E}$.
Poniendo $w=uv$, se tendrá
$$
u=\alpha_f(\xi)u'\quoi,\quad v=\alpha_g(\eta)v'\quoi,\quad
w=\alpha_{fg}(\xi)w'.
$$
\marginpar{192}%
con
$$
w'=c_{f,g}(\xi)g^*(u')v'\quad \hbox{(\cf \No \Ref{VI.8})}.
$$

Con esta notación, hay que demostrar la conmutatividad del contorno
exterior del diagrama siguiente:
$$
\xymatrix{ F_U(\zeta) \ar@/^1.5pc/[rrrrrr]^{F(w')}
\ar[rr]^{\ifthenelse{\boolean{orig}}{F_u(v')}{F_U(v')}}
\ar[drr]_{F(v)} & & F_Ug^*(\eta)
\ar[rr]^{\ifthenelse{\boolean{orig}}{F_ug^*(u')}{F_Ug^*(u')}}
\ar[d]^{\varphi_g(\eta)} & & F_Ug^*f^*(\xi)
\ar[rr]^{F_U(c_{f,g}(\xi))}\ar[d]^{\varphi_g(f^*(\xi))}& &
F_U(fg)^*(\xi) \ar[d]^{\varphi_{fg}(\xi)} \\ & & F_T(\eta)
\ar@/_1.2pc/[rrrr]_{F(u)} \ar[rr]_{F_T(u')} & & F_Tf^*(\xi)
\ar[rr]_{\varphi_f(\xi)} & & F_S(\xi) }
$$
Ahora bien, el triángulo izquierdo es conmutativo por definición de
$F(v)$; el cuadrado central es conmutativo porque se deduce del morfismo
\ifthenelse{\boolean{orig}}
{$u'\colon \xi\to f(\eta)$}
{$u'\colon \eta\to f^*(\xi)$}
mediante la transformación natural $\alpha_g$;
por último, el cuadrado derecho es conmutativo en virtud de la condición
b). De ello resulta la conclusión buscada.

Supongamos ahora que $\cal{C}$ es también una categoría con clivaje
normalizado sobre~$\cal{E}$, que en adelante llamaremos $\cal{G}$, y que
nos interesan los $\cal{E}$-funtores de~$\cal{F}$ en~$\cal{G}$. Si $F$ es
uno de tales funtores, induce funtores
$$
F_S\colon \cal{F}_S\to\cal{G}_S
$$
entre las categorías fibra. Por otra parte, para todo morfismo
$f\colon T\to S$ de~$\cal{E}$ y todo objeto $\xi$ de~$\cal{F}_S$, el
$f$-morfismo $F(\alpha_f(\xi))$ se factoriza de manera única por un
$T$-morfismo
$$
\varphi_f(\xi)\colon F_T(f^*_{\cal{F}}(\xi))\to
f^*_{\cal{G}}(F_S(\xi))
$$
(donde el subíndice $\cal{F}$ o $\cal{G}$ indica la categoría con
clivaje para la que se toma el funtor de imagen inversa), de donde se
obtiene una transformación natural entre funtores de~$\cal{F}_S$ en
$\cal{G}_T$:
$$
\varphi_f\colon F_T f^*_{\cal{F}}\to f^*_{\cal{G}}F_S.
$$
\marginpar{193}%
Los dos sistemas $(F_S)$ y $(\varphi_f)$ satisfacen las condiciones
siguientes:

a') $\varphi_{\id_S}=\id_{F_S}$.

b') Para dos morfismos $f\colon T\to S$ y $g\colon U\to T$ de
$\cal{E}$, es conmutativo el diagrama siguiente de transformaciones
naturales:
$$
\xymatrix{ F_U
\ifthenelse{\boolean{orig}}{g_F^*f_F^*}{g_{\cal{F}}^*f_{\cal{F}}^*}
\ar[rr]^{F_U *c_{f,g}^\cal{F}}
\ar[d]_{\ifthenelse{\boolean{orig}}{\varphi_{g^*}f^*_{\cal{F}}}{\varphi_g*f^*_{\cal{F}}}} & &
F_U\ifthenelse{\boolean{orig}}{(fg)^*_F}{(fg)^*_{\cal{F}}}
\ar[dd]^{\varphi_{fg}} \\ g^*_{\cal{G}}F_T
\ifthenelse{\boolean{orig}}{f^*_F}{f^*_{\cal{F}}}
\ar[d]_{g^*_{\cal{G}}*\varphi_f} & & \\ g^*_{\cal{G}}f^*_{\cal{G}}F_S
\ar[rr]^{c_{f,g}^\cal{G}*F_S} & & (fg)^*_{\cal{G}}F_S. }
$$
Dejamos al lector su verificación, así como el enunciado y la demostración
del análogo de la Proposición~\Ref{VI.12.1}; este implica que se obtiene
así una correspondencia biunívoca entre el conjunto de los
$\cal{E}$-funtores de~$\cal{F}$ en~$\cal{G}$ y el conjunto de los
sistemas $(F_S)$,$(\varphi_f)$ que satisfacen las condiciones a') y b')
anteriores. Por supuesto, en esta correspondencia, los funtores
cartesianos se caracterizan por la propiedad de que las transformaciones
$\varphi_f$ son isomorfismos.

\begin{remarquestar}
Por supuesto, casi siempre conviene razonar directamente con categorías
fibradas sin utilizar clivajes explícitos, lo cual evita, en particular,
tener que recurrir, para la noción sencilla de $\cal{E}$-funtor o de
$\cal{E}$-funtor cartesiano, a una interpretación pesada como la anterior.
Para evitar pesadeces insoportables y obtener enunciados más intrínsecos,
\marginpar{194}
hemos tenido que renunciar a partir, como en~\cite{VI.2}, de la noción de
categoría con clivaje (llamada «categoría fibrada» en \loccit), que pasa a
segundo plano en favor de la de categoría fibrada. Es, por lo demás,
probable que, contrariamente al uso todavía predominante, ligado a antiguos
hábitos de pensamiento, acabe por resultar más cómodo, en los problemas
universales, no poner el acento en \emph{una} solución supuestamente elegida
de una vez por todas, sino situar todas las soluciones en pie de igualdad.
\end{remarquestar}

\begin{sgabibliography}{0}{VI.13}
\bibitem[1]{VI.1} A\ptbl \textsc{Grothendieck}, Sur quelques points d'algèbre
homologique,
\ifthenelse{\boolean{orig}}
{Tohoku Math.\ Journal, pp\ptbl 119--221, Vol\ptbl 9, 1957.}
{T\^ohoku Math.\ J. \textbf{9} (1957) p\ptbl 119--221.}

\bibitem[2]{VI.2} A\ptbl \textsc{Grothendieck}, Technique de descente et
Théorèmes d'existence,~I, Séminaire Bourbaki~190, diciembre de~1959.
\end{sgabibliography}


\refstepcounter{chapter}
\addtocontents{toc}{\protect\medskip\protect\par\protect\noindent\protect\textbf{VII: no existe}}

\chapterspace{-2}
\chapter{Descenso fielmente plano}
\label{VIII}
\marginpar{195}

\section{Descenso de módulos cuasicoherentes}
\label{VIII.1}

Sea $\Sch$
\label{indnot:hb}\oldindexnot{$\Sch$|hyperpage}%
la categoría de los preesquemas. Procediendo como en~\Ref{VI.11}.b,
se obtiene que la categoría de los pares~$(X,F)$ formados por un
preesquema~$X$ y un módulo~$F$ sobre~$X$ (donde los morfismos se definen,
como en \loccit, mediante la noción de imagen directa de un módulo por un
morfismo de espacios anillados) puede considerarse como una categoría
fibrada sobre~$\Sch$; el funtor de cambio de base relativo a un morfismo
$f\colon X\to Y$ de~$\Sch$ es el funtor de imagen inversa de módulos
por~$f$. (Obsérvese que la categoría fibra en~$X\in\Ob(\Sch)$ de la
categoría fibrada anterior es la categoría \emph{opuesta} a la categoría
de los módulos sobre~$X$). Como la imagen inversa de un módulo
cuasicoherente es cuasicoherente, se ve que la subcategoría plena de la
categoría de los pares~$(X,F)$ formada por aquellos para los que $F$ es
cuasicoherente es una subcategoría fibrada de la categoría fibrada
anterior. (En cambio, si no se imponen hipótesis sobre~$f$, la imagen
directa de un módulo cuasicoherente no es, en general, un módulo
cuasicoherente). Llamaremos simplemente a esta categoría fibrada la
\emph{categoría fibrada de los módulos cuasicoherentes sobre los
preesquemas}.

Recordemos, por otra parte, que un morfismo~$f\colon X\to Y$ de espacios
anillados se llama \emph{fielmente plano}
\index{fielmente plano (morfismo)|hyperpage}%
\index{plano (morfismo fielmente)|hyperpage}%
si es \emph{plano} (\ie si, para todo~$x\in X$,
$\mathcal{O}_{X,x}$ es un módulo plano sobre~$\mathcal{O}_{Y,f(x)}$,
\eqref{IV}) y \emph{sobreyectivo}. Se dice que $f$ es un morfismo
\emph{cuasicompacto} si la imagen inversa por~$f$ de toda parte
cuasicompacta es cuasicompacta; cuando $f$ es un morfismo de
preesquemas, esto significa también que la imagen inversa por~$f$ de un
abierto afín de~$Y$ es reunión \emph{finita} de abiertos afines de~$X$.

\begin{theoreme}
\label{VIII.1.1}
Sea
\marginpar{196}
$\cal{F}$ la categoría fibrada de los módulos cuasicoherentes sobre los
preesquemas. Sea
\ifthenelse{\boolean{orig}}
{$g\colon S''\to S$}
{$g\colon S'\to S$}
un morfismo de preesquemas fielmente plano y cuasicompacto. Entonces $g$
es un morfismo de descenso efectivo para~$\cal{F}$.
\index{descenso efectivo (morfismo de)|hyperpage}%
\end{theoreme}

Recordemos\footnote{Admitiremos aquí la teoría general del descenso
\ifthenelse{\boolean{orig}}
{expuesta expuesta con detalle}
{expuesta con detalle}
en el artículo de J\ptbl\ \textsc{Giraud} citado en la nota
\eqref{footnotegiraud} al pie de la página del Aviso, trabajo que en lo
sucesivo citaremos como \cite{VIII.D}. \Cf también \cite{VIII.2} para una
exposición sucinta.} que esto significa dos cosas:

\begin{corollaire}[Descenso de homomorfismos de módulos]
\label{VIII.1.2}
\index{descenso de homomorfismos de módulos|hyperpage}%
Sean $g\colon S'\to S$ un morfismo de preesquemas fielmente plano y
cuasicompacto, $F$ y $G$ dos módulos cuasicoherentes sobre~$S$, $F'$ y
$G'$ sus imágenes inversas sobre~$S'$, y, finalmente, $F''$ y $G''$ sus
imágenes inversas sobre~$S''=S'\times_{S}S'$. Consideremos el diagrama de
aplicaciones de conjuntos definido por los funtores de cambio de base
por~$g$, $p_{1}$ y $p_{2}$ (donde $\xymatrix@C=.5cm{p_{1},p_{2}\colon
S'\times_{S}S'\ar@/^2pt/[r]\ar@/_2pt/[r]&S'}$ son las dos proyecciones):
$$
\xymatrix@C=.5cm{\Hom_{S}(F,G)\ar[r] &
\Hom_{S'}(F',G')\ar@/^2pt/[r]\ar@/_2pt/[r] &
\Hom_{S''}(F'',G'').}
$$
Este diagrama es exacto, \ie identifica biyectivamente el primer conjunto
con el igualador de las dos aplicaciones escritas del segundo conjunto
en el tercero.
\end{corollaire}

En otros términos, el funtor de cambio de base
\ifthenelse{\boolean{orig}}
{por~$f$,}
{por~$g$,}
$F\mto F'$, define un funtor \emph{plenamente fiel} de la categoría de
los módulos cuasicoherentes sobre~$S$ en la categoría de los módulos
cuasicoherentes sobre~$S'$ provistos de un dato de descenso relativo
\ifthenelse{\boolean{orig}}
{a~$f$.}
{a~$g$.}
Además:

\begin{corollaire}[Descenso de módulos]
\label{VIII.1.3}
\index{descenso de módulos|hyperpage}%
Para todo módulo cuasicoherente~$F'$ sobre~$S'$, todo dato de descenso
sobre~$F'$ relativo a~$g$ es \emph{efectivo},
\index{dato de descenso efectivo|hyperpage}%
\ie $F'$, con su dato de descenso, es isomorfo a la imagen inversa por~$g$
de un módulo cuasicoherente sobre~$S$ (determinado salvo isomorfismo único
en virtud de~\Ref{VIII.1.2}).
\end{corollaire}

En otros términos, el funtor plenamente fiel anterior es incluso una
\emph{equivalencia}. En la práctica, esto significa que dar un módulo
cuasicoherente sobre~$S$ equivale a dar un módulo cuasicoherente sobre
$S'$ provisto de un dato de descenso relativo a~$g$.

\subsubsection*{Demostración de~\Ref{VIII.1.1}}
Sea
\marginpar{197}
primero $T$ un $S$-preesquema $S$-isomorfo a la suma disjunta de una
familia
\ifthenelse{\boolean{orig}}
{de abiertos}
{de abiertos}
inducidos~$S_{i}$ de~$S$ que recubren~$S$. Es entonces evidente que el
morfismo estructural~$T\to S$ es un morfismo de descenso efectivo
para~$\cal{F}$ (esto significa precisamente que dar un módulo
cuasicoherente~$F$ sobre~$S$ equivale a dar módulos cuasicoherentes~$F_{i}$
sobre los~$S_{i}$ e isomorfismos de pegado~$\varphi_{ji}\colon
F_{i}|S_{i}\cap S_{j}\to F_{j}|S_{i}\cap S_{j}$ que satisfacen la
\ifthenelse{\boolean{orig}}
{conocida condición de cocadenas}
{conocida condición de cociclo}.
En virtud de~VII,~8
se sigue que, para comprobar que $g\colon S'\to S$ es un morfismo de
descenso efectivo para~$\cal{F}$, basta comprobarlo para el morfismo
$g_{T}\colon T'=T\times_{S}S'\to T$ deducido de~$g$ mediante el cambio de
base~$T\to S$. (Obsérvese que la hipótesis sobre~$T\to S$ permanece
estable bajo cualquier cambio de base, de modo que $T\to S$ es, de
hecho, un morfismo \emph{universal} de descenso efectivo para~$\cal{F}$).
Tomando como~$S_{i}$ abiertos afines que recubren~$S$, quedamos reducidos al
caso en que~$S$ es afín.

Entonces $S'$ es reunión finita de abiertos afines y, tomando el
$S$-esquema suma disjunta de estos, se obtiene un $S$-esquema afín~$S_{1}$
y un $S$-morfismo~$S_{1}\to S'$ plano y sobreyectivo. Por tanto, $S_{1}$ es
también fielmente plano sobre~$S$. Así, si se demuestra que un morfismo
afín y fielmente plano es un morfismo de descenso efectivo para~$\cal{F}$
y, por tanto, un morfismo universal de descenso estricto para~$\cal{F}$
(pues la hipótesis es estable bajo cambio de base), se concluye, en
particular, que el morfismo estructural~$S_{1}\to S$ es un morfismo
universal de descenso estricto para~$\cal{F}$; y, como existe un
$S$-morfismo~$S_{1}\to S'$, resultará, por \cite{VIII.D}, que
$g\colon S'\to S$ es, en efecto, un morfismo de descenso estricto
para~$\cal{F}$.

\ifthenelse{\boolean{orig}}{}
{\enlargethispage{.5cm}}%
Quedamos, pues, reducidos al caso en que~$g$ es un morfismo afín y, como
se ha visto, se puede suponer además que~$S$ es afín; por tanto,
\emph{podemos suponer que $S$ y $S'$ son afines}. En este caso,
\Ref{VIII.1.2} equivale al

\begin{lemme}
\label{VIII.1.4}
Sean $A$ un anillo, $A'$ una $A$-álgebra fielmente plana, $M$ y $N$ dos
$A$-módulos, $M'$ y $N'$ los $A'$-módulos obtenidos mediante el cambio
de anillos~$A\to A'$, y $M''$,$N''$ los
$A''=A'\otimes_{A}A'$-módulos obtenidos mediante el cambio de
anillos~$A\to A''$. Entonces la sucesión de aplicaciones de conjuntos
$$
\xymatrix@C=.5cm{\Hom_{A}(M,N)\ar[r] &
\Hom_{A'}(M',N')\ar@/^2pt/[r]\ar@/_2pt/[r] &
\Hom_{A''}(M'',N'')}
$$
es exacta.
\end{lemme}

Como
\marginpar{198}
el homomorfismo~$N\to N'$ es inyectivo (pues $A'$ es fielmente plano
sobre~$A$), se ve que la primera flecha es inyectiva. Queda demostrar que
si un
\ifthenelse{\boolean{orig}}
{$A'$ homomorfismo}
{$A'$-homomorfismo}
$u'\colon M'\to N'$ es compatible con los datos de descenso, entonces
proviene de un $A$-homomorfismo~$u\colon M\to N$. Ahora bien, esto
significa sencillamente que $u'$ aplica el subconjunto~$M$ de~$M'$ en el
subconjunto~$N$ de~$N'$ (la aplicación inducida~$u\colon M\to N$ será
entonces automáticamente $A$-lineal, pues $u'$ es $A'$-lineal, y se ve
igualmente que $u'$ es necesariamente igual a~$u\otimes_{A}A'$). Si
$x\in M$, entonces $u'(x)$ pertenece al igualador del par de aplicaciones
$\xymatrix@C=.5cm{N'\ar@/^2pt/[r]\ar@/_2pt/[r] & N''}$. Quedamos, pues,
reducidos, para demostrar~\Ref{VIII.1.4},
al caso particular siguiente (correspondiente al caso~$M=A$):

\begin{corollaire}
\label{VIII.1.5}
Sea $N$ un~$A$-módulo. Entonces la sucesión de aplicaciones entre
conjuntos
$$
\xymatrix@C=.5cm{N\ar[r] & N'\ar@/^2pt/[r]\ar@/_2pt/[r] & N''}
$$
es exacta.
\end{corollaire}

Sea, en efecto, $A_{1}$ una~$A$-álgebra fielmente plana. Para demostrar
que la sucesión considerada es exacta, basta probar que lo es la sucesión
que se deduce de ella por el cambio de anillos~$A\to A_{1}$. Ahora bien,
esta última, como se ve inmediatamente, es la correspondiente al
\ifthenelse{\boolean{orig}}
{$A$-módulo}
{$A_1$-módulo}
$N_{1}=N\otimes_{A}A_{1}$ y a la $A_{1}$-álgebra
$A_{1}'=A_{1}\otimes_{A}A'$. Basta, pues, hallar un $A_{1}$ fielmente
plano sobre~$A$ tal que $\Spec(A_{1}')\to \Spec(A_{1})$ sea un morfismo
de descenso estricto para~$\cal{F}$. En efecto, basta tomar $A_{1}=A'$,
pues entonces el morfismo anterior admite un morfismo inverso por la derecha y,
por tanto, en virtud de~\cite{VIII.D}, es un morfismo de descenso efectivo
para cualquier categoría fibrada sobre~$\Sch$.

Queda, por último, demostrar que, si $N'$ es un~$A'$-módulo provisto de
un dato de descenso para~$A\to A'$, es decir, de un isomorfismo
$$
\varphi\colon N_{1}'\isomto N_{2}'
$$
entre los dos módulos deducidos de
\ifthenelse{\boolean{orig}}
{~$N$}
{~$N'$}
por los cambios de anillos
$\xymatrix@C=.5cm{A'\ar@/^2pt/[r]\ar@/_2pt/[r] & A'\otimes_{A}A'}$,
entonces
\marginpar{199}
$N'$ es isomorfo, junto con su dato de descenso, a un módulo de la forma
$N\otimes_{A}A'$. Teniendo en cuenta~\Ref{VIII.1.5}, se ve fácilmente
que este enunciado equivale al siguiente:

\begin{lemme}
\label{VIII.1.6}
Sea $N'$ un~$A'$-módulo provisto de un dato de descenso relativamente
a~$A\to A'$ (donde $A'$ es una $A$-álgebra). Sea $N$ el sub-$A$-módulo
de~$N'$ formado por los~$x$ tales que
$\varphi(x\otimes_{A'}1_{A'})=1_{A'}\otimes_{A'}x$, y consideremos el
homomorfismo canónico
$$
N\otimes_{A}A'\to N'\quoi,
$$
(que es entonces compatible con los datos de descenso). Si $A'$ es
fielmente plana sobre~$A$, este homomorfismo es un isomorfismo.
\end{lemme}

Demostremos el lema. Sea de nuevo $A_{1}$ una $A$-álgebra fielmente
plana. Para demostrar que el morfismo considerado es un isomorfismo,
basta probar que pasa a serlo después del cambio de anillos
\ifthenelse{\boolean{orig}}
{$A_{1}\to A$.}
{$A\to A_1$.}
Ahora bien, usando la planitud de~$A_{1}$ sobre~$A$, se ve que el
homomorfismo así obtenido no es otro que el que se obtendría directamente
en términos del módulo~$N'\otimes_{A}A_{1}$
sobre~$A_{1}'=A'\otimes_{A}A_{1}$, provisto del dato de descenso relativamente
a~$A_{1}\to A_{1}'$ que se deduce canónicamente, por cambio de anillos,
del que estaba dado sobre~$N'$. Así, basta hallar un~$A_{1}$ fielmente
plano sobre~$A$ tal que $\Spec(A_{1}')\to \Spec(A_{1})$ sea un morfismo
de descenso efectivo para~$\cal{F}$. Se toma entonces, como antes,
$A_{1}=A'$. Esto concluye la demostración de~\Ref{VIII.1.6} y, con ella,
la demostración de~\Ref{VIII.1.1}.

\begin{corollaire}[Descenso de secciones de módulos]
\label{VIII.1.7}
\index{descenso de secciones de módulos|hyperpage}%
Sea $g\colon S'\to S$ un morfismo de preesquemas fielmente plano y
cuasicompacto. Para todo módulo cuasicoherente~$G$ sobre~$S$, sean $G'$ y
$G''$ sus imágenes inversas sobre~$S'$ y $S''=S'\times_{S}S'$, y
consideremos el diagrama de homomorfismos de módulos sobre~$S$:
$$
\xymatrix@C=.5cm{G\ar[r] & g_*(G')\ar@/^2pt/[r]\ar@/_2pt/[r] & h_*(G'')}
$$
(donde $h\colon S''\to S$ es el morfismo estructural). Este diagrama es
\emph{exacto}.
\end{corollaire}

En
\marginpar{200}
efecto, esto significa que, para todo abierto~$U$ de~$S$, el diagrama
correspondiente formado por las secciones sobre~$U$ es exacto. Puede
suponerse, evidentemente, que $U=S$, y la exactitud en cuestión es
entonces un caso particular de~\Ref{VIII.1.2}, obtenido tomando
$F=\mathcal{O}_{S}$.

Como el funtor imagen inversa de módulos es exacto por la derecha, se
deduce formalmente de~\Ref{VIII.1.1}:

\begin{corollaire}[Descenso de módulos cociente]
\label{VIII.1.8}
\index{descenso de módulos cociente|hyperpage}%
Con las notaciones de~\Ref{VIII.1.7}, sea además, para todo módulo
cuasicoherente~$F$ sobre un preesquema, $\Quot(F)$ el conjunto de los
módulos cociente cuasicoherentes de~$F$. Con esta convención, el diagrama
de aplicaciones entre conjuntos
$$
\xymatrix@C=.5cm{\Quot(G)\ar[r] & \Quot(G')\ar@/^2pt/[r]\ar@/_2pt/[r] &
\Quot(G'')}
$$
es exacto.
\end{corollaire}

(Evidentemente, se tendría el mismo enunciado con los submódulos en lugar
de los módulos cociente, pues unos y otros se corresponden
biunívocamente). Tomando, en particular, $G=\mathcal{O}_{S}$, se obtiene:

\begin{corollaire}[Descenso de subpreesquemas cerrados]
\label{VIII.1.9}
\index{descenso de subpreesquemas cerrados|hyperpage}%
Para todo preesquema~$X$, sea $H(X)$ el conjunto de los subpreesquemas
cerrados de~$X$. Con esta notación y bajo las condiciones
de~\Ref{VIII.1.7}, el siguiente diagrama de aplicaciones entre conjuntos
$$
\xymatrix@C=.5cm{H(S)\ar[r] & H(S')\ar@/^2pt/[r]\ar@/_2pt/[r] & H(S'')}
$$
\emph{es exacto}.
\end{corollaire}

Conviene completar el teorema~\Ref{VIII.1.1} con el resultado siguiente:

\begin{proposition}[Descenso de propiedades de módulos]
\label{VIII.1.10}
\index{descenso de propiedades de módulos|hyperpage}%
Sean $g\colon S'\to S$ un morfismo fielmente plano y cuasicompacto, y
$F$ un módulo cuasicoherente sobre~$S$. Para que~$F$ sea de tipo finito,
\resp de presentación finita, \resp localmente libre y de tipo finito,
es necesario y suficiente que su imagen inversa $F'$ sobre~$S'$ lo sea.
\end{proposition}

Solo hay que demostrar la suficiencia. Puede suponerse, evidentemente,
que~$S$ es afín,
\marginpar{201}
y, reemplazando entonces~$S'$ por
\ifthenelse{\boolean{orig}}
{la}
{una}
suma de abiertos afines que recubren~$S'$, se queda reducido al caso en
que~$S'$ es también afín. Nuestro enunciado equivale entonces al
siguiente:

\begin{corollaire}
\label{VIII.1.11}
Sean $A$ un anillo, $A'$ una~$A$-álgebra fielmente plana, $M$
un~$A$-módulo y $M'$ el~$A'$-módulo $M\otimes_{A}A'$. Para que~$M$ sea de
tipo finito (\resp de presentación finita, \resp localmente libre de tipo
finito), es necesario y suficiente que~$M'$ lo sea.
\end{corollaire}

En efecto, se tiene $M=\varinjlim_i M_{i}$, donde los~$M_{i}$ son los
submódulos de tipo finito de~$M$. Por consiguiente,
$M'=\varinjlim_i M_{i}'$, y, si $M'$ es de tipo finito, $M'$ es igual a
uno de los~$M_{i}'$; por planitud fiel, $M$ es entonces igual a~$M_{i}$ y,
por tanto, $M$ es de tipo finito. En consecuencia, existe una sucesión
exacta
$$
0\to R\to L\to M\to 0\quoi,
$$
con $L$ libre de tipo finito, de donde se obtiene una sucesión exacta
$$
0\to R'\to L'\to M'\to 0\quoi,
$$
con $L'$ libre de tipo finito. Así pues, si $M'$ es de presentación
finita, $R'$ es de tipo finito; por lo que precede, $R$ es de tipo finito
y, por tanto, $M$ es de presentación finita. Por último, decir que~$M$ es
localmente libre y de tipo finito significa que es de presentación
finita y plano (\cf \Ref{IV} en el caso noetheriano; el caso general se
deja al lector). Como cada una de estas propiedades desciende, lo mismo
ocurre con su conjunción, lo que concluye la demostración.

\begin{remarque}
\label{VIII.1.12}
La conjunción de~\Ref{VIII.1.1} y~\Ref{VIII.1.10} muestra que el
enunciado~\Ref{VIII.1.1} sigue siendo válido cuando se reemplaza la
categoría fibrada~$\cal{F}$ por la subcategoría fibrada formada por los
módulos cuasicoherentes de tipo finito, \resp de presentación finita,
\resp localmente libres de tipo finito, \resp localmente libres de rango
dado~$n$.
\end{remarque}

\section{Descenso de los preesquemas afines sobre otro}
\label{VIII.2}
\marginpar{202}

Como el funtor imagen inversa de módulos es compatible con el producto
tensorial y con otras operaciones tensoriales, el teorema~\Ref{VIII.1.1}
implica diversas variantes, obtenidas al considerar, en lugar de un solo
módulo cuasicoherente, un módulo cuasicoherente o un sistema de módulos
cuasicoherentes dotado de diversas estructuras adicionales que se expresan
mediante operaciones tensoriales. Por ejemplo, el dato de tres módulos
cuasicoherentes~$F$, $G$, $H$ sobre~$S$ y de una aplicación bilineal
$$
F\otimes G\to H\quoi,
$$
equivale al dato de tres módulos cuasicoherentes~$F'$, $G'$, $H'$ sobre~$S'$,
provistos de datos de descenso con respecto a~$g\colon S'\to S$, y de una
aplicación bilineal
$$
F'\otimes G'\to H'\quoi,
$$
«compatible» con dichos datos de descenso, en el sentido evidente del
término. Por ejemplo, si $F=G=H$, se ve que el dato de un módulo
cuasicoherente~$F$ sobre~$S$, dotado de una ley de álgebra (de la que por el
momento no suponemos que satisfaga axioma alguno de asociatividad,
conmutatividad o existencia de una sección unidad), equivale al mismo dato
sobre~$S'$, provisto además de un dato de descenso. Utilizando los resultados
del número anterior, se comprueba de inmediato que~$F$ satisface uno de los
axiomas usuales que acabamos de mencionar si y solo si~$F'$ lo satisface. Por
ejemplo, el dato de un álgebra cuasicoherente~$\mathcal{A}$ sobre~$S$ (por lo
que entenderemos en adelante: asociativa, conmutativa y provista de sección
unidad) equivale al dato de un álgebra cuasicoherente~$\mathcal{A}'$ sobre~$S'$,
provista de un dato de descenso con respecto a
\ifthenelse{\boolean{orig}}
{$G\colon S'\to S$.}
{$g\colon S'\to S$.}
Si se recuerda la equivalencia entre
\ifthenelse{\boolean{orig}}
{la categoría opuesta de las álgebras cuasicoherentes sobre~$S$, y de la
categoría}
{la categoría opuesta de las álgebras cuasicoherentes sobre~$S$ y la
categoría}
de los $S$-preesquemas afines sobre~$S$ (EGA~II, §1), se obtiene de inmediato:

\begin{theoreme}
\label{VIII.2.1}
Sea $\cal{F}'$ la categoría fibrada de los morfismos afines de preesquemas
$f\colon X\to S$, considerada como subcategoría fibrada de la categoría
fibrada
\marginpar{203}
de las flechas de la categoría de los preesquemas $\Sch$
\textup{(\Ref{VI.11}.a)}. Sea $g\colon S'\to S$ un morfismo de preesquemas
fielmente plano y cuasicompacto. Entonces $g$ es un morfismo de descenso
efectivo para~$\cal{F}'$.
\end{theoreme}

\section[Descenso de propiedades conjuntistas]{Descenso de propiedades
conjuntistas y de propiedades de finitud de morfismos\protect\footnotemark}
\label{VIII.3}


\footnotetext{Para otros resultados como los tratados en los números 3
y~4, \Cf EGA~IV~2.3, 2.6 y~2.7.}

\begin{proposition}
\label{VIII.3.1}
Sean $f\colon X\to Y$ un~$S$-morfismo, $g\colon S'\to S$ un morfismo
sobreyectivo y $f'\colon X'=X\times_{S}S'\to Y'=Y\times_{S}S'$ el morfismo
deducido de~$f$ mediante el cambio de base por~$g\colon S'\to S$. Para que
$f$ sea sobreyectivo (\resp radicial), es necesario y suficiente que lo
sea~$f'$.
\end{proposition}

Obsérvese que $f'$ también puede obtenerse mediante el cambio de base
$Y'\to Y$, que asimismo es sobreyectivo, pues se deduce del morfismo
sobreyectivo~$g\colon S'\to S$. Por otra parte, para todo~$y\in Y$ y
todo~$y'\in Y'$ situado sobre~$y$, se tiene un isomorfismo
$$
X'_{y'}\isomfrom X_{y}\otimes_{\kres(y)}\kres(y')\quoi,
$$
donde $X_{y}$ denota la fibra de~$X$ en~$y$, y $X'_{y'}$ la de~$X'$ en
$y'$. Se sigue que~$X_{y}$ es no vacía (\resp tiene a lo sumo un punto, y
este corresponde a una extensión residual radicial) si y solo si
\ifthenelse{\boolean{orig}}
{$X_{y'}$}
{$X'_{y'}$}
tiene la misma propiedad. Esto demuestra~\Ref{VIII.3.1}.

\begin{corollaire}
\label{VIII.3.2}
Bajo las condiciones de~\Ref{VIII.3.1}, si $f'$ es inyectivo
(\resp biyectivo), entonces $f$ también lo es.
\end{corollaire}

Esto se debe a que, si $X'_{y'}$ tiene a lo sumo un punto
(\resp exactamente un punto), lo mismo ocurre con~$X_{y}$; así es, en
efecto, puesto que el morfismo $X'_{y'}\to X_{y}$ es sobreyectivo (por
deducirse del morfismo sobreyectivo
$\Spec(\kres(y'))\to \Spec(\kres(y))$).

\begin{proposition}
\label{VIII.3.3}
Con las notaciones de~\Ref{VIII.3.1}, supongamos que~$g\colon S'\to S$ es
sobreyectivo y cuasicompacto (\resp fielmente plano y cuasicompacto). Para
que~$f$ sea cuasicompacto (\resp de tipo finito), es necesario y suficiente
que lo sea~$f'$.
\end{proposition}

\marginpar{204}
Solo hay que demostrar la suficiencia. Puede suponerse, evidentemente,
$S=Y$, pues la hipótesis hecha sobre~$g\colon S'\to S$ se conserva para
$Y'\to Y$. Además, puede suponerse que~$Y$ es afín. Entonces $Y'$ es
cuasicompacto y, por tanto, $X'$ es cuasicompacto (pues~$f'$ lo es por
hipótesis). Sea $(X_{i})_{i\in I}$ una familia de abiertos afines de~$X$ que
recubre~$X$; los~$X_{i}'$ son abiertos de~$X'$ que recubren~$X'$, por lo que
una subfamilia finita recubre~$X'$. Como $X'\to X$ es sobreyectivo, se sigue
que los~$X_{i}$ correspondientes ya recubren~$X$; por tanto, $X$ es
cuasicompacto, es decir, $f$ es cuasicompacto. Supongamos ahora que~$f'$ es
de tipo finito y demostremos que~$f$ lo es, suponiendo~$g$ fielmente plano.
Reemplazando $Y'$ por la suma de una familia de abiertos afines que lo
recubre, puede suponerse que~$Y'$ es afín. Por último, como, por lo que
precede, $X$ está recubierto por un número finito de abiertos afines
$X_{i}$, hay que demostrar que estos son de tipo finito sobre~$Y$, sabiendo
que $X_{i}'$ es de tipo finito sobre~$Y'$. Quedamos reducidos entonces al

\begin{corollaire}
\label{VIII.3.4}
Sean $B$ una $A$-álgebra, $A'$ una~$A$-álgebra fielmente plana y
$B'=B\otimes_{A}A'$ la
\ifthenelse{\boolean{orig}}
{$A$-álgebra}
{$A'$-álgebra}
deducida de~$B$ por cambio de anillos. Para que~$B$ sea de tipo finito, es
necesario y suficiente que lo sea~$B'$.
\end{corollaire}

Solo hay que demostrar la suficiencia. Se tiene
$B=\varinjlim_i B_{i}$, donde los~$B_{i}$ son las subálgebras de tipo finito
de~$B$. Por tanto, $B'=\varinjlim_i B_{i}'$ y, si $B'$ es de tipo finito
sobre~$A'$, entonces $B'$ es igual a uno de los~$B_{i}'$; por consiguiente,
$B$ es igual a~$B_{i}$ y, por tanto, es de tipo finito.

\begin{corollaire}
\label{VIII.3.5}
Supongamos de nuevo que el morfismo de cambio de base~$g\colon S'\to S$ es
fielmente plano y cuasicompacto. Para que~$f$ sea cuasifinito, es necesario y
suficiente que lo sea~$f'$.
\end{corollaire}

En efecto, la propiedad «ser cuasifinito» es, por definición, la conjunción
de «ser de tipo finito» y «tener fibras finitas», cada una de las cuales
desciende bien por~$g$: la primera en virtud de~\Ref{VIII.3.3}, y la segunda
por el razonamiento de~\Ref{VIII.3.1} (que solo utiliza la sobreyectividad
de~$g$).

\begin{remarques}
\label{VIII.3.6}
Sean $A$ un anillo y $X$ un~$A$-preesquema. Se ve fácilmente que
\marginpar{205}
las condiciones siguientes son equivalentes:

\begin{enumerate}
\item[(i)] Existe un anillo noetheriano~$A_{0}$ (que, si se desea, puede
suponerse un subanillo de tipo finito de~$A$), un $A_{0}$-preesquema de tipo
finito~$X_{0}$, un homomorfismo
\ifthenelse{\boolean{orig}}
{$A\to A_{0}$}
{$A_0\to A$}
y un $A$-isomorfismo~$X\isomto X_{0}\times_{A_{0}}A$.
\item[(ii)] El morfismo diagonal~$X\to X\times_{\Spec(A)}X$ es cuasicompacto
(condición vacía si $X$ está separado sobre~$A$), y $X$ es reunión finita de
abiertos afines~$X_{i}$ cuyos anillos~$B_{i}$ son álgebras de presentación
finita sobre~$A$, es decir, cocientes de álgebras de polinomios en un número
finito de indeterminadas por ideales de tipo finito.
\end{enumerate}

Si $X$ es él mismo afín, de anillo~$B$, estas condiciones significan
simplemente que~$B$ es un álgebra de presentación finita sobre~$A$.

Un morfismo~$f\colon X\to Y$ se llama \emph{morfismo de presentación
finita},
\index{presentacion finita@presentación finita (morfismo de)|hyperpage}%
y se dice también que~$X$ es de presentación finita sobre~$Y$ si $Y$ es
reunión de abiertos afines~$Y_{i}$
\ifthenelse{\boolean{orig}}
{tal que}
{tales que}
$X|Y_{i}$, como $Y_{i}$-preesquema, satisface las condiciones equivalentes
anteriores. Lo mismo ocurre entonces con $X|Y'$ para \emph{todo} abierto
afín~$Y'$ de~$Y$. Esta propiedad es estable por cambio de base y, además, el
compuesto de dos morfismos de presentación finita es de presentación finita.

Sentadas estas nociones, se ve en~(ii), procediendo como
en~\Ref{VIII.1.10}, que este enunciado sigue siendo válido al reemplazar en
él las palabras «de tipo finito» por «de presentación finita».
\end{remarques}

\section{Descenso de propiedades topológicas}
\label{VIII.4}

\begin{theoreme}
\label{VIII.4.1}
Sean $g\colon Y'\to Y$ un morfismo y $Z$ un subconjunto de~$Y$.
Supongamos que~$g$ es plano y que existe un morfismo
cuasicompacto~$f\colon X\to Y$ tal que~$Z=f(X)$ (N.B.: si $Y$ es
noetheriano, esta última condición está implicada por
\ifthenelse{\boolean{orig}}
{$Z$ es constructible}
{«$Z$ es constructible»}). Entonces se tiene
$$
g^{-1}(\overline{Z})=\overline{g^{-1}(Z)}\quoi.
$$
\end{theoreme}

\ifthenelse{\boolean{orig}}
{Puede suponerse~$Y$ afín y después $Y'$ afín. Como $Y$ es afín,
$X$ es una \emph{reunión
\marginpar{206}
finita de abiertos afines~$X_{i}$ y, reemplazando~$X$ por la suma de
los~$X_{i}$, puede suponerse asimismo que~$X$ es afín. Sean~$A,A',B$ los
anillos de~$Y,Y',X$, sea $B'=B\otimes_{A}A'$ el de
$X'=X\times_{Y}Y'$, sea $I$ el núcleo de~$A\to B$ y sea $I'$ el núcleo
de~$A'\to B'$. Así, los subconjuntos cerrados de~$Y$ y~$Y'$ definidos
por estos ideales son, respectivamente, la adherencia de~$Z=f(X)$ y la
adherencia de~$Z'=f'(X')=g^{-1}(Z)$. Se quiere probar que esta última es
igual a~$g^{-1}(\overline{Z})$, lo que resultará de~$I'=IA'$; a su vez,
esta igualdad es consecuencia de la planitud de~$A'$ sobre~$A$.}}
{Puede suponerse~$Y$ afín y después $Y'$ afín. Como $Y$ es afín,
$X$ es una reunión
\marginpar{206}
finita de abiertos afines~$X_{i}$ y, reemplazando~$X$ por la suma de
los~$X_{i}$, puede suponerse asimismo que~$X$ es afín. Sean~$A,A',B$ los
anillos de~$Y,Y',X$, sea $B'=B\otimes_{A}A'$ el de
$X'=X\times_{Y}Y'$, sea $I$ el núcleo de~$A\to B$ y sea $I'$ el núcleo
de~$A'\to B'$. Así, los subconjuntos cerrados de~$Y$ y~$Y'$ definidos
por estos ideales son, respectivamente, la adherencia de~$Z=f(X)$ y la
adherencia de~$Z'=f'(X')=g^{-1}(Z)$. Se quiere probar que esta última es
igual a~$g^{-1}(\overline{Z})$, lo que resultará de~$I'=IA'$; a su vez,
esta igualdad es consecuencia de la planitud de~$A'$ sobre~$A$.}

\begin{corollaire}
\label{VIII.4.2}
Sean $g\colon Y'\to Y$ un morfismo plano y cuasicompacto y $Z'$ un
subconjunto cerrado de~$Y'$ saturado para la relación de equivalencia
conjuntista definida por~$g$. Entonces
$Z'=g^{-1}(\overline{g(Z')})$.
\end{corollaire}

\ifthenelse{\boolean{orig}}{}
{\enlargethispage{.5cm}}%
En efecto, se tiene $Z'=g^{-1}(Z)$, con~$Z=g(Z')$. Puede aplicarse
entonces~\Ref{VIII.4.1}, observando que la condición impuesta a~$Z$ en
\Ref{VIII.4.1} se satisface tomando como~$X$ el preesquema~$Z'$ provisto de
la estructura reducida inducida por~$Y'$. (El hecho de que~$g$ sea
cuasicompacto asegura entonces que el morfismo inducido
$f\colon Z'\to Y$ sea cuasicompacto).

El enunciado~\Ref{VIII.4.2} significa también que \emph{la topología de
~$g(Y')$ inducida por~$Y$ es la topología cociente de la de~$Y'$}. En
particular:

\begin{corollaire}
\label{VIII.4.3}
Sea $g\colon Y'\to Y$ un morfismo fielmente plano y cuasicompacto. Entonces
$g$ hace de~$Y$ un espacio topológico cociente de~$Y'$; es decir, para un
subconjunto~$Z$ de~$Y$, $Z$ es cerrado (\resp abierto) si y solo si lo es
$Z'=g^{-1}(Z)$.
\end{corollaire}

Recordemos ahora que dos elementos
\ifthenelse{\boolean{orig}}
{$a,b$, de~$Y'$}
{$a,b$ de~$Y'$}
tienen la misma imagen en~$Y$ si y solo si
\ifthenelse{\boolean{orig}}
{son}
{son}
de la forma $p_{1}(c),p_{2}(c)$ para un elemento conveniente~$c$ de
$Y''=Y'\times_{Y}Y'$. Se sigue que, si $g$ es sobreyectivo, se tiene un
diagrama \emph{exacto} de conjuntos
$$
\xymatrix@C=.5cm{\cal{P}(Y)\ar[r] & \cal{P}(Y')\ar@/^2pt/[r]\ar@/_2pt/[r] &
\cal{P}(Y'')}\quoi,
$$
donde, para todo conjunto~$E$, se denota por~$\cal{P}(E)$ el conjunto de
sus subconjuntos. Sentado esto, \Ref{VIII.4.3} puede interpretarse también
del modo siguiente:

\begin{corollaire}[Descenso de los subconjuntos abiertos \resp cerrados]
\label{VIII.4.4}
Sea
\marginpar{207}%
\index{descenso de los subconjuntos abiertos (\resp cerrados)|hyperpage}%
$g\colon Y'{\to} Y$ como en~\Ref{VIII.4.3}. Para todo preesquema~$X$, sea
$\Ouv(X)$ \resp $\Fer(X)$ el conjunto de sus subconjuntos abiertos \resp
el conjunto de sus subconjuntos cerrados. Entonces se tienen diagramas
\emph{exactos} de aplicaciones entre conjuntos (deducidos de~$g$ y de las
dos proyecciones de~$Y''=Y'\times_{Y}Y'$):
\vspace*{-5pt}
\begin{gather*}
\xymatrix@C=.5cm{\Ouv(Y)\ar[r] & \Ouv(Y')\ar@/^2pt/[r]\ar@/_2pt/[r] &
\Ouv(Y'')}\\[-3pt]
\xymatrix@C=.5cm{\Fer(Y)\ar[r] & \Fer(Y')\ar@/^2pt/[r]\ar@/_2pt/[r] &
\Fer(Y'')}
\end{gather*}
\end{corollaire}

Se tiene el complemento siguiente de~\Ref{VIII.4.3}:

\begin{corollaire}
\label{VIII.4.5}
Sea $g\colon Y'\to Y$ como en~\Ref{VIII.4.3}, y sea~$Z$ un subconjunto
de~$Y$ tal que exista un morfismo cuasicompacto~$f\colon X\to Y$ de
imagen~$Z$ (por ejemplo, $Z$ constructible e~$Y$ noetheriano). Para que
$Z$ sea un subconjunto localmente cerrado de~$Y$, es necesario y suficiente
que $Z'=g^{-1}(Z)$ sea un subconjunto localmente cerrado de~$Y'$.
\end{corollaire}

Basta probar la suficiencia. Sea $Y_{1}$ el subpreesquema cerrado de~$Y$
dado por la adherencia de~$Z$ provista de la estructura reducida inducida,
y sea~$Y_{1}'=Y_{1}\times_{Y}Y'$ el subpreesquema cerrado de~$Y'$ imagen
inversa de~$Y_{1}$. Su conjunto subyacente es la imagen inversa
$g^{-1}(Y_{1})=g^{-1}(\overline{Z})$ y, por tanto, es igual, en virtud
de~\Ref{VIII.4.1}, a~$\overline{Z'}$. Como $Z'$ es localmente cerrado
en~$Y'$, es abierto en~$\overline{Z'}$ y, por tanto, abierto en~$Y_{1}'$.
Ahora bien, este último es fielmente plano y cuasicompacto sobre~$Y_{1}$;
por~\Ref{VIII.4.3} se concluye que~$Z$ es abierto en~$Y_{1}$, es decir,
en~$\overline{Z}$, lo que significa que $Z$ es localmente cerrado.

\begin{corollaire}
\label{VIII.4.6}
Sea $g\colon S'\to S$ un morfismo fielmente plano y cuasicompacto,
$f\colon X\to Y$ un~$S$-morfismo y $f'\colon X'\to Y'$ el
$S'$-morfismo que se deduce de él por cambio de base. Supongamos que~$f'$
es una aplicación abierta (\resp una aplicación cerrada, \resp
\ifthenelse{\boolean{orig}}
{cuasicompacta}
{cuasicompacta}
y un homeomorfismo sobre su imagen, \resp un homeomorfismo); entonces~$f$
tiene la misma propiedad.
\end{corollaire}

Como $Y'$ es fielmente plano y cuasicompacto sobre~$Y$, puede suponerse
$Y=S$. Sea~$Q$
\marginpar{208}
un subconjunto de~$X$. Se tiene entonces (denotando por~$h$ el morfismo de
proyección $X'\to X$):
$$
g^{-1}(f(Q))=f'(h^{-1}(Q))\quoi.
$$

Si $Q$ es abierto (\resp cerrado), lo mismo sucede con~$h^{-1}(Q)$ y, por
tanto, también con~$f'(h^{-1}(Q))$ si se supone que~$f'$ es una aplicación
abierta (\resp cerrada); por la fórmula precedente y~\Ref{VIII.4.3}, lo
mismo sucede entonces con~$f(Q)$. Esto demuestra las dos primeras
afirmaciones de~\Ref{VIII.4.6}. Queda por examinar el caso en que~$f'$ es
un homeomorfismo sobre su imagen, y demostrar entonces que~$f$ es un
homeomorfismo sobre su imagen. (El caso de un homeomorfismo resultará
entonces de~\Ref{VIII.3.1}). En virtud
\ifthenelse{\boolean{orig}}
{de~\Ref{VIII.3.1}}
{de~\Ref{VIII.3.2}}
$f$ es inyectivo; queda por probar que la aplicación $X\to f(X)$ es
abierta. Se sabe ya que~$f$ es cuasicompacto en virtud
de~\Ref{VIII.3.3}. Basta probar que, para todo subconjunto cerrado~$Z$
de~$X$, se tiene $Z=f^{-1}(\overline{f(Z)})$; esto equivale a la fórmula
análoga para las imágenes inversas por la aplicación sobreyectiva
$h\colon X'\to X$, es decir, a
$$
\ifthenelse{\boolean{orig}}
{Z'=f'^{-1}(g^{-1}(\overline{f(Z)})}
{Z'=f'^{-1}(g^{-1}(\overline{f(Z)}))}\quoi,
$$
donde se pone $Z'=h^{-1}(Z)$. Ahora bien, en virtud de~\Ref{VIII.4.1}
aplicado al subconjunto~$f(Z)$ de~$Y$, se tiene
$g^{-1}(\overline{f(Z)})=\overline{g^{-1}(f(Z))}$, y la fórmula que hay
que demostrar equivale a
$$
Z'=f'^{-1}(\overline{f'(Z')})\quoi,
$$
la cual resulta de la hipótesis de que~$f'$ es un homeomorfismo sobre su
imagen.

N.B.: en este último razonamiento, suponiendo ya que~$f$ es cuasicompacto,
no se ha utilizado que~$g$ sea cuasicompacto, sino solamente que
\ifthenelse{\boolean{orig}}
{$g'$}
{$g$}
es fielmente plano. Por tanto, bajo esta hipótesis puede descenderse la
propiedad «homeomorfismo sobre su imagen» o «homeomorfismo», o también,
gracias al razonamiento precedente, la propiedad «$f'$ es cuasicompacto y
hace de~$f'(X')$ un espacio topológico cociente de~$X'$».

Diremos que un morfismo~$f\colon X\to Y$ de preesquemas es
\emph{universalmente abierto} (\resp \emph{universalmente cerrado},
\resp \emph{universalmente bicontinuo},
\ifthenelse{\boolean{orig}}{etc...}{etc.})
\index{universalmente abierto (\resp universalmente cerrado, \resp
universalmente bicontinuo, etc.)|hyperpage}%
si, para todo cambio de base~$Y'\to Y$, $f'\colon X'\to Y'$ es un
morfismo abierto (\resp cerrado, \resp un
\ifthenelse{\boolean{orig}}
{homeom.}
{homeomorfismo}
sobre su espacio imagen). Se deduce entonces de~\Ref{VIII.4.6}:

\begin{corollaire}
\label{VIII.4.7}
Bajo
\marginpar{209}
las condiciones de~\Ref{VIII.4.6}, para que~$f$ sea universalmente
abierto (\resp universalmente cerrado, \resp un homeomorfismo universal
sobre su imagen, \resp un homeomorfismo universal), es necesario y
suficiente que~$f'$ lo sea.
\end{corollaire}

\begin{corollaire}
\label{VIII.4.8}
Bajo las condiciones de~\Ref{VIII.4.6}, para que~$f$ sea separado
(\resp propio), es necesario y suficiente que~$f'$ lo sea.
\end{corollaire}

Decir que~$f$ es separado significa que el morfismo diagonal
$X\to X\times_{Y}X$ es cerrado, o también universalmente cerrado, y, por
tanto, la primera afirmación
\ifthenelse{\boolean{orig}}
{}
{de}
~\Ref{VIII.4.8} resulta de~\Ref{VIII.4.7}. Decir que~$f$ es propio
significa que~$f$ satisface las condiciones: a)~$f$ es de tipo finito;
b)~$f$ es separado; c)~$f$ es universalmente cerrado. La condición a)
desciende bien en virtud de~\Ref{VIII.3.3}; la b), también, por lo que
acaba de verse; finalmente, la c), asimismo, por~\Ref{VIII.4.7}.

\begin{remarques}
\label{VIII.4.9}
Recordemos que, cuando $g\colon Y'\to Y$ es un morfismo plano de tipo
finito, con~$Y$ localmente noetheriano, entonces~$g$ es un morfismo
abierto
(\ifthenelse{\boolean{orig}}{VI~6.6}{\Ref{IV.6.6}}),
lo que es un resultado más preciso que~\Ref{VIII.4.3}. Obsérvese, sin
embargo, que, si~$f$ es un morfismo fielmente plano y cuasicompacto de
preesquemas noetherianos, entonces~$f$ no es, en general, un morfismo
abierto. Sea, por ejemplo, $Y$ un esquema irreducible cuyo punto
genérico~$y$ no sea abierto (por ejemplo, una curva algebraica), y tomemos
como~$Y'$ el esquema suma~$Y\amalg\Spec(\kres(y))$. Entonces la imagen por
el morfismo estructural~$Y'\to Y$ del subconjunto
abierto~$\Spec(\kres(y))$ no es un subconjunto abierto de~$Y$. El lector
observará asimismo que diversos enunciados de la presente exposición se
vuelven falsos si se abandona en ellos la hipótesis de que el morfismo
fielmente plano considerado sea también cuasicompacto. El caso típico que
refuta los enunciados es aquel en que se toma como~$Y'$ el esquema suma de
los espectros de los anillos locales de los puntos de~$Y$. Por ejemplo,
tomando de nuevo como~$Y$ una curva algebraica irreducible y como~$Z$ el
subconjunto de~$Y$ reducido al punto genérico, su imagen inversa en~$Y'$ es
abierta sin que~$Z$ sea abierto.
\end{remarques}

\setcounter{subsection}{9}
\subsection{}
\label{VIII.4.10}
Diversos
\marginpar{210}
enunciados de la presente exposición siguen siendo válidos reemplazando la
hipótesis de que~$Y'$ sea plano sobre~$Y$ por la siguiente: existe un
\ifthenelse{\boolean{orig}}
{módulo}
{módulo}
de tipo finito~$F$ sobre~$Y'$, de soporte~$Y'$, plano con respecto a~$Y$;
la hipótesis de planitud fiel se reemplazará entonces por la precedente,
más la hipótesis de que~$Y'\to Y$ sea sobreyectivo. Esto se aplica a las
dos primeras afirmaciones de~\Ref{VIII.1.10}, a~\Ref{VIII.3.3},
\Ref{VIII.3.5}, \Ref{VIII.4.1} y, por consiguiente, a todos los resultados
del presente número.

\section{Descenso de morfismos de preesquemas}
\label{VIII.5}

\begin{proposition}
\label{VIII.5.1}
Sea $g\colon S'\to S$ un morfismo de preesquemas.
\begin{enumerate}
\item[a)] Supongamos que~$g$ es sobreyectivo y que el homomorfismo
$$
g^*\colon \mathcal{O}_{S}\to g_*(\mathcal{O}_{S'})
$$
es inyectivo. Entonces~$g$ es un epimorfismo en la categoría de los
preesquemas, e incluso en la categoría de los espacios anillados.
\item[b)] Supongamos que~$g$ es sobreyectivo y hace de~$S$ un espacio
topológico cociente de~$S'$. Sea~$S''=S'\times_{S}S'$ y sea
$h\colon S''\to S$ el morfismo estructural; consideremos el diagrama de
homomorfismos canónico:
$$
\xymatrix@C=.5cm{\mathcal{O}_{S} \ar[r] &
g_*(\mathcal{O}_{S'})\ar@/^2pt/[r]\ar@/_2pt/[r] &
h_*(\mathcal{O}_{S''})}\quoi,
$$
y supongamos que este diagrama es \emph{exacto}. Entonces~$g$ es un
epimorfismo efectivo
\index{epimorfismo efectivo|hyperpage}%
en la categoría de los preesquemas (y también en la categoría de los
espacios anillados), es decir, el diagrama
$$
\xymatrix@C=.5cm{S & S'\ar[l] & S''\ar@/^2pt/[l]\ar@/_2pt/[l]}
$$
es exacto.
\end{enumerate}
\end{proposition}

\subsubsection*{Demostración} a) Hay que demostrar que un morfismo de
espacios anillados~$f\colon S\to Z$ queda determinado cuando se conoce
$fg$. Ahora bien, como~$g$ es sobreyectivo, se conoce la aplicación entre
conjuntos~$f_{0}$ subyacente a~$f$; queda por determinar el homomorfismo
de haces de anillos~$\mathcal{O}_{Z}\to \mathcal{O}_{S}$ o, lo que es lo
mismo, el homomorfismo
\marginpar{211}
de haces de anillos
$$
u\colon f_{0}^{-1}(\mathcal{O}_{Z})\to \mathcal{O}_{S}
$$
definido por~$f$. Se conoce ya el homomorfismo
$$
(fg)_{0}^{-1}(\mathcal{O}_{Z})=
g_{0}^{-1}(f_{0}^{-1}(\mathcal{O}_{Z}))\to \mathcal{O}_{S'}
$$
definido por~$fg$ o, lo que es lo mismo, se dispone de un homomorfismo
$$
f_{0}^{-1}(\mathcal{O}_{Z})\to
g_{0*}(\mathcal{O}_{S'})=g_*(\mathcal{O}_{S'})\quoi.
$$
Se comprueba inmediatamente que este último no es sino el compuesto de
$g^*\colon \mathcal{O}_{S}\to g_*(\mathcal{O}_{S'})$ y de~$u$; como
$g^*$ es inyectivo, $u$ queda determinado cuando se conoce~$g^*u$.
[N.B.: evidentemente, no se ha utilizado que~$g\colon S'\to S$ sea un
morfismo de preesquemas; el enunciado valdría para un morfismo cualquiera
de espacios anillados. La misma observación vale para b), tanto en la
categoría de los espacios anillados como en la de los espacios localmente
anillados. Obsérvese también que, si~$g$ es un morfismo de
preesquemas no necesariamente sobreyectivo, pero tal que
$g^*\colon \mathcal{O}_{S}\to g_*(\mathcal{O}_{S'})$ sea inyectivo,
entonces, para dos morfismos~$f_{1},f_{2}$ de~$S$ en un
\emph{esquema}~$Z$ tales que~$f_{1}g=f_{2}g$, se tiene $f_{1}=f_{2}$. En
efecto, si~$I$ es el ideal sobre~$S$ que define el subpreesquema de~$S$ de
las coincidencias de~$f_{1},f_{2}$ (imagen inversa del subpreesquema
diagonal de~$Z\times Z$ por~$(f_{1},f_{2})$), se ve que~$I$ está contenido
en~$\Ker(g^*)$].

b) Hay que demostrar que, para todo espacio anillado~$Z$, el diagrama
siguiente de aplicaciones
$$
\xymatrix@C=.5cm{\Hom(S,Z)\ar[r] & \Hom(S',Z)\ar@/^2pt/[r]\ar@/_2pt/[r] &
\Hom(S'',Z)}
$$
es exacto, y que lo mismo sucede cuando~$Z$ es un espacio localmente
anillado y uno se limita a los homomorfismos de espacios localmente
anillados. Como ya se sabe por a) que la primera aplicación es
inyectiva, queda por ver que, si~$f'\colon S'\to Z$ es un homomorfismo de
espacios anillados tal que~$f'p_{1}=f'p_{2}$,
\marginpar{212}
entonces~$f'$ es de la forma~$fg$, donde~$f\colon S\to Z$ es un
homomorfismo de espacios anillados. Como~$g$ es sobreyectivo, resulta
entonces evidente que, si~$f'$ es un morfismo de espacios localmente
anillados, lo mismo sucederá con~$f$.

De la hipótesis sobre~$f'$ resulta que la aplicación entre conjuntos
subyacente~$f_{0}'$ es constante sobre las fibras de la
aplicación~$g_{0}$; como esta última es sobreyectiva, $f_{0}'$ se
factoriza de manera única como~$f_{0}'=f_{0}g_{0}$, donde
$f_{0}\colon S\to Z$ es una aplicación, necesariamente continua porque
$g_{0}$ identifica~$S$ con un espacio topológico cociente de~$S'$.
Consideremos ahora el homomorfismo
$$
f_{0}^{-1}(\mathcal{O}_{Z})\to g_*(\mathcal{O}_{S'})
$$
deducido del homomorfismo~$(f_{0}g_{0})^{-1}(\mathcal{O}_{Z})\to
\mathcal{O}_{S'}$ correspondiente a~$f'$. La hipótesis
$f'p_{1}=f'p_{2}$ se interpreta entonces diciendo que los compuestos del
homomorfismo precedente con los dos homomorfismos
$\xymatrix@C=.5cm{g_*(\mathcal{O}_{S'})\ar@/^2pt/[r]\ar@/_2pt/[r] &
h_*(\mathcal{O}_{S''})}$ son iguales; por la hipótesis b), aquel se
factoriza, por tanto, a través de un morfismo
$$
f_{0}^{-1}(\mathcal{O}_{Z})\to \mathcal{O}_{S}\quoi.
$$
Este último define un morfismo de espacios anillados~$f\colon S\to Z$,
que es el morfismo buscado.

\begin{theoreme}
\label{VIII.5.2}
Sea $\cal{F}$ la categoría fibrada de las flechas en la categoría~$\Sch$
de los preesquemas \textup{(\Ref{VI.11}.a)}. Entonces todo
morfismo~$g\colon S'\to S$ fielmente plano y cuasicompacto es un morfismo
de descenso para~$\cal{F}$ (o también, como se dice, un morfismo de
descenso en~$\Sch$).
\index{descenso (morfismo de)|hyperpage}%
\end{theoreme}

Esto significa lo siguiente: sea~$S''=S'\times_{S}S'$ y, para dos
preesquemas~$X,Y$ sobre~$S$, consideremos sus imágenes inversas~$X',Y'$ sobre
$S'$ y sus imágenes inversas~$X'',Y''$ sobre~$S''$; se obtiene un diagrama
de aplicaciones
$$
\xymatrix@C=.5cm{\Hom_{S}(X,Y)\ar[r] &
\Hom_{S'}(X',Y')\ar@/^2pt/[r]\ar@/_2pt/[r] & \Hom_{S''}(X'',Y'')}\quoi;
$$
sentadas estas notaciones, \Ref{VIII.5.2} significa que este diagrama es
exacto. Obsérvese que
\marginpar{213}
no es cierto, en general, que~$g$ sea un morfismo de descenso efectivo,
es decir, que para todo preesquema~$X'$ sobre~$S'$, todo dato de descenso
sobre~$X'$ con respecto a~$g\colon S'\to S$ sea efectivo. La cuestión de
la efectividad, a menudo delicada, se examinará en el
\No~\Ref{VIII.7}.

Se ha visto en~\cite{VIII.D} (teniendo en cuenta que en~$\Sch$ existen los
productos fibrados) que el enunciado~\Ref{VIII.5.2} equivale al siguiente:

\begin{corollaire}
\label{VIII.5.3}
Todo morfismo fielmente plano y cuasicompacto de preesquemas es un
epimorfismo efectivo universal.
\end{corollaire}

Como un morfismo fielmente plano y cuasicompacto sigue siéndolo tras toda
extensión de la base, se reduce uno a demostrar que es un epimorfismo
efectivo. Se aplica entonces el criterio~\Ref{VIII.5.1} b), que da el
resultado buscado teniendo en cuenta~\Ref{VIII.4.3} y~\Ref{VIII.1.7}.

\begin{corollaire}
\label{VIII.5.4}
Sea $g\colon S'\to S$ un morfismo fielmente plano y cuasicompacto,
$f\colon X\to Y$ un~$S$-morfismo y $f'\colon X'\to Y'$ el
$S'$-morfismo que se deduce de él por el cambio de base~$S'\to S$. Para
que~$f$ sea un isomorfismo, es necesario y suficiente que~$f'$ lo sea.
\end{corollaire}

En efecto, si $f'$ es un isomorfismo, es también un isomorfismo para las
estructuras de descenso naturales sobre~$X',Y'$; como el
funtor~$X\mto X'$ de~$\Sch_{/S}$ en la categoría de los preesquemas
sobre~$S'$ provistos de un dato de descenso con respecto a~$g$ es plenamente
fiel por~\Ref{VIII.5.2}, se obtiene la conclusión buscada.

\begin{corollaire}
\label{VIII.5.5}
Bajo las condiciones de~\Ref{VIII.5.4}, para que~$f$ sea una inmersión
cerrada (\resp una inmersión abierta, \resp una inmersión cuasicompacta),
\ifthenelse{\boolean{orig}}
{es necesario y suficiente}
{es necesario y suficiente}
que~$f'$ lo sea.
\end{corollaire}

Puede suponerse, como de costumbre, $Y=S$, y solo hay que demostrar la
suficiencia. Obsérvese que el hecho de que~$X'/Y'$ esté provisto de un dato
de descenso con respecto a~$g\colon Y'\to Y$ y que el morfismo
estructural~$f'\colon X'\to Y'$ sea una inmersión y, por tanto, un
monomorfismo, implica que los dos subobjetos de~$Y''$ obtenidos como
imágenes
\marginpar{214}
inversas de~$X'/Y'$ por una y otra proyección de~$S''$ en~$S'$ son los
mismos. Si~$f'$ es una inmersión cerrada, de~\Ref{VIII.1.9} resulta que
existe un subpreesquema cerrado~$X_{1}$ de~$Y$ cuya imagen inversa por
$g\colon Y'\to Y$ es~$X'$. Por la unicidad de la solución de un problema
de descenso con respecto a un morfismo de descenso para~$\cal{F}$,
\ifthenelse{\boolean{orig}}
{}
{se}
sigue que~$X_{1}$ es $Y$-isomorfo a~$X$; por tanto,
$f\colon X\to Y$ es una inmersión cerrada. Se procede del mismo modo para
una inmersión abierta, utilizando~\Ref{VIII.4.4}. Finalmente, si~$f'$ es
una inmersión cuasicompacta, $f$ es cuasicompacto en virtud
de~\Ref{VIII.3.3}; puede aplicarse entonces al subconjunto~$f(X)$ de~$Y$
el criterio~\Ref{VIII.4.5}, que demuestra que~$f(X)$ es localmente cerrado,
pues su imagen inversa~$f'(X')$ en~$Y'$ lo es. Reemplazando entonces~$Y$
por un subconjunto abierto en el que~$f(X)$ sea cerrado, se reduce uno al
caso en que~$f'$ es una inmersión cerrada; por lo anterior, también lo
es~$f$.

\begin{corollaire}
\label{VIII.5.6}
Bajo las condiciones de~\Ref{VIII.5.4}, para que~$f$ sea afín, es necesario
y suficiente que~$f'$ lo sea.
\end{corollaire}

Se procede como en~\Ref{VIII.5.5}, utilizando~\Ref{VIII.2.1} (también
puede utilizarse el criterio cohomológico de Serre [EGA II 5.2], que
demuestra~\Ref{VIII.5.6} sin utilizar ninguna técnica de descenso).

\begin{corollaire}
\label{VIII.5.7}
Bajo las condiciones de~\Ref{VIII.5.4}, para que~$f$ sea entero (\resp
finito, \resp finito y localmente libre), es necesario y suficiente
que~$f'$ lo sea.
\end{corollaire}

Solo hay que demostrar la suficiencia y, como de costumbre, puede
suponerse $Y=S$, con $Y$ afín e~$Y'$ afín. Como la hipótesis implica que
$f'$ es afín, lo mismo sucede con~$f$ por~\Ref{VIII.5.6}; por tanto, $X$
y, en consecuencia, $X'$ son afines. Sean~$A,A',B,B'=B\otimes_{A}A'$ los
anillos de~$Y,Y',X,X'$. Se tiene $B=\varinjlim_i B_{i}$, donde~$B_{i}$
recorre las sub-$A$-álgebras de~$B$ que son de tipo finito sobre~$A$; de
aquí $B'=\varinjlim_i B_{i}'$, donde las~$B_{i}'$ son subálgebras de tipo
finito de la $A'$-álgebra~$B'$. Si $B'$ es entero
\ifthenelse{\boolean{orig}}
{sobre~$A$,}
{sobre~$A'$,}
las~$B_{i}'$ son módulos de tipo finito sobre~$A'$; como $A'$ es fielmente
plano sobre~$A$, las~$B_{i}$ son módulos de tipo finito sobre~$A$, es decir,
$B$ es entero sobre~$A$. Se ve del mismo modo que, si~$B'$ es finito
sobre~$A'$, $B$ lo es sobre~$A$. La misma conclusión vale para «localmente
libre de tipo finito»; \cf \Ref{VIII.1.11}.

\begin{corollaire}
\label{VIII.5.8}
Bajo
\marginpar{215}
las condiciones de~\Ref{VIII.5.4}, supongamos~$f$ cuasicompacto y sean
$\mathcal{L}$ un haz invertible sobre~$X$ y~$\mathcal{L}'$ su imagen
inversa sobre~$X'$. Para que~$\mathcal{L}$ sea amplio (\resp muy amplio)
con respecto a~$f$, es necesario y suficiente
\ifthenelse{\boolean{orig}}
{que~$f'$}
{que~$\cal{L}'$}
sea amplio (\resp muy amplio) con respecto a~$f'$.
\end{corollaire}

Solo hay que demostrar la suficiencia. La hipótesis sobre
\ifthenelse{\boolean{orig}}
{$\mathcal{L}$}
{$\mathcal{L}'$}
implica, en todo caso, que~$f'$ es separado; por~\Ref{VIII.4.8}, $f$ es
separado. Como~$f$ es cuasicompacto y~$g\colon Y'\to Y$ es plano, el
cálculo de las imágenes directas mediante recubrimientos afines demuestra
que se tienen isomorfismos
$$
g^*(f_*(\mathcal{L}^{\otimes n}))\isomto
f_*'(\mathcal{L}'^{\otimes n})
$$
para todo entero~$n$ y, por tanto, un isomorfismo
$$
g^*(\mathcal{S})\isomto \mathcal{S}'\quoi,
$$
donde~$\mathcal{S}$ (\resp $\mathcal{S}'$) denota el álgebra graduada
cuasicoherente sobre~$Y$ (\resp sobre~$Y'$) suma directa de
los~$f_*(\mathcal{L}^{\otimes n})$ (\resp de
los~$f_*'(\mathcal{L}'^{\otimes n})$) para~$n\geq 0$. Obsérvese que, para
todo~$n\geq 0$, el conúcleo del homomorfismo
canónico~$f'^*(\mathcal{S}_{n}')\to \mathcal{L}'^{\otimes n}$ es la
imagen inversa por~$X'\to X$ del conúcleo de
$f^*(\mathcal{S}_{n})\to \mathcal{L}^{\otimes n}$; por tanto, su
soporte~$Z_{n}'$ es la imagen inversa del soporte~$Z_{n}$. Si
\ifthenelse{\boolean{orig}}
{$f$}
{$\cal{L}'$}
es amplio, la intersección de los~$Z_{n}'$ es vacía; como~$X'\to X$ es
sobreyectivo, la intersección de los~$Z_{n}$ es vacía, es decir, se tiene
un morfismo canónico
$$
j\colon X\to \Proj(\mathcal{S})
$$
(EGA~II 3). Además, el morfismo análogo
$$
j'\colon X'\to \Proj(\mathcal{S}')
$$
no es sino el que se deduce del precedente mediante el cambio de
base~$Y'\to Y$ (\loccit). Sentado esto, decir que~$\mathcal{L}'$ es amplio
con respecto a~$f'$ significa que~$j'$ es una inmersión, necesariamente
cuasicompacta, además, porque~$f'$ es cuasicompacto. Por
\Ref{VIII.5.5}, $j$ es, por tanto, una inmersión, es decir,
$\mathcal{L}$ es amplio con respecto a~$f$. --- Se procede de manera
enteramente análoga en el caso «muy
\marginpar{216}
amplio», limitándose arriba a~$n=1$ y reemplazando la consideración de
$\Proj(\mathcal{S})$ por la del fibrado proyectivo
$\mathcal{P}(\mathcal{S}_{1})$ asociado a~$\mathcal{S}_{1}$.

Recordemos (EGA~II 5.1.1) que un morfismo cuasicompacto~$f$ se llama
\emph{cuasiafín} si, para todo abierto afín~$U$ de~$Y$, $f^{-1}(U)$ es un
preesquema isomorfo a un subesquema abierto de un esquema afín. Se
demuestra (\loccit) que equivale a decir que~$\mathcal{O}_{X}$ es amplio
(o también muy amplio) con respecto a~$f$. Por tanto, \Ref{VIII.5.8}
implica:

\begin{corollaire}
\label{VIII.5.9}
Bajo las condiciones de~\Ref{VIII.5.4} y suponiendo~$f$ cuasicompacto,
para que~$f$ sea cuasiafín, es necesario y suficiente que~$f'$ lo sea.
\end{corollaire}

\begin{remarques}
\label{VIII.5.10}
El ejemplo de variedad no proyectiva de Hironaka muestra que puede tenerse
un morfismo propio~$f\colon X\to Y$ de variedades algebraicas no singulares
(con~$Y$ proyectiva), tal que~$Y$ sea reunión de dos abiertos~$Y_{i}$ para
los cuales~$X_{i}=X\times_{Y}Y_{i}$ sea proyectivo sobre~$Y_{i}$, pero sin
que~$f$ sea proyectivo. Poniendo entonces $Y'=Y_{1}\amalg Y_{2}$, $Y'$ es
fielmente plano y cuasicompacto (e incluso cuasifinito) sobre~$Y$,
$f'\colon X'\to Y'$ es proyectivo, pero~$f$ no lo es. Hay que tener
cuidado, por tanto: para aplicar~\Ref{VIII.5.8} y deducir del hecho de
que~$f'$ sea proyectivo la misma conclusión sobre~$f$, hay que disponer ya
sobre~$X'$ de un haz invertible~$\mathcal{L}'$ amplio para~$f'$,
\emph{provisto de un dato de descenso con respecto a}~$X'\to X$ (lo que
permite considerar~$\mathcal{L}'$ como la imagen inversa de un haz
invertible~$\mathcal{L}$ sobre~$X$, que será entonces amplio para~$f$
gracias a~\Ref{VIII.5.8}). No obstante, cuando~$g\colon S'\to S$ es finito
y localmente libre, véase~\Ref{VIII.7.7}.
\end{remarques}

\section[Aplicación a los morfismos finitos y cuasifinitos]{Aplicación a los morfismos finitos y cuasifinitos\protect\footnotemark}
\label{VIII.6}
\footnotetext{\Cf EGA IV 18.12 para generalizaciones a preesquemas no necesariamente localmente noetherianos}

Vamos a demostrar los dos teoremas siguientes:

\begin{theoreme}
\label{VIII.6.1}
Sea~$f\colon X\to Y$ un morfismo \emph{propio con fibras finitas}, con~$Y$
localmente noetheriano. Entonces~$f$ es finito.
\end{theoreme}

\begin{theoreme}
\label{VIII.6.2}
Sea~$f\colon X\to Y$ un morfismo \emph{cuasifinito} y \emph{separado},
con~$Y$ localmente noetheriano. Entonces~$f$ es cuasiafín y, a fortiori,
cuasiproyectivo.
\end{theoreme}

\begin{remarques}
\label{VIII.6.3}
El
\marginpar{217}
teorema~\Ref{VIII.6.1} es bien conocido y se debe a Chevalley en el caso de
las variedades algebraicas; también se encontrará una demostración sencilla
en [EGA III 4], que utiliza el «teorema de las funciones holomorfas». La
demostración dada aquí no utiliza este último teorema, sino la teoría del
descenso; la damos como obsequio al lector, pues se obtiene «gratis» al mismo
tiempo que la de~\Ref{VIII.6.2}. Recordemos también ([EGA III 4] o
\cite{VIII.1}) que la forma global del \emph{Main Theorem} de Zariski,
deducida del «teorema de las funciones holomorfas», afirma que, si
$f\colon X\to Y$ es cuasifinito y \emph{cuasiproyectivo}, siendo~$Y$
noetheriano, entonces~$X$ es $Y$-isomorfo a un subpreesquema abierto de un
$Y$-preesquema \emph{finito}~$Z$. La conjunción del \emph{Main Theorem} y
de~\Ref{VIII.6.2} se enuncia, por tanto, así:
\end{remarques}

\begin{corollaire}
\label{VIII.6.4}
Sea~$f\colon X\to Y$ un morfismo cuasifinito y separado, con~$Y$
noetheriano. Entonces~$X$ es $Y$-isomorfo a un subpreesquema abierto de un
$Y$-preesquema finito~$Z$.
\end{corollaire}

Otra consecuencia interesante de~\Ref{VIII.6.2} para la teoría del descenso
se dará junto con~\Ref{VIII.7.9}.

\subsubsection*{Demostración de~\Ref{VIII.6.1} y~\Ref{VIII.6.2}}
Admitiremos el hecho siguiente, cuya demostración es
fácil\footnote{\Cf EGA IV 8}:

\begin{lemme}
\label{VIII.6.5}
Sea~$X$ un preesquema de tipo finito sobre~$Y$ localmente noetheriano, y
sea~$y\in Y$. Para que exista un entorno abierto~$U$ de~$y$ tal que~$X|U$
sea finito (\resp cuasiafín, \resp\ ...) sobre~$U$, es necesario y suficiente
que~$X\times_{Y}\Spec(\mathcal{O}_{y})$ sea finito (\resp cuasiafín,
\resp\ ...) sobre~$\Spec(\mathcal{O}_{y})$.
\end{lemme}

Como, por otra parte, la propiedad de~$f\colon X\to Y$ de ser finito,
\resp cuasiafín, es local sobre~$Y$, para demostrar~\Ref{VIII.6.1}
y~\Ref{VIII.6.2} se reduce uno al caso en que~$Y$ es el espectro de un anillo
local y, a fortiori, tiene dimensión finita. (N.B.: se llama
\emph{dimensión de un preesquema}~$Y$ al supremo de las dimensiones de Krull
de sus anillos locales). Procedemos por inducción sobre
$$
n=\dim(Y)\quoi,
$$
siendo la afirmación
\marginpar{218}
trivial para~$n<0$. Podemos suponer, por tanto, $n\geq 0$ y demostrada la
afirmación para las dimensiones~$n'<n$. Puede suponerse de nuevo que~$Y$ es
el espectro de un anillo local noetheriano~$A$, de dimensión~$n$. Obsérvese
que las hipótesis hechas en~\Ref{VIII.6.1} y~\Ref{VIII.6.2} son estables por
cambio de base (ya se ha hecho uso de ello en la reducción inicial); seguirán
siendo ciertas tras el cambio de base~$\Spec(\widehat{A})\to \Spec(A)$. Como
este último es fielmente plano y cuasicompacto, los
enunciados~\Ref{VIII.5.7} y~\Ref{VIII.5.9} reducen la cuestión al caso en
que~$A$ es, además, completo. Utilizando entonces el hecho de que todo anillo
local noetheriano~$B$ sobre~$A$ que es cuasifinito sobre~$A$ es finito
sobre~$A$, y el hecho de que~$X$ es separado sobre~$Y$ y la fibra sobre~$y$ está formada por
puntos aislados, se obtiene una descomposición
$$
X=X'\amalg X''
$$
en la que~$X'$ es \emph{finito} sobre~$Y$ y la fibra de~$X''$ en~$y$ es
vacía. Si~$X$ es propio sobre~$Y$, lo mismo sucede con~$X''$; su imagen en
$Y$ es, por tanto, cerrada y, como no contiene a~$y$, es vacía. Así,
$X''=\emptyset$ y, por consiguiente, $X=X'$, lo que demuestra que~$X$ es
finito sobre~$Y$ y prueba~\Ref{VIII.6.1} (N.B.: aquí es inútil la hipótesis
de inducción). Si~$X$ es cuasifinito sobre~$Y$, $X''$ también lo es; ahora
bien, $X''$ se encuentra de hecho sobre
\ifthenelse{\boolean{orig}}
{el abierto~$Y-(y)$}
{el abierto~$Y-\{ y\}$}
de~$Y$, \emph{que tiene dimensión} $<n$. En virtud de la hipótesis de
inducción, $X''$ es cuasiafín
\ifthenelse{\boolean{orig}}
{sobre~$Y-(y)$,}
{sobre~$Y-\{ y\}$,}
y, por tanto, también sobre~$Y$; evidentemente, lo mismo sucede con~$X'$ y,
por consiguiente, con su suma~$X$, lo que prueba~\Ref{VIII.6.2}.

\begin{remarque}
\label{VIII.6.6}
Los teoremas~\Ref{VIII.6.1} y~\Ref{VIII.6.2} siguen siendo válidos si ya no
se supone~$Y$ localmente noetheriano, con tal de precisar que se supone~$f$
de presentación finita
\ifthenelse{\boolean{orig}}
{(\cf \Ref{VIII.3.6}).}
{(\cf \Ref{VIII.3.6}).}
En efecto, todavía puede suponerse~$Y$ afín; se comprueba entonces sin
dificultad que la situación~$f\colon X\to Y$
\ifthenelse{\boolean{orig}}
{se deduce,}
{se deduce,}
mediante un cambio de base~$Y\to Y_{0}$, de una situación
$f_{0}\colon X_{0}\to Y_{0}$ que satisface las mismas hipótesis que~$f$,
con~$Y_{0}$ \emph{noetheriano}. Por tanto, según el
resultado~\Ref{VIII.6.1} \resp \Ref{VIII.6.2}, $f_{0}$ es finito
\resp cuasiafín y, en consecuencia, lo mismo sucede con~$f$. Este género de
razonamiento resulta útil a menudo para librarse de hipótesis noetherianas
\marginpar{219}
(que siempre acaban siendo molestas en las aplicaciones).
\end{remarque}

\section{Criterios de efectividad para un dato de descenso}
\label{VIII.7}

Consideremos, como de costumbre, un morfismo de preesquemas
$$
g\colon S'\to S
$$
y un $S'$-preesquema~$X'$. De acuerdo con los hechos
\ifthenelse{\boolean{orig}}
{generales~VII,~9,}
{generales~VII,~9\footnote{Este número no existe; véase la exposición 212
del Seminario Bourbaki (1962)},}
dar un dato de descenso sobre~$X'$ con respecto a~$g$ equivale a dar un par
de equivalencia \cite{VIII.3}:
$$
\xymatrix@C=.5cm{ q_{1},q_{2}\colon X''\ar@/^2pt/[r]\ar@/_2pt/[r] & X'}
$$
tal que el morfismo estructural~$X'\to S'$ sea compatible con este par y con
el par de equivalencia
$$
\xymatrix@C=.5cm{p_{1},p_{2}\colon
S''=S'\times_{S}S'\ar@/^2pt/[r]\ar@/_2pt/[r] & S'}
$$
definido por~$g$, y tal que los dos cuadrados (o uno de ellos, lo que viene
a ser lo mismo por simetría) extraídos del diagrama correspondiente
$$
\xymatrix{X'\ar[d] & X''\ar@/^2pt/[l]\ar@/_2pt/[l]\ar[d]\\ S' &
S''\ar@/^2pt/[l]\ar@/_2pt/[l]}
$$
(utilizando bien~$p_{1},q_{1}$, bien~$p_{2},q_{2}$)
\ifthenelse{\boolean{orig}}
{sean \emph{cartesianos}}
{sean \emph{cartesianos}}.
Una solución del problema de descenso planteado por este
dato de descenso, \ie un objeto~$X$ sobre~$S$ provisto de un
isomorfismo~$X\times_{S}S'\isomfrom X'$ compatible con los datos de descenso,
equivale a dar un cuadrado \emph{cartesiano}
$$
\xymatrix{X\ar[d] & X'\ar[l]_{h}\ar[d]\\ S & S'\ar[l]_{g}}
$$
que satisface~$hq_{1}=hq_{2}$.

Como
\marginpar{220}
la clase de los morfismos fielmente planos y cuasicompactos es estable por
cambio de base, y un morfismo fielmente plano y cuasicompacto es un
epimorfismo efectivo en virtud de~\Ref{VIII.5.3}, se sigue de la teoría
general \cite{VIII.D}:

\begin{proposition}
\label{VIII.7.1}
Supongamos que~$g\colon S'\to S$ es fielmente plano y cuasicompacto. Para que
un dato de descenso sobre~$X'$ con respecto a~$g$ sea efectivo, es necesario
y suficiente que la relación de equivalencia~$R=(q_{1},q_{2})$ que este define
sea efectiva (\ie que exista el cociente~$X'/R$ y que~$X''$ se convierta en
el cuadrado fibrado de~$X'$ sobre~$X'/R$), y que el morfismo
canónico~$X'\to X'/R$ sea fielmente plano y cuasicompacto.
\end{proposition}

Así, la cuestión de la efectividad de un dato de descenso es un caso
particular de la cuestión de la efectividad de un grafo de equivalencia, y
por esta vía pueden obtenerse diversos criterios de efectividad dados en
este número. Sin embargo, en el contexto del descenso se dispone del
teorema~\Ref{VIII.2.1}, que implica que \emph{si $X'$ es afín sobre~$S'$,
todo dato de descenso sobre~$X'$ con respecto a~$g$ es efectivo}; este
enunciado no tiene análogo para el paso al cociente por un grafo de
equivalencia plano general. Todos los criterios de efectividad que damos
aquí pueden considerarse asimismo deducidos del anterior.

Sea~$U'$ un subpreesquema de~$X'$ (o, más generalmente, un subobjeto de~$X'$
en la categoría~$\Sch$); se dice que~$U'$ es \emph{estable por el dato de
descenso} sobre~$X'$ si puede dotarse a~$U'$ de un dato de descenso con
respecto a~$g$ tal que la inmersión~$U'\to X'$ sea compatible con los datos
de descenso. Esto significa también que las imágenes inversas de~$U'$ en
~$X''$ por~$q_{1}$ y~$q_{2}$ son iguales (o también, como se dice, que~$U'$
es \emph{estable por la relación de equivalencia}~$R$); desde luego, el dato
de descenso en cuestión sobre~$U'$ es entonces único y se llama \emph{dato
de descenso inducido}
\index{dato de descenso inducido|hyperpage}%
por el de~$X'$. Sentado esto:

\begin{proposition}
\label{VIII.7.2}
Sea~$(X_{i}')$ un recubrimiento de~$X'$ por abiertos~$X_{i}'$ estables
\marginpar{221}
por el dato de descenso. Para que el dato de descenso sobre~$X'$ sea
efectivo, es necesario y suficiente que lo mismo suceda con los datos de
descenso inducidos en los~$X_{i}'$.
\end{proposition}

Esta es, por ejemplo, una consecuencia fácil de~\Ref{VIII.7.1}; se deja al
lector el detalle de la demostración.

\begin{corollaire}
\label{VIII.7.3}
Sea~$(S_{i})$ un recubrimiento abierto de~$S$, y para todo~$i$ sean~$S_{i}'$
y~$X_{i}'$ deducidos de~$S'$
\ifthenelse{\boolean{orig}}
{de}
{y}
$X'$ mediante el cambio de base~$S_{i}\to S$. Para que el dato de descenso
sobre~$X'$ sea efectivo, es necesario y suficiente que, para todo~$i$, el
dato de descenso sobre~$X_{i}'$ con respecto a
~$g_{i}\colon S_{i}'\to S_{i}$ sea efectivo.
\end{corollaire}

Este criterio nos reduce siempre, en la práctica, al caso en que~$S$ es
afín. En el caso en que~$S'$ también es afín, que es el más frecuente en las
aplicaciones, se tiene:

\begin{corollaire}
\label{VIII.7.4}
Supongamos que~$S$ y~$S'$ son afines. Para que el dato de descenso sobre~$X'$
sea efectivo, es necesario y suficiente que~$X'$ sea reunión de
abiertos~$X_{i}'$ afines y estables por el dato de descenso.
\end{corollaire}

La suficiencia proviene de~\Ref{VIII.7.2} y del hecho de que, si~$X_{i}'$ es
afín, es afín sobre~$S'$ y puede aplicarse~\Ref{VIII.2.1}. Para la necesidad,
obsérvese que, si~$X'$ proviene de~$X$ y~$X$ está recubierto por abiertos
afines~$X_{i}$, entonces los~$X_{i}'=X_{i}\times_{S}S'$ son abiertos afines
estables por los datos de descenso y recubren~$X'$.

\begin{corollaire}
\label{VIII.7.5}
Sea~$g\colon S'\to S$ un morfismo fielmente plano, cuasicompacto y
\emph{radicial}. Entonces~$g$ es un morfismo de descenso \emph{efectivo},
\ie para todo~$X'$ sobre~$S'$, todo dato de descenso sobre~$X'$ con respecto
a~$g\colon S'\to S$ es efectivo.
\end{corollaire}

En efecto, en virtud de~\Ref{VIII.7.3} puede suponerse que~$S$ es afín; como
~$S'$ es radicial sobre~$S$ y, por tanto, separado, $S'$ es separado. Por
otra parte, para todo~$x'\in X'$, la fibra
~$R(x')=q_{2}(q_{1}^{-1}(x'))$ de la relación de equivalencia conjuntista
definida
\marginpar{222}
por la relación de equivalencia~$R$ se reduce a un punto. En efecto, como~$g$
es radicial, lo mismo sucede con~$p_{1},p_{2}$, que se deducen de aquel por
el cambio de base~$S'\to S$, y, por tanto, también con~$q_{1},q_{2}$, que se
deducen de los anteriores por el cambio de base~$X'\to S''$. Así,
\emph{todo abierto de~$X'$ es estable} por el dato de descenso. Recubramos
entonces~$X'$ por abiertos afines~$X_{i}'$; son afines
\ifthenelse{\boolean{orig}}
{sobre~$S$}
{sobre~$S'$}
puesto que~$S'$ es separado, por lo que el dato de descenso inducido es
efectivo según~\Ref{VIII.2.1}. Se concluye entonces mediante~\Ref{VIII.7.2}.

Obsérvese que~\Ref{VIII.7.5} da el único caso conocido de un morfismo de
descenso efectivo en la categoría de los preesquemas, y probablemente sea
en efecto el único caso, incluso limitándose a los esquemas noetherianos o a
los esquemas de tipo finito sobre un cuerpo.

Cuando se supone~$S$ localmente noetheriano y~$S'$ de tipo finito sobre~$S$,
el enunciado~\Ref{VIII.7.5} es también un caso particular del siguiente (que
generaliza el descenso galoisiano de Weil y el descenso inseparable de
Cartier):

\begin{corollaire}
\label{VIII.7.6}
Sea~$g\colon S'\to S$ un morfismo finito localmente libre (\ie definido por
un álgebra sobre~$S$ que es un módulo localmente libre de tipo finito) y
sobreyectivo (por tanto, $g$ es fielmente plano y cuasicompacto, luego es un
morfismo de descenso). Sea~$X'$ un~$S'$-preesquema provisto de un dato de
descenso. Para que este dato sea efectivo, es necesario y suficiente que,
para todo~$x'\in X'$, la fibra~$R(x')=q_{2}(q_{1}^{-1}(x'))$ esté contenida
en un abierto afín (condición que se satisface automáticamente si~$X'$ es
cuasiproyectivo sobre~$S'$).
\end{corollaire}

La observación entre paréntesis proviene del hecho de que, si~$s$ es el
punto de~$S$ que es imagen de~$x'$, entonces~$R(x')$ es finito y está
contenido en la
\ifthenelse{\boolean{orig}}
{fibra~$X_{s}$,}
{fibra~$X'_{s}$,}
por otra parte, como~$X'$ es cuasiproyectivo sobre~$S'$ y~$S'$ es finito
sobre~$S$, $X'$ es cuasiproyectivo sobre~$S$, lo cual implica que una fibra
\ifthenelse{\boolean{orig}}
{de~$X/S$}
{de~$X'$ sobre~$S$}
está contenida en un abierto afín.

Como todo subconjunto finito de un esquema afín admite un sistema
fundamental de entornos afines, se ve que no se pierde la hipótesis al
restringirse sobre un abierto afín de~$S$; esto, en virtud
de~\Ref{VIII.7.3}, nos reduce al caso en que~$S$ es afín. En virtud
de~\Ref{VIII.7.4}, queda por demostrar
\ifthenelse{\boolean{orig}}
{que~$X'$}
{que~$x'$}
está contenido
\marginpar{223}
en un abierto afín \emph{estable} por el dato de descenso. Sea, en efecto,
~$U$ un abierto afín que contiene~$R(x')$; entonces el saturado
$$
\ifthenelse{\boolean{orig}}
{R(X-U)=q_{2}(q_{1}^{-1}(X-U))}
{R(X'-U)=q_{2}(q_{1}^{-1}(X'-U))}
$$
no corta a~$R(x')$. Por otra parte, como~$q_{2}$ es finito (pues~$g$ y, por
tanto, $p_{2}$ lo son) y, por consiguiente, cerrado, el segundo miembro es
un subconjunto cerrado de~$X'$. Sea~$U'$ su complementario en~$X'$; es,
pues, un abierto \emph{saturado}, y se tiene
$$
R(x')\subset U'\subset U
$$
con~$U$ afín, aunque~$U'$ no es afín \emph{a priori}. Como el subconjunto
finito~$R(x')$ del esquema afín~$U$ posee un sistema fundamental de entornos
afines de la forma~$U_{f}$, se ve, reemplazando~$f$ por su restricción a
~$U'$, que existe una sección~$f$ de~$\mathcal{O}_{U}$ tal que:
$$
R(x')\subset U_{f}',\quad U_{f}'\textrm{ es afín}.
$$
Sea entonces~$U''=q_{1}^{-1}(U')=q_{2}^{-1}(U')$; denotemos todavía
por~$q_{1},q_{2}$ los morfismos inducidos~$U''\to U'$ y consideremos
$$
f'=\Norm_{q_{2}}(q_{1}^*(f))\quoi,
$$
donde~$\Norm_{q_{2}}$ designa la \emph{norma} con respecto al morfismo
finito localmente libre~$q_{2}\colon U''\to U'$. La compatibilidad de la
formación de la norma con el cambio de base implica fácilmente que~$f'$ es
una sección \emph{invariante}:
$$
q_{1}^*(f')=q_{2}^*(f')
$$
lo cual implica que~$U_{f'}'$ es un abierto saturado de~$U'$. Más
precisamente, si~$Z(f')$ designa el conjunto de ceros de una sección~$f'$,
se obtiene, en virtud de las propiedades de las normas:
$$
Z(f')=q_{2}(Z(q_{1}^*(f)))=q_{2}(q_{1}^*(Z(f)))=
R(U'-U_{f}')\quoi,
$$
lo cual
\marginpar{224}
implica que~$U_{f'}'=U'-Z(f')$ es saturado, contiene~$R(x')$ y está
contenido en~$U_{f}'$. Como este último es afín, se sigue que~$U_{f'}'$ lo
es también (pues es igual a~$(U_{f}')_{f''}$, con
~$f''=f'|U_{f'}'$). Es, por tanto, un abierto afín saturado que contiene
~$R(x')$ y, por consiguiente, a~$x'$, lo que concluye la demostración.
\par

Obsérvese que este razonamiento se aplica siempre que se tiene una relación
de equivalencia (o incluso solo una preequivalencia, \cf \cite{VIII.3}) en un
preesquema~$X'$, 
\ifthenelse{\boolean{orig}}
{finita y localmente libre;}
{finita y localmente libre;}
por lo demás, \Ref{VIII.7.6} es también un caso particular del resultado
análogo para las preequivalencias finitas y localmente libres, \cf \loccit
La misma observación vale para~\Ref{VIII.7.7} más abajo.

Una vez obtenida la existencia de un abierto cuasiafín saturado~$U'$ que
contiene~$x'$, también puede recurrirse a~\Ref{VIII.7.9}
y~\Ref{VIII.7.2}, con lo cual se evita el uso de las normas.

Obsérvese además que, en las condiciones de~\Ref{VIII.7.6}, si el dato de
descenso sobre~$X'$ es efectivo y~$X'$ proviene de~$X$ sobre~$S$, entonces
el morfismo~$X'\to X$ es finito, localmente libre y sobreyectivo, pues se
deduce de~$g$ por el cambio de base~$X\to S$. Se sigue (EGA~II 6.6.4) que,
si~$X'$ es cuasiproyectivo sobre~$S'$ y, por tanto, sobre~$S$, entonces~$X$
es cuasiproyectivo sobre~$S$ (se obtiene un haz invertible relativamente
amplio sobre~$X$ tomando la \emph{norma} de un haz invertible sobre~$X'$
relativamente amplio sobre~$S$, o sobre~$S'$, lo que viene a ser lo mismo).
Se obtiene así:

\begin{corollaire}
\label{VIII.7.7}
Un morfismo~$g\colon S'\to S$ finito localmente libre y sobreyectivo es un
morfismo de descenso efectivo para la categoría fibrada de los preesquemas
cuasiproyectivos sobre otros; \ie para todo~$X'$ cuasiproyectivo sobre~$S'$,
todo dato de descenso sobre~$X'$ con respecto a~$g$ es efectivo, y el
$S$-preesquema descendido~$X$ es cuasiproyectivo sobre~$S$.
\end{corollaire}

\begin{proposition}
\label{VIII.7.8}
Sea~$g\colon S'\to S$ un morfismo fielmente plano y cuasicompacto. Entonces
~$g$ es un morfismo de descenso efectivo para la categoría fibrada de los
preesquemas~$Z$ cuasicompactos sobre un preesquema~$T$, provistos de un haz
invertible amplio
\marginpar{225}
con respecto a~$T$. En particular, para todo preesquema~$X'$ sobre~$S'$,
provisto de un dato de descenso con respecto a~$g\colon S'\to S$, y todo haz
invertible~$\mathcal{L}'$ sobre~$X'$ amplio con respecto a~$S'$, provisto
también de un dato de descenso con respecto al dado sobre~$X'$ (\ie provisto
de un isomorfismo entre~$q_{1}^*(\mathcal{L}')$ y~$q_{2}^*(\mathcal{L}')$ que
satisface la condición habitual de transitividad). Entonces el dato de
descenso sobre~$X'$ es efectivo, y el haz invertible~$\mathcal{L}$ sobre el
preesquema descendido~$X$, deducido de~$\mathcal{L}'$ por descenso, es
amplio con respecto a~$S$.
\end{proposition}

La demostración es enteramente análoga a la de~\Ref{VIII.5.8}. Obsérvese
que, en el álgebra graduada cuasicoherente~$\mathcal{S}'$ sobre~$S'$
definida por~$\mathcal{L}'$, hay un dato de descenso que permite construir,
gracias a~\Ref{VIII.1.1}, un álgebra graduada
cuasicoherente~$\mathcal{S}$ sobre~$S$, de donde resulta un
~$P=\Proj(\mathcal{S})$ sobre~$S$ tal que~$P'=\Proj(\mathcal{S}')$, junto con
su dato de descenso, se identifica con~$P\times_{S}S'$. Como, por hipótesis,
~$X'$ se identifica con un abierto de~$P'$, necesariamente estable por el
dato de descenso sobre~$P'$, el dato de descenso sobre~$X'$ es también
efectivo y se obtiene el preesquema descendido como un abierto de~$P$. Se
deja el detalle al lector. --- En particular, tomando
~$\mathcal{L}'=\mathcal{O}_{X'}$, se obtiene:

\begin{corollaire}
\label{VIII.7.9}
Sea~$g\colon S'\to S$ un morfismo fielmente plano y cuasicompacto, y sea
~$X'$ un preesquema \emph{cuasiafín} sobre~$S'$. Entonces todo dato de
descenso sobre~$X'$ con respecto a~$g$ es efectivo, y el preesquema
descendido~$X$ es cuasiafín sobre~$S$.
\end{corollaire}

En virtud de~\Ref{VIII.6.2}, este resultado se aplica, en particular, si~$S'$
es localmente noetheriano y~$X'$ es cuasifinito y separado sobre~$S'$; más
generalmente, si~$S'$ es cualquiera y~$X'$ es de presentación finita,
cuasifinito y separado sobre~$S'$ (\cf \Ref{VIII.6.6}).

\begin{remarques}
\label{VIII.7.10}
Los resultados dados en este número agotan los criterios de efectividad
conocidos actualmente, y probablemente incluso los criterios útiles
existentes\footnote{Esta opinión resultó parcialmente errónea; véanse, por
ejemplo, J.P\ptbl \textsc{Murre}, Sém. Bourbaki 294 (Apéndice), mayo de 1965,
y resultados especiales (en particular, de \textsc{Néron} y
\textsc{Raynaud}). Para el descenso de los esquemas en grupos, \cf
M\ptbl Raynaud, tesis (1968).}. Obsérvense los contraejemplos siguientes en
apoyo de esta afirmación:
\begin{enumerate}
\item[(i)]
\marginpar{226}
Si~$S$ es el espectro de un cuerpo y~$S'$ es el espectro de una extensión
cuadrática galoisiana, puede encontrarse un~$X'$ sobre~$S'$, propio y liso
sobre~$S'$, de dimensión 3, provisto de un dato de descenso que no es
efectivo (Serre).
\item[(ii)] Puede encontrarse un~$S$ que sea el espectro de un anillo local regular de
dimensión 3 (si se quiere, el anillo local de un esquema algebraico sobre un
cuerpo de característica dada), un recubrimiento principal~$T$ de~$S$ con
grupo~$\ZZ/2\ZZ$, tal que, si~$t$ designa uno de los puntos de~$T$ situados
sobre el punto cerrado~$s$ de~$S$, y
\ifthenelse{\boolean{orig}}
{$S'=T-s$,}
{$S'=T-\{t\}$,}
pueda encontrarse un~$X'$ \emph{proyectivo} sobre~$S'$, regular, provisto de
un dato de descenso con respecto a~$g\colon S'\to S$, sin que este dato de
descenso sea efectivo.
\end{enumerate}

Para estas construcciones se utiliza el ejemplo de Hironaka de variedades
no proyectivas. Para~(i), basta utilizar el hecho de que puede encontrarse
sobre~$k$ un esquema propio y liso~$X_{0}$ de dimensión 3 sobre el cual
~$G=\ZZ/2\ZZ$ actúa sin inercia y en el que existen dos puntos~$a,b$
racionales
\ifthenelse{\boolean{orig}}
{sobre~$P$,}
{sobre~$S$,}
conjugados por la acción de~$G$, que no están contenidos en un abierto afín. Se pone
entonces~$X'=X_{0}\times_{k}k'$ y se hace actuar~$G$ sobre~$X'$ mediante las
acciones de~$G$ sobre los dos factores, lo que da un dato de descenso
sobre~$X'$ con respecto a~$g\colon \Spec(k')\to \Spec(k)$. Por encima de~$a$
\resp de~$b$, hay exactamente un punto~$a'$ \resp $b'$ (con extensión residual
cuadrática), y~$a'$ y~$b'$ son conjugados por la acción de~$G$, pues~$X'\to X_{0}$ es
compatible con las acciones de~$G$. Entonces~$a'$ y~$b'$ no pueden estar
contenidos en un abierto afín~$U'$, pues, de lo contrario,
~$U=X_{0}-\Im(X'-U')$ sería un abierto de~$X_{0}$ que contendría~$(a,b)$ y
cuya imagen inversa en~$X'$ estaría contenida en~$U'$ y sería, por tanto,
cuasiafín; luego~$U$ sería cuasiafín y, en consecuencia,~$(a,b)$ tendría un
entorno afín en~$U$.

Para~(ii), se utiliza el hecho de que, en el ejemplo de Hironaka, $X_{0}$ se
obtiene como preesquema propio sobre un~$k$-esquema proyectivo~$Y$, liso
sobre~$k$ (el morfismo~$f\colon X_{0}\to Y$ es, por lo demás, birracional,
pero poco importa); el grupo~$G$ actúa asimismo sobre~$Y$ de manera
compatible con sus acciones sobre~$X_{0}$; finalmente, poniendo
\ifthenelse{\boolean{orig}}
{$S'=Y-f(b)$,}
{$S'=Y-\{f(b)\}$,}
~$X'=X_{0}|S'$, $X'$ es proyectivo sobre~$S'$. Entonces~$X_{0}$ está
provisto de un dato de descenso natural con respecto al morfismo
\marginpar{227}
canónico~$Y\to S=Y/G$, gracias a las acciones de~$G$ sobre~$X_{0}$
compatibles con sus acciones sobre~$Y$. Este dato de descenso no es efectivo,
pues~$(a,b)$ no está contenido en un abierto afín.
\ifthenelse{\boolean{orig}}
{El dato de descenso inducido}
{El dato de descenso inducido}
sobre~$X'$ con respecto a~$g\colon S'\to S$ tampoco es entonces efectivo,
como se comprueba fácilmente.
\end{remarques}


\begin{sgabibliography}{D}{VIII.8}

\bibitem[D]{VIII.D} J\ptbl \textsc{Giraud}, Méthode de la descente, Mémoire \no 2
de la Soc.\ Math.\ Fr., 1964.

\bibitem[1]{VIII.1} A\ptbl \textsc{Grothendieck}, Séminaire Bourbaki:
Géométrie formelle et Géométrie algébrique, mayo de~1959, \No 182.

\bibitem[2]{VIII.2} A\ptbl \textsc{Grothendieck}, Séminaire Bourbaki: Technique de
descente et Théorèmes d'existence~I, diciembre de~1959, \No 190.

\bibitem[3]{VIII.3} A\ptbl \textsc{Grothendieck}, Séminaire Bourbaki: Technique de
descente et Théorèmes d'existence~III, febrero de~1961, \No 212.
\end{sgabibliography}

\chapter[Descenso de los morfismos étales]{Descenso de los morfismos étales.\\ Aplicación al grupo fundamental}
\label{IX}
\marginpar{228}

\section{Recordatorio sobre los morfismos étales}
\label{IX.1}

Vamos a repasar aquí las propiedades de los morfismos étales (desarrolladas
en~\Ref{I}) que vamos a utilizar, aprovechando la ocasión para eliminar de
la teoría las hipótesis noetherianas superfluas. El lector observará que,
aun cuando solo interesen los esquemas noetherianos, la técnica del descenso
lleva a introducir esquemas no noetherianos (tales como
$\Spec(\widehat{A}\otimes_{A}\widehat{A})$, donde~$A$ es un anillo local
noetheriano); para poder aplicar el lenguaje de las categorías fibradas,
importa definir nociones como la de morfismo étale,
\ifthenelse{\boolean{orig}}{etc...}{etc.}, sin introducir en ellas ninguna
restricción noetheriana.
El lector que no desee comprobar ni admitir que los enunciados siguientes
son verdaderos sin hipótesis noetherianas podrá limitarse a admitirlos bajo
las hipótesis noetherianas de~\Ref{I}, con tal de introducir esas mismas
hipótesis noetherianas en los enunciados de los números siguientes y de
utilizar la definición~\Ref{IX.1.1} que se da a continuación para los
esquemas no noetherianos que aparecen en los razonamientos.

\begin{definition}
\label{IX.1.1}
Sean~$f\colon X\to S$ un morfismo de preesquemas y~$x$ un punto de~$X$.
Se dice que~$f$ es \emph{étale en}~$x$,
\index{étale (morfismo, álgebra)|hyperpage}%
o que~$X$ es étale sobre~$S$ en~$x$, si existen un entorno abierto
afín~$U$ de~$s=f(x)$, un entorno abierto afín~$V$ de~$x$ situado sobre~$U$,
un esquema afín \emph{noetheriano}~$U_{0}$, un~$U_{0}$-esquema
afín~$V_{0}$, \emph{étale} en el sentido de~\eqref{I}, un
morfismo~$U\to U_{0}$ y un
$U$-isomorfismo
$$
V\isomto V_{0}\times_{U_{0}}U\quoi.
$$
\end{definition}

Se observará
\marginpar{229}
que, cuando~$S$ es localmente noetheriano, esta terminología coincide con
la de~\loccit Se dirá asimismo que~$f$ \emph{es étale}, o que~$X$
\emph{es étale sobre}~$S$, si~$f$ es étale en todo punto~$x$ de~$X$.
Con estas definiciones, las proposiciones siguientes se reducen sin
dificultad al caso noetheriano, donde se demuestran en~\Ref{I.4},
\Ref{I.5} y~\Ref{I.7}. Para más detalles, el lector podrá consultar
EGA~IV\footnote{Más precisamente, EGA~IV 17,~18.}.

\begin{remarques}
\label{IX.1.2}
Si~$f$ es étale en~$x$, entonces~$f$ es «\emph{de presentación finita
en}~$x$» (\Ref{VIII.3.5}), el anillo local de~$x$ en la fibra
$f^{-1}(s)$ es una \emph{extensión finita separable} de~$\kres(s)$ y,
finalmente, $f$ es \emph{plano en}~$x$. Puede demostrarse que la recíproca
es cierta; por tanto, la definición~\Ref{IX.1.1} es la misma que cuando~$S$
es localmente noetheriano, salvo que debe reemplazarse la condición «de
tipo finito en~$x$» por «de presentación finita en~$x$». Como la
demostración de este resultado es delicada, no hemos querido adoptar aquí
esta definición de la noción de morfismo étale, que no se presta
directamente a la demostración de las propiedades que siguen.
\end{remarques}

Observemos primero que se tiene trivialmente:
\begin{proposition}
\label{IX.1.3}
Si~$f\colon X\to S$ es étale, entonces todo morfismo~$f'\colon X'\to S'$
que se deduce de él mediante un cambio de base~$S'\to S$ es también étale.
\end{proposition}

Puede decirse, por tanto, que los morfismos étales forman una
\emph{subcategoría fibrada} de la categoría de las flechas en~$\Sch$
\ifthenelse{\boolean{orig}}
{(\cf \Ref{VI.11}~a)}
{(\cf \Ref{VI.11}~a))}. El objeto de la presente exposición es estudiar
las propiedades de exactitud de esta categoría fibrada sobre~$\Sch$.
\begin{proposition}
\label{IX.1.4}
Sea~$f\colon X\to S$ un morfismo de preesquemas. Para que sea una inmersión
abierta es necesario y suficiente que sea étale y radicial.
\end{proposition}

\Cf \Ref{I.5.1}. Se deduce que, si~$X$ es étale sobre~$S$, toda sección
de~$X$ sobre~$S$ es una inmersión abierta; utilizando de nuevo~\Ref{IX.1.4},
se obtiene:
\begin{corollaire}
\label{IX.1.5}
Sea~$X$ un $S$-preesquema étale. Entonces existe una correspondencia
biunívoca entre el conjunto de las secciones de~$X$ sobre~$S$ y el conjunto
de las partes
\marginpar{230}
abiertas~$\Gamma$ de~$X$ tales que el morfismo~$\Gamma\to S$ inducido por
el morfismo estructural sea \emph{radicial} y \emph{sobreyectivo}.
\end{corollaire}

Si, además, $X$ está separado sobre~$S$, $\Gamma$ será una parte de~$X$ a la
vez abierta y cerrada, pero esto poco importa. --- Efectuando un cambio de
base evidente, puede ponerse~\Ref{IX.1.5} bajo la forma aparentemente más
general siguiente:
\begin{corollaire}
\label{IX.1.6}
Sean~$X$ e~$Y$ dos $S$-preesquemas, siendo~$Y$ étale sobre~$S$. Entonces la
aplicación~$f\mto\Gamma_f$ que asocia a todo $S$-morfismo~$f$ de~$X$ en~$Y$
la parte de~$X\times_S Y$ subyacente al grafo de~$f$ es una biyección de
$\Hom_S(X,Y)$ sobre
\ifthenelse{\boolean{orig}}
{el conjunto de las partes abiertas de~$\Gamma$ de~$X\times_S Y$}
{el conjunto de las partes abiertas~$\Gamma$ de~$X\times_S Y$}
tales que el morfismo~$\Gamma\to X$ inducido por~$\pr_1$ sea
\emph{radicial} y \emph{sobreyectivo}.
\end{corollaire}

\begin{proposition}
\label{IX.1.7}
Sea~$S_0$ el subpreesquema de~$S$ definido por un ideal cuasicoherente
formado por elementos nilpotentes, es decir, tal que~$S_0$ tenga el mismo
conjunto subyacente que~$S$. Entonces el funtor~$X\mto X\times_S S_0$ de
la categoría de los preesquemas étales sobre~$S$ a la categoría de los
preesquemas étales sobre~$S_0$ es una equivalencia de categorías.
\end{proposition}

El hecho de que este funtor sea plenamente fiel es una consecuencia
inmediata de~\Ref{IX.1.6}. El hecho de que sea esencialmente sobreyectivo
está contenido en~\Ref{I.8.3}. Se observará que, en la equivalencia
anterior, $X$ es de tipo finito, es decir, cuasifinito sobre~$S$
(respectivamente finito, es decir, un recubrimiento étale de~$S$), si y
solo si~$X_0$ satisface la condición análoga sobre~$S_0$; la misma
observación vale para la condición de separación. Estos hechos son
inmediatos y están contenidos también en~\Ref{IX.2.4} más abajo.

\begin{corollaire}
\label{IX.1.8}
Sea~$A$ un anillo local noetheriano completo de cuerpo residual~$k$.
Entonces el funtor~$B\mto B\otimes_A k$ es una equivalencia entre la
categoría de las álgebras finitas y étales sobre~$A$ y la categoría de las
álgebras finitas y étales sobre~$k$ (es decir, productos finitos de
extensiones finitas separables de~$k$).
\end{corollaire}

\begin{proposition}
\label{IX.1.9}
Para que~$X$ sea un \emph{recubrimiento étale} de~$S$, es decir, finito y
étale sobre~$S$, es necesario y suficiente que~$X$ sea $S$-isomorfo al
espectro de un álgebra~$\cal{A}$ sobre~$S$ que sea un módulo localmente
libre de tipo finito y tal que, para todo
\marginpar{231}
$s\in S$, $\cal{A}_s\otimes_{\cal{O}_s}\kres(s)$ sea un álgebra separable
sobre~$\kres(s)$ y, por tanto, en este caso un producto directo de
extensiones finitas separables de~$\kres(s)$.
\end{proposition}

Por último, el resultado siguiente es de naturaleza menos elemental, pues
es la conjunción de~\Ref{I.8.4} y del \emph{teorema de existencia para
haces en geometría algebraica} (EGA~III~5; \cf también \cite{IX.1},
teor\ptbl 3).

\begin{theoreme}
\label{IX.1.10}
Sean~$S$ el espectro de un anillo local noetheriano completo, $X$ un
$S$-esquema propio y~$X_0$ la fibra de~$X$ en el punto cerrado de~$S$ (de
modo que~$X_0$ es un subesquema cerrado de~$X$). Entonces el funtor de
restricción~$X'\mto X'\times_X X_0$ es una equivalencia entre la categoría
de los recubrimientos étales de~$X$ y la categoría de los recubrimientos
étales de~$X_0$.
\end{theoreme}

\section{Morfismos submersivos y universalmente submersivos}
\label{IX.2}

\begin{definition}
\label{IX.2.1}
Un morfismo
\ifthenelse{\boolean{orig}}
{$g\colon S'\to S$}
{$f\colon S'\to S$}
de preesquemas se dice \emph{submersivo}
\index{submersivo (morfismo)|hyperpage}%
si es sobreyectivo y hace de~$S$ un espacio topológico cociente
de~$S'$ (es decir, toda parte~$U$ de~$S$ tal que $f^{-1}(U)$ sea
abierta es abierta). Se dice que~$f$ es \emph{universalmente
submersivo}
\index{universalmente submersivo (morfismo)|hyperpage}%
si, para todo morfismo~$T\to S$, el morfismo
$f'\colon T'=S'\times_S T\to T$ deducido de~$f$ por cambio de base es
submersivo.
\end{definition}

Es inmediato que el compuesto de dos morfismos submersivos
(respectivamente universalmente submersivos) es submersivo
(respectivamente universalmente submersivo), y que un cambio de base de un
morfismo universalmente submersivo da un morfismo universalmente submersivo
(pues la definición está hecha precisamente para ello). Si~$fg$ es
submersivo (respectivamente universalmente submersivo), también lo es~$f$.

\begin{exemples}
\label{IX.2.2}
a) Un morfismo sobreyectivo que sea abierto o cerrado es submersivo; por
tanto, un morfismo sobreyectivo universalmente cerrado o universalmente
abierto es universalmente submersivo. Por ejemplo, \emph{un morfismo propio
sobreyectivo es universalmente submersivo}. Por otra parte, \emph{un
morfismo fielmente plano y cuasicompacto es universalmente submersivo}
(\Ref{VIII.4.3}). Estos serán para nosotros los dos casos más importantes.
\end{exemples}

Pueden aplicarse
\marginpar{232}
a un morfismo submersivo o universalmente submersivo~$g\colon S'\to S$ los
razonamientos de~\Ref{VIII.4.3}; se obtiene, en particular:
\begin{proposition}
\label{IX.2.3}
Supongamos que~$g\colon S'\to S$ es submersivo. Entonces el siguiente
diagrama de aplicaciones es exacto:
$$
\xymatrix@C=.5cm{\Ouv (S)\ar[r] &
\Ouv (S')\ar@/^2pt/[r]\ar@/_2pt/[r] & \Ouv(S''),}
$$
donde~$S''=S'\times_S S'$ y~$\Ouv(X)$ designa el conjunto de las partes
abiertas del preesquema~$X$.
\end{proposition}

\begin{proposition}
\label{IX.2.4}
Sean~$g\colon S'\to S$ un morfismo universalmente submersivo,
$f\colon X\to Y$ un $S$-morfismo y~$f'\colon X'\to Y'$ el
$S'$-morfismo que se deduce de él por cambio de base. Para que~$f$ sea
\ifthenelse{\boolean{orig}}
{abierta (respectivamente cerrada)}
{abierto (respectivamente cerrado)}, basta con que~$f'$ lo sea. Para que~$f$
sea \ifthenelse{\boolean{orig}}
{universalmente abierta, respectivamente universalmente cerrada,
respectivamente separada}
{universalmente abierto, respectivamente universalmente cerrado,
respectivamente separado}, es necesario y suficiente que~$f'$ lo sea. Si,
además, $g$ es cuasicompacto y~$f$ localmente de tipo finito, para que~$f$
sea propio es necesario y suficiente que~$f'$ lo sea.
\end{proposition}

Para este último punto, obsérvese que, si~$f'$ es propio y, por tanto,
cuasicompacto, entonces~$f$ es cuasicompacto (\Ref{VIII.3.3}) y, por
consiguiente, de tipo finito, puesto que es localmente de tipo finito. Por
otra parte, es separado y universalmente cerrado por lo que precede; luego
es propio.

\begin{proposition}
\label{IX.2.5}
Sea~$S'$ un preesquema de tipo finito sobre el espectro~$S$ de un anillo
local noetheriano completo. Supongamos que la fibra del punto cerrado~$s$
de~$S$ sea finita, de modo que los anillos locales de~$S'$ en los puntos
$s'$ de esta fibra sean finitos sobre~$A=\cal{O}_s$. Sea~$S''$ el esquema
suma disjunta de los espectros de los~$\cal{O}_{S',s'}$ en cuestión,
considerado como $S$-esquema finito. Para que~$g\colon S'\to S$ sea
universalmente submersivo es necesario y suficiente que el morfismo
estructural~$S''\to S'$ sea sobreyectivo.
\end{proposition}

Como existe un $S$-morfismo natural~$S''\to S'$ y todo morfismo finito
sobreyectivo es universalmente submersivo por~\Ref{IX.2.2}, la condición
enunciada es suficiente. Demostremos, pues, que, si~$S''\to S'$ no es
sobreyectivo, entonces~$g$ no es universalmente submersivo. En efecto, sea
$t$ un punto de~$S$ que no pertenece a la
\marginpar{233}
imagen de~$S''$; existe entonces
un $S$-esquema~$T$, espectro de un anillo de valoración discreta, cuya
imagen en~$S$ es~$\{s,t\}$. Obsérvese que la imagen de~$S''$ en~$S'$ es
abierta, pues el morfismo~$S''\to S'$ es un isomorfismo local; además, esta
imagen contiene~$S'_s$ y no corta a~$S'_t$. Se sigue que la preimagen de
esta última en~$T'=S'\times_S T$ es \emph{abierta} y coincide con la
preimagen del punto cerrado de~$T$. Esto
demuestra que~$T'\to T$ no es submersivo y, por tanto, que~$S'\to S$ no es
universalmente submersivo.

\begin{remarque}
\label{IX.2.6}
Utilizando el criterio~\Ref{IV.6.3} para que una parte constructible de un
espacio noetheriano sea abierta, se obtiene fácilmente el siguiente criterio
valuativo para que un morfismo~$g\colon S'\to S$ \emph{de tipo finito}, con
$S$ localmente noetheriano, sea universalmente submersivo: es necesario y
suficiente que, para todo $S$-esquema~$T$ que sea el espectro de un anillo de
valoración discreta, poniendo~$T'=S'\times_S T$, la preimagen en~$T'$ del
punto cerrado de~$T$ no sea abierta.
\end{remarque}

\section{Descenso de morfismos de preesquemas étales}
\label{IX.3}

\begin{proposition}
\label{IX.3.1}
Sean~$g\colon S'\to S$ un morfismo \emph{sobreyectivo} de preesquemas,
$X$ e~$Y$ dos preesquemas sobre~$S$, y~$X'$, $Y'$ sus preimágenes
sobre~$S'$. Si~$Y$ es no ramificado sobre~$S$, entonces la aplicación
canónica
$$
\Hom_S(X,Y) \to \Hom_{S'}(X',Y')
$$
es inyectiva.
\end{proposition}

En efecto, en virtud de~\Ref{IX.1.6}, un $S$-morfismo
$f\colon X\to Y$ queda determinado cuando se conoce el conjunto subyacente
a su grafo~$\Gamma$, que es una parte de~$Z=X\times_SY$. Como
$$
Z'=Z\times_SS'=X'\times_{S'}Y' \to Z
$$
es sobreyectivo (pues~$S'\to S$ lo es), esta parte~$\Gamma$ queda
determinada cuando se conoce su preimagen en~$X'\times_{S'}Y'$, que no es
otra que el conjunto subyacente al grafo de~$f'$. De ahí la conclusión.

Una parte~$\Gamma$ de~$Z$ es el grafo de un $S$-morfismo
$f\colon X\to Y$ si y solo si es abierta y si el morfismo inducido por
$\pr_1$ de~$\Gamma$ en
\ifthenelse{\boolean{orig}}
{$Y$}
{$X$}
es radicial
\marginpar{234}
y sobreyectivo; \cf \Ref{IX.1}. Cuando se verifica la primera propiedad,
la segunda se verifica si y solo si la preimagen~$\Gamma'$ de~$\Gamma$ en
$Z'$ satisface la misma condición (\Ref{VIII.3.1}). Si, por último, se sabe
que~$Z'\to Z$ es submersivo, lo que sucede en particular si~$S'\to S$ es
universalmente submersivo, entonces~$\Gamma$ es abierta si y solo si lo
es~$\Gamma'$.

Así, el conjunto~$\Hom_S(X,Y)$ está entonces en correspondencia biunívoca
con el conjunto de las partes abiertas~$\Gamma'$ de~$Z'$ tales que el
morfismo de proyección
\ifthenelse{\boolean{orig}}
{$\pr_1\colon Z'\to X'$}
{$\pr_1\colon \Gamma'\to X'$}
sea radicial y sobreyectivo (es decir, que correspondan a un $S'$-morfismo
$f'\colon X'\to Y'$), y que estén saturadas para la relación de
equivalencia definida por~$Z'\to Z$; es decir, cuyas dos preimágenes en
$Z''=Z'\times_Z Z'=Z\times_S S''$ (donde~$S''=S'\times_S S'$), por una y
otra proyección, sean iguales. Ahora bien, estas últimas son los grafos de
los dos $S''$-morfismos~$X''\to Y''$ deducidos de~$f'$ por cambio de base,
mediante una y otra proyección~$S''\to S'$. Hemos obtenido así:

\begin{proposition}
\label{IX.3.2}
Sean~$g\colon S'\to S$ un morfismo \emph{universalmente submersivo} de
preesquemas, $S''=S'\times_S S'$, $X$ e~$Y$ dos $S$-preesquemas, $X'$ e
$Y'$ sus preimágenes sobre~$S'$, y~$X''$, $Y''$ sus preimágenes
sobre~$S''$. Si~$Y$ es étale sobre~$S$, el siguiente diagrama canónico de
aplicaciones es exacto:
$$
\xymatrix@C=.5cm{ {\Hom_S(X,Y)} \ar[r] & {\Hom_{S'}(X',Y')}
\ar@/^2pt/[r]\ar@/_2pt/[r] & {\Hom_{S''}(X'',Y'')}. }
$$
\end{proposition}

Tomando~$X$ e~$Y$ étales sobre~$S$, se obtiene el enunciado siguiente, que,
por lo demás, vuelve a dar~\Ref{IX.3.2} (incluso restringiéndose a~$X=S$,
caso al que siempre puede reducirse~\Ref{IX.3.2} mediante el cambio de base
$X\to S$):

\begin{corollaire}
\label{IX.3.3}
Un morfismo universalmente submersivo de preesquemas es un morfismo de
descenso para la categoría fibrada de los preesquemas étales sobre otros.
\end{corollaire}

Ignoro, por lo demás, si es necesariamente un morfismo de descenso
\emph{efectivo} para la categoría fibrada en cuestión, incluso añadiendo la
hipótesis de que~$S$ sea noetheriano y~$g$ de tipo finito, y limitándose a
los recubrimientos étales. No obstante, en el apartado siguiente daremos
criterios útiles de efectividad.

\begin{corollaire}
\label{IX.3.4}
Sea
\marginpar{235}
$g\colon S'\to S$ un morfismo universalmente submersivo cuyas fibras
$g^{-1}(s)$ sean \og geométricamente conexas\fg,
\index{geometricamente conexo@geométricamente conexo|hyperpage}%
es decir, tales que, para toda extensión~$K/\kres(s)$,
$g^{-1}(s)\otimes_{\kres(s)}K$ sea conexo. Entonces~$S'$ es conexo si~$S$
lo es. El funtor de la categoría de los preesquemas étales sobre~$S$ en la
\ifthenelse{\boolean{orig}}
{categorías}
{categoría}
de los preesquemas étales sobre~$S'$ definido por~$g$ es plenamente fiel.
\end{corollaire}

Una parte de~$S'$ que sea a la vez abierta y cerrada está saturada para la
relación de equivalencia conjuntista definida por~$g$, puesto que las
fibras son conexas; por tanto, es la preimagen de una parte de~$S$, que es
necesariamente abierta y cerrada porque~$g$ es submersivo. Así pues, si~$S$
es conexo, $S'$ lo es. Esto implica también el resultado siguiente: el
compuesto~$fg$ de dos morfismos con fibras universalmente conexas,
\index{universalmente conexo: sinónimo de geométricamente conexo|hyperpage}%
siendo~$f$ universalmente submersivo, tiene fibras universalmente conexas;
si~$S_1'$ y~$S_2'$ sobre~$S$ tienen fibras universalmente conexas, lo mismo
ocurre con~$S_1'\times_S S_2'$. En particular, bajo las condiciones
de~\Ref{IX.3.4}, $S''$ tiene fibras universalmente conexas sobre~$S$. Sean
entonces~$X$ e~$Y$ étales sobre~$S$, y sea~$u'$ un $S'$-morfismo de~$X'$ en
$Y'$; demostremos que es compatible con los datos de descenso (lo cual
entraña la conclusión buscada gracias a~\Ref{IX.3.3}). Sean, en efecto,
$u_1''$ y~$u_2''$ los dos $S''$-morfismos~$X''\to Y''$ deducidos de~$u'$.
El subpreesquema de~$S''$ de coincidencia de~$u_1''$ y~$u_2''$ es un
subpreesquema abierto inducido, cerrado fibra a fibra, como preimagen del
preesquema diagonal de~$Y''$ sobre~$S''$\kern1pt\footnote{Obsérvese que las
fibras de
\ifthenelse{\boolean{orig}}
{$S'$ sobre~$S$}
{$Y$ sobre~$S$}
están separadas.}. Por tanto, es la preimagen de una parte de~$S$. Como
contiene la diagonal en~$S''$, coincide con~$S''$, de donde
$u_1''=u_2''$, cqfd.

\section{Descenso de preesquemas étales: criterios de efectividad}
\label{IX.4}

\ifthenelse{\boolean{orig}}{}
{\enlargethispage{.7cm}}%
\begin{proposition}
\label{IX.4.1}
Sea $g\colon S'\to S$ un morfismo fielmente plano y cuasicompacto.
Entonces~$g$ es un morfismo de descenso efectivo para la categoría
fibrada de los preesquemas étales, separados y de tipo finito sobre otros.
\end{proposition}

En efecto, es un morfismo de descenso para la categoría fibrada en
cuestión, en virtud de~\Ref{IX.3.3} o de~\Ref{VIII.5.2}, a elección.
Queda por demostrar que, si~$X'$ es étale, separado y de tipo finito
sobre~$S'$, y está provisto de un dato de descenso respecto de
$g\colon S'\to S$, este dato es efectivo en la categoría fibrada en
cuestión. Ahora bien, se ve fácilmente que, si~$X$ es un preesquema
sobre~$S$, entonces
\marginpar{236}
es étale sobre~$S$ si y solo si lo es sobre~$S'$ (en virtud de la
Definición~\Ref{IX.1.1} y de \loccit \Ref{VIII.3.6}). Por tanto, es
étale, separado y de tipo finito sobre~$S$ si y solo si~$X'$ lo es
sobre~$S'$; \cf por ejemplo,~\Ref{IX.2.4}. Basta, pues, cerciorarse de
la efectividad del dato de descenso sobre
\ifthenelse{\boolean{orig}}
{$X$}
{$X'$}
para la categoría fibrada de las flechas de~$\Sch$. Ahora bien,~$X'$ es
cuasiafín sobre~$S'$ en virtud de~\Ref{VIII.6.2}
y~\Ref{VIII.6.6}. Se puede concluir entonces utilizando
\Ref{VIII.7.9}. El lector observará, además, que la demostración
exige menos si uno se limita a los preesquemas étales y \emph{finitos}
sobre otros, pues entonces se puede invocar directamente
\Ref{VIII.2.1}.

\begin{corollaire}
\label{IX.4.2}
Sean $g\colon S'\to S$ un morfismo universalmente submersivo, $X'$
\ifthenelse{\boolean{orig}}
{unos $S'$-preesquemas}
{un $S'$-preesquema}
étale, separado y de tipo finito, provisto de un dato de descenso respecto
de~$g$; $S_1\to S$ un morfismo fielmente plano y cuasicompacto; y~$S_1'$
y~$X_1'$ los preesquemas deducidos de~$S'$ y~$X'$ por cambio de base, de
modo que $S_1'\to S_1$ es universalmente submersivo, $X_1'$ es étale,
separado y de tipo finito sobre~$S_1'$, y está provisto de un dato de
descenso respecto de $g_1\colon S_1'\to S_1$. El dato de descenso
sobre~$X'$ es efectivo si y solo si lo es el dato de descenso sobre~$X_1'$.
\end{corollaire}

Esto resulta de la teoría del descenso en las categorías~\cite{IX.D},
teniendo en cuenta~\Ref{IX.4.1} y~\Ref{IX.3.3}.

De manera análoga se demuestra:

\begin{corollaire}
\label{IX.4.3}
Sean $g\colon S'\to S$ un morfismo universalmente submersivo, $X'$ un
$S'$-preesquema étale provisto de un dato de descenso respecto de~$g$, y
$(S_i)$ un recubrimiento de~$S$ por abiertos. El dato de descenso es
efectivo si y solo si, para todo~$i$, lo es el dato de descenso
correspondiente sobre
\ifthenelse{\boolean{orig}}
{$X_i'=X\times_S S_i$}
{$X_i'=X'\times_S S_i$},
respecto del morfismo
$g_i\colon S_i'=S'\times_S S_i\to S_i$.
\end{corollaire}

Este último resultado conduce a aislar un criterio local de efectividad:

\enlargethispage{.3cm}
\begin{proposition}
\label{IX.4.4}
Sea $g\colon S'\to S$ un morfismo de presentación finita en el sentido
de~\Ref{VIII.3.6},
\marginpar{237}
y universalmente submersivo; sea $X'$ un preesquema étale y de
presentación finita sobre~$S'$, provisto de un dato de descenso respecto
de~$g$; y sea, finalmente,~$a$ un punto de~$S$. Para que exista un entorno
abierto~$U$ de~$a$ tal que el dato de descenso correspondiente
sobre~$X_U'=X'\times_S U$ respecto del morfismo
$$g_U\colon S_U'=S'\times_S U\to S_U=U$$
sea efectivo, es necesario y suficiente que el dato de descenso
correspondiente sobre $X_a'=X'\times_S\Spec(\cal{O}_a)$ respecto del
morfismo
$$g_a\colon S_a'=S'\times_S\Spec(\cal{O}_a)\to
S_a=\Spec(\cal{O}_a),$$
sea efectivo.
\end{proposition}

La necesidad es trivial; demostremos la suficiencia. Se dispone, pues, de
un preesquema étale de tipo finito~$X_a$ sobre~$S_a$, y de un isomorfismo
\begin{equation*}
\label{eq:IX.4.*}
\tag{$*$} {X_a'\isomto X_a\times_{S_a}S_a'}
\end{equation*}
compatible con los datos de descenso. De acuerdo con un argumento general
elemental sobre los preesquemas definidos sobre un límite inductivo de
anillos (aquí, los anillos~$A_f$, donde~$A$ es el anillo de un entorno
abierto afín de~$a$ y~$f$ recorre los elementos de~$A$ que no pertenecen
al ideal primo correspondiente a~$a$), se pueden encontrar un entorno
abierto~$U$ de~$a$, un preesquema étale de tipo finito~$X_U$ sobre
$U=S_U$, y un $S_a$-isomorfismo
$X_a\isomto X_U\times_{S_U}S_a$. Además, tomando~$U$ suficientemente
pequeño, se puede suponer que el isomorfismo~\eqref{eq:IX.4.*} proviene de
un isomorfismo
$$
X_U'\isomto X_U\times_{S_U} S_U';
$$
Este último podría no ser compatible con los datos de descenso; sin
embargo, tras restringir~$U$, será compatible con ellos. Esto concluye la
demostración.

\begin{corollaire}
\label{IX.4.5}
Bajo las condiciones de~\Ref{IX.4.4}, el dato de descenso sobre~$X'$ es
efectivo si y solo si, para todo $a\in S$, lo es el dato de descenso
correspondiente sobre~$X_a'$ respecto del morfismo
$S_a'=S'\times_S\Spec(\cal{O}_a)\to\Spec(\cal{O}_a)$. Cuando~$S$ es
localmente noetheriano y~$X'$ está separado sobre~$S'$, en el criterio
precedente se puede sustituir también~$\cal{O}_a$ por su completación.
\end{corollaire}

La
\marginpar{238}
primera afirmación resulta de~\Ref{IX.4.4} y~\Ref{IX.4.3}; la segunda es
entonces consecuencia de~\Ref{IX.4.2}. Utilizando de nuevo~\Ref{IX.4.2}
y el hecho de que, para todo anillo local noetheriano~$A$, se puede
encontrar un anillo local noetheriano completo~$B$ y un homomorfismo local
$A\to B$ tales que~$B$ sea plano sobre~$A$ y que $B/\goth{m}B$ sea una
extensión dada del cuerpo residual $k=A/\goth{m}$ de~$A$, se obtiene:

\begin{corollaire}
\label{IX.4.6}
Bajo las condiciones de~\Ref{IX.4.4}, supongamos además que~$X'$ está
separado sobre~$S'$ y que~$S$ es localmente noetheriano. El dato de
descenso sobre~$X'$ es efectivo si y solo si, para todo preesquema~$S_1$
sobre~$S$ que sea espectro de un anillo local completo con cuerpo residual
algebraicamente cerrado, lo es el dato de descenso correspondiente sobre
$X_1'=X'\times_S S_1$ respecto del morfismo
$g_1\colon S_1'\to S_1$.
\end{corollaire}

\begin{theoreme}
\label{IX.4.7}
Sea $g\colon S'\to S$ un morfismo finito, sobreyectivo y de presentación
finita (esta última hipótesis es consecuencia de las demás si~$S$ es
localmente noetheriano)\footnote{Se puede demostrar que, de hecho, basta
que~$g$ sea un morfismo \emph{entero}, reduciéndose al caso del texto
mediante un procedimiento de paso al límite al estilo de EGA IV~8.}.
Entonces~$g$ es un morfismo de descenso efectivo para la categoría fibrada
de los preesquemas étales, separados y de tipo finito sobre otros.
\end{theoreme}

Hay que demostrar que, si~$X'$ es étale, separado y de tipo finito
sobre~$S'$, y está provisto de un dato de descenso respecto de~$g$, este
dato es efectivo. Utilizando~\Ref{IX.4.3}, se reduce fácilmente al caso en
que~$S$ es noetheriano. Gracias a~\Ref{IX.4.5}, se puede suponer, por
tanto, que~$S$ es el espectro de un anillo local noetheriano y, con mayor
razón, que
$$
\dim S=n<+\infty.
$$
Se razona entonces por inducción sobre~$\dim S=n$; la afirmación es trivial
para $n<0$. Supongamos, pues, que $n\geq 0$ y que el teorema está
demostrado para las dimensiones $n'<n$. En virtud de~\Ref{IX.4.6}, se
reduce al caso en que~$S$ es el espectro de un anillo local completo; por
tanto,~$S'$ es una unión finita de espectros de anillos locales completos.
Se tiene entonces
$$
X'=X_1'\sqcup X_2'
$$
donde~$X_1'$ es \emph{finito} sobre~$S'$ y~$X_2'$ no tiene ningún punto
sobre ninguno de los puntos cerrados
\marginpar{239}
de~$S'$. Consideremos los morfismos
$$
\xymatrix@C=.5cm{ q_1,q_2\colon X'' \ar@/^2pt/[r] \ar@/_2pt/[r] & X' }
$$
correspondientes al dato de descenso, compatibles con
$\xymatrix@C=.5cm{p_1,p_2\colon S'' \ar@/^2pt/[r] \ar@/_2pt/[r] & S'}$.
Se ve inmediatamente que
$$
X''=q_i^{-1}(X_1')\sqcup q_i^{-1}(X_2')\quad i=1,2
$$
es la descomposición canónica análoga de~$X''$ sobre~$S''$, lo que implica
$q_1^{-1}(X_1')=\allowbreak q_2^{-1}(X_1')$; por consiguiente,~$X_1'$ y
$X_2'$ están provistos de datos de descenso inducidos. Sea ahora~$T$ el
abierto de~$S$ complementario de su punto cerrado; entonces
$T'=S'\times_S T$ es la parte de~$S'$ complementaria del conjunto de
puntos cerrados, y~$X_2'$, que se encuentra íntegramente sobre~$T'$, está
provisto de un dato de descenso respecto del morfismo $T'\to T$ inducido
por~$g$. Como este último es finito y sobreyectivo, y
$\dim T<\dim S=n$, este dato de descenso es efectivo por la hipótesis de
inducción. Se ve, pues, que basta demostrar que el dato de descenso
sobre~$X_1'$ es efectivo; por tanto, se puede suponer ahora que~$X'$ es
étale y \emph{finito} sobre~$S'$. (N.B.: el razonamiento por inducción es
innecesario si en el enunciado~\Ref{IX.4.7} se limita uno a los
recubrimientos étales). Sea entonces~$S_0$ el espectro del cuerpo residual
de~$A$, sea $S_0'=S'\times_S S_0$, y definamos análogamente~$S_0''$ y
$S_0'''$ a partir del cuadrado y el cubo fibrados~$S''$ y~$S'''$ de~$S'$
sobre~$S$. En virtud de~\Ref{IX.1.8}, los morfismos $S_0\to S$,
$S_0'\to S'$, etc., inducen equivalencias para las categorías de los
recubrimientos étales de~$S$ y~$S_0$, por una parte, y de~$S'$ y~$S_0'$,
por otra, \ifthenelse{\boolean{orig}}{etc...}{etc.} De los argumentos
formales de la teoría del descenso en las categorías~\cite{IX.D} se deduce
que $g\colon S'\to S$ es un morfismo de descenso efectivo para la
categoría fibrada de los recubrimientos étales si y solo si lo es
$g_0\colon S_0'\to S_0$. Pero este es, en efecto, el caso, por ejemplo
como caso particular de~\Ref{IX.4.1}. Esto concluye la demostración.

\begin{corollaire}
\label{IX.4.8}
La conclusión de~\Ref{IX.4.7} sigue siendo válida si solo se supone que
$S'\to S$ es universalmente submersivo, de tipo finito y cuasifinito,
siempre que~$S$ sea localmente noetheriano.
\end{corollaire}

En
\marginpar{240}
efecto, en virtud de~\Ref{IX.4.6}, se puede suponer que~$S$ es el espectro
de un anillo local noetheriano completo. Entonces, en virtud
de~\Ref{IX.2.5}, existen un morfismo finito y sobreyectivo $S_1\to S$ y un
$S$-morfismo $S_1\to S'$. Como $S_1\to S$ es un morfismo universal de
descenso estricto para la categoría fibrada considerada, en virtud
de~\Ref{IX.4.7}, y $S'\to S$ es un morfismo universal de descenso para
dicha categoría, \Ref{IX.4.8} resulta de los argumentos formales
generales~\cite{IX.D}.

\begin{corollaire}
\label{IX.4.9}
Sea $g\colon S'\to S$ un morfismo de tipo finito, sobreyectivo y
universalmente abierto, con~$S$ localmente noetheriano. Entonces~$g$ es un
morfismo de descenso efectivo para la categoría fibrada de los
preesquemas étales, separados y de tipo finito sobre otros.
\end{corollaire}

Procediendo como en~\Ref{IX.4.7}, se reduce al caso en que~$S$ es el
espectro de un anillo local noetheriano completo~$A$. Sea~$A_1$ una
$A$-álgebra finita, de espectro~$S_1$, tal que $S_1\to S$ sea finito y
\emph{sobreyectivo}, y por tanto un morfismo universal de descenso efectivo
para la categoría fibrada considerada, gracias a~\Ref{IX.4.7}. Se deduce
entonces de los teoremas generales~\cite{IX.D} que~$g$ es un morfismo de
descenso efectivo para la categoría fibrada considerada si y solo si lo es
el morfismo correspondiente
$g_1\colon S_1'=S'\times_S S_1\to S_1$. Como este último satisface las
mismas hipótesis que~$g$, se reduce la cuestión a demostrar~\Ref{IX.4.9}
para~$S_1$ en lugar de~$S$. Tomando primero como~$A_1$ el producto directo
de los $A/\goth{p}_i$, para los ideales primos mínimos~$\goth{p}_i$
de~$A$, se reduce al caso en que~$A$ es un \emph{dominio de integridad}.
Se demuestra entonces\footnote{\Cf EGA IV~14.3.13 y~14.5.4.} que existe un
subesquema íntegro~$S_1$ de~$S'$, cuasifinito sobre~$S$ y dominante
sobre~$S$, que pasa por un punto de la fibra de~$S'$ sobre el punto
cerrado~$y$ de~$S$ (gracias a que~$S'$ es universalmente abierto y de tipo
finito sobre el esquema local noetheriano íntegro~$S$, y a que
$S_y'\neq\emptyset$). Como~$A$ es completo, $S_1$ es finito sobre~$S$ y,
como domina~$S$, el morfismo $S_1\to S$ es sobreyectivo. Sustituyendo una
vez más~$S$ por~$S_1$, se reduce al caso en que~$S'$ tiene una sección
sobre~$S$, caso en el cual el enunciado es trivial.

\begin{theoreme}
\label{IX.4.10}
Sea $g\colon S'\to S$ un morfismo finito, radicial, sobreyectivo y de
presentación finita (esta última condición es superflua si~$S$ es
localmente noetheriano\footnote{Basta incluso que~$g$ sea entero, radicial
y sobreyectivo, como se ve mediante una reducción fácil al caso del texto,
al estilo de EGA~IV~8; \cf SGA~4 VIII~1.1.}).
\marginpar{241}
Entonces el funtor de imagen inversa induce una equivalencia de la categoría
de los preesquemas étales sobre~$S$ con la categoría de los preesquemas
étales sobre~$S'$.
\end{theoreme}

Como los morfismos diagonales de~$S'$ en $S'\times_S S'$ y
$S'\times_S S'\times_S S'$ son inmersiones sobreyectivas y, por tanto,
inducen en virtud de~\Ref{IX.1.9} equivalencias de las categorías de los
preesquemas étales sobre~$S'\times_S S'$ \resp
$S'\times_S S'\times_S S'$ con la categoría de los preesquemas étales
sobre~$S'$, se deduce de los argumentos formales del descenso~\cite{IX.D}
que todo~$X'$ étale sobre~$S'$ está provisto de un único dato de descenso
respecto de $g\colon S'\to S$. Por tanto,~\Ref{IX.3.3} implica que el
funtor de imagen inversa por~$g$, de la categoría de los preesquemas étales
sobre~$S$ en la categoría de los preesquemas étales sobre~$S'$, es
\emph{plenamente fiel}. Queda por demostrar que es esencialmente
sobreyectivo, es decir, que todo~$X'$ étale sobre~$S'$ es isomorfo a la
imagen inversa de un~$X$ étale sobre~$S$. Como la cuestión es evidentemente
local sobre~$S$ \emph{y sobre}~$X'$, se puede suponer que $S$,~$S'$ y~$X'$
son afines. Pero entonces~$X'$ está separado y es de tipo finito
sobre~$S'$, y se puede aplicar el criterio de efectividad~\Ref{IX.4.7}.

\begin{corollaire}
\label{IX.4.11}
La conclusión de
\ifthenelse{\boolean{orig}}
{~\Ref{IX.4.9}}
{~\Ref{IX.4.10}}
sigue siendo válida si se sustituye la hipótesis sobre~$g$ por la
siguiente: $g$ es fielmente plano, cuasicompacto y radicial.
\end{corollaire}

La misma demostración, invocando~\Ref{IX.4.1} en lugar
de~\Ref{IX.4.7}.

Se observará que la demostración de~\Ref{IX.4.7} es
\og elemental\fg{} en cuanto no utiliza los teoremas
\ifthenelse{\boolean{orig}}
{}
{de}
finitud y de comparación para los morfismos propios (EGA~III~3,~4,~5).
No sucede ya lo mismo con el resultado siguiente:

\begin{theoreme}
\label{IX.4.12}
Sea $g\colon S'\to S$ un morfismo propio, sobreyectivo y de presentación
finita (esta última hipótesis es consecuencia de la primera si~$S$ es
localmente noetheriano). Entonces~$g$ es un morfismo de descenso efectivo
para la categoría fibrada de los recubrimientos étales de preesquemas.
\end{theoreme}

En virtud de~\Ref{IX.3.3} y de~\Ref{IX.2.2}, se reduce la cuestión a
demostrar que, para todo recubrimiento étale~$X'$ sobre~$S'$, provisto de
un dato de descenso respecto de $g\colon S'\to S$, este dato es efectivo.
Utilizando~\Ref{IX.4.3}, se reduce fácilmente
\marginpar{242}
al caso en que~$S$ es noetheriano y, utilizando~\Ref{IX.4.6}, se puede
suponer, por tanto, que~$S$ es el espectro de un anillo local noetheriano
\emph{completo}~$A$. Introduzcamos~$S''$ y~$S'''$ como de costumbre; sea
$S_0$ el espectro del cuerpo residual de~$A$, y sean~$S_0'$,~$S_0''$ y
$S_0'''$ los preesquemas deducidos de~$S'$,~$S''$ y~$S'''$ por el cambio
de base $S_0\to S$, es decir, las fibras de~$S'$,~$S''$ y~$S'''$ en el
punto cerrado de~$S$. Según~\Ref{IX.1.10}, los morfismos $S_0\to S$,
$S_0'\to S'$, etc., inducen equivalencias de la categoría de los
recubrimientos étales sobre el esquema de llegada con la categoría de los
recubrimientos étales sobre el esquema de origen. En consecuencia,
$g\colon S'\to S$ es un morfismo de descenso estricto para la categoría
fibrada de los recubrimientos étales de preesquemas si y solo si lo es
$g_0\colon S_0'\to S_0$, lo que, en efecto, ocurre en virtud
de~\Ref{IX.4.1}. Esto concluye la demostración de~\Ref{IX.4.12}. (En este
razonamiento solo se necesitaba de~\Ref{IX.1.10} el hecho de que el funtor
considerado allí es \emph{plenamente fiel}, lo cual \emph{no} utiliza el
teorema de existencia de haces coherentes en geometría algebraica.)

\section{Traducción en términos del grupo fundamental}
\label{IX.5}

Sea
$$
g\colon S'\to S
$$
un \emph{morfismo de descenso efectivo} para la categoría fibrada de los
\emph{recubrimientos étales} de preesquemas; por ejemplo, un morfismo propio,
sobreyectivo y de presentación finita~\eqref{IX.4.12},
o un morfismo fielmente plano y cuasicompacto. Introduzcamos, como de
costumbre,~$S''$ y~$S'''$, y denotemos por~$\cal{C}$,~$\cal{C}'$,
~$\cal{C}''$ y~$\cal{C}'''$ las categorías de los recubrimientos étales de
~$S$,~$S'$,~$S''$ y~$S'''$, respectivamente. Se tiene entonces un diagrama
$2$-exacto de categorías
\begin{equation*}
\label{eq:IX.5.*}
\tag{$*$}
\xymatrix@C=.5cm{
{\cal{C}} \ar[r]^-{~p^*} & {\cal{C}'}
\ar@/^2pt/[rr]^{p_1^*,p_2^*} \ar@/_2pt/[rr] && {\cal{C}''}
\ar@/^4pt/[rrr]^{p_{21}^*,p_{32}^*,p_{31}^*} \ar[rrr] \ar@/_4pt/[rrr] &&&
{\cal{C}'''}
}
\end{equation*}
correspondiente al diagrama
$$
\xymatrix@C=.5cm{ S & \ar[l]_-{~p} S' && \ar@/^2pt/[ll] \ar@/_2pt/[ll]_-{p_1,p_2}
S'' &&& \ar@/^4pt/[lll] \ar[lll] \ar@/_4pt/[lll]_{p_{21},p_{32},p_{31}} S'''.
}
$$
Supongamos
\marginpar{243}
que los preesquemas~$S$,~$S'$,~$S''$ y~$S'''$ son sumas disjuntas de
preesquemas conexos, como ocurre, en particular, si son localmente conexos y,
con mayor razón, si son localmente noetherianos (por ejemplo, si~$S'$ es de
tipo finito sobre~$S$ localmente noetheriano). Las categorías~$\cal{C}$,
~$\cal{C}'$, …, de~\eqref{eq:IX.5.*} son entonces categorías
multigaloisianas~\Ref{V.9}, descritas cada una por una familia de grupos
topológicos compactos totalmente disconexos, a saber, los grupos
fundamentales de las componentes conexas de los preesquemas~$S$,~$S'$,
~$S''$ y~$S'''$. Para simplificar, supondremos que~$S$ es conexo y daremos un
procedimiento de cálculo de su grupo fundamental en términos de la categoría
fibrada formada por~$\cal{C}'$,~$\cal{C}''$ y~$\cal{C}'''$, explicitada
convenientemente mediante los grupos fundamentales que describen estas
categorías. El lector observará que el procedimiento esbozado vale, de hecho,
en el marco general de las categorías multigaloisianas (que no tienen por qué
proceder de preesquemas dados~$S$,~$S'$,~$S''$ y~$S'''$). Es, por lo demás,
el análogo del procedimiento bien conocido para calcular el grupo fundamental
de un espacio topológico~$S$, unión localmente finita de subespacios
cerrados~$S_i$ (o unión arbitraria de subespacios abiertos~$S_i$), mediante
los grupos fundamentales de las componentes
\ifthenelse{\boolean{orig}}
{de los~$X_i$}
{de los~$S_i$}
y de las componentes de los~$S_i\cap S_j$. Naturalmente, la situación análoga
para los preesquemas cae dentro del marco general del descenso, introduciendo
el preesquema~$S'$, suma de los~$S_i$, y el morfismo canónico
$g\colon S'\to S$.

Pongamos
$$
E'=\pi_0(S'), \quad E''=\pi_0(S''), \quad E'''=\pi_0(S'''),
$$
donde~$\pi_0$ designa el funtor «conjunto de las componentes conexas». Como
los productos fibrados de~$S'$ sobre~$S$ forman un objeto simplicial de
~$\Sch$, el funtor~$\pi_0$ lo transforma en un conjunto simplicial del cual
~$E'$,~$E''$ y~$E'''$ son las componentes de dimensiones~$0$,~$1$ y~$2$.
Habremos
de utilizar las aplicaciones simpliciales
$$
q_i=\pi_0(p_i),\quad(i=1,2)\quad\text{y}
\quad q_{ij}=\pi_0(p_{ij}),\quad (i,j)=(2,1),(3,2),(3,1),
$$
puestas de manifiesto en el diagrama
\begin{equation*}
\label{eq:IX.5.1}
\tag{1} {\xymatrix@C=.5cm{ E' && \ar@/^2pt/[ll] \ar@/_2pt/[ll]_{q_1,q_2} E'' &&&
\ar@/^4pt/[lll] \ar[lll] \ar@/_4pt/[lll]_{q_{21},q_{32},q_{31}} E'''. }}
\end{equation*}
Los elementos
\marginpar{244}
de~$E'$ se denotarán con una prima, como~$s'$; los de~$E''$ \resp los de
~$E'''$ se denotarán con~$''$ \resp con~$'''$. Que~$S$ sea conexo se expresa por
$\pi_0(K)=0$, donde~$K$ es el conjunto simplicial definido por
$g\colon S'\to S$, o también por el hecho de que la relación de equivalencia
en~$E'$ engendrada por el par de aplicaciones~$(q_1,q_2)$ es transitiva.

Elegiremos, de una vez por todas, un elemento~$s_0'$ de~$E'$ y, para cada
~$s'$ de~$E'$, un elemento~$\overline{s'}\in E''$ tal que\footnote{Debe
advertirse que el elemento~$\overline{s'}$, cuya existencia se admite
implícitamente, no existe en todos los casos. Por consiguiente, el
teorema~\Ref{IX.5.1}, tal como está enunciado, no se aplica en todos los casos.
Sin embargo, no es difícil, inspirándose en el texto, modificar su enunciado
de modo que proporcione un método de cálculo aplicable en todos los casos. En
particular, los corolarios de dicho teorema son válidos tal cual.}
$$
q_1(\overline{s'})=s_0', \quad q_2(\overline{s'})=s',
$$
poniendo así de manifiesto la conexidad de~$S$. Para cada~$s'\in E'$, elijamos
un punto geométrico~$\underline{s}'$ en la componente conexa~$s'$ de~$S'$;
este punto intervendrá, de hecho, mediante el funtor fibra~$F_{s'}'$ asociado
sobre la categoría multigaloisiana~$\cal{C}'$. El grupo de automorfismos de
este funtor, es decir, el grupo fundamental de~$S'$ en~$\underline{s}'$, se
denotará por~$\pi_{s'}$. Elijamos del mismo modo puntos~$\underline{s}''$ y
~$\underline{s}'''$, y por tanto funtores~$F_{s''}''$ y~$F_{s'''}'''$, de
donde se obtienen los grupos fundamentales~$\pi_{s''}$ y~$\pi_{s'''}$. Así,
$$
\pi_{s'}=\pi_1(S',\underline{s}'), \quad
\pi_{s''}=\pi_1(S'',\underline{s}''), \quad
\pi_{s'''}=\pi_1(S''',\underline{s}''').
$$
Para todo~$s''\in E''$, $p_1(\underline{s}'')$ pertenece a la misma
componente conexa que~$\underline{q_1(s'')}$; existe, pues, un isomorfismo de
funtores
\ifthenelse{\boolean{orig}}
{$F''_{s''}\circ p_1^* F'_{s'}$}
{$F''_{s''}\circ p_1^*\isomto F'_{q_1(s'')}$}
(es decir, una «clase de caminos» de~$p_1(\underline{s}'')$ a
~$\underline{q_1(s'')} $). La misma observación se aplica a~$q_2$ y a los
~$q_{ij}$. Elijamos todas estas clases de caminos:
$$
F''_{s''}\circ p_i^*\isomto F'_{q_i(s'')},\quad F_{s'''}'''\circ
p_{ij}^*\isomto
\ifthenelse{\boolean{orig}}{F'''_{q_{ij}(s''')}}{F''_{q_{ij}(s''')}},
$$
(para
\marginpar{245}
$i=1,2$ y $(i,j)=(2,1),(3,2),(3,1)$). Se obtienen, en particular,
homomorfismos de grupos
\begin{equation*}
\label{eq:IX.5.2}
\tag{2} {q_i^{s''}\colon \pi_{s''}\to\pi_{q_i(s'')},\quad
q_{ij}^{s'''}\colon \pi_{s'''}\to\pi_{q_{ij}(s''')},}
\end{equation*}
(con los mismos valores de~$i$ y~$(i,j)$). Recordemos, finalmente, que en la
estructura con clivaje de la categoría fibrada con fibras~$\cal{C}'$,
~$\cal{C''}$ y~$\cal{C'''}$ figuran también isomorfismos de funtores
$$
p_{21}^*p_1^*\isomto p_{31}^*p_1^*,\quad p_{21}^*p_2^*\isomto
p_{32}^*p_1^*,\quad p_{31}^*p_2^*\isomto p_{32}^*p_2^*,
$$
deducidos de sendos isomorfismos de ambos miembros con los~$u_i^*$
~($i=1,2,3$), donde los~$u_i$ son las tres proyecciones de~$S'''$ en~$S'$.
Al explicitar estos datos, para cada~$s'''$ se obtiene un elemento bien
determinado
\begin{equation*}
\label{eq:IX.5.3}
\tag{3} {a_i^{s'''}\in\pi_{v_i(s''')},}
\end{equation*}
(donde los~$v_i$, $i=1,2,3$, son las tres aplicaciones~$E'''\to E'$ definidas
por~$v_i=\pi_0(u_i)$), sometido además a las condiciones
$$
q_1^{s_1''}q_{21}^{s'''}=\mathrm{int}(a_1^{s'''})
q_1^{s_2''}q_{31}^{s'''}\qquad (s_1''=q_{21}(s'''),\
s_2''=q_{31}(s''')),
$$
y a las dos condiciones análogas, en las que intervienen~$a_2$ y~$a_3$. El
lector observará, por lo demás, que los datos~\eqref{eq:IX.5.1},
~\eqref{eq:IX.5.2} y~\eqref{eq:IX.5.3} permiten reconstruir, salvo equivalencia
de categorías fibradas, la categoría fibrada considerada, con fibras
~$\cal{C}'$,~$\cal{C}''$ y~$\cal{C}'''$. Deben permitir, pues, reconstruir en
principio~$\cal{C}$ salvo equivalencia y, por tanto, su grupo fundamental salvo
isomorfismo. Determinaremos, en realidad, el grupo fundamental en el punto
geométrico~$p(\underline{s}_0')$ de~$S$, es decir, el grupo de automorfismos
de~$F'_{s_0'}\circ p^*$.

Obsérvese que dar un objeto~$X'$ de~$\cal{C}'$ equivale esencialmente a dar
conjuntos finitos~$X'_{s'}$ ($s'\in E'$) sobre los que los~$\pi_{s'}$ actúan
continuamente. Un
\marginpar{246}
dato de pegado sobre tal objeto equivale entonces a dar, para cada
~$s''\in E''$, una biyección
$$
\varphi_{s''}\colon X'_{q_1(s'')}\isomto X'_{q_2(s'')}
$$
compatible con las acciones de~$\pi_{s''}$ sobre ambos miembros por medio de
los homomorfismos~$q_i^{s''}\colon \pi_{s''}\to\pi_{q_i(s'')}$. Tomando
primero los~$s''$ de la forma~$\overline{s'}$, se ve que tal dato define
biyeciones
$$
\psi_{s'}\colon X'_{s_0'}=F_0'(X')\to X_{s'}'
$$
lo que permite identificar todos los~$X'_{s'}$ con el mismo conjunto
~$F_0'(X')=X'_{s_0'}$, sobre el cual actuarán desde ese momento todos los
grupos~$\pi_{s'}$. Hecho esto, las biyecciones~$\varphi_{s''}$ corresponderán
a biyecciones
$$
g_{s''}\colon F_0'(X')\isomto F_0'(X'),
$$
sometidas, por una parte, a las relaciones de conmutación con~$\pi_{s''}$
\begin{equation*}
\label{eq:IX.5.a}
\tag*{a)} {g_{s''}q_1^{s''}(g'')=q_2^{s''}(g'')g_{s''}\qquad (s''\in E'',\ g''\in\pi_{s''}),}
\end{equation*}
y, por otra, a las relaciones
\begin{equation*}
\label{eq:IX.5.b}
\tag*{b)} {g_{\overline{s'}}=g_{\overline{s_0'}}\qquad (s'\in E'),}
\end{equation*}
que expresan cómo habíamos identificado entre sí los
~$\ifthenelse{\boolean{orig}}{X_{s'}}{X_{s'}'}$. Al explicitar
la condición para que tal dato de pegado sea, en efecto, un dato de descenso,
se obtienen las relaciones
\begin{equation*}
\label{eq:IX.5.c}
\tag*{c)} {a_3^{s'''}g_{q_{31}(s''')}a_1^{s'''}=
g_{q_{32}(s''')}a_2^{s'''}g_{q_{21}(s''')}\qquad (s'''\in E''').}
\end{equation*}

Esto proporciona una equivalencia entre la categoría de los objetos de
~$\cal{C}'$ provistos de un
\marginpar{247}
dato de descenso y la categoría de los conjuntos finitos sobre los que actúan
continuamente los grupos~$\pi_{s'}$, provistos además de biyecciones
~$g_{s''}$ que satisfacen las relaciones~\Ref{eq:IX.5.a},
~\Ref{eq:IX.5.b} y~\Ref{eq:IX.5.c}. Sea entonces~$G$ el grupo engendrado por
los grupos~$\pi_{s'}$ y los nuevos generadores~$g_{s''}$, sometidos a las
relaciones~\Ref{eq:IX.5.a},~\Ref{eq:IX.5.b} y~\Ref{eq:IX.5.c}, y sea~$\pi$
el grupo límite proyectivo de los cocientes de~$G$ por los subgrupos de índice
finito cuyas imágenes inversas en los grupos~$\pi_{s'}$ sean subgrupos
abiertos. También se dice que~$\pi$ es el \emph{grupo de tipo galoisiano
engendrado por los~$\pi_{s'}$ y los~$g_{s''}$, sometido a las
relaciones~\Ref{eq:IX.5.a},~\Ref{eq:IX.5.b} y~\Ref{eq:IX.5.c}}. Se comprueba
de inmediato que la categoría considerada es asimismo equivalente a la
categoría de conjuntos finitos sobre los que el grupo topológico~$\pi$ actúa
continuamente. Esto demuestra el enunciado siguiente:

\begin{theoreme}
\label{IX.5.1}
Sea $g\colon S'\to S$ un morfismo de preesquemas que sea un morfismo de
descenso efectivo para la categoría fibrada de recubrimientos étales de
preesquemas (\cf \Ref{IX.4.9} y~\Ref{IX.4.12}). Supongamos que~$S$ es conexo
y que~$S'$, su cuadrado fibrado~$S''$ y su cubo fibrado~$S'''$ son sumas de
preesquemas conexos (como ocurre, por ejemplo, si~$S'$ es de tipo finito
sobre~$S$ localmente noetheriano y conexo). Elijamos, como antes, un punto
geométrico en cada componente conexa de~$S'$,~$S''$ y~$S'''$, ciertas clases
de caminos, un~$s_0'\in E'$ y, para cada~$s'\in E'$, un
\ifthenelse{\boolean{orig}}
{$s'\in E''$}
{$s''\in E''$}
cuyas dos imágenes en~$E'$ sean~$s_0'$ y~$s'$. ($E'$,~$E''$ y~$E'''$
designan, respectivamente, los conjuntos de componentes conexas de~$S'$,
~$S''$ y~$S'''$). Entonces el grupo fundamental de~$S$ en el punto geométrico
imagen de~$s_0'$ es canónicamente isomorfo al grupo de tipo galoisiano
engendrado por los~$\pi_{s'}=\pi_1(S',\underline{s}')$ ($s'\in E'$) y por
generadores~$g_{s''}$ ($s''\in E''$), sometidos a las relaciones
~\Ref{eq:IX.5.a},~\Ref{eq:IX.5.b} y~\Ref{eq:IX.5.c} anteriores, en las que
intervienen los elementos de los grupos
$\pi_{s''}=\pi_1(S'',\underline{s}'')$ y los elementos~$a_i^{s'''}$
($i=1,2,3$, $s'''\in E'''$) introducidos más arriba.
\end{theoreme}

\begin{corollaire}
\label{IX.5.2}
Supongamos que~$S'$ y~$S''$ no tienen más que un número finito de componentes
conexas y que los grupos fundamentales de las componentes conexas de~$S'$
son topológicamente finitamente generados. Entonces el grupo fundamental de
~$S$ es topológicamente finitamente generado.
\end{corollaire}

Así,
\marginpar{248}
demostraremos más adelante que el grupo fundamental de un esquema proyectivo
normal sobre un cuerpo algebraicamente cerrado es topológicamente finitamente
generado. Utilizando el lema de Chow y la normalización de los esquemas
algebraicos, se seguirá que el mismo resultado es cierto para todo esquema
propio sobre un cuerpo algebraicamente cerrado.

\begin{corollaire}
\label{IX.5.3}
Supongamos que~$S'$,~$S''$ y~$S'''$ no tienen más que un número finito de
componentes conexas, que los grupos fundamentales de las componentes conexas
de~$S'$ son topológicamente de \emph{presentación finita}, y que los grupos
fundamentales de las componentes conexas de~$S''$ son topológicamente
\emph{finitamente generados}. Entonces el grupo fundamental de~$S$ es
topológicamente de \emph{presentación finita}.
\end{corollaire}

Obsérvese que
\ifthenelse{\boolean{orig}}{\Ref{IX.4.9}}{\Ref{IX.4.10}}
puede expresarse, al restringirlo a los \emph{recubrimientos} étales,
diciendo que \emph{un morfismo finito, radicial y sobreyectivo de preesquemas
noetherianos induce un isomorfismo de los grupos fundamentales}; en términos
figurados, se puede decir, pues, que el grupo fundamental es un
\emph{invariante topológico} para los preesquemas.
Más generalmente, mediante~\Ref{IX.5.1} se puede explicitar el efecto sobre
el grupo fundamental de operaciones sobre los preesquemas, tales como el
«pinzamiento» de un preesquema a lo largo de un conjunto finito de puntos,
que tiene un significado topológico sencillo. Se obtiene, por ejemplo:

\begin{corollaire}
\label{IX.5.4}
Sean $g\colon S'\to S$ un morfismo finito de presentación finita y~$T$ un
subconjunto discreto de~$S$. Para cada~$s\in S$, sea~$n(s)$ el «número geométrico de
puntos» de la fibra~$g^{-1}(s)$ (que puede describirse también como el grado
separable de~$g^{-1}(s)$ sobre~$\kres(s)$, suma de los grados separables de sus
extensiones residuales). Supongamos que, para~$s\in S-T$, se tiene~$n(s)=1$. Para cada
~$s\in T$, sea~$K_s$ una extensión algebraicamente cerrada de~$\kres(s)$,
sea~$I_s$ el conjunto de los puntos geométricos de~$S'$ con valores en~$K_s$
(un conjunto con~$n(s)$ elementos), sea~$I_s'$ el complementario de un punto
elegido de~$I_s$ y, finalmente, sea~$I'$ la unión de los~$I_s'$. Supongamos
que~$S'$ es conexo. Entonces el grupo fundamental de~$S$ es isomorfo al grupo
de tipo galoisiano engendrado por el grupo fundamental de~$S'$ y por
generadores~$g_i$ ($i\in I'$), sin relaciones suplementarias.
\end{corollaire}

Los
\marginpar{249}
detalles de la demostración se dejan al lector; el enunciado obtenido no es
más que la traducción, al lenguaje de la teoría de grupos, del hecho de que la
categoría~$C$ de los recubrimientos étales de~$S$ es equivalente a la
categoría de los recubrimientos étales~$X'$ de~$S'$, \emph{provistos}, para
cada~$s\in T$, de un sistema transitivo de biyecciones entre las~$n(s)$ fibras
de~$X'$ en los puntos de~$g^{-1}(s)$ con valores en~$K_s$. (En esta forma
intrínseca, naturalmente, ya no es necesario suponer que~$S'$ es conexo).

\begin{exemple}
\label{IX.5.5}
Se demuestra fácilmente que la curva racional~$\PP^1_k$ sobre un cuerpo
algebraicamente cerrado~$k$ es simplemente conexa\footnote{\Cf
\Ref{XI.1.1}.}. Por tanto, el grupo fundamental de una curva racional
completa que tiene exactamente un punto doble con~$n$ ramas analíticas es el
grupo de tipo galoisiano libre engendrado por~$n-1$ generadores. Por ejemplo,
en el caso de un punto doble ordinario se obtiene el grupo fundamental
~$\hat\ZZ$, como se anunció en~(\Ref{I.11}~a)). En cambio, la existencia de
una cúspide (que es un punto «geométricamente unirramoso») no influye en el
grupo fundamental.
\end{exemple}

\begin{corollaire}
\label{IX.5.6}
Sea $g\colon S'\to S$ un morfismo de preesquemas universalmente submersivo,
con fibras geométricamente conexas, siendo~$S$ conexo. Entonces~$S'$ es
conexo y, eligiendo un punto geométrico~$s'$ de~$S'$ y denotando por~$s$ su
imagen en~$S$, el homomorfismo
$$
\pi_1(S',s')\to \pi_1(S,s)
$$
es \emph{sobreyectivo}. Si~$g$ es un morfismo de descenso efectivo para la
categoría fibrada de los recubrimientos étales de preesquemas
(\cf \emph{\Ref{IX.4.12}}), introduzcamos el punto geométrico
~$s^{\prime\prime}=\diag(s')$ de~$S^{\prime\prime}=S'\times_S S'$ y los dos
homomorfismos
$$
p_{1*},p_{2*}\colon \pi_1(S^{\prime\prime},s^{\prime\prime}) \to
\pi_1(S',s')
$$
inducidos por las dos proyecciones. Entonces~$\pi_1(S,s)$ es isomorfo al
conúcleo de este par de homomorfismos en la categoría de los grupos de tipo
galoisiano, es decir, al cociente de~$\pi_1(S',s')$ por el subgrupo normal
cerrado engendrado por los
\marginpar{250}
elementos de la forma
$p_{1*}(g^{\prime\prime})p_{2*}(g^{\prime\prime})^{-1}$, con
$g^{\prime\prime}\in \pi_1(S^{\prime\prime},s^{\prime\prime})$.
\end{corollaire}

En efecto, por~\Ref{IX.3.4} se sabe que el funtor
$X\mto X\times_SS'$ de los recubrimientos étales de~$S$ en los
recubrimientos étales de~$S'$ es plenamente fiel, lo que equivale a que el
homomorfismo de grupos fundamentales sea un epimorfismo (\Ref{V.6.9}). La
última afirmación es consecuencia inmediata de la descripción de
~\Ref{IX.5.1}.

\begin{remarque}
\label{IX.5.7}
No se sabe actualmente si el grupo fundamental de un esquema propio sobre un
cuerpo algebraicamente cerrado~$k$ es topológicamente de presentación
finita\footnote{Esto parece muy improbable en el caso de las curvas lisas de
género~$g\geq 2$ en característica~$p>0$. En cambio, cuando se sustituye
~$\pi_1$ por su mayor cociente primo con~$p$, parece que las técnicas bien
conocidas permiten dar una respuesta afirmativa, incluso sin suponer que el
esquema sea propio. \Cf un trabajo en preparación de J.P\ptbl Murre.}. Utilizando
~\Ref{IX.5.3}, una técnica bien conocida de secciones hiperplanas y la
desingularización de las superficies normales, se reduce al caso de una
\emph{superficie lisa sobre}~$k$. Esto permite, al menos, demostrar por
métodos trascendentes que la respuesta es afirmativa en característica~$0$
(sin necesidad de admitir la triangulabilidad de las variedades algebraicas
singulares). En característica~$p>0$, la dificultad principal parece residir
en el caso de las curvas, cuyo grupo fundamental solo se sabe que es un
cociente del que aparece en el caso clásico (\cf el exposé siguiente), pero
el núcleo por el que se toma el cociente se conoce muy mal.
\end{remarque}

\begin{remarque}
\label{IX.5.8}
Además de~\Ref{IX.5.4} y~\Ref{IX.5.6}, se podrían explicitar otros casos
particulares en los que~\Ref{IX.5.1} adopta una forma especialmente sencilla.
Un caso interesante es aquel en que~$S$ es el cociente de~$S'$ por un grupo
finito de automorfismos~$\Gamma$. \emph{Entonces la categoría de los
recubrimientos étales
\ifthenelse{\boolean{orig}}{de~$S'$}{de~$S$}
es equivalente a la categoría de los recubrimientos étales~$X'$ de~$S'$ en
los que el grupo~$\Gamma$ actúa de modo compatible con su acción sobre~$S'$,
de manera que para cada~$s'\in S'$ y cada~$g\in\Gamma_{s'}$} (donde
~$\Gamma_{s'}$ designa \emph{el grupo de inercia} de~$s'$ en~$\Gamma$),
\emph{$g$ actúa trivialmente sobre la fibra}~$X_{s'}^\prime$. Si~$S'$ es conexo,
este enunciado se interpreta del modo siguiente. Sea~$C'_0$ la categoría de
los recubrimientos étales de~$S'$ en los que~$\Gamma$ actúa de modo compatible
con su acción sobre~$S'$ (pero sin satisfacer necesariamente la condición
anterior sobre los grupos de inercia de los puntos de~$S'$). Se ve fácilmente
que es una categoría galoisiana (\Ref{V.5}) y que, para
\marginpar{251}
cada punto geométrico~$a'$ de~$S'$, el funtor fibra
$X'\mto X_{a'}'$ sobre~$C'_0$ es un funtor fundamental. Sea
~$\pi_1(S',\Gamma;a')=G$ el grupo de automorfismos de este funtor, provisto de
su topología habitual. Se tiene entonces una sucesión exacta canónica
$$
e\to \pi_1(S',a')\to G\to \Gamma\to e
$$
(caso particular de~\Ref{V.6.13}, donde se toma por~$S$ el recubrimiento
trivial~$S'\times\Gamma$ de~$S'$ definido por~$\Gamma$, sobre el cual se hace
actuar~$\Gamma$ de modo evidente). Además, para cada punto geométrico~$b'$ de
~$S'$, se tiene un isomorfismo
$\pi_1(S',\Gamma;b')\to G=\pi_1(S',\Gamma;a')$, definido salvo un
automorfismo interior procedente de~$\pi_1(S',a')$; y, como hay un
homomorfismo evidente de~$\Gamma_{b'}$ al primer miembro, se obtiene un homomorfismo
$$
u_{b'}\colon \Gamma_{b'}\to G \quoi,
$$
definido salvo un automorfismo interior (procedente
\ifthenelse{\boolean{orig}}
{de~$\pi_1(S',a')$,}
{de~$\pi_1(S',a')$),}
cuyo compuesto con el homomorfismo canónico~$G\to\Gamma$ es la inclusión
canónica~$\Gamma_{b'}\to\Gamma$. Hecho esto, \emph{el grupo fundamental
~$\pi_1(S,a)$ es canónicamente isomorfo al grupo cociente de
~$G=\pi_1(S',\Gamma;a')$ por el subgrupo normal cerrado engendrado por las
imágenes de los homomorfismos~$\Gamma_{b'}\to G$}. En particular, la imagen de
~$\pi_1(S',a')$ en~$\pi_1(S,a)$ es un subgrupo normal, y el cociente
correspondiente es isomorfo a un cociente de~$\Gamma$. Además, se puede
reducir el número de «relaciones» introducidas considerando, para cada
~$g\in\Gamma$, $g\neq e$, el subpreesquema~$S'_g$ de coincidencias de los
automorfismos
\ifthenelse{\boolean{orig}}{$\id_S$ y~$g$ de~$S$}{$\id_{S'}$ y~$g$ de~$S'$},
eligiendo un punto geométrico~$b'_{g,i}$ en cada componente conexa de~$S'_g$
y, después, uno de los homomorfismos correspondientes
~$\pi_1(S',\Gamma;b'_{g,i})\to G$, de donde se obtienen levantamientos
~$\bar{g}_i$ de~$g$ en~$G$. Basta entonces tomar el cociente de~$G$ por el
subgrupo normal cerrado de~$G$ engendrado por los~$\bar{g}_i$.

Cuando~$a'$ es fijo por~$\Gamma$, se ve fácilmente que~$\Gamma$ actúa de modo
natural sobre~$\pi_1(S',a')$ y que~$G$ se identifica con el producto
semidirecto correspondiente. Identificando entonces~$\Gamma$
\ifthenelse{\boolean{orig}}
{a}
{con}
un subgrupo de~$G$, se ve que entre las relaciones
\marginpar{252}
introducidas más arriba, al tomar~$b'=a'$, aparecen «$g=e$» para
~$g\in\ifthenelse{\boolean{orig}}{I}{\Gamma}$. Por tanto, \emph{si~$S'$ tiene
un punto geométrico~$a'$ fijo por}~$\ifthenelse{\boolean{orig}}{I}{\Gamma}$
(es decir, un punto~$s'$ cuyo grupo de inercia es~$\Gamma$), \emph{entonces
~$\pi_1(S,a)$ es un grupo cociente del grupo cociente de tipo galoisiano de
~$\pi_1(S',a')$ obtenido «haciendo triviales» las acciones de~$\Gamma$ sobre
~$\pi_1(S',a')$; y es incluso isomorfo a este último grupo si se supone que,
para cada~$g\in\ifthenelse{\boolean{orig}}{G}{\Gamma}$, el lugar de inercia
~$S_g'$ es conexo y, por tanto, pasa por el punto subyacente de~$a'$}.
Esta última afirmación está, en efecto, contenida en la segunda descripción
dada más arriba de las relaciones que se han de imponer a~$G$.

Este último resultado se aplica, en particular, si se toma por~$S'$ la
potencia cartesiana~$X^n$ de un preesquema conexo sobre un cuerpo
algebraicamente cerrado, y por~$\Gamma$ el grupo simétrico
~$\Gamma={\goth {S}}_n$, que actúa de la manera habitual, con lo que~$S$ es la
producto simétrico \ifthenelse{\boolean{orig}}{\nieme}{de grado~$n$}
\ifthenelse{\boolean{orig}}{de~$S$}{de~$X$}.
Tomando por~$a'$ un punto geométrico cuyo punto subyacente pertenece a la
diagonal, se cumplen las condiciones anteriores, pues todos los lugares de
inercia~$S_g'$ contienen la diagonal. Utilizando el hecho, demostrado en el
exposé siguiente, de que si~$X$ es propio y conexo sobre~$k$, el grupo
fundamental de~$X^n$ se identifica con~$\pi_1(X)^n$, se obtiene el siguiente
resultado curioso: \emph{si~$X$ es propio y conexo sobre un cuerpo
algebraicamente cerrado~$k$, el grupo fundamental de su producto simétrico
\ifthenelse{\boolean{orig}}{\nieme}{de grado~$n$}, $n\geq 2$, es isomorfo a la abelianización del grupo fundamental
de~$X$}. (Ignoro si se conoce el hecho análogo en topología algebraica; debería
poder establecerse mediante el mismo método de descenso). Tomemos, por
ejemplo, por~$X$ la curva racional~$X=\PP^1_k$; se obtiene una
\ifthenelse{\boolean{orig}}{\Nieme}{enésima} demostración de
que~$\PP^r_k$ es simplemente conexo, utilizando que~$\PP^1_k$ lo es. Tomemos
ahora por~$X$ una curva no singular sobre~$k$ y sea~$n\geq 2g-1$, de modo que
~$\Symm^n(X)$ está fibrado sobre la jacobiana~$J$, con fibras espacios
proyectivos, y tiene, por tanto (como se verá mediante los resultados de los
dos exposés siguientes), el mismo grupo fundamental que~$J$. Se recupera así,
sin una reducción por etapas, el conocido hecho de que \emph{el grupo fundamental de la
jacobiana de~$X$ es isomorfo a la abelianización del grupo fundamental de
~$X$}.
\end{remarque}

\section[Una sucesión exacta fundamental]{Una sucesión exacta fundamental. Descenso por morfismos con
fibras relativamente conexas}
\label{IX.6}
\marginpar{253}

\begin{theoreme}
\label{IX.6.1}
Sean~$S$ el espectro de un anillo artiniano~$A$ con cuerpo residual~$k$,
~$\overline{k}$ una clausura algebraica de~$k$,~$X$ un~$S$-preesquema,
~$X_0=X\otimes_A k$,~$\overline{X}_0=X\otimes_A\overline{k}$,
~$\overline{a}$ un punto geométrico
\ifthenelse{\boolean{orig}}{de~$\overline{X}$}{de~$\overline{X}_0$},
~$a$ su imagen en~$X$ y~$b$ su imagen en~$S$. Supongamos que~$X_0$ es
cuasicompacto y \emph{geométricamente conexo} sobre~$k$ (N.B.: si~$X$ es
propio sobre~$S$, esto significa que~$\H^0(X_0,\cal{O}_{X_0})$ es un anillo
artiniano \emph{local} cuyo cuerpo residual es \emph{radicial} sobre~$k$).
Entonces la sucesión de homomorfismos canónicos
$$
e\mto \pi_1(\overline{X}_0,\overline{a})\to
\pi_1(X,a)\to\pi_1(S,b)\to e
$$
es exacta, y se tiene
$$
\pi_1(S,b)\isomfrom \pi_1(k,\overline{k})=\text{grupo de Galois
de~$\overline{k}$ sobre~$k$.}
$$
\end{theoreme}

Como los grupos fundamentales no cambian al eliminar los elementos
nilpotentes, podemos suponer que~$A=k$, lo cual hace ya evidente el último
isomorfismo. Sea~$k'$ la clausura separable de~$k$ en~$\overline{k}$,
consideremos~$X'=X\otimes_k k'$ y sea~$a'$ la imagen de~$\overline{a}$
en~$X'$. Se tiene una sucesión de homomorfismos canónicos
$$
e\mto \pi_1(\overline{X}_0,\overline{a})\to
\pi_1(X',a')\to\pi_1(S',b')\to e
$$
(donde~$S'=\Spec(k')$). Por último, el diagrama
$$
\xymatrix{ S & X \ar[l] & \overline{X}_0 \ar[l] \\ S' \ar[u] &
X'\ar[u] \ar[l] & \overline{X}_0 \ar[l] \ar@{=}[u]. }
$$
da un homomorfismo canónico de esta sucesión en la correspondiente a~$X/k$.
Por otra parte, se ve que este homomorfismo de sucesiones de grupos es un
isomorfismo, como resulta de~\Ref{IX.4.11}. Quedamos, pues, reducidos a probar
que la segunda sucesión es exacta, es decir, podemos suponer que~$k$ es
\emph{perfecto}. Sean entonces~$k_i$ las subextensiones galoisianas finitas
de~$k$ en~$\overline{k}$, pongamos~$X_i=X\otimes_k k_i$ y sea~$a_i$ la
imagen de~$\overline{a}$ en~$X_i$. Dejamos al lector
\ifthenelse{\boolean{orig}}
{\ignorespaces}
{la tarea}
de comprobar que el homomorfismo natural
$$
\pi_1(\overline{X}_0,\ifthenelse{\boolean{orig}}{a}{\overline{a}})
\isomto \varprojlim\pi_1(X_i,a_i)
$$
\marginpar{254}%
es un isomorfismo. Esto significa simplemente que todo recubrimiento étale
de~\ifthenelse{\boolean{orig}}{$\overline{X}$}{$\overline{X}_0$} procede de
un recubrimiento étale de algún~$X_i$, y que este último es esencialmente
único, salvo pasar a un~$X_j$,~$j\geq i$.
Por otra parte, sea~$\pi_i$ el grupo de Galois de~$k_i$ sobre~$k$, es decir,
el grupo opuesto al grupo de los~$S$-automorfismos de~$S_i=\Spec(k_i)$.
Como el funtor~$S'\mto X\times_S S'$ de los recubrimientos étale de~$S$ en
los recubrimientos étale de~$X$ es plenamente fiel~\eqref{IX.3.4}, se
sigue que~$\pi_i$ es también isomorfo al grupo opuesto al grupo de los
~$X$-automorfismos del recubrimiento principal conexo~$X_i$ de~$X$. Resulta,
por tanto, de~\Ref{V.6.13} que se tiene una sucesión exacta
$$
e\mto \pi_1(X_i,a_i)\to \pi_1(X,a)\to\pi_i\to e.
$$
Pasando al límite proyectivo sobre~$i$ en estas sucesiones exactas, se obtiene
una sucesión exacta (pues estamos en la categoría de los grupos de tipo
galoisiano), que no es otra que la sucesión considerada en~\Ref{IX.6.1}.
Esto termina la demostración.

La traducción de la exactitud por la derecha en~\Ref{IX.6.1} al lenguaje
geométrico es la siguiente:

\begin{corollaire}
\label{IX.6.2}
Con las notaciones precedentes, sea~$X'$ un recubrimiento étale de~$X$ y sea
~$\overline{X'}_0$ el recubrimiento étale correspondiente de
~$\overline{X}_0$. Las condiciones siguientes son equivalentes:
\begin{enumerate}
\item[(i)] Existe un~$S'$ étale sobre~$S$ y un~$X$-isomorfismo
$X'\isomto X\times_S S'$ (entonces~$S'$ queda determinado salvo isomorfismo
único en virtud de~\Ref{IX.3.4}).
\item[(ii)] $\overline{X}_0'$ está completamente descompuesto sobre
~$\overline{X}_0$.

Si~$X'$ es conexo, estas condiciones equivalen también a:
\item[(ii~bis)] $\overline{X}_0'$ posee una sección sobre~$\overline{X}_0$.
\end{enumerate}
\end{corollaire}

(N.B.: Este último complemento es esencial; la equivalencia de~(i) y~(ii)
solo significa que~$\pi_1(S,b)$ es el grupo cociente de~$\pi_1(X,a)$ por el
subgrupo
\marginpar{255}
normal cerrado engendrado por la imagen de
~$\pi_1(\overline{X}_0,\overline{a})$, y no por esta imagen misma). Bajo las
condiciones precedentes, diremos que~$X'$ es un recubrimiento
\emph{geométricamente trivial}
\index{geométricamente trivial (recubrimiento)|hyperpage}%
de~$X$.

\begin{remarque}
\label{IX.6.3}
En el enunciado~\Ref{IX.6.1} no puede sustituirse~$\overline{k}$ por una
extensión algebraicamente cerrada cualquiera de~$k$, ni siquiera si se supone
ya que~$k$ es algebraicamente cerrado. En otras palabras, no es cierto en
general que, si~$X$ es un esquema algebraico conexo sobre un cuerpo
algebraicamente cerrado~$k$, su grupo fundamental permanezca inalterado al
sustituir~$k$ por una extensión algebraicamente cerrada. Esto es falso ya, por
ejemplo, en característica~$p>0$ para la recta afín sobre~$k$, a causa de los
fenómenos de «ramificación superior» en el punto del infinito, que introducen
una estructura «continua» en el grupo fundamental. Veremos, sin embargo, en el
exposé siguiente que tales fenómenos no pueden producirse si~$X$ es
\emph{propio} sobre~$k$. Demostraremos también por métodos trascendentes que
ocurre lo mismo si~$k$ es de característica cero.
\end{remarque}

\begin{corollaire}
\label{IX.6.4}
Supongamos que~$a$ está localizado en un punto~$x\in X$ racional sobre~$k$ (o,
más generalmente, cuyo cuerpo residual es radicial sobre~$k$). Entonces la
sucesión exacta~\Ref{IX.6.1} es escindida.
\end{corollaire}

Podemos suponer~$S=\Spec(k)$. Si~$x$ es racional sobre~$k$, corresponde a una
sección~$S\to X$ de~$X$ sobre~$S$, que lleva~$b$ en~$a$ y define un
homomorfismo~$\pi_1(S,b)\to\pi_1(X,a)$ que es la escisión buscada. Si
~$\kres(x)$ es radicial sobre~$k$, nos reducimos al caso precedente mediante
la extensión de base~$\Spec(\kres(x))\to\Spec(k)$.

\begin{theoreme}
\label{IX.6.5}
Sean~$f\colon X\to S$ un morfismo propio y sobreyectivo de presentación
finita, con fibras geométricamente conexas,~$X'$ un preesquema de presentación
finita y propio sobre~$X$,~$s$ un punto de~$S$,~$F=X_s$ la fibra de~$X$ en
~$s$, y~$F_1'$ una componente conexa de la fibra~$F'=X'_s$ de~$X'$ en~$s$.
Para que existan un entorno abierto~$X_1'$ de~$F_1'$ en~$X'$, un
~$S$-esquema étale~$S_1'$ y un~$X$-isomorfismo
$X_1'\isomto S_1'\times_S X$, es necesario y suficiente que~$X'$ sea étale
sobre~$X$ en los puntos de~$F_1'$ y que~$F_1'$ sea un recubrimiento
geométricamente trivial de~$F$.
\end{theoreme}

La
\marginpar{256}
necesidad de la condición es trivial, por lo que resta probar la suficiencia.
Nos reducimos fácilmente al caso en que~$S$ es noetheriano. Consideremos la
factorización de Stein~$X\to T\to S$ de~$f$, donde~$T$ es el espectro del
álgebra~$f_*(\cal{O}_X)$ sobre~$S$. Como las fibras de~$X$ sobre~$S$ son
geométricamente conexas y~$f$ es sobreyectivo, el morfismo~$T\to S$ es finito,
sobreyectivo y radicial; por consiguiente, según~\eqref{IX.4.10}, todo~$T'$
étale sobre~$T$ procede por imagen inversa de un~$S'$ étale sobre~$S$.
Esto nos reduce a
\ifthenelse{\boolean{orig}}
{probar~\Ref{IX.6.5}.}
{probar~\Ref{IX.6.5}}
reemplazando allí~$S$ por~$T$, es decir, al caso en que suponemos
~$f_*(\cal{O}_X)=\cal{O}_S$. Consideremos entonces la factorización de Stein
~$X'\to S'\to S$ del morfismo propio~$h\colon X'\to S$, donde~$S'$ es el
espectro del álgebra~$h_*(\cal{O}_{X'})$. Los morfismos~$X'\to X$ y
~$X'\to S'$ definen un morfismo canónico
$$
X'\to X\times_S S',
$$
y nuestra afirmación está contenida en la siguiente:

\begin{corollaire}
\label{IX.6.6}
Sea~$f\colon X\to S$ un morfismo propio de preesquemas localmente
noetherianos tal que~$f_*(\cal{O}_X)=\cal{O}_S$, y sea~$X'$ un preesquema
propio sobre~$X$. Consideremos la factorización de \emph{Stein}
~$X'\to S'\to S$ de~$X'\to S$ y el morfismo canónico
~$X'\to X\times_S S'$. Sean~$s$ un punto de~$S$ y~$s'$ un punto de~$S'$
situado sobre~$s$, correspondiente a una componente conexa~$F_1'$ de la
fibra~$X'_s$ de~$X'$ en~$s$. Para que el morfismo
~$X'\to X\times_S S'$ sea un isomorfismo sobre un entorno abierto~$U'$ de
~$s'$ étale sobre~$S$, es necesario y suficiente que~$X'$ sea étale sobre
~$X$ en los puntos de~$F_1'$ y que~$F_1'$ sea un recubrimiento
geométricamente trivial de la fibra~$F=X_s$.
\end{corollaire}

La necesidad sigue siendo trivial, por lo que resta probar la suficiencia.
La conclusión significa también que a) el morfismo deducido de
~$X'\to X\times_S S'$ por el cambio de base
~$\Spec(\widehat{\cal{O}}_{s'})\to S'$ es un isomorfismo; b)~$S'$ es étale
sobre~$S$ en~$s'$, es decir,~$\widehat{\cal{O}}_{s'}$ es étale sobre
~$\widehat{\cal{O}}_s$. Bajo esta forma, se ve que la conclusión es
invariante por el cambio de base~$\Spec(\widehat{\cal{O}}_s)\to S$. Como las
hipótesis son también estables por este cambio de base, podemos suponer que
~$S$ es el espectro de un anillo local noetheriano completo. Podemos además
suponer evidentemente que~$X'$ es conexo, lo cual implica aquí que
~$S'=\Spec(\cal{O}_{s'})$ y~$F'=F_1'$.

Como
\marginpar{257}
el conjunto de los puntos de~$X'$ en los que~$X'$ es étale sobre~$X$ es
abierto y contiene la fibra~$X'_{s'}=F'$, y como~$X'$ es propio sobre~$S$, se
sigue que~$X'$ es étale sobre~$X$. Puesto que induce sobre~$F=X_s$ un
recubrimiento étale isomorfo a un~$F\otimes_{\kres(s)}L$, donde~$L$ es étale
sobre~$\kres(s)$, resulta de~\Ref{IX.1.10} que es isomorfo a un
recubrimiento de la forma~$X\times_S T$, con~$T$ étale sobre~$S$ (N.B.: aquí
también basta utilizar que el funtor de~\Ref{IX.1.10} es plenamente fiel,
lo cual resulta de que todo isomorfismo formal de haces coherentes sobre~$X$
procede de un isomorfismo de esos haces). Por tanto, si~$T$ está definido por
el álgebra~$B$ finita sobre~$A$,~$X'$ se identifica con el espectro del
álgebra~$\cal{O}_X\otimes_A B$ sobre~$X$, de donde resulta inmediatamente,
puesto que~$f_*(\cal{O}_X)=\cal{O}_S$, que~$h_*(\cal{O}_{X'})$ está definido
por~$B$. Así, el homomorfismo canónico~$X'\to X\times_S S'$ no es otro que el
isomorfismo considerado~$X'\isomto X\times_S T$. Esto termina la
demostración.

\begin{corollaire}
\label{IX.6.7}
Bajo las condiciones de~\Ref{IX.6.5}, para que existan un preesquema~$S'$
étale sobre~$S$ y un~$X$-isomorfismo~$X'\isomto X\times_S S'$, es
necesario y suficiente que~$X'$ sea étale sobre~$X$ y que, para todo
~$s\in S$,
\ifthenelse{\boolean{orig}}{$X'_{s'}$}{$X'_s$}
sea un recubrimiento geométricamente trivial de~$X_s$.
\end{corollaire}

En efecto, si esto ocurre,~$X'$ es unión de abiertos~$X_i'$ isomorfos a
imágenes inversas de~$S_i'$ étale sobre~$S$. Se ve entonces fácilmente que
estos~$S_i'$ se pegan en un~$S'$ étale sobre~$S$ y que se obtiene un
isomorfismo~$X'\isomto X\times_S S'$. Puede decirse, por ejemplo, que los
~$X_i'$ están provistos de datos de descenso relativos a~$X\to S$, que se
pegan necesariamente en un dato de descenso sobre~$X'$ entero relativo a
~$X\to S$. Y, como este último es efectivo sobre los~$X_i'$, se sigue
fácilmente (gracias a un sorites omitido en el número~\Ref{IX.4}) que es
efectivo. También puede enunciarse~\Ref{IX.6.7} en la forma siguiente:

\begin{corollaire}
\label{IX.6.8}
Sea~$f\colon X\to S$ un morfismo propio y sobreyectivo de presentación
finita, con fibras geométricamente conexas. Entonces~$f$ es un morfismo de
descenso efectivo para la categoría fibrada de los preesquemas étale y
finitos sobre otros. El funtor~$S'\mto X\times_S S'$ induce una equivalencia
de la categoría de los preesquemas étale y finitos sobre~$S$ con la
categoría de los preesquemas
\marginpar{258}
étale y finitos sobre~$X$ que inducen sobre cada fibra~$X_s$ un
recubrimiento geométricamente trivial.
\end{corollaire}

\begin{remarque}
\label{IX.6.9}
Sea~$f\colon X\to S$ un morfismo propio y sobreyectivo, con
\ifthenelse{\boolean{orig}}{$Y$}{$S$}
localmente noetheriano. Entonces~$f$ se factoriza en un morfismo~$X\to S'$
que satisface la hipótesis de~\Ref{IX.6.8} y un morfismo finito y sobreyectivo
~$S'\to S$ al que se aplica~\Ref{IX.4.7}; por tanto, se expresa como
composición de dos
morfismos que son \emph{morfismos de descenso efectivo universales} para la
categoría fibrada de los preesquemas étale y finitos sobre otros. Puede
concluirse que el propio~$f$ es un morfismo de descenso efectivo universal
para la categoría fibrada considerada. Se recupera así~\Ref{IX.4.12} por un
método diferente.
\end{remarque}

\begin{remarque}
\label{IX.6.10}
La conclusión de~\Ref{IX.6.7} deja de ser válida si se sustituye la hipótesis
de que~$f$ es propio por la siguiente:~$X$ es de tipo finito sobre~$S$ y
admite una sección sobre~$S$ (por tanto,~$f$ es universalmente submersivo y
un morfismo de descenso para la categoría fibrada de los preesquemas étale
sobre otros), incluso cuando~$S$ es el espectro de un anillo de valoración
discreta y~$X'$ es un recubrimiento étale de~$X$. Para verlo, se parte de un
~$Z$ propio sobre~$S$, cuya fibra genérica
es una curva racional no singular y cuya fibra especial~$Z_0$ consta de dos
rectas secantes. Por ejemplo, si~$t$ es un uniformizante del anillo de
valoración~$A$, se toma el subesquema cerrado~$Z$ de~$\PP_A^2$ definido por
la ecuación homogénea~$x^2+y^2+tz^2=0$ (coordenadas homogéneas~$x$,~$y$,
~$z$). Se toma por~$X$ el complementario del punto singular~$a$ de~$Z_0$ en
la unión~$Z\cup\PP_k^2$. Las fibras de~$X$ son~$\PP_k^1$ y
~$\PP_k^2-a$; son, por tanto, geométricamente simplemente conexas (es decir,
todo recubrimiento étale de una de estas fibras es geométricamente trivial).
Sin embargo,
\marginpar{259}
procediendo como en el número~\Ref{IX.4}, se construyen fácilmente
recubrimientos étale de~$X$ que no proceden de recubrimientos étale de~$S$,
pegando recubrimientos triviales de~$Z-a$ y de~$\PP_k^2-a$. En cambio, es
posible que la conclusión de~\Ref{IX.6.7} subsista si se sustituye la
hipótesis de que el morfismo sea propio por la de que~$X$ sea
\emph{universalmente abierto} y
de presentación finita sobre~$S$\kern1pt\footnote{Esto ya está demostrado,
siendo~\ifthenelse{\boolean{orig}}{$g$}{$f$} solamente universalmente abierto
y sobreyectivo; \cf SGA~4 XV~1.15.}.
Esto es cierto al menos si se supone que las fibras de~$X$ sobre~$S$ son
geométricamente irreducibles, y no solo geométricamente conexas. Señalemos
únicamente que, en esta cuestión, podemos reducirnos al caso en que~$S$ es
el espectro de un anillo de valoración discreta completo con cuerpo residual
algebraicamente cerrado.
\end{remarque}

La interpretación de~\Ref{IX.6.7} en términos del grupo fundamental es la
siguiente:

\begin{corollaire}
\label{IX.6.11}
Sea~$f\colon X\to S$ un morfismo propio y sobreyectivo de presentación
finita, con fibras geométricamente conexas. Supongamos que~$X$, y por tanto
~$S$, es conexo. Sean~$a$ un punto geométrico de~$X$,~$b$ su imagen en~$S$ y,
para todo~$s\in S$, elijamos una clausura algebraica
~$\overline{\kres(s)}$ de~$\kres(s)$, un punto geométrico~$a_s$ de
\ifthenelse{\boolean{orig}}{$S_s$}{$X_s$}
con valores en esta extensión y una clase de caminos desde~$a_s$ hasta~$a$, de donde
se obtiene un homomorfismo~$\pi_1(\overline X_s,a_s)\to\pi_1(X,a)$, donde
~$\overline X_s=X_s\otimes_{\kres(s)}\overline{\kres(s)}$. Entonces el
homomorfismo~$\pi_1(X,a)\to\pi_1(S,b)$ es sobreyectivo, y su núcleo es el
subgrupo normal cerrado de~$\pi_1(X,a)$ engendrado por las imágenes de los
~$\pi_1(\overline X_s,a_s)$.
\end{corollaire}

\begin{remarque}
\label{IX.6.12}
Bajo las condiciones de~\Ref{IX.6.7}, suponiendo~$S$ noetheriano, se ve
fácilmente que el conjunto de los puntos~$s\in S$ tales que
\ifthenelse{\boolean{orig}}{$S'_{s'}$}{$X'_s$}
sea geométricamente trivial sobre~$X_s$ es constructible; si~$X'$ es propio
sobre~$X$, es incluso abierto, como se ve en~\Ref{IX.6.6}. Esto permite, por
tanto, si~$S$
\ifthenelse{\boolean{orig}}{y}{es}
un preesquema de Jacobson (por ejemplo, de tipo finito sobre un cuerpo), o
\marginpar{260}
si~$X'$ es propio sobre~$X$, limitarse, para verificar las condiciones
de~\Ref{IX.6.7}, a los puntos~$s$ de~$S$ que son \emph{cerrados}. Del mismo
modo, en~\Ref{IX.6.11} basta entonces tomar los
~$\pi_1(\overline X_s,a_s)$ para los puntos de~$S$ que son cerrados.
\end{remarque}

\begin{sgabibliography}{D}{IX.7}

\bibitem[D]{IX.D} J\ptbl \textsc{Giraud}, Méthode de la
\ifthenelse{\boolean{orig}}
{d'escente,}
{descente,}
Mémoire \no 2 de la Soc.\ Math.\ Fr.,~1964.

\bibitem[1]{IX.1} A\ptbl \textsc{Grothendieck}, Géométrie Algébrique et
Géométrie Formelle, Séminaire Bourbaki t\ptbl 11, 1959, \No 182.
\end{sgabibliography}

\chapter{Teoría de la especialización del grupo fundamental}
\label{X}
\marginpar{261}
En la presente exposición nos limitamos al estudio del grupo fundamental de
las fibras geométricas en un morfismo \emph{propio}, \ie del grupo fundamental
de un esquema algebraico \emph{propio} variable. En una exposición posterior
generalizaremos la técnica empleada a los recubrimientos étale
\emph{moderadamente ramificados} «en el infinito». Esto nos dará, por ejemplo,
una solución del «problema de los tres puntos» en el caso de los recubrimientos
galoisianos de orden primo con la característica (\ie una determinación de los
recubrimientos galoisianos de la recta proyectiva~$\PP_k^1$, ramificados a lo
sumo en tres puntos dados y moderadamente ramificados en esos puntos), así como
de sus variantes evidentes.

\section[La sucesión exacta de homotopía para un morfismo propio]
{La sucesión exacta de homotopía para un morfismo propio y separable}
\label{X.1}

\begin{definition}
\label{X.1.1}
Un preesquema~$X$ sobre un cuerpo~$k$ se dice \emph{separable}, o
\emph{separable sobre~$k$}, si para toda extensión~$K$ de~$k$,
$X\otimes_k K$ es reducido. Si~$f\colon X\to Y$ es un morfismo de
preesquemas, se dice que~$f$ es \emph{separable}, o que~$X$ \emph{es separable
sobre~$Y$}, si~$X$ es plano sobre~$Y$ y si, para todo~$y\in Y$, la fibra
$X\otimes_Y \kres(y)$ es separable sobre~$\kres(y)$.
\end{definition}

Si~$X$ es un preesquema sobre un cuerpo~$k$, decir que es separable significa
también que es \emph{reducido} y que los cuerpos~$\kres(x)$, para~$x$ punto
genérico de una componente irreducible de~$X$, son extensiones separables
de~$k$. Si~$k$ es perfecto, equivale, pues, a decir que~$X$ es reducido.
Obsérvese que, si~$X$ es separable sobre~$Y$, entonces, para todo cambio de
base~$Y'\to Y$,~$X'=X\times_Y Y'$ es separable sobre~$Y'$. Puede probarse
también, bajo hipótesis de finitud convenientes, que el compuesto de dos
morfismos separables es
\marginpar{262}
un morfismo separable. Solo lo necesitaremos en la forma siguiente:
\emph{si~$X$ es separable sobre~$Y$ y~$X'$ es étale sobre~$X$, entonces~$X'$
es separable sobre~$Y$.} Es, en efecto, una consecuencia inmediata de las
definiciones y de~\Ref{I.9.2}. Por otra parte, la hipótesis «morfismo
separable» nos servirá mediante la proposición siguiente:

\begin{proposition}
\label{X.1.2}
Sea~$f\colon X\to Y$ un morfismo propio y separable, con~$Y$ localmente
noetheriano, y consideremos su factorización de \emph{Stein}
$X\lto{f'}Y'\dpl\to Y$, donde~$f_*'(\cal{O}_X)=\cal{O}_{Y'}$, siendo~$Y'$
finito sobre~$Y$ e isomorfo al espectro del álgebra~$f_*(\cal{O}_X)$. Entonces
~$Y'$ es un \emph{recubrimiento étale} de~$Y$.
\end{proposition}

Esta proposición figurará en (EGA~III 7)\footnote{\Cf EGA~III 7.8.10 (i)}.
Indiquemos el principio de la demostración. Se reduce fácilmente al caso en que
~$Y$ es el espectro de un anillo local completo~$A$ y, haciendo además una
extensión finita plana conveniente de este último (correspondiente a una
extensión residual conveniente), puede suponerse que las componentes conexas
de la fibra del punto cerrado~$y$ son geométricamente conexas, lo que significa
también que~$\H^0(X_y,\cal{O}_{X_y})$ se descompone en un producto de cuerpos
idénticos a~$k=\kres(y)$. Suponiendo entonces~$X$ conexo, como es lícito, se
tendrá~$\H^0(X_y,\cal{O}_{X_y})=k$, por lo que el homomorfismo
$A\to \H^0(X_y,\cal{O}_{X_y})$ es \emph{sobreyectivo}. Se concluye, mediante
una proposición general (de tipo Künneth), que~$f_*(\cal{O}_X)$ está definido
por un módulo~$B$ sobre~$A$ que es libre sobre~$A$, y que
$B/\goth{m} B\to \H^0(X_y,\cal{O}_{X_y})=k$ es biyectivo. Por tanto, en este
caso~$B$ es un álgebra étale sobre~$A$, lo cual termina la demostración.

\begin{theoreme}
\label{X.1.3}
Sea~$f\colon X\to Y$ un morfismo propio y separable, con~$Y$ localmente
noetheriano y conexo, y supongamos que~$f_*(\cal{O}_X)=\cal{O}_Y$ (lo que
implica que las fibras de~$X$ sobre~$Y$ son geométricamente conexas, y
recíprocamente gracias a~\Ref{X.1.2}). Sean~$y$ un punto de~$Y$,
$\overline{\kres(y)}$ una clausura algebraica de~$\kres(y)$ y
$\overline X_y=X_y\otimes_{\kres(y)}\overline{\kres(y)}$. Sean, por último,
~$X'$ un recubrimiento étale \emph{conexo} de~$X$ y
$\overline X_y'=X_y'\otimes_{\kres(y)}\overline{\kres(y)}$. Para que exista
un recubrimiento étale~$Y'$ de~$Y$ y un~$X$-isomorfismo
\marginpar{263}
$X'\isomto X\times_Y Y'$, es necesario y suficiente que~$\overline{X'}_y$
admita una sección sobre~$\overline X_y$.
\end{theoreme}

Poniendo~$Y'=\Spec(h_*(\cal{O}_{X'}))$ (donde~$h\colon X'\to Y$ es el
morfismo compuesto~$X'\to X\to Y$), basta probar que el~$Y$-morfismo canónico
$$
X'\to X\times_Y Y'
$$
es un \emph{isomorfismo} y que~$Y'$ es étale sobre~$Y$. Ahora bien, ya sabemos
por~\Ref{X.1.2} que~$Y'$ es étale sobre~$Y$; por tanto,
~$X\times_Y Y'$ es étale sobre~$X$, y por ello el morfismo
~$X'\to X\times_Y Y'$ es también étale (\Ref{I.4.8}). Además,~$Y'$ es
conexo como imagen de~$X'$, que lo es; por tanto,~$X\times_Y Y'$ es conexo,
pues~$X$ tiene fibras conexas sobre~$Y$ (\Ref{IX.3.4} y
\Ref{V.6.9}~(iii)). Así, para probar que~$X'\to X\times_Y Y'$ es un
isomorfismo, basta ver que su grado de proyección en \emph{un} punto de
~$X\times_Y Y'$ es igual a~1. Esto resulta fácilmente de la hipótesis de que
~$\overline{X'}_y$ admite una sección sobre~$\overline X_y$, bien utilizando
\Ref{IX.6.6}, bien, más simplemente, observando que basta probar la
existencia de tal punto en~$X\times_Y Y'$ después del cambio de base
\ifthenelse{\boolean{orig}}
{$\Spec(\overline k)\to Y$}
{$\Spec(\overline{\kres(y)})\to Y$}, donde esto es evidente.
Esto termina la demostración de~\Ref{X.1.3}.

Teniendo en cuenta~\Ref{IX.3.4} y el diccionario~\Ref{V.6.9} y
\Ref{V.6.11}, puede ponerse~\Ref{X.1.3} en la forma equivalente siguiente:

\begin{corollaire}
\label{X.1.4}
Con las notaciones precedentes para~$f\colon X\to Y$ y~$\overline X_y$, sea
~$\bar{a}$ un punto geométrico de~$\overline X_y$,~$a$ su imagen en~$X$ y
~$b$ su imagen en~$Y$. Entonces la sucesión siguiente de homomorfismos de
grupos es \emph{exacta}:
$$
\pi_1(\overline X_y,\bar{a})\to \pi_1(X,a)\to \pi_1(Y,b)\to e.
$$
\end{corollaire}

\begin{remarques}
\label{X.1.5}
Obsérvese que la demostración de~\Ref{X.1.3} hace intervenir de manera esencial
~\Ref{X.1.2} y, por ello, el «primer teorema de comparación» en geometría
algebraico-formal. En cambio, la teoría del descenso de la exposición~\Ref{IX}
\marginpar{264}
solo ha intervenido mediante~\Ref{IX.3.4}, cuya demostración directa es
fácil en el caso de un morfismo \emph{propio}~$f\colon X\to Y$ tal que
~$f_*(\cal{O}_X)=\cal{O}_Y$. Sea, en efecto,~$Y'$ étale sobre~$Y$ y
supongamos que~$X'=X\times_Y Y'$ sea unión disjunta de dos abiertos no vacíos;
probemos que ocurre lo mismo con~$Y'$. En efecto, se tendrá
~$Y'=\Spec(\cal{A})$, y por tanto~$X'=\Spec(\cal{B})$, con
~$\cal{B}=\cal{A}\otimes_{\cal{O}_Y}\cal{O}_X$; la descomposición de~$X'$ en
unión disjunta corresponde a una descomposición de~$\cal{B}$ en producto de
dos álgebras no nulas~$\cal{B}_1$ y~$\cal{B}_2$. Como
~$f_*(\cal{O}_X)=\cal{O}_Y$, se concluye fácilmente que
~$f_*(\cal{B})=\cal{A}$; por tanto,~$\cal{A}$ será producto de dos álgebras
(también no nulas, pues sus secciones unidad no son nulas)
~$f_*(\cal{B}_1)$ y~$f_*(\cal{B}_2)$, c.q.d.
\end{remarques}

\setcounter{subsection}{5}% Preserve the source-visible paragraph number X.1.6 in the standalone class.
\subsection{}
\ifthenelse{\boolean{orig}}
{\label{X.I.6}}
{\label{X.1.6}}
\setcounter{proposition}{6}% The source paragraph occupies item 1.6 in the shared theorem sequence.
Supongamos todavía que~$f$ es propio y separable, pero no hagamos ya ninguna
hipótesis sobre~$f_*(\cal{O}_X)$, que corresponderá a un recubrimiento étale
~$Y'$ bien determinado de~$Y$, provisto además, sobre~$b$, del punto~$b'$
imagen de~$a$. Aplicando entonces~\Ref{X.1.4} al morfismo canónico~$X\to Y'$,
y suponiendo~$f$ sobreyectivo, la sucesión exacta~\Ref{X.1.4} se sustituye por
la siguiente, análoga a la sucesión exacta de homotopía de los espacios
fibrados en topología algebraica:
$$
\pi_1(\overline X_y,\overline a)\to \pi_1(X,a)\to \pi_1(Y,b)\to
\pi_0(\overline X_y,\overline a)\to \pi_0(X,a)\to \pi_0(Y,b)\to e.
$$
Por supuesto, en~\Ref{X.1.4} no puede afirmarse en general que el
homomorfismo~$\pi_1(\overline X_y,\overline a)\to \pi_1(X,a)$ sea inyectivo;
en topología algebraica su núcleo es la imagen de~$\pi_2(Y,b)$, y también en
geometría algebraica habría que introducir grupos de homotopía en todas las
dimensiones y la sucesión exacta de homotopía completa para un morfismo propio
que satisfaga hipótesis convenientes (por ejemplo, ser un morfismo
\ifthenelse{\boolean{orig}}
{liso.}
{liso).}
En la actualidad no se dispone de ningún resultado en este sentido, salvo una
definición razonable (si no definitiva) de los grupos de homotopía superiores.

\begin{corollaire}
\label{X.1.7}
Sean~$k$ un cuerpo algebraicamente cerrado y~$X$ e~$Y$ dos preesquemas conexos
sobre~$k$; supongamos~$X$ propio sobre~$k$ e~$Y$ localmente
\marginpar{265}
noetheriano. Sean~$a$ un punto geométrico de~$X$ y~$b$ un punto geométrico
de~$Y$, con valores en una misma extensión algebraicamente cerrada~$K$ de~$k$;
consideremos el punto geométrico~$c=(a,b)$ de~$X\times_k Y$ y el homomorfismo
~$\pi_1(X\times_k Y,c)\to \pi_1(X,a)\times\pi_1(Y,b)$ deducido de los
homomorfismos de grupos fundamentales asociados a las dos proyecciones
~$X\times_k Y\to X$ y~$X\times_k Y\to Y$. El homomorfismo precedente es un
\emph{isomorfismo}.
\end{corollaire}

Supongamos primero~$K=k$. Pongamos~$Z=X\times_k Y$, consideremos la proyección
~$f\colon Z\to Y$ y el punto subyacente~$y$ del punto geométrico~$b$ de~$Y$,
y apliquemos a esta situación el resultado~\Ref{X.1.4}. Obsérvese para ello
que, tras pasar a~$X_\red$ (lo que no cambia los grupos fundamentales
considerados), puede suponerse~$X$ reducido y, por tanto, separable sobre~$k$;
así,
\ifthenelse{\boolean{orig}}
{$Z$ es separable sobre~$k$}
{$Z$ es separable sobre~$Y$},
y tiene evidentemente fibras geométricamente conexas (pues~$X$ es conexo). La
fibra geométrica de~$Z$ en~$b$ es canónicamente isomorfa a
~$X\otimes_k K=X$. Por otra parte, como el compuesto de los morfismos
~$X\to Z\to X$ es la identidad, se ve que
~$\pi_1(X,a)\to \pi_1(Z,c)$ es inyectivo, y~\Ref{X.1.4} nos da una sucesión
exacta:
$$
e\to \pi_1(X,a)\to \pi_1(Z,c)\to \pi_1(Y,b)\to e.
$$
Por otra parte, se tiene la sucesión exacta canónica:
$$
e\to \pi_1(X,a)\to \pi_1(X,a)\times \pi_1(Y,b)\to \pi_1(Y,b)\to e,
$$
donde los dos homomorfismos escritos son la inyección canónica y la proyección
canónica. Por último, se tiene un homomorfismo de la primera sucesión exacta en
la segunda mediante los morfismos identidad en los términos extremos y el
homomorfismo canónico~$\pi_1(Z,c)\to \pi_1(X,a)\times\pi_1(Y,b)$ en los
términos medios. La conmutatividad del diagrama así obtenido se comprueba de
inmediato. Como los homomorfismos en los términos extremos son isomorfismos,
ocurre lo mismo con los términos medios, lo que prueba~\Ref{X.1.7} en este
caso.

Cuando ya no se supone~$K=k$, solo se obtiene un isomorfismo
$$
\pi_1(Z,c)\to \pi_1(X\otimes_k K,a)\times \pi_1(Y,b),
$$
\marginpar{266}%
y~\Ref{X.1.7} equivale entonces al caso particular siguiente:

\begin{corollaire}
\label{X.1.8}
Sean~$X$ un esquema propio y conexo sobre un cuerpo algebraicamente cerrado
~$k$,~$k'$ una extensión algebraicamente cerrada de~$k$,~$a'$ un punto
geométrico de~$X\otimes_k k'$ y~$a$ su imagen en~$X$. Entonces el
homomorfismo canónico~$\pi_1(X\otimes_k k',a')\to \pi_1(X,a)$ es un
\emph{isomorfismo}.
\end{corollaire}

Que este homomorfismo sea sobreyectivo equivale a decir que, si~$X'$ es un
recubrimiento étale conexo de~$X$, entonces~$X'\otimes_k k'$ es también
conexo, y resulta inmediatamente del hecho de que~$k$ es algebraicamente
cerrado; es también un caso particular
\ifthenelse{\boolean{orig}}
{\Ref{IX.3.4}.}
{de~\Ref{IX.3.4}.}
La hipótesis de que~$X$ sea propio no ha intervenido todavía. Dicho esto, la
inyectividad del homomorfismo considerado significa también lo siguiente:
\emph{todo recubrimiento étale de~$X\otimes_k k'$ es isomorfo a la imagen
inversa de un recubrimiento étale de~$X$.} Es esencialmente un sorites comprobar
que puede encontrarse una subálgebra sobre~$k$,~$A$, de
\ifthenelse{\boolean{orig}}
{$K$}
{$k'$},
de tipo finito sobre~$k$, y un recubrimiento étale de~$X\otimes_k A$ cuya
imagen inversa sobre~$X\otimes_k k'$ sea isomorfa al recubrimiento dado. Sea,
pues,~$Y=\Spec(A)$, que es un~$k$-esquema íntegro de tipo finito y tiene, por
tanto, puntos racionales sobre~$k$. Apliquemos entonces~\Ref{X.1.7} al grupo
fundamental de~$X\times Y$ en un punto~$(a,b)$ racional sobre~$k$: se obtiene
que todo recubrimiento étale conexo de~$X\times Y$ es isomorfo a un cociente
de un recubrimiento~$X'\times Y'$, donde~$X'$ e~$Y'$ son recubrimientos
galoisianos étale de~$X$ y~$Y$, de grupos~$G$ y~$G'$, por un subgrupo~$H$
de~$G\times G'$. Esto implica que la imagen inversa de este recubrimiento de
~$X\times Y$ sobre~$X\times Y'$ es isomorfa a un recubrimiento de la forma
~$X_1'\times Y'$, donde~$X_1'$ es un recubrimiento étale de~$X$. Si~$L$ es
el cuerpo de funciones de~$Y$, igual al cuerpo de fracciones de~$A$ en~$k'$,
el recubrimiento étale de~$X\otimes_k L$ inducido por el recubrimiento dado de
~$X\times_k Y$ es tal que existe una extensión finita separable~$L'$ de~$L$
para la cual la imagen inversa de dicho recubrimiento sobre~$X\otimes_k L'$
es isomorfa a~$X_1'\otimes_k L'$. Ahora bien, como~$k'$ es algebraicamente
cerrado, puede suponerse que la extensión~$L'$ de~$L$ está contenida en~$k'$.
Esto prueba que el recubrimiento étale dado de~$X\otimes_k k'$ es isomorfo a
~$X_1'\otimes_k k'$, c.q.d.

La forma explícita señalada de paso para los recubrimientos étale de un
producto
\marginpar{267}
~$X\times_k Y$ implica de inmediato el resultado siguiente:

\begin{corollaire}
\label{X.1.9}
Sean~$k$ un cuerpo algebraicamente cerrado,~$X$ e~$Y$ dos preesquemas
localmente noetherianos sobre~$k$,~$Z=X\times_k Y$ su producto y~$Z'$ un
recubrimiento étale de~$Z$. Para todo punto~$y\in Y$ racional sobre~$k$, sea
~$i_y \colon \Spec(k)\to Y$ el morfismo canónicamente asociado,
~$j_y=\id_X\times_k i_y$ el morfismo correspondiente~$X\to Z$. Sea, por
último,~$X_y'$ el recubrimiento étale de~$X$ imagen inversa de~$Z'$ por~$j_y$.
Supongamos~$Y$ conexo y~$X$ o~$Y$ propio sobre~$k$. Entonces los
recubrimientos~$X_y'$ de~$X$ son todos isomorfos.
\end{corollaire}

De manera figurada, puede decirse que \emph{una familia de recubrimientos
étale de~$X$, parametrizada por un preesquema conexo~$Y$, es constante si~$X$
o el preesquema de parámetros~$Y$ es propio sobre~$k$.}

\begin{remarques}
\label{X.1.10}
Los corolarios~\Ref{X.1.7} a~\Ref{X.1.9} se deben a
Lang--Serre~\cite{X.2} en el caso de los esquemas algebraicos normales. (Su
trabajo fue la motivación inicial de la teoría del grupo fundamental
desarrollada en este Seminario). Como observaron estos autores, estos
resultados dejan de ser ciertos cuando se abandona la hipótesis de que sean
propios, al menos en característica~$p>0$. Tomando, por ejemplo, como~$X$ la
recta afín~$X=\Spec(k[t])$, no es difícil ver que los recubrimientos de~$X$,
parametrizados por la recta afín~$Y=\Spec(k[s])$, definidos por las ecuaciones
$$
x^p - x = st,
$$
son étale y no isomorfos dos a dos. Esto contradice~\Ref{X.1.9} y, con mayor
razón,~\Ref{X.1.7}; se ve igualmente que, si~$s$ se considera como un elemento
trascendente sobre~$k$ en una extensión algebraicamente cerrada~$K$ de~$k$,
se obtiene un recubrimiento étale~$X'$ de~$X\otimes_kK$ que no procede de un
recubrimiento étale de~$X$.
\end{remarques}

\vfill
\section[Aplicación del teorema de existencia de haces]
{Aplicación del teorema de existencia de haces: teorema de semicontinuidad
para los grupos fundamentales de las fibras de un morfismo propio y separable}
\label{X.2}
\marginpar{268}

\begin{theoreme}
\label{X.2.1}
\ifthenelse{\boolean{orig}}
{Sean~$Y$ el espectro de un anillo local noetheriano \emph{completo}, de
cuerpo residual~$\kres$,~$X$ un~$Y$-esquema propio,
$X_0=X\otimes_A\kres$,~$a_0$ un punto geométrico de~$X_0$ y~$a$ el punto
geométrico correspondiente de~$X$.}
{Sean~$A$ un anillo local noetheriano \emph{completo} de cuerpo residual~$k$,
$Y=\Spec(A)$,~$X$ un~$Y$-esquema propio, $X_0=X\otimes_Ak$,~$a_0$ un punto
geométrico de~$X_0$ y~$a$ el punto geométrico correspondiente de~$X$.}
Entonces el homomorfismo canónico~$\pi_1(X_0,a_0)\to\pi_1(X,a)$ es un
\emph{isomorfismo}.
\end{theoreme}

No es sino una traducción, al lenguaje del grupo fundamental, del resultado
recordado en~\Ref{IX.1.10}. Es aquí donde el teorema de existencia
\ifthenelse{\boolean{orig}}
{de}
{de}
haces en geometría algebraico-formal entra de manera esencial en la teoría del
grupo fundamental.

Introduzcamos ahora una clausura algebraica~$\bar{k}$ del cuerpo residual~$k$ y
la fibra geométrica~$\overline X_0=X_0\otimes_k\bar{k}$.
\ifthenelse{\boolean{orig}}
{}
{Elijamos un punto geométrico~$\overline a_0$ de~$\overline X_0$ sobre~$a_0$.}
Se tiene, pues, la sucesión exacta (\Ref{IX.6.1})
$$
\ifthenelse{\boolean{orig}}
{e\to\pi_1(\overline X_0,\overline a_0)\to\pi_1(X_0,a_0)\to
\pi_1(k,\bar{k})\to e,}
{e\to\pi_1(\overline X_0,\overline a_0)\to\pi_1(X_0,a_0)\to
\pi_1(k,\bar{k})\to e,}
$$
por otra parte, se tienen el isomorfismo~\Ref{X.2.1} y el isomorfismo análogo,
más elemental,
$$
\pi_1(k,\bar{k})\to\pi_1(Y,b),
$$
donde~$b$ es la imagen de~$a$ en~$Y$. Se obtiene así:

\begin{corollaire}
\label{X.2.2}
\ifthenelse{\boolean{orig}}
{Con las notaciones anteriores, supongamos conexo~$\overline X_0$ y sean
$\overline a_0$ un punto geométrico de
$\overline X_0=X_0\otimes_k\bar{k}$,~$a_0$ su imagen en~$X$,}
{Con las notaciones anteriores, supongamos conexo
$\overline X_0=X_0\otimes_k\bar{k}$ y sean~$\overline a_0$ un punto geométrico
de~$\overline X_0$,~$a_0$ su imagen en~$X$,}
~$b_0$ su imagen en~$Y$. Entonces la sucesión siguiente de homomorfismos
canónicos
$$
\ifthenelse{\boolean{orig}}
{e\to\pi_1(\overline X_0,\overline a_0)\to\pi_1(X,a_0)\to
\pi_1(Y,b_0)\to e}
{e\to\pi_1(\overline X_0,\overline a_0)\to\pi_1(X,a_0)\to
\pi_1(Y,b_0)\to e}
$$
es exacta.
\end{corollaire}

Compárese esta sucesión con la sucesión exacta~\Ref{X.1.4}, pero obsérvese que
a)~no ha sido necesario imponer ninguna hipótesis de planitud ni de
separabilidad sobre las fibras para~$X\to Y$; y b)~se obtiene el complemento
importante de que \emph{el morfismo
$\pi_1(\overline X_0,\overline a_0)\to\pi_1(X,a_0)$ es inyectivo}.

Este último hecho permitirá comparar el grupo fundamental de las demás
\marginpar{269}
fibras geométricas de~$X$ sobre~$Y$ con el de~$\overline X_0$. Sea, en efecto,
~$y_1$ un punto cualquiera de~$Y$,~$X_1$ su fibra y~$\overline X_1$ su fibra
geométrica respecto de una extensión algebraicamente cerrada de~$\kres(y_1)$.
Sea~$\overline a_1$ un punto geométrico de~$\overline X_1$,~$a_1$ su imagen
en~$X$ y~$b_1$ su imagen en~$Y$. Elijamos una «clase de caminos» de~$a_1$
a~$a_0$, de donde se obtiene una clase de caminos de~$b_1$ a~$b_0$ y, por
tanto, un diagrama conmutativo de homomorfismos:
$$
\xymatrix@=.5cm{
&\pi_1(\overline X_1,\overline a_1) \ar[r] & \pi_1(X,a_1) \ar[r]\ar[d] & \pi_1(Y,b_1) \ar[r]\ar[d] & e \\
e \ar[r]& \pi_1(\overline X_0,\overline a_0) \ar[r] & \pi_1(X,a_0) \ar[r] &\pi_1(Y,b_0) \ar[r] & e,
}
$$
en el cual las dos flechas verticales indicadas son isomorfismos. Como la
segunda fila es exacta, se obtiene un homomorfismo canónico, que llamaremos
\emph{homomorfismo de especialización para el grupo fundamental} (que solo
depende de la clase de caminos elegida de~$a_1$ a~$a_0$ y está, por tanto,
\emph{definido módulo un automorfismo interior de}~$\pi_1(X,a_0)$):
$$
\pi_1(\overline X_1,\overline a_1)\to
\pi_1(\overline X_0,\overline a_0).
$$
Cuando la primera fila anterior es también exacta, se sigue inmediatamente que
el homomorfismo de especialización es sobreyectivo. Teniendo en cuenta
~\Ref{X.1.4}, se obtiene, pues:

\begin{corollaire}
\label{X.2.3}
Bajo las hipótesis de~\Ref{X.2.1}, supongamos además que el morfismo
~$f\colon X\to Y$ es \emph{separable} \eqref{X.1.1} y que~$\overline X_0$ es
conexo (por tanto, en virtud de~\Ref{X.1.2}, se tiene
$f_*(\cal{O}_X)=\cal{O}_Y$). Entonces, para toda fibra geométrica
~$\overline X_1$ de~$X$ sobre~$Y$, provista de un punto geométrico
~$\overline a_1$, el homomorfismo de especialización
\index{homomorfismo de especialización para el grupo fundamental|hyperpage}%
\index{especialización para el grupo fundamental (homomorfismo de)|hyperpage}%
definido anteriormente es \emph{sobreyectivo}.
\end{corollaire}

Este es un resultado de \emph{semicontinuidad} para el grupo fundamental, que
no parece tener todavía un análogo en topología algebraica. Puede enunciarse,
por lo demás, bajo hipótesis más generales:

\begin{corollaire}
\label{X.2.4}
Sean
\marginpar{270}
~$f\colon X\to Y$ un morfismo propio de fibras geométricamente conexas,
con~$Y$ localmente noetheriano;~$y_0$ e~$y_1$ dos puntos de~$Y$ tales que
$y_0\in\overline{\{y_1\}}$;~$\overline X_0$ y~$\overline X_1$ las fibras
geométricas de~$X$ correspondientes a extensiones algebraicamente cerradas
elegidas de~$\kres(y_0)$ y~$\kres(y_1)$; y~$\overline a_0$,
\resp~$\overline a_1$, un punto geométrico de~$\overline X_0$,
\resp~$\overline X_1$. Entonces puede definirse de manera natural un
homomorfismo de especialización
$$
\pi_1(\overline X_1,\overline a_1)\to
\pi_1(\overline X_0,\overline a_0),
$$
definido salvo automorfismo interior, y este homomorfismo es \emph{sobreyectivo}
si~$f$ es un morfismo separable \eqref{X.1.1}.
\end{corollaire}

En efecto, de~\Ref{X.1.8} se sigue primero que~\Ref{X.2.4} es esencialmente
independiente de las extensiones algebraicamente cerradas elegidas para los
cuerpos residuales~$\kres(y_0)$ y~$\kres(y_1)$. Esto permite sustituir~$Y$ por
un preesquema~$Y'$ sobre~$Y$ que tenga un punto~$y_0'$ (\resp~$y_1'$) sobre
~$y_0$ (\resp~$y_1$). Se toma entonces por~$Y'$ el espectro del completado del
anillo local de~$y_0$ en~$Y$, y se aplica~\Ref{X.2.3}.

\begin{remarques}
\label{X.2.5}
La conclusión final de~\Ref{X.2.4} sobre la sobreyectividad del homomorfismo de
especialización, y a fortiori los resultados~\Ref{X.1.3} y~\Ref{X.1.4} de los
que se deduce, deja de ser cierta si ya no se supone que~$f\colon X\to Y$ es
un morfismo separable, incluso para esquemas proyectivos sobre un cuerpo
algebraicamente cerrado de característica~$0$. Más adelante veremos ejemplos
tanto cuando~$f$ es plano pero~$f$ admite una fibra no separable (aunque~$X$
e~$Y$ sean lisos sobre~$k$), como cuando las fibras de~$f$ son separables pero~$f$ no
es plano (por ejemplo, cuando~$f\colon X\to Y$ es un morfismo birracional de
esquemas íntegros normales); véase~\Ref{XI.3}. En estos ejemplos puede
ocurrir que el grupo fundamental de la fibra geométrica genérica sea trivial,
pero no el de una fibra geométrica especial conveniente.

Por otra parte, incluso si~$f\colon X\to Y$ es un morfismo propio y separable
como en~\Ref{X.2.4}, es frecuente que el homomorfismo de especialización no sea
\marginpar{271}
un isomorfismo. Por ejemplo, es fácil dar casos en los que~$\overline X_1$ es
una curva elíptica no singular (por tanto, su grupo fundamental es conmutativo
y su componente~$\ell$-primaria, para un número primo~$\ell$ distinto de la
característica, es isomorfa a~$\ZZ_\ell^2$; véase~\Ref{XI}), mientras que
~$\overline X_0$ está formada ya por dos curvas racionales no singulares que se
cortan en dos puntos, ya por dos curvas racionales tangentes en un punto, ya,
por último, por una curva racional con un punto singular que es una cúspide.
(Para la clasificación completa de las curvas elípticas degeneradas, véanse
los trabajos recientes de Kodaira~\cite{X.1} y Néron). En estos casos, el grupo
fundamental de~$\overline X_0$ es, respectivamente,~$\hat\ZZ$,~$e$,~$e$, y es,
por tanto, «estrictamente menor» que el de~$\overline X_1$.

Veremos, sin embargo, más adelante, cuando~$f$ es un morfismo \emph{liso}, una
cota para el núcleo del homomorfismo de especialización que implica, en
particular, que si~$\kres(y_0)$ tiene característica~0, el homomorfismo de
especialización es un isomorfismo. Pero incluso para un morfismo liso, si la
característica de~$\kres(y_0)$ es~$>0$, el homomorfismo de especialización
puede no ser un isomorfismo, como se ve, por ejemplo, cuando~$X$ es un esquema
abeliano sobre~$Y$ (incluso de dimensión relativa~$1$),
\ifthenelse{\boolean{orig}}
{\cf~\Ref{XI.2}).}
{\cf~\Ref{XI.2}.}
Una teoría satisfactoria de la especialización del grupo fundamental debe
tener en cuenta la «componente continua» del «verdadero» grupo fundamental,
correspondiente a la clasificación de los recubrimientos principales cuyos
grupos estructurales son grupos infinitesimales. Con ello cabría esperar que
los «verdaderos» grupos fundamentales de las fibras geométricas de un morfismo
liso y propio~$f\colon X\to Y$ formasen un sistema local bien comportado
sobre~$X$, límite proyectivo de esquemas en grupos finitos y planos sobre~$X$
\kern1pt\footnote{Esta conjetura sumamente seductora queda, por desgracia,
refutada por un ejemplo inédito de M\ptbl Artin, ya cuando las fibras de~$f$
son curvas algebraicas de un género dado~$g\ge 2$.}. Volveremos más adelante
sobre este punto de vista; nuestro propósito presente es, por el contrario,
llevar tan lejos como sea posible los fenómenos comunes a la teoría topológica
y a la teoría esquemática del grupo fundamental.
\end{remarques}

Sea ahora~$X_0$ una curva propia, lisa y conexa de género~$g$ sobre un cuerpo
algebraicamente cerrado~$k$. Si~$k$ tiene característica cero, su grupo
fundamental puede determinarse por métodos trascendentes como sigue. Se sabe
que~$X_0$ procede por extensión de la base de una curva definida
\marginpar{272}
sobre una extensión algebraicamente cerrada de grado de trascendencia finito
del cuerpo primo~$\QQ$ y, teniendo en cuenta~\Ref{X.1.8}, puede suponerse que
~$k$ mismo tiene grado de trascendencia finito sobre~$\QQ$. Puede suponerse,
por tanto, que~$k$ es un subcuerpo del cuerpo~$\CC$ de los números complejos,
y una nueva aplicación de~\Ref{X.1.8} permite suponer que~$k=\CC$. No es
difícil comprobar entonces que el grupo fundamental de
\ifthenelse{\boolean{orig}}
{$X$}
{$X_0$}
es isomorfo a la completación profinita del grupo fundamental del espacio
topológico asociado
\ifthenelse{\boolean{orig}}
{$\tilde X$}
{$\widetilde X_0$}
(una superficie compacta orientada de género~$g$), para la topología definida
por los subgrupos de índice finito\footnote{Esta deducción se explicó en una de
las exposiciones orales que no fueron redactadas.}. Por otra parte, es clásico
que el grupo fundamental topológico está engendrado por~$2g$ generadores
~$s_i,t_i$ ($1\le i\le g$), sujetos a la única relación
$$
(s_1t_1s_1^{-1}t_1^{-1})\cdots(s_gt_gs_g^{-1}t_g^{-1})=1.
$$
Por tanto, el grupo fundamental de
\ifthenelse{\boolean{orig}}
{$X$}
{$X_0$}
admite~$2g$ generadores \emph{topológicos}~$s_i,t_i$ ($1\le i\le g$), ligados
por la única relación anterior.

Supongamos ahora que la característica de~$k$ es~$p>0$. Sea~$A$ el anillo de
vectores de Witt construido a partir de~$k$, y~$K$ una extensión
algebraicamente cerrada de su cuerpo de fracciones. Se vio en
~\Ref{III.7.4} que existe un esquema~$X$ sobre~$Y=\Spec(A)$, propio y liso
sobre~$Y$, cuya reducción es~$X_0$. Aplicando~\Ref{X.2.3}, se obtiene un
homomorfismo \emph{sobreyectivo}
$$
\pi_1(X_1)\to\pi_1(X_0),
$$
donde~$X_1=X\otimes_AK$. Es inmediato\footnote{\Cf EGA~IV 12.2.} que~$X_1$
es liso sobre~$K$, conexo por~\eqref{X.1.2}, de dimensión~$1$, y que su género
es igual a~$g$ por la invariancia de la característica de Euler--Poincaré;
véase EGA~III 7. Como~$K$ tiene característica~$0$, puede aplicársele el
resultado anterior. Se ha probado así, por \emph{métodos trascendentes}:

\begin{theoreme}
\label{X.2.6}
Sea~$X_0$ una curva algebraica lisa, propia y conexa sobre un
\marginpar{273}
cuerpo algebraicamente cerrado~$k$, y sea~$g$ su género. Entonces
~$\pi_1(X_0)$ admite un sistema de~$2g$ generadores topológicos, ligados por la
relación escrita anteriormente. Si~$k$ tiene característica~$0$,
~$\pi_1(X_0)$ es además el grupo libre de tipo galoisiano sobre esos generadores
sujeto a dicha relación.
\end{theoreme}

\begin{remarques}
\label{X.2.7}
En la actualidad, que sepa el redactor, no existe una demostración por métodos
puramente algebraicos del resultado anterior, salvo para los géneros~$0,1$.
Para empezar, no se ve en absoluto cómo distinguir en~$\pi_1(X_0)$
~$2g$ elementos de los que después pudiera esperarse que formasen un sistema de
generadores topológicos. A este respecto, es más favorable la situación de la
recta racional privada de~$n$ puntos y el estudio de sus recubrimientos
moderadamente ramificados en esos puntos, pues los grupos de ramificación en
los~$n$ puntos proporcionan~$n$ elementos del grupo fundamental estudiado y se
prueba que lo engendran topológicamente, como veremos más adelante
\footnote{\Cf Exposición~\Ref{XII}. Incluso estos elementos solo quedan
determinados módulo conjugación, y conviene elegir \emph{simultáneamente} y con
cuidado representantes de sus clases.}. Pero ni siquiera en este caso
particularmente concreto parece conocerse una demostración puramente
algebraica. Una demostración tal sería, desde luego, sumamente interesante. El
único hecho relativo al grupo fundamental de una curva que parece admitir una
demostración puramente algebraica, aparte del teorema débil de finitud
~\Ref{X.2.12} que Lang--Serre~\cite{X.2} probaron algebraicamente, es la
determinación de su abelianización mediante la jacobiana, señalada en la última
línea de~\Ref{IX.5.8}.
\end{remarques}

\setcounter{subsection}{7}
\subsection{}
\label{X.2.8}
La última afirmación
\ifthenelse{\boolean{orig}}
{\Ref{X.2.6}}
{de~\Ref{X.2.6}}
deja de ser válida en característica~$p>0$, como se ve ya en el caso de las
curvas elípticas. Como se ha señalado anteriormente, no sabemos si el grupo
fundamental de~$X_0$ es topológicamente de presentación finita; parece por
completo improbable.

\setcounter{proposition}{8}
\begin{theoreme}
\label{X.2.9}
Sean~$k$ un cuerpo algebraicamente cerrado y~$X$ un esquema propio y conexo
sobre~$k$. Entonces el grupo fundamental de~$X$ es topológicamente de
generación finita.
\end{theoreme}

Procedemos
\marginpar{274}
por inducción sobre~$n=\dim X$; la afirmación es trivial para~$n\le 0$.
Supongamos, pues, que~$n>0$ y que el teorema está demostrado en dimensiones
~$n'<n$. Por el lema de Chow (EGA~II 5.6.2), existe un esquema proyectivo~$X'$
sobre~$k$ y un morfismo sobreyectivo~$X'\to X$. Puede suponerse claramente
que~$X'$ es reducido y, pasando a su normalización, normal. Gracias a la teoría
del descenso, basta probar que los grupos fundamentales de las componentes
conexas de~$X'$ son topológicamente de generación finita
(\Ref{IX.5.2}). Quedamos reducidos así al caso en que~$X$ es \emph{proyectivo}
y \emph{normal}. Si~$n=1$, basta aplicar~\Ref{X.2.6}. Si~$n\ge 2$, considérese
una inmersión proyectiva~$X\to\PP_k^r$ y una sección de hiperplano~$Y=X\cdot H$
\ifthenelse{\boolean{orig}}
{(dotado}
{(dotada}
de la estructura reducida inducida), tal que
\ifthenelse{\boolean{orig}}
{$Y\ne X$}
{$Y\ne X$,}
\ie $H\not\supset X$. Se tiene entonces~$\dim Y<n$ y, teniendo en cuenta la
hipótesis de inducción, basta probar que~$\pi_1(Y)\to\pi_1(X)$ es
\emph{sobreyectivo}. Más generalmente:

\begin{lemme}
\label{X.2.10}
Sean~$X$ un preesquema propio sobre un cuerpo algebraicamente cerrado~$k$ y
~$g\colon X\to\PP_k^r$ un morfismo. Supongamos que~$X$ es irreducible y normal
y que~$\dim g(X)\ge 2$. Sean~$H$ un hiperplano de~$\PP_k^r$ y
~$Y=X\times_{\PP_k^r}H$. Entonces~$Y$ es conexo y el homomorfismo
~$\pi_1(Y)\to\pi_1(X)$ es sobreyectivo.
\end{lemme}

Estas afirmaciones se deducen, en efecto, de la siguiente:

\begin{corollaire}
\label{X.2.11}
Bajo las hipótesis anteriores, sean~$X'$ un recubrimiento étale conexo de~$X$ y
~$Y'=X'\times_XY=X'\times_{\PP_k^r}H$ el recubrimiento inducido sobre~$Y$.
Entonces~$Y'$ es conexo.
\end{corollaire}

Como~$X$ es normal,~$X'$ es normal y, por ser conexo,~$X'$ es irreducible;
además, su imagen en~$\PP_k^r$ tiene dimensión~$\ge 2$. Un conocido lema de
Zariski (llamado «\emph{teorema de Bertini}») implica, por tanto, que si
~$H_1'$ es el hiperplano genérico de~$\PP_k^r$, definido sobre una extensión
~$K$ de~$k$, entonces
\ifthenelse{\boolean{orig}}
{$X'\times_{\PP^r}H_1$}
{$X'\times_{\PP^r}H_1'$}
es universalmente irreducible y, por tanto, universalmente conexo sobre~$K$.
El teorema de conexidad de Zariski (EGA~III 4) implica entonces que para
\emph{todo} hiperplano~$H$ (definido sobre una extensión cualquiera de~$k$),
~$X'\times_{\PP^r}H$ es geométricamente conexo. Esto termina la demostración
de~\Ref{X.2.11} y, por consiguiente, la de~\Ref{X.2.9}.

\begin{corollaire}[Lang--Serre]
\label{X.2.12}
Bajo
\marginpar{275}
las hipótesis de~\Ref{X.2.9}, para todo grupo finito~$G$, el conjunto de clases,
salvo isomorfismo, de recubrimientos principales de~$X$ con grupo~$G$ es finito.
\end{corollaire}

\begin{remarque}
\label{X.2.13}
Bajo las hipótesis de~\Ref{X.2.10}, cuando~$\dim g(X)\ge 3$ (al menos cuando
~$g$ es una inmersión y~$X$ es regular), probaremos un resultado más preciso,
conocido en geometría algebraica con el nombre de «\emph{teorema de Lefschetz}»:
~$\pi_1(Y)\to\pi_1(X)$ \emph{es un isomorfismo}.\footnote{\Cf el seminario
SGA~2 (1962), que siguió a este%
\ifthenelse{\boolean{orig}}
{.}
{, teorema X~3.10.}}
En los casos clásicos hay enunciados análogos para los grupos de homología y
los grupos de homotopía superiores, que tarde o temprano habrán de quedar
abarcados por la geometría algebraica abstracta. Ni siquiera para la
cohomología de Hodge~$\H^p(X,\Omega^q)$ parece haberse estudiado aún la
cuestión; por lo demás, es poco probable que para esta última los teoremas de
Lefschetz subsistan sin cambio en característica~$p>0$.
\end{remarque}

\iforig
\else
\begin{remarqueMR}
\label{X.2.14}
Sea~$R$ un anillo de valoración discreta completo, de cuerpo residual
algebraicamente cerrado~$\kres$ de característica~$p>0$ y cuerpo de
fracciones~$K$, y sea~$Y$ una curva propia, lisa y conexa de género~$g$
sobre~$R$. Se tiene, pues, un morfismo de especialización sobreyectivo
~$\sp : \pi_1(Y_{\overline{K}})\to\pi_1(Y_{\kres})$.
Ya se ha observado que, si~$K$ tiene característica~0,~$\sp$ no es un
isomorfismo cuando~$g\geq1$. Supongamos que~$K$ tiene característica~$p$, de
modo que~$R$ es isomorfo al anillo de series formales~$\kres[[T]]$.

En igual característica~$p>0$, el grupo fundamental no queda determinado por
el género~$g$, como se ve ya en las curvas elípticas, que pueden ser ordinarias
o supersingulares. Citemos el resultado reciente de A\ptbl Tamagawa (todavía
inédito). Si~$G$ es un grupo profinito, se denota por~$G^{\res}$ el grupo
profinito cociente de~$G$ que es límite proyectivo de los cocientes topológicos
finitos resolubles de~$G$.

\begin{theoremstar}[A\ptbl Tamagawa]
Supongamos que~$g\geq2$, que la fibra especial~$Y_{\kres}$ es definible sobre
un cuerpo finito y que el morfismo
~$\sp^{\res}:\pi_1(Y_{\overline{K}})^{\res}\to
\pi_1(Y_{\kres})^{\res}$ inducido por~$\sp$ al pasar al cociente es biyectivo.
Entonces la curva~$Y$ es constante sobre~$R$.
\end{theoremstar}

Obsérvese que el grupo de Galois de~$\overline{K}/K$ es resoluble. El enunciado
anterior también puede expresarse así: supongamos que todo recubrimiento finito
étale~$X_K\to Y_K$ de la fibra genérica~$Y_K$, galoisiano con grupo de Galois
resoluble, tiene reducción potencialmente buena (es decir, se extiende a un
recubrimiento finito étale de~$Y$ después de sustituir eventualmente~$R$ por
su normalización en una extensión finita de~$K$). Entonces~$Y$ es constante
sobre~$R$.
\end{remarqueMR}
\fi

\section[Aplicación del teorema de pureza]{Aplicación del teorema de pureza:
teorema de continuidad para los grupos fundamentales de las fibras de un
morfismo propio y liso}
\label{X.3}

Recordemos sin demostración\footnote{Para una demostración, \cf SGA~2
\Ref{X.3.4}.} el

\begin{theoremedepurete}[Zariski--Nagata]
\label{X.3.1}
\index{teorema de pureza|hyperpage}%
Sea~$f\colon X\to Y$ un morfismo cuasifinito y dominante de preesquemas
íntegros, con~$X$ normal e~$Y$ regular y localmente noetheriano, y sea~$Z$ el
conjunto de los puntos de~$X$ en los que~$f$ no es étale, \ie en los que~$f$
está ramificado (ambas condiciones son equivalentes,
\textup{\Ref{I.9.5}~(ii)}). Si~$Z\ne X$, entonces~$Z$ tiene codimensión~$1$
en~$X$ en todos sus puntos, \ie para toda componente irreducible~$Z'$ de~$Z$, de
punto genérico~$z$, la dimensión de Krull de~$\cal{O}_{X,z}$ es igual a~$1$.
\end{theoremedepurete}

Recordemos que un preesquema se dice \emph{normal}, respectivamente
\emph{regular}, si sus anillos locales son normales, respectivamente regulares,
y que la relación~$Z\ne X$ significa también que la extensión finita
\ifthenelse{\boolean{orig}}
{$R(Z)/R(X)$}
{$R(X)/R(Y)$}
(donde~$R$ denota el cuerpo de funciones racionales) es \emph{separable}.
Situándonos en el punto genérico~$z$ de una componente~$Z'$ de~$Z$ y
localizando en el punto~$y$ de~$Y$ situado debajo de~$z$, se obtiene
\marginpar{276}
el enunciado equivalente:

\begin{corollaire}
\label{X.3.2}
Sean~$A$ un anillo local noetheriano regular y~$A\to B$ un homomorfismo local
inyectivo tal que~$B$ sea normal, localizado de un álgebra de tipo finito
sobre~$A$, y \emph{cuasifinito} sobre~$A$; supongamos además que
$\dim A (=\dim B)\ge 2$ y que, para todo ideal primo~$\goth{p}$ de~$B$
distinto del ideal maximal,~$B$ sea étale sobre~$A$ en~$\goth{p}$, \ie que
$B_\goth{p}$ sea étale sobre~$A_\goth{q}$ (donde
$\goth{q}=A\cap\goth{p}$). Entonces~$B$ es étale sobre~$A$.
\end{corollaire}

Por lo demás, no es difícil reducir este último enunciado al caso en que~$A$
es un anillo local \emph{completo}, y por tanto~$B$ es \emph{finito} sobre~$A$.
Zariski~\cite{X.5} da una demostración sencilla de este resultado, válida en
el caso de igual característica; el caso general se debe a
Nagata~\cite{X.3}, quien se apoya en un resultado delicado de Chow. Este último
no fue verificado por ninguno de los participantes del Seminario y debería ser
objeto de una exposición ulterior. Señalemos aquí solamente la demostración muy
sencilla en el caso particular~$\dim A=2$, suficiente para la aplicación más
importante que haremos de él en este número. Como~$B$ es normal, es un
$B$-módulo de profundidad (terminología antigua: codimensión cohomológica)
$\ge 2$; por tanto, es un~$A$-módulo de profundidad~$\ge 2$. Como~$A$ es
regular de dimensión~$2$, se deduce que~$B$ es un \emph{módulo libre}
sobre~$A$\kern1pt\footnote{\Cf EGA $0_{\textup{IV}}$ 17.3.4}. Resulta entonces
de \Ref{I.4.10} que el conjunto de los ideales primos~$\goth{q}$ de~$A$ en
los que~$B$ está ramificado sobre~$A$ es la parte de~$\Spec(A)$ definida por un
ideal principal (engendrado por el discriminante de una base de~$B$ sobre~$A$),
y por consiguiente es vacía si está contenida en el punto cerrado de~$\Spec(A)$.
Esto demuestra~\Ref{X.3.2} cuando~$\dim A=2$.

Utilizaremos sobre todo~\Ref{X.3.1} bajo la forma equivalente siguiente:

\begin{corollaire}
\label{X.3.3}
Sean~$X$ un preesquema regular localmente noetheriano y~$U$ una parte abierta
de~$X$, complementaria de una parte cerrada~$Z$ de~$X$ de
codimensión~$\ge 2$. Entonces el funtor~$X'\mto X'\times_XU$ de la categoría
de los recubrimientos étale de~$X$ en la categoría de los recubrimientos
étale de~$U$ es una equivalencia
\marginpar{277}
de categorías. En particular, si~$a$ es un punto geométrico de~$U$, el
homomorfismo canónico~$\pi_1(U,a)\to\pi_1(X,a)$ es un isomorfismo.
\end{corollaire}

La última afirmación es evidentemente consecuencia de la primera; para esta
última, podemos suponer evidentemente que~$X$ es conexo y, por tanto,
irreducible. De la normalidad de~$X$ se deduce ya que el funtor
$X'\mto X'\times_XU$ de la categoría de los recubrimientos localmente libres
(no necesariamente étale)
\ifthenelse{\boolean{orig}}
{de~$X'$}
{de~$X$}
en la categoría de los recubrimientos de~$U$ es plenamente fiel, pues lo es el
funtor~$\cal{E}\mto\cal{E}\vert U$ de la categoría de los Módulos localmente
libres sobre~$X$ en la categoría de los Módulos localmente libres sobre~$U$.
Queda, pues, probar que para todo recubrimiento \emph{étale}~$U'$ de~$U$
existe un recubrimiento étale~$X'$ de~$X$ (necesariamente único por lo que
precede) tal que~$U'$ sea isomorfo a~$X'\times_XU$. Podemos suponer
evidentemente que~$U'$ es conexo y, por tanto, irreducible, pues al ser~$U$
normal,~$U'$ es normal. Sea~$K$ el cuerpo de funciones racionales sobre~$X$ o
sobre~$U$ (es el mismo), y sea~$K'$ el de~$U'$; entonces~$U'$ se identifica con
el normalizado de~$U$ en~$K'$, por \Ref{I.10.3}. Sea~$X'$ el normalizado
de~$X$ en~$K'$ (EGA~II~6.3); entonces
\ifthenelse{\boolean{orig}}
{$X'=X\times_XU\cong U'$}
{$X'\times_XU\cong U'$}.
Por otra parte,~$X'$ es normal e íntegro, y el morfismo estructural
$f\colon X'\to X$ es \emph{finito} y dominante (pues~$X$ es normal y~$K'/K$
es una extensión finita separable). Es étale en
\ifthenelse{\boolean{orig}}
{$U'=f^{-1}(U)=X'=f^{-1}(Z)$}
{$U'=f^{-1}(U)$, y $X'\setminus U'=f^{-1}(Z)$};
como~$Z$ tiene codimensión~$\ge 2$ en~$X$,~$f^{-1}(Z)$ tiene
codimensión~$\ge 2$ en~$X'$. Se concluye entonces de~\Ref{X.3.1} que~$X'$ es
étale sobre~$X$, lo cual termina la demostración.

\medskip
Sea ahora~$f\colon X\to Y$ una aplicación racional de un preesquema regular
localmente noetheriano~$X$ en un preesquema~$Y$, y supongamos que~$f$ está
definida en un abierto~$U$ complementario de una parte cerrada de
codimensión~$\ge 2$. Entonces~\Ref{X.3.3} proporciona un funtor, definido
salvo isomorfismo, de la categoría de los recubrimientos étale de~$Y$ en la
categoría de los recubrimientos étale de~$X$; de ahí, para todo punto
geométrico~$a$ de~$U$, de imagen~$b$ en~$Y$, un homomorfismo canónico
$$
\pi_1(X,a)\to \pi_1(Y,b)
$$
(deducido
\marginpar{278}
del homomorfismo canónico~$\pi_1(U,a)\to\pi_1(Y,b)$ mediante el isomorfismo
$\pi_1(U,a)\isomto \pi_1(X,a))$.
Cuando~$f$ es un morfismo dominante, siendo~$X$ e~$Y$ íntegros de cuerpos de
funciones~$K$ y~$L$, de modo que~$K$ es una extensión de~$L$, y siendo~$Y$
normal, estas correspondencias se precisan en términos de extensiones de
cuerpos observando que, para toda extensión finita~$L'$ de~$L$ no ramificada
sobre~$Y$, el álgebra~$K'=L'\otimes_LK$ sobre~$K$ no está ramificada
sobre~$X$.

En particular, estas observaciones muestran que el grupo fundamental de los
preesquemas regulares, conexos y localmente noetherianos, punteados por puntos
geométricos situados en codimensión~$\le 1$, es un \emph{funtor} cuando se
toman como morfismos en esta categoría las aplicaciones racionales dominantes
definidas en complementarios de partes cerradas de codimensión~$\ge 2$.
Recordando, por ejemplo, que una aplicación racional de un esquema normal sobre
un cuerpo~$k$ en un esquema propio sobre~$k$ está definida en el complementario
de un conjunto de codimensión~$\ge 2$, se obtiene:

\begin{corollaire}[Invariancia birracional del grupo fundamental]
\label{X.3.4}
\index{invariancia birracional del grupo fundamental|hyperpage}%
\index{teorema de invariancia birracional del grupo fundamental|hyperpage}%
Sean~$k$ un cuerpo,~$X$ e~$Y$ dos esquemas regulares y propios sobre~$k$,
$f\colon X\to Y$ una aplicación birracional de~$X$ en~$Y$, y~$\Omega$ una
extensión algebraicamente cerrada del cuerpo de funciones~$K$ de~$X$ que
permita definir el grupo fundamental de~$X$ y el grupo fundamental de~$Y$.
Estos grupos son entonces canónicamente isomorfos.
\end{corollaire}

Esto significa asimismo que, para una extensión finita~$K'$ de~$K$, si no está
ramificada sobre un «modelo» propio no singular~$X$ de~$K$, tampoco lo está
sobre ningún otro modelo propio no singular.

\begin{remarque}
\label{X.3.5}
Para otras aplicaciones del teorema de pureza, véanse los trabajos
\ifthenelse{\boolean{orig}}
{de \textsc{Abhyankar}}
{de \textsc{Abhyankar}}
expuestos en~\cite{X.4}, inspirados en los resultados de
\textsc{Zariski}~\cite[cap\ptbl VIII]{X.6}, demostrados por vía topológica.
Estos últimos distan mucho de haber sido asimilados por la geometría algebraica
«abstracta» y merecen nuevos esfuerzos.
\end{remarque}

Necesitaremos
\marginpar{279}
algunos hechos elementales de la teoría de la ramificación. Sean~$V$ un
anillo de valoración discreta de cuerpo de fracciones~$K$ y cuerpo
residual~$k$,~$L$ una extensión galoisiana de~$K$ de grupo~$G$, y~$V'$ el
normalizado de~$V$ en~$L$, que es un módulo libre de rango
$n=[L{\colon }K]$ sobre~$V$. Sean~$\goth{m'}$ un ideal maximal de~$V'$,
$G_d$ el subgrupo de~$G$ formado por los elementos que dejan invariante
$\goth{m'}$, de modo que~$G_d$ opera en la extensión residual
$k'=V'/\goth{m'}$ de~$k$, y~$G_i$ el subgrupo de los elementos de~$G_d$ que
operan trivialmente (recordemos que~$G_d$ y~$G_i$ se llaman, respectivamente,
subgrupos de \emph{descomposición} y de \emph{inercia} de~$G$). Se dice que~$L$
está \emph{moderadamente ramificada} sobre~$V$
\ifthenelse{\boolean{orig}}
{si~$n_i=[G_i{\colon }e]$ es de orden~$p$ primo a la característica de~$k$}
{si~$n_i=[G_i{\colon }e]$ es primo a la característica~$p$ de~$k$}
(condición siempre satisfecha si~$k$ tiene característica~0).
Es bien sabido que~$G_i$ se sumerge entonces canónicamente en el grupo
${k'}^*$ y, por tanto, es isomorfo al grupo de las raíces
\ifthenelse{\boolean{orig}}
{$n_i$-esimas}
{$n_i$-esimas}
de la unidad en~${k'}^*$. Esto implica, en particular, que~$G_i$ \emph{es
cíclico}. El caso típico de esta situación es aquel en que se toma
$L=K[t]/(t^n-u)$, siendo~$u$ un uniformizante de~$V$ y~$n$ un entero primo
a~$p$: si~$K$ contiene las raíces \niemes de la unidad,~$L$ es una extensión
galoisiana totalmente ramificada de~$K$, de grupo de Galois~$G=G_i$ isomorfo
a~$\ZZ/n\ZZ$.

\begin{lemme}[%
\ifthenelse{\boolean{orig}}
{Lema de \textsc{Abhyankar}}
{Lema de \textsc{Abhyankar}}]
\label{X.3.6}
\index{Abhyankar (lema de)|hyperpage}%
\index{lema de Abhyankar|hyperpage}%
Sea~$V$ un anillo de valoración discreta de cuerpo de fracciones~$K$, sean~$L$
y~$K'$ dos extensiones galoisianas de~$K$ \emph{moderadamente ramificadas}
sobre~$V$, y sean~$n$ y~$m$ los órdenes de los grupos de inercia
correspondientes
\ifthenelse{\boolean{orig}}
{.}
{,}
Sea~$L'$ una extensión compuesta de~$L$ y~$K'$ sobre~$K$. Si~$m$ es múltiplo
de~$n$, entonces~$L'$ no está ramificada sobre los localizados de la clausura
normal~$V'$ de~$V$ en~$K'$.
\end{lemme}

Sean, en efecto,~$W'$ el normalizado de~$V'$ en~$L'$,~$\goth{m'}$ un ideal
maximal de~$V'$,~$\goth{n'}$ un ideal maximal de~$W'$ situado sobre
\ifthenelse{\boolean{orig}}
{$\goth{n'}$,}
{$\goth{m'}$,}
~$\goth{n}$ el ideal maximal que induce en el normalizado~$W$ de~$V$ en~$L$;
sean~$G,\ H,\ M$ los grupos de Galois de~$L,\ K',\ L'$ sobre~$K$, y sean
\ifthenelse{\boolean{orig}}
{$G_i,\ H_i,\ L_i'$}
{$G_i,\ H_i,\ M_i$}
los grupos de inercia correspondientes a los ideales maximales elegidos.
Entonces~$M$ se sumerge en el producto~$G\times H$ y~$M_i$ en el producto
$G_i\times H_i$, de modo que las proyecciones~$M\to G$ y~$M\to H$,
$M_i\to G_i$
\marginpar{280}
y~$M_i\to H_i$ sean sobreyectivas (sorites del cuerpo intermedio). Se deduce
ya, puesto que~$G_i$ y~$H_i$ son por hipótesis cíclicos de órdenes
\ifthenelse{\boolean{orig}}
{$m$ y~$n$}
{$n$ y~$m$}
primos a~$p$, que~$M_i$ tiene orden primo a~$p$ y, por tanto, es cíclico.
Como~$m$ es múltiplo de~$n$, los elementos de~$G_i\times H_i$ tienen
\ifthenelse{\boolean{orig}}
{potencia $m$-esima nula}
{potencia $m$-esima igual a la identidad};
por consiguiente,~$M_i$ tiene orden que divide a~$m$, y por tanto orden igual
a~$m$, pues~$M_i\to H_i$ es sobreyectivo. Este último homomorfismo es también
inyectivo. Ahora bien, su núcleo es el grupo de inercia de~$\goth{n'}$ sobre
$\goth{m'}$, lo cual demuestra que~$L'$ no está ramificada sobre~$K'$ en
$\goth{n'}$. De ahí el lema.

Situémonos ahora bajo las condiciones de~\Ref{X.2.4}, donde se tiene un
homomorfismo de especialización
$$
\pi_1(\overline X_1,\overline a_1)\to\pi_1(\overline X_0,\overline
a_0)
$$
que es \emph{sobreyectivo}, relativo a un morfismo propio y separable
$f\colon X\to Y$. Queremos precisar el núcleo de este homomorfismo. Procediendo
como en la demostración de~\Ref{X.2.4}, se ve que en esta cuestión siempre
podemos suponer que~$Y$ es el espectro de un \emph{anillo de valoración
discreta~$V$, completo y de cuerpo residual algebraicamente cerrado} (pues
siempre puede encontrarse un anillo semejante y un morfismo de su espectro~$Y'$
en~$Y$ cuya imagen sea~$\{y_0,y_1\}$). Se tiene entonces
$X_0=\overline X_0$,~$\kres(y_0)=k=$ cuerpo residual de~$V$ y
$\kres(y_1)=K=$ cuerpo de fracciones de~$V$. Sea~$K_s$ la clausura separable
de~$K$, sea~$\overline K$ su clausura algebraica y, para todo subanillo~$W$ de
$\overline K$ que contenga a~$V$, pongamos~$X_W=X\otimes_VW$. En particular,
se tiene
$$
X_V=\ifthenelse{\boolean{orig}}{V}{X}, \qquad X_K=X_1, \qquad X_{\overline K}=\overline X_1.
$$
Por otra parte, el morfismo canónico~$\overline X_1=X_{\overline K}\to X_{K_s}$
induce un isomorfismo en los grupos fundamentales (\Ref{IX.4.11}); por
consiguiente, teniendo en cuenta el isomorfismo~\Ref{X.2.1}
$\pi_1(X_0)\to\pi_1(X)$, quedamos reducidos a estudiar el homomorfismo
sobreyectivo
$$
\pi_1(X_{K_s})\to\pi_1(X)
$$
asociado
\marginpar{281}
al morfismo canónico~$X_{K_s}\to X$. Determinar el núcleo de este último
equivale a resolver el problema siguiente: \emph{se tiene un recubrimiento
principal conexo~$Z_{K_s}$ de~$X_{K_s}$, de grupo~$G$} (asociado, pues, a un
homomorfismo de~$\pi_1(X_{K_s})$ en~$G$); \emph{determinar en qué condiciones
es isomorfo a la imagen inversa de un recubrimiento principal~$Z$ de~$X$ de
grupo~$G$.}

Observemos primero que~$K_s$ es la unión filtrante creciente de sus
subextensiones finitas~$K'$ sobre~$K$ y que, por tanto,~$Z_{K_s}$ es isomorfo a
la imagen inversa de un recubrimiento principal~$Z_{K'}$ de~$X_{K'}$ para un
$K'$ conveniente (adviértase, sin embargo, que para~$K'$ fijado,~$Z_{K'}$ no
queda determinado de manera única). Decir que~$Z_{K_s}$ es isomorfo a la
imagen inversa de un recubrimiento principal~$Z$ de~$X$ significa que existe una
subextensión finita~$K''\supset K'$ de~$K_s$ tal que
$Z_{K''}=Z_{K'}\otimes_{K'}K''$ sea isomorfo a~$Z\otimes_VK''$. Para una
subextensión finita~$K'$ de~$K_s$, designemos ahora por~$V'$ el normalizado
de~$V$ en~$K'$, que es un anillo de valoración discreta completo, de cuerpo
residual~$k$. Así, el morfismo canónico~$X_{V'}\to X_V$ induce un isomorfismo
en las fibras sobre los puntos cerrados de~$Y=\Spec(V)$ y
$Y'=\Spec(V')$. Resulta entonces de~\Ref{X.2.1}, aplicado a~$X_V$ y~$X_{V'}$,
que el homomorfismo inducido en los grupos fundamentales
$\pi_1(X_{V'})\to\pi_1(X_V)$ es un isomorfismo; o, de manera equivalente, que
todo recubrimiento principal de~$X_{V'}$ es la imagen inversa de un
recubrimiento principal de~$X_V$, determinado salvo isomorfismo. Esto implica
el lema siguiente:

\begin{lemme}
\label{X.3.7}
Sea~$Z_{K'}$ un recubrimiento principal conexo de~$X_{K'}$ de grupo~$G$, y sea
\ifthenelse{\boolean{orig}}
{$Z_{K}$}
{$Z_{K_s}$}
su imagen inversa sobre~$X_{K_s}$\quoi. Para que este último sea isomorfo a la
imagen inversa de un recubrimiento principal~$Z$ de~$X$, es necesario y
suficiente que exista una extensión finita~$K''\supset K'$ de~$K$ en~$K_s$ tal
que el recubrimiento principal~$Z_{K''}$ de~$X_{K''}$ esté inducido por un
recubrimiento principal de~$X_{V''}$.
\end{lemme}

Supongamos, en particular, que los~$X_{V''}$ sean normales (para ello basta,
por ejemplo, que~$X_0$ sea normal y, a fortiori, que~$X_0$ sea simple,
\cf \Ref{I.9.1}).
\marginpar{282}
Como son conexos, son, pues, irreducibles. Sean~$L$ el cuerpo de funciones
racionales de~$X$ y~$X_{K}$,~$L'$ el de~$X_{V'}$ y~$X_{K'}$, y~$L''$ el de
$X_{V''}$ y~$X_{K''}$\quoi. Bajo las condiciones de~\Ref{X.3.7},~$Z_{K'}$
define una extensión finita separable~$R'$ de~$L'$, y~$Z_{K''}$ define la
extensión~$R''=R'\otimes_{L'}L'' = R'\otimes_{K'}K''$. La condición
considerada en~\Ref{X.3.7} significa también que existe una extensión finita
separable~$K''$ de~$K'$ tal que~$R''=R'\otimes_{K'}K''$ no esté
\emph{ramificada} sobre el esquema normal~$X_{V''}$, de cuerpo
$L''=L'\otimes_{K'}K''$, y no solamente sobre la parte abierta~$X_{K''}$
de~$X_{V''}$.

Supondremos en adelante que~$f\colon X\to Y$ es un morfismo \emph{liso}; por
tanto, los morfismos~$X_{V'}\to\Spec(V')$ son lisos y los esquemas~$X_{V'}$
son \emph{regulares}. Obsérvese que la fibra del punto cerrado de~$\Spec(V')$
en~$X_{V'}$ es irreducible y de codimensión~$1$. Sea~$\goth{o'}$ su anillo
local, que es, pues, un anillo de valoración discreta de cuerpo~$L'$ y de
cuerpo residual isomorfo al cuerpo de funciones racionales de~$X_0$, y por
tanto de la misma característica que~$k$. Definamos análogamente~$\goth{o''}$
en~$L''$; es evidentemente el normalizado de
\ifthenelse{\boolean{orig}}
{$\goth{o''}$}
{$\goth{o'}$}
en~$L''$. Se deduce entonces del teorema de pureza~\Ref{X.3.1} o~\Ref{X.3.3}
que, para que~$R''$ no esté ramificada sobre~$X_{V''}$, es necesario y
suficiente que~$R''$ no esté ramificada sobre~$\goth{o''}$, normalizado de
$\goth{o'}$ en~$L''$.

Observemos ahora que si~$u'$ es un uniformizante de~$V'$, también lo es
de~$\goth{o'}$. Si~$n$ es un entero primo a la característica~$p$ de~$k$,
\ifthenelse{\boolean{orig}}
{}
{tras ampliar primero~$K'$ de modo que contenga las raíces $n$-esimas de la
unidad,}
y si se toma
\ifthenelse{\boolean{orig}}
{$K''=K'[t]/(t^n-u)$}
{$K''=K'[t]/(t^n-u')$},
entonces
\ifthenelse{\boolean{orig}}
{$K'$}
{$K''$}
es una extensión galoisiana finita de~$K'$ y~$L''$ es isomorfo a
$L'[t]/(t^n-u')$; por tanto, está moderadamente ramificada sobre~$\goth{o'}$ y
su grupo de inercia tiene orden~$n$. Supongamos entonces que~$G$ \emph{tiene
orden primo a~$p$}; esto implica que~$R'$ está moderadamente ramificada
sobre~$\goth{o'}$. Tomemos para~$n$ un múltiplo, primo a~$p$, del orden del
grupo de inercia de~$R'$ sobre~$\goth{o'}$ (por ejemplo,
$n=[G\colon e]$). Aplicando el
\ifthenelse{\boolean{orig}}
{lema de \textsc{Abhyankar}}
{lema de \textsc{Abhyankar}}~\Ref{X.3.6},
se ve que se satisface la condición considerada en~\Ref{X.3.7}.

Esto demuestra el teorema siguiente:

\begin{theoreme}
\label{X.3.8}
Sean~$f\colon X\to Y$ un morfismo propio y liso, de fibras
\marginpar{283}%
geométricamente conexas, con~$Y$ localmente noetheriano, e~$y_0$ e~$y_1$ dos
puntos de~$Y$ tales que
\ifthenelse{\boolean{orig}}
{$y_0\in\overline{y_1}$}
{$y_0\in\overline{\{y_1\}}$};
sean~$\overline{X}_0$ y~$\overline{X}_1$ las fibras geométricas
correspondientes y considérese el homomorfismo de especialización
\eqref{X.2.4}~$\pi_1(\overline X_1)\to\pi_1(\overline X_0)$. Este
homomorfismo es sobreyectivo, y todo homomorfismo continuo de
$\pi_1(\overline X_1)$ en un grupo finito~$G$ de orden primo a la
característica~$p$ de~$\kres(y_0)$ proviene de un homomorfismo de
$\pi_1(\overline{X}_0)$ en~$G$.
\end{theoreme}

En otros términos:

\begin{corollaire}
\label{X.3.9}
Si~$\kres(y_0)$ tiene característica cero, el homomorfismo de especialización
es un isomorfismo. Si
\ifthenelse{\boolean{orig}}
{$p>0$}
{$\kres(y_0)$ tiene característica $p>0$},
entonces el núcleo del homomorfismo de especialización está contenido en la
intersección de los núcleos de los homomorfismos continuos de
$\pi_1(\overline X_1)$ en grupos finitos de orden primo a~$p$ (o, de manera
equivalente, en el subgrupo normal cerrado engendrado por un $p$-subgrupo de
Sylow del grupo de tipo galoisiano~$\pi_1(\overline X_1)$). Por consiguiente,
si
\ifthenelse{\boolean{orig}}
{$\pi_1(X_1)^{(p)}$}
{$\pi_1(\overline X_1)^{(p)}$}
designa el grupo cociente de~$\pi_1(\overline X_1)$ por el subgrupo cerrado
precedente, y si se define análogamente~$\pi_1(\overline X_0)^{(p)}$, el
homomorfismo de especialización induce un \emph{isomorfismo}
$$
\pi_1(\overline X_1)^{(p)}\isomto \pi_1(\overline X_0)^{(p)}
$$
\end{corollaire}

Obsérvese que la demostración de~\Ref{X.3.8} es puramente algebraica.
Procediendo como en~\Ref{X.2.6}, se concluye \emph{por vía trascendente}:

\begin{corollaire}
\label{X.3.10}
Sea~$X_0$ una curva propia, lisa y conexa de género~$g$ sobre un cuerpo
algebraicamente cerrado de característica~$p$. Con la notación introducida
en~\Ref{X.3.9}, el grupo~$\pi_1(X_0)^{(p)}$ es isomorfo a~$\Gamma^{(p)}$,
donde~$\Gamma$ es el grupo de tipo galoisiano engendrado por los generadores
$s_i,~t_i$ ($1\le i\le g$), ligados por la relación
$$
(s_1t_1s_1^{-1}t_1^{-1})\cdots (s_gt_gs_g^{-1}t_g^{-1})=1
$$
\end{corollaire}

\begin{remarques}
\label{X.3.11}
\marginpar{284}
Cuando~$\kres(y_0)$ tiene característica cero, el resultado~\Ref{X.3.9} es bien
conocido por vía trascendente. Obsérvese que la demostración de~\Ref{X.3.10}
recurre al teorema de pureza en el caso de característica mixta, pero
solamente para anillos de dimensión~2, caso en el que la demostración de dicho
teorema es fácil y se ha recordado en el texto.
\end{remarques}

\begin{sgabibliography}{0}{X.4}

\bibitem[1]{X.1} K\ptbl \textsc{Kodaira}, On compact analytic surfaces,
\ifthenelse{\boolean{orig}}
{Princeton University Press, 1960.}
{Ann.\ of Math.\ \textbf{71} (1960), p\ptbl 111--152.}

\bibitem[2]{X.2} S\ptbl \textsc{Lang} y
\ifthenelse{\boolean{orig}}{J.P\ptbl}{J-P\ptbl}%
\textsc{Serre}, Sur les rev\^etements non ramifi\'es des vari\'et\'es
alg\'ebriques,
\ifthenelse{\boolean{orig}}
{Amer.\ Journ. of Math., pp\ptbl 319--330, 1957.}
{Amer.\ J.\ Math.\ \textbf{79} (1957), p\ptbl319--330.}

\bibitem[3]{X.3} M\ptbl \textsc{Nagata},
On the purity of branch loci in regular local rings,
\ifthenelse{\boolean{orig}}
{Illinois Journal of Math., p\ptbl 328--333, 1959.}
{Illinois J.\ Math. \textbf{3} (1959), p\ptbl 328--333.}

\bibitem[4]{X.4}
\ifthenelse{\boolean{orig}}{J.P\ptbl}{J-P\ptbl}%
\textsc{Serre}, Rev\^etements ramifi\'es du plan projectif (d'apr\`es
S\ptbl Abhyankar), S\'eminaire Bourbaki, Mai 1960.

\bibitem[5]{X.5} O\ptbl \textsc{Zariski},
On the purity of the branch locus of algebraic functions,
\ifthenelse{\boolean{orig}}
{Proc.\ Nat.\ Acad.\ Sc.\ USA, t\ptbl 44, p\ptbl 791--796, 1958.}
{Proc.\ Nat.\ Acad.\ Sc.\ USA  \textbf{44} (1958), p\ptbl 791--796.}

\bibitem[6]{X.6} O\ptbl \textsc{Zariski},
Algebraic surfaces, Ergebnisse... 1948, Chelsea (New York).
\end{sgabibliography}

\chapterspace{-2}
\chapter{Ejemplos y complementos}
\label{XI}
\marginpar{285}

\section{Espacios proyectivos, variedades unirracionales}
\label{XI.1}

\begin{proposition}
\label{XI.1.1}
Sean $k$ un cuerpo algebraicamente cerrado y $X=\PP_k^r$ el espacio
proyectivo de dimensión $r$ sobre~$k$. Entonces $X$ es \emph{simplemente
conexo,} es decir, $\pi_1(X)=0$.
\end{proposition}

Para $r=0$, esto es trivial. Si $r=1$, hay que demostrar que, si $X'$ es
un recubrimiento étale conexo no vacío de~$X=\PP_k^1$, entonces
$X'\isomto X$. La fórmula del género nos da aquí, si~$g$ y $g'$ son
los géneros de~$X$ y~$X'$:
$$
1-g' = d(1-g)\quoi,
$$
donde $d$ es el grado de~$X'$ sobre~$X$. Como $g=0$, se tendrá entonces
$1-g'=d$, lo cual exige $d=1$, puesto que $g'\ge 0$; esto demuestra
$X'\isomto X$. Cuando $r\ge 2$, se procede por inducción
sobre~$r$, suponiendo que $\PP^{r'}$ es simplemente conexo para
$r'<r$. Aplicando esto a un hiperplano de~$\PP^r$ y utilizando
(\Ref{X.2.10}), resulta que $\PP^r$ es simplemente conexo. Otra
demostración: en virtud de (\Ref{X.1.7}), se tiene
$\pi_1(\PP^1\times\cdots\times \PP^1) = \pi_1(\PP^1)\times\cdots
\times \pi_1(\PP^1)$; por tanto, $(\PP^1)^r$ es simplemente conexo,
pues $\PP^1$ lo es, y entonces $\PP^r$ es simplemente conexo en virtud
de la invariancia birracional del grupo fundamental (\Ref{X.3.4}). Esta
demostración prueba, más generalmente:

\begin{corollaire}
\label{XI.1.2}
Sea $X$ un esquema propio y normal sobre un cuerpo algebraicamente
cerrado~$k$; si $X$ es una variedad racional, es decir, si es íntegro y
su cuerpo de funciones es una extensión puramente trascendente de~$k$,
entonces $X$ es simplemente conexo.
\end{corollaire}

Este resultado se aplica, en particular, a las variedades grassmannianas y,
más
\marginpar{286}%
generalmente, a las variedades~$G/H$, donde $G$ es un grupo lineal conexo
sobre~$k$ y $H$ un subgrupo algebraico que contiene un subgrupo de Borel
\ifthenelse{\boolean{orig}}
{de~$B$.}
{de~$G$.}

Recordemos que se llama variedad \emph{unirracional sobre~$k$}
\index{variedad unirracional|hyperpage}%
a un esquema propio e íntegro sobre~$k$ cuyo cuerpo de funciones $K$
está contenido en una extensión puramente trascendente~$K'$ de~$k$,
finita sobre~$K$ (es decir, que tiene el mismo grado de trascendencia
sobre~$k$ que~$K$), o, equivalentemente, si existe una aplicación racional
dominante $f\colon\PP_k^r\to X$, con $r=\dim X$. Si $X$ es normal,
se deduce de las consideraciones que preceden a \Ref{X.3.4} que, para todo
recubrimiento étale conexo $X'$ de~$X$, de cuerpo~$L/K$, el álgebra
$L\otimes_K K'$ sobre~$K'$ no está ramificada sobre el modelo~$\PP^r$ y,
por tanto, está completamente descompuesta en virtud de~\Ref{X.1.1}; esto
muestra que $L$ es $K$-isomorfo a una subextensión de~$K'/K$. Teniendo en
cuenta \Ref{V.8.2}, queda demostrado:

\begin{corollaire}
\label{XI.1.3}
El grupo fundamental de una variedad unirracional normal sobre un cuerpo
algebraicamente cerrado es finito.
\end{corollaire}

(N.B. Obsérvese que, en la definición de «unirracional», no era necesario
suponer que $K'/K$ fuera finita).

\ifthenelse{\boolean{orig}}
{\begin{remarques}
\label{XI.1.4}
Naturalmente, los resultados de este número son bien conocidos. Por otra
parte,
\ifthenelse{\boolean{orig}}{J.P\ptbl}{J-P\ptbl}%
Serre demostró~\cite{XI.10} que, cuando $X$ es una variedad proyectiva
unirracional lisa sobre un cuerpo algebraicamente cerrado de
\emph{característica cero,} $X$ es simplemente conexa. Su demostración es
trascendente, en el sentido de que utiliza el teorema de simetría de Hodge,
y se ignora si este resultado se extiende a la característica~$p>0$.
Parece, además, que no se conoce ningún ejemplo de variedad unirracional
lisa sobre~$k$ que no sea ya racional.
\end{remarques}}
{\begin{remarques}
\label{XI.1.4}
Naturalmente, los resultados de este número son bien conocidos. Por otra
parte,
\ifthenelse{\boolean{orig}}{J.P\ptbl}{J-P\ptbl}%
Serre demostró~\cite{XI.10} que, cuando $X$ es una variedad proyectiva
unirracional lisa sobre un cuerpo algebraicamente cerrado de
\emph{característica cero,} $X$ es simplemente conexa. Su demostración es
trascendente, en el sentido de que utiliza el teorema de simetría de Hodge.

\lcrochetbf Añadido en 2003 (MR) Limitémonos al caso de un cuerpo
algebraicamente cerrado $k$ de característica $p\geq 0$. En característica
$p>0$, la definición de $k$-variedad unirracional dada más arriba es la de
$k$-variedad débilmente unirracional, por oposición a la de $k$-variedad
fuertemente unirracional, en la que se supone además que $K'$ es una
extensión separable de~$K$.

En dimensión $2$, existen superficies proyectivas y lisas, débilmente
unirracionales, cuyo grupo fundamental (finito por \Ref{XI.1.3}) no es
trivial y que, por tanto, no son racionales (T\ptbl Shioda, On
unirationality of supersingular surfaces, Math.~Ann. \textbf{225} (1977),
p\ptbl155--159.). En cambio, una superficie fuertemente unirracional es
racional: esto resulta del criterio de racionalidad de Castelnuovo,
extendido a la característica $p>0$ por O\ptbl Zariski (\cf J-P\ptbl Serre,
Sém. Bourbaki \no146, 1957 y {\OE}uvres--Collected papers, vol\ptbl 1,
p\ptbl 467).

Sobre el cuerpo $\CC$ de los números complejos, se conocen ejemplos de
variedades proyectivas y lisas de \hbox{dimensión $\geq 3$} que son
unirracionales y no racionales (\cf P\ptbl Deligne, Variétés
unirationnelles non rationnelles [d'après M\ptbl Artin et D\ptbl Mumford],
Sém.\ Bourbaki \no402, 1971-72, Lecture Notes vol\ptbl 317). Este es, en
particular, el caso de una hipersuperficie cúbica lisa del espacio
proyectivo $\PP^4$. Algunos de estos ejemplos se extienden a característica
$>0$ y proporcionan variedades fuertemente unirracionales no racionales
(\cf J.P\ptbl Murre, Reduction of the proof of the non-rationality of a non
singular cubic threefold to a result of Mumford, Compositio Math.
\textbf{27} (1973), p\ptbl63--82).

Sea $V$ una $k$-variedad proyectiva, normal e íntegra. Se dice que $V$ es
racionalmente conexa si existen un $k$-esquema íntegro de tipo finito $T$ y
un $k$-morfismo $F : T\times\PP^1\to V$ tales que el morfismo
\begin{equation}\tag*{(1)}
\begin{split}
T \times\PP^1\times\PP^1&\to V\times V\\
(t, u, u')&\mto(F(t,u), F(t,u'))
\end{split}
\end{equation}
sea dominante. En particular, por dos puntos racionales suficientemente
generales de~$V$ pasa una curva racional. La terminología se justifica por
el hecho de que, si $V$ es racionalmente conexa, se pueden unir dos puntos
racionales mediante una cadena finita de curvas racionales. Si $k$ tiene
característica $p>0$, se refuerza la definición anterior exigiendo que la
aplicación (1) sea genéricamente lisa. Se obtiene así la noción de variedad
racionalmente conexa separable. Por tanto, las variedades unirracionales son
racionalmente conexas y, en característica $p>0$, las variedades
fuertemente unirracionales son racionalmente conexas separables. J\ptbl
Kollár demostró que las variedades racionalmente conexas separables tienen
grupo fundamental algebraico trivial (\cf O\ptbl Debarre, Variétés
rationnellement connexes [d'après T\ptbl Graber, J\ptbl Harris, J\ptbl
Starr et A.J\ptbl de~Jong], Sém. Bourbaki \no906, 2001-2002).\rcrochetbf
\end{remarques}}

\section{Variedades abelianas}
\label{XI.2}
Sean $k$ un cuerpo algebraicamente cerrado, $A$ una variedad abeliana
sobre~$k$, es decir, un esquema en grupos sobre~$k$, propio sobre~$k$,
liso sobre~$k$ y conexo, y, finalmente, $G$ un esquema en grupos conmutativos de tipo
finito sobre~$k$. Denotemos por $\Ext(A,G)$ el grupo de las clases de
extensiones conmutativas de~$A$ por~$G$, por $\H^1(A,G)$ el grupo de las
clases de fibrados principales sobre~$A$ de grupo~$G$ (comparar
\No \Ref{XI.4} más adelante), y consideremos el homomorfismo canónico
$$
\Ext(A,G)\to \H^1(A,G)\quoi.
$$
Un
\marginpar{287}
razonamiento de Serre~\cite[cap\ptbl VII, teor\ptbl 5]{XI.5} muestra que
este es un homomorfismo inyectivo cuya imagen es el conjunto de los
\og elementos primitivos\fg{} de~$\H^1(A,G)$, es decir, de los elementos
$\xi$ para los cuales se tiene:
$$
\pi^*(\xi) = \pr_1^*(\xi) + \pr_2^*(\xi)\quoi,
$$
donde los $\pr_i$ son las dos proyecciones de~$A\times A$ sobre~$A$, y
$\pi\colon A\times A\to A$ la ley de composición de~$A$ (N.B. Serre solo
enuncia su teorema para $G$ lineal y conexo y, por supuesto, liso sobre~$k$,
pero, simplificando la primera parte de su razonamiento, se ve que estas
restricciones son innecesarias: basta observar que todo morfismo de~$A$ a un
esquema en grupos $E$ de tipo finito sobre~$k$ que transforma la unidad en
la unidad es un homomorfismo de grupos, y aplicar esto a las secciones sobre
$A$ de una extensión $E$ de~$A$ por~$G$).

Vamos a aplicar este resultado al caso en que $G$ es un grupo finito
separable sobre~$k$, es decir, un grupo finito ordinario, que supondremos
conmutativo. Utilizando entonces
$\pi_1(A\times A)\isomto\allowbreak\pi_1(A)\times\pi_1(A)$
(\Ref{X.1.7}) e interpretando $\H^1(X,G)$ como
$\Hom(\pi_1(X),G)$ para todo esquema algebraico~$X$, en particular para
$X=A$ o $X=A\times A$, se ve que toda clase de~$\H^1(A,G)$ es primitiva;
por tanto, se tiene un isomorfismo
$$
\Ext(A,G)\isomto \H^1(A,G)\quoi,
$$
en otras palabras, \emph{todo recubrimiento principal de~$A$ de grupo
estructural conmutativo~$G$, punteado sobre el origen de~$A$, queda dotado
de manera única de una estructura de grupo algebraico que tiene al punto
marcado como origen y tal que $A'\to A$ es un homomorfismo de grupos
algebraicos.} En particular, si $A'$ es conexa, también es una variedad
abeliana isógena a~$A$.

Por otra parte, como el funtor $X\mto\pi_1(X)$ de los esquemas algebraicos
punteados $X$ en los grupos conmuta con el producto (\Ref{IX.1.7}),
transforma un grupo de la primera categoría en un grupo de la categoría de
grupos, es decir, en un grupo \emph{conmutativo}. Por tanto, \emph{si $A$ es
una variedad abeliana, $\pi_1(A)$
\marginpar{288}
es un grupo conmutativo.} Así, para conocer $\pi_1(A)$, basta conocer el
funtor $G\mto\H^1(A,G)=\Hom(\pi_1(A),G)$ cuando $G$ recorre los grupos
finitos \emph{conmutativos}. Finalmente, recordemos que, para todo entero
$n>0$, el homomorfismo de multiplicación por $n$ en~$A$,
$$
A\lto{n} A
$$
es sobreyectivo y, por tanto, de núcleo finito, es decir, es una isogenia,
y que de ello resulta que toda isogenia $A'\to A$ es cociente de una
isogenia del tipo anterior. De esto y de razonamientos estándar
(\cf por ejemplo,~\cite{XI.6}) se deduce:

\begin{theoreme}[Serre--Lang]
\label{XI.2.1}
Sea $A$ una variedad abeliana sobre un cuerpo algebraicamente cerrado~$k$.
Para todo entero $n>0$, consideremos el grupo finito ordinario $K_n$
subyacente al núcleo ${}_nA$ de la multiplicación por $n$ en~$A$ y, para
todo número primo~$\ell$, pongamos
$$
T_\ell(A)=\varprojlim_r K_{\ell^r}
$$
y
$$
T_{.}(A)=\underset{\ell}{\prod}T_\ell(A)=\varprojlim_n K_n
$$
(donde, si $m$ es múltiplo de~$n$, $m=ns$, se envía $K_m$ en $K_n$
mediante la multiplicación por~$s$). Entonces el grupo $\pi_1(A)$ es
canónicamente isomorfo a $T_{.}(A)$; por consiguiente, para todo número
primo~$\ell$, la componente $\ell$-primaria de~$\pi_1(A)$ es canónicamente
isomorfa a $T_\ell(A)$.
\end{theoreme}

Obsérvese que estos isomorfismos son funtoriales respecto de~$A$. El módulo
$T_\ell(A)$ se llama \emph{módulo $\ell$-ádico de Tate} de la variedad
abeliana~$A$. Es un funtor aditivo en~$A$; en particular, da lugar a una
representación del anillo $\Hom(A,A)$ de los endomorfismos de~$A$ sobre
$T_\ell(A)$, llamada \emph{representación $\ell$-ádica de Weil}, y que
desempeña un papel importante en la teoría de las variedades abelianas
(\cf por ejemplo, \hbox{\cite[cap.~VII]{XI.4}}). El
teorema~\Ref{XI.2.1} le da una interpretación como la representación
natural en el \emph{grupo de homología $\ell$-ádica} de~$A$,
$\H_1(A,\ZZ_\ell)=\pi_1(A)_\ell$, lo cual
\marginpar{289}
es evidentemente más satisfactorio a priori, en particular desde el punto
de vista de la fórmula de Lefschetz, \cite[cap.~V]{XI.4}. Recordemos aquí
los resultados de Weil sobre la estructura de~$T_\ell(A)$:
\begin{enumerate}
\item[a)] Si $n$ es primo con $\car(k)$, entonces $K_n$ es un módulo libre
de rango $2\dim A$ sobre~$\ZZ/n\ZZ$; por tanto, si $\ell$ es un número
primo $\neq\car(k)$, $T_\ell(A)$ es un módulo libre de rango $2\dim A$
sobre el anillo $\ZZ_\ell$ de los enteros $\ell$-ádicos;
\item[b)] Si $n$ es una potencia de~$\car(k)=p$, entonces $K_n$ es un
módulo libre de rango $\nu\leq\dim A$ sobre~$\ZZ/n\ZZ$, con $\nu$
independiente de~$n$; por tanto, $T_p(A)$ es un módulo libre de rango
$\nu\leq\dim A$ sobre el anillo $\ZZ_p$ de los enteros $p$-ádicos.
\end{enumerate}

Esto muestra que, en la teoría del grupo fundamental desarrollada aquí, el
grupo fundamental de una variedad abeliana variable no varía de manera
regular con el parámetro: su componente $p$-primaria puede disminuir
bruscamente para valores del parámetro $t$ correspondientes a una
característica residual~$p$; el caso más conocido de este fenómeno es el de
las curvas elípticas. Obsérvese, sin embargo, que, cualquiera que sea $n$
(primo o no con la característica), el verdadero núcleo ${}_nA$ en~$A$ de
la multiplicación por $n$ es un esquema en grupos finito sobre~$k$ de grado
$n^{2g}$, donde $g=\dim A$, que no será separable sobre~$k$ si $n$ es
múltiplo de~$p=\car(k)$. Además, cuando $A$ varía en una familia de
variedades abelianas, es decir, cuando se tiene un esquema abeliano $A$
sobre un esquema base $S$, se demuestra más generalmente que ${}_nA$ es un
esquema en grupos finito y plano sobre~$S$, de grado $n^{2g}$ sobre~$S$;
esto es, a condición de tener en cuenta las partes infinitesimales de los
núcleos ${}_nA$, estos se comportan de manera regular cualquiera que
sea~$n$. Esto sugiere que el \og verdadero\fg{} grupo fundamental de una
variedad abeliana $A$ es el progrupo algebraico (límite proyectivo formal
de grupos finitos sobre~$k$) $\varprojlim_n{{}_nA}$, donde por \og verdadero
grupo fundamental\fg{} de un esquema algebraico~$X$ hay que entender el
progrupo que clasifica los recubrimientos principales de~$X$ cuyo grupo
estructural es un grupo finito cualquiera $G$ sobre~$k$ (no necesariamente
separable sobre~$k$). De este modo, por ejemplo, las representaciones de
$\Hom(A,A)$ en la componente $p$-primaria del verdadero grupo fundamental
de~$A$ recuperan el polinomio característico de Weil, definido por este
último mediante los primos $\ell\neq p$, de forma más natural que la
construcción de Serre~\cite{XI.8}.

\section{Conos proyectivos, ejemplo de Zariski}
\label{XI.3}
\marginpar{290}

Sea, como antes y para simplificar, $k$ un cuerpo algebraicamente cerrado,
y sea $V$ un $k$-esquema proyectivo conexo, subesquema cerrado
de~$\PP^r_k$, que, si se desea, puede suponerse no singular. Sean
$Y=\hat C$ el cono proyectivo sobre~$V$, $y_0$ su vértice,
$X=\hat C_V$ la clausura proyectiva usual del fibrado vectorial
$C_V=\VV(\cal{O}_V(1))$ asociado a~$\cal{O}_V(1)$ y, por último,
$$
f\colon X\to Y
$$
el morfismo canónico que contrae a un punto la sección nula $X_0$
de~$C_V$ en~$X$ (EGA~II~8.6.4). Como $X$ es un fibrado localmente trivial
sobre~$V$, con fibras $\PP^1$ y, por tanto, simplemente conexas, el
morfismo $p\colon X\to V$ induce, en virtud
\ifthenelse{\boolean{orig}}
{de~(\Ref{XI.4}.)}
{de~(\Ref{X.1.4})}
un isomorfismo:
$$
\pi_1(X)\isomto\pi_1(V)\quoi.
$$
Como $p$ induce un isomorfismo $X_0\to V$, se concluye que
\emph{un recubrimiento étale de~$X$ está completamente descompuesto si y
solo si su restricción a~$X_0$ lo está}. Ahora bien, para todo recubrimiento
étale $Y'$ de~$Y$, la imagen inversa $X'=X\times_Y Y'$ es un recubrimiento
étale de~$X$ completamente descompuesto sobre la fibra~$X_0$ y, por tanto,
trivial. Como el homomorfismo $\pi_1(X)\to\pi_1(Y)$ es sobreyectivo
(\Ref{IX.3.4}), se concluye que
$$
\pi_1(Y)=(e)
$$
en otras palabras, \emph{todo cono proyectivo es simplemente conexo.}
(N.B.: en característica~$0$, el mismo resultado sigue siendo válido si se
toma como~$Y$ el cono afín.)

Supongamos ahora que $V$ es regular, \ie lisa sobre~$k$; entonces $X$ es
regular y, para una inmersión proyectiva conveniente de~$V$, se obtiene un
cono proyectivo~$Y$ \emph{normal}. Si $V$ no es simplemente conexa, y por
tanto tampoco lo es~$X$, sea $X'$ un recubrimiento étale conexo no trivial
de~$X$. Como las fibras de~$X$ sobre los puntos $y\in Y$ distintos de~$y_0$
se reducen a un punto, se ve que la restricción de~$X'$ a estas fibras
(en particular, a la fibra
\marginpar{291}
genérica) es trivial; sin embargo,
\ifthenelse{\boolean{orig}}
{$X$}
{$X'$}
no proviene por imagen inversa de un recubrimiento étale de~$Y$, pues $Y$
es simplemente conexa y $X'$ estaría completamente descompuesto. Esto
muestra que (\Ref{X.1.3} y~\Ref{X.1.4}) se vuelven falsos si se sustituye
la hipótesis de que $f$ es separable por la más débil de que sus fibras son
esquemas algebraicos separables (o incluso lisos) sobre los~$\kres(s)$.
Obsérvese asimismo que los grupos fundamentales de las fibras geométricas
$\overline X_y$ para $y\neq y_0$ se reducen evidentemente a~$(e)$, puesto
que esas fibras se reducen a un punto, mientras que $\pi_1(X_0)\neq e$;
por tanto, el teorema de semicontinuidad (\Ref{X.2.4}) también falla
para~$f$.

Indiquemos, por último, el ejemplo señalado por Zariski que hace fallar
estos mismos teoremas cuando se sustituye en ellos la hipótesis de que $f$
es separable por la de que $f$ es plano. Sea $f\colon X\to Y$ un morfismo
de una superficie proyectiva no singular a la recta racional $Y=\PP^1$,
tal que
\ifthenelse{\boolean{orig}}
{$K=k(x)$}
{$K=k(X)$}
sea una extensión «regular», \ie primaria y separable, de~$k(f)$ (\ie la
fibra genérica geométrica es conexa y separable), y tal que el divisor
$(f)=X_0-X_\infty$ sea
\ifthenelse{\boolean{orig}}
{un múltiplo $n$.º de un divisor}
{$n$ veces un divisor}
(donde $n$ es un entero primo con la característica). Es posible construir
tales ejemplos en toda característica. Sea $X'$ la normalización de~$X$ en
$K(f^{1/n})$, donde $K=k(X)$ es el cuerpo de funciones de~$X$. De la
hipótesis sobre~$(f)$ resulta que $X'$ es étale sobre~$X$. Sea $Y'$ el
normalización de~$Y$ en $k(t)(t^{1/n})$; está ramificada sobre~$Y$ exactamente
en los puntos $t=0$ y $t=\infty$,
\ifthenelse{\boolean{orig}}
{y la restricción $X'|f^{-1}(U)$ es isomorfa a la imagen inversa de~$Y'|U$.}
{y, poniendo $U=Y-\{0,\infty\}$, la restricción $X'|f^{-1}(U)$ es isomorfa a la imagen inversa de~$Y'|U$.}
En particular, la restricción de~$X'$ a la fibra genérica \emph{geométrica}
de~$X$ sobre~$Y$ se descompone completamente. Sin embargo, $X'$ no es
isomorfo a la imagen inversa de un recubrimiento étale de~$Y$, pues se ve
enseguida que este último sería necesariamente~$Y'$, lo cual es absurdo
porque $Y'$ está ramificado sobre~$Y$\kern2pt\footnote{Puede observarse,
desde el punto de vista de la «topología étale» (SGA~4~VII), que en este
ejemplo $\R^1(f_{\et})_*(\ZZ/n\ZZ)$ es «no separado» sobre
\ifthenelse{\boolean{orig}}{$S$}{$Y$}.
}.

He aquí (siguiendo a Serre) una manera sencilla de realizar las condiciones
de este ejemplo, inspirada en~\cite[\No 20]{XI.5}: se toma como~$n$ un
número primo $\geq5$, distinto de la característica, y se hace actuar
$G=\ZZ/n\ZZ$ sobre $\underline{k}^4$%
\ifthenelse{\boolean{orig}}
{}
{{\renewcommand{\thefootnote}{*}\addtocounter{footnote}{-1}\kern1pt\footnote{\lcrochetbf Añadido en 2003: $\underline{k}^4$ designa el espacio afín de dimensión~$4$ sobre~$k$.\rcrochetbf}}}
multiplicando las coordenadas por cuatro caracteres distintos de~$G$ (lo
cual es posible porque $n\geq5$). Entonces $G$ actúa sobre el espacio
proyectivo $\PP^3_k$, y los únicos puntos fijos por~$G$ son los cuatro
puntos correspondientes a los ejes coordenados. La superficie $X'$ de
ecuación
\ifthenelse{\boolean{orig}}
{$x^4+y^4+z^4+t^4=0$}
{$x^n+y^n+z^n+t^n=0$}
es lisa sobre
\marginpar{292}
$k$ (criterio jacobiano) y no contiene ninguno de los puntos fijos; por
tanto, como $G$ es de orden primo,
\ifthenelse{\boolean{orig}}
{actúa sobre~$X$ «sin puntos fijos», \ie $X$ es un recubrimiento principal de~$X=X'/G$ de grupo~$G$.}
{actúa sobre~$X'$ «sin puntos fijos», \ie $X'$ es un recubrimiento principal de~$X=X'/G$ de grupo~$G$.}
Sea $g=x/y$ en $k(X')=K'$; es un generador de Kummer de~$K'$ sobre
$K=k(X)$ si los caracteres elegidos eran~$\chi^i$, $i=0,1,2,3$, con
$\chi$ un carácter primitivo. Sea $f$ su potencia $n$-ésima, que es un
elemento de~$K$. Se ve enseguida que $K'$ es una extensión regular
de~$k(g)$; esto resulta de que la curva plana de ecuación homogénea en
$U,T,Z$,
$T^n+Z^n+(1+g^n)U^n=0$, es lisa sobre~$k(g)$ (criterio jacobiano), y de que
se sabe que toda curva plana es conexa. Por otra parte,
$k(f)=K\cap k(g)$, pues el segundo miembro es una extensión de~$k(f)$
contenida en la extensión $k(g)$ de grado primo y
\ifthenelse{\boolean{orig}}{distinto}{distinto}
de~$k(g)$ (pues $g\not\in K$). Esto implica que $K$ es una extensión regular
de~$k(f)$. Finalmente, el divisor de~$f$ sobre~$X$ es $n$ veces un divisor,
pues su imagen inversa sobre~$X'$ es el divisor de~$g^n$, luego es $n$
veces un divisor, y puede hacerse el descenso porque $X'$ es étale
sobre~$X$. Se habría terminado si la aplicación racional
$f\colon X\to\PP^1$ fuese un morfismo, es decir, si los divisores de los
ceros y de los polos de~$f$ no se intersecaran. De hecho, se comprueba
fácilmente (mirando de nuevo sobre~$X'$) que los dos divisores en cuestión
son respectivamente $n$ veces sendas curvas lisas sobre~$k$ que se cortan transversalmente en
un punto~$a$. Sustituyendo ahora $X$ por el esquema obtenido mediante una
explosión en~$a$, siguen verificándose las condiciones anteriores
($\divisor(f)$ divisible por~$n$, y $k(X_1)=k(X)$ extensión regular de
\ifthenelse{\boolean{orig}}{$k(X)$}{$k(f)$})
pero, además, $f$ es un \emph{morfismo} $X_1\to\PP^1$; por tanto, se está
en las condiciones buscadas.


\section{La sucesión exacta de cohomología}
\label{XI.4}

Sea~$S$ un preesquema, de modo que queda determinada la categoría
$\Sch_{/S}$ de los preesquemas sobre
\ifthenelse{\boolean{orig}}
{$X$}
{$S$},
y, por tanto, también la noción de grupo en ella, al que se llamará asimismo
\emph{preesquema en grupos sobre~$S$}, o simplemente \emph{$S$-grupo}.
Para simplificar la exposición y fijar las ideas, en lo sucesivo nos
limitaremos la mayoría de las veces a grupos \emph{afines} y \emph{planos}
sobre~$S$\kern1pt\footnote{En realidad, para lo que sigue bastaría la
hipótesis de cuasiafinidad en lugar de afinidad; \cf la nota al pie
\eqref{note-296}, p\ptbl\pageref{note-296}.}, lo cual bastará para las
aplicaciones que tenemos en mente. (Naturalmente, se presentan muchos casos
en los que no se verifica ninguna de las dos hipótesis.) Sea
\marginpar{293}
$G$ uno de tales $S$-grupos, y sea~$P$ un preesquema sobre~$S$ en el que~$G$
actúa por la derecha, lo cual implica, en particular, un morfismo
$$
\pi:P\times_SG\to P
$$
en la categoría~$\Sch_{/S}$, que satisface los axiomas conocidos. Se dice que
$P$ es \emph{formalmente homogéneo principal bajo}~$G$ si el morfismo
$$
P\times_S G\to P\times_S P
$$
de componentes~$\pr_1$ y~$\pi$ es un isomorfismo; equivalentemente, para todo
objeto~$S'$ de
\ifthenelse{\boolean{orig}}
{$\Sch_{/S'}$}
{$\Sch_{/S}$},
$P(S')=\Hom_S(S',P)$, considerado como conjunto con grupo de operadores
$G(S')=\Hom_S(S',G)$, es vacío u homogéneo principal (\ie vacío o isomorfo
a~$G(S')$, sobre el cual el grupo~$G(S')$ actúa por traslaciones a la
derecha). Se dice que~$P$ es \emph{trivial} si~$P$ es isomorfo a~$G$, sobre
el cual~$G$ actúa por traslaciones a la derecha o, lo que es equivalente, si
cada uno de los conjuntos con operadores~$P(S')$ bajo~$G(S')$ es trivial.
Se comprueba, por ejemplo mediante el procedimiento patentado de reducción al
caso conjuntista, que~$P$ \emph{es trivial si y solo
\ifthenelse{\boolean{orig}}
{si}
{si}
es formalmente homogéneo principal y admite una sección sobre~$S$} (este
último hecho se expresa en términos categóricos diciendo que~$P$ tiene una
sección sobre el objeto final~$e=S$ de~$\Sch_{/S}$, \ie que existe un
morfismo de~$e$ en~$P$). Para definir la noción de fibrado homogéneo
principal~$P$ bajo~$G$, más fuerte que la de fibrado formalmente homogéneo
principal, es preciso especificar primero en~$\Sch_{/S}$ un conjunto de
morfismos que se usarán para el «descenso» y desempeñarán el papel de
«morfismos de localización» para «trivializar» fibrados. La elección más
adecuada varía según el contexto, y ninguna contiene a todas las demás%
\footnote{\label{note-293}Véase al respecto SGA~3~IV, en particular~\S 4.}.
Aquí convendrá adoptar la definición siguiente:

\begin{definition}
\label{XI.4.1}
Sea~$G$ un $S$-grupo. Se llama \emph{fibrado homogéneo principal} (por la
derecha) bajo~$G$ a un $S$-preesquema~$P$ con acción por la derecha del
$S$-grupo~$G$, tal que exista un recubrimiento de~$S$ por abiertos~$U_i$ y,
para todo~$i$, un morfismo de cambio de base~$S'_i\to U_i$ fielmente plano y
cuasicompacto, tal que
\ifthenelse{\boolean{orig}}
{$P'_i=P\times_S S'_i$ sea un preesquema con operadores trivial bajo
$G'=G\times_S S'$.}
{$P'=P\times_S S'$ sea un preesquema con operadores trivial bajo
$G'=G\times_S S'$, donde~$S'$ es el $S$-preesquema suma disjunta de los~$S'_i$.}
\end{definition}

(Obsérvese que el funtor de cambio de base~$X\mto X'=X\times_S S'$,
\marginpar{294}
por ser exacto por la izquierda, transforma grupos en grupos y objetos con
grupos de operadores en objetos con grupos de operadores.) Obsérvese que
\Ref{XI.4.1} es \emph{estable por cambio de base}. Obsérvese también lo
siguiente:

\begin{proposition}
\label{XI.4.2}
Sean~$G$ un $S$-grupo plano y cuasicompacto sobre~$S$, y~$P$ un
$S$-preesquema en el que~$G$ actúa por la derecha. Son equivalentes las
condiciones siguientes:
\begin{enumerate}
\item[(i)] $P$ es un fibrado homogéneo principal
\ifthenelse{\boolean{orig}}
{bajo~$G$.}
{bajo~$G$;}
\item[(ii)] $P$ es formalmente homogéneo principal bajo~$G$,
\index{formalmente homogéneo principal (preesquema)|hyperpage}%
y el morfismo estructural~$P\to S$ es fielmente plano y cuasicompacto.
\end{enumerate}
\end{proposition}

Si~$P$ es homogéneo principal bajo~$G$, entonces, con las notaciones
de~\Ref{XI.4.1}, $P'$ es fielmente plano y cuasicompacto sobre~$S'$ (pues
$G'$ lo es y~$P'$ es $S'$-isomorfo a él), luego~$P$ tiene las mismas
propiedades sobre~$S$ (para «sobreyectivo» y «cuasicompacto», \cf
\Ref{VIII.3.1}; para «plano», se trata de una omisión en los sorites del
Exposé~\Ref{VIII}). Recíprocamente, si se verifica~(ii), tomemos el cambio de
base~$S'=P$, que es fielmente plano y cuasicompacto sobre~$S$; entonces~$P'$
será formalmente homogéneo principal sobre~$S'$, puesto que~$P$ lo es
sobre~$S$ y el funtor de cambio de base es exacto por la izquierda; por otra
parte, $P'$ tiene una sección sobre~$S'$,
\ifthenelse{\boolean{orig}}
{a saber}
{a saber}
la sección diagonal, por lo que es un fibrado principal trivial, lo cual
concluye la demostración.

\begin{corollaire}
\label{XI.4.3}
Si~$G$ es afín y plano sobre~$S$, todo fibrado homogéneo principal~$P$
bajo~$G$ es afín y plano sobre~$S$.
\end{corollaire}

En efecto, adquiere esas propiedades tras una extensión fielmente plana y
cuasicompacta de la base, y se aplica~\Ref{VIII.5.6}.

La utilidad de la definición~\Ref{XI.4.1} para los $S$-grupos \emph{planos}
y \emph{afines} sobre~$S$ se debe a~\Ref{VIII.2.1}, \ie al hecho de que los
morfismos~$S'\to S$ considerados en~\Ref{XI.4.1} son morfismos de descenso
efectivo para la categoría de los preesquemas afines sobre otros
preesquemas. Gracias a este hecho, la comprobación de los hechos esbozados a
continuación se efectúa de manera esencialmente «categórica»%
\footnote{\Cf \loccit en la nota al pie~\eqref{note-293},
p\ptbl\pageref{note-293}.}. Sea~$E$ un $S$-preesquema en el que el
$S$-grupo~$G$ actúa por la izquierda, y sea~$P$ un fibrado homogéneo principal
(por la derecha) bajo~$G$; queremos definir un fibrado asociado~$E^{(P)}$,
«localmente» isomorfo a~$E$. Para ello, hagamos actuar
\marginpar{295}
$G$ por la derecha sobre~$P\times_S E$ según la ley~$(x,y)\mto
(xg,g^{-1}y)$, que describe tales acciones en el contexto conjuntista y se
extiende a las categorías mediante el procedimiento patentado. A reserva de
su existencia, pondremos:
$$
E^{(P)}=(P\times_S E)/G
$$
Con ello se comprueba que~$P\times_S E$ será un preesquema sobre
$T=E^{(P)}$, con grupo de operadores por la derecha~$G_T=G\times_S T$; para
trabajar con comodidad, querríamos además que~$P\times_S E$ fuera un fibrado
homogéneo principal sobre~$T$ de grupo~$G_T$. Para comprobar la existencia
de~$E^{(P)}$ y la propiedad precedente, retomemos el~$S'$ de la
definición~\Ref{XI.4.1} y consideremos sobre~$S'$ la imagen inversa de la
situación inicial. Puesto que~$P'$ es
\ifthenelse{\boolean{orig}}
{trivial}
{trivial,}
\ie isomorfo a~$G'_d$, se ve inmediatamente que~$E^{\prime(P')}$ existe y
posee la propiedad de exactitud requerida. En efecto,
$E'\times_{S'} P'$ es $G'$-isomorfo al producto~$E'\times_{S'} G'$, luego
$E^{\prime(P')}$ es isomorfo a~$E'$. Además, la formación del «fibrado
asociado» en el caso de un espacio con operadores trivial conmuta con toda
extensión de la base y, tomando
\ifthenelse{\boolean{orig}}
{en este caso}
{en este caso}
las diversas extensiones de la base
\begingroup
\shorthandoff{<>}% Spanish babel trata de otro modo los delimitadores de curvatura de Xy-pic.
$\xymatrix@C=.5cm{S''\ar@<2pt>[r]\ar@<-2pt>[r]&S'}$
y
$\xymatrix@C=.5cm{S'''\ar@<4pt>[r]\ar@<0pt>[r]\ar@<-4pt>[r]&S'}$%
\endgroup, donde
$S''$ y~$S'''$ son los productos fibrados doble y triple de~$S'$ sobre~$S$,
se comprueba que~$E^{\prime(P')}$ \emph{está provisto de un dato de descenso
relativo al morfismo~$S'\to S$, y que~$E^{(P)}$ existe con las propiedades
requeridas si y solo si este dato de descenso es efectivo}; naturalmente,
$E^{(P)}$ no es entonces sino el objeto descendido. (Úsese el hecho de que
$S'\to S$ es un morfismo de descenso en la categoría de los
$S$-preesquemas, \cf \Ref{VIII.5.2}.) Se sigue que \emph{el fibrado asociado
existe si~$E$ es afín sobre~$S$}. Aplicaremos esta construcción al caso en
que se tiene un homomorfismo de $S$-grupos~$G\to H$ y se toma como~$E$ el
$S$-preesquema~$H$, dotado de las acciones por la izquierda de~$G$ sobre~$H$
inducidas por el morfismo dado; como~$H$ actúa por la derecha sobre sí mismo
de manera que conmuta con las acciones de~$G$ sobre~$H$, y como la formación
del fibrado asociado (cuando existe sobre~$S$) conmuta con la extensión de la
base, se comprueba fácilmente que~$H$ actúa por la derecha sobre
\ifthenelse{\boolean{orig}}
{$P^{(H)}$}
{$H^{(P)}$},
que es, por consiguiente, un fibrado homogéneo principal bajo~$H$ en el
sentido de~\Ref{XI.4.1}, y
\marginpar{296}
más precisamente queda
\ifthenelse{\boolean{orig}}
{trivializado}
{trivializado}
por el mismo morfismo~$S'\to S$ que~$P$. En particular, \emph{a todo fibrado
homogéneo principal~$P$ bajo~$G$ y a todo homomorfismo de $S$-grupos
$G\to H$, con~$H$ afín sobre~$S$, se asocia un fibrado homogéneo principal de
grupo~$H$}, de manera funtorial en~$(G\to H)$ y compatible con los cambios
de base cualesquiera~$T\to S$.


\begin{definition}
\label{XI.4.4}
Sea~$G$ un $S$-preesquema. Se denota por~$\H^0(S,G)$ el conjunto de las
secciones de~$G$ sobre~$S$, que se considerará como grupo cuando~$G$ sea un
$S$-grupo. En tal caso, se denota por~$\H^1(S,G)$ el conjunto de las clases,
salvo isomorfismo, de fibrados homogéneos principales sobre~$S$ de grupo
\ifthenelse{\boolean{orig}}
{$S$}
{$G$},
considerando~$\H^1(S,G)$ dotado del «punto marcado» que corresponde a los
fibrados triviales\footnote{\label{note-296}Esta notación solo es coherente
\ifthenelse{\boolean{orig}}
{con respecto a}
{con respecto a}
las notaciones cohomológicas generales (SGA~4~V) cuando se dispone de
criterios de efectividad del descenso, que apenas están garantizados salvo si
$G$ es afín (o tan solo cuasiafín, \cf (\Ref{VIII.7.9})).}.
\end{definition}

Así, $\H^0(S,G)$ es un funtor de la variable $S$-preesquema~$G$, con valores
en la categoría de los conjuntos. Este funtor es exacto por la izquierda y,
con mayor razón, conmuta con los productos finitos, lo cual implica, en
efecto, que transforma grupos en grupos y grupos conmutativos en grupos
conmutativos. De manera análoga, $\H^1(S,G)$ es un funtor de la variable
$S$-grupo \emph{afín}~$G$, con valores en la categoría de los conjuntos por
medio de la formación de fibrados asociados; se comprueba fácilmente que
este funtor conmuta con los productos finitos. En particular, transforma los
grupos de la categoría de los $S$-grupos afines, \ie los $S$-grupos
\emph{afines conmutativos}, en grupos de la segunda categoría, e incluso en
grupos conmutativos (pues los grupos de la primera categoría son
conmutativos). Así, \emph{si~$G$ es un $S$-grupo afín conmutativo,
$\H^1(S,G)$ es un grupo conmutativo}, y un homomorfismo~$G\to H$ de
$S$-grupos afines conmutativos da lugar a un homomorfismo de grupos
$\H^1(S,G)\to\H^1(S,H)$.

Para simplificar, en lo sucesivo nos limitaremos a considerar $S$-grupos
\emph{afines y conmutativos}. Sea
$$
0\to G'\lto{u} G\lto{v}G''\to 0
$$
una sucesión de morfismos de tales grupos; \emph{diremos que esta sucesión es
exacta si~$vu=0$} (lo que permite considerar~$G$ como un preesquema
sobre~$G''$, con grupo de operadores por la derecha~$G'_{G''}$) \emph{y
si~$G$ es un fibrado homogéneo principal
\marginpar{297}%
sobre~$G''$ de grupo~$G'_{G''}=G'\times_S G''$}. Esto implica, en particular,
que~$u\colon G'\to G$ es un núcleo de~$v$ y, con mayor razón, implica la
exactitud de la sucesión
$$
\ifthenelse{\boolean{orig}}
{0\to\H^0(X,G')\to\H^0(X,G)\to\H^0(X,G'').}
{0\to\H^0(S,G')\to\H^0(S,G)\to\H^0(S,G'').}
$$
Esto implica, además, la posibilidad de definir una aplicación
$$
\ifthenelse{\boolean{orig}}
{\partial\colon\H^0(X,G'')\to\H^1(X,G')\quoi,}
{\partial\colon\H^0(S,G'')\to\H^1(S,G')\quoi,}
$$
asociando a toda sección de~$G''$ sobre~$S$, \ie a todo $S$-morfismo
$f\colon S\to G''$, el fibrado homogéneo principal~$P_f$ de grupo
$G'\simeq f^*(G'_{G''})$ sobre~$S$, imagen inversa del fibrado homogéneo
principal~$G$ sobre~$G''$. Desde el punto de vista de los $S$-preesquemas, no
es, pues, sino la imagen inversa por~$v\colon G\to G''$ del subpreesquema
imagen de~$S$ por la inmersión~$f$, y las acciones de~$G'$ sobre~$P_f$ son
inducidas por las acciones por la derecha de~$G'$
\ifthenelse{\boolean{orig}}
{sobre~$G$).}
{sobre~$G$.}
Dejamos asimismo al lector la comprobación de la proposición siguiente, que
no presenta más dificultad que la de su redacción:

\begin{proposition}
\label{XI.4.5}
La aplicación
\ifthenelse{\boolean{orig}}
{$\partial\colon \H^0(X,G'')\to\H^1(X,G')$}
{$\partial\colon \H^0(S,G'')\to\H^1(S,G')$}
es un homomorfismo de grupos. La sucesión de homomorfismos siguiente es
exacta:
\begin{multline*}
\ifthenelse{\boolean{orig}}
{0\to\H^0(X,G')\lto{u^0}\H^0(X,G)\lto{v^0}\H^0(X,G')\lto{\partial}\H^1(X,G'')\\
\lto{u^1}\H^1(X,G) \lto{v^1}\H^1(X,G'')}
{0\to\H^0(S,G')\lto{u^0}\H^0(S,G)\lto{v^0}\H^0(S,G'')\lto{\partial}\H^1(S,G')\\
\lto{u^1}\H^1(S,G) \lto{v^1}\H^1(S,G'')}
\end{multline*}
(donde los homomorfismos distintos de~$\partial$ proceden de la ley
funtorial de~$\H^0$ \resp de~$\H^1$).
\end{proposition}

\begin{remarques}
\label{XI.4.6}
El punto de vista expuesto aquí para el estudio de los fibrados homogéneos
principales está visiblemente inspirado en Serre~\cite{XI.7}, cuya consulta
será de gran provecho para el lector. Cuando se desea un formalismo que se
aplique también a $S$-grupos estructurales cuasiproyectivos sobre~$S$ (de
manera que abarque, en particular, los esquemas abelianos proyectivos),
conviene modificar~\Ref{XI.4.1} exigiendo que~$S'$ sea suma de
preesquemas~$S'_i$ finitos y localmente libres sobre abiertos~$S_i$ de~$S$
que recubran~$S$. Los desarrollos precedentes siguen siendo entonces
válidos, incluido en particular~\Ref{XI.4.5}, sustituyendo en todas partes la
hipótesis de afinidad por la de cuasiproyectividad e interpretando de manera
correspondiente la definición dada
\ifthenelse{\boolean{orig}}
{anteriormente}
{anteriormente}
de una sucesión exacta de $S$-grupos. En efecto, basta sustituir la referencia
a~(\Ref{VIII.2.1})
\marginpar{298}
por~(\Ref{VIII.7.7}): los
\ifthenelse{\boolean{orig}}
{morfismos utilizados}
{morfismos utilizados}
$S'\to S$ son morfismos de descenso efectivo para la categoría fibrada de
los preesquemas cuasiproyectivos sobre otros preesquemas. Sin embargo, debe
advertirse que esta segunda noción de fibrado homogéneo principal es más
restrictiva que la primera de~\Ref{XI.4.1}.
\end{remarques}

\sganum
\label{XI.4.7}
Se obtiene una noción aún más restrictiva de fibrado homogéneo principal
exigiendo que~$S$ esté recubierto por abiertos~$S_i$ tales que, para todo~$i$,
$P|S_i$ sea un fibrado con operadores trivial bajo~$G|S_i$; se dirá entonces
que~$P$ es un fibrado homogéneo principal \emph{localmente trivial}.
\index{homogéneo principal localmente trivial (fibrado)|hyperpage}%
Las clases de estos fibrados, para~$G$ dado, forman una parte de
\ifthenelse{\boolean{orig}}
{$\H^1(X,G)$}
{$\H^1(S,G)$},
que está en correspondencia biyectiva con
\ifthenelse{\boolean{orig}}
{$\H^1(X,\cal{O}_X(G))$}
{$\H^1(S,\cal{O}_S(G))$},
donde
\ifthenelse{\boolean{orig}}
{$\cal{O}_X(G)$}
{$\cal{O}_S(G)$}
es el haz (en el sentido ordinario) de las secciones
de~$G$ sobre~$S$, \cf \cite{XI.2}. Para que estos~$\H^1$ sigan dando lugar a
una sucesión exacta de cohomología~\Ref{XI.4.5}, es preciso suponer,
evidentemente, que la sucesión~$0\to G'\to G\to G''\to 0$ sea exacta en el
sentido razonable para este nuevo contexto, \ie que~$G$ sea un fibrado
localmente trivial sobre~$G''$, de grupo~$G'_{G''}$; esto significa también
que~$u\colon G'\to G$ es un núcleo de~$v\colon G\to G''$ y que~$G$ admite
localmente una sección sobre~$G''$.

\sganum
\label{XI.4.8}
Es, evidentemente, muy deseable prolongar la sucesión exacta~\Ref{XI.4.5}
introduciendo los grupos de cohomología superiores
\ifthenelse{\boolean{orig}}
{$\H^i(X,G)$}
{$\H^i(S,G)$}.
Esto es posible situándose en el marco de la «cohomología de Weil»: se
considera la categoría~$\cal{B}$ de los preesquemas cuasicompactos sobre~$S$,
\ifthenelse{\boolean{orig}}
{dotada}
{dotada}
del conjunto~$\cal{M}$ de los morfismos fielmente planos y cuasicompactos,
que se llamarán morfismos localizantes. Un «haz de Weil» abeliano sobre~$S$
(o, mejor, sobre~$(\cal{B},\cal{M})$) es entonces un funtor contravariante
$\cal{F}$ de~$\cal{B}$ en la categoría de los grupos abelianos, que transforma
sumas en productos y una sucesión
\begingroup
\shorthandoff{<>}% Spanish babel trata de otro modo los delimitadores de curvatura de Xy-pic.
${\def\labelstyle{\scriptstyle}\xymatrix@C=.5cm{T''=T'\times_TT'\ar@<+2pt>[rr]^-{\pr_1,\pr_2}\ar@<-2pt>[rr] &&T'\ar[r]^f & T}}$,
con~$f\in\cal{M}$, en un diagrama \emph{exacto} de conjuntos
$\xymatrix@C=.5cm{\cal{F}(T)\ar[r] & \cal{F}(T')
\ar@<2pt>[r]\ar@<-2pt>[r] & \cal{F}(T'')}$.%
\endgroup\ 
Los haces de Weil forman una categoría abeliana con límites inductivos
exactos que admite un generador y, por tanto, que posee suficientes objetos
inyectivos~\cite{XI.1}. Los funtores derivados por la derecha del funtor
$\Gamma(\cal{F})=\cal{F}(S)$ se denotan entonces por~$\H^i(S,\cal{F})$. Por
otra parte, todo $S$-grupo conmutativo define evidentemente un haz de Weil
(\Ref{VIII.5.2}), cuyos $\H^0$ y~$\H^1$ no son sino~$\H^0(S,G)$ y
$\H^1(S,G)$, lo que permite definir de manera razonable los demás
$\H^i(S,G)$. Se
\marginpar{299}
demuestra, además, que una sucesión exacta de $S$-grupos define una sucesión
exacta de haces de Weil, lo que permite recuperar y prolongar la sucesión
exacta~\Ref{XI.4.5}\footnote{Para un estudio sistemático de este punto de
vista, \cf SGA~4~I a~IX.}.

\sganum
\label{XI.4.9}
Sería conveniente desarrollar las variantes no conmutativas de~\Ref{XI.4.5},
como en~\cite{XI.2}. Para un desarrollo sistemático, en el marco adecuado,
de las diversas nociones cohomológicas esbozadas en el presente número,
remitimos a un trabajo entonces en preparación de
J\ptbl\ \textsc{Giraud}\footnote{\Cf J\ptbl\ \textsc{Giraud},
\ifthenelse{\boolean{orig}}
{\emph{Álgebra homológica no abeliana}, de próxima aparición en}
{\emph{Cohomología no abeliana},}
Springer-Verlag 1971.}.


\section{Casos particulares de fibrados principales}
\label{XI.5}

Supongamos ahora que~$S$ es conexo y está dotado de un punto
geométrico~$a$, de donde se obtiene un grupo fundamental~$\pi_1(S,a)$ que
permite clasificar los recubrimientos étale de~$S$: la categoría de los
recubrimientos étale de~$S$ es equivalente a la categoría de los conjuntos
finitos sobre los que~$\pi_1$ actúa continuamente. Se sigue que un esquema
en grupos~$G$, finito y étale sobre~$S$, está determinado esencialmente por
un grupo finito ordinario~$\cal{G}$, sobre el cual~$\pi_1$ actúa continuamente
por automorfismos de grupo. Un recubrimiento étale~$P$ de~$S$, sobre el que
$G$ actúa por la derecha, está determinado esencialmente por un conjunto
finito~$\cal{P}$ sobre el que~$\pi_1$ actúa continuamente (por la izquierda),
y sobre el que~$\cal{G}$ actúa por la derecha de manera compatible con las
acciones de~$\pi_1$:
$$
s(p\cdot g)=(sp)\cdot(sg)\qquad \text{para }s\in\pi_1,\
p\in\cal{P},\ g\in\cal{G}.
$$
Se comprueba que~$P$ es un fibrado homogéneo principal en el sentido
de~\Ref{XI.4.1} si y solo si~$\cal{P}$ es un conjunto homogéneo principal
bajo~$\cal{G}$ (úsese, por ejemplo, el criterio~\Ref{XI.4.2}). En otros
términos, \emph{la categoría de los fibrados homogéneos principales sobre~$S$
de grupo~$G$ es equivalente a la categoría de los fibrados homogéneos
principales de grupo
\ifthenelse{\boolean{orig}}
{$G$}
{$\cal G$}
en la categoría de los conjuntos finitos sobre los que~$\pi_1$ actúa
continuamente.} Se deduce, en particular, una biyección canónica, funtorial
en~$G$:
\begin{equation*}
\label{eq:XI.5.etoile}
\tag{$*$} \H^1(S,G)\isomto \H^1(\pi_1,\cal{G})\quoi,
\end{equation*}
donde
\marginpar{300}
el segundo miembro denota el conjunto de las clases, salvo isomorfismo, de
los fibrados homogéneos principales bajo~$\cal{G}$ en la categoría de los
conjuntos finitos sobre los que~$\pi_1$ actúa (por lo demás, no hace falta
precisar: continuamente), conjunto que, como es bien sabido, se explicita
como el conjunto cociente del conjunto
\ifthenelse{\boolean{orig}}
{$Z^1(\pi_1,G)$}
{$Z^1(\pi_1,\cal G)$}
de los $1$-cociclos~$\varphi:\pi_1\to \cal{G}$ (que satisfacen
$\varphi(1)=1$, $\varphi(st)=\varphi(s)(s.\varphi(t))$) por el grupo~$\cal{G}$,
que actúa sobre él de manera
\ifthenelse{\boolean{orig}}
{natural).}
{natural.}

\label{page-300}
Un caso importante es aquel en que~$\pi_1$ actúa trivialmente sobre~$\cal{G}$,
\ie cuando~$G$ es un recubrimiento completamente descompuesto de~$S$,
isomorfo a la suma de~$\cal{G}$ ejemplares de~$S$; se escribe entonces también
$\H^1(S,\cal{G})$ en lugar de~$\H^1(S,G)$, y este conjunto está en
correspondencia biyectiva~\eqref{eq:XI.5.etoile} con
$\H^1(\pi_1,\cal G)= \Hom(\pi_1,\cal{G})/$automorfismos interiores
de~$\cal{G}$. Obsérvese, además, que en este caso un fibrado homogéneo
principal sobre~$S$ de grupo~$G$ no es sino un \emph{recubrimiento principal}
de~$S$ con grupo~$\cal{G}$ (\Ref{V.2.7}), y la correspondencia biyectiva
precedente es la que se deduce de la correspondencia entre recubrimientos
principales de~$S$ con grupo~$\cal{G}$, \emph{punteados} sobre~$a$, y los
homomorfismos continuos de~$\pi_1(S,a)$ en~$\cal{G}$
(\Ref{V}, al final de~\Ref{V.5}).

El interés de relacionar la teoría de los recubrimientos étale con la de los
fibrados principales (ya implícita en \textsc{A\ptbl\ Weil},
\emph{Generalización de las funciones abelianas}, y explicitada por
\textsc{S\ptbl\ Lang} en su teoría geométrica del cuerpo de clases, \cf
Serre~\cite{XI.5}) proviene del hecho de que todo $S$-grupo finito y étale
sobre~$S$ puede sumergirse en un $S$-grupo~$H$, afín y liso sobre~$S$, de
fibras conexas,
\ifthenelse{\boolean{orig}}
{conmutativo cuando~$G$ lo es de modo que}
{conmutativo cuando~$G$ lo es; así,}
mediante la sucesión
exacta~\Ref{XI.4.5} (y, eventualmente, sus variantes no conmutativas), la
clasificación «discreta» de los recubrimientos principales de grupo~$G$ puede
estudiarse con ayuda de la clasificación «continua» de los fibrados
principales de grupo~$H$, y también recíprocamente. Para la idea de la
construcción general de la inmersión de~$G$ en~$H$ (al parecer bastante poco
utilizada en la práctica), véase~\cite[VI~2.8]{XI.5}. Nos limitaremos a
desarrollar en el número siguiente dos casos particulares importantes, por
lo demás clásicos. Necesitaremos allí un resultado auxiliar:

\begin{proposition}
\label{XI.5.1}
Sea
\marginpar{301}
$S$ un preesquema y~$G$ un $S$-grupo isomorfo a~$\Gl(n)_S$ (por ejemplo,
$\GG_{m\,S}$) o a~$\GG_{a\,S}$; entonces todo fibrado homogéneo principal
bajo~$G$ es localmente trivial.
\end{proposition}

Precisemos que~$\Gl(n)_S$ ($n$ entero~$\geq 0$) denota el $S$-grupo que
representa el funtor contravariante~$T\mto
\Gl(n,\Gamma(T,\cal{O}_T))$ de la variable $S$-preesquema~$T$; en particular,
$\GG_{m\,S}$ («grupo multiplicativo sobre~$S$») representa el funtor
contravariante~$T\mto \Gamma(T,\cal{O}_T^*)$ y, por tanto, como preesquema
sobre~$S$, es isomorfo a~$\Spec {\cal{O}_S[t,t^{-1}]}$, donde~$t$ es una
indeterminada. Del mismo modo, $\GG_{a\,S}$ representa el funtor
contravariante~$T\mto \Gamma(T,\cal{O}_T)$ y es, por tanto, isomorfo como
$S$-preesquema a~$\Spec(\cal{O}_S[t])$, donde~$t$ es una indeterminada.
Obsérvese que, mediante una reducción por etapas, \Ref{XI.5.1} recupera el
resultado de
\ifthenelse{\boolean{orig}}
{trivialidad local}
{trivialidad local}
de Rosenlicht, relativo al caso en que~$G$ admite una «serie de composición»
cuyos factores sucesivos son grupos del tipo considerado aquí. (Para un
estudio más fino de las cuestiones de
\ifthenelse{\boolean{orig}}
{trivialidad local}
{trivialidad local}
de los fibrados homogéneos principales, \cf
\cite{XI.7}~y~\cite{XI.3}).

La primera afirmación se demuestra observando que
$G(T)=\Aut(\cal{O}_T^n)$ y que los morfismos~$S'\to S$ que intervienen
en~\Ref{XI.4.1} (\ie que son fielmente planos y cuasicompactos) son morfismos
de descenso efectivo para la categoría fibrada de los
\ifthenelse{\boolean{orig}}
{módulos}
{módulos}
localmente
isomorfos a~$\cal{O}_T^n$, \ie localmente libres de rango~$n$
(\Ref{VIII.1.12}). La segunda se demuestra de manera análoga,
observando que en este caso se tiene~$G(T)=\Aut(\cal{E}_T)$, donde
$\cal{E}_T$ es la \emph{extensión} trivial de~$\cal{O}_T$ por~$\cal{O}_T$ (y
donde, naturalmente, los automorfismos deben respetar la estructura de
extensión), y que los morfismos~$S'\to S$ que intervienen en~\Ref{XI.4.1}
son morfismos de descenso efectivo para la categoría fibrada de las
extensiones de~$\cal{O}_T$ por~$\cal{O}_T$ (como se deduce fácilmente de
\Ref{VIII.1.1}), y que tales extensiones son automáticamente localmente
triviales.

\begin{remarque}
\label{XI.5.2}
Obsérvese que el mismo tipo de demostración se aplica al grupo
simpléctico~$\Symp(2n)_S$, teniendo en cuenta que una forma alternada sobre un
módulo localmente isomorfo a~$\cal{O}_S^{2n}$ que sea «no degenerada», \ie que
defina un isomorfismo de este módulo con su dual, es localmente isomorfa a la
forma estándar. En cambio, el resultado para el grupo ortogonal es falso, ya
cuando~$S$ es el espectro de un cuerpo, pues puede haber formas
\marginpar{302}
cuadráticas sobre un cuerpo que no sean isomorfas a la forma estándar. Por
otra parte, se demuestra esencialmente en~\cite{XI.3} que los grupos~$\Gl$,
$\Symp$, $\GG_a$ y los que se descomponen por etapas en tales grupos son, con
pequeñas salvedades, los únicos para los cuales se tiene un resultado de
trivialidad local del tipo considerado aquí.
\end{remarque}

\begin{corollaire}
\label{XI.5.3}
Se tienen biyecciones canónicas
$$
\H^1(S,\GL(n)_S)\isomfrom \H^1(S,\GL(n,\cal{O}_S))\quoi,
$$
en particular
$$
\H^1(S,\GG_{m\,S})\isomfrom \H^1(S,\cal{O}_S^*)\quoi,
$$
y
$$
\H^1(S,\GG_{a\,S})\isomfrom \H^1(S,\cal{O}_S)\quoi,
$$
donde los segundos miembros denotan cohomologías del espacio topológico~$S$
con coeficientes en haces ordinarios.
\end{corollaire}

En particular, $\H^1(S,\GL(n)_S)$ se identifica con el conjunto de las
clases, salvo isomorfismo, de módulos localmente libres de rango~$n$ sobre~$S$,
y~$\H^1(S,\GG_{a\,S})$ se identifica con el conjunto de las clases de
extensiones del módulo~$\cal{O}_S$ por sí mismo.


\section[Aplicación a los recubrimientos principales]{Aplicación a los
recubrimientos principales: teorías de Kummer y de Artin--Schreier}
\label{XI.6}


\begin{proposition}
\label{XI.6.1}
Sean~$S$ un preesquema y~$n$ un entero~$>0$; sea
$u_n\colon\GG_{m\,S}\to\GG_{m\,S}$ el homomorfismo de elevación a la
potencia $n$-ésima, y sea~$\bbmu_{n\,S}$
\label{indnot:kb}\oldindexnot{$\bbmu_{n\,S}$|hyperpage}%
su núcleo. Entonces~$\bbmu_{n\,S}$ es finito y localmente libre de rango~$n$
sobre~$S$, y es étale sobre~$S$ si y solo si, para todo~$s\in S$, la
característica de~$s$ es prima con~$n$. La sucesión de homomorfismos
$$
0\to \bbmu_{n\,S}\to\GG_{m\,S}\lto{u_n}\GG_{m\,S}\to 0
$$
es exacta en el sentido de~\Ref{XI.4}. (La llamaremos la \emph{sucesión
exacta de Kummer}
\index{Kummer (sucesión exacta de)|hyperpage}%
sobre~$S$, relativa al entero~$n$).
\end{proposition}

Se tiene
\marginpar{303}
$$
\GG_m=\Spec\cal{O}_S[t,t^{-1}],
$$
y~$u_n$ corresponde al homomorfismo~$u_n$ de las
$\cal{O}_S$-álgebras afines, dado por
$$
u_n(t)=t^n,
$$
por otra parte, la sección unidad de~$\GG_{m\,S}$ corresponde al
homomorfismo de aumento de $\cal{O}_S$-álgebras dado por
$$
\varepsilon(t)=1,
$$
cuyo núcleo es, por tanto, el
\ifthenelse{\boolean{orig}}
{ideal}
{ideal}
principal~$(t-1)$. La imagen de este último por~$u_n$ es, pues, el ideal
principal~$(1-t^n)$, y se obtiene
$$
\bbmu_{n\,S}=\Spec\cal{O}_S[t]/(1-t^n),
$$
lo que muestra, en particular, que~$\bbmu_{n\,S}$ es finito sobre~$S$ y está
definido por un álgebra sobre~$S$ que es libre de rango~$n$, con base formada
por los~$t^i$ ($0\leq i\leq n-1$). Para que sea étale en~$s\in S$, es
necesario y suficiente que el álgebra reducida~$k[t]/(1-t^n)$, donde
$k=\kres(s)$, obtenida adjuntando formalmente las raíces $n$-ésimas de la
unidad a~$k$, sea separable sobre~$k$, \,\ie que las raíces de~$1-t^n$ en una
clausura algebraica de~$k$ sean todas distintas, lo cual equivale a que~$n$
sea primo con la característica. Por último, para demostrar que la sucesión
de homomorfismos de~\Ref{XI.6.1} es exacta, en virtud del
criterio~\Ref{XI.4.2} se reduce uno a probar que
\ifthenelse{\boolean{orig}}
{$v$}
{$u_n$}
es fielmente plano. Podemos suponer, evidentemente, que~$S$ es afín de
anillo~$A$, de modo que~$\GG_{m\,S}$ es afín de anillo
$B=A[t,t^{-1}]$; basta comprobar que~$u_n$ hace de~$B$ un módulo libre de
rango~$n$ sobre~$B$, o, lo que es lo mismo, que~$u_n$ es inyectivo y que
$A[t,t^{-1}]$ es un módulo libre de rango~$n$ sobre~$A[t^n,t^{-n}]$. En
efecto, se comprueba fácilmente que los~$t^i$ ($0\leq i\leq n-1$) forman
una base del primero sobre el segundo, lo que termina la demostración.

\begin{definition}
\label{XI.6.2}
Se
\marginpar{304}
llama a~$\bbmu_{n\,S}$ el \emph{grupo de Kummer de rango~$n$ sobre~$S$}, y
se llama \emph{recubrimiento principal de Kummer de rango
\ifthenelse{\boolean{orig}}
{$n\,S$}
{$n$ sobre~$S$}}
a todo fibrado homogéneo principal sobre~$S$ cuyo grupo es el grupo de
Kummer de rango~$n$.
\end{definition}

El conjunto de estos recubrimientos es un grupo, denotado por
$\H^1(S,\bbmu_{n\,S})$ o, simplemente, por~$\H^1(S,\bbmu_n)$. Obsérvese que
la formación del grupo de Kummer de rango~$n$ sobre~$S$ es compatible con la
extensión de la base, de modo que~$\bbmu_{n\,S}$ proviene por extensión de la
base del \emph{grupo de Kummer absoluto~$\bbmu_n$} sobre~$\Spec(\ZZ)$.

Denotemos por~$(\ZZ/n\ZZ)_S$ el $S$-grupo definido por el grupo finito
ordinario~$\ZZ/n\ZZ$. Si~$G$ es un $S$-grupo cualquiera, los homomorfismos
de $S$-grupos~$u$ de~$(\ZZ/n\ZZ)_S$ en~$G$ corresponden biyectivamente, y
de manera compatible con el cambio de base, a las secciones de~$G$ sobre~$S$
cuya potencia $n$-ésima es la sección unidad, haciendo corresponder a~$u$ la
imagen por~$u$ de la sección de~$(\ZZ/n\ZZ)_S$ sobre~$S$ definida por el
generador~$1\bmod n\ZZ$ de~$\ZZ/n\ZZ$. Dicho esto:

\begin{corollaire}
\label{XI.6.3}
Si~$\bbmu_{n\,S}$ es étale sobre~$S$, se obtiene así una correspondencia
biyectiva entre los isomorfismos de $S$-grupos
$(\ZZ/n\ZZ)_S\isomto\bbmu_{n\,S}$ y las secciones de~$\cal{O}_S$ que son
de orden exactamente~$n$ en cada componente conexa de~$S$ (una sección tal
se llamará «raíz primitiva $n$-ésima de la unidad sobre~$S$»). Por tanto,
para que~$\bbmu_{n\,S}$ sea isomorfo, como $S$-grupo, a~$(\ZZ/n\ZZ)_S$, es
necesario y suficiente que sea étale
\ifthenelse{\boolean{orig}}
{sobre~$S$}
{sobre~$S$,}
\ie que las características residuales de~$S$ sean primas con~$n$, y que
exista una raíz primitiva $n$-ésima de la unidad sobre~$S$.
\end{corollaire}

Esto explica el papel que desempeñan en la teoría clásica de Kummer la
hipótesis de que el cuerpo base (que desempeña el papel de~$S$) tenga
característica prima con~$n$ y contenga las raíces $n$-ésimas de la unidad,
y la elección de una raíz primitiva $n$-ésima de la unidad. Una vez que se
dispone del lenguaje de los esquemas, ya no hay por qué imponerse estas
hipótesis, y conviene razonar directamente sobre~$\bbmu_n$ en vez
de~$\ZZ/n\ZZ$. Así, la conjunción de~\Ref{XI.6.1}, \Ref{XI.4.5}
y~\Ref{XI.5.3} nos da la relación general siguiente entre la teoría de los
recubrimientos principales de Kummer y la de los grupos de Picard:

\begin{proposition}
\label{XI.6.4}
Sea
\marginpar{305}
$S$ un preesquema. Se tiene una sucesión exacta canónica
$$
\textstyle
0\to\H^0(S,\bbmu_n)\to\H^0(S,\cal{O}^*_S)\to
\H^0(S,\cal{O}^*_S)\to\H^1(S,\bbmu_n)\to
\H^1(S,\cal{O}^*_S)\to\H^1(S,\cal{O}^*_S),
$$
de donde, poniendo
\ifthenelse{\boolean{orig}}
{$\H^1(S,\cal{O}_S)=\Pic(S)$,}
{$\H^1(S,\cal{O}_S^*)=\Pic(S)$,}
y denotando, para todo grupo abeliano~$A$, por~${_n}A$ y~$A_n$ el núcleo y
el conúcleo de la multiplicación por~$n$ en~$A$, se obtiene la sucesión exacta
$$
0\to
\ifthenelse{\boolean{orig}}
{\H^0(S,\cal{O}_S)^*_n}
{\H^0(S,\cal{O}_S^*)_n}
\to\H^1(S,\bbmu_n)\to
{\sideset{_n}{}\Pic}(S)\to 0.
$$
\end{proposition}

Vamos a explicitar dos casos importantes en los que uno u otro término
extremo de esta sucesión exacta es nulo:

\begin{corollaire}
\label{XI.6.5}
Supongamos que~$\sideset{_n}{}\Pic(S)=0$ (por ejemplo, que~$S$ sea el
espectro de un anillo local o de un anillo factorial), y sea~$A$ el anillo
$\H^0(S,\cal{O}_S)$. Entonces se tiene un isomorfismo canónico
$$
\H^1(S,\bbmu_n)\isomto A^*/{A^*}^n.
$$
\end{corollaire}

Este es esencialmente el enunciado clásico de la teoría de Kummer cuando~$S$
es el espectro de un cuerpo.

\begin{corollaire}
\label{XI.6.6}
Supongamos que todo elemento de~$\H^0(S,\cal{O}_S)$ es una potencia
$n$-ésima; por ejemplo, que~$\H^0(S,\cal{O}_S)$ sea un
\ifthenelse{\boolean{orig}}
{compuesto}
{producto}
de cuerpos algebraicamente cerrados, o que~$S$ sea reducido y propio sobre
un cuerpo algebraicamente cerrado~$k$. Entonces se tiene un isomorfismo
canónico
$$
\H^1(S,\bbmu_n)\isomto\sideset{_n}{}\Pic(S).
$$
\end{corollaire}

En particular, cuando~$S$ es propio y conexo sobre un cuerpo algebraicamente
cerrado~$k$, esto relaciona el grupo fundamental de~$S$ con los puntos de
orden finito del esquema de Picard~$P$ de~$S$ sobre~$k$; así, para~$n$ primo
con la característica, se tiene un isomorfismo
$$
\Hom(\pi_1(S),\ZZ/n\ZZ)\simeq {_n P}(k)
$$
que se utiliza a menudo en geometría
\marginpar{306}
algebraica. Como aplicación, cuando la componente conexa~$P^0$ de~$P$ es un
esquema en grupos completo, de dimensión~$g$, se ve, utilizando los
resultados recordados en~\Ref{XI.2} y la finitud del grupo de torsión de
\ifthenelse{\boolean{orig}}
{Néron--Sévéri}
{Néron--Severi}
que, para todo número primo~$\ell$ distinto de la característica, la
componente $\ell$-primaria del grupo fundamental abelianizado~$\pi_1(S)$ es
un módulo de tipo finito y de rango~$2g$ sobre el anillo~$\ZZ_\ell$ de los
enteros $\ell$-ádicos (y, además, libre excepto, a lo sumo, para un número
finito de valores de~$\ell$). Como observó Serre, esto permite demostrar,
bajo ciertas condiciones, que cuando~$X$ es un esquema plano y proyectivo
sobre~$S$ conexo, los esquemas de Picard de las fibras de~$X$ tienen todos
la misma dimensión, aplicando el teorema de semicontinuidad
(\Ref{X.2.3}); el argumento de Serre se aplica siempre que exista el esquema
de Picard de~$X$ sobre~$S$ y los esquemas de Picard conexos de las fibras
de~$X$ sobre~$S$ sean esquemas en grupos propios, por ejemplo cuando las
fibras geométricas de~$X$ sobre~$S$ sean normales (manteniéndose~$X$ plano y proyectivo
sobre~$S$), en particular si~$X$ es liso y proyectivo sobre~$S$.

Sea ahora~$p$ un número primo, y supongamos que~$S$ es un preesquema de
característica~$p$, \,\ie tal que~$p\cdot\cal{O}_S=0$. Entonces el
homomorfismo de elevación a la potencia $p$-ésima en~$\cal{O}_S$ es aditivo,
y el morfismo correspondiente, obtenido reemplazando~$S$ por un~$T$ variable
sobre~$S$,
$$
F\colon\GG_{a\,S}\to\GG_{a\,S}
$$
es, por tanto, un homomorfismo de $S$-grupos, llamado el \emph{homomorfismo
de Frobenius}. (N.B. Un morfismo tal está definido para todo $S$-preesquema
$G$ que provenga por extensión de la base de un preesquema~$G_0$ sobre el
cuerpo primo~$\ZZ/p\ZZ$, y este morfismo es un homomorfismo de grupos si~$G_0$
es un preesquema en grupos). Pondremos
$$
\wp=\id-F\colon\GG_{a\,S}\to\GG_{a\,S}.
$$


Consideremos, por otra parte, el $S$-grupo~$(\ZZ/p\ZZ)_S$ definido por el
grupo finito ordinario~$\ZZ/p\ZZ$. Hemos dicho que, para todo $S$-grupo~$G$,
los homomorfismos de $S$-grupos de~$(\ZZ/p\ZZ)_S$ en~$G$ corresponden
biyectivamente a las secciones de~$G$ sobre~$S$ cuya potencia $p$-ésima es la
sección unidad. Cuando~$G=\GG_{a\,S}$, corresponden
\marginpar{307}
por tanto a secciones cualesquiera de~$G$ sobre~$S$. Tomando, en particular,
la sección de~$\GG_{a\,S}$ sobre~$S$ que corresponde a la sección unidad del
haz de anillos~$\cal{O}_S$, se obtiene un homomorfismo de $S$-grupos
$$
\ifthenelse{\boolean{orig}}
{i\colon(\ZZ/p\ZZ)\to\GG_{a\,S}}
{i\colon(\ZZ/p\ZZ)_S\to\GG_{a\,S}}
$$

\begin{proposition}
\label{XI.6.7}
La sucesión de homomorfismos de $S$-grupos
$$
\ifthenelse{\boolean{orig}}
{0\to(\ZZ/p\ZZ)_S\to\GG_{a\,S}\to G_{a\,S}\to 0}
{0\to(\ZZ/p\ZZ)_S\to\GG_{a\,S}\to \GG_{a\,S}\to 0}
$$
es exacta (en el sentido de~\Ref{XI.4}). (Se la llama la \emph{sucesión
exacta de Artin--Schreier} sobre~$S$).
\index{Artin--Schreier (sucesión exacta de)|hyperpage}%
\end{proposition}

Basta demostrarlo sobre el cuerpo primo~$k=\ZZ/p\ZZ$. Basta observar que el
homomorfismo~$\wp^*\colon k[t]\to k[t]$ definido por
$\wp^*(t)=t-t^p$ hace de~$k[t]$ un módulo libre de rango~$p$ sobre~$k[t]$;
con mayor precisión, $k[t]$ es un módulo libre sobre~$k[s]$, donde
$s=t-t^p$, con base formada por los~$t^i$ ($0\leq i\leq p-1$).

Utilizando~\Ref{XI.4.5} y~\Ref{XI.5.3}, se concluye:

\begin{proposition}
\label{XI.6.8}
Se tiene una sucesión exacta canónica
\begin{multline*}
0\to\H^0(S,\ZZ/p\ZZ)\to\H^0(S,\cal{O}_S)\to\H^0(S,\cal{O}_S)\to
\H^1(S,\ZZ/p\ZZ)\\
\to\H^1(S,\cal{O}_S)\to\H^1(S,\cal{O}_S),
\end{multline*}
de donde se obtiene una sucesión exacta
$$
0\to\H^0(S,\cal{O}_S)/\wp\H^0(S,\cal{O}_S)\to
\H^1(S,\ZZ/p\ZZ)\to\H^1(S,\cal{O}_S)^F\to 0,
$$
(donde el exponente~$F$ del último término denota el subgrupo de los
invariantes por el endomorfismo~$F$, igual al núcleo de~$\wp=\id-F$).
\end{proposition}

Explicitemos además dos casos extremos:

\begin{corollaire}
\label{XI.6.9}
Supongamos que~$\H^1(S,\cal{O}_S)^F=0$, por ejemplo, que~$S$ sea un esquema
afín. Entonces, poniendo~$A=\H^0(S,\cal{O}_S)$, se tiene un isomorfismo
canónico
$$
\H^1(S,\ZZ/p\ZZ)\isomto A/\wp A.
$$
\end{corollaire}

Esta es la \emph{teoría de Artin--Schreier} en su forma clásica, al menos
cuando
\ifthenelse{\boolean{orig}}
{$A$}
{$S$}
es el espectro de un cuerpo.

\begin{corollaire}
\label{XI.6.10}
Supongamos que se verifica la igualdad
\marginpar{308}%
$\wp\H^0(S,\cal{O}_S)\mkern1mu{=}\mkern1mu\H^0(S,\cal{O}_S)$; por
ejemplo, que~$\H^0(S,\cal{O}_S)$ sea un
\ifthenelse{\boolean{orig}}
{compuesto}
{producto}
de cuerpos algebraicamente cerrados, o que~$S$ sea propio sobre un cuerpo
algebraicamente cerrado. Entonces se tiene un isomorfismo canónico
$$
\H^1(S,\ZZ/p\ZZ)\isomto\H^1(S,\cal{O}_S)^F
$$
\end{corollaire}

\begin{remarques}
\label{XI.6.11}
El último enunciado se debe a
\ifthenelse{\boolean{orig}}{J.P\ptbl}{J-P\ptbl}%
Serre~\cite{XI.9}. También es posible desarrollar una teoría análoga para
el grupo estructural~$\ZZ/p^n\ZZ$, con~$n$ cualquiera, utilizando en lugar
de~$\GG_a$ el esquema en grupos de Witt~$\WW_n$, \cf \loccit Obsérvese que,
en característica~$p>0$, la teoría de Kummer ya no da información sobre los
recubrimientos principales de orden~$p$, pues~$\bbmu_p$ es entonces un grupo
«infinitesimal», \,\ie radicial sobre la base, y, por tanto, sin relación
directa con~$\ZZ/p\ZZ$; así, a primera vista, la teoría de estos recubrimientos
ya no puede abordarse (cuando, para fijar las ideas,~$S$ es un esquema propio
sobre un cuerpo algebraicamente cerrado) mediante la teoría del esquema de
Picard como en~\Ref{XI.6.6}. Sin embargo, si se recuerda que el espacio
tangente de Zariski en el origen de
\ifthenelse{\boolean{orig}}
{$\mathbf{Pic}_{S/K}$}
{$\mathbf{Pic}_{S/k}$}
\kern1pt\footnote{Para la definición de
\ifthenelse{\boolean{orig}}
{$\mathbf{Pic}_{S/K}$}
{$\mathbf{Pic}_{S/k}$},
\cf A. Grothendieck, Sém. Bourbaki~\No 232 (febrero de 1962).} se identifica
con~$\H^1(S,\cal{O}_S)$, se constata que \emph{el conocimiento del esquema en
grupos~${_p}\mathbf{Pic}_{S/k}$, núcleo de la multiplicación por~$p$ en
$\mathbf{Pic}_{S/k}$, implica el de~$\H^1(S,\ZZ/p\ZZ)$ tanto como el
de~$\H^1(S,\bbmu_p)$; obsérvese que implica también el
de~$\H^1(S,\bbalpha_p)$}, donde~$\bbalpha_p$ denota el esquema en grupos
infinitesimal sobre el cuerpo primo, núcleo de~$F\colon\GG_a\to\GG_a$ (que
también puede describirse como el espectro del álgebra envolvente restringida
de la $p$-álgebra de Lie trivial de dimensión~$1$): en efecto, la sucesión
exacta~\Ref{XI.4.5} da aquí
$$
\H^1(S,\bbalpha_p)\simeq\Ker
(F\colon\H^1(S,\cal{O}_S)\to\H^1(S,\cal{O}_S)),
$$
y, más generalmente, denotando por~$\bbalpha_{p^n}$ el núcleo en~$\GG_a$ del
$n$-ésimo iterado de~$F$, se tendrá
$$
\H^1(S,\bbalpha_{p^n})\simeq\Ker(F^n\colon\H^1(S,\cal{O}_S)\to
\H^1(S,\cal{O}_S)).
$$

De hecho, el conocimiento de~${_p}\mathbf{Pic}_{S/k}$ equivale al
de~$\H^1(S,G)$ para todo grupo algebraico
\marginpar{309}
conmutativo finito anulado por~$p$; más generalmente, el conocimiento
de~${_{p^n}}\mathbf{Pic}_{S/k}$ equivale al de~$\H^1(S,G)$ para todo grupo
algebraico conmutativo finito~$G$ anulado por~$p^n$, en virtud del teorema
siguiente, que engloba, en el caso considerado, tanto la teoría de Kummer
como
\ifthenelse{\boolean{orig}}
{la de Artin--Schreier}
{la de Artin--Schreier}
:

Sea~$G$ un grupo algebraico finito sobre~$k$, y sea
$D(G)=\SheafHom_{k\textup{-grupos}}(G,\GG_m)$ su \emph{dual de Cartier} (cuya
álgebra afín tiene como espacio subyacente el espacio vectorial dual del
álgebra afín de~$G$, \,\ie la hiperálgebra de~$G$ en el sentido de
Dieudonné--Cartier); entonces se tiene un isomorfismo canónico
\begin{equation}\label{eqXI.6.11}\tag{$*$}
\H^1(S,G)\simeq\Hom_{k\textup{-grupos}}(D(G),\mathbf{Pic}_{S/k}).
\end{equation}
(N.B. $S$ es un esquema propio sobre~$k$ algebraicamente cerrado, tal que
$\H^0(S,\cal{O}_S)=k$). Esta fórmula también puede expresarse diciendo que
el \og verdadero grupo fundamental\fg{} de~$S$ al que se aludía en
\Ref{XI.2}, una vez abelianizado, es isomorfo al límite proyectivo de
los~$D(P_i)$, donde~$P_i$ recorre los subgrupos algebraicos \emph{finitos}
de~$\mathbf{Pic}_{S/k}$, que se denotará por
\ifthenelse{\boolean{orig}}
{$T^\bbullet(\mathbf{Pic}_{X/k})$}
{$T^\bbullet(\mathbf{Pic}_{S/k})$}.
Cuando~$S$ es una variedad abeliana, se vio en~\Ref{XI.2.1} que este grupo
es también isomorfo al \og verdadero\fg{} módulo de Tate
$T_{\bbullet}(S)=\varprojlim\, {{_n}S}$, y el
isomorfismo~\eqref{eqXI.6.11} se escribe entonces de manera más llamativa
$$
\Ext^1(A,G)\simeq\Hom(D(G),B),
$$
siendo~$A$ una variedad abeliana,~$B$ su dual y~$G$ un grupo algebraico
finito sobre~$k$. Los resultados que acabamos de indicar pueden generalizarse,
además, al caso en que~$k$ se sustituye por un preesquema base cualquiera, y
a otros grupos de coeficientes~$G$ distintos de los grupos finitos.
\end{remarques}



\begin{sgabibliography}{10}{XI.7}
\marginpar{310}

\bibitem[1]{XI.1} A\ptbl \textsc{Grothendieck},
Sur quelques points d'algèbre homologique,
\ifthenelse{\boolean{orig}}
{Tohoku Math.\ Journal, pp\ptbl 119--221, Vol\ptbl 9, 1957.}
{T\^ohoku Math.~J.\ \textbf{9} (1957) p\ptbl 119--221.}

\bibitem[2]{XI.2} A\ptbl \textsc{Grothendieck}, A general theory of fibre spaces with
structure sheaf, University of Kansas, 1955.

\bibitem[3]{XI.3} A\ptbl \textsc{Grothendieck}, Torsion homologique et sections
rationnelles, Séminaire \hbox{Chevalley}, 16 de junio de 1958.

\bibitem[4]{XI.4} S\ptbl \textsc{Lang}, Abelian varieties, Interscience tracts in pure
and applied mathematics, \No 7, New York.

\bibitem[5]{XI.5}
\ifthenelse{\boolean{orig}}{J.P\ptbl}{J-P\ptbl}%
\textsc{Serre}, Groupes algébriques et corps de classes,
Actualités Scient. et Ind.\ \No 1264, Hermann, Paris, 1959.

\bibitem[6]{XI.6}
\ifthenelse{\boolean{orig}}{J.P\ptbl}{J-P\ptbl}%
\textsc{Serre}, Groupes proalgébriques,
\ifthenelse{\boolean{orig}}
{Publications Mathématiques de l'IHES \No 7, pp\ptbl 1--67, 1960.}
{Publications Mathématiques de l'IHES \textbf{7} (1960), p\ptbl 1--67.}

\bibitem[7]{XI.7}
\ifthenelse{\boolean{orig}}{J.P\ptbl}{J-P\ptbl}%
\textsc{Serre}, Espaces fibrés algébriques,
Séminaire Chevalley, 21 de abril de 1958.

\bibitem[8]{XI.8}
\ifthenelse{\boolean{orig}}{J.P\ptbl}{J-P\ptbl}%
\textsc{Serre}, Quelques propriétés des
variétés abéliennes en
\ifthenelse{\boolean{orig}}
{car.~$p$, American Journal of Math., vol\ptbl 80, pp\ptbl 715--739, 1958.}
{caractéristique~$p$, Amer.\ J.\ Math.\ \textbf{80} (1958), p\ptbl 715--739.}

\bibitem[9]{XI.9}
\ifthenelse{\boolean{orig}}{J.P\ptbl}{J-P\ptbl}%
\textsc{Serre}, Sur la topologie des variétés
algébriques en
\ifthenelse{\boolean{orig}}
{car.~$p$}
{caractéristique~$p$},
\ifthenelse{\boolean{orig}}
{Symposium Internacional de Topologia Algebraica}
{Symposium Internacional de Topología Algebraica},
1958.

\bibitem[10]{XI.10}
\ifthenelse{\boolean{orig}}{J.P\ptbl}{J-P\ptbl}%
\textsc{Serre}, On the fundamental group of a unirational
variety,
\ifthenelse{\boolean{orig}}
{Journal London Math.\ Society vol\ptbl 34, pp\ptbl 481--484, 1959.}
{J.~London Math.\ Soc.\ \textbf{34} (1959), p\ptbl 481--484.}
\end{sgabibliography}


\chapterspace{-4}
\chapter{Geometría algebraica y~geometría~analítica}
\label{XII}
\marginpar{311}

\begin{center}
{Sra. M.\ \textsc{Raynaud}\footnote{A partir de notas inéditas de
A\ptbl\ Grothendieck.}}
\end{center}
\vspace*{1cm}

Siguiendo el procedimiento de~\cite{XII.10}, a todo esquema~$X$ localmente
de tipo finito sobre el cuerpo de los números complejos~$\CC$ se le asocia
un espacio analítico~$X^\an$
\label{indnot:lb}\oldindexnot{$X^\an$|hyperpage}%
cuyo conjunto subyacente es~$X(\CC)$.

En los números~\Ref{XII.2} y~\Ref{XII.3} de esta exposición damos un
\og diccionario\fg entre las propiedades usuales de~$X$ y de~$X^\an$, y
entre las propiedades de un morfismo $f\colon X\to Y$ y las del morfismo
asociado $f^\an\colon X^\an\to Y^\an$.
\label{indnot:lc}\oldindexnot{$f^\an$|hyperpage}%

Demostramos después que los teoremas de comparación entre haces coherentes
sobre~$X$ y sobre~$X^\an$, establecidos en~\cite[\No 12]{XII.10} para una
variedad proyectiva, siguen siendo válidos cuando~$X$ es un esquema propio.

Por último, en el número~\Ref{XII.5} demostramos la equivalencia entre la
categoría de los recubrimientos étales finitos de~$X$ y la categoría de los
recubrimientos étales finitos de~$X^\an$. Como añadido para el lector, damos
una nueva demostración del teorema de Grauert--Remmert~\cite{XII.6}, que
utiliza la resolución de singularidades~\cite{XII.8}.

\section{Espacio analítico asociado a un esquema}
\label{XII.1}
\marginpar{312}

Sea~$X$ un esquema localmente de tipo finito sobre~$\CC$. Sea~$\Phi$ el
funtor de la categoría de los espacios analíticos~\cite[\No 9]{XII.4} en la
categoría de los conjuntos que asocia a un espacio analítico~$\othercal{X}$
el conjunto de los morfismos de espacios anillados en $\CC$-álgebras
$\Hom_\CC(\othercal{X},X)$. Se tiene el teorema siguiente:

\begin{theoremedefinition}
\label{XII.1.1}
El funtor~$\Phi$ es representable por un espacio analítico~$X^\an$ y un
morfismo $\varphi\colon X^\an\to X$. Se dice que~$X^\an$ es el espacio
analítico asociado a~$X$.
\index{espacio analítico asociado|hyperpage}%

Si~$|X^\an|$ es el conjunto subyacente a~$X^\an$, entonces~$\varphi$ induce
una biyección de~$|X^\an|$ sobre el conjunto~$X(\CC)$ de los puntos de~$X$
con valores en~$\CC$. Además, para cada punto~$x$ de~$X^\an$, el morfismo
$$
\varphi_x\colon\cal{O}_{X,\varphi(x)}\to\cal{O}_{X^\an,x},
$$
que es necesariamente local, induce al pasar a las completaciones un
isomorfismo
$$
\widehat{\varphi}_x\colon\widehat{\cal{O}}_{X,\varphi(x)}\isomto
\widehat{\cal{O}}_{X^\an,x}
\ifthenelse{\boolean{orig}}
{,}
{.}%
$$
En particular, el morfismo~$\varphi$ es plano.
\end{theoremedefinition}

Obsérvese que el hecho de que~$\varphi$ induzca una biyección de~$X^\an$
sobre~$X(\CC)$ resulta de la propiedad universal de~$X^\an$. Por otra parte,
se tienen las afirmaciones siguientes:

\begin{enumerate}
\item[a)] Si el teorema es cierto para un esquema~$Y$, también lo es para
todo subesquema~$X$ de~$Y$. Supongamos primero que~$X$ es un subesquema
abierto de~$Y$; si $\psi\colon Y^\an\to Y$ es el morfismo canónico,
$\psi^{-1}(X)$ es un abierto
\marginpar{313}
de~$Y^\an$
\ifthenelse{\boolean{orig}}
{al que se dota}
{al que se dota}%
~de la estructura de espacio analítico inducida por la de~$Y^\an$. Como todo
morfismo de un espacio analítico~$\othercal{X}$ en~$X$ se factoriza a través
de~$Y^\an$ por la propiedad universal de este último, y por tanto a través
de~$X^\an$, que es el producto fibrado~$Y^\an\times_YX$, resulta que~$X^\an$
es el espacio analítico asociado a~$X$. Por último, la afirmación relativa a
los~$\varphi_x$ es evidente.

Basta ahora considerar el caso en que~$X$ es un subesquema cerrado de~$Y$.
Sea~$I$ el $\cal{O}_Y$-ideal coherente que define~$X$; entonces
$I\cdot\cal{O}_{Y^\an}$ es un haz coherente de ideales de~$\cal{O}_{Y^\an}$
que define un subespacio analítico cerrado~$X^\an$ de~$Y^\an$; como en el
caso de un subesquema abierto, se ve que~$X^\an$ es el espacio analítico
asociado a~$X$. Sea $\varphi\colon X^\an\to X$ el morfismo canónico. Para
todo punto~$x$ de~$X^\an$, el morfismo~$\varphi_x$ no es sino el morfismo
$$
\cal{O}_{Y,\psi(x)}/I_{\psi(x)}\to
\cal{O}_{Y^\an,x}/I_{\psi(x)}\cdot\cal{O}_{Y^\an,x}
$$
inducido por~$\psi_x$; su completación
$$
\widehat{\varphi}_x\colon
\widehat{\cal{O}}_{Y,\psi(x)}/I_{\psi(x)}\cdot\widehat{\cal{O}}_{Y,\psi(x)}
\to\widehat{\cal{O}}_{Y^\an,x}/I_{\psi(x)}\cdot\widehat{\cal{O}}_{Y^\an,x}
$$
es un isomorfismo porque~$\widehat\psi_x$ lo es, lo que demuestra~a).

\item[b)] Si se tienen dos $\CC$-esquemas~$X_1$ y~$X_2$ tales que existen
$X_1^\an$ y~$X_2^\an$, entonces también existe~$(X_1\times X_2)^\an$. Sean,
en efecto, $\varphi_1\colon X_1^\an\to X_1$ y
$\varphi_2\colon X_2^\an\to X_2$ los morfismos canónicos, y sean~$p_1$
y~$p_2$ las dos proyecciones de~$X_1^\an\times X_2^\an$. Se deduce
formalmente de EGA~I~1.8.1 que~$X_1\times X_2$ es el producto de~$X_1$ y~$X_2$
en la categoría de los espacios localmente anillados; de ello resulta que
los morfismos~$\varphi_1\cdot p_1$ y~$\varphi_2\cdot p_2$ definen un
\marginpar{314}
morfismo $\varphi\colon X_1^\an\times X_2^\an\to X_1\times X_2$, y que el
par~$(X_1^\an\times X_2^\an,\varphi)$ representa el funtor
$\othercal{X}\mto\Hom_\CC(\othercal{X},X_1\times X_2)$.

\item[c)] Si~$\othercal{E}^1$ denota el espacio afín de dimensión~1, es
decir, el espacio topológico~$\CC$ provisto del haz de funciones holomorfas,
el funtor%
\ifthenelse{\boolean{orig}}
{}
{{\renewcommand{\thefootnote}{*}\addtocounter{footnote}{-1}\footnote{\lcrochetbf Añadido en 2003: $E^1_\CC$ denota la recta afín (algebraica) sobre~$\CC$.\rcrochetbf}}}
$\othercal{X}\mto\Hom_\CC(\othercal{X},E^1_\CC)$ es representable por
$\othercal{E}^1$, siendo el morfismo canónico
$\varphi\colon\othercal{E}^1\to E^1_\CC$ el morfismo evidente. En efecto,
darse un morfismo de un espacio analítico~$\othercal{X}$ en~$E^1_\CC$
equivale a darse un elemento de~$\Gamma(\othercal{X},\cal{O}_\othercal{X})$,
lo que equivale también a darse un morfismo de~$\othercal{X}$
en~$\othercal{E}^1$. Se tiene evidentemente una biyección
$|\othercal{E}^1|\simeq E^1(\CC)$ y, para cada punto
$x\in\othercal{E}^1$, el morfismo~$\widehat\varphi_x$ no es sino el
morfismo identidad de un anillo de series formales en una variable
sobre~$\CC$.
\end{enumerate}

\ifthenelse{\boolean{orig}}{}
{\enlargethispage{.5cm}}%
De b) y c) se deduce que el teorema es cierto para el espacio afín
$E^n_\CC$, $n\geq0$. Utilizando~a), se ve que también lo es para todo esquema
afín~$X$ localmente de tipo finito sobre~$\CC$. Si ya no se supone que~$X$
sea afín y si~$(X_i)$ es un recubrimiento de~$X$ por abiertos afines, se
deduce de la propiedad universal y de~a) que los~$X_i^\an$ se pegan y definen
así el espacio analítico~$X^\an$ asociado a~$X$.

\setcounter{subsection}{1}% Preserve the source-visible paragraph number XII.1.2.
\subsection{}
\label{XII.1.2}
Sea $f\colon X\to Y$ un morfismo de $\CC$-esquemas localmente de tipo finito.
Si $\varphi\colon X^\an\to X$ y $\psi\colon Y^\an\to Y$ son los morfismos
canónicos, de la propiedad universal de~$Y^\an$ resulta que existe un único
morfismo $f^\an\colon X^\an\to Y^\an$ tal que el diagrama
$$
\xymatrix{
X^\an \ar[r]\ar[d]_-{f^\an}& X\ar[d]^-f\\ Y^\an \ar[r]& Y
}
$$
sea
\marginpar{315}
conmutativo. Se ha definido, pues, un funtor~$\Phi$ de la categoría de los
$\CC$-esquemas localmente de tipo finito en la categoría de los espacios
analíticos.

El funtor~$\Phi$ conmuta con los límites proyectivos finitos. Basta, en
efecto, ver que~$\Phi$ conmuta con los productos fibrados. Ahora bien, si
$X$, $Y$ y~$Z$ son esquemas localmente de tipo finito sobre~$\CC$, del hecho
de que~$X\times_ZY$ sea el producto fibrado de~$X$ y~$Y$ sobre~$Z$ en la
categoría de los espacios localmente anillados resulta que
$X^\an\times_{Z^\an}Y^\an$ satisface la propiedad universal que caracteriza
a~$(X\times_ZY)^\an$.

\subsection{}
\label{XII.1.3}
Sean~$X$ un $\CC$-esquema localmente de tipo finito, $X^\an$ el espacio
analítico asociado y $\varphi\colon X^\an\to X$ el morfismo canónico. Si~$F$
es un $\cal{O}_X$-módulo, la imagen inversa
$\varphi^*F=F^\an$ es un haz de módulos sobre~$\cal{O}_{X^\an}$. Se define
así un funtor de la categoría de los $\cal{O}_X$-módulos en la categoría de
los módulos sobre~$X^\an$. Este funtor conmuta con los límites inductivos
(EGA~0~4.3.2). Como el haz~$\cal{O}_{X^\an}$ es coherente
\cite[\no18~\S 2~th\ptbl 2]{XII.4}, transforma haces coherentes en haces
coherentes (EGA~0~5.3.11). Se tiene además:

\begingroup
\setcounter{proposition}{0}% The source numbers this statement inside paragraph XII.1.3.
\renewcommand{\theproposition}{\thesubsection.\arabic{proposition}}
\renewcommand{\theHproposition}{\theHsubsection.\arabic{proposition}}
\begin{subproposition}
\label{XII.1.3.1}
El funtor que a un
\ifthenelse{\boolean{orig}}
{$\cal{O}_{X^\an}$-módulo}
{$\cal{O}_X$-módulo}%
~$F$ asocia su imagen inversa~$F^\an$ sobre~$X^\an$ es exacto, fiel y
conservativo.
\end{subproposition}
\endgroup

La exactitud resulta de que el morfismo $\varphi\colon X^\an\to X$ es plano
\ifthenelse{\boolean{orig}}
{\eqref{XIII.1.1}}
{\eqref{XII.1.1}}.%
\ Demostremos que el funtor~$F\mto F^\an$ es fiel. Habida cuenta de la
exactitud, basta demostrar que, si~$F^\an$ es nulo, también lo es~$F$. Ahora
bien, para todo punto~$x$ de~$X^\an$, se tiene entonces
$F_{\varphi(x)}\otimes_{\cal{O}_{X,\varphi(x)}}\cal{O}_{X^\an,x}=0$. Como el
morfismo $\cal{O}_{X,\varphi(x)}\to\cal{O}_{X^\an,x}$ es fielmente plano, se
tiene $F_{\varphi(x)}=0$ para todo punto cerrado~$\varphi(x)$ de~$X$ y, como
$X$ es de Jacobson (EGA~IV~10.4.8), esto implica que~$F$ es nulo.

El
\marginpar{316}
hecho de que el funtor~$F\mto F^\an$ sea conservativo resulta formalmente de
la exactitud y la fidelidad.


\section[Comparación de las propiedades de un esquema]{Comparación de las propiedades de un esquema y del espacio
analítico \hbox{asociado}}
\label{XII.2}

\begin{proposition}
\label{XII.2.1}
Sean~$X$ un $\CC$-esquema localmente de tipo finito, $X^{\an}$ el espacio
analítico asociado y~$n$ un entero. Consideremos la propiedad~$P$ de ser
\begin{itemize}
\item[(i)] no vacío
\item[(i')] discreto
\item[(ii)] de Cohen--Macaulay
\item[(iii)] $(\mathrm{S}_n)$
\item[(iv)] regular
\item[(v)] $(\mathrm{R}_n)$
\item[(vi)] normal
\item[(vii)] reducido
\item[(viii)] de dimensión~$n$.
\end{itemize}
Entonces, para que~$X$ posea la propiedad~$P$, es necesario y suficiente que
también la posea~$X^{\an}$.
\end{proposition}

Sea $\varphi\colon X^{\an}\to X$ el morfismo canónico. La afirmación~(i)
resulta de la igualdad $|X^{\an}|=X(\CC)$
\ifthenelse{\boolean{orig}}
{\eqref{XIII.1.1}}
{\eqref{XII.1.1}}%
~y del hecho de que~$X$ es de Jacobson (EGA~IV~10.4.8). Decir que~$X$
(\resp $X^{\an}$) es discreto equivale a decir que $\dim X=0$
(\resp $\dim X^{\an}=0$ según~\cite[\no19~\S 4~cor\ptbl 6]{XII.4}); por
tanto, (i') resulta de~(viii).

Sea~$P$ una de las propiedades (ii) a~(vii). Para que~$X$ posea la
propiedad~$P$, es necesario y suficiente que~$P$ se verifique en cada punto
cerrado de~$X$; en efecto, como~$X$ es excelente (EGA~IV~7.8.6~(iii)), el
conjunto de los puntos
\marginpar{317}
en los que~$X$ verifica~$P$ es abierto (\loccit) y, si este abierto contiene
todos los puntos cerrados, es igual a todo~$X$. Decir que~$X$
(\resp $X^{\an}$) posee la propiedad~$P$ equivale, pues, a decir que, para
todo punto~$x$ de~$X^{\an}$, el anillo local
$\cal{O}_{X,\varphi(x)}$ (\resp $\cal{O}_{X^{\an},x}$) posee la propiedad~$P$.
Como el hecho de que un anillo local excelente posea la propiedad~$P$ puede
comprobarse tras pasar a la completación, la proposición resulta de los
isomorfismos
$\widehat{\cal{O}}_{X,\varphi(x)}\isomto\widehat{\cal{O}}_{X^{\an},x}$ en
los casos (ii) a~(vii). Lo mismo sucede en el caso~(viii), teniendo en cuenta
las relaciones
$$
\dim X=\sup_x~\dim \cal{O}_{X,\varphi(x)}\qquad \dim
X^{\an}=\sup_x~\dim \cal{O}_{X^{\an},x}
$$
donde $x\in X^{\an}$. Esto concluye la demostración.

\begin{proposition}
\label{XII.2.2}
Sean~$X$ un $\CC$-esquema localmente de tipo finito,
$\varphi\colon X^{\an}\to X$ el morfismo canónico y~$T$ un subconjunto
localmente constructible de~$X$. Entonces se tiene la relación
$$
\varphi^{-1}(\overline{T})=\overline{\varphi^{-1}(T)}.
$$
\end{proposition}

Podemos suponer que~$T$ es un abierto denso de~$X$. Sea~$H$ el subesquema
cerrado reducido de~$X$ cuyo espacio subyacente es $X-T$; el espacio
asociado~$H^{\an}$ es un subespacio analítico cerrado de~$X^{\an}$ cuyo
espacio subyacente es $X^{\an}-\varphi^{-1}(T)$. Debemos demostrar que todo
punto~$x$ de~$H^{\an}$ pertenece a~$\overline{\varphi^{-1}(T)}$. Ahora bien,
en tal punto~$x$, el germen de espacio analítico~$(X^{\an},x)$ contiene el
subgermen~$(H^{\an},x)$, y este está definido por un
\ifthenelse{\boolean{orig}}
{Ideal}
{ideal}%
~no nilpotente de~$\cal{O}_{X^{\an},x}$. Del teorema de los ceros
\cite[\no19~\S 4~cor\ptbl 3]{XII.4} resulta entonces que todo entorno abierto
de~$x$ contiene puntos de~$X^{\an}$ que no pertenecen a~$H^{\an}$, lo cual
demuestra que $x\in\overline{\varphi^{-1}(T)}$.

\begin{corollaire}
\label{XII.2.3}
Sean
\marginpar{318}
$X$ un $\CC$-esquema localmente de tipo finito,
$\varphi\colon X^\an\to X$ el morfismo canónico y~$T$ un subconjunto
localmente constructible de~$X$. Para que~$T$ sea abierto (\resp cerrado,
\resp denso), es necesario y suficiente que lo mismo ocurra con
$\varphi^{-1}(T)$.
\end{corollaire}

El corolario resulta de~\Ref{XII.2.2} y del hecho de que, como~$X$ es un
esquema de Jacobson (EGA~IV~10.4.8), dos subconjuntos localmente
constructibles de~$X$ que tienen la misma traza sobre el subconjunto muy
denso~$X(\CC)$ son iguales.

\begin{proposition}
\label{XII.2.4}
Sea~$X$ un $\CC$-esquema localmente de tipo finito. Para que~$X$ sea conexo
(\resp irreducible), es necesario y suficiente que también lo sea~$X^\an$.
\end{proposition}

Supongamos que~$X^\an$ es conexo (\resp irreducible). La imagen~$X(\CC)$
de~$X^\an$ en~$X$ es entonces conexa (\resp irreducible). De ello resulta
que~$X$ es conexo (\resp irreducible), pues los subconjuntos cerrados de~$X$
y de~$X(\CC)$ se corresponden biyectivamente (EGA~IV~10.1.2).

Recíprocamente, supongamos que~$X$ es conexo (\resp irreducible) y demostremos
que lo mismo ocurre con~$X^\an$. Podemos limitarnos al caso en que~$X$ es
irreducible. Supongamos, en efecto, que~$X$ es conexo. Dado un punto~$x$
de~$X$, el conjunto de los puntos $y\in X$ para los cuales existe una
sucesión finita de subesquemas cerrados irreducibles $X_1,\dots,X_n$ de~$X$,
con $x\in X_1$, $y\in X_n$ y $X_i\cap X_{i+1}\neq\emptyset$ para
$1\leq i\leq n-1$, es a la vez abierto y cerrado y, por tanto, igual a todo
~$X$. Para una sucesión $X_1,\dots,X_n$ como la precedente, se tiene también
$X_{i}^\an\cap X_{i+1}^\an\neq\emptyset$ para $1\leq i\leq n-1$; si se supone
demostrado que los~$X_{i}^\an$ son conexos, lo mismo ocurre entonces
con~$X^\an$.

Supondremos
\marginpar{319}
en adelante que~$X$ es irreducible. Podemos suponer además que~$X$ es afín.
En efecto, si $(U_i)_{i\in I}$ es un recubrimiento de~$X$ por abiertos
afines, dos cualesquiera de estos abiertos tienen intersección no vacía y,
por tanto, la misma propiedad se verifica para el recubrimiento
$(U_{i}^\an)_{i\in I}$ de~$X^\an$; si se supone demostrado que los~$U_{i}^\an$
son irreducibles, lo mismo ocurre entonces con~$X^\an$.

Podemos suponer además que~$X$ es normal. Sea, en efecto, $\tilde X$ la
normalización de~$X$; como el morfismo $\tilde X\to X$ es sobreyectivo,
también lo es $\tilde X^\an\to X^\an$, lo cual demuestra que, si
$\tilde X^\an$ es irreducible, también lo es~$X^\an$.

Supondremos en adelante que~$X$ es afín y normal. Como los anillos locales
de~$X^\an$ son íntegros, decir que~$X^\an$ es irreducible equivale a decir
que es conexo. En efecto, si~$\cal{F}$ es un subconjunto analítico cerrado
de~$X^\an$, el conjunto de los puntos~$x$ de~$X^\an$ en los que
$\codim_x(\cal{F},X^\an)=0$ es un subconjunto analítico cerrado de~$X^\an$
\cite[\no20 A cor\ptbl 1]{XII.4} que también es abierto; si~$X^\an$ es
conexo, esto demuestra que, si $\cal{F}\neq X^\an$, entonces~$\cal{F}$ es
nunca denso y, por tanto, que~$X^\an$ es irreducible. Quedamos así reducidos
a demostrar que~$X^\an$ es conexo.

Sea
$$
i\colon X\to P
$$
una compactificación de~$X$, donde~$P$ es un $\CC$-esquema proyectivo normal
e~$i$ una inmersión abierta dominante. De~\cite[\no12~th\ptbl 1]{XII.10}
resulta entonces que~$P^\an$ es conexo. Como~$X^\an$ se obtiene quitando
de~$P^\an$ un subconjunto analítico cerrado nunca denso, de~\Ref{XII.2.5}
más abajo resulta que~$X^\an$ también es conexo.

\begin{lemme}
\label{XII.2.5}
Sean
\marginpar{320}
$\cal{P}$ un espacio analítico normal conexo y~$\cal{Y}$ un subconjunto
analítico cerrado nunca denso. Entonces $\cal{X}=\cal{P}-\cal{Y}$ es conexo.
\end{lemme}

Cuando~$\cal{Y}$ es de codimensión~$\geq2$, la afirmación resulta
de~\cite[\no3~prop\ptbl 4]{XII.11}. En el caso general, quitando si es
necesario de~$\cal{P}$ un subconjunto analítico cerrado de
codimensión~$\geq2$, podemos suponer que~$\cal{P}$ y~$\cal{Y}$ (considerado
como subespacio analítico reducido de~$\cal{P}$) son regulares. Según el
teorema de las funciones implícitas, todo punto~$y$ de~$\cal{Y}$ posee un
entorno~$\cal{U}$ isomorfo a una bola de un espacio afín~$\cal{E}^{n}$, de modo
que~$\cal{U}\cap\cal{Y}$ esté definido por la anulación de cierto número de
funciones coordenadas. Esto demuestra que
$\cal{U}-\cal{U}\cap\cal{Y}$ es conexo y, por tanto, también lo
es~$\cal{X}$.

\begin{corollaire}
\label{XII.2.6}
Sea~$X$ un $\CC$-esquema localmente de tipo finito; el morfismo
$$
\pi_0(X^{\an})\to\pi_0(X)
$$
inducido por el morfismo canónico $X^\an\to X$ es biyectivo.
\end{corollaire}


\section{Comparación de las propiedades de los morfismos}
\label{XII.3}

\begin{proposition}
\label{XII.3.1}
Sean $f\colon X\to Y$ un morfismo de $\CC$-esquemas localmente de tipo
finito y $f^\an\colon X^\an\to Y^\an$ el morfismo deducido de~$f$ sobre los
espacios analíticos asociados. Sea~$P$ la propiedad de ser
\begin{enumerate}
\item[(i)] plano
\item[(ii)] neto (\ie no ramificado)
\item[(iii)] étale
\item[(iv)] liso
\item[(v)] normal
\item[(vi)]
\marginpar{321}
reducido
\item[(vii)] inyectivo
\item[(viii)] separado
\item[(ix)] un isomorfismo
\item[(x)] un monomorfismo
\item[(xi)] una inmersión abierta.
\end{enumerate}
Entonces, para que~$f$ posea la propiedad~$P$, es necesario y suficiente que
también la posea~$f^\an$.
\end{proposition}

Denotemos por $\varphi\colon X^\an\to X$ y $\psi\colon Y^\an\to Y$ los
morfismos canónicos. Sean~$x$ un punto de~$X^\an$ e $y=f^\an(x)$. Los
morfismos $\cal{O}_{Y^\an,y}\to{\mathcal O}_{X^\an,x}$ y
$\cal{O}_{Y,\psi(y)}\to\cal{O}_{X,\varphi(x)}$,
\ifthenelse{\boolean{orig}}
{inducidos por~$f$ y~$f^\an$, respectivamente,}
{inducidos por~$f^\an$ y~$f$, respectivamente,}%
dan el mismo morfismo al pasar a las completaciones~\eqref{XII.1.1}. Según
\cite[ch\ptbl 3 \S 5 prop\ptbl 4]{XII.2} (\resp EGA~IV~17.4.4), equivale,
pues, decir que~$f^\an$ verifica la propiedad~(i) (\resp (ii)) o decir que~$f$
verifica~(i) (\resp (ii)) en cada punto cerrado de~$X$. Como el conjunto de
los puntos de~$X$ en los que se verifica~(i) (\resp (ii)) es abierto
(EGA~IV~11.1.1 e I~3.3), esto demuestra~(i) y~(ii) y, por tanto,
también~(iii).

\ifthenelse{\boolean{orig}}{}
{\enlargethispage{.5cm}}%
Sea~$P$ la propiedad~(iv) (\resp (v), \resp (vi)). Teniendo en cuenta
\Ref{XII.2.1}
\ifthenelse{\boolean{orig}}
{((v), (vi), (vii)),}
{((iv), (vi), (vii)),}%
~equivale decir que las fibras geométricas de~$f^\an$ en los distintos
puntos~$y$ de~$Y^\an$ son regulares (\resp normales, \resp reducidas) o que
lo mismo ocurre con las fibras geométricas de~$f$ en los distintos puntos
cerrados~$\psi(y)$ de~$Y$. Los casos~(iv) (\resp (v), \resp (vi)) resultan
entonces de~(i) y del hecho de que el conjunto de los puntos de~$Y$ en los
que las fibras geométricas de~$f$ son regulares es abierto
(EGA~IV~12.1.7).

(vii). Si~$f$ es inyectivo, también lo es~$f^\an$. Recíprocamente,
supongamos que~$f^\an$ es inyectivo y demostremos que~$f$ también lo es.
Podemos suponer
\marginpar{322}
que~$f$ es de tipo finito. Como el morfismo~$f^\an$ es inyectivo, las fibras
de~$f$ en los puntos cerrados de~$Y$ son radiciales; como el conjunto de los
puntos de~$Y$ cuya fibra es radicial es localmente constructible
(EGA~IV~9.6.1) y como~$Y$ es un esquema de Jacobson, todas las fibras de~$f$
son radiciales y, por tanto, es inyectivo.

(viii). Sean $\Delta\colon X\to X\times_YX$ y
$\Delta^\an\colon X^\an\to X^\an\times_{Y^\an}X^\an$ las inmersiones
diagonales, y sea
$\Theta\colon X^\an\times_{Y^\an}X^\an\to X\times_YX$ el morfismo canónico.
En virtud de~\Ref{XII.2.3}, equivale decir que~$\Delta(X)$ es cerrado en
$X\times_YX$ o que~$\Delta^\an(X^\an)$ es cerrado en
$X^\an\times_{Y^\an}X^\an$.

Como una inmersión abierta no es sino un morfismo étale inyectivo
(EGA~IV~17.9.1 y~\cite[\no13 \S 1]{XII.4}), (xi) resulta de~(iii) y~(vii).
Como un isomorfismo es lo mismo que una inmersión abierta sobreyectiva, (ix)
resulta de~(xi) y de (\Ref{XII.3.2}~(i)) más abajo. Decir que~$f$ es un
monomorfismo equivale a decir que
\ifthenelse{\boolean{orig}}
{la morfismo}
{el morfismo}%
~diagonal $\Delta\colon X\to X\times_YX$ es un isomorfismo; por tanto, (x)
resulta de~(ix).

\begin{proposition}
\label{XII.3.2}
Sean~$X$ e~$Y$ dos $\CC$-esquemas localmente de tipo finito,
$f\colon X\to Y$ un morfismo de tipo finito y
$f^\an\colon X^\an\to Y^\an$ el morfismo deducido de~$f$ sobre los espacios
analíticos asociados. Sea~$P$ la propiedad de ser
\begin{enumerate}
\item[(i)] sobreyectivo
\item[(ii)] dominante
\item[(iii)] una inmersión cerrada
\item[(iv)] una inmersión
\item[(v)] propio\footnote{Diremos que un morfismo de espacios analíticos es
propio si lo es en el sentido de~\cite[ch\ptbl 1 \S 10 \no1]{XII.1} y si es
separado.}
\item[(vi)]
\marginpar{323}
finito.
\end{enumerate}
Entonces, para que~$f$ posea la propiedad~$P$, es necesario y suficiente que
también la posea~$f^\an$.
\end{proposition}

Sean $\varphi\colon X^\an\to X$ y $\psi\colon Y^\an\to Y$ los morfismos
canónicos.

(i). Si~$f$ es sobreyectivo, para todo punto~$y$ de~$Y^\an$,
$f^{-1}(\psi(y))$ es un subconjunto cerrado no vacío de~$X$; contiene, pues,
al menos un punto cerrado, lo cual demuestra que~$f^\an$ es sobreyectivo.
Recíprocamente, si~$f^\an$ es sobreyectivo, $f(X)$ es un subconjunto
localmente constructible de~$Y$ (EGA~IV~1.8.4) que contiene todos los puntos
cerrados de~$Y$; por tanto, $f(X)=Y$.

(ii) resulta de~\Ref{XII.2.2}.

(iii). Si~$f$ es una inmersión cerrada, también lo es~$f^\an$ según
\Ref{XII.1.1}~a). Recíprocamente, si~$f^\an$ es una inmersión cerrada, también
lo es~$f$ según~\Ref{XII.3.1}~(x) y~\Ref{XII.3.2}~(v), pues esto equivale a
decir que~$f$ es un monomorfismo propio (EGA~IV~8.11.5).

(iv). Es claro que, si~$f$ es una inmersión, también lo es~$f^\an$.
Recíprocamente, supongamos que~$f^\an$ es una inmersión, y sean~$T$ la imagen
de~$X$ en~$Y$ y~$\overline{T}$ la clausura esquemática de~$f$. Se tiene una
factorización de~$f$,
$$
X \lto{i} \overline{T} \lto{j}
\ifthenelse{\boolean{orig}}
{J}
{Y}%
\quoi,
$$
donde~$j$ es una inmersión cerrada e~$i$ el morfismo canónico; de ella se
deduce la factorización siguiente de~$f^\an$:
$$
X^\an \lto{i^\an} \overline{T}^{\mkern1mu \an} \lto{j^\an} Y^\an\ .
$$
Como
\marginpar{324}
$T=f(X)$ es un subconjunto localmente constructible de~$Y$ (EGA~IV~1.8.4),
se tiene, según~\Ref{XII.2.2},
$\overline{T}^{\mkern1mu \an}=\overline{f^\an(X^\an)}$. De ello resulta que
$i^\an(X^\an)$ es un abierto de~$\overline{T}^{\mkern1mu \an}$ y, por tanto,
que~$i(X)$ es un abierto de~$\overline{T}$. Consideremos la factorización
canónica de~$i$:
$$
X \lto{i_1} i(X) \lto{i_2} \overline{T}.
$$
El morfismo~$i_1^\an$ es un monomorfismo propio; por tanto, también lo es
$i_1$ según~\Ref{XII.3.2}~(v) y~\Ref{XII.3.1}~(x). Esto demuestra que~$i_1$
y, por consiguiente, también~$f$ es una inmersión.

(v). Supongamos que~$f$ es propio y demostremos que~$f^\an$ también lo es.
Como el hecho de que~$f^\an$ sea propio es local sobre~$Y^\an$, podemos
suponer que~$Y$ es afín. Según el lema de Chow (EGA~II~5.6.1), puede
encontrarse un $Y$-esquema proyectivo~$X'$ y un morfismo proyectivo
sobreyectivo
$$
g\colon X'\to X.
$$
El morfismo $(fg)^\an=f^\an g^\an$ es proyectivo y, por tanto, propio;
$g^\an$ es sobreyectivo, y de~\cite[ch\ptbl 1 \S 10]{XII.1} resulta que
$f^\an$ es propio.

Recíprocamente, supongamos que~$f^\an$ es propio y demostremos que~$f$
también lo es. Según~\Ref{XII.3.1}~(viii), $f$ es separado. Falta probar que
$f$ es universalmente cerrado, e incluso basta demostrar que~$f$ es cerrado;
en efecto, para todo $Y$-esquema~$Y'$ localmente de tipo finito, el morfismo
$$
f_{(Y')}=h\colon X\times_YY'\to Y'
$$
también será cerrado, pues~$h^\an$ es propio. Sea~$T$ un subconjunto cerrado
de~$X$; $f(T)$ es un conjunto localmente constructible y se tiene
$$
f^\an(\varphi^{-1}(T))=\psi^{-1}(f(T)).
$$
Como
\marginpar{325}
$f^\an$ es propio, $\psi^{-1}(f(T))$ es un subconjunto cerrado de~$Y^\an$;
de~\Ref{XII.2.2} resulta, pues, que
$$
\psi^{-1}(\overline{f(T)})=\psi^{-1}(f(T))\quoi.
$$
Esto implica que $f(T)=\overline{f(T)}$, \ie que~$f$ es
\ifthenelse{\boolean{orig}}
{cerrada}
{cerrado}%
~y, por tanto, que~$f$ es propio.

(vi). Decir que un morfismo es finito equivale a decir que es propio de
fibras finitas (EGA~III~4.4.2 y~\cite[\no19 \S 5]{XII.4}). Como el conjunto
de los puntos en los que las fibras de~$f$ son finitas es localmente
constructible (EGA~IV~9.7.9), las fibras de~$f$ son finitas
\ifthenelse{\boolean{orig}}
{si y solo si}
{si y solo si}%
~lo son las fibras de~$f^\an$; por tanto, (vi) resulta de~(v).

\begin{remarque}
\label{XII.3.3}
\begin{enumerate}
\item[a)] sea $f\colon X\to Y$ un morfismo de $\CC$-esquemas localmente de
tipo finito. El hecho de que~$f^\an$ sea un isomorfismo local no implica que
lo mismo ocurra con~$f$. En efecto, si~$f$ es étale, $f^\an$ es étale y, por
tanto, un isomorfismo local~\hbox{\cite[\no13 \S 1]{XII.4}}, pero~$f$ no lo es
necesariamente.
\item[b)]
\ifthenelse{\boolean{orig}}
{el enunciado~\Ref{XII.3.2}}
{El enunciado~\Ref{XII.3.2}}%
~no es cierto si no se supone que~$f$ es de tipo finito. Demostremos, por
ejemplo, que~$f^\an$ puede ser una inmersión cerrada sin que~$f$ lo sea. Basta
tomar como~$X$ la suma disjunta de~$\ZZ$ copias de~$\Spec\CC$, como~$Y$ la
recta afín y como~$f$ el morfismo obtenido enviando los puntos de~$X$ a
puntos distintos de~$Y$ que forman un subconjunto discreto.
\end{enumerate}
\end{remarque}


\section[Teoremas de comparación cohomológica]{Teoremas de comparación cohomológica y teoremas de existencia}
\label{XII.4}
\marginpar{326}

El objeto de este número es volver a demostrar los resultados de
\cite[\no 2 th\ptbl 5 y th\ptbl 6]{XII.3}; estos generalizan al caso de un
esquema propio los teoremas establecidos en \cite[\no 12]{XII.10} cuando
$X$ es proyectivo, y los extienden al caso relativo. En
\cite[ch\ptbl VIII \no 3]{XII.7} se demuestran resultados más generales
relativos a esquemas relativos propios sobre un espacio analítico.

Recordemos que la cohomología de {\v{C}}ech utilizada en
\cite[\no 12]{XII.10} coincide con la cohomología usual tanto en el caso
algebraico como en el analítico (EGA~III 1.4.1 y
\cite[II 5.10]{XII.5}).
\subsection{}
\label{XII.4.1}
\setcounter{proposition}{1}% The first theorem in this source section is XII.4.2.
Sea $f\colon X\to Y$ un morfismo de~$\CC$-esquemas localmente de tipo
finito y consideremos el diagrama conmutativo
$$
\xymatrix{ X^\an \ar[r]^\varphi \ar[d]_{f^\an} & X \ar[d]^f
\\ Y^\an \ar[r]^\psi & Y \quoi.}
$$
Si $F$ es un $\cal{O}_X$-módulo, para todo entero $p\geq0$ se tienen
morfismos
$$
\R^p f_* F \lto{i} \R^p f_* (\varphi_* F^\an )
\lto{j} \R^p (f\cdot\varphi)_* F^\an \lto{k} \psi_*
(\R^p f_*^\an F^\an ) \quoi,
$$
donde $i$ se deduce del morfismo canónico $F\to\varphi_*F^\an$, y $j,k$
son ``homomorfismos de borde'' de sucesiones espectrales de Leray. Al
compuesto $k.j.i$ se le asocia un morfismo canónico
\begin{equation*}
\label{eq:XII.4.1.1}
\tag{4.1.1} {\theta_p \colon (\R^p f_* F)^\an \to \R^p f_*^\an
(F^\an )}
\end{equation*}

\begin{theoreme}
\label{XII.4.2}
Sean
\marginpar{327}
$f\colon X\to Y$ un morfismo propio de~$\CC$-esquemas localmente de tipo
finito y $F$ un $\cal{O}_X$-módulo coherente. Entonces, para todo entero
$p\geq0$, el morfismo~\eqref{eq:XII.4.1.1}
$$
\theta_p \colon (\R^p f_* F)^\an \to \R^p f_*^\an (F^\an)
$$
es un isomorfismo.
\end{theoreme}

1) \emph{Caso en que $f$ es proyectivo}. La demostración es análoga a la de
\cite[\no 13]{XII.10}. Recordémosla brevemente. Nos reducimos al caso en
que $X$ es un espacio proyectivo
\ifthenelse{\boolean{orig}}
{tipo $\PP^r_Y$}
{de tipo $\PP^r_Y$}%
~sobre~$Y$. Sean $\cal{Y}=Y^\an$ y $\cal{P}=\PP^r_\cal{Y}$; se demuestra
primero que
$$
f_*^\an \cal{O}_\cal{P} = \cal{O}_\cal{Y} \quoi, \qquad \R^p f_*^\an
(\cal{O}_\cal{P}) = 0 \qquad \text{para } p>0
$$
En efecto, para verificar las relaciones anteriores podemos reducirnos al
caso en que $\cal{Y}$ es una bola $\cal{B}$ de un espacio afín
$\cal{E}^n$. Consideremos el ``recubrimiento estándar'' $\{\cal{U}_i\}$
de~$\cal{P}$ por $r+1$ abiertos isomorfos a
$\cal{B}\times\cal{E}^r$. Como estos abiertos son de Stein, para todo
entero $p\geq0$ se tienen isomorfismos
$$
\H^p (\{ \cal{U}_i \}, \cal{O}_\cal{P}) \isomto \H^p (\cal{P},
\cal{O}_\cal{P})
$$
Podemos expresar entonces las secciones del haz estructural
$\cal{O}_\cal{P}$ sobre los abiertos $\cal{U}_i$ y sobre sus intersecciones
en términos de series de Laurent; un cálculo sencillo demuestra que
$$
\H^0 (\cal{P}, \cal{O}_\cal{P} ) \isomto \H^0 (\cal{Y},
\cal{O}_\cal{Y} ) \quoi, \qquad \H^p (\cal{P}, \cal{O}_\cal{P} ) = 0
\qquad \text{para } p>0
$$
La demostración concluye repitiendo \cite[\no 12 lema~5]{XII.10}, con los
grupos de cohomología sustituidos por los haces de cohomología.

2)
\marginpar{328}
\emph{Caso en que $f$ es propio}. Utilizamos EGA~III 3.1.2 para reducirnos
al caso proyectivo. Sea $\cal{K}$ la categoría de los
$\cal{O}_X$-módulos coherentes para los que $\theta_p$ es un isomorfismo
para todo $p\geq0$. Basta demostrar que, para toda sucesión exacta
$0\to F'\to F\to F''\to0$ dos de cuyos términos pertenecen a~$\cal{K}$,
lo mismo ocurre con el tercero; que todo sumando directo de un objeto
de~$\cal{K}$ pertenece a~$\cal{K}$; y que, para todo punto $x$ de~$X$,
puede encontrarse un objeto $F$ de~$\cal{K}$ tal que $F_x\neq0$.

\ifthenelse{\boolean{orig}}
{El primera condición}
{La primera condición}%
~resulta de aplicar el lema de los cinco al siguiente diagrama conmutativo,
cuyas filas son exactas:
$$
\xymatrix@=.7cm{ \ar[r] & (\R^p f_* F')^\an \ar[d] \ar[r] &(\R^p f_*
F)^\an \ar[d] \ar[r] & (\R^p f_* F'')^\an \ar[d] \ar[r] &
(\R^{p+1} f_* F' )^\an \ar[r] \ar[d] &\, \\ \ar[r] & \R^p
f_*^\an F'^\an \ar[r] &\R^p f_*^\an F^\an \ar[r] & \R^p
f_*^\an F''^\an \ar[r] & \R^{p+1} f_*^\an F'^\an \ar[r] &
\quoi,}
$$
y la segunda condición se verifica de manera análoga.

Para verificar la tercera condición, podemos limitarnos al caso en que $X$
es un esquema irreducible con punto genérico~$x$. Podíamos haber supuesto
desde el principio que~$Y$
era noetheriano. Por el lema de Chow (EGA~II 5.6.1), existe un
$Y$-esquema proyectivo~$X'$ y un morfismo proyectivo sobreyectivo
$g\colon X'\to X$. Por otra parte, existe un entero~$n$ tal que
$\R^p g_* (\cal{O}_{X'}(n))=0$ para todo $p>0$ y que el morfismo canónico
$g^*g_*(\cal{O}_{X'}(n))\to\cal{O}_{X'}(n)$ es sobreyectivo
(EGA~III 2.2.1). Si ponemos $F=g_*(\cal{O}_{X'}(n))$, el haz~$F$ posee la
propiedad requerida. En efecto, $F_x\neq0$; además, como la sucesión
espectral de Leray
$$
\R^p f_* (\R^q g_* (\cal{O}_{X'} (n))) \To \R^{p+q}
(\ifthenelse{\boolean{orig}}{f.g}{f\cdot g})_* (\cal{O}_{X'} (n))
$$
es degenerada, se tiene un isomorfismo
$$
\R^p f_* F \isomto \R^p (\ifthenelse{\boolean{orig}}{f.g}{f\cdot g})_* (\cal{O}_{X'} (n))
$$
\marginpar{329}%
Como en el caso algebraico, se tiene un isomorfismo canónico
$$\R^p f_*^\an F^\an \isomto
\R^p (\ifthenelse{\boolean{orig}}{f.g}{f\cdot g})_*^\an (\cal{O}_{X'}(n)^\an ),$$
y el diagrama
$$
\xymatrix@C=1cm{
(\R^p f_* F)^\an \ar[d]_{\theta_p} \ar[r]^-{\sim} & (\R^p (\ifthenelse{\boolean{orig}}{f.g}{f\cdot g})_* (\cal{O}_{X'} (n)))^\an \ar[d]_{\psi_p} \\
\R^p f_*^\an F^\an \ar[r]^-{\sim} & \R^p (\ifthenelse{\boolean{orig}}{f.g}{f\cdot g})_*^\an
(\cal{O}_{X'}(n)^\an)}
$$
es conmutativo. Por~1), $\psi_p$ es un isomorfismo; por tanto, también lo
es~$\theta_p$, lo cual concluye la demostración.

\begin{corollaire}
\label{XII.4.3}
Sean $X$ un $\CC$-esquema propio y $F$ un $\cal{O}_X$-módulo coherente.
Entonces, para todo entero $p\geq0$, el morfismo canónico
$$
\H^p(X,F)\to \H^p(X^{\an},F^{\an})
$$
es un isomorfismo.
\end{corollaire}

\begin{theoreme}
\label{XII.4.4}
Sea $X$ un $\CC$-esquema propio. El funtor que asocia a cada
$\cal{O}_X$-módulo coherente~$F$ su imagen inversa $F^\an$ sobre~$X^\an$
es una equivalencia de categorías.
\end{theoreme}

1) \emph{El funtor es plenamente fiel}. Sean, en efecto, $F$ y $G$ dos
$\cal{O}_X$-módulos coherentes. El morfismo canónico
$$
\Hom_{\cal{O}_X}(F,G)\to \Hom_{\cal{O}_{X^{\an}}}(F^{\an},G^{\an})
$$
se identifica con el morfismo canónico
$$
\H^0(X,\SheafHom_{\cal{O}_X}(F,G))\to
\H^0(X^{\an},
\ifthenelse{\boolean{orig}}
{\SheafHom_{\cal{O}_X}(F,G)}
{\SheafHom_{\cal{O}_{X^\an}}(F^\an,G^\an)})
$$
\marginpar{330}%
(EGA~$0_{\textup{I}}$~6.7.6). Como $\SheafHom_{\cal{O}_X}(F,G)$ es
coherente, de~\Ref{XII.4.3} resulta que este morfismo es biyectivo.

2) \emph{El funtor es esencialmente sobreyectivo}. Cuando $X$ es
proyectivo, la afirmación resulta de \cite[\no 12~th\ptbl 3]{XII.10}. El
caso general se reduce al anterior mediante el lema de Chow
(EGA~II~5.6.1). Sean, en efecto, $X'$ un $\CC$-esquema proyectivo,
$f\colon X'\to X$ un morfismo proyectivo sobreyectivo y $U$ un abierto
denso de~$X$ tal que $f$ induce un isomorfismo $f^{-1}(U)\simeq U$.
Razonamos por inducción noetheriana sobre~$X$; podemos, pues, suponer que,
para todo haz coherente $\cal{G}$ sobre~$X^\an$ para el cual existe un
subconjunto cerrado $Y$ de~$X$, distinto de~$X$, que satisface
$Y^\an\supset\Supp\cal{G}$, existe un haz coherente $G$ sobre~$X$ tal que
$G^\an\simeq\cal{G}$.

Sea $\cal{F}$ un haz coherente de $\cal{O}_{X^\an}$-módulos, y sean
$\cal{K}$ y $\cal{L}$ los haces coherentes definidos por la exactitud de la
sucesión
$$
0\to\cal{K}\to\cal{F}\to f_*^\an f^{\an *}\cal{F}\to\cal{L}\to 0
$$
Como $X'$ es proyectivo, existe un $\cal{O}_{X'}$-módulo coherente $F'$ tal
que $F'^\an\simeq{f^\an}^*\cal{F}$; de~\Ref{XII.4.2} se deduce entonces
un isomorfismo $(f_*F')^\an\simeq f_*^\an{f^\an}^*\cal{F}$. Como
$\cal{K}|U^\an$ y $\cal{L}|U^\an$ son nulos, existen
$\cal{O}_X$-módulos coherentes $K$ y $L$
\ifthenelse{\boolean{orig}}
{tal que}
{tales que}%
~$K^\an\simeq\cal{K}$ y $L^\an\simeq\cal{L}$. Por~1), el morfismo
$f_*^\an{f^\an}^*\cal{F}\to\cal{L}$ procede de un único morfismo
$f_*F'\to L$; sea $I=\Ker(f_*F'\to L)$. El haz $\cal{F}$ es entonces una
extensión de~$I^\an$ por~$K^\an$, y basta ver que esta extensión procede
\marginpar{331}
por imagen inversa de una extensión de~$I$ por~$K$. Basta, por tanto,
demostrar que el morfismo canónico
\begin{equation*}
\label{eq:XII.4.*}
\tag{$*$}
\ifthenelse{\boolean{orig}}
{\Ext_{\cal{O}_X}^q(I,K)^\an}
{\Ext_{\cal{O}_X}^q(I,K)}
\isomto
\Ext_{\cal{O}_{X^{\an}}}^q(I^{\an},K^{\an})\quad
\ifthenelse{\boolean{orig}}{q\neq 1}{q=1}
\end{equation*}
es biyectivo. Ahora bien, se tienen isomorfismos
$\mathbf{Ext}_{\cal{O}_X}^q(I,K)^\an\isomto
\mathbf{Ext}_{\cal{O}_{X^{\an}}}^q(I^{\an},K^{\an})$
para todo entero $q\geq0$ (EGA~$0_{\textup{III}}$~12.3.5), y un morfismo
de sucesiones espectrales
\begingroup
\shorthandoff{>}% Spanish babel otherwise treats the Xy-pic arrow tip as a shorthand.
$$
\xymatrix@R=.7cm@C=.5cm{
\H^p(X,\mathbf{Ext}_{\cal{O}_X}^q(I,K))\ar@{=>}[r]\ar[d] &
\Ext_{\cal{O}_X}^{p+q}(I,K)\ar[d]\\
\H^p(X^{\an},\mathbf{Ext}_{\cal{O}_{X^{\an}}}^q(I^{\an},K^{\an}))\ar@{=>}[r] & \Ext_{\cal{O}_{X^{\an}}}^{p+q}(I^{\an},K^{\an}).
}
$$
\endgroup
Este morfismo es un isomorfismo porque, por~\Ref{XII.4.3}, lo es en los
términos $E^{pq}_2$, y esto demuestra la biyectividad de~$(*)$.

\begin{corollaire}
\label{XII.4.5}
El funtor que asocia a cada $\CC$-esquema propio $X$ el espacio analítico
$X^\an$ es plenamente fiel.
\end{corollaire}

Debemos demostrar que, si $X$ e $Y$ son dos $\CC$-esquemas propios, la
aplicación canónica
$$
\Hom_{\CC}(X,Y)\to\Hom(X^{\an},Y^{\an})
$$
es biyectiva. Ahora bien, dar un morfismo de~$X$ en~$Y$ (\resp de
$X^\an$ en $Y^\an$) equivale a dar su grafo, \ie un subesquema cerrado
$Z$ de~$X\times Y$ (\resp un subespacio analítico cerrado $\othercal{Z}$
de~$X^\an\times Y^\an$), tal que la restricción de la primera proyección
$X\times Y\to X$ a~$Z$ (\resp la de
$X^\an\times Y^\an\to X^\an$ a~$\othercal{Z}$) sea
\marginpar{332}
un isomorfismo. Como dar un subesquema cerrado de~$X\times Y$ (\resp un
subespacio analítico cerrado de~$X^\an\times Y^\an$) equivale a dar un haz
coherente de ideales de~$\cal{O}_{X\times Y}$ (\resp de
$\cal{O}_{X^\an\times Y^\an}$), el corolario resulta de~\Ref{XII.4.4}.

\begin{corollaire}
\label{XII.4.6}
Sea $X$ un $\CC$-esquema propio. El funtor que asocia a cada esquema
finito (\resp finito étale) $X'$ sobre~$X$ el espacio analítico $X'^\an$
es una equivalencia de la categoría de los esquemas finitos (\resp finitos
étale) sobre~$X$ con la categoría de los espacios analíticos finitos
(\resp finitos étale) sobre~$X^\an$.
\end{corollaire}

En efecto, dar un morfismo finito $X'\to X$ (\resp
$X'^\an\to X^\an$) equivale a dar un haz coherente de álgebras sobre
$\cal{O}_X$ (\resp sobre $\cal{O}_{X^\an}$)
\cite[\no 19~\S 5~th\ptbl 2]{XII.4}. El corolario resulta, pues,
de~\Ref{XII.4.4} en el caso sin paréntesis, y el caso entre paréntesis se
deduce de él teniendo en cuenta~\Ref{XII.3.1}~(iii).


\section{Teoremas de comparación de los recubrimientos étales}
\label{XII.5}
\setcounter{subsection}{-1}
\subsection{}
\label{XII.5.0}
Precisemos la noción de recubrimiento finito de un espacio analítico.
\index{recubrimiento finito de un espacio analítico|hyperpage}%
Si $\othercal{X}$ es un espacio analítico, se dice que un espacio analítico
$\othercal{X}'$ finito sobre~$\othercal{X}$ es un recubrimiento finito de
$\othercal{X}$ si cada componente irreducible de~$\othercal{X}'$ domina
una componente irreducible de~$\othercal{X}$.


\begin{theoreme}[«Teorema de existencia de Riemann»]
\label{XII.5.1}
\index{Riemann (teorema de existencia de)|hyperpage}%
\index{teorema de existencia de Riemann|hyperpage}%
Sea $X$ un $\CC$-esquema localmente de tipo finito y sea $X^{\an}$
el espacio analítico asociado a~$X$. El funtor $\Psi$ que asigna
a cada recubrimiento étale finito $X'$ de~$X$ el espacio
$X^{\prime\an}$ es una equivalencia
\marginpar{333}
de la categoría de los recubrimientos étale finitos de~$X$ en
la categoría de los recubrimientos étale finitos de~$X^{\an}$.
\end{theoreme}

1) \emph{El funtor $\Psi$ es plenamente fiel.} Sean $X'$ y
$X''$ dos recubrimientos étale finitos de~$X$ y demostremos que
la aplicación canónica
\begin{equation*}
\label{eq:XII.5.*}
\tag{$*$} { \Hom_{X}(X',X'') \to
\Hom_{X^{\an}}(X^{\prime\an},X^{\prime\prime\an}) }
\end{equation*}
es biyectiva. Podemos suponer que $X'$ es conexo. Dar un
$X$-morfismo de~$X'$ en $X''$ equivale a dar una
componente conexa $X_i$ de~$X'\times_XX''$ tal que el morfismo
$X_i\to X'$ inducido por la primera proyección sea un
isomorfismo. Como las componentes conexas de~$X'\times_XX''$
corresponden biyectivamente a las componentes conexas de
$X^{\prime\an}\times_{X^{\an}}X^{\prime\prime\an}$~\eqref{XII.2.6}, y un
morfismo $X_i\to X'$ es un isomorfismo
\ifthenelse{\boolean{orig}}
{si y solo si}
{si y solo si}%
~lo es $X_i^{\an}\to X^{\prime\an}$, queda demostrada la
biyectividad de~\eqref{eq:XII.5.*}.

2) \emph{El funtor $\Psi$ es esencialmente sobreyectivo.} Sea
$\othercal{X}'$ un recubrimiento étale finito de~$X^{\an}$ y demostremos
que existe un recubrimiento étale $X'$ de~$X$ junto con
un isomorfismo $X^{\prime\an}\isomto\othercal{X}'$. En vista de~1), la
cuestión es local sobre~$X$ y podemos suponer, por tanto, que $X$ es afín.

a) Reducción al caso en que $X$ es normal. Podemos suponer que $X$
es reducido. Supongamos, en efecto, demostrado el teorema
para $X_{\red}$. El funtor que a un recubrimiento étale
finito $X'$ de~$X$ le asigna el recubrimiento étale finito
$X^{\prime\an}_{\red}$ de~$X_{\red}^{\an}$ es entonces una equivalencia.
Como se obtiene componiendo $\Psi$ con el funtor $\Theta$ que
a un recubrimiento étale finito de~$X^{\an}$ le asigna su imagen
inversa sobre~$X_{\red}^{\an}$, y $\Theta$ es plenamente
fiel, se sigue que $\Psi$ es una equivalencia de categorías.

Podemos
\marginpar{334}
suponer que $X$ es normal. Sea, en efecto, $\widetilde{X}$ la
normalización de~$X$ y sea $p\colon\widetilde{X}\to X$ el morfismo
canónico. Como $p$ es finito, $p$ es un morfismo de descenso
efectivo para la categoría de los recubrimientos
étales~(\Ref{IX.4.7}). Supuesto demostrado el teorema
para $\widetilde{X}$, pongamos
$\widetilde{\othercal{X}'}=\othercal{X}'\times_{X^{\an}}\widetilde{X}^{\an}$.
Existen un recubrimiento étale $\widetilde{X}'$ de
$\widetilde{X}$ y un isomorfismo $\widetilde{X}'{}^{\an}\simeq
\widetilde{\othercal{X}'}$. Resulta entonces de~1) que el dato
de descenso natural sobre~$\widetilde{\othercal{X}'}$ se levanta
a un dato de descenso sobre~$\widetilde{X}'$ relativo a
$\widetilde{X}\to X$; esto prueba la existencia de un recubrimiento
étale $X'$ de~$X$ y de un isomorfismo
$i\colon X^{\prime\an}\times_{X^{\an}}\widetilde{X}^{\an}\simeq
\widetilde{\othercal{X}'}$, cuyas imágenes inversas por las dos
proyecciones de
$\widetilde{X}^{\an}\times_{X^{\an}}\widetilde{X}^{\an}$ son iguales.
Según~\Ref{IX.3.2}, cuya demostración vale en el caso analítico,
el morfismo $\widetilde{X}^{\an}\to X^{\an}$ es un morfismo de descenso
para la categoría de los recubrimientos étales y, por consiguiente,
$i$ procede de un isomorfismo $X^{\prime\an}\simeq\othercal{X}'$.

b) Reducción al caso en que $X$ es regular. Sean $U$
el abierto de los puntos regulares de~$X$, $i\colon U\to X$ e
$i^{\an}\colon U^{\an}\to X^{\an}$ los morfismos canónicos; como
$X$ es normal, se tiene $\codim(X-U,X)\geq2$. Supongamos que existe un
recubrimiento étale $U'$ de~$U$ tal que
$U^{\prime\an}\simeq\othercal{X}'|U^{\an}$ y demostremos que entonces $U'$ se
prolonga en un recubrimiento étale $X'$ de~$X$ tal que
$X^{\prime\an}\simeq\othercal{X}'$. Basta ver que $U'$ se prolonga
en un recubrimiento étale $X'$ de~$X$; se tendrá entonces un
isomorfismo $X^{\prime\an}|U^{\an}\simeq\othercal{X}'|U^{\an}$. Pero si
$\cal{F}$ y $\cal{G}$ son los haces coherentes de álgebras
sobre~$\cal{O}_{X^{\an}}$ que definen, respectivamente, $\othercal{X}'$ y
$X^{\prime\an}$, la normalidad de~$X$ y la desigualdad
$\codim(X-U,X)\geq2$ implican que los morfismos canónicos
$$
\cal{F}\to i_*^{\an}(\cal{F}|U^{\an}) \qquad
\cal{G}\to i_*^{\an}(\cal{G}|U^{\an})
$$
son
\marginpar{335}
isomorfismos \cite[\no 3~prop\ptbl 4]{XII.11}. Se deduce que $\cal{F}$
y $\cal{G}$, y por tanto también $X^{\prime\an}$ y $\othercal{X}'$, son
isomorfos.

Sea $\varphi\colon X^{\an}\to X$ el morfismo canónico. Como el problema
de prolongar $U'$ a~$X$ es local sobre~$X$, basta demostrar que, para
cada punto $y$ de~$X^{\an}-U^{\an}$, el recubrimiento étale
$U'_{\varphi(y)}=U'\times_X\Spec\cal{O}_{X,\varphi(y)}$ de
$U_{\varphi(y)}=U\times_X\Spec\cal{O}_{X,\varphi(y)}$ se prolonga a
$\Spec\cal{O}_{X,\varphi(y)}$. Sea $H$ la
$\cal{O}_U$-álgebra coherente que define $U'$. El morfismo canónico
$$
\alpha\colon(i_*H)^{\an}\to i_*^{\an}(H^{\an})=\cal{F}
$$
define un morfismo de haces de módulos sobre
$\Spec\cal{O}_{X^{\an},y}$:
$$
\alpha_y\colon(i_*H)^{\an}_y\to\cal{F}_y\quoi,
$$
cuya restricción a
$U_y=U_{\varphi(y)}\times_{\Spec\cal{O}_{X,\varphi(y)}}
\Spec\cal{O}_{X^{\an},y}$ es un isomorfismo. Esto demuestra que
$H|U_y$ es trivial y, por tanto, que $U'_{\varphi(y)}$ se prolonga a
$\Spec\cal{O}_{X,\varphi(y)}$.

c) Caso en que $X$ es afín y regular. Sea
$$
j\colon X\to P
$$
una compactificación de~$X$, donde $P$ es un $\CC$-esquema
proyectivo y $j$ una inmersión abierta dominante. Gracias al
teorema de resolución de singularidades~\cite{XII.8}, se pueden
encontrar un esquema regular $R$ y un morfismo proyectivo
$r\colon R\to P$ tales que $r$ induzca un isomorfismo
$r^{-1}(X)\simeq X$ y $r^{-1}(X)$ sea el complementario en~$R$
de un divisor con cruzamientos normales. Sea
\marginpar{336}
$$
k\colon X\to R
$$
la inmersión canónica. Vamos a demostrar que existe un recubrimiento
finito normal~\eqref{XII.5.0} $\othercal{R}'$ de~$R^{\an}$ que prolonga
el recubrimiento étale $X^{\prime\an}$. Según la
Proposición~\Ref{XII.5.3} siguiente, tal recubrimiento es único; el
problema de prolongar $X^{\prime\an}$ es, por tanto, local sobre~$R^{\an}$
en un entorno de~$R^{\an}-X^{\an}$. Ahora bien, cada punto de
$R^{\an}-X^{\an}$ tiene un entorno abierto $\othercal{V}$ isomorfo a una
bola de un espacio afín $\othercal{E}^n$, tal que
$\othercal{V}-\othercal{V}\cap X^{\an}$ está definido por la anulación de
las primeras $p$ funciones coordenadas $z_1,\dots,z_p$, con
$0\leq p\leq n$. El grupo fundamental de
$\othercal{U}=\othercal{V}\cap X^{\an}$ es isomorfo a~$\ZZ^p$, y todo
recubrimiento étale de~$\othercal{U}$ es cociente de un recubrimiento de
la forma
$$
\othercal{U}''=\othercal{U}[T_1,\dots,T_p]/
(T_1^{n_1}-z_1,\dots,T_p^{n_p}-z_p)\quoi,
$$
donde los $n_i$ son enteros $>0$, por un subgrupo $H$ del grupo de
Galois $\ZZ/n_1\ZZ\times\cdots\times\ZZ/n_p\ZZ$ de~$\othercal{U}''$.
Ahora bien, $\othercal{U}''$ se prolonga en el recubrimiento regular
$$
\othercal{V}''=\othercal{V}[T_1,\dots,T_p]/
(T_1^{n_1}-z_1,\dots,T_p^{n_p}-z_p)\quoi,
$$
de~$\othercal{V}$ sobre el cual actúa $H$, y el cociente de
$\othercal{V}''$ por $H$ es la prolongación buscada.

La demostración concluye entonces gracias a~\Ref{XII.4.6}.
El recubrimiento $\othercal{R}'$ procede de un recubrimiento finito $R'$
de~$R$; la restricción de~$R'$ a~$X$ es un recubrimiento $X'$ de
$X$ tal que $X^{\prime\an}\simeq\othercal{X}'$, y según
\Ref{XII.3.1}(iii), $X'$ es un recubrimiento étale de~$X$.

\begin{corollaire}
\label{XII.5.2}
Sea
\marginpar{337}%
$X$ un $\CC$-esquema conexo localmente de tipo finito,
sea $\varphi\colon X^{\an}{\to}X$ el morfismo canónico y sea $x$ un
punto de~$X^{\an}$. Sea $\pi_1(X^{\an},x)$ el grupo fundamental del
espacio topológico $X^{\an}$ en~$x$, y sea $\pi_1(X,\varphi(x))$ el
grupo fundamental del esquema $X$ en~$\varphi(x)$
\textup{(\Ref{V.7})}. Entonces $\pi_1(X,\varphi(x))$ es canónicamente
isomorfo a la completación de~$\pi_1(X^{\an},x)$ para la topología de los
subgrupos de índice finito.
\end{corollaire}

En efecto, sea $\cal{C}$ la categoría de los recubrimientos étale finitos
de~$X^{\an}$, sea $F$ el funtor de~$\cal{C}$ en~$\Ens$ que a cada
recubrimiento étale finito $\othercal{X}'$ de~$X^{\an}$ le asigna el
conjunto de los puntos de~$\othercal{X}'$ situados sobre~$x$, y sea
$\widehat{\pi}_1(X^{\an},x)$ el grupo profinito asociado a~$\cal{C}$ y~$F$
como en~\Ref{V.4}. Como todo recubrimiento étale finito de~$X^{\an}$
es cociente del recubrimiento universal por un subgrupo de índice finito,
$\widehat{\pi}_1(X^{\an},x)$ no es sino la completación de
$\pi_1(X^{\an},x)$ para la topología de los subgrupos de índice finito.
El corolario resulta, pues, de~\Ref{XII.5.1} y~\Ref{V.6.10}.


\begin{proposition}
\label{XII.5.3}
Sean $\othercal{X}$ un espacio analítico normal y $\othercal{Y}$ un
subconjunto analítico cerrado tal que
$\othercal{U}=\othercal{X}-\othercal{Y}$ sea denso en~$\othercal{X}$.
Entonces el funtor que a cada recubrimiento finito
normal~\eqref{XII.5.0} $\othercal{X}'$ de~$\othercal{X}$ le asigna su
restricción a~$\othercal{U}$ es plenamente fiel.
\end{proposition}

Sean $\othercal{X}'$
\ifthenelse{\boolean{orig}}
{y}
{y}%
$\othercal{X}''$ dos recubrimientos finitos normales de~$\othercal{X}$.
Debemos demostrar que la aplicación canónica
$$
\Hom_{\othercal{X}}(\othercal{X}',\othercal{X}'') \to
\Hom_{\othercal{U}}(\othercal{X}'|\othercal{U},\othercal{X}''|\othercal{U})
$$
es biyectiva. Sean $u$, $v$ dos $\othercal{X}$-morfismos de
$\othercal{X}'$ en~$\othercal{X}''$ cuyas restricciones a~$\othercal{U}$
son iguales y demostremos que $u=v$. Los morfismos $u$ y
\marginpar{338}
$v$ coinciden sobre el abierto denso
$\othercal{U}\times_{\othercal{X}}\othercal{X}'$ y, por tanto, sobre los
espacios topológicos subyacentes. Según
\hbox{\cite[\no 19~\S 4~cor\ptbl 5]{XII.4}}, esto demuestra que $u=v$.

Sea ahora $u$ un $\othercal{U}$-morfismo de
$\othercal{X}'|\othercal{U}$ en~$\othercal{X}''|\othercal{U}$ y
demostremos que se prolonga a todo~$\othercal{X}'$. Podemos suponer que
$\othercal{X}'$ es regular. En efecto, como $\othercal{X}'$ es normal,
se puede encontrar un abierto $\othercal{V}$ de~$\othercal{X}$ cuyo
complementario sea un subconjunto analítico de codimensión $\geq2$, tal
que $\othercal{X}'\times_{\othercal{X}}\othercal{V}=\othercal{V}'$ sea
regular. Sea
$\othercal{V}''=\othercal{X}''\times_{\othercal{X}}\othercal{V}$ y
supongamos demostrada la proposición para~$\othercal{V}$. Consideremos el
diagrama conmutativo
\begingroup
\shorthandoff{>}% Spanish babel otherwise treats the Xy-pic arrow tip as a shorthand.
$$
\xymatrix{ \othercal{V}' \ar[dd]_{g'}\ar[rd]\ar[rrr]^{i'} && &\othercal{X}'
\ar[dd]^{\hbox{\raise10mm\hbox{$f'$}}} \cr &\othercal{V}'' \ar[ld]_{g''\mkern-10mu}\ar[rrr]^-{i''~~} & &
&\othercal{X}'' \ar[ld]_{f''\mkern-12mu} \cr \othercal{V} \ar[rrr]^{i} && &\othercal{X} \cr
}
$$
\endgroup
A~$u$ se le asocia un morfismo de
$\cal{O}_{\othercal{V}}$-álgebras
$g''_*\cal{O}_{\othercal{V}''}\to g'_*\cal{O}_{\othercal{V}'}$, de donde
se deduce un morfismo
$$
i_*g''_*\cal{O}_{\othercal{V}''}\to
i_*g'_*\cal{O}_{\othercal{V}'}\quoi.
$$
Teniendo en cuenta los isomorfismos
$i'_*\cal{O}_{\othercal{V}'}\simeq\cal{O}_{\othercal{X}'}$ e
$i''_*\cal{O}_{\othercal{V}''}\simeq\cal{O}_{\othercal{X}''}$
\cite[\no 3~prop\ptbl 4]{XII.11}, se deduce un morfismo de
$\cal{O}_{\othercal{X}}$-álgebras
$$
f''_*\cal{O}_{\othercal{X}''}\to
f'_*\cal{O}_{\othercal{X}'}\quoi,
$$
de donde se obtiene el morfismo buscado
$\othercal{X}'\to\othercal{X}''$.

Supongamos
\marginpar{339}
en adelante que $\othercal{X}'$ es regular. Sean
$\othercal{U}'=\othercal{U}\times_{\othercal{X}}\othercal{X}'$ y
$\othercal{Y}'=\othercal{X}'-\othercal{U}'$. Consideremos
$\othercal{Y}'$ como subespacio analítico reducido de~$\othercal{X}'$;
si $\othercal{Y}'_1$ es el lugar singular de~$\othercal{Y}'$, se tiene
$\dim\othercal{Y}'_1<\dim\othercal{Y}'$
\cite[\no 20~D\ptbl Th\ptbl 3]{XII.4}. Por recurrencia sobre la dimensión
de~$\othercal{Y}'$, podemos suponer, pues, que $\othercal{Y}'$ es liso.
Como basta prolongar $u$ a un entorno abierto de cada punto de
$\othercal{Y}'$, por el teorema de las funciones implícitas podemos
suponer que $\othercal{X}'$ es una bola de un espacio afín
$\othercal{E}^n$ y que $\othercal{Y}'$ es el cerrado definido por la
anulación de las primeras $p$ funciones coordenadas
$z_1,\dots,z_p$, con $0\leq p\leq n$.

A~$u$ le asociamos una sección $s$ de
$p\colon\othercal{X}'\times_{\othercal{X}}\othercal{X}''\to\othercal{X}'$
sobre~$\othercal{U}'$; restringiendo $\othercal{X}'$ si es necesario,
podemos suponer que
$p_*(\cal{O}_{\othercal{X}'\times_{\othercal{X}}\othercal{X}''})$ está
engendrado por elementos $x_1,\dots,x_q$ de
$\Gamma(\othercal{X}',p_*\cal{O}_{\othercal{X}'\times_{\othercal{X}}\othercal{X}''})$;
sean $u_1,\dots,u_q\in\Gamma(\othercal{U}',\cal{O}_{\othercal{X}'})$ las
imágenes por~$s$ de~${x_1}|\othercal{U}',\dots,{x_q}|\othercal{U}'$.
Decir que $s$ se prolonga a~$\othercal{X}'$ equivale a decir que
$u_1,\dots,u_q$ se prolongan en secciones de
$\Gamma(\othercal{X}',\cal{O}_{\othercal{X}'})$. Pero como
\ifthenelse{\boolean{orig}}
{$f$}
{$p$}%
~es finito, cada $u_i$ es una serie de Laurent en
$z_1,\dots,z_p$, con coeficientes que son series de potencias en
$z_{p{+}1},\dots,z_n$, y satisface relaciones de dependencia integral.
Se deduce que $u_i$ es acotado, luego es una serie de potencias en
$z_1,\dots,z_n$ y, por consiguiente, se prolonga a~$\othercal{X}'$.

Cabe preguntarse si el funtor introducido en~\Ref{XII.5.3} es una
equivalencia de categorías. El teorema de
\textsc{Grauert}--\textsc{Remmert} \cite{XII.6} da una respuesta a esta
pregunta; a continuación presentamos una demostración que utiliza la
resolución de singularidades. También se habría podido emplear el teorema
de \textsc{Grauert}--\textsc{Remmert} para demostrar~\Ref{XII.5.1}; así
se hacía antes de disponer de~\cite{XII.8}.

\ifthenelse{\boolean{orig}}
{}
{\enlargethispage{.5cm}}
\begin{theoreme}[Teorema de \textsc{Grauert}--\textsc{Remmert}]
\label{XII.5.4}
Sean
\marginpar{340}
\index{Grauert-Remmert (teorema de)|hyperpage}%
\index{teorema de Grauert-Remmert|hyperpage}%
$\othercal{X}$ un espacio analítico normal y $\othercal{Y}$ un
subconjunto analítico cerrado tal que
$\othercal{U}=\othercal{X}-\othercal{Y}$ sea denso en~$\othercal{X}$.
Sea $\othercal{U}'$ un recubrimiento finito normal de~$\othercal{U}$;
supongamos que existe un subconjunto analítico cerrado de interior vacío
$\othercal{S}$ de~$\othercal{X}$ tal que la restricción de
$\othercal{U}'$ a~$\othercal{U}-\othercal{U}\cap\othercal{S}$ sea étale.
Entonces existe un recubrimiento finito normal $\othercal{X}'$ de
$\othercal{X}$ que prolonga $\othercal{U}'$, y $\othercal{X}'$ es único
salvo isomorfismo.
\end{theoreme}

La unicidad resulta de~\Ref{XII.5.3}. El problema de prolongar
$\othercal{U}'$ es, por tanto, local sobre~$\othercal{X}$. Podemos
suponer que $\othercal{U}$ es regular y que $\othercal{U}'$ es étale
sobre~$\othercal{U}$. En efecto, el conjunto de los puntos regulares de
$\othercal{U}$ es un abierto $\othercal{V}$ denso en~$\othercal{X}$ cuyo
complementario es un subconjunto analítico
\hbox{\cite[\no20~D th\ptbl 2]{XII.4}}, y basta sustituir
$\othercal{U}$ por el abierto
$\othercal{V}-\othercal{V}\cap\othercal{S}$.

Sea $y$ un punto de~$\othercal{X}-\othercal{U}$ y demostremos que se
puede prolongar $\othercal{U}'$ a un entorno de~$y$. Restringiendo
$\othercal{X}$ a un entorno abierto de~$y$, resulta del teorema de
resolución de singularidades \cite{XII.8} que se pueden encontrar un
espacio analítico regular $\othercal{X}_1$ y un morfismo proyectivo
$f\colon\othercal{X}_1\to\othercal{X}$ que induce por restricción a
$\othercal{U}$ un isomorfismo
$\othercal{U}_1=f^{-1}(\othercal{U})\simeq\othercal{U}$, tales que
$\othercal{U}_1$ sea el complementario en~$\othercal{X}_1$ de un divisor
con cruzamientos normales. Demostremos que $\othercal{U}'$ se prolonga
en un recubrimiento finito normal de~$\othercal{X}_1$. Como la cuestión
es local sobre~$\othercal{X}_1$, podemos suponer que $\othercal{X}_1$ es
una bola de un espacio afín $\othercal{E}^n$ y que
$\othercal{X}_1-\othercal{U}_1$ está definido por la anulación de las
primeras $p$ funciones coordenadas
\ifthenelse{\boolean{orig}}
{$z_1,\dots,z_q$,}
{$z_1,\dots,z_p$,}
con $0\leq p\leq n$. El recubrimiento étale $\othercal{U}'$ de
$\othercal{U}_1$ es cociente de un recubrimiento de la forma
$$
\othercal{U}_2=\othercal{U}_1[T_1,\dots,T_p]/\left(T_1^{n_1}-z_1,\dots,
T_p^{n_p}-z_p\right)
$$
por un subgrupo $H$ del grupo de Galois de~$\othercal{U}_2$. El
recubrimiento $\othercal{U}_2$ se prolonga
\marginpar{341}
en el recubrimiento
$$
\othercal{X}_2=\othercal{X}_1[T_1,\dots,T_p]/\left(T_1^{n_1}-z_1,\dots,
T_p^{n_p}-z_p\right)
$$
de~$\othercal{X}_1$, sobre el cual actúa $H$, y
$\othercal{X}_2/H$ prolonga $\othercal{U}'$ a~$\othercal{X}_1$.

Denotemos por~$\othercal{X}_1'$ el recubrimiento finito normal de
$\othercal{X}_1$ que prolonga $\othercal{U}'$ y por~$\othercal{F}_1$ la
$\cal{O}_{\othercal{X}_1}$-álgebra coherente
\ifthenelse{\boolean{orig}}
{definida por~$\othercal{F}_1$}
{que define~$\othercal{X}_1'$}%
. Según el teorema de finitud de
\textsc{Grauert}--\textsc{Remmert}
\hbox{\cite[\no15 th\ptbl 1.1]{XII.4}}, $f_*\othercal{F}_1$ es una
$\cal{O}_{\othercal{X}}$-álgebra coherente. Le corresponde, pues, un
recubrimiento finito $\othercal{X}'$ de~$\othercal{X}$, que además es
normal porque $\othercal{X}_1'$ lo es, y $\othercal{X}'$ es la
prolongación buscada de~$\othercal{U}'$.

\begin{remarque}
\label{XII.5.5}
En el enunciado~\Ref{XII.5.4} no se puede suprimir la hipótesis sobre el
lugar de los puntos en los que el morfismo
$\othercal{U}'\to\othercal{U}$ no es étale. Sea, por ejemplo,
$\othercal{X}$ el disco unidad del plano complejo, $\othercal{U}$ el
complementario del origen en~$\othercal{X}$ y
$\othercal{U}'=\othercal{U}[T]/(T^2-\sin 1/z)$, donde $z$ es la función
coordenada sobre~$\othercal{X}$. Entonces $\othercal{U}'$ es un
recubrimiento finito normal de~$\othercal{U}$ que no se prolonga a
$\othercal{X}$. Supongamos, en efecto, que $\othercal{U}'$ se prolonga
en un recubrimiento finito $\othercal{X}'$ de~$\othercal{X}$; el lugar de
los puntos de~$\othercal{X}$ en los que el morfismo
$\othercal{X}'\to\othercal{X}$ no es étale sería entonces un cerrado
analítico que contendría todos los puntos $z$ tales que
$\sin 1/z=0$, lo cual es absurdo.

Sin embargo, se puede suprimir la hipótesis sobre el lugar singular del
morfismo $\othercal{U}'\to\othercal{U}$ cuando
$\codim(\othercal{X}-\othercal{U},\othercal{X})\geq2$. Podemos suponer,
en efecto, que $\othercal{U}$ es regular. El lugar de los puntos de
$\othercal{U}$ en los que $\othercal{U}'\to\othercal{U}$ no es étale es
un divisor de~$\othercal{U}$, y del teorema de
\textsc{Remmert}--\textsc{Stein} \cite[th\ptbl 3]{XII.9} resulta que es
\marginpar{342}
la traza sobre~$\othercal{U}$ de un divisor de~$\othercal{X}$. Ahora
bien, en este caso, si $\othercal{A}$ es una
$\cal{O}_{\othercal{U}}$-álgebra coherente tal que
$\othercal{U}'=\Spec\an(\othercal{A})$ y si
$i:\othercal{U}\to\othercal{X}$ es el morfismo canónico, basta
tomar $\othercal{X}'=\Spec\an(i_*\othercal{A})$; se sabe, en efecto, que
$i_*\othercal{A}$ es
\ifthenelse{\boolean{orig}}
{coherentes}
{coherente}
\cite[\no 1 th\ptbl 1]{XII.11}.
\end{remarque}

\ifthenelse{\boolean{orig}}
{}
{\begin{remarqueMR}
\label{XII.5.6}
Existen grupos no triviales finitamente presentados $G$ que no poseen
subgrupos de índice finito distintos de~$G$, por ejemplo el grupo de
G\ptbl Higman (\cf J-P\ptbl Serre, \emph{Arbres et amalgames}, Astérisque
N$^\circ$46, prop\ptbl 6, cap\ptbl I, \S 1). Por consiguiente, si un grupo
de este tipo se realiza como grupo fundamental topológico de un esquema
$V$ sobre~$\CC$, digamos liso y proyectivo, el espacio topológico
$V_{\mathrm{top}}$ subyacente a~$V$ no es simplemente conexo, pero $V$
es algebraicamente simplemente conexo. En la actualidad no se conoce
ningún $V$ de este tipo. Señalemos, no obstante, que D\ptbl Toledo ha
construido esquemas lisos y proyectivos $V$ sobre~$\CC$ cuyo grupo
fundamental topológico no es Hausdorff para la topología de los
subgrupos de índice finito (el morfismo natural
$\pi_1(V_{\mathrm{top}})\to\pi_1(V_{\alg})$ no es inyectivo).
[D\ptbl Toledo, Projective varieties with non-residually finite
fundamental group, Publ.\ Math.\ IHES \textbf{77} (1993),
p\ptbl 103--119].
\end{remarqueMR}}


\begin{sgabibliography}{10}{XII.6}
\marginpar{343}

\bibitem[1]{XII.1} N\ptbl \textsc{Bourbaki},
Topologie Générale, Hermann, Paris (1960).

\bibitem[2]{XII.2} N\ptbl \textsc{Bourbaki},
Algèbre Commutative, Hermann, Paris (1961).

\bibitem[3]{XII.3} H\ptbl \textsc{Cartan}. Séminaire E.N.S., Paris (1956-57).

\bibitem[4]{XII.4} H\ptbl \textsc{Cartan}. Séminaire E.N.S., Paris (1960-61).

\bibitem[5]{XII.5} R\ptbl \textsc{Godement}, Théorie des Faisceaux, Hermann, Paris
(1958).

\bibitem[6]{XII.6} H\ptbl \textsc{Grauert} y R\ptbl \textsc{Remmert},
\ifthenelse{\boolean{orig}}
  {Komplexe R\"aume, Math.\ Ann., 136 (1958).}
  {Komplexe R\"aume, Math.\ Ann. \textbf{136} (1958), p\ptbl 245--318.}

\bibitem[7]{XII.7} M\ptbl \textsc{Hakim}, Schémas relatifs, tesis, Paris (1967).

\bibitem[8]{XII.8} H\ptbl \textsc{Hironaka},
\ifthenelse{\boolean{orig}}
{Resolution of singularities of an algebraic variety over a field of characteristic zero, Annals of Math., 39 (1964).}
{Resolution of singularities of an algebraic variety over a field of characteristic zero, Ann.\ of Math. \textbf{39} (1964), p\ptbl 109--236.}

\bibitem[9]{XII.9} R\ptbl \textsc{Remmert} y K\ptbl \textsc{Stein},
Ueber die wesentlichen Singularit\"aten analytischer Mengen,
\ifthenelse{\boolean{orig}}
{Math.\ Ann., 126 (1953) p\ptbl 263--306.}
{Math.\ Ann. \textbf{126} (1953), p\ptbl 263--306.}

\bibitem[10]{XII.10}
\ifthenelse{\boolean{orig}}{J.P\ptbl}{J-P\ptbl}%
\textsc{Serre},
\ifthenelse{\boolean{orig}}
{Géométrie Algébrique et Géométrie Analytique, Annales de l'Institut Fourier, VI (1956), p\ptbl 1--42.}
{Géométrie algébrique et géométrie analytique, Ann.\ Inst.\ Fourier (Grenoble) \textbf{6} (1956), p\ptbl 1--42.}

\bibitem[11]{XII.11}
\ifthenelse{\boolean{orig}}{J.P\ptbl}{J-P\ptbl}%
\textsc{Serre},
\ifthenelse{\boolean{orig}}
{Prolongement des Faisceaux analytiques
cohérents, Annales de l'Institut Fourier, XVI (1966), p\ptbl 363--374.}
{Prolongement de faisceaux analytiques
cohérents, Ann.\ Inst.\ Fourier (Grenoble) \textbf{16} (1966), p\ptbl 363--374.}

\end{sgabibliography}


\chapterspace{-4}
\chapter[Propiedad de ser cohomológicamente propios los haces...]{Propiedad de ser cohomológicamente propios\\ los haces de conjuntos y\\ los haces de grupos no conmutativos}
\label{XIII}
\marginpar{344}

\begin{center}
{por la Sra. M.\ \textsc{Raynaud}\footnote{A partir de notas inéditas de
A\ptbl\ Grothendieck.}}
\end{center}
\vspace*{1cm}

La presente exposición se propone utilizar la cohomología étale para
generalizar ciertos resultados de~\Ref{IX} y~\Ref{X}. Muestra asimismo cómo
se pueden extender a los haces de grupos no necesariamente conmutativos los
resultados de~SGA~5~II que siguen teniendo sentido para tales haces. Se dan
por conocidas las nociones de cohomología étale expuestas en~SGA~4.

El resultado principal \eqref{XIII.2.4} proporciona un ejemplo importante de
un morfismo no propio $f\colon~U\to~S$ que es \og cohomológicamente propio en
dimensión~$\leq1$\fg, es decir, tal que, para ciertos haces de grupos~$F$
sobre~$U$ (en el sentido de la topología étale), la formación de~$f_*F$ y
$\R^1f_*F$ conmuta con todo cambio de base~$S'\to~S$. En efecto, esta
propiedad la satisface el abierto~$U$ de un esquema~$X$ propio sobre~$S$ que
es el complementario de un divisor~$D$ con cruzamientos normales
relativamente a~$S$, al menos si se exige que~$F$ sea constante finito y de
orden primo con las características residuales de~$S$. Si ya no se supone que
$F$ sea de orden primo con las características residuales de~$S$, se obtiene
un resultado análogo sustituyendo~$\R^1f_*F$ por el subhaz
$\R^1_{\tame}f_*F$ obtenido al limitarse a considerar los torsores bajo~$F$
\og moderadamente ramificados sobre~$X$ relativamente a~$S$\fg. En
particular, esto permite demostrar que el grupo fundamental moderadamente
ramificado de una curva algebraica propia y lisa sobre un cuerpo
separablemente cerrado, privada de un número finito de puntos cerrados, está
topológicamente generado por un número finito de elementos
\eqref{XIII.2.12}.

El
\marginpar{345}
\No \Ref{XIII.4} está dedicado a la sucesión exacta de homotopía y a la
fórmula de Künneth.

Por último, un apéndice ofrece variantes útiles del lema de Abhyankar
demostrado en~\Ref{X.3.6}.

\setcounter{section}{-1}


\section{Recordatorio sobre la teoría de stacks}
\label{XIII.0}

En lo que sigue utilizaremos la teoría de stacks expuesta en
\cite{XIII.1} y~\cite{XIII.2}. Nos limitamos al caso del sitio étale de un
esquema. Dado un esquema~$X$, denotemos por~$X_\et$ el sitio étale de~$X$.
Recordemos que un
\ifthenelse{\boolean{orig}}
{stack}
{stack $F$}
\index{stack|hyperpage}%
sobre~$X$ es una categoría fibrada sobre~$X_\et$ tal que, para todo
esquema~$X'$ étale sobre~$X$ y para todo par
\ifthenelse{\boolean{orig}}
{de objeto}
{de objetos}
$x$, $y$ de la fibra~$F_{X'}$, el prehaz~$\SheafHom_{X'}(x,y)$ sea un haz,
y tal que, para todo morfismo étale sobreyectivo $X''\to X'$, todo objeto
de~$F_{X''}$ provisto de un dato de descenso relativamente a $X''\to X'$
sea imagen inversa de un objeto de~$F_{X'}$.

Se denota por~$F(X')$ la categoría de las secciones cartesianas de~$F/X'$.
Más generalmente, si~$\Sch_X$ es la categoría de los esquemas sobre~$X$
provista de la topología étale, el stack~$F$ puede extenderse a un
stack~$\cal{F}$ sobre~$\Sch_X$ y, para todo morfismo $f\colon X'\to X$, se
denota también por~$F(X')$ la categoría de las secciones cartesianas de este
stack~$\cal{F}$ sobre~$X'$.

Un gerbe
\index{gerbe|hyperpage}%
es un stack tal que, para todo esquema~$X'$ étale sobre~$X$ y para todo par
de objetos $x$, $y$ de~$F_{X'}$, todo morfismo de~$x$ a~$y$ sea un
isomorfismo, que $x$ e~$y$ sean localmente isomorfos y que la familia de
objetos~$X'$ de~$X_\et$ tales que la fibra~$F_{X'}$ no sea vacía constituya
un recubrimiento de~$X_\et$. Por ejemplo, el stack de los torsores bajo un haz
de grupos es un gerbe que, además, posee una sección cartesiana.
Recíprocamente, un gerbe que posee una sección, \ie tal que existe un
objeto~$x$ de~$F_X$, es equivalente al stack de los torsores bajo el haz
\marginpar{346}
de grupos~$\SheafAut_X(x)$.

Existe una noción evidente de subgerbe y de subgerbe maximal de un
stack~$F$. Dada una sección cartesiana~$x$ de~$F(X)$, existe un único
subgerbe maximal~$G_x$ de~$F$ tal que~$x$ se factoriza a través de~$G_x$.
$G_x$ recibe el nombre de subgerbe generado por~$x$; por definición, es un
gerbe trivial.
\index{gerbe trivial|hyperpage}%
El prehaz~$SF$
\label{indnot:mb}\oldindexnot{$SF$ o $S(F)$|hyperpage}%
definido por
$$
SF(X') = \{ \text{subgerbes maximales de }F|X' \}
$$
es un haz, llamado haz de los subgerbes maximales de~$F$.
\index{haz de los subgerbes maximales|hyperpage}%
Sea~$O$ el prehaz definido por
$$
O(X') = \{ \text{clases de objetos de }F_{X'}\text{ salvo isomorfismo}
\} \text{.}
$$
Asociando a todo objeto~$x$ de~$F_{X'}$ el subgerbe maximal de~$F|X'$
generado por~$x$, se obtiene un morfismo
$$
O \to SF \text{;}
$$
según~\cite[III 2.1.4]{XIII.2}, este morfismo hace de~$SF$ un haz asociado
a~$O$.

Se dice que un stack~$F$ es \emph{constructible}
\index{constructible (stack)|hyperpage}%
(\emph{\resp ind-$\LL$-finito},
\index{stack ind-finito (\resp ind-$l$-finito)|hyperpage}%
donde~$\LL$ es un conjunto de números primos) si, para todo esquema~$X'$
étale sobre~$X$ y para todo objeto~$x$ de~$F_{X'}$, el haz
$\SheafAut_{X'}(x)$ posee la propiedad correspondiente
\cite[VII 2.2.1]{XIII.2}. Se dice que un stack es $1$-constructible
\index{stack $1$-constructible (\resp constructible)|hyperpage}%
si es constructible y si el haz de los subgerbes maximales es constructible.


\section{Propiedad de ser cohomológicamente propio}
\label{XIII.1}
\setcounter{subsection}{-1}
\subsection{}
\label{XIII.1.0}
Sean~$S$ un esquema y $f \colon X \to Y$ un morfismo de
$S$-esquemas. Si~$S'$ es un $S$-esquema, consideremos el diagrama
siguiente, en el que todos los cuadrados son cartesianos:
\marginpar{347}
\begin{equation*}
\label{eq:XIII.1.0.1}
\tag{\thesubsection.1}
\begin{array}{c}
\xymatrix{ X \ar[d]_f & \ar[l]_h X'
\ar[d]^{f'} \\ Y \ar[d] & \ar[l]_g Y' \ar[d] \\ S & \ar[l] \,S'. }
\end{array}
\end{equation*}

Si~$Y_1$ es un esquema étale sobre~$Y$, pongamos $X_1 =
X \times_Y Y_1$, $Y_1' = Y' \times_Y Y_1$, y consideremos el
cuadrado cartesiano
\begin{equation*}
\label{eq:XIII.1.0.2}
\tag{\thesubsection.2}
\begin{array}{c}
\xymatrix{ X_1 \ar[d]_{f_1} & \ar[l]_{h_1}
X_1' \ar[d]^{f_1'} \\ Y_1 & \ar[l]^{g_1} \,Y_1'. }
\end{array}
\end{equation*}

\begin{definition}
\label{XIII.1.1}
Sea~$F$ un stack sobre~$X$. Se dice que~$(F,f)$ es
\emph{cohomológicamente propio relativamente a~$S$ en dimensión
$\leqslant -1$}
\index{cohomológicamente propio|hyperpage}%
\index{propiedad de ser cohomológicamente propio|hyperpage}%
(\resp en dimensión $\leqslant 0$, \resp en dimensión $\leqslant 1$)
si, para todo $S$-esquema~$S'$, el funtor canónico (definido de manera
evidente mediante la propiedad universal de la imagen inversa de stacks)
$$
g^* f_* F \to f'_* h^* F \qquad\text{ (\cf 1.0.1)}
$$
es fiel (\resp plenamente fiel, \resp una equivalencia de categorías).
\end{definition}

Si no hay confusión posible respecto de~$S$, en particular si~$S=Y$, se
dice cohomológicamente propio en lugar de cohomológicamente propio
relativamente a~$S$.

\setcounter{subsection}{1}% Conservar el número visible XIII.1.2 de la fuente.
\subsection{}
\label{XIII.1.2}
Sea~$F$ un haz de conjuntos sobre~$X$; sea~$\Phi$ el stack en categorías
discretas asociado a~$F$, \ie el stack cuya fibra sobre todo esquema~$X_1$
étale sobre~$X$ es la categoría discreta cuyo conjunto de objetos es
$F(X_1)$. Se dice que~$(F,f)$ es cohomológicamente propio relativamente
a~$S$ en dimensión $\leqslant -1$ (\resp en dimensión $\leqslant 0$) si
$(\Phi,f)$ es cohomológicamente propio relativamente a~$S$ en dimensión
$\leqslant 0$ (\resp en dimensión $\leqslant 1$).

El
\marginpar{348}
morfismo canónico
\begin{equation*}
\label{eq:XIII.1.2.1}
\tag{\thesubsection.1} { g^* f_* F \to f'_* h^* F }
\end{equation*}
induce, al pasar a los stacks en categorías discretas asociados, el
morfismo canónico
$$
g^* f_* \Phi \to f'_* h^* \Phi \text{.}
$$
Por consiguiente, decir que~$(F,f)$ es cohomológicamente propio
relativamente a~$S$ en dimensión $\leqslant -1$ (\resp en dimensión
$\leqslant 0$) equivale a decir que, para todo $S$-esquema~$S'$, el
morfismo~\eqref{eq:XIII.1.2.1} es inyectivo (\resp biyectivo).

\subsection{}
\label{XIII.1.3}
Sea~$F$ un haz de grupos sobre~$X$ y~$\Phi$ el stack de los torsores
sobre~$X$ bajo~$F$ \cite[II 2.3.2]{XIII.1}. Se dice que~$(F,f)$ es
cohomológicamente propio relativamente a~$S$ en dimensión
$\leqslant -1$ (\resp $\leqslant 0$, \resp $\leqslant 1$) si
$(\Phi,f)$ es cohomológicamente propio relativamente a~$S$ en dimensión
$\leqslant -1$ (\resp $\leqslant 0$, \resp $\leqslant 1$). La condición
de ser cohomológicamente propio puede explicitarse como sigue.

\begingroup
\setcounter{proposition}{0}% La fuente numera esta proposición dentro del párrafo XIII.1.3.
\renewcommand{\theproposition}{\thesubsection.\arabic{proposition}}
\renewcommand{\theHproposition}{\theHsubsection.\arabic{proposition}}
\begin{subproposition}
\label{XIII.1.3.1}
Las notaciones son las de~\eqref{eq:XIII.1.0.1} y
\eqref{eq:XIII.1.0.2}. Sea~$F$ un haz de grupos sobre~$X$. Se denota por
$F'$ (\resp $F_1$, \resp $F_1'$, etc.) la imagen inversa de~$F$ sobre~$X'$
(\resp sobre~$X_1$, \resp sobre~$X_1'$, etc.). Entonces las condiciones
siguientes son equivalentes:
\begin{enumerate}
\item[(i)] $(F,f)$ es cohomológicamente propio relativamente a~$S$ en
dimensión $\leqslant -1$ (\resp $\leqslant 0$, \resp $\leqslant 1$).
\item[(ii)] Para todo morfismo $S' \to S$, para todo esquema~$Y_1$ étale
sobre~$Y$ y para todo torsor~$P$ sobre~$X_1$ bajo~$F_1$, si
$\leftexp{P\mkern-4mu}{F}_1$ denota el grupo torcido de~$F_1$ por~$P$
\textup{\cite[II 4.1.2.3]{XIII.1}}, el morfismo canónico
$$
a_0 \colon g_1^*(f_{1 *}(\leftexp{P\mkern-4mu}{F}_1)) \to f'_{1
*}(\leftexp{P'\mkern-6mu}{F}'_1)
$$
es inyectivo (\resp $a_0$ es biyectivo y el morfismo canónico
$$
a_1 \colon g^*(\R^1 f_* F) \to \R^1 f'_* F'
$$
es
\marginpar{349}
inyectivo, \resp $a_0$ y~$a_1$ son biyectivos).
\item[(ii~bis)] Para todo morfismo $S' \to S$, para todo esquema~$Y_1$
étale sobre~$Y$, para todo torsor~$P$ sobre~$X_1$ bajo~$F_1$ y para todo
torsor~$R$ bajo~$\leftexp{P\mkern-4mu}{F}_1$, el morfismo canónico
$$
\alpha_0 \colon g_1^*(f_{1 *}R) \to f'_{1 *} R'
$$
es inyectivo (\resp $\alpha_0$ es biyectivo, \resp los morfismos
$\alpha_0$ y
$$
\alpha_1 \colon g_1^*(\R^1 f_{1 *}(\leftexp{P\mkern-4mu}{F}_1)) \to \R^1
f'_{1 *}(\leftexp{P'\mkern-6mu}{F}'_1)
$$
son biyectivos).
\end{enumerate}
\end{subproposition}
\endgroup

\subsubsection*{Demostración}
(i)$\To$(ii~bis). Según~\cite[II 4.2.5]{XIII.1}, todo torsor~$R$
bajo~$\leftexp{P\mkern-4mu}{F}_1$ es de la forma $R = Q
\overset{F_1}{\wedge} P^\circ$, donde~$Q$ es un torsor bajo~$F_1$ y
$P^\circ$ el opuesto de~$P$. Se tiene entonces $R' \simeq Q'
\overset{F'_1}{\wedge} {P'}^\circ$. Sea~$\Phi$ el stack de los torsores
bajo~$F$, y sean~$x$, $y$ (\resp $x'$, $y'$) los objetos de la categoría
fibra~$(g^* f_* \Phi)_{Y_1'}$ (\resp $(f'_*
\Phi')_{Y_1'}$) asociados a~$P$, $Q$ (\resp $P'$, $Q'$). Se tiene la
relación
$$
Q\overset{F_1}{\wedge} P^\circ \simeq \SheafHom_{F_1}(P,Q) \text{,}
$$
y de ella se deducen los isomorfismos canónicos
$$
\SheafHom_{Y'_1}(x,y) \simeq g_1^* f_{1 *}(Q
\overset{F_1}{\wedge} P^\circ) \text{,}\quad \SheafHom_{Y_1'}(x',y')
\simeq f'_{1 *}(Q' \overset{F'_1}{\wedge}{P'}^\circ)\text{.}
$$
Por consiguiente, el morfismo~$\alpha_0$ se identifica con el morfismo
$$
\SheafHom_{Y'_1}(x,y) \to \SheafHom_{Y'_1}(x',y') \text{,}
$$
de donde se sigue que, si~$(F,f)$ es cohomológicamente propio
relativamente a~$S$ en dimensión $\leqslant -1$, $\alpha_0$ es inyectivo,
y que, si~$(F,f)$ es cohomológicamente propio en dimensión $\leqslant 0$,
$\alpha_0$ es biyectivo.

Supongamos ahora que~$(F,f)$ sea cohomológicamente propio relativamente
a~$S$ en dimensión $\leqslant 1$, \ie que el morfismo canónico
\marginpar{350}
$$
\varphi \colon g^* f_* \Phi \to f'_* \Phi'
$$
sea una equivalencia. Sea~$G$ el haz de los subgerbes maximales del stack
$f_* \Phi$ \hbox{\cite[III~2.1.8]{XIII.1}}; se tiene entonces un
isomorfismo $G \simeq \R^1 f_* F$. Como~$g^* G$ es el haz de los subgerbes
maximales de~$g^* f_* \Phi$ \cite[III 2.1.5.5]{XIII.2}, el morfismo
$\alpha_1$ se obtiene a partir de~$\varphi|Y'_1$ tomando los haces de los
subgerbes maximales, y por tanto es un isomorfismo.

(ii~bis)$\To$(ii). Basta demostrar que, si los morfismos~$\alpha_0$ son
biyectivos, entonces los morfismos~$a_1$ son inyectivos. Sean~$Y'_1$ un
esquema étale sobre~$Y'$, y~$s$ y~$t$ dos elementos de
$g^*(\R^1 f_* F)(Y'_1)$ que tienen la misma imagen en
$\R^1 f'_* F'(Y'_1)$; demostremos que~$s=t$. La afirmación es local para
la topología étale de~$Y'_1$ y, teniendo en cuenta la definición de la
imagen inversa $g^*(\R^1 f_* F)$, puede suponerse que~$Y'_1$ es la imagen
inversa de un esquema~$Y_1$ étale sobre~$Y$ y que~$s$ y~$t$ provienen de
torsores~$P$ y~$Q$ sobre~$X_1$. La hipótesis sobre~$s$ y~$t$ significa
entonces que las imágenes inversas~$P'$ y~$Q'$ de~$P$ y~$Q$ sobre~$X'_1$
son localmente isomorfas para la topología étale de~$Y'_1$. Si se pone
$R = \SheafHom_{F_1}(P,Q)$, el hecho de que el morfismo
$$
g_1^* f_{1 *} R \to f'_{1 *} R'
$$
sea biyectivo demuestra que~$P$ y~$Q$ son localmente isomorfos para la
topología étale de~$Y_1$, y por tanto que~$s=t$.

(ii)$\To$(i). Para demostrar que~$\varphi$ es fiel (\resp plenamente fiel),
basta demostrar que, si~$Y_1$ es un esquema étale sobre~$Y$, si~$P$, $Q$
son dos torsores sobre~$X_1$ bajo~$F_1$, y si~$x$, $y$ (\resp $x'$, $y'$)
son los objetos de~$(g^* f_* \Phi)_{Y'_1}$ (\resp $(f'_* \Phi')_{Y'_1}$)
asociados a~$P$, $Q$ (\resp $P'$, $Q'$), el morfismo
$$
a \colon \Hom(x,y) \to \Hom(x',y')
$$
es inyectivo (\resp biyectivo). Ahora bien, $a$ se identifica con el
morfismo canónico
$$
\H^0(Y'_1,g^*_1 f_{1 *} (Q \overset{F_1}{\wedge} P^\circ)) \to
\H^0(Y'_1, f'_{1 *}(Q' \overset{F'_1}{\wedge} {P'}^\circ))\text{.}
$$
Si
\marginpar{351}
se tiene $\Hom(x,y) \neq \emptyset$, entonces $Q \overset{F_1}{\wedge}
P^\circ$ es un torsor bajo~$\leftexp{P\mkern-4mu}{F}_1$ localmente trivial
sobre~$Y_1$; de ello se sigue que $f_{1 *}(Q \overset{F_1}{\wedge}
P^\circ)$ es un torsor bajo~$f_{1*}(\leftexp{P\mkern-4mu}{F}_1)$ y que
$g^*_1 f_{1 *} (Q \overset{F_1}{\wedge} P^\circ)$ es un torsor trivial.
El morfismo~$a$ se identifica entonces con el morfismo canónico
$$
\H^0(\ifthenelse{\boolean{orig}}{Y_1}{Y'_1},
g^*_1 f_{1 *}(\leftexp{P\mkern-4mu}{F}_1)) \to \H^0 (Y'_1,
f'_{1 *}(\leftexp{P'\mkern-6mu}{F}'_1)) \text{.}
$$
Lo mismo sucede si~$\Hom(x',y') \neq \emptyset$ y~$a_1$ es inyectivo,
pues entonces $Q' \overset{F'_1}{\wedge} {P'}^\circ$ es trivial, y de la
inyectividad de~$a_1$ se sigue que~$P$ y~$Q$ son localmente isomorfos
sobre~$Y_1$. Se concluye que, si~$a_0$ es inyectivo (\resp si~$a_0$ es
biyectivo y~$a_1$ inyectivo), $\varphi$ es fiel (\resp plenamente fiel).

Queda por demostrar que, si~$a_0$ y~$a_1$ son biyectivos, el funtor
$\varphi$ es esencialmente sobreyectivo. Sean~$Y''$ un esquema étale
sobre~$Y'$, $X'' = X' \times_{Y'} Y''$, y~$P''$ un torsor sobre~$X''$ bajo
$F'' = F'|X''$. Se demostrará que existe un elemento~$x$ de
$(g^* f_* \Phi)_{Y''}$ cuya imagen en~$(f'_* \Phi')_{Y''}$ es isomorfa
a~$P''$. Sea~$p''$ la clase de~$P''$. Como~$a_1$ es sobreyectivo, pueden
hallarse un morfismo étale sobreyectivo $Y''_1 \to Y''$, un morfismo étale
$Y_1 \to Y$ tal que exista un morfismo $Y''_1 \to Y'_1$, y un torsor~$P_1$
sobre~$X_1$ bajo~$F_1$ cuya imagen inversa~$P''_1$ sobre~$X''_1$ sea
isomorfa a la imagen inversa de~$P''$. Utilizando que~$\varphi$ es
plenamente fiel, se ve que el objeto~$x_1$ de~$(g^* f_*
\Phi)_{Y''_1}$ que corresponde a~$P''_1$ está provisto de un dato de
descenso relativamente a~$Y''_1 \to Y''$, y por tanto proviene de un
elemento~$x$ de~$(g^* f_* \Phi)_{Y''}$. Como la imagen de~$x$ en
$(f'_* \Phi')_{Y''}$ es~$P''$, esto demuestra que~$\varphi$ es
esencialmente sobreyectivo y concluye la demostración.

\setcounter{proposition}{3}% El ejemplo siguiente ocupa el ítem XIII.1.4 de la secuencia fuente.
\begin{exemple}
\label{XIII.1.4}
Sea $f \colon X \to Y$ un morfismo propio. Se sigue de~\cite[VII
2.2.2]{XIII.2} que, para todo stack ind-finito~$F$ sobre~$X$, el par
$(F,f)$ es cohomológicamente propio (relativamente a~$Y$) en dimensión
$\leqslant 1$. En particular, para todo haz de conjuntos (\resp todo haz de
grupos, \resp todo haz de grupos ind-finito) $F$ sobre~$X$, $(F,f)$ es
cohomológicamente
\marginpar{352}
propio en dimensión $\leqslant 0$ (\resp en dimensión $\leqslant 0$,
\resp en dimensión $\leqslant 1$).
\end{exemple}


\begin{remarques}
\label{XIII.1.5}
a) Sea~$F$ un haz de grupos sobre~$X$ tal que~$(F,f)$ sea
cohomológicamente propio relativamente a~$S$ en dimensión
$\leqslant -1$ (\resp $\leqslant 0$). Si se considera~$F$ como haz de
conjuntos, $(F,f)$ es cohomológicamente propio relativamente a~$S$ en
dimensión $\leqslant -1$ (\resp $\leqslant 0$), pero el recíproco es
falso.

Sea, por ejemplo, $Y$ el espectro de un anillo de valoración discreta
estrictamente local, con punto cerrado~$t$ y punto genérico~$s$;
sea $f\colon X\to Y$ un esquema no vacío sobre~$Y$ cuya fibra cerrada
sea vacía; sea~$F$ un haz constante no trivial de grupos sobre~$X$; y
sea~$P$ un torsor bajo~$F$ tal que
$\H^0(X_s,\leftexp{P\mkern-4mu}{F}|X_s)=1$. Entonces
$(\leftexp{P\mkern-4mu}{F},f)$ es cohomológicamente propio relativamente
a~$Y$ en dimensión $\leqslant -1$ cuando se considera
$\leftexp{P\mkern-4mu}{F}$ como haz de conjuntos. Si se considera
$\leftexp{P\mkern-4mu}{F}$ como haz de grupos, se tiene un isomorfismo
${\vphantom{F}}^{P^\circ\mkern-7mu}(\leftexp{P\mkern-4mu}{F})\simeq F$;
como el morfismo canónico
$$
\H^0(X,F)\to\H^0(X_t,F|X_t)=1
$$
no es inyectivo, esto demuestra que
$(\leftexp{P\mkern-4mu}{F},f)$ no es cohomológicamente propio
relativamente a~$Y$ en dimensión $\leqslant -1$.

b) Supongamos que~$f$ sea coherente (\ie cuasicompacto y
cuasiseparado). Sea~$F$ un stack sobre~$X$. Para todo punto geométrico
$\overline{y}$ de~$Y'$, se denota por~$\overline{Y}$ (\resp $\overline{Y'}$)
la localización estricta de~$Y$ (\resp $Y'$) en~$\overline{y}$, y se pone
$\overline{X}=X\times_Y\overline{Y}$,
$\overline{X'}=X'\times_{Y'}\overline{Y'}$, etc. Para que~$(F,f)$ sea
cohomológicamente propio relativamente a~$S$ en dimensión
$\leqslant -1$ (\resp $\leqslant 0$, \resp $\leqslant 1$), es necesario
y suficiente que, para todo $S$-esquema~$S'$ y todo punto geométrico
$\overline{y}$ de~$Y'$, el funtor canónico
$$
\overline{F}(\overline{X})\to\overline{F'}(\overline{X'})
$$
sea fiel (\resp plenamente fiel, \resp una equivalencia).

En
\marginpar{353}
efecto, si~$S'$ es un $S$-esquema, para que el funtor
$$
g^*f_*F\to f'_*F'
$$
sea fiel (\resp plenamente fiel, \resp una equivalencia), es necesario
y suficiente que lo sea el funtor inducido sobre las fibras en los
distintos puntos geométricos $\overline{y}'$ de~$\overline{Y'}$
\cite[III 2.1.5.9]{XIII.2}. La afirmación resulta, pues, del cálculo de
las fibras geométricas de la imagen directa de un stack por un morfismo
coherente \cite[VII 2.1.5]{XIII.2}.

c) Sea~$F$ un stack sobre~$X$. El hecho de que~$(F,f)$ sea
cohomológicamente propio relativamente a~$S$ en dimensión
$\leqslant -1$ (\resp $\leqslant 0$, \resp $\leqslant 1$) es local
sobre~$Y$ para la topología étale.

Sea~$S'$ un $S$-esquema, y sea~$F'$ la imagen inversa de~$F$ sobre~$X'$
(\cf \eqref{eq:XIII.1.0.1}). Si~$(F,f)$ es cohomológicamente propio
relativamente a~$S$ en dimensión $\leqslant 1$, también lo es~$(F',f')$.
Pero, si~$(F,f)$ es cohomológicamente propio relativamente a~$S$ en
dimensión $\leqslant -1$ (\resp $\leqslant 0$), no necesariamente sucede
lo mismo con~$(F',f')$.

Sea, por ejemplo, $S'$ un anillo de valoración discreta,
$f'\colon E_{S'}\to S'$ el espacio afín sobre~$S'$, $x$ un punto cerrado
de~$E_{S'}$ situado sobre el punto genérico de~$S'$, y~$F'$ el haz de
conjuntos sobre~$E_{S'}$ cuya restricción a~$E_{S'}-\{x\}$ es el haz
constante con un elemento y cuya fibra en un punto geométrico situado
sobre~$x$ tiene dos elementos. Entonces~$(F',f')$ no es
cohomológicamente propio relativamente a~$S'$ en dimensión
$\leqslant -1$. Sean $S=S'[Z]$, $f\colon E_S\to S$ el espacio afín
sobre~$S$ y~$T$ una parte cerrada de~$X=E_S$ que no interseca el cerrado
$Z=0$ y tal que~$f(T)$ contiene el punto genérico de~$S$.
Sean~$G$ la imagen inversa de~$F'$ sobre~$X$ y~$F$ el haz sobre~$X$
obtenido prolongando~$G|X-T$ por el vacío. Entonces~$(F,f)$ es
cohomológicamente propio relativamente a~$S$ en dimensión
$\leqslant -1$, pero deja de serlo tras el cambio de base $S'\to S$
definido por~$Z=0$.

d)
\marginpar{354}
Sea~$F$ un stack sobre~$X$ tal que~$(F,f)$ sea cohomológicamente propio
relativamente a~$Y$ en dimensión $\leqslant -1$ (\resp $\leqslant 0$,
\resp $\leqslant 1$). Entonces, para todo punto geométrico
$\overline{y}$ de~$Y$, el funtor canónico
$$
(f_*F)_{\overline{y}}\to F(X_{\overline{y}})
$$
es fiel (\resp plenamente fiel, \resp una equivalencia de categorías).
\end{remarques}


\begin{proposition}
\label{XIII.1.6}
Sean $f\colon X\to Y$ y $g\colon Y\to Z$ dos $S$-morfismos,
y sea~$\Phi$ un stack sobre~$X$.
\begin{enumerate}
\item[1)] Supongamos que $(\Phi,f)$ y $(f_*\Phi,g)$ son
cohomológicamente propios relativamente a~$S$ en dimensión
$\leqslant -1$ (\resp $\leqslant 0$, \resp $\leqslant 1$). Entonces
lo mismo ocurre con~$(\Phi,gf)$.
\item[2)] Supongamos que $(\Phi,gf)$ es cohomológicamente propio
relativamente a~$S$ en dimensión $\leqslant -1$ (\resp que $(\Phi,gf)$
es cohomológicamente propio relativamente a~$S$ en dimensión
$\leqslant 0$ y $(\Phi,f)$ cohomológicamente propio relativamente
a~$S$ en dimensión $\leqslant -1$, \resp que $(\Phi,gf)$ es
cohomológicamente propio relativamente a~$S$ en dimensión
$\leqslant 1$ y $(\Phi,f)$ cohomológicamente propio relativamente
a~$S$ en dimensión $\leqslant 0$). Entonces $(f_*\Phi,g)$ es
cohomológicamente propio relativamente a~$S$ en dimensión
$\leqslant -1$ (\resp en dimensión $\leqslant 0$, \resp en dimensión
$\leqslant 1$).
\end{enumerate}
\end{proposition}

Para todo $S$-esquema~$S'$, considérese el diagrama siguiente, en el
que todos los cuadrados son cartesianos:
\begingroup
\renewcommand{\thesubsection}{\theproposition}% La etiqueta fuente pertenece al ítem XIII.1.6.
\begin{equation*}
\label{eq:XIII.1.6.1}
\tag{\thesubsection.1}
\begin{array}{c}
\xymatrix{ X' \ar[r]^{f'} \ar[d]_h & Y'
\ar[r]^{g'} \ar[d]_k & Z' \ar[r] \ar[d]_{m} & S' \ar[d] \\ X \ar[r]^f
& Y \ar[r]^g & Z \ar[r] & \,S. }
\end{array}
\end{equation*}
\endgroup
\ifthenelse{\boolean{orig}}
{el\space}
{El\space}%
morfismo canónico
$$
m^*(g_*f_*\Phi)\to g'_*f'_*(h^*\Phi)
$$
se identifica
\marginpar{355}
con el compuesto de los morfismos canónicos
$$
m^*(g_*f_*\Phi)\lto{i}g'_*(k^*f_*\Phi)
\lto{j}g'_*f'_*(h^*\Phi).
$$
\begin{enumerate}
\item[1)] La hipótesis implica que~$i$ y~$j$ son fieles
(\resp plenamente fieles, \resp equivalencias); por tanto, lo mismo
ocurre con~$ji$.
\item[2)] La hipótesis implica que~$ji$ es fiel (\resp que~$ji$ es
\ifthenelse{\boolean{orig}}
{plenamnete\space}
{plenamente\space}%
fiel y~$j$ fiel, \resp que~$ji$ es una equivalencia y~$j$ plenamente
fiel); de ello resulta que~$i$ es fiel (\resp plenamente fiel, \resp una
equivalencia).
\end{enumerate}


\begin{corollaire}
\label{XIII.1.7}
Sean $f\colon X\to Y$ y $g\colon Y\to Z$ dos $S$-morfismos, y sea~$F$
un haz de grupos sobre~$X$. Supongamos que $(F,gf)$ es cohomológicamente
propio relativamente a~$S$ en dimensión $\leqslant -1$ (\resp que
$(F,gf)$ es cohomológicamente propio relativamente a~$S$ en dimensión
$\leqslant 0$ y que $(F,f)$ es cohomológicamente propio relativamente
a~$S$ en dimensión $\leqslant -1$). Entonces $(f_*F,g)$ es
cohomológicamente propio relativamente a~$S$ en dimensión
$\leqslant -1$ (\resp en dimensión $\leqslant 0$).
\end{corollaire}

Retomemos las notaciones de~\eqref{eq:XIII.1.6.1} y, para todo esquema
$Y_1$ étale sobre~$Y$, denotemos por~$f_1$ y~$F_1$ las respectivas
imágenes inversas de~$f$ y~$F$ por el morfismo $Y_1\to Y$. Sean~$\Phi$
el stack de los torsores bajo~$F$ y~$\Psi$ el stack de los torsores
bajo~$f_*F$. Se tiene un funtor canónico
$$
\varphi\colon\Psi\to f_*\Phi\text{,}
$$
obtenido asociando a todo esquema~$Y_1$ étale sobre~$Y$ y a todo torsor
$P$ sobre~$Y_1$ bajo~$f_{1*}F_1$ el torsor~$\tilde{P}$ sobre~$X_1$
obtenido a partir de~$f_1^*P$ por extensión del grupo estructural a lo
largo de $f_1^*f_{1*}F_1\to F_1$. El funtor~$\varphi$ es plenamente
fiel. En efecto, si~$P$ y~$Q$ son dos torsores sobre~$Y_1$ bajo
$f_{1*}F_1$, se tiene un morfismo canónico
$$
\SheafIsom_{f_{1*}F_1}(P,Q)\to
f_{1*}(\SheafIsom_{F_1}(\tilde{P},\tilde{Q}))
$$
que
\marginpar{356}
es un isomorfismo porque lo es localmente. De ello se deduce que el
morfismo canónico
$$
\Isom_{f_{1*}F_1}(P,Q)\to\Isom_{F_1}(\tilde{P},\tilde{Q})
$$
es un isomorfismo y, por tanto, que~$\varphi$ es plenamente fiel.

Se tiene un diagrama conmutativo
$$
\ifthenelse{\boolean{orig}}
{\xymatrix{ g'_*k^*\Psi \ar[r]^i \ar[d]_{g'_*k^*(\varphi)} &
m^*(g_*\Psi) \ar[d]^{m^*g_*(\varphi)} \\ g'_*k^*(f_*\Phi) \ar[r]^j &
m^*(g_*f_*\Phi)\text{,} }}
{\xymatrix{ m^*(g_*\Psi) \ar[r]^i \ar[d]_{m^*g_*(\varphi)} &
g'_*k^*\Psi \ar[d]^{g'_*k^*(\varphi)} \\ m^*(g_*f_*\Phi) \ar[r]^j &
g'_*k^*(f_*\Phi)\text{,} }}
$$
donde~$i$ y~$j$ son los morfismos de cambio de base. De
\Ref{XIII.1.6}~2) se deduce que~$j$ es fiel (\resp plenamente fiel).
Como $g'_*k^*(\varphi)$ y $m^*g_*(\varphi)$ son plenamente fieles, del
diagrama anterior se deduce que~$i$ es fiel (\resp plenamente fiel).


\begin{corollaire}
\label{XIII.1.8}
Sean $f\colon X\to Y$ un $S$-morfismo coherente, $g\colon Y\to Z$ un
$S$-morfismo propio y~$\Phi$ un stack ind-finito sobre~$X$
\textup{\cite[VII 2.2.1]{XIII.2}}. Supongamos que $(\Phi,f)$ es
cohomológicamente propio relativamente a~$S$ en dimensión
$\leqslant -1$ (\resp $\leqslant 0$, \resp $\leqslant 1$); entonces
lo mismo sucede con~$(\Phi,gf)$.
\end{corollaire}

Como $f$ es coherente, $f_*\Phi$ es un stack ind-finito (SGA~4 IX
1.6~(ii)). Por tanto, el corolario se deduce de~\Ref{XIII.1.6}~1)
y de~\Ref{XIII.1.4}.


\begin{corollaire}
\label{XIII.1.9}
Sean $f\colon X\to Y$ un $S$-morfismo entero y $g\colon Y\to Z$ un
$S$-morfismo. Si~$F$ es un haz de conjuntos sobre~$X$, entonces
$(f_*F,g)$ es cohomológicamente propio relativamente a~$S$ en dimensión
$\leqslant -1$ (\resp $\leqslant 0$) si y solo si lo mismo sucede con
$(F,gf)$. Si~$F$ es un haz de grupos sobre~$X$, entonces $(f_*F,g)$ es
cohomológicamente propio relativamente a~$S$ en dimensión
$\leqslant -1$ (\resp $\leqslant 0$, \resp $\leqslant 1$) si y solo si
lo mismo sucede con~$(F,gf)$.
\end{corollaire}

La afirmación
\marginpar{357}
relativa al caso de un haz de conjuntos se deduce de~\Ref{XIII.1.6} y
del hecho de que $(F,f)$ es cohomológicamente propio relativamente a~$S$
en dimensión $\leqslant 0$. Sean~$F$ un haz de grupos sobre~$X$ y~$\Phi$
el stack de los torsores bajo~$F$. Según SGA~4 VIII~5.8, todo torsor
bajo~$F$ es localmente trivial sobre~$Y$. De ello se sigue que el stack
$f_*\Phi$ es equivalente al stack de los torsores bajo~$f_*F$; la
equivalencia se obtiene asociando a todo esquema~$Y_1$ étale sobre~$Y$ y
a todo torsor~$P$ sobre~$X_1=X\times_Y Y_1$ bajo~$F|X_1$ el torsor
$f_*P$ bajo~$f_*F|Y_1$. Como $(F,f)$ es cohomológicamente propio
relativamente a~$S$ en dimensión $\leqslant 1$, el corolario se deduce
de~\Ref{XIII.1.6}.


\enlargethispage{.5cm}
\skpt
\begingroup
\shorthandoff{<>}% Babel español trata de otro modo los delimitadores de curvatura de Xy-pic.
\setcounter{secnumdepth}{3}
\setcounter{subsubsection}{0}
\renewcommand{\thesubsubsection}{\arabic{section}.\arabic{proposition}.\arabic{subsubsection}}
\begin{definitions}
\label{XIII.1.10}

\subsubsection{}
\label{XIII.1.10.1}
Sea $E$ una categoría y consideremos un diagrama
$$
\xymatrix@C=.6cm{ \Phi \ar[r]^p & \Phi_1 \ar@<2pt>[r]^{p_1}
\ar@<-2pt>[r]_{p_2} & \Phi_2 } \text{,}
$$
donde $\Phi$, $\Phi_1$, $\Phi_2$ son categorías fibradas sobre~$E$ y las
flechas son morfismos de categorías fibradas, y sea
$$
a \colon p_1 p \isomto p_2 p
$$
un isomorfismo de funtores.

Se dice que el diagrama anterior es exacto si se satisface la condición
siguiente:

a) Para todo par de objetos $x$, $y$ de~$\Phi$ y todo morfismo $f
\colon p(x) \to p(y)$ tal que $p_1(f) = p_2(f)$ (identificando $p_1p$ y
$p_2p$ mediante~$a$), existe un único morfismo $g \colon x \to y$ tal que
$p(g)=f$.

\subsubsection{}
\label{XIII.1.10.2}
Consideremos el diagrama
\ifthenelse{\boolean{orig}}
{$$
\xymatrix{ \Phi \ar[r]^p & \Phi_1 \ar@<0.6ex>[r]^{p_1}
\ar@<-0.6ex>[r]_{p_2} & \Phi_2 \ar@<2ex>[r]^{p_{23}} \ar[r]^{p_{31}}
\ar@<-0.6ex>[r]_{p_{12}} & \Phi_3 } \text{,}
$$}
{$$
\xymatrix@C=.5cm{ \Phi \ar[r]^p & \Phi_1 \ar@<0.6ex>[rr]^{p_1}
\ar@<-0.6ex>[rr]_{p_2} && \Phi_2 \ar@<2ex>[rrr]^{p_{23}} \ar[rrr]^{p_{31}}
\ar@<-0.6ex>[rrr]_{p_{12}} &&& \Phi_3 } \text{,}
$$}
donde
\marginpar{358}
$\Phi$, $\Phi_i$, $1 \leqslant i \leqslant 3$, son categorías fibradas
sobre~$E$ y las flechas son morfismos de categorías fibradas. Supongamos
dados unos isomorfismos de funtores
$$
a \colon p_1 p \isomto p_2 p
$$
$$
a_1 \colon p_{31} p_2 \isomto p_{12} p_1 \text{,}\quad a_2 \colon p_{12}
p_2 \isomto p_{23} p_1 \text{,}\quad a_3 \colon p_{23} p_2 \isomto
p_{31}p_1
$$
tales que el diagrama siguiente sea conmutativo:
$$
\xymatrix@C=1.2cm{ p_{23} p_1 p \ar[r]^{\id.a} & p_{23} p_2 p \ar[r]^{a_3.\id}
& p_{31} p_1 p \ar[d]^{\id.a} \\ p_{12} p_2 p \ar[u]^{a_2.\id} &
p_{12} p_1 p \ar[l]^{\id.a} & p_{31} p_2 p \ar[l]^{a_1.\id} } \text{.}
$$
Identificamos $p_1 p$ con $p_2 p$, $p_{31} p_2$ con $p_{12} p_1$, etc.

Se dice que el diagrama anterior es exacto si se satisfacen las condiciones
siguientes:
\begin{enumerate}
\item[a)] Análoga a la condición a) de~\Ref{XIII.1.10.1}.
\item[b)] Para todo objeto $x_1$ de~$\Phi_1$ y todo isomorfismo
$u \colon p_1(x_1) \isomto p_2(x_1)$ tal que
\begin{equation*}
\label{eq:XIII.1.10.2.1}
\tag{\thesubsubsection.1} { \ifthenelse{\boolean{orig}}
{p_{23}(u) p_{31}(u)}
{p_{31}(u) p_{23}(u)} = p_{12}(u)^{-1}
\text{,} }
\end{equation*}
existe un objeto $x$ de~$\Phi$ y un isomorfismo $i
\colon p(x) \isomto x_1$ que hace conmutativo el diagrama
\begin{equation*}
\label{eq:XIII.1.10.2.2}
\tag{\thesubsubsection.2}
\begin{array}{c}
{ \xymatrix{ p_1p(x) \ar@{=}[r]
\ar[d]_{p_1(i)} & p_2p(x) \ar[d]^{p_2(i)} \\ p_1(x_1) \ar[r]^{\sim}_u
& p_2(x_1) } }
\end{array}
\end{equation*}
\end{enumerate}

\subsubsection{}
\label{XIII.1.10.3}
Se define de manera evidente la noción de morfismo de diagramas exactos de
categorías fibradas sobre una categoría~$E$.
\end{definitions}

\subsubsection{}
\label{XIII.1.10.4}
Utilizaremos
\marginpar{359}%
en particular la noción de diagrama exacto cuando~$E$ es un sitio y
$\Phi$, $\Phi_i$, $1 \leqslant i \leqslant 3$, son stacks sobre~$E$.

Sean $f \colon E \to E'$ un morfismo de sitios y
\ifthenelse{\boolean{orig}}
{\begin{equation*}
\label{eq:XIII.1.10.4.1}
\tag{\thesubsubsection.1} { \xymatrix{ \Phi \ar[r]^p & \Phi_1
\ar@<0.6ex>[r]^{p_1} \ar@<-0.6ex>[r]_{p_2} & \Phi_2
\ar@<2ex>[r]^{p_{23}} \ar[r]^{p_{31}} \ar@<-0.6ex>[r]_{p_{12}} &
\Phi_3 } }
\end{equation*}}
{\begin{equation*}
\label{eq:XIII.1.10.4.1}
\tag{\thesubsubsection.1} { \xymatrix@C=.5cm{ \Phi \ar[r]^p & \Phi_1
\ar@<0.6ex>[rr]^{p_1} \ar@<-0.6ex>[rr]_{p_2} && \Phi_2
\ar@<2ex>[rrr]^{p_{23}} \ar[rrr]^{p_{31}} \ar@<-0.6ex>[rrr]_{p_{12}} &&&
\Phi_3 } }
\end{equation*}}
un diagrama exacto de stacks sobre~$E$. Por imagen directa se obtiene un
diagrama
$$
\xymatrix@C=.5cm{ f_* \Phi \ar[r] & f_* \Phi_1 \ar@<2pt>[r]
\ar@<-2pt>[r] & f_* \Phi_2 \ar@<4pt>[r]\ar[r] \ar@<-4pt>[r] &
f_* \Phi_3 }
$$
que evidentemente es exacto.

Si $h \colon E'' \to E$ es un morfismo de sitios, se tiene asimismo un
diagrama exacto
\ifthenelse{\boolean{orig}}
{$$
\xymatrix{ h^* \Phi \ar[r]^{p''} & h^* \Phi_1
\ar@<0.6ex>[r]^{p''_1} \ar@<-0.6ex>[r]_{p''_2} & h^* \Phi_2
\ar@<2ex>[r]^{p''_{23}} \ar[r]^{p''_{31}} \ar@<-0.6ex>[r]_{p''_{12}} &
h^* \Phi_3 }
$$}
{$$
\xymatrix@C=.5cm{ h^* \Phi \ar[r]^{p''} & h^* \Phi_1
\ar@<0.6ex>[rr]^{p''_1} \ar@<-0.6ex>[rr]_{p''_2} && h^* \Phi_2
\ar@<2ex>[rrr]^{p''_{23}} \ar[rrr]^{p''_{31}} \ar@<-0.6ex>[rrr]_{p''_{12}} &&&
h^* \Phi_3 }
$$}

Verifiquemos primero la condición~a) de~\Ref{XIII.1.10.2}. Sean
$F''$ un objeto de~$E''$, $x''$ e~$y''$ dos objetos de~$(h^*
\Phi)_{F''}$, $x''_1$ e~$y''_1$ sus imágenes respectivas en
$h^* \Phi_1$, y $x''_2$, $y''_2$ sus imágenes en~$h^*
\Phi_2$. Sea $u''_1 \colon x''_1 \to y''_1$ un morfismo tal que
$p''_1(u''_1) = p''_2(u''_1)$, y demostremos que~$u''_1$ proviene de un
único morfismo $u'' \colon x'' \to y''$. Como la cuestión es local
sobre~$F''$, podemos suponer que se tienen un objeto $F_1$ de~$E$, un
morfismo de~$F''$ en la imagen inversa $F''_1$ de~$F_1$ por~$h$, y unos
objetos $x$, $y$ de~$\Phi_{F_1}$ cuyas imágenes inversas sobre~$F''$
son $x''$ e~$y''$. Sean $x_1$, $y_1$ (\resp $x_2$, $y_2$) las imágenes
de~$x$, $y$ en~$\Phi_1$ (\resp en~$\Phi_2$). Podemos suponer que
$u''_1$ proviene de un morfismo $u_1 \colon x_1 \to y_1$ tal que
$p_1(u_1) = p_2(u_1)$. Por la exactitud de
\eqref{eq:XIII.1.10.4.1}, se obtiene un único morfismo $u \colon x
\to y$ cuya imagen inversa por~$h$ es el morfismo buscado~$u''$.

La
\marginpar{360}
condición~b) de~\Ref{XIII.1.10.2} se verifica de manera análoga. Sean
$x''_1$ un objeto de~$(h^* \Phi_1)_{F''}$ y $u''
\colon p''_1(x''_1) \to p''_2(x''_1)$ un morfismo que satisface la
relación
$$
\ifthenelse{\boolean{orig}}
{p''_{23}(u'')p''_{31}(u'')}
{p''_{31}(u'')p''_{23}(u'')}=p''_{12}(u'')^{-1} \text{,}
$$
y demostremos que existe un objeto $x''$ de~$(h^* \Phi)_{F''}$ y
un isomorfismo $i'' \colon p''(x'') \simeq x''_1$ que hace conmutativo
un diagrama análogo a~\eqref{eq:XIII.1.10.2.2}. Como la cuestión es
local sobre~$F''$, podemos suponer que se tienen un objeto $F_1$, un
morfismo $F'' \to F''_1$ como antes, y un objeto $x_1$
de~$(\Phi_1)_{F_1}$ cuya imagen inversa en~$(h^* \Phi_1)_{F''}$
es~$x''_1$. De igual manera, podemos suponer que~$u''$ proviene de un
morfismo $u \colon p_1(x_1) \to p_2(x_1)$ que satisface
\ifthenelse{\boolean{orig}}
{\eqref{eq:XIII.1.10.2.2}}
{\eqref{eq:XIII.1.10.2.1}}.
La existencia de un objeto $x$ de~$\Phi_{F_1}$ cuya imagen inversa por~$h$
sea un objeto~$x''$ que satisfaga lo requerido resulta entonces de la
exactitud de~\eqref{eq:XIII.1.10.4.1}.
\endgroup


\begingroup
\shorthandoff{<>}% Babel español trata de otro modo los delimitadores de curvatura de Xy-pic.
\renewcommand{\thesubsection}{\arabic{section}.\arabic{proposition}}
\begin{exemples}
\label{XIII.1.11}
\begin{enumerate}
\item[1)] Sea $f \colon X_1 \to X$ un \emph{morfismo de descenso}
para la categoría de los haces étales sobre esquemas variables (por
ejemplo, un morfismo universalmente submersivo (SGA~4 VIII~9.3)). Sean
$X_2 = X_1 \times_X X_1$, $g \colon X_2 \to X$ la proyección canónica
y~$F$ un haz de conjuntos sobre~$X$. Se deduce entonces de \loccit que
se tiene una sucesión exacta de haces de conjuntos
\begin{equation*}
\label{eq:XIII.1.11.1}
\tag{\thesubsection.1} { \xymatrix@C=.5cm{ F \ar[r] & f_* f^* F
\ar@<2pt>[r] \ar@<-2pt>[r] & g_* g^* F \text{.}}
}
\end{equation*}
Si~$\Phi$ es el stack en categorías discretas asociado a~$F$ y~$\Phi_3$
el stack final sobre~$X$, es decir, el stack cuyas categorías fibra tienen
un único objeto y un único morfismo, la identidad de ese objeto, decir que
la sucesión~\eqref{eq:XIII.1.11.1} es exacta equivale a decir que también
lo es el diagrama de stacks
$$
\xymatrix@C=.5cm{ \Phi \ar[r] & f_* f^* \Phi \ar@<2pt>[r]
\ar@<-2pt>[r] & g_* g^* \Phi \ar[r] \ar@<4pt>[r]
\ar@<-4pt>[r] & \Phi_3 } \text{.}
$$

\item[2)]
Sea
\marginpar{361}
$f \colon X_1 \to X$ un \emph{morfismo de descenso efectivo} para la
categoría de los haces étales sobre esquemas variables (por ejemplo, un
morfismo propio sobreyectivo, un morfismo entero sobreyectivo o un morfismo
fielmente plano y localmente de presentación finita (SGA~4 VIII~9.4)). Sean
$X_2 = X_1 \times_X X_1$, $g \colon X_2 \to X$ la proyección canónica,
$X_3 = X_1 \times_X X_1 \times_X X_1$, $h \colon X_3 \to X$ el morfismo
canónico. Sean $\Phi$ un stack sobre~$X$, $\Phi_1 = f_* f^* \Phi$,
$\Phi_2 = g_* g^* \Phi$ y $\Phi_3 = h_* h^* \Phi$. Se tiene entonces un
diagrama exacto
$$
\xymatrix@C=.5cm{ \Phi \ar[r] & \Phi_1 \ar@<2pt>[r] \ar@<-2pt>[r] &
\Phi_2 \ar[r] \ar@<4pt>[r] \ar@<-4pt>[r] & \Phi_3 \text{,} }
$$
donde las flechas son los morfismos canónicos asociados a las proyecciones.

En efecto, consideremos~$\Phi$ como un stack sobre la categoría~$\Sch_X$
de los esquemas sobre~$X$, provista de la topología étale. Entonces, según
\cite[VII 2.2.8]{XIII.2}, $\Phi$ es también un stack para la topología más
fina sobre~$\Sch_X$ en la que los morfismos cubrientes son los morfismos de
descenso efectivo para la categoría de los haces étales. De ello se sigue
inmediatamente la exactitud del diagrama anterior.
\end{enumerate}
\end{exemples}
\endgroup


\begingroup
\shorthandoff{<>}% Babel español trata de otro modo los delimitadores de curvatura de Xy-pic.
\begin{proposition}
\label{XIII.1.12}
Sean $S$ un esquema y $f\colon X \to Y$ un $S$-morfismo.
\begin{enumerate}
\item[1)] Sea
$$
\xymatrix@C=.5cm{ \Phi \ar[r] & \Phi_1 \ar@<2pt>[r] \ar@<-2pt>[r] &
\Phi_2 }
$$
un diagrama exacto de stacks sobre~$X$. Supongamos que $(\Phi_1, f)$ es
cohomológicamente propio relativamente a~$S$ en dimensión
$\leqslant 0$ y que $(\Phi_2, f)$ es cohomológicamente propio relativamente
a~$S$ en dimensión $\leqslant -1$. Entonces $(\Phi,f)$ es
cohomológicamente propio relativamente a~$S$ en dimensión $\leqslant 0$.

\item[2)] Sea
$$
\xymatrix@C=.5cm{ \Phi \ar[r] & \Phi_1 \ar@<2pt>[r] \ar@<-2pt>[r] &
\Phi_2 \ar@<4pt>[r] \ar[r] \ar@<-4pt>[r] & \Phi_3 }
$$
un diagrama exacto de stacks sobre~$X$. Supongamos que
\ifthenelse{\boolean{orig}}
{$(\Phi,f)$}
{$(\Phi_1,f)$}
es cohomológicamente propio relativamente a~$S$ en dimensión
$\leqslant 1$, que $(\Phi_2, f)$ es cohomológicamente propio
\marginpar{362}
relativamente a~$S$ en dimensión $\leqslant 0$ y que
$(\Phi_3, f)$ es cohomológicamente propio
\ifthenelse{\boolean{orig}}
{}
{relativamente}
a~$S$ en dimensión $\leqslant -1$. Entonces $(\Phi,f)$ es
cohomológicamente propio relativamente a~$S$ en dimensión $\leqslant 1$.
\end{enumerate}
\end{proposition}

Para todo $S$-esquema~$S'$, consideremos el diagrama conmutativo siguiente,
en el que todos los cuadrados son cartesianos:
$$
\xymatrix{ X' \ar[r]^{f'} \ar[d]_h & Y' \ar[r] \ar[d]_g & S' \ar[d] \\
X \ar[r]^f & Y \ar[r] & S }
$$

Demostremos 2), pues la demostración de 1) es análoga. Como los funtores de
imagen directa y de imagen inversa transforman
\ifthenelse{\boolean{orig}}
{diagrama exacto de stacks en diagrama exacto}
{un diagrama exacto de stacks en un diagrama exacto}
\eqref{XIII.1.10.4}, se tiene el morfismo siguiente
de diagramas exactos de stacks:
$$
\xymatrix{ g^* f_* \Phi \ar[r]^\pi \ar[d]_{\varphi} & g^*
f_* \Phi_1 \ar@<2pt>[r]^{\pi_1} \ar@<-2pt>[r]_{\pi_2}
\ar[d]_{\varphi_1} & g^* f_* \Phi_2 \ar@<4pt>[r] \ar[r]
\ar@<-4pt>[r] \ar[d]_{\varphi_2} & g^* f_* \Phi_3
\ar[d]_{\varphi_3} \\ f'_* h^* \Phi \ar[r] & f'_* h^*
\Phi_1 \ar@<2pt>[r] \ar@<-2pt>[r] & f'_* h^* \Phi_2
\ar@<4pt>[r] \ar[r] \ar@<-4pt>[r] & f'_* h^* \Phi_3 }
\text{.}
$$
Por hipótesis, $\varphi_1$ es una equivalencia de categorías,
$\varphi_2$ es plenamente fiel y $\varphi_3$ es fiel. Se deduce, por tanto,
del diagrama anterior que~$\varphi$ es una equivalencia.
\endgroup


\begingroup
\shorthandoff{<>}% Babel español trata de otro modo los delimitadores de curvatura de Xy-pic.
\begin{proposition}
\label{XIII.1.13}
Sea $f \colon X \to Y$ un $S$-morfismo.
\begin{enumerate}
\item[1)] Sea
$$
\xymatrix@C=.5cm{ F \ar[r] & G \ar@<2pt>[r] \ar@<-2pt>[r] & H }
$$
un diagrama exacto de haces de conjuntos sobre~$X$. Supongamos que
$(G,f)$ es cohomológicamente propio relativamente a~$S$ en dimensión
$\leqslant 0$ y que $(H,f)$ es cohomológicamente propio relativamente
a~$S$ en dimensión $\leqslant -1$. Entonces $(F,f)$ es
cohomológicamente propio relativamente a~$S$ en dimensión
$\leqslant 0$.

\item[2)] Sea
\marginpar{363}
$F \to G$ un monomorfismo de haces de grupos sobre~$X$. Si $Y_1$ es un
esquema étale sobre~$Y$, se pone $X_1 = Y_1 \times_Y X$, y se denota por
$f_1$ (\resp $F_1$, \resp $G_1$) la imagen inversa de~$f$ (\resp $F$,
\resp $G$) sobre~$Y_1$ (\cf \Ref{eq:XIII.1.0.2}). Supongamos que
$(G,f)$ es cohomológicamente propio relativamente a~$S$ en dimensión
$\leqslant 0$ (\resp en dimensión $\leqslant 1$) y que, para todo
esquema $Y_1$ étale sobre~$Y$ y todo torsor $Q$ bajo $G_1$,
$(Q/F_1,f_1)$ es cohomológicamente propio relativamente a~$S$ en
dimensión $\leqslant -1$ (\resp en dimensión $\leqslant 0$). Entonces
$(F,f)$ es cohomológicamente propio relativamente a~$S$ en dimensión
$\leqslant 0$ (\resp en dimensión $\leqslant 1$).

\item[3)] Sea $F \to G$ un monomorfismo de haces de grupos sobre~$X$.
Supongamos que $(F,f)$ es cohomológicamente propio relativamente a~$S$
en dimensión $\leqslant 1$ y que $(G,f)$ es cohomológicamente propio
relativamente a~$S$ en dimensión $\leqslant 0$. Entonces, para todo
torsor $Q$ bajo~$G$, $(Q/F,f)$ es cohomológicamente propio relativamente
a~$S$ en dimensión $\leqslant 0$.
\end{enumerate}
\end{proposition}

\subsubsection*{Demostración}
1) Sea $\Phi$ (\resp $\Phi_1$, \resp $\Phi_2$) el stack en categorías
discretas asociado a~$F$ (\resp $G$, \resp $H$), y sea $\Phi_3$ el stack
final sobre~$X$. Se tiene entonces un diagrama exacto
$$
\xymatrix@C=.5cm{ \Phi \ar[r] & \Phi_1 \ar@<2pt>[r] \ar@<-2pt>[r] &
\Phi_2 \ar@<4pt>[r] \ar[r] \ar@<-4pt>[r] & \Phi_3 }
\text{.}
$$
Por hipótesis, $(\Phi_1,f)$ es cohomológicamente propio relativamente
a~$S$ en dimensión $\leqslant 1$ y $(\Phi_2,f)$ es cohomológicamente
propio relativamente a~$S$ en dimensión $\leqslant 0$
\eqref{XIII.1.2}. Como $(\Phi_3,f)$ es evidentemente cohomológicamente
propio relativamente a~$S$ en dimensión $\leqslant -1$, se deduce
de~\Ref{XIII.1.12} que $(\Phi,f)$ es cohomológicamente propio
relativamente a~$S$ en dimensión $\leqslant 1$, \ie que $(F,f)$ es
cohomológicamente propio relativamente a~$S$ en dimensión
$\leqslant 0$.

\smallskip
2) Demostremos primero que, si $(G,f)$ es cohomológicamente propio
relativamente a~$S$ en dimensión $\leqslant 0$ y si los
$(Q/F_1,f_1)$ son cohomológicamente propios relativamente a~$S$ en
\marginpar{364}
dimensión $\leqslant -1$, entonces $(F,f)$ es cohomológicamente propio
relativamente a~$S$ en dimensión $\leqslant 0$. Según
\Ref{XIII.1.3.1}, basta demostrar que, para todo esquema $Y_1$ étale
sobre~$Y$ y todo torsor $P$ sobre~$X_1$ bajo~$F_1$,
$(\leftexp{P\mkern-4mu}{F}_1,f_1)$ es cohomológicamente propio
relativamente a~$S$ en dimensión $\leqslant 0$ cuando se considera
$\leftexp{P\mkern-4mu}{F}_1$ como haz de conjuntos, y que el morfismo
canónico
$$
d \colon g^*(\R^1 f_* F) \to \R^1 f'_* F'
$$
es inyectivo. La primera afirmación se deduce inmediatamente de 1),
pues, si $Q$ denota el torsor obtenido de~$P$ por extensión del grupo
estructural $F_1 \to G_1$, se tiene un isomorfismo
$\leftexp{Q}{G}_1 / \leftexp{P\mkern-4mu}{F}_1 \isomto Q/F_1$.

Demostremos que $d$ es inyectivo. Basta demostrar que, si $Y_1$ es un
esquema étale sobre~$Y$, si $P$ y $\tilde{P}$ son dos torsores bajo
$F_1$ cuyas imágenes inversas $P'$ y~$\tilde{P}'$ sobre~$X'_1$ son
isomorfas, entonces, tras efectuar una extensión étale sobreyectiva
de~$Y_1$, $P$ y~$\tilde{P}$ se vuelven isomorfos. Elijamos un
isomorfismo $p' \colon P' \isomto \tilde{P}'$. Sean $Q$
(\resp~$\tilde{Q}$) el torsor obtenido de~$P$ (\resp $\tilde{P}$) por
extensión del grupo estructural $F_1 \to G_1$. Las imágenes inversas
$Q'$ (\resp $\tilde{Q}'$) de~$Q$ (\resp $\tilde{Q}$) sobre~$X'_1$ se
obtienen de~$P'$ (\resp $\tilde{P}'$) por
\ifthenelse{\boolean{orig}}
{extensión}
{la extensión}
del grupo estructural $F'_1 \to G'_1$; sea
$q' \colon Q' \isomto \tilde{Q}'$ el isomorfismo que se obtiene del
mismo modo a partir de~$p'$. Como $(G,f)$ es cohomológicamente propio
relativamente a~$S$ en dimensión $\leqslant 0$, puede suponerse, tras
efectuar una extensión étale sobreyectiva de~$Y_1$, que $q'$ es la
imagen de un isomorfismo $q \colon Q \isomto \tilde{Q}$. Al torsor~$P$
(\resp~$\tilde{P}$) se asocia una sección $x$ de~$Q/F_1$ (\resp una
sección $\tilde{x}$ de~$\tilde{Q}/F_1$), y, para que $P$ y
$\tilde{P}$ sean isomorfos, es necesario y suficiente que exista un
isomorfismo $Q \isomto \tilde{Q}$ tal que el isomorfismo
$$
e \colon \H^0(X_1,Q/F_1) \to \H^0(X_1,\tilde{Q}/F_1)
$$
que induce lleve $x$ a~$\tilde{x}$. Tomamos el isomorfismo~$q$. Las
secciones $e(x)$ y $\tilde{x}$ de~$\H^0(X_1,\tilde{Q}/F_1)$ tienen la
misma imagen en~$\H^0(X'_1,\tilde{Q}'/F'_1)$. Como
$(\tilde{Q}/F_1,f_1)$ es cohomológicamente propio relativamente a~$S$
en dimensión $\leqslant -1$,
\marginpar{365}
tras efectuar una extensión étale sobreyectiva de~$Y_1$, se tiene
efectivamente $e(x)=\tilde{x}$, lo que demuestra la inyectividad de~$d$.

Para concluir la demostración, queda demostrar que, si $(G,f)$ es
cohomológicamente propio relativamente a~$S$ en dimensión
$\leqslant 1$, y si, para todo esquema $Y_1$ étale sobre~$Y$ y todo
torsor $Q$ sobre~$X_1$ bajo
\ifthenelse{\boolean{orig}}
{$F_1$}
{$G_1$},
$(Q/F_1,f_1)$ es cohomológicamente propio relativamente a~$S$ en
dimensión $\leqslant 0$, entonces el morfismo $d$ es sobreyectivo. Sea
$P'$ un torsor sobre~$X'_1$ bajo~$F'_1$, y sea $Q'$ el torsor bajo
$G'_1$ obtenido a partir de~$P'$ por extensión del grupo estructural.
Dar~$P'$ equivale a dar~$Q'$ y una sección
$x'$ de~$\H^0(X'_1,Q'/F'_1)$. Se deduce entonces de la sobreyectividad
del morfismo
$$
g^*(\R^1 f_* G) \to \R^1 f'_* G'
$$
que, tras efectuar una extensión étale sobreyectiva de~$Y_1$, existe
un torsor $Q$ bajo~$G_1$ cuya imagen inversa sobre~$X'_1$ es isomorfa
a~$Q'$. Utilizando que $(Q/F_1,f_1)$ es cohomológicamente propio
relativamente a~$S$ en dimensión $\leqslant 0$, puede suponerse
asimismo que existe un elemento $x$ de~$\H^0(X_1,Q/F_1)$ cuya imagen
en~$\H^0(X'_1,Q'/F'_1)$ es~$x'$. Los datos de~$Q$ y~$x$ determinan un
torsor $P$ bajo~$F_1$, cuya imagen inversa sobre
\ifthenelse{\boolean{orig}}
{$X_1$}
{$X'_1$}
es isomorfa a~$P'$, lo que demuestra la sobreyectividad de~$d$.

\smallskip
3) Demostremos que $(Q/F,f)$ es cohomológicamente propio relativamente
a~$S$ en dimensión $\leqslant -1$, \ie que, para todo $S$-esquema
$S'$, para todo esquema $Y_1$ étale sobre~$Y$, si $x$ y~$\tilde{x}$
son dos elementos de~$\H^0(X_1,Q_1/F_1)$ cuyas imágenes $x'$ y
$\tilde{x}'$ en~$\H^0(X'_1,Q'_1/F'_1)$ son iguales, entonces, tras
una extensión sobreyectiva de~$Y_1$, se tiene $x=\tilde{x}$. A~$x$
(\resp~$\tilde{x}$) se asocia un torsor $P$ (\resp~$\tilde{P}$) bajo
$F_1$, tal que $Q_1$ se obtiene de~$P$ (\resp~$\tilde{P}$) por la
extensión $F_1 \to G_1$ del grupo estructural. De la igualdad
$x'=\tilde{x}'$ se deduce que existe un isomorfismo
$u' \colon P' \isomto \tilde{P}'$ tal que el isomorfismo inducido
sobre~$Q'_1$ por~$u'$ es la identidad. Como $(F,f)$ es
cohomológicamente propio relativamente a~$S$ en dimensión
$\leqslant 0$, se deduce que,
\marginpar{366}
tras una extensión étale sobreyectiva de~$Y_1$, existe un isomorfismo
$u \colon P \to \tilde{P}$ que levanta~$u'$; el hecho de que $(G,f)$
sea cohomológicamente propio relativamente a~$S$ en dimensión
$\leqslant -1$ implica entonces que $x=\tilde{x}$.

Demostremos que $(Q/F,f)$ es cohomológicamente propio relativamente
a~$S$ en dimensión $\leqslant 0$. Sean $Y''$ un esquema étale sobre
$Y'$ y $x''$ un elemento de~$\H^0(X'',Q''/F'')$. A~$x''$ se asocia un
torsor $P''$ sobre~$X''$ bajo~$F''$. Como $(F,f)$ es
cohomológicamente propio relativamente a~$S$ en dimensión
$\leqslant 1$, pueden encontrarse morfismos étale sobreyectivos
$Y''_1 \to Y''$ y $Y_1 \to Y$, tales que existe un morfismo
$Y''_1 \to Y'_1$, y un torsor $P$ sobre~$X_1$ bajo~$F_1$ cuya imagen
inversa sobre~$X''_1$ es isomorfa a la imagen inversa de~$P''$. Del
hecho de que $(G,f)$ es cohomológicamente propio relativamente a~$S$
en dimensión $\leqslant 0$ se deduce entonces que pueden elegirse
$Y''_1$ y $Y_1$ de modo que el torsor obtenido de~$P$ por extensión
del grupo estructural $F_1 \to G_1$ sea isomorfo a~$Q_1$; a~$P$
corresponde un elemento $x$ de~$\H^0(X_1,Q_1/F_1)$, cuya imagen en
$\H^0(X''_1,Q''_1/F''_1)$
\ifthenelse{\boolean{orig}}
{es isomorfa a}
{es igual a}
la imagen inversa de~$x''$,
lo que concluye la demostración.
\endgroup


\begin{proposition}
\label{XIII.1.14}
Sea $f\colon X \to S$ el morfismo estructural de un $S$-esquema, y sea
$F$ un haz de conjuntos o de grupos sobre~$X$ (\resp un haz de
ind-$\LL$-grupos, donde $\LL$ es un conjunto de números primos).
Supongamos que $F$ es localmente constante, que $(F,f)$ es
cohomológicamente propio en dimensión $\leqslant 0$ (\resp en dimensión
$\leqslant 1$) y que $f$ es localmente $0$-acíclico (\resp localmente
$1$-asférico para $\LL$) \textup{(SGA~4 XV~1.11)}. Entonces, para toda
especialización $\overline{s}_1 \to \overline{s}_2$ de puntos geométricos
de~$S$, el morfismo de especialización \textup{(SGA~4 VIII~7.1)}
$$
a_0 \colon (f_* F)_{\overline{s}_2} \to (f_*
F)_{\overline{s}_1}
$$
es un isomorfismo y, si $F$ es un haz de grupos, el morfismo
$$
a_1 \colon (\R^1 f_* F)_{\overline{s}_2} \to (\R^1 f_*
F)_{\overline{s}_1}
$$
es inyectivo (\resp los morfismos $a_0$ y $a_1$ son isomorfismos).
\end{proposition}

La
\marginpar{367}
demostración se obtiene reproduciendo palabra por palabra la de SGA~4
XVI~2.3, pero reemplazando la expresión \og propio\fg{} por la expresión
\og cohomológicamente propio\fg.

\ifthenelse{\boolean{orig}}
{}
{\enlargethispage{.5cm}}


\begingroup
\shorthandoff{<>}% Babel español trata de otro modo ciertos delimitadores de Xy-pic.
\begin{corollaire}
\label{XIII.1.15}
Sean $f \colon X \to S$ un morfismo, $\Phi$ un stack sobre~$X$ y $\LL$
un conjunto de números primos. Supongamos que, para todo esquema $X_1$
étale sobre~$X$ y para toda pareja de objetos $x$, $y$ de
$\Phi_{X_1}$, el haz $\SheafHom_{X_1}(x,y)$ es localmente constante,
que el haz $\SheafAut_{X_1}(x)$ es un ind-$\LL$-grupo localmente
constante y que el haz de los subgerbes maximales $S\Phi$ de~$\Phi$
\textup{\cite[III 2.1.7]{XIII.1}} es localmente constante. Supongamos que
$(\Phi,f)$ es cohomológicamente propio en dimensión $\leqslant 1$ y que
$f$ es localmente $1$-asférico para $\LL$. Entonces, para toda
especialización $\overline{s}_1 \to \overline{s}_2$ de puntos geométricos
de~$S$, el morfismo de especialización
$$
a \colon (f_* \Phi)_{\overline{s}_2} \to (f_*
\Phi)_{\overline{s}_1}
$$
es una equivalencia de categorías.
\end{corollaire}

Sean $\overline{S}_1$ (\resp $\overline{S}_2$) la localización estricta
de~$S$ en $\overline{s}_1$ (\resp la localización estricta de~$S$ en
$\overline{s}_2$), $\overline{X}_2$, $\overline{\Phi}_2$
(\resp $\overline{X}_1$, $\overline{\Phi}_1$) las imágenes inversas de
\ifthenelse{\boolean{orig}}
{$X_2$, $\Phi_2$}
{$X$, $\Phi$}
sobre~$\overline{S}_2$ (\resp de
\ifthenelse{\boolean{orig}}
{$X_1$, $\Phi_1$}
{$X$, $\Phi$}
sobre~$\overline{S}_1$), y consideremos el cuadrado cartesiano
$$
\xymatrix{ \overline{X}_1 \ar[r]^h \ar[d]_{\overline{f}_1} &
\overline{X}_2 \ar[d]^{\overline{f}_2} \\ \overline{S}_1 \ar[r]^g &
\overline{S}_2 }
\text{.}
$$
Debe demostrarse que el funtor
$$
\varphi \colon \overline{\Phi}_2(\overline{X}_2) \to \overline{\Phi}_1
(\overline{X}_1)
$$
es una equivalencia. El funtor $\varphi$ es plenamente fiel;
\ifthenelse{\boolean{orig}}
{en}
{}
en
\marginpar{368}
efecto, sean $x$, $y$ dos objetos de~$(\overline{\Phi}_2)_{\overline{X}_2}$;
el morfismo canónico
$$
\Hom_{\overline{X}_2}(x,y) \to
\Hom_{\overline{X}_1}(\varphi(x),\varphi(y))
$$
se identifica con el morfismo canónico
$$
\H^0(\overline{X}_2, \SheafHom_{\overline{X}_2}(x,y)) \to
\H^0(\overline{X}_1,h^*(\SheafHom_{\overline{X}_2}(x,y))
\ifthenelse{\boolean{orig}}
{}
{)}
\text{.}
$$
Este morfismo es un isomorfismo por~\Ref{XIII.1.14}.

Demostremos que $\varphi$ es una equivalencia. Sean $x_1$ un objeto
de~$\overline{\Phi}_1(\overline{X}_1)$ y $G_1$ el subgerbe maximal
de~$\overline{\Phi}_1$ generado por $x_1$. El morfismo
$$
\H^0(\overline{X}_2, S\overline{\Phi}_2) \to \H^0(\overline{X}_1,
h^*(S \overline{\Phi}_2)) =
\H^0(\overline{X}_1,S\overline{\Phi}_1)
$$
es biyectivo y existe, por tanto, un subgerbe maximal $G_2$ de
$\overline{\Phi}_2$ tal que $h^* G_2$ es isomorfo a~$G_1$. Basta entonces
probar que el funtor
$$
G_2 \to h_* h^* G_2
$$
es una equivalencia. Pero, bajo esta forma, la cuestión es local para la
topología étale sobre~$\overline{X}_2$. Puede suponerse, pues, que $G_2$ es
un gerbe de torsores bajo el grupo de automorfismos de un objeto de~$G_2$,
caso en el que la afirmación se deduce de~\Ref{XIII.1.14}.
\endgroup


\begin{corollaire}
\label{XIII.1.16}
Las hipótesis son las de~\Ref{XIII.1.14}. Si se supone además que $F$ es
un haz de conjuntos (\resp de ind-$\LL$-grupos) y que $f_* F$ (\resp
$\R^1 f_* F$) es constructible, entonces $f_* F$ (\resp $\R^1 f_* F$)
es localmente constante.
\end{corollaire}

El corolario resulta de~\Ref{XIII.1.14} en virtud de SGA~4 IX~2.11.

\begin{remarque}
\label{XIII.1.17}
Recordemos que la condición de que $f$ sea localmente $0$-acíclico se
satisface si $f$ es plano con fibras separables y
\ifthenelse{\boolean{orig}}
{$X$ e $Y$}
{$X$ y $S$}
son localmente noetherianos (SGA~4 XV~4.1), y que la condición de que $f$ sea localmente
$1$-asférico
\marginpar{369}
para $\LL$ se satisface si $f$ es liso, siendo $\LL$ el conjunto de los
números primos distintos de las características residuales de~$S$
(SGA~4 XV~2.1).
\end{remarque}


\section[Un caso particular: ser cohomológicamente propio]{Un caso particular de la
propiedad de ser cohomológicamente propio: divisores de cruzamientos normales relativos}
\label{XIII.2}

\setcounter{subsection}{-1}

\subsection{}
\label{XIII.2.0}
Sean $R$ un anillo de valoración discreta de cuerpo de fracciones $K$ y $L$
una $K$-álgebra étale; entonces $L$ es producto directo de un número finito
de cuerpos $L_i$, donde $L_i$ es una extensión étale de~$K$. Si $L'_i$
designa la extensión galoisiana generada por $L_i$ dentro de una clausura
algebraica de~$L_i$, se dice que $L$ está \emph{moderadamente ramificada}
\index{moderadamente ramificada (álgebra)|hyperpage}%
sobre~$R$ si los $L'_i$ son extensiones moderadamente ramificadas en el
sentido de~\Ref{X.3}, \ie si un grupo de inercia~$I_i$ de~$L'_i|K$ tiene
orden primo con la característica residual $p$ de~$R$.

Se sabe que $I_i$ es, en todo caso, extensión de un grupo cíclico de orden
primo con~$p$ por un $p$-grupo. (Esto resulta de
\cite[ch\ptbl IV prop\ptbl 7 cor\ptbl 4]{XIII.5} cuando la extensión residual
de~$R$ es separable. La demostración dada en \loccit se extiende al caso
general de la manera siguiente. Retomemos las hipótesis y la notación de
\loccit, pero sin suponer separable la extensión residual. Sea $H_i$ el
subgrupo del grupo de inercia $G_0$ formado por los elementos $s$ de~$G_0$
tales que $s\pi/\pi\in U^i$ para todo uniformizante $\pi$ de~$A_L$. Se
comprueba entonces que $G_0/H_1$ es un grupo de orden primo con~$p$ y que, para
$i\geqslant 1$, los
\ifthenelse{\boolean{orig}}
{$H^i/H^{i+1}$}
{$H_i/H_{i+1}$}
son $p$-grupos, de donde se deduce el resultado anunciado.)

\subsubsection{}
\label{XIII.2.0.1}
Cuando $R$ es estrictamente local se tiene la caracterización sencilla
siguiente: la $K$-álgebra $L$ está moderadamente ramificada sobre~$R$ si y
solo si los $[L_i:K]$ son primos con~$p$. Además, si $L$ está moderadamente
ramificada sobre~$R$, los $L_i$ son extensiones cíclicas de~$K$. En efecto,
cuando $R$ es estrictamente local, $I_i$ es igual al grupo de Galois
\marginpar{370}
de~$L'_i$ sobre~$K$. Como acaba de recordarse, $I_i$ es extensión de un grupo
cíclico de orden primo con~$p$ por un $p$-grupo. Si se supone que $L'_i$ está
moderadamente ramificada sobre~$R$, entonces $I_i$ es un grupo cíclico de
orden primo con~$p$. Se sigue que $[L_i:K]$ es primo con~$p$ y que $L_i=L'_i$.
Recíprocamente, si $[L_i:K]$ es primo con~$p$, $I_i$ no puede contener ningún
$p$-subgrupo normal no trivial; por tanto, $I_i$ es un grupo cíclico de orden
primo con~$p$, lo que demuestra que $L_i$ está moderadamente ramificada
sobre~$R$.

\subsubsection{}
\label{XIII.2.0.2}
Sean $R$ un anillo de valoración discreta de cuerpo de fracciones $K$, $L$
una $K$-álgebra étale, $\overline{R}$ una localización estricta de~$R$,
$\overline{K}$ su cuerpo de fracciones y
$\overline{L}=L\otimes_K\overline{K}$. Entonces $L$ está moderadamente
ramificada sobre~$R$ si y solo si $\overline{L}$ está moderadamente
ramificada sobre~$\overline{R}$. En efecto, se reduce al caso en que $L$ es
un cuerpo. Escribamos $\overline{L}=\prod_i\overline{L}_i$, donde los
$\overline{L}_i$ son extensiones de~$\overline{K}$; si $L'$ es la
extensión galoisiana generada por $L$ y si
\ifthenelse{\boolean{orig}}
{$L'=L'\otimes_K\overline{K}$,}
{$\overline{L'}=L'\otimes_K\overline{K}$,}
se tiene asimismo una descomposición de~$\overline{L'}$ en producto de
cuerpos,
\ifthenelse{\boolean{orig}}
{$L'=\prod_j\overline{L'_j}$}
{$\overline{L'}=\prod_j\overline{L'_j}$}
y cada $\overline{L}_i$ es una subextensión de al menos uno de los
$\overline{L'_j}$. Como $L'$ es una extensión galoisiana de~$K$, los
$\overline{L'_j}$ son extensiones galoisianas de~$\overline{K}$. Supongamos
que $L$ está moderadamente ramificada sobre~$R$; como el grupo de Galois de
\ifthenelse{\boolean{orig}}
{$\overline{L'_j}|K$}
{$\overline{L'_j}|\overline{K}$}
es isomorfo al grupo de inercia de~$L'|K$, los
\ifthenelse{\boolean{orig}}
{$L'_j$}
{$\overline{L'_j}$}
están también moderadamente ramificados sobre~$\overline{R}$ y, por tanto,
lo mismo sucede con los $\overline{L}_i$. Recíprocamente, supongamos que
$\overline{L}$ está moderadamente ramificada sobre~$\overline{R}$. Para cada
$j$, sea $v_j$ la valoración discreta de~$\overline{L'_j}$ que prolonga la
valoración de~$\overline{K}$, y denotemos también por $v_j$ la valoración
inducida sobre~$L'$. Al variar~$j$, $v_j$ recorre el conjunto de las
valoraciones de
\ifthenelse{\boolean{orig}}
{$L$}
{$L'$}
que prolongan la valoración de~$K$. Sean $G=\Gal(L'|K)$,
$H=\Gal(L'|L)$, $I_j$ el grupo de inercia de~$L'|K$ en $v_j$ y $J_j$ el
grupo de inercia de~$L'|L$ en $v_j$. El grupo
\ifthenelse{\boolean{orig}}
{$I_i$}
{$I_j$}
es extensión de un grupo cíclico de orden primo con~$p$ por un $p$-grupo
$P_j$. Como los $\overline{L}_i$ están moderadamente ramificados
sobre~$\overline{R}$, $I_j/J_j$ tiene orden primo con~$p$ y, por tanto,
$P_j\subset J_j$. En consecuencia, el grupo $H$ contiene todos los $P_j$ y,
por tanto, también el grupo $P$ engendrado por
\marginpar{371}
los $P_j$ al variar~$j$. Pero el grupo~$P$ es invariante en~$G$, pues un
automorfismo interior de~$G$ permuta los~$I_j$ entre sí y, por tanto, también
los $P_j$ entre sí. Se sigue que $P$ es un subgrupo de~$H$ normal en~$G$ y,
por consiguiente, puesto que $L'$ es la extensión galoisiana generada por
$L$, que $P=1$, lo que demuestra que $L$ está moderadamente ramificada
sobre~$R$.

Más generalmente, sea $R\to R'$ un homomorfismo de anillos de valoración
discreta tal que la imagen de un uniformizante $\pi$ de~$R$ sea un
uniformizante $\pi'$ de~$R'$ y que la extensión residual $\kres(R')$ sea una
extensión separable de~$\kres(R)$. Sean $K$ el cuerpo de fracciones de~$R$,
$K'$ el cuerpo de fracciones de~$R'$, $L$ una $K$-álgebra étale y
$L'=L\otimes_KK'$. Entonces $L$ está moderadamente ramificada sobre~$R$ si y
solo si $L'$ está moderadamente ramificada sobre~$R'$. En efecto, puede
suponerse que $R$ y $R'$ son estrictamente locales. Según~\Ref{XIII.2.0.1},
basta probar que, cuando $L$ es un cuerpo, también lo es~$L'$. Sean
$\tilde{R}$ la normalización de~$R$ en~$L$, $\tilde{\pi}$ un
uniformizante de~$\tilde{R}$ y
$\tilde{R}'=\tilde{R}\otimes_RR'$. Como la extensión
$\kres(\tilde{R})|\kres(R)$ es radicial y la extensión
$\kres(R')|\kres(R)$ es étale,
$\kres(\tilde{R})\otimes_{\kres(R)}\kres(R')$ es un cuerpo [EGA IV 4.3.2
y 4.3.5]. Esto demuestra que
\ifthenelse{\boolean{orig}}
{$R'$}
{$\tilde{R}'$}
es un anillo local y, como $\pi$ tiene por imagen $\pi'$ en~$R'$, se tiene
\ifthenelse{\boolean{orig}}
{$\kres(\tilde{R}')=R'/(\tilde{\pi})$;}
{$\kres(\tilde{R}')=\tilde{R}'/(\tilde{\pi})$;}
por consiguiente,
\ifthenelse{\boolean{orig}}
{$R'$}
{$\tilde{R}'$}
es un anillo de valoración discreta
\cite[ch\ptbl I \S 2 prop\ptbl 2]{XIII.5} y, por tanto, $L'$ es un cuerpo.

\subsubsection{}
\label{XIII.2.0.3}
Reduciendo al caso estrictamente local, se ve que una subálgebra de un
álgebra moderadamente ramificada está moderadamente ramificada; que el
producto tensorial de dos álgebras moderadamente ramificadas está
moderadamente ramificado; que un álgebra moderadamente ramificada sigue
siéndolo tras extender el anillo de valoración discreta; y que un álgebra
que pasa a estar moderadamente ramificada tras una extensión moderadamente
ramificada está moderadamente ramificada.


\begingroup
\shorthandoff{<>}% Babel español trata de otro modo ciertos delimitadores de Xy-pic.
\subsection{}
\label{XIII.2.1}
Sea
\marginpar{372}
$X$ un $S$-esquema y sea $D$ un divisor $\geq 0$ sobre~$X$. Recordemos
(SGA~5 II~4.2) que se dice que $D$ es \emph{de cruzamientos normales
estrictos relativamente a}~$S$
\index{cruzamientos normales (divisor de)|hyperpage}%
\index{divisor de cruzamientos normales|hyperpage}%
\index{cruzamientos normales estrictos (divisor de)|hyperpage}%
si existe una familia finita $(f_i)_{i\in I}$ de elementos de
$\Gamma(X,\cal{O}_X)$ tal que $D=\sum_{i\in I}\divisor(f_i)$ y se satisface
la condición siguiente:

\setcounter{subsubsection}{-1}
\subsubsection{}
\label{XIII.2.1.0}
Para todo punto $x$ de~$\Supp D$, $X$ es liso sobre~$S$ en~$x$ y, si se
denota por $I(x)$ el conjunto de los $i\in I$ tales que $f_i(x)=0$, el
subesquema $V((f_i)_{i\in I(x)})$ es liso sobre~$S$ de codimensión
\ifthenelse{\boolean{orig}}
{$\mathrm{card}.\,I(x)$}
{$\card I(x)$}
en~$X$.

Se dice que $D$ es \emph{de cruzamientos normales relativamente a}~$S$ si,
localmente sobre~$X$ para la topología étale, es de cruzamientos normales
estrictos.

Sea $D$ un divisor de cruzamientos normales relativamente a~$S$. Pongamos
$Y=\Supp D$, $U=X-Y$, y denotemos por $i\colon U\to X$ la inmersión
canónica. Para todo punto geométrico $\overline{s}$ de~$S$ y para todo punto
maximal $y$ de la fibra geométrica $Y_{\overline{s}}$, el anillo
$R=\cal{O}_{X_{\overline{s}},y}$ es un anillo de valoración discreta.

En lo que resta de este número utilizaremos la definición técnica siguiente:

\begin{subdefinition}
\label{XIII.2.1.1}
Sea $F$ un haz de conjuntos sobre~$U$. Se dice que $F$ está
\emph{moderadamente ramificado sobre}~$X$
\index{moderadamente ramificado (haz)|hyperpage}%
(\emph{a lo largo de}~$D$) \emph{relativamente a}~$S$ si, para todo punto
geométrico $\overline{s}$ de~$S$, se satisface la condición siguiente:

Para todo punto maximal $y$ de~$Y_{\overline{s}}$, la restricción de~$F$ al
cuerpo de fracciones $K$ de
\ifthenelse{\boolean{orig}}
{$\cal{O}_{X_{\overline{s}}}$}
{$\cal{O}_{X_{\overline{s}},y}$}
es representable por el espectro de una $K$-álgebra étale $L$ moderadamente
ramificada sobre
\ifthenelse{\boolean{orig}}
{$\cal{O}_{X_{\overline{s},y}}$}
{$\cal{O}_{X_{\overline{s}},y}$}.
\end{subdefinition}

La mayoría de las veces, cuando no pueda resultar confusión, omitiremos la
mención de~$D$ en la terminología.

\begin{subdefinition}
\label{XIII.2.1.2}
Si
\marginpar{373}
$F$ es un haz de grupos sobre~$U$, moderadamente ramificado sobre~$X$
relativamente a~$S$, se denota por
$$
\H^1_{\tame}(U,F)
$$
\label{indnot:mc}\oldindexnot{$\H^1_{\tame}(U,F)$|hyperpage}%
el subconjunto de~$\H^1(U,F)$ formado por las clases de torsores bajo
$F$ que están moderadamente ramificados sobre~$X$ relativamente a~$S$.
\end{subdefinition}

Sea
$$
\xymatrix{ U \ar[r]^i \ar[d]_g & X \ar[dl]^f \\ T}
$$
un diagrama conmutativo de $S$-esquemas, con $i$ como en~\Ref{XIII.2.1};
se denota por
$$
\R^1_{\tame} g_* F
$$
\label{indnot:md}\oldindexnot{$\R^1_{\tame} g_* F$|hyperpage}%
el haz sobre~$T$ asociado al prehaz $T'\mto\H^1_{\tame}(U',F)$, donde $T'$
recorre los esquemas étale sobre~$T$ y $U'=U\times_TT'$;
$\R^1_{\tame}g_*F$ es un subhaz de~$\R^1g_*F$.

Obsérvese que, si $g$ es coherente, si $\overline{t}$ es un punto geométrico
de~$T$, si $\overline{T}$ es la localización estricta de~$T$ en
$\overline{t}$ y si $\overline{U}=U\times_T\overline{T}$, se tiene un
isomorfismo
\begin{equation*}
\label{eq:XIII.2.1.2.1}
\tag{\thesubsubsection.1} (\R^1_{\tame} g_*F)_{\overline{t}}\simeq
\H^1_{\tame}(\overline{U},\overline{F}).
\end{equation*}

\subsubsection{}
\label{XIII.2.1.3}
Sea $C_t((U,X)/S)$, o simplemente $C_t$,
\label{indnot:me}\oldindexnot{$C_t((U,X)/S)$ o $C_t$|hyperpage}%
la categoría de los recubrimientos étale de~$U$ moderadamente ramificados
sobre~$X$ relativamente a~$S$. Supongamos que $U$ es conexo y sea $a$ un
punto geométrico de~$U$. Sea $\Gamma_t$ el funtor que asocia a todo
recubrimiento étale~$U'$ de~$U$ moderadamente ramificado sobre~$X$
relativamente a~$S$ el conjunto de los puntos geométricos de~$U'$ situados
sobre~$a$. De~\Ref{XIII.2.0} resulta que el par $(C_t,\Gamma_t)$ satisface
los axiomas ($G_1$) a ($G_6$) de~\Ref{V.4}. Por consiguiente, $\Gamma_t$
es representable por un proobjeto denominado \emph{recubrimiento universal
moderadamente ramificado
\marginpar{374}
de~$(U,X)$ relativamente a~$S$, punteado en~$a$}.
\index{moderadamente ramificado (recubrimiento universal)|hyperpage}%
\index{recubrimiento universal moderadamente ramificado|hyperpage}%
El grupo opuesto al grupo de los $U$-automorfismos del recubrimiento universal
moderadamente ramificado se denomina \emph{grupo fundamental moderadamente
ramificado}
\index{grupo fundamental moderadamente ramificado|hyperpage}%
\index{moderadamente ramificado (grupo fundamental)|hyperpage}%
y se denota por
$$
\pi_1^{\tame}((U,X)/S,a)\text{ o simplemente }\pi_1^{\tame}(U,a)\text{ o
incluso }\pi_1^{\tame}(U).
$$
\label{indnot:mf}\oldindexnot{$\pi_1^{\tame}((U,X)/S,a)$ o $\pi_1^{\tame}(U,a)$ o $\pi_1^{\tame}(U)$|hyperpage}%
Es evidentemente un cociente del grupo fundamental $\pi_1(U,a)$
(\Ref{V.6.9}).

\subsubsection{}
\label{XIII.2.1.4}
Sea $F$ un haz de grupos sobre~$U$, sea $P$ un torsor por la derecha bajo~$F$
y sea $Q$ un torsor por la izquierda bajo~$F$. Supongamos que $P$ y $Q$ están
moderadamente ramificados sobre~$X$ relativamente a~$S$. Entonces lo mismo
sucede con~$P\stackrel{F}{\wedge}Q$. En efecto, se reduce a demostrar que,
si $R$ es un anillo de valoración discreta de cuerpo de fracciones $K$, si
$F$ es un esquema en grupos finito étale sobre~$K$ y si $P$ y $Q$ son dos
torsores bajo~$F$ moderadamente ramificados sobre~$R$, entonces lo mismo
sucede con~$P\stackrel{F}{\wedge}Q$. Ahora bien,
$T=P\stackrel{F}{\wedge}Q$ es un cociente de~$P\times_KQ$. Si $L$, $M$ y
$N$ designan las $K$-álgebras que representan, respectivamente, $T$, $P$ y
$Q$, entonces $L$ es una subálgebra de $M\otimes_KN$, y de
\Ref{XIII.2.0.3} resulta que $L$ está moderadamente ramificada sobre~$R$.

De lo anterior se sigue que, si $F$ es un haz de grupos sobre~$U$ y existe
un torsor bajo~$F$ moderadamente ramificado sobre~$X$ relativamente a~$S$,
entonces $F$ está moderadamente ramificado sobre~$X$ relativamente a~$S$.
En efecto, el torsor $P^{\circ}$, opuesto a~$P$, está moderadamente ramificado
sobre~$X$ relativamente a~$S$, pues es isomorfo a~$P$ como haz de conjuntos.
Si $\leftexp{P\mkern-4mu}{F}$ es el grupo torcido de~$F$ por~$P$,
existe un isomorfismo
$$
F\simeq P^{\circ}\stackrel{\leftexp{P\mkern-4mu}{F}}{\wedge}P
$$
y, por consiguiente, $F$ está moderadamente ramificado sobre~$X$
relativamente a~$S$.

Como antes, se ve que, si $F\to F'$ es un morfismo de haces de grupos sobre
$U$, moderadamente ramificados sobre~$X$ relativamente a~$S$, y si $P$ es un
torsor bajo $F$ moderadamente ramificado sobre~$X$ relativamente
\marginpar{375}
a~$S$, entonces el torsor $P'$ obtenido de~$P$ por extensión del grupo
estructural $F\to F'$ está moderadamente ramificado sobre~$X$ relativamente
a~$S$.

En particular, el morfismo canónico
$$
\H^1(U,F)\to\H^1(U,F')
$$
induce, por restricción a $\H^1_{\tame}(U,F)$, un morfismo canónico
$$
\H^1_{\tame}(U,F)\to\H^1_{\tame}(U,F').
$$

\subsubsection{}
\label{XIII.2.1.5}
Sea $S'\to S$ un morfismo, y denotemos por $U'$ (\resp $X'$, etc.) la
imagen inversa de~$U$ (\resp $X$, etc.) sobre~$S'$. Si $F$ es un haz de
conjuntos sobre~$U$ moderadamente ramificado sobre~$X$ relativamente a~$S$,
de la Definición~\Ref{XIII.2.1.1} y de~\Ref{XIII.2.0.3} resulta que $F'$ está
moderadamente ramificado sobre~$X'$ relativamente a~$S'$.

Si ahora $F$ es un haz de grupos sobre~$U$, la imagen inversa sobre~$S'$ de
un torsor bajo $F$ moderadamente ramificado sobre~$X$ relativamente a~$S$ es
un torsor bajo $F'$ moderadamente ramificado sobre~$X'$ relativamente
a~$S'$. En particular, se tiene un funtor canónico
\begin{equation*}
\label{eq:XIII.2.1.5.1}
\tag{\thesubsubsection.1} C_t((U,X)/S)\to C_t((U',X')/S').
\end{equation*}
Supongamos que $U$ y $U'$ son conexos, y sean $a$ un punto geométrico de~$U$
y $a'$ un punto geométrico de~$U'$ situado sobre~$a$; lo anterior proporciona
un morfismo canónico
\begin{equation*}
\label{eq:XIII.2.1.5.2}
\tag{\thesubsubsection.2} \pi_1^{\tame}(U',a')\to\pi_1^{\tame}(U,a).
\end{equation*}

Si $S'\to S$ es un morfismo y $h\colon T'\to T$ es la proyección canónica,
el morfismo
$$
h^*(\R^1g_*F)\to\R^1g'_*F'
$$
se restringe a un morfismo canónico
\begin{equation*}
\label{eq:XIII.2.1.5.3}
\tag{\thesubsubsection.3} h^*(\R^1_{\tame}g_*F)\to
\R^1_{\tame}g'_*F'.
\end{equation*}

\subsubsection{}
\label{XIII.2.1.6}
Sea
\marginpar{376}
$F$ un haz de grupos sobre~$U$, moderadamente ramificado sobre~$X$
relativamente a~$S$. Con la notación de~\Ref{XIII.2.1.2}, se tienen las
sucesiones exactas canónicas:
\begin{equation*}
\label{eq:XIII.2.1.6.1}
\tag{\thesubsubsection.1}
\begin{array}{ccccccc}
1 & \longrightarrow & \H^1(X,i_*F)& \longrightarrow & \H^1_{\tame}(U,F) &
\longrightarrow & \H^0(X,\R^1_{\tame} i_*F)\\[10pt] 1 &
\longrightarrow & \R^1f_*(i_*F) & \longrightarrow & \R^1_{\tame}g_*F &
\longrightarrow & f_*(\R^1_{\tame}i_*F).
\end{array}
\end{equation*}

La primera se obtiene de la sucesión exacta (SGA~4 III~3.2):
$$
1\to\H^1(X,i_*F)\to\H^1(U,F)\to\H^0(X,\R^1i_*F).
$$
En efecto, basta demostrar que la imagen de~$\H^1(X,i_*F)$ en
$\H^1(U,F)$ está contenida de hecho en~$\H^1_{\tame}(U,F)$ y que la imagen
de~$\H^1_{\tame}(U,F)$ en~$\H^0(X,\R^1i_*F)$ está contenida de hecho en
$\H^0(X,\R^1_{\tame}i_*F)$. Ahora bien, la imagen inversa sobre~$U$ de un
torsor bajo $i_*F$ es un torsor bajo $i^*i_*F$ que está evidentemente
moderadamente ramificado sobre~$X$ relativamente a~$S$; por tanto, lo mismo
ocurre tras la extensión del grupo estructural $i^*i_*F\to F$, lo que demuestra
la existencia de la flecha
$\H^1(X,i_*F)\to\H^1_{\tame}(U,F)$. El hecho de que la imagen de
$\H^1_{\tame}(U,F)$ en~$\H^0(X,\R^1i_*F)$ esté contenida en
$\H^0(X,\R^1_{\tame}i_*F)$ resulta inmediatamente de la definición de
$\R^1_{\tame}i_*F$. Esto demuestra la existencia de la primera sucesión
exacta; la segunda se deduce de ella por localización.
\endgroup


\begingroup
\shorthandoff{<>}% Babel español trata de otro modo ciertos delimitadores de Xy-pic.
\subsection{}
\label{XIII.2.2}
Conservemos las notaciones de~\Ref{XIII.2.1}. Vamos a definir una noción de
objeto moderadamente ramificado de un stack~$\Phi$ sobre~$U$ cuando este viene
dado, localmente \emph{sobre}~$X$ y sobre~$S$ para la topología étale, como
\emph{imagen inversa de un stack}~$\Psi$ \emph{sobre}~$S$.

Sea primero $G$ un gerbe sobre~$U$ y supongamos dados un morfismo étale
sobreyectivo $S_1\to S$, un morfismo étale sobreyectivo
$X_2\to X\times_S S_1$, un gerbe trivial $H$ sobre~$S_1$ y un isomorfismo
$$
G|{U_2}\to H|{U_2},
$$
donde
\marginpar{377}
\ifthenelse{\boolean{orig}}
{$U_1=U\times_X X_1$, $U_2=U\times_X X_2$.}
{$X_1=X\times_S S_1$, $U_1=U\times_X X_1$,
$U_2=U\times_X X_2$.}
Al elegir una trivialización de~$H|{X_2}$, el isomorfismo anterior identifica
$G|{U_2}$ con el stack de torsores bajo un haz de grupos~$F$. Se dice que un
objeto $x$ de~$G_U$ está moderadamente ramificado sobre~$X$ relativamente
\ifthenelse{\boolean{orig}}
{a~$S$,}
{a~$S$}
si la restricción de~$x$ a~$U_2$ es un torsor moderadamente ramificado
\ifthenelse{\boolean{orig}}
{sobre~$X$ relativamente a~$S$.}
{sobre~$X_2$ relativamente a~$S_1$.}
Según~\Ref{XIII.2.1.4}, esta noción no depende de la trivialización elegida
de~$H|X_2$.

Sea ahora $\Phi$ un stack sobre~$U$ y supongamos dados un morfismo étale
sobreyectivo $S_1\to S$, un morfismo étale sobreyectivo
$X_2\to X\times_S S_1$, un stack $\Psi$ sobre~$S_1$ y un isomorfismo
$$
i\colon \Phi|U_2\to \Psi|U_2.
$$
Sea $x$ un objeto de~$\Phi_U$, sea $G_x$ el subgerbe maximal de~$\Phi$
generado por~$x$
\ifthenelse{\boolean{orig}}
{[III~2.1.7]}
{\cite[III~2.1.7]{XIII.1}},
y $S\Phi$ el haz de los subgerbes maximales de~$\Phi$. El isomorfismo~$i$
induce un isomorfismo
$$
S\Phi|U_2\to S\Psi|U_2.
$$
De~\Ref{XIII.5.7} resulta que, tras reemplazar si es preciso $S_1$ por una
extensión étale sobreyectiva, existe un único subgerbe maximal~$H$ de~$\Psi$,
que puede suponerse trivial, tal que~$i$ define un isomorfismo
$$
G_x|U_2\to H|U_2.
$$
Se dice que el objeto~$x$ está moderadamente ramificado sobre~$X$
relativamente a~$S$ si lo está como objeto de~$G_x$, provisto del
isomorfismo anterior.

\subsubsection{}
\label{XIII.2.2.1}
Sea $\Phi$ un stack sobre~$U$ dado, localmente sobre~$X$ y~$S$, como imagen
inversa de un stack sobre~$S$, y sea
$$
\xymatrix{ U \ar[r]^i
\ar[d]_g & X \ar[dl]^f \\ T }
$$
un diagrama como en~\Ref{XIII.2.1.2}. Para todo esquema $T'$ étale sobre~$T$,
\ifthenelse{\boolean{orig}}
{si $U'=U\times_TT'$,}
{si $U'=U\times_TT'$ y $X'=X\times_TT'$,}
consideremos el subconjunto
$(g_*^{\tame}\Phi)_{T'}$
\label{indnot:mg}\oldindexnot{$(g_*^{\tame}\Phi)_{T'}$|hyperpage}%
de~$(g_*\Phi)_{T'}=\Phi_{U'}$ formado
\marginpar{378}
por los objetos de~$\Phi_{U'}$ moderadamente ramificados
\index{objeto moderadamente ramificado de un stack|hyperpage}%
\ifthenelse{\boolean{orig}}
{sobre~$X$ relativamente a~$S$.}
{sobre~$X'$ relativamente a~$S$.}
Se denomina \emph{imagen directa
moderadamente ramificada}
\index{imagen directa moderadamente ramificada de un stack|hyperpage}%
\index{moderadamente ramificada (imagen directa de un stack)|hyperpage}%
de~$\Phi$ por~$g$, y se denota por
$$
g_*^{\tame}\Phi
$$
la subcategoría plena de~$g_*\Phi$ cuyos objetos sobre un esquema $T'$ étale
sobre~$T$ son los objetos de~$(g_*^{\tame}\Phi)_{T'}$. Es claro que
$g_*^{\tame}\Phi$ es un substack de~$g_*\Phi$.

\subsubsection{}
\label{XIII.2.2.2}
\ifthenelse{\boolean{orig}}
{}
{Retomemos las notaciones de~\Ref{XIII.2.1.5}. }
Reduciendo al caso de un stack de torsores, se ve que, si $\Phi$ es un stack
sobre~$U$ que, localmente para la topología étale sobre~$S$ y sobre~$X$, es
imagen inversa de un stack sobre~$S$, el morfismo canónico
$$
h^*(g_*\Phi)\to g'_*\Phi'
$$
induce por restricción un morfismo canónico
$$
h^*(g_*^{\tame}\Phi)\to g_*^{\prime\tame}\Phi'
$$

\setcounter{proposition}{2}% La observación siguiente es XIII.2.3; los §§2.0--2.2
\begin{remarques}
\label{XIII.2.3}

a) Si $F$ es un haz de conjuntos localmente constante y constructible
sobre~$U$, para que $F$ esté moderadamente ramificado sobre~$X$ relativamente
a~$S$, basta que la condición de~\Ref{XIII.2.1.1} se satisfaga para los puntos
geométricos de~$S$ situados sobre los puntos maximales de~$S$. Para verlo,
puede suponerse que $D$ es de cruzamientos normales estrictos. El haz~$F$ es
representable por un recubrimiento étale~$V$ de~$U$. Si $\overline{s}$ es un
punto geométrico de~$S$ y $y$ un punto maximal de~$Y_{\overline{s}}$,
\ifthenelse{\boolean{orig}}
{denotemos por $\overline{S}$ la localización estricta de~$S$ en
$\overline{s}$ y por $\overline{X}$ la localización estricta de~$X$ en
$\overline{y}$,}
{elijamos un punto geométrico $\overline y$ situado sobre~$y$; denotemos por
$\overline{S}$ la localización estricta de~$S$ en~$\overline{s}$ y por
$\overline{X}$ la localización estricta de~$X$ en~$\overline{y}$,}
$\overline{U}=U\times_X\overline{X}$ y
$\overline{V}=V\times_X\overline{X}$. Si la condición
de~\Ref{XIII.2.1.1} se satisface en los puntos geométricos situados sobre los
puntos maximales de~$S$, de~\Ref{XIII.5.5} más abajo resulta que
$\overline{V}$ es un recubrimiento de~$\overline{U}$ moderadamente ramificado
sobre~$\overline{X}$ relativamente a~$\overline{S}$; por tanto, $V$ es un
recubrimiento étale de~$U$ moderadamente ramificado sobre~$X$ relativamente
a~$S$.

b)
\marginpar{379}
Sea $F$ un haz de grupos sobre~$U$ moderadamente ramificado sobre~$X$
relativamente a~$S$. Si $\overline{s}$ es un punto geométrico de~$S$ y $y$
un punto maximal de~$Y_{\overline{s}}$, denotemos por $K$ el cuerpo de
fracciones de~$\cal{O}_{X_{\overline{s}},y}$. Supongamos que, para todo punto
$\overline{s}$ y todo punto~$y$, la $K$-álgebra~$L$ cuyo espectro representa
$F|K$ es de rango primo con la característica residual~$p$ de
\ifthenelse{\boolean{orig}}
{$\cal{O}_{X_{\overline{s}}}$.}
{$\cal{O}_{X_{\overline{s}},y}$.}
Diremos a veces, por abuso de lenguaje, que $F$ es primo con las
características residuales de~$S$.
\index{primo con las características residuales (haz)|hyperpage}%
En tal caso, todo torsor~$P$ bajo~$F$ está moderadamente ramificado sobre~$X$
relativamente a~$S$.

En efecto,
\ifthenelse{\boolean{orig}}
{sean}
{elijamos un punto geométrico $\overline y$ situado sobre~$y$; sean}
$\overline{R}$ la localización estricta de
$\cal{O}_{X_{\overline{s}},y}$ en~$\overline{y}$, $\overline{K}$ su cuerpo de
fracciones y $\overline{F}$ la imagen inversa de~$F$ sobre~$\overline{K}$.
Demostremos que puede suponerse constante
\ifthenelse{\boolean{orig}}
{$F$.}
{$\overline{F}$.}
Como $F$ está moderadamente ramificado sobre~$X$ relativamente a~$S$,
$\overline{F}$ es representable por el espectro de una $\overline{K}$-álgebra
$L=\prod L_i$, donde los $L_i$ son extensiones de~$\overline{K}$ de grado
primo con~$p$. Puede encontrarse, por tanto, una extensión~$K'$ de
$\overline{K}$ de grado primo con~$p$ tal que $\overline{F}|K'$ sea un haz
constante. Según~\Ref{XIII.2.0.3}, para demostrar que $P|\overline{K}$ está
moderadamente ramificado sobre~$\overline{R}$, basta ver que $P|K'$ está
moderadamente ramificado sobre la clausura integral de~$\overline{R}$ en
$K'$; esto reduce al caso en que $\overline{F}$ es constante. Supongamos en
adelante que $\overline{F}$ es constante. La $\overline{K}$-álgebra~$H$ que
representa $P|\overline{K}$ es entonces producto de extensiones~$H_i$ de
$\overline{K}$ isomorfas entre sí. Como el rango de~$H$ es primo con~$p$,
también lo es $[H_1:\overline{K}]$, lo que demuestra que $H$ está
moderadamente ramificada
\ifthenelse{\boolean{orig}}
{sobre~$X$ relativamente a~$S$.}
{sobre~$\overline{R}$.}

c) Sean $X$ un esquema regular, $D$ un divisor de cruzamientos normales
sobre~$X$ (SGA~5 I~3.1.5), $U=X-\Supp D$ y $F$ un haz de conjuntos sobre~$U$.
Si $y$ es un punto maximal de~$\Supp D$, denotemos por $K$ el cuerpo de
fracciones de~$\cal{O}_{X,y}$. Se dice que $F$ está
\emph{moderadamente ramificado relativamente a}~$D$
\index{moderadamente ramificado (haz)|hyperpage}%
si, para todo punto maximal $y$ de~$\Supp D$, $F|K$ es representable por el
espectro de una $K$-álgebra moderadamente ramificada sobre~$\cal{O}_{X,y}$.
\end{remarques}
\endgroup


\begingroup
\shorthandoff{<>}% Babel español trata de otro modo ciertos delimitadores de Xy-pic.
\begin{theoreme}
\label{XIII.2.4}
Sean
\marginpar{380}
$f\colon X\to S$ un $S$-esquema, $D$ un divisor sobre~$X$
de cruzamientos normales relativamente a~$S$ \eqref{XIII.2.1},
$Y=\Supp D$, $U=X-Y$, e $i\colon U\to X$ la inmersión canónica. Sea $F$
un haz de conjuntos (\resp de grupos) sobre~$U$ que satisfaga una de las
condiciones siguientes:
\begin{enumerate}
\item[a)] $F$ es, localmente para la topología étale \emph{sobre}
$X$ y sobre~$S$, la imagen inversa de un haz de conjuntos (\resp de un
haz de grupos constructible) sobre~$S$.
\item[b)] $F$ es localmente constante y constructible sobre~$U$, y está
moderadamente ramificado sobre~$X$ relativamente a~$S$.
\end{enumerate}
Entonces se tienen las conclusiones siguientes:
\begin{enumerate}
\item[1)]$(F,i)$ es cohomológicamente propio relativamente a~$S$
en dimensión $\leq 0$ (\resp para todo morfismo $h\colon S'\to S$, si
$i'\colon U'\to X'$ es la imagen inversa de~$i$ sobre~$S'$, si $F'=F|U'$
y si $k=h_{(X)}$, el morfismo canónico
$$
\Psi\colon k^*(\R^1_{\tame}i_*F)\to\R^1_{\tame}i_*'F'
$$
es un isomorfismo).

Si $F$ es un haz de grupos primo con las características
residuales de~$S$
\ifthenelse{\boolean{orig}}
{(\Ref{XIII.2.3}.b))}
{(\Ref{XIII.2.3}.b)}
(moderadamente
ramificado sobre~$X$ relativamente a~$S$), entonces $(F,i)$ es
cohomológicamente propio relativamente a~$S$ en dimensión $\leq 1$.
\item[2)]Si $F$ es un haz de conjuntos (\resp de grupos)
constructible, $i_*F$ (\resp $\R^1_{\tame}i_*F$) es constructible.
\end{enumerate}
\end{theoreme}

\subsubsection*{Demostración}
Para todo $S$-esquema $S'$, consideremos el diagrama siguiente,
en el que todos los cuadrados son cartesianos:
$$
\xymatrix{ U \ar[dd]_-g \ar[dr]^-i & & & U'\ar[lll] \ar[dd]_{\hbox{\raise10mm\hbox{$g'$}}}
\ar[dr]^-{i'}& \\ & X \ar[dl]_{f\mkern-8mu} & & & X'\ar[lll]_{k\hspace*{1cm}} \ar[dl]_{f'\mkern-10mu} \\ S
& & & S'\ar[lll]_-h & & }
$$
Como
\marginpar{381}
la cuestión es local sobre~$X$ para la topología étale, puede suponerse
que $D$ es un divisor de cruzamientos normales estrictos relativamente
a~$S$ \eqref{XIII.2.1}; además, tras restringir si es preciso $X$ a un
entorno de~$Y$, puede suponerse que $X$ es liso sobre~$S$.

\subsubsection*{Demostración de~\Ref{XIII.2.4}~\textup{1)}}
\setcounter{subsection}{4}% Los párrafos numerados pertenecen al Teorema XIII.2.4.
\setcounter{subsubsection}{0}% Comenzar visiblemente por XIII.2.4.1.
\subsubsection{}
\label{XIII.2.4.1}
\emph{Caso de un haz de conjuntos que satisface} a). Puede suponerse que
$F=g^*G$, donde $G$ es un haz sobre~$S$. Se deduce entonces de
(SGA~4 XVI~3.2) que el morfismo canónico
\begin{equation*}
\label{eq:XIII.2.4.1.1}
\tag{\thesubsubsection.1} f^*G\to i_*F
\end{equation*}
es un isomorfismo. Para todo $S$-esquema $h\colon S'\to S$, se tiene
asimismo un isomorfismo $f'^*G'\to i_*'F'$; por consiguiente, el morfismo
canónico
$$
\varphi\colon k^*(i_*F)\to i_*'F'
$$
se identifica con el isomorfismo natural
$$
k^*f^*G\simeq f'^*G'
$$

\subsubsection{}
\label{XIII.2.4.2}
\emph{Caso de un haz de conjuntos que satisface} b). Hay que demostrar
que $\varphi$ es un isomorfismo, y para ello basta comprobarlo en cada
punto geométrico $\overline{x}'$ de~$X'$. Sea $\overline{S}$
(\resp $\overline{X}$, \resp $\overline{S'}$, \resp $\overline{X'}$)
la localización estricta de~$S$ (\resp de~$X$, \resp de~$S'$,
\resp de~$X'$) en $\overline{x}'$, y pongamos
$\overline{U}=U_{(\overline{X})}$,
$\overline{U'}=U_{(\overline{X'})}$, etc. El morfismo
\ifthenelse{\boolean{orig}}
{$\varphi_{\overline{x}}$}
{$\varphi_{\overline{x}'}$}
se identifica con el morfismo canónico
$$
\overline{\varphi}\colon \H^0(\overline{U},\overline{F})\to
\H^0(\overline{U'}, \overline{F'})
$$
Puede encontrarse un recubrimiento principal $V$ de~$\overline{U}$, del
tipo que figura en~\Ref{XIII.5.4}, tal que la imagen inversa de
$\overline{F}$ sobre~$V$ sea un haz constante de valor $C$. Si
$\Pi$ es el grupo de Galois de~$V$ sobre~$\overline{U}$, $\Pi$
actúa sobre~$\overline{F}|V$, y se tiene
\begin{equation*}
\label{eq:XIII.2.4.2.1}
\tag{\thesubsubsection.1}
\ifthenelse{\boolean{orig}}
{\H^0(U,F))\simeq \H^0(V,C_V)^{\Pi}}
{\H^0(\overline{U},\overline{F})\simeq \H^0(V,C_V)^{\Pi}}
\end{equation*}
\label{indnot:mh}\oldindexnot{$\H^0(V,C_V)^{\Pi}$|hyperpage}%
donde el segundo miembro designa el conjunto de los elementos
de~$\H^0(V,C_V)$
\marginpar{382}
invariantes bajo la acción de~$\Pi$. Como
$V'=V\times_{\overline{U}}\overline{U'}$ es un recubrimiento
principal de~$\overline{U'}$, de grupo de Galois $\Pi'\simeq \Pi$, se ve que
el morfismo $\overline{\varphi}$ se obtiene, tomando los invariantes
bajo la acción de~$\Pi$, a partir del morfismo canónico
$$
\H^0(V,C_V)\to \H^0(V',C_{V'}).
$$
Como $V$ y $V'$ son conexos \eqref{XIII.5.4}, este morfismo, y por tanto
también $\overline{\varphi}$, es un isomorfismo.

Obsérvese que, si además $F$ es un haz de grupos y $P$ es un torsor
sobre~$\overline{U}$ bajo $\overline{F}$, moderadamente ramificado
relativamente a~$D$, se deduce de la demostración precedente y
de~\Ref{XIII.2.2} que
$(\leftexp{P}{\overline{F}},\overline{i})$ es cohomológicamente
propio relativamente a~$S$ en dimensión $\leq 0$.

\subsubsection{}
\label{XIII.2.4.3}
\emph{Caso de un haz de grupos.} Para demostrar que $\Psi$ es un
isomorfismo, basta probar que, para todo punto geométrico
$\overline{y}'$ de~$Y'$, el morfismo
$$
\Psi_{\overline{y}'}\colon
(k^*(\R^1_{\tame}i_*F))_{\overline{y}'}\to (\R^1_{\tame}
i_*'F')_{\overline{y}'}
$$
es un isomorfismo. Ahora bien, según~\Ref{XIII.2.1.2},
$\Psi_{\overline{y}'}$ se identifica con el morfismo canónico
$$
\overline{\Psi}\colon
\H^1_{\tame}(\overline{U},\overline{F})\to
\H^1_{\tame}(\overline{U'},\overline{F'}).
$$
Sean $\widetilde{U}$ el recubrimiento universal moderadamente
ramificado de~$\overline{U}$ \eqref{XIII.2.1.3} y $\widetilde{F}$
la imagen inversa de~$\overline{F}$ sobre~$\widetilde{U}$. De
\Ref{XIII.5.7} en el caso a) y de~\Ref{XIII.5.5} en el caso b) resulta
que $\H^1_{\tame}(\overline{U},\overline{F})$ se identifica con el
subconjunto de~$\H^1(\overline{U},\overline{F})$ formado por las clases
de torsores bajo~$\overline{F}$ cuya imagen inversa sobre~$\widetilde{U}$
es trivial. Por otra parte, un razonamiento
\ifthenelse{\boolean{orig}}
{clásico}
{clásico (\cf \Ref{IX.5}, p\ptbl\pageref{page-300})}
muestra que el conjunto de los elementos
de~$\H^1(\overline{U},\overline{F})$ cuya imagen inversa
sobre~$\widetilde{U}$ es trivial se identifica con
$\H^1(\pi_1^{\tame}(\overline{U}),\H^0(\widetilde{U},\widetilde{F})).$
Se obtiene así un isomorfismo canónico
\begin{equation*}
\label{eq:XIII.2.4.3.1} \tag{\thesubsubsection.1} {}
\H^1_{\tame}(\overline{U},\overline{F})\isomto
\H^1(\pi_1^{\tame}(\overline{U}),\H^0(\widetilde{U},\widetilde{F})).
\end{equation*}

Por consiguiente, el morfismo $\overline{\Psi}$ se identifica con el
morfismo canónico
$$
\H^1(\pi_1^{\tame}(\overline{U}),\H^0(\widetilde{U},\widetilde{F}))
\to
\H^1(\pi_1^{\tame}(\overline{U'}),\H^0(\widetilde{U}',\widetilde{F}')).
$$
Demostremos
\marginpar{383}
que este morfismo es un isomorfismo. El morfismo
$\pi_1^{\tame}(\overline{U'})\to\pi_1^{\tame}(\overline{U})$
es un isomorfismo por~\Ref{XIII.5.6}, y lo mismo ocurre con el morfismo
$\H^0(\widetilde{U},\widetilde{F})\to\allowbreak
\H^0(\widetilde{U}',\widetilde{F}')$. En efecto, esto es evidente en el
caso b), pues $\widetilde{F}$ es constante y $\widetilde{U}$ y
$\widetilde{U}'$ son conexos. En el caso a), sea $\overline{G}$ un haz
de grupos constructible sobre~$\overline{S}$ tal que
$\overline{F}=\overline{g}^*\overline{G};$ como los morfismos
$\widetilde{U}\to\overline{S}$ y
$\widetilde{U}'\to\overline{S'}$ son $0$-acíclicos
\eqref{XIII.5.7}, se tiene
$$
\H^0(\widetilde{U},\widetilde{F})\simeq
\H^0(\overline{S},\overline{G}) \simeq
\H^0(\overline{S'},\overline{G'})\simeq
\H^0(\widetilde{U}',\widetilde{F}'),
$$
lo que implica que $\Psi$ es un isomorfismo. La última afirmación
de~\Ref{XIII.2.4}~1) se deduce de lo anterior, teniendo en cuenta
\Ref{XIII.2.3}~b).

\subsubsection*{Demostración de~\Ref{XIII.2.4}~\textup{2)}}

El caso de un haz de conjuntos constructible que satisface a) se deduce
inmediatamente de~\eqref{eq:XIII.2.4.1.1}. Sea $F=g^*G$ un haz de grupos
que satisface a), donde $G$ es un haz constructible; puede suponerse que
$S$ es afín; sea $(S_j)_{j\in J}$
\ifthenelse{\boolean{orig}}
{una familia finita de subesquemas cerrados reducidos de~$S$}
{una familia finita de subesquemas localmente cerrados reducidos de~$S$}
cuya unión recubre~$S$, y tal que la imagen inversa de~$G$ sobre~$S_j$
sea un haz localmente constante. Teniendo en cuenta
\Ref{XIII.2.4}~1), para establecer que $\R^1_{\tame}i_*F$ es
constructible basta comprobarlo después del cambio de base
$S_j\to S$, para cada $j\in J$. Quedamos así reducidos al caso~b),
en el que $F$ es localmente constante.

Supongamos en adelante que $F$ es un haz de conjuntos o de grupos que
satisface b). Como la cuestión es local para la topología étale sobre~$X$,
puede suponerse que $X$ es de presentación finita sobre~$S$ y, por paso
al límite, que $X$ y $S$ son noetherianos.

Sea $D=\sum_{1\leq i\leq r}\divisor f_i$, donde, para cada punto
$x$ de~$\Supp D$, si $I(x)$ es el conjunto de los $i$ tales que
$f_i(x)=0$, el subesquema $V((f_i)_{i\in I(x)})$ es liso sobre~$S$ de
codimensión $\card I(x)$ en~$X$. Sea $\mathcal{P}$ el conjunto de las
partes de~$[ 1,r ]$ y, para cada $I\in\mathcal{P}$, pongamos
$$
X_I=\textstyle\Big(\bigcap_{i\in I}V(f_i)\Big)\cap\Big(\bigcap_{i\not\in I}X_{f_i}\Big).
$$
Sea
\marginpar{384}
$z$ un punto de~$X_I$. Tras restringirnos primero a un entorno étale
de~$z$, pueden encontrarse un abierto $W$ de~$X$ que contiene a~$z$ y un
recubrimiento principal $V$ de~$U\cap W$, moderadamente ramificado
sobre~$W$ relativamente a~$S$, del tipo considerado
en~\Ref{XIII.5.6.1}, tal que la imagen inversa de~$F$ sobre~$V$ sea
un haz constante de valor~$C$. Sea $\pi$ el grupo de Galois del
recubrimiento~$V$. Para todo punto geométrico $\overline{x}$
\ifthenelse{\boolean{orig}}
{de~$X_I$}
{de~$X_I\cap W$}
se tiene entonces, por~\eqref{eq:XIII.2.4.2.1},
$$
\H^0(\overline{U},\overline{F})\simeq \H^0(V,C_V)^{\pi}.
$$
De ello resulta que $i_*F|X_I\cap W$ es localmente constante
(SGA~4 IX~2.13) y, por consiguiente, $i_*F$ es constructible.

Demostremos, por último, que, si $F$ es un haz de grupos localmente
constante, $\R^1_{\tame}i_*F|_{X_I}$ es constructible. Si
$\overline{x}$ es un punto geométrico
\ifthenelse{\boolean{orig}}
{de~$X$}
{de~$X_I$},
se obtuvo en~\eqref{eq:XIII.2.4.3.1} la expresión
$$
\H^1_{\tame}(\overline{U},\overline{F})\isomto
\H^1(\pi^{\tame}_1(\overline{U}),\H^0(\widetilde{U},\widetilde{F})).
$$
Si $p$ es la característica residual de~$\overline{X}$, se tiene,
por~\Ref{XIII.5.6},
$\pi_1^{\tame}(\overline{U})=\prod_{l\neq p }\ZZ_l(1)^{\card I}$.
Denotemos por $\LL$ el conjunto de los números primos que dividen el
orden del grupo finito $\H^0(\widetilde{U},\widetilde{F})$, y sea
$K=\prod_{l\in\LL -\{p\}\cap\LL }\ZZ_l(1)^{\card I}$. De
\cite[I \S5 ex.2]{XIII.4} resulta que se tiene
$$
\H^1_{\tame}(\overline{U},\overline{F})\simeq
\H^1(K,\H^0(\widetilde{U},\widetilde{F})).
$$
Como $K$ está topológicamente generado por un número finito de elementos y
$\H^0(\widetilde{U},\widetilde{F})$ es finito, se deduce en primer lugar
que las fibras del haz $\R^1_{\tame}i_*F|_{X_I}$ son finitas. Por otra
parte, el conjunto $\LL$ no depende del punto $\overline{x}$. Para todo
$q\in\LL$, sea $X_{I,q}$ el cerrado de~$X_I$ definido por la ecuación
$q=0$, y sea
\ifthenelse{\boolean{orig}}
{$X_{I'}$}
{$X_I'$}
el abierto de~$X_I$ complementario de la unión de los $X_{I,q}$. Entonces
$\R^1_{\tame}i_*F|_{X_{I,q}}$ y $\R^1_{\tame}i_*F|_{X_I'}$ son
localmente constantes; en efecto, una flecha de especialización de puntos
geométricos de~$X_{I,q}$ (\resp de
\ifthenelse{\boolean{orig}}
{$X_{I'}$}
{$X_I'$})
induce un isomorfismo sobre los grupos $K$ \eqref{XIII.5.6.1}, y por tanto
también sobre los conjuntos
$\H^1_{\tame}(\overline{U},\overline{F})$; puede aplicarse entonces
SGA~4 IX~2.13.
\endgroup


\begingroup
\shorthandoff{<>}% Babel español trata de otro modo ciertos delimitadores de Xy-pic.
\begin{corollaire}
\label{XIII.2.5}
Sean
\marginpar{385}
$f\colon X\to S$ un morfismo, $D$ un divisor sobre~$X$
de cruzamientos normales relativamente a~$S$ \eqref{XIII.2.1},
$Y=\Supp D$, $U=X-Y$, e $i\colon U\to X$ la inmersión
canónica. Sea $\Phi$ un stack sobre~$U$ y supongamos dados
morfismos étale sobreyectivos $S_1\to S$ y $X_2\to
X\times_{S}S_1$, un stack $\Psi$ sobre~$S_1$ y un isomorfismo
$\Phi|U_2\simeq\Psi|U_2$ (\cf \Ref{XIII.2.2}).

Entonces, para todo morfismo $h\colon S'\to S$, si $k=h_{(X)}$,
el funtor canónico
$$
\varphi\colon k^*i_*\Phi\to i_*'\Phi'
$$
es plenamente fiel. Si $\Psi$ es constructible, el funtor
canónico
$$
\psi\colon k^*i_*^{\tame}\Phi\to i_*^{\prime\tame}\Phi'
$$
es una equivalencia de categorías.

Además, si el stack $\Psi$ es constructible (\resp si $\Psi$ es
$1$-constructible \eqref{XIII.0} y $S$ localmente noetheriano),
$i_*\Phi$ es constructible (\resp
$i_*^t\Phi$
es $1$-constructible).
\end{corollaire}

Demostremos que $\varphi$ es plenamente fiel. Basta ver
que, para todo punto geométrico $\overline{x}'$ de~$X'$, lo
mismo ocurre con el funtor
$$
\overline{\varphi}\colon
\Phi(\overline{U})\to\Phi'(\overline{U'})
$$
(se retoman las notaciones de~\Ref{XIII.2.4.2}). Sean $a$, $b$ dos
objetos de~$\Phi(\overline{U})$, y sean $a',\ b'$ sus imágenes por
$\overline{\varphi}$. Como el morfismo
$\overline{U}\to\overline{S}$ es localmente $0$-acíclico
\eqref{XIII.5.7}, se tiene un isomorfismo
$$
\H^0(\overline{U},\overline{S\Phi})
\overset{\sim}{\to}
\ifthenelse{\boolean{orig}}
{\H^0(S,\overline{S\Psi})}
{\H^0(\overline{S},\overline{S\Psi})}.
$$
Por consiguiente, $a$ y $b$ provienen por imagen inversa de objetos
de~$\Psi$, y lo mismo ocurre, por tanto, con
$F=\SheafHom_{\overline{U}}(a,b)$. Como el haz
$F'=\SheafHom_{\overline{U'}}(a',b')$ es la imagen inversa
sobre~$\overline{U'}$ de~$F$, de~\Ref{XIII.2.4.1} resulta que el
morfismo canónico
$$
\H^0(\overline{U},F)\to \H^0(\overline{U'},F')
$$
es un isomorfismo, lo que demuestra que $\overline{\varphi}$ es
plenamente fiel.

Demostremos
\marginpar{386}
que, para todo punto geométrico $\overline{x}'$ de~$X'$, el
funtor
$$
\overline{\psi}\colon
i^{\tame}_*\Phi(\overline{X'})
\to i_*^{\prime\tame}\Phi'(\overline{X'})
$$
es una equivalencia. Por lo anterior,
$\overline{\psi}$ es plenamente fiel. Demostremos que
$\overline{\psi}$ es esencialmente sobreyectivo. Sea $a'$ un
objeto de~$\Phi'(\overline{U'})$ moderadamente
ramificado sobre~$X'$ relativamente a~$S'$ \eqref{XIII.2.2}, y
demostremos que es la imagen de un objeto moderadamente
ramificado de~$\Phi(\overline{U})$. De~\Ref{XIII.2.4}
1) resulta que el morfismo canónico
$$
\H^0(\overline{U},\overline{S}\overline{\Phi})\to
\H^0(\overline{U'},\overline{S}\overline{\Phi}')
$$
es un isomorfismo. Sea $G'$ el subgerbe maximal de~$\Phi'$
generado por~$a'$; existe entonces un subgerbe maximal~$G$ de
$\overline{\Phi}$, imagen inversa de un gerbe sobre~$\overline{S}$, tal
que se tiene
$$
\overline{m}^*G\simeq G',
$$
donde
\ifthenelse{\boolean{orig}}
{$m$}
{$\overline{m}$}
es el morfismo $\overline{U'}\to
\overline{U}$. El funtor canónico
\begin{equation*}
\label{eq:XIII.2.5.*}
\tag{$*$} {\overline{k}^*\overline{i}_*^{\tame}G\to
\overline{i}_*^{\prime\tame}G'}
\end{equation*}
es una equivalencia, pues $G$ se identifica con un gerbe de
torsores bajo un haz de grupos constructible proveniente de
$\overline{S}$, y puede aplicarse \Ref{XIII.2.4}~1). De
\eqref{eq:XIII.2.5.*} resulta entonces que existe un
objeto $a$ de~$G(\overline{U})$, moderadamente
ramificado sobre~$X$ relativamente a~$S$, cuya imagen inversa sobre
$\overline{U'}$ es $a'$, lo que demuestra que $\psi$ es una
equivalencia.

Si $\Psi$ es constructible, lo mismo ocurre con~$i_*\Phi$; en
efecto, un objeto $x$ de~$i_*\Phi$ es, localmente para la topología
étale sobre~$S$, imagen inversa de un objeto $y$ de~$\Psi$; de
\Ref{eq:XIII.2.4.1.1} resulta, por tanto, que $\SheafAut (x)$ es
la imagen inversa de~$\SheafAut (y)$, luego es constructible. Finalmente, si
$\Psi$ es $1$-constructible, lo mismo ocurre con
$i^{\tame}_*\Phi$ por~\Ref{XIII.6.3} más abajo.

\begin{corollaire}
\label{XIII.2.6}
Las
\marginpar{387}
notaciones son las de~\Ref{XIII.2.4}. Supongamos que $S$ tiene
característica cero en todo punto $s$ tal que $Y_s\neq
\emptyset.$ Entonces, si $\cal{ F}$ es un haz de grupos localmente
constante y constructible sobre~$U$ (\resp un stack constructible
sobre~$U$ que, localmente sobre~$X$ y~$S$, es imagen inversa de un
stack constructible sobre~$S$), $(\cal{ F},i)$ es cohomológicamente
propio relativamente a~$S$ en dimensión $\leq 1.$
\end{corollaire}

Como todo haz de conjuntos constructible sobre~$U$ está
moderadamente ramificado sobre~$X$ relativamente a~$S,$ el
corolario resulta de~\Ref{XIII.2.4} (\resp \Ref{XIII.2.5}).

\begin{corollaire}
\label{XIII.2.7}
Las notaciones son las de~\Ref{XIII.2.4}, pero se da además
un $S$-esquema $T$, un morfismo propio $p\colon X\to T,$
y se supone que $X$ y $T$ son de presentación finita sobre~$S$; sea
$q=pi$. Sea $\cal{ F}$ un haz de conjuntos constructible sobre~$U$
que satisfaga una de las condiciones \textup{a)} o \textup{b)}
de~\Ref{XIII.2.4} (\resp un haz de grupos que satisfaga una de las
condiciones \textup{a)} o \textup{b)} de~\Ref{XIII.2.4},
\resp un stack sobre~$U$ que, localmente sobre~$X$ y~$S$, es imagen
inversa de un stack constructible $G$ sobre~$S$). Entonces
se tienen las conclusiones siguientes:
\begin{enumerate}
\item[1)] $(\cal{ F},q)$ es cohomológicamente propio relativamente
a~$S$ en dimensión $\leq 0$ (\resp para todo morfismo $h\colon
S^{'}\to S$, si $m=h_{(T)}$, el morfismo canónico
$$
\Theta\colon m^*(\R^1_{\tame}q_*\cal{ F})\to
\R^1_{\tame}q_*^{'}\cal{ F}^{'}
$$
es un isomorfismo, \resp
\ifthenelse{\boolean{orig}}
{para todo morfismo $S^{'}\to S$,}
{para todo morfismo $h\colon S^{'}\to S$, si $m=h_{(T)}$,}
el morfismo canónico
$$
\xi \colon m^*(q_*^{\tame}\cal{ F})\to
{q'}_*^{{\tame}}\cal{F}^{'}
$$
es una equivalencia).
\item[2)] El haz $q_*\cal{ F}$ (\resp el haz
$\R_{\tame}^1q_*\cal{ F}$, \resp el stack $q_*^{\tame}\cal{ F}$)
es constructible. En el último caso, si se supone que $S$ es localmente
noetheriano y que $G$ es $1$-constructible, lo mismo ocurre con
$q_*^{\tame}\cal{ F}.$
\end{enumerate}
\end{corollaire}

La
\marginpar{388}
primera parte resulta inmediatamente de~\Ref{XIII.2.4},
\Ref{XIII.2.5} y de la demostración
de~\Ref{XIII.1.8}. Demostremos~2). Si $\cal{ F}$ es un haz
de conjuntos constructible sobre~$U$ que satisface \Ref{XIII.2.4}~a)
o \Ref{XIII.2.4}~b), de \Ref{XIII.2.4}~2) resulta que
$i_*\cal{ F}$ es constructible; lo mismo ocurre, por tanto, con
$q_*\cal{ F}=p_*(i_*\cal{ F})$ (SGA~4 XIV~1.1).

Sea $\cal{F}$ un haz de grupos constructible sobre~$U$
que satisfaga \Ref{XIII.2.4}~a) o \Ref{XIII.2.4}~b), y demostremos
que $\R_{\tame}^1q_*\cal{ F}$ es constructible. Por paso al
límite (EGA~IV 8.10.5 y 17.7.8) y usando~1), puede
suponerse que $S$ es noetheriano. Sea entonces $\Phi$ el stack sobre~$X$
cuya fibra sobre todo esquema $X^{'}$ étale sobre~$X$ está formada
por los torsores sobre~$U^{'}=U\times_{X}X^{'}$, de grupo
\ifthenelse{\boolean{orig}}
{$\cal{F}|U$}
{$\cal{F}|U'$},
que están moderadamente ramificados sobre
\ifthenelse{\boolean{orig}}
{~$X$}
{~$X'$}
relativamente a~$S$; se tiene, por tanto,
$$
\ifthenelse{\boolean{orig}}
{S(i_*^{\tame}\Phi)}
{S(\Phi)}
\simeq \R_{\tame}^1i_*\cal{ F}
$$
y este haz es constructible por~\Ref{XIII.2.4}~2).
De ello resulta, por
\ifthenelse{\boolean{orig}}
{\Ref{XIII.6.3}}
{\Ref{XIII.6.2}}
más abajo, que $S(p_*\Phi)$
es constructible, \ie que $\R^1_{\tame}q_*\cal{ F}$ es
constructible.

Finalmente, si $\cal{ F}$ es un stack sobre~$U$ que, localmente
sobre~$X$ y $S$, es imagen inversa de un stack constructible sobre~$S,$
$i_*^{\tame}\cal{ F}$ es constructible, y lo mismo ocurre, por tanto,
con $q_*^{\tame}\cal{ F}=p_*(i_*^{\tame}\cal{ F})$. Si
además $S$ es localmente noetheriano y $SG$ es constructible,
$S(i_*^{\tame}\cal{ F})$ es constructible por~6.3; lo mismo ocurre,
por tanto, con
\ifthenelse{\boolean{orig}}
{$S(q_*i_*^{\tame}\cal{ F})$}
{$S(p_*i_*^{\tame}\cal{ F})$},
\ie con
$S(q_*^{\tame}\cal{ F})$~\Ref{XIII.6.2}.

\begin{corollaire}
\label{XIII.2.8}
Sea
$$
\xymatrix{ U \ar[r]^{i} \ar[d]_{g} & X \ar[ld]^{f} \\ S }
$$
un diagrama conmutativo de esquemas, en el que $U$ es el abierto
complementario
\marginpar{389}
en $X$ de un divisor de cruzamientos normales relativamente a~$S$, y
$f$ un morfismo \emph{propio} de presentación finita. Sea
$\LL$ un conjunto de números primos. Supongamos que~$g$ es localmente
$0$-acíclico (\resp localmente $1$-asférico para $\LL$). Entonces, si
$\cal{ F}$ es un haz de conjuntos sobre~$U$ (\resp un haz de $\LL$-grupos)
sobre~$U,$ localmente constante y constructible,
moderadamente ramificado sobre~$X$ relativamente a~$S$,
\ifthenelse{\boolean{orig}}
{$f_*\cal{ F}$ (\resp $\R^1_{\tame}f_*\cal{ F}$)}
{$g_*\cal{ F}$ (\resp $\R^1_{\tame}g_*\cal{ F}$)}
es localmente constante y constructible, y
\ifthenelse{\boolean{orig}}
{$(\cal{ F},f)$}
{$(\cal{ F},g)$}
es cohomológicamente propio
relativamente a~$S$ en dimensión $\leq 0$ (\resp la formación de
\ifthenelse{\boolean{orig}}
{$\R^1_{\tame}f_*\cal{ F}$}
{$\R^1_{\tame}g_*\cal{ F}$}
conmuta con todo cambio de base
$S^{'}\to S$). En el caso sin paréntesis, si $\cal{ F}$ es un
haz de grupos, para toda especialización $\overline{s}_1
\to \overline{s}_2$ de puntos geométricos de~$S$, el
morfismo de especialización
$$
\ifthenelse{\boolean{orig}}
{(\R^1_{\tame}f_*\cal{ F})_{\overline{s}_2}\to
(\R^1_{\tame}f_*\cal{ F})_{\overline{s}_1}}
{(\R^1_{\tame}g_*\cal{ F})_{\overline{s}_2}\to
(\R^1_{\tame}g_*\cal{ F})_{\overline{s}_1}}
$$
es inyectivo.
\end{corollaire}
El corolario resulta inmediatamente de~
\ifthenelse{\boolean{orig}}
{\Ref{XIII.2.6}}
{\Ref{XIII.2.7}}
y
de~\Ref{XIII.1.14} (\resp del análogo de~\Ref{XIII.1.14} para
\ifthenelse{\boolean{orig}}
{el $\R^1_{\tame}f_*\cal{ F}$}
{el $\R^1_{\tame}g_*\cal{ F}$},
que se demuestra como en \loccit).

\begin{corollaire} \label{XIII.2.9}
Sea
$$
\xymatrix{ U \ar[r]^{i} \ar[d]_{g} & X \ar[ld]^{f} \\ S }
$$
un diagrama conmutativo de esquemas, en el que $U$ es el abierto
complementario en~$X$ de un divisor de cruzamientos normales
relativamente a~$S$, y $f$~un morfismo propio y liso de
presentación finita. Sea~$\LL $ el conjunto de los números primos
distintos de las características residuales de~$S$. Sea~$\cal{F}$ un
haz de~$\LL $-grupos localmente constante y constructible sobre~$U$,
moderadamente ramificado sobre~$X$ relativamente a~$S$. Entonces
\ifthenelse{\boolean{orig}}
{$\R^1 f_*\cal{ F}$}
{$\R^1 g_*\cal{ F}$}
es localmente constante y constructible, y
\ifthenelse{\boolean{orig}}
{$(\cal{ F},f)$}
{$(\cal{ F},g)$}
es cohomológicamente propio
relativamente
\marginpar{390}
a~$S$ en dimensión $\leq 1$.
\end{corollaire}
El corolario resulta de~\Ref{XIII.2.8} y del hecho de que se tiene
\ifthenelse{\boolean{orig}}
{$\R_{\tame}^1f_*\cal{ F}=\R^1f_*\cal{ F}$}
{$\R_{\tame}^1g_*\cal{ F}=\R^1g_*\cal{ F}$}
(\Ref{XIII.2.3}~b)).

\endgroup


\begingroup
\shorthandoff{<>}% Babel español trata de otro modo ciertos delimitadores de Xy-pic.
\setcounter{subsection}{9}% Conservar el número visible XIII.2.10 de la fuente.
\subsection{}
\label{XIII.2.10}
Si $U$ es un esquema conexo, $a$ un punto geométrico de~$U$ y
$\LL$ un conjunto de números primos, se denota por
\begin{equation*}
\label{eq:XIII.2.10.0}
\tag{2.10.0} \pi_1^{\LL}(U,a)
\end{equation*}
\label{indnot:mi}\oldindexnot{$\pi_1^{\LL}(U,a)$|hyperpage}%
el límite proyectivo de los cocientes finitos de~$\pi_1(U,a)$ cuyos
órdenes tienen todos sus factores primos en~$\LL$.

Vamos a definir homomorfismos de especialización para el grupo
fundamental que generalizan X.\Ref{X.2}.

Sea $g\colon U\to S$ un morfismo coherente de fibras geométricamente
conexas (\resp un morfismo de la forma $g=fi$, donde $f\colon X\to S$
es un morfismo propio de presentación finita e $i\colon U\to X$ es una
inmersión abierta tal que $U$ sea el complementario en~$X$ de un divisor
de cruzamientos normales relativamente a~$S$ (\cf \Ref{XIII.2.8})). Sea
$\LL$ un conjunto de números primos y supongamos, en el caso sin
paréntesis, que, para todo $\LL$-grupo constante finito $C$, $(C_U,g)$
sea cohomológicamente propio relativamente a~$S$ en dimensión $\leq 1$.
Sea $\overline{s}_1\to\overline{s}_2$ un morfismo de especialización de
puntos geométricos de~$S$, sea $\overline{S}$ la localización estricta
de~$S$ en $\overline{s}_2$ y sea
$\overline{U}=U\times_S\overline{S}$. Se tiene un diagrama conmutativo
$$
\ifthenelse{\boolean{orig}}
{\xymatrix{ U_{\bar{s}_1} \ar[r]^{h_1} \ar[d] & \bar{U}
\ar[d]^{\bar{g}} & U_{\bar{s}_1} \ar[l]_{h_2} \ar[d] \\ \bar{s}_1
\ar[r] & \bar{S} & \,\bar{s}_2\,. \ar[l] }}
{\xymatrix{ U_{\bar{s}_1} \ar[r]^{h_1} \ar[d] & \bar{U}
\ar[d]^{\bar{g}} & U_{\bar{s}_2} \ar[l]_{h_2} \ar[d] \\ \bar{s}_1
\ar[r] & \bar{S} & \,\bar{s}_2\,. \ar[l] }}
$$
Si $a_1$ es un punto geométrico de~$U_{\bar{s}_1}$ y $a_2$ un
punto geométrico de~$U_{\bar{s}_2}$, los morfismos $h_1$ y $h_2$
definen homomorfismos canónicos
\marginpar{391}%
\begin{align*}
\pi_1\colon \pi_1^{\LL}(U_{\bar{s}_1},a_1)&\to
\pi_1^{\LL}(\bar{U},a_1) & \pi_2\colon
\pi_1^{\LL}(U_{\bar{s}_2},a_2)&\to \pi_1^{\LL}(\bar{U},a_2)\\
(\text{\resp} \pi_1\colon \pi_1^\tame (U_{\bar{s}_1},a_1)&\to
\pi_1^\tame (\bar{U},a_1) & \pi_2\colon
\pi_1^\tame(U_{\bar{s}_2},a_2)&\to \pi_1^\tame(\bar{U},a_2))
\end{align*}
(\Ref{V.7} y~\Ref{eq:XIII.2.1.5.2}). Las hipótesis relativas a la
propiedad de ser cohomológicamente propio (\resp
\ifthenelse{\boolean{orig}}
{\Ref{XIII.2.8}}
{\Ref{XIII.2.7}})
demuestran que $\pi_2$ es un isomorfismo. Si se elige una clase de
caminos de~$a_1$ a~$a_2$,
se obtiene un isomorfismo
$$
\pi_{12}\colon \pi_1^{\LL}(\bar{U},a_1)\isomto
\pi_1^{\LL}(\bar{U},a_2) \quad (\text{\resp} \pi_{12}\colon
\pi_1^\tame (\bar{U},a_1)\isomto \pi_1^\tame(\bar{U},a_2))
$$
de donde resulta un homomorfismo $\pi=\pi_2^{-1}\pi_{12}\pi_1$
$$
\pi\colon \pi_1^{\LL}(U_{\bar{s}_1},a_1)\to
\pi_1^{\LL}(U_{\bar{s}_2},a_2) \quad (\text{\resp}\pi\colon
\pi_1^\tame (U_{\bar{s}_1},a_1)\to \pi_1^\tame
(U_{\bar{s}_2},a_2)).
$$
Cambiar la clase de caminos de~$a_1$ a~$a_2$ equivale a modificar
$\pi$ por un automorfismo interior de
\ifthenelse{\boolean{orig}}
{$\pi_1^{\LL}(X_{\bar{s}_2},a_2)$ (\resp de
$\pi_1^\tame(X_{\bar{s}_2},a_2)$)}
{$\pi_1^{\LL}(U_{\bar{s}_2},a_2)$ (\resp de
$\pi_1^\tame(U_{\bar{s}_2},a_2)$)}. Se llama
\emph{homomorfismo de especialización para el grupo fundamental}
\index{homomorfismo de especialización para el grupo fundamental|hyperpage}%
\index{especialización para el grupo fundamental (homomorfismo de)|hyperpage}%
asociado al morfismo $\bar{s}_1\to\bar{s}_2$ y se denota simplemente por
$$
\ifthenelse{\boolean{orig}}
{\pi\colon \pi_1^{\LL}(X_{\bar{s}_1})\to
\pi_1^{\LL}(X_{\bar{s}_2}) \quad \text{(\resp} \pi\colon
\pi_1^{\tame}(X_{\bar{s}_1})\to \pi_1^{\tame}(X_{\bar{s}_2}))}
{\pi\colon \pi_1^{\LL}(U_{\bar{s}_1})\to
\pi_1^{\LL}(U_{\bar{s}_2}) \quad \text{(\resp} \pi\colon
\pi_1^{\tame}(U_{\bar{s}_1})\to \pi_1^{\tame}(U_{\bar{s}_2}))}
$$
uno de los homomorfismos definidos anteriormente.

\endgroup


\begingroup
\shorthandoff{<>}% Babel español trata de otro modo ciertos delimitadores de Xy-pic.
\setcounter{proposition}{10}% Conservar el número visible XIII.2.11 de la fuente.
\begin{lemme}
\label{XIII.2.11}
Sean $f\colon X\to S$ un morfismo propio de presentación finita,
$D$ un divisor sobre~$X$ de cruzamientos normales relativamente a~$S$,
$Y=\Supp D$, $U=X-Y$, $i\colon U\to X$ el morfismo canónico,
$\bar{s}_1\to\bar{s}_2$ un morfismo de especialización de puntos
geométricos de~$S$, $y_1$ un punto geométrico de
$Y_{\bar{s}_1}$, $y_2$ un punto geométrico de
$Y_{\bar{s}_2}$,
\ifthenelse{\boolean{orig}}
{de modo que la proyección $z_1$ de~$y_1$ sobre~$X$ sea una
generalización de la proyección $z_2$ de~$y_2$.}
{de modo que la proyección $z_1$ de~$y_1$ sobre~$X$ sea una
generalización de la proyección $z_2$ de~$y_2$; elijamos además un
morfismo de especialización $y_1\to y_2$ compatible con ella y
supongamos que $z_1$ y $z_2$ pertenecen a los mismos componentes
irreducibles de~$D$.}
Sea $I_{y_1}$
un subgrupo de inercia de~$\pi_1^{\tame}(\bar{U}_{\bar{s}_1})$ en
$y_1$. Entonces la imagen de~$I_{y_{1}}$ por el homomorfismo de
especialización
$$
\pi\colon \pi_1^{\tame}(U_{\bar{s}_1})\to
\pi_1^{\tame}(U_{\bar{s}_2})
$$
es
\marginpar{392}
un subgrupo de inercia de~$\pi_1^{\tame}(U_{\overline{s}_2})$
en~$y_2$.
\end{lemme}

En efecto, sean $\overline{X}$ (\resp $\widetilde{X}$) la localización
estricta de~$X$ en $y_2$ (\resp en $y_1$),
$\overline{U}=U \times_X \overline{X}$
(\resp $\widetilde{U}=U\times_X \widetilde{X}$). Hay un morfismo
canónico $\widetilde{U}\to \overline{U}$, y de~
\ifthenelse{\boolean{orig}}
{\Ref{XIII.1.10}}
{\Ref{XIII.2.10}}
resulta el diagrama conmutativo
$$
\xymatrix{ \pi_1^{\tame} \left(\widetilde{U}_{\overline{s}_1} \right)
\ar[r]^{\pi '} \ar[d] & \pi_1^{\tame} \left(
\overline{U}_{\overline{s}_2} \right) \ar[d] \\ \pi_1^{\tame} \left(
U_{\overline{s}_1} \right) \ar[r]^{\pi } & \pi_1^{\tame} \left(
U_{\overline{s}_2} \right) }
$$
donde $\pi'$ es el compuesto del homomorfismo canónico
$\pi_1^{\tame}(\widetilde{U}_ {\overline{s}_1})\to
\pi_1^{\tame}(\overline{U}_{\overline{s}_1})$, seguido del homomorfismo de
especialización. Como
$\pi_1^{\tame}(\widetilde{U}_{\overline{s}_1})$
(\resp
\ifthenelse{\boolean{orig}}
{$\pi_1^{\tame}(\overline{U}_{\overline{s}_1})$}
{$\pi_1^{\tame}(\overline{U}_{\overline{s}_2})$})
es un grupo de inercia de
$\pi_1^{\tame}(U_{\overline{s}_1})$ en~$y_1$
(\resp de~$\pi_1^{\tame}(U_{\overline{s}_2})$ en~$y_2$), basta
demostrar que $\pi'$ es sobreyectivo. Pero esto resulta de la expresión
obtenida en~\Ref{XIII.5.6}.

\endgroup


\begingroup
\setcounter{proposition}{11}% Conservar el número visible XIII.2.12 de la fuente.
\begin{corollaire}
\label{XIII.2.12}
Sea $X$ una curva propia, lisa y conexa de género $g$ sobre un cuerpo
separablemente cerrado $k$ de característica $p\geq 0$. Sea $U$ el abierto
obtenido al quitar de~$X$ $n$ puntos cerrados distintos
$a_1,\dots,a_n$. Entonces el grupo fundamental moderadamente ramificado
$\pi_1^{\tame}(U)$ \eqref{XIII.2.1.3} puede ser engendrado por $2g+n$
elementos $x_i,y_i,\sigma_j$, con $1\leq i\leq g$ y
$1\leq j\leq n$, de modo que $\sigma_j$ sea un generador de un grupo
de inercia correspondiente a~$a_j$, y se tenga la relación
\begin{equation*}
\label{eq:XIII.2.12.*}
\tag{$*$} \prod_{1\leq i\leq
g}(x_iy_ix^{-1}_iy^{-1}_i)\cdot\prod_{1\leq j\leq n}\sigma_j=1.
\end{equation*}
Para todo grupo finito $G$ \emph{de orden primo con} $p$,
engendrado por elementos $\overline{x}_i,\overline{y}_i,
\overline{\sigma}_j$ que satisfacen la relación
\eqref{eq:XIII.2.12.*}, existe un recubrimiento étale de~$U$, de grupo
$G$, correspondiente a un homomorfismo $\pi_1^{\tame}(U)\to G$ que
\marginpar{393}%
envía $x_i,y_i,\sigma_j$ a~$\overline{x}_i,\overline{y}_i,
\overline{\sigma}_j$, respectivamente. En otros términos, si $p'$ denota
el conjunto de los números primos distintos de~$p$, $\pi_1^{p'}(U)$ es el
grupo pro-$p'$ engendrado por los generadores $x_i,y_i,\sigma_j$, sujetos
a la única relación~\eqref{eq:XIII.2.12.*}.
\end{corollaire}

\subsubsection*{Demostración}
Podemos suponer que $k$ es algebraicamente cerrado. Supongamos primero que
$k$ tiene característica cero. Existe entonces una subextensión
algebraicamente cerrada $k'$ de~$k$, de grado de trascendencia finito
sobre~$\QQ$, tal que $X$ provenga por extensión de escalares de una curva
propia y lisa $X'$ definida sobre~$k'$, y podemos suponer que los puntos
$a_1,\dots,a_n$ provienen de puntos racionales $a'_1,\dots,a'_n$ de~$X'$.
Como $k'$ tiene grado de trascendencia finito sobre~$\QQ$, se puede hallar
una inmersión de~$k'$ en el cuerpo de los números complejos $\CC$; sea
$\widetilde{U}=U'\times_{k'}\CC$. Sea $k''$ una extensión algebraicamente
cerrada de~$k'$ tal que existan morfismos sobre~$k'$ de~$k$ y de~$\CC$
en~$k''$. Si $g'\colon U'\to k'$ es el morfismo estructural y
\ifthenelse{\boolean{orig}}
{si $\cal F$ es un haz de grupos constante finito sobre~$U''$}
{si $\cal F'$ es un haz de grupos constante finito sobre~$U'$},
se deduce de~\Ref{XIII.2.9} que los morfismos de especialización
$$
(\R^1g'_*\cal{F}')_{\CC}\to(\R^1g'_*\cal{F}')_{k''}
\from(\R^1g'_*\cal{F}')_k
$$
son isomorfismos. En términos de grupos fundamentales, esto demuestra que
se tiene un isomorfismo, definido salvo automorfismo interior,
$$
\pi_1(U)\to\pi_1(\widetilde{U}),
$$
y está claro que este isomorfismo transforma un grupo de inercia relativo
a un punto de~$X'-U'$ en un grupo de inercia relativo al mismo punto.
Por tanto, podemos suponer que $k=\CC$. En este último caso, del teorema
de existencia de Riemann (\Ref{XII.5.2}) resulta que el grupo fundamental
$\pi_1(U)$ no es sino la completación, para la topología de los subgrupos
de índice finito, del grupo fundamental del espacio analítico asociado
a~$U$. Ahora bien, este último puede calcularse por métodos trascendentes
\cite[ch\ptbl 7
\S 47]{XIII.3};
\marginpar{394}%
puede ser engendrado por $2g+n$ elementos $x_i,y_i,\sigma_j$ tales que
$\sigma_j$ sea la imagen de un generador
\ifthenelse{\boolean{orig}}
{del grupo fundamental local $\pi_1(D_j)$ de un pequeño disco centrado
en~$a_j$}
{del grupo fundamental local $\pi_1(D_j-\{a_j\})$, donde $D_j$ es un
pequeño disco centrado en~$a_j$},
\ie un generador de un grupo de inercia correspondiente al punto $a_j$,
y estos elementos satisfacen la única relación \eqref{eq:XIII.2.12.*}.

Si ahora $k$ tiene característica $p>0$, se puede hallar un anillo de
valoración discreta completo $A$, con cuerpo residual $k$ y cuerpo de
fracciones $K$ de característica cero, y un esquema conexo $X_1$, propio
y liso sobre~$S=\Spec A$, tal que
$X_1\times_S\Spec k\simeq X$ (\Ref{III.7.4}). Los puntos $a_j$ se
levantan entonces a secciones $s_j$ de~$X_1$ sobre~$S$; sea $Y_{1j}$ el
subesquema cerrado reducido cuyo espacio subyacente es $s_j(S)$, sea $Y_1$
la unión de los $Y_{1j}$, sea $U_1=X_1-Y_1$ y sea
$g_1\colon U_1\to S$ el morfismo estructural. Sean $\overline K$ una
extensión algebraicamente cerrada de~$K$,
\ifthenelse{\boolean{orig}}
{}
{$\overline X=X_1\times_S\overline K$ y}
$\overline U=U_1\times_S\overline K$. Si $C$ es un grupo constante finito,
se deduce de~\Ref{XIII.2.8} que el morfismo de especialización
$$
(\R^1g_{1*}C_{U_1})_k\to(\R^1g_{1*}C_{U_1})_{\overline K}
$$
es inyectivo, e incluso biyectivo si $C$ es de orden primo con~$p$.
Ahora bien, esto significa, en términos de grupos fundamentales, que el
morfismo de especialización
\ifthenelse{\boolean{orig}}
{\eqref{XIII.1.10}}
{\eqref{XIII.2.10}}
$$
\pi\colon \pi_1(\overline U)\to\pi^{\tame}_1(U)
$$
es sobreyectivo, y que el morfismo de especialización
$$
\pi^{p'}_1(\overline U)\to\pi^{p'}_1(U)
$$
es biyectivo. Por último, si $x_i,y_i,\sigma_j$ son generadores de
$\pi_1(\overline U)$ tales que $\sigma_j$ sea un generador de un grupo de
inercia correspondiente al punto $b_j=Y_{1j}(\overline K)$ de~$\overline X$,
entonces, según
\ifthenelse{\boolean{orig}}
{\Ref{XIII.1.11}}
{\Ref{XIII.2.11}},
$\pi(\sigma_j)$ es un generador de un grupo de inercia correspondiente
a~$a_j$, lo cual concluye la demostración.

\endgroup


\ifthenelse{\boolean{orig}}
{}
{\begin{remarqueMR}
\label{XIII.2.13}
Sean $k$ un cuerpo algebraicamente cerrado de característica $p>0$, $X$ una
curva algebraica propia y lisa sobre~$k$, conexa, de género $g$, y $U$ un
abierto afín de~$X$ cuyo complemento consta de $r\geq1$ puntos racionales
de~$X$. Se dispone del grupo fundamental $\pi_1(U)$, de su cociente
$\pi_1^{\tame}(U)$, que clasifica los recubrimientos étales finitos de~$U$
moderadamente ramificados en los puntos de~$X-U$, y del cociente
$\pi_1^{p'}(U)$ de~$\pi_1^{\tame}(U)$, que clasifica los recubrimientos
étales galoisianos de~$U$ cuyo grupo de Galois es de orden primo con~$p$.

En (Coverings of algebraic curves, Amer.\ J.~Math.\ \textbf{79}
(1957), p\ptbl825--856), S\ptbl Abhyankar formuló varias
conjeturas sobre la estructura de los grupos finitos que son grupos de
Galois de recubrimientos étales finitos de~$U$.

Las conjeturas relativas a los grupos finitos que aparecen como grupos de
Galois de recubrimientos étales conexos de~$U$ y cuyo orden es primo con~$p$
quedan demostradas y precisadas en el corolario \Ref{XIII.2.12}. Antes de abordar
los cocientes finitos de~$\pi_1(U)$, comencemos por dar algunas indicaciones
sobre el \og tamaño\fg de este grupo. Sea
\ifthenelse{\boolean{orig}}
{$X = \Spec(A)$}
{$U = \Spec(A)$}.
Se sabe que $\Hom_{\mathrm{cont}}(\pi_1(U), \ZZ/p\ZZ)$ se describe mediante
la teoría de Artin--Schreier (cor\ptbl \Ref{XI.6.9}). Este grupo es
isomorfo a $A/\wp(A)$, donde $\wp$ es la aplicación de~$A$ en~$A$ que envía
$a$ a~$a^p-a$. Supongamos primero que $U$ es la recta afín $\AA^1$, de
anillo $k[T]$. Sea~$E$ el conjunto de los elementos de~$k[T]$ de la forma
$\sum_i a_iT^i$, con $a_i = 0$ cuando $p$ divide a~$i$. Es inmediato que la
aplicación compuesta $E\to k[T]\to k[T]/\wp k[T]$ es biyectiva. Por tanto,
los coeficientes $a_i$, $(i,p)=1$, se comportan como coordenadas de un
espacio que parametriza los recubrimientos cíclicos de grado $p$ de la
recta afín. En particular, se deduce que $\pi_1(\AA^1)$ no está
topológicamente generado por un número finito de elementos y que, si se
efectúa una extensión $k\to k'$ de cuerpos algebraicamente cerrados, el
morfismo natural $\pi_1(\AA^1\times_kk')\to\pi_1(\AA^1)$, que es
sobreyectivo, no es biyectivo, a diferencia del caso propio. En el caso
general, la curva $U$ puede realizarse como esquema finito sobre la recta
afín, y se concluye que los mismos fenómenos se producen para $\pi_1(U)$
(y, de hecho, más generalmente, para $\pi_1(V)$, para todo $k$-esquema
afín conexo $V$, de tipo finito y de dimensión $\geq1$).

Dicho esto, si $G$ es un grupo finito, denotemos por $G^{(p')}$ el mayor
grupo cociente de~$G$ cuyo orden es primo con~$p$. Para que un grupo finito $G$
sea un cociente topológico de~$\pi_1(U)$, es necesario que $G^{(p')}$ sea
un cociente topológico de~$\pi_1^{p'}(U)$, condición que, en principio,
puede decidirse gracias al corolario \Ref{XIII.2.12}.
En el artículo citado, Abhyankar conjetura que esta condición necesaria es
también suficiente. Esta conjetura fue demostrada por M\ptbl Raynaud en el
caso de la recta afín y por D\ptbl Harbater en el caso general
(Invent. Math.\ \textbf{116} (1994), p\ptbl425--462 y \textbf{117},
p\ptbl1--25). Por ejemplo, en el caso de la recta afín
$\pi_1^{p'}(\AA^1)=1$, y se concluye que un grupo finito $G$ es grupo de
Galois de un recubrimiento étale conexo de~$\AA^1$ si y solo si
$G^{(p')} = 1$, es decir, si y solo si $G$ está engendrado por sus
$p$-subgrupos de Sylow. Así, todo grupo finito simple de orden múltiplo
de~$p$ cumple esta condición.
\end{remarqueMR}}


\section[Ser cohomológicamente propio y aciclicidad local]{Propiedad de ser cohomológicamente propio y aciclicidad local genérica}
\label{XIII.3}
\marginpar{395}

\begin{theoreme}
\label{XIII.3.1}
Sean $S$ un esquema irreducible con punto genérico~$s$, $X$ e~$Y$ dos
$S$-esquemas de presentación finita y $f\colon X\to Y$ un
$S$-morfismo. Para todo $S$-esquema~$S'$, se denotan por $Y'$, $X'$, etc.
las imágenes inversas de~$Y$, $X$, etc. por el morfismo $S'\to S$.
Se tienen las propiedades siguientes:

\begin{enumerate}
\item[1)]
{\upshape a)} Puede encontrarse un abierto no vacío~$S'$ de~$S$ tal que,
para todo haz de conjuntos constante finito~$F'$ sobre~$X'$, $f'_*F'$ sea
constructible y~$(F',f')$ sea cohomológicamente propio relativamente
a~$S'$ en dimensión $\leqslant 0$.

{\upshape b)} Sea $F$ un haz de conjuntos constructible sobre~$X$.
Entonces puede encontrarse un abierto no vacío~$S'$ de~$S$ (dependiente
de~$F$) tal que $f'_*F'$ sea constructible y~$(F',f')$ sea
cohomológicamente propio relativamente a~$S'$ en dimensión
$\leqslant 0$.

\item[2)] Supongamos que los esquemas de tipo finito de dimensión
$\leqslant\dim X_s$ sobre una clausura algebraica~$\overline{k}$
de~$\kres(s)$ son fuertemente desingularizables
\textup{(SGA~5 I~3.1.5)}. Se tienen entonces, además, las propiedades
siguientes:

{\upshape a)} Puede encontrarse un abierto no vacío~$S'$ de~$S$ tal que,
para todo haz de grupos constante finito~$F'$ sobre~$X'$, de orden primo
con las características residuales de~$S$, si~$\Phi'$ es el stack de los
torsores bajo~$F'$, $f'_*\Phi'$ sea $1$-constructible y~$(F',f')$ sea
cohomológicamente propio relativamente a~$S'$ en dimensión
$\leqslant 1$.

{\upshape b)} Sean $\LL$ el conjunto de los números primos distintos de
las características residuales de~$S$ y~$\Phi$ un stack ind-$\LL$,
$1$-constructible, sobre~$X$ \eqref{XIII.0}, tal que, para todo
esquema~$X_1$ étale sobre~$X$ y para todo par de objetos $x,x_1$
de~$\Phi_{X_1}$, el haz $\SheafHom_{X_1}(x,x_1)$ sea constructible.
Supongamos además que~$S$ es localmente noetheriano. Entonces puede
encontrarse un abierto no vacío~$S'$ de~$S$ tal que $f'_*\Phi'$ sea
$1$-constructible, que,
\ifthenelse{\boolean{orig}}
{}
{para todo esquema~$Y_1$ étale sobre~$Y'$ y}
para todo par de objetos $y,y_1$ de la fibra $(f'_*\Phi')_{Y_1}$,
el haz $\SheafHom_{Y_1}(y,y_1)$ sea constructible,
\marginpar{396}
y que~$(\Phi',f')$ sea cohomológicamente propio relativamente a~$S'$
en dimensión $\leqslant 1$.
\end{enumerate}
\end{theoreme}

\subsubsection*{Demostración}
Puede suponerse que~$S$ es afín; según SGA~4 VIII~1.1, puede suponerse
que~$S$ es íntegro; finalmente, por paso al límite, puede suponerse que
$S$ es el espectro de un álgebra de tipo finito sobre~$\ZZ$; en
particular, $S$ es entonces noetheriano.
Como la cuestión es local sobre~$Y$, puede suponerse que~$Y$ es afín.
Además, para demostrar el teorema, basta hacerlo tras una extensión finita
de la base $S'\to S$, donde~$S'$ es un esquema íntegro y~$S'\to S$ es un
compuesto de morfismos étales y de morfismos finitos, radiciales y
sobreyectivos.

\shorthandoff{<>}% Babel español trata como abreviaturas los delimitadores de curvatura de Xy-pic.

\Subsection*{1) Caso de los haces de conjuntos constantes}

1) 1. Reducción al caso en que $X$ es normal sobre~$S$. Sea
$X_{1\overline{s}}$ la normalización de~$(X_{\overline{s}})_{\red}$;
tras restringir $S$ a un abierto no vacío y efectuar una extensión
radicial de~$S$, puede suponerse que $X_{1\overline{s}}$ procede de un
esquema $X_1$ normal sobre~$S$, y que el morfismo
$X_{1\overline{s}}\to X_{\overline{s}}$ procede de un morfismo finito
sobreyectivo $p\colon X_1\to X$ (EGA~IV 8.8.2 y 9.6.1). Supongamos
demostrado el teorema para $fp$. Tras restringir $S$ a un abierto,
puede suponerse que, para todo haz de conjuntos constante finito $F$
sobre~$X$, $(p^*F,fp)$ es cohomológicamente propio relativamente a~$S$
en dimensión $\leqslant 0$ y que $f_*p_*(p^*F)$ es constructible. Según
\Ref{XIII.1.9}, $(p_*p^*F,f)$ es entonces cohomológicamente propio
relativamente a~$S$ en dimensión $\leqslant 0$. El morfismo
$$
F \to p_* p^* F=G
$$
es un monomorfismo. De ello resulta ya, pues $f_*F$ es un subhaz
de~$f_*G$, que $f_*F$ es constructible (SGA~4 IX~2.9~(ii)) y que
$(F,f)$ es cohomológicamente propio relativamente a~$S$ en dimensión
$\leqslant -1$.

Sean $X_2=X_1 \times_{X} X_1$ y $q \colon X_2 \to X$ el morfismo
canónico. Según \Ref{XIII.1.11}~1),
\marginpar{397}
se tiene una sucesión exacta
$$
\xymatrix@C=.5cm{F \ar[r] & G \ar@<2pt>[r]\ar@<-2pt>[r] & q_*q^*F }.
$$

Por lo que acaba de demostrarse, aplicado a $fq$ en lugar de~$f$, puede
suponerse, tras restringir $S$ a un abierto no vacío, que, para todo haz
de conjuntos constante finito $F$ sobre~$X$, $(q^*F,fq)$ es
cohomológicamente propio relativamente a~$S$ en dimensión $\leqslant
-1$, y por tanto que $(q_*q^*F,f)$ es cohomológicamente propio
relativamente a~$S$ en dimensión $\leqslant -1$. Resulta entonces de
\Ref{XIII.1.13}~1) que $(F,f)$ es cohomológicamente propio relativamente
a~$S$ en dimensión $\leqslant 0$.

1) 2. Reducción al caso en que $X$ es normal y afín sobre~$S$.

Sea $U_s$ un abierto afín de~$X_s$, denso en $X_s$. Tras restringir $S$
a un abierto no vacío, puede suponerse que $U_s\to X_s$ se levanta a una
inmersión abierta $i\colon U\to X$, esquemáticamente dominante
relativamente a~$S$ (EGA~IV 8.9.1). Como el morfismo $X\to S$ es
normal, según SGA~2 XIV~1.18 se tiene:
$$
\ifthenelse{\boolean{orig}}
{\prof \et_{S- U} (X) \ge 2 \quoi ;}
{\prof \et_{X- U} (X) \ge 2 \quoi ;}
$$
por consiguiente, para todo haz de conjuntos constante finito $F$
sobre~$X$, el morfismo canónico
$$
F \to i_*i^*F
$$
es un isomorfismo. De ello resulta que, si se supone demostrado el
teorema para $fi$ y para~$i$, entonces, tras restringir $S$ a un abierto
no vacío, $(i^*F,i)$ y $(i^*F,fi)$ son cohomológicamente propios
relativamente a~$S$ en dimensión $\leqslant 0$. Lo mismo vale, pues,
para~$(F,f)$ (\Ref{XIII.1.6}~2)). Como, además,
$f_*F=(fi)_*(i^*F)$ es constructible, queda terminada la reducción.

1) 3. Fin de la demostración.

Puede suponerse que $S$ es normal (EGA~IV 7.8.3). Puede encontrarse una
compactificación de~$X_s$:
\marginpar{398}
$$
\xymatrix{ X_s \ar[r]^{j_s} \ar[d]_{f_s} & P_s \ar[ld]^{g_s} \\ Y_s }
$$
donde $j_s$ es una inmersión abierta dominante y $g_s$ un morfismo
propio; tras efectuar una extensión radicial de~$\kres(s)$ y sustituir
$P_s$ por su normalización, lo que no modifica $X_s$, puede suponerse
que $P_s$ es geométricamente normal. Tras restringir $S$ a un abierto
no vacío y efectuar una extensión radicial sobreyectiva, puede suponerse
que el diagrama anterior procede de un diagrama
$$
\xymatrix{ X \ar[r]^{j} \ar[d]_{f} & P \ar[ld]^{g} \\ Y }
$$
donde $P$ es un esquema normal sobre~$S$, $j$ una inmersión abierta
esquemáticamente dominante relativamente a~$S$ y $g$ un morfismo propio
(EGA~IV 6.9.1, 9.9.4 y 9.6.1). Para todo haz de conjuntos constante
finito $F$ sobre~$X$, de valor $C$, $j_*F$ es el haz constante de
valor~$C$ (SGA~4 \ifthenelse{\boolean{orig}}{}{IX~}2.14.1), y lo mismo
ocurre tras todo cambio de base $S'\to S$;
de ello resulta que $(F,j)$ es cohomológicamente propio relativamente
a~$S$ en dimensión $\leqslant 0$. Lo mismo vale, pues, para~$(F,f)$,
puesto que $g$ es propio \eqref{XIII.1.8}. Como $g$ es propio,
$f_*F=g_*C_P$ es constructible, lo que termina la demostración de
1)\kern2pt a).

\shorthandoff{<>}% Babel español trata como abreviaturas los delimitadores de curvatura de Xy-pic.

\Subsection*{2) Caso de un haz de conjuntos constructible}

Sea $F$ un haz de conjuntos constructible sobre~$X$. Según
SGA~4 IX~2.14~(ii), puede encontrarse una familia finita de morfismos
$p_i \colon Z_i \to X$ y, sobre cada $Z_i$, un haz de conjuntos
constante finito $C_i$, de modo que se tenga un monomorfismo
\marginpar{399}
$$
j \colon F \to \prod_{i} {p_i}_* C_i =G \quoi.
$$
Según~1)\kern2pt a), puede suponerse, tras restringir $S$ a un abierto
no vacío, que los $(C_i,fp_i)$ son cohomológicamente propios relativamente
a~$S$ en dimensión $\leqslant 0$ y que los $f_*{p_i}_*C_i$ son
constructibles. De ello se concluye ya que $(G,f)$ es
cohomológicamente propio relativamente a~$S$ en dimensión $\leqslant 0$
\eqref{XIII.1.9}, y por tanto que $(F,f)$ es cohomológicamente propio
relativamente a~$S$ en dimensión $\leqslant -1$, y que $f_*F$ es
constructible. Sea~$K$ la suma amalgamada $K= G \coprod_{F} G$;
como $F$ y $G$ son constructibles, también lo es~$K$. De lo anterior
se concluye, pues, que, tras restringir $S$ a un abierto no vacío, puede
suponerse que $(K,f)$ es cohomológicamente propio relativamente a~$S$
en dimensión $\leqslant -1$. Resulta entonces de
\Ref{XIII.1.13}~1) que $(F,f)$ es cohomológicamente propio relativamente
a~$S$ en dimensión $\leqslant 0$.

\Subsection*{3) Caso de los haces de grupos constantes}

Si $F$ es un haz de grupos constante sobre~$X$, se denota por $\Phi$ el
stack de los torsores bajo~$F$.

3) 1. Demostremos primero que, tras restringir $S$ a un abierto no
vacío, para todo haz de grupos constante $F$ sobre~$X$, de orden primo
con las características residuales de~$S$, $(F,f)$ es cohomológicamente
propio relativamente a~$S$ en dimensión $\leqslant 0$, y que
$f_*\Phi$ es constructible.

Para ello se reduce al caso en que $X$ es liso sobre~$S$. Tras efectuar
una extensión finita de~$\kres(s)$, lo cual es lícito porque puede
factorizarse como una extensión étale seguida de una extensión
radicial, puede encontrarse un morfismo propio sobreyectivo
$p_s \colon X_{1s}\to X_s$, donde $X_{1s}$ es un esquema liso
\ifthenelse{\boolean{orig}}
{sobre~$S$}
{sobre~$s$}
de la misma dimensión que $X_s$ y, tras restringir $S$ a un abierto no
vacío, puede suponerse que $p_s$ procede de un morfismo propio
sobreyectivo $p \colon X_1 \to X$, donde $X_1$ es un esquema liso
sobre~$S$
\marginpar{400}
(EGA~IV 9.6.1 y 12.1.6). Sean $X_2=X_1 \times_X X_1$ y
$q \colon X_2 \to X$ el morfismo canónico. Se tiene un diagrama exacto
de stacks sobre~$X$
$$
\xymatrix@C=.5cm{\Phi \ar[r] & p_*p^* \Phi \ar@<2pt>[r]\ar@<-2pt>[r] &
q_*q^*\Phi }
$$
(\Ref{XIII.1.11}~2)). Como se supone demostrado el teorema en el caso
liso, se ve ante todo que puede suponerse que $f_*p_*p^*\Phi$ es
constructible; lo mismo vale, pues, para $f_*\Phi$
(\Ref{XIII.3.1.1} más adelante). Además, según
\Ref{XIII.1.6}~2), puede suponerse que $(p_*p^*\Phi,f)$ es
cohomológicamente propio relativamente a~$S$ en dimensión
$\leqslant 0$; de ello resulta que $(\Phi,f)$ es cohomológicamente
propio relativamente a~$S$ en dimensión $\leqslant -1$. Puede
suponerse entonces que $(q_*q^*\Phi,f)$ es cohomológicamente propio
relativamente a~$S$ en dimensión $\leqslant -1$, y ello implica que
$(\Phi,f)$ es cohomológicamente propio relativamente a~$S$ en dimensión
$\leqslant 0$ (\Ref{XIII.1.12}~1)).

Se reduce después, como en 1)\kern2pt2, al caso en que $X$ es liso y
afín sobre~$S$. Sea entonces
$$
\xymatrix{ X \ar[r]^{i} \ar[d]_{f} & P \ar[ld]^{q} \\ Y }
$$
una compactificación de~$X$, donde $i$ es una inmersión abierta
dominante y $q$ un morfismo propio. Como $\dim P_s=\dim X_s$, puede
aplicarse la hipótesis de resolución de singularidades a
$P_{\overline s}$. Tras efectuar una extensión étale y una extensión
radicial de~$S$, puede encontrarse un morfismo propio
$r \colon Z \to P$, donde $Z$ es liso sobre~$S$,
$r^{-1}(X)\simeq X$, y donde $r^{-1}(X)$ es el complementario en~$Z$
de un divisor de cruzamientos normales relativamente a~$S$. Todo
torsor bajo~$F$ está entonces moderadamente ramificado sobre~$Z$
relativamente a~$S$ (\Ref{XIII.2.3}~b)). Resulta, pues, de
\Ref{XIII.2.7} que $(F,f)$ es cohomológicamente propio relativamente
a~$S$ en dimensión $\leqslant 0$, lo que demuestra la afirmación.

\ifthenelse{\boolean{orig}}
{3) \textbf{2. Reducción\marginpar{401} al caso en que $X$ es liso sobre~$S$}}%
{\newpage 3) 2. Reducción\marginpar{401} al caso en que $X$ es liso sobre~$S$.}%
\par\nopagebreak
Tras efectuar una extensión finita de~$\kres(s)$, puede encontrarse un
morfismo propio sobreyectivo
\ifthenelse{\boolean{orig}}
{$p_s\colon X_{1s}\to X$}
{$p_s\colon X_{1s}\to X_s$},
donde~$X_{1s}$ es un esquema liso sobre~$s$, y puede suponerse que
$p_s$ procede de un morfismo propio sobreyectivo
$p\colon X_1 \to X$, donde $X_1$ es liso sobre~$S$.
Supongamos demostrado el teorema para $fp$ y demostrémoslo para~$f$.
Sea $F$ un haz de grupos constante finito sobre~$X$, de orden primo con
las características residuales de~$S$, y sea $\Phi$ el stack de los
torsores bajo~$F$. Sean $X_2=X_1 \times_X X_1$,
$X_3=X_1 \times_X X_1 \times_X X_1$, y sean
$q\colon X_2 \to X$ y $r\colon X_3 \to X$ los morfismos canónicos.
Según \Ref{XIII.1.11}~2), se tiene un diagrama exacto de stacks
$$
\xymatrix@C=.5cm{\Phi \ar[r] & p_*p^*\Phi \ar@<2pt>[r]\ar@<-2pt>[r] &
q_*q^*\Phi \ar@<4pt>[r] \ar[r] \ar@<-4pt>[r] & r_*r^*\Phi}\quoi.
$$
Para demostrar que $(\Phi,f)$ es cohomológicamente propio relativamente
a~$S$ en dimensión $\leqslant 1$, basta demostrar que lo mismo vale
para~$(p_*p^*\Phi,f)$, que $(q_*q^*\Phi, f )$ es cohomológicamente
propio relativamente a~$S$ en dimensión $\leqslant 0$ y que
$(r_*r^*\Phi,f)$ es cohomológicamente propio relativamente a~$S$ en
dimensión $\leqslant -1$ (\Ref{XIII.1.12}~2)). Según~3)\kern.3em 1
anterior, puede suponerse que, para todo haz de grupos constante
finito~$F$, $(q^*\Phi,fq)$, $(q^*\Phi,q)$, $(r^*\Phi,fr)$ y
$(r^*\Phi,r)$ son cohomológicamente propios relativamente a~$S$ en
dimensión $\leqslant 0$. Resulta entonces de \Ref{XIII.1.6}~2) que
$(q_*q^*\Phi,f)$ y $(r_*r^*\Phi,f)$ son cohomológicamente propios
relativamente a~$S$ en dimensión $\leqslant 0$. Como se supone
demostrado el teorema en el caso liso, $(p^*\Phi,fp)$ y
$(p^*\Phi,p)$, y por tanto también $(p_*p^*\Phi,f)$
(\Ref{XIII.1.6}~2)), son cohomológicamente propios relativamente a~$S$
en dimensión $\leqslant 1$. Esto demuestra que $(\Phi,f)$ es
cohomológicamente propio relativamente a~$S$ en dimensión
$\leqslant 1$.

Además, $f_*p_*p^*\Phi$ es $1$-constructible por hipótesis;
\ifthenelse{\boolean{orig}}
{según \Ref{XIII.3.1},}
{según~3)\kern.3em 1 anterior,}
puede suponerse que $f_*q_*q^*\Phi$ es constructible; resulta, pues,
de \Ref{XIII.3.1.1} más adelante que $f_*\Phi$ es $1$-constructible.

\shorthandoff{<>}% Babel español trata como abreviaturas los delimitadores de curvatura de Xy-pic.

3) 3. Reducción al caso en que $X$ es liso y afín sobre~$S$.
\marginpar{402}

Según 3)\kern2pt2, puede suponerse que $X$ es liso sobre~$S$. Sea
$(U_i)_{i\in I}$ un recubrimiento finito de~$X$ por abiertos afines, y
sea $X_1$ la suma disjunta de los $U_i$ y
$p\colon X_1 \to X$ el morfismo canónico. Como $p$ es un morfismo de
descenso efectivo para la categoría de los haces étales de tipo finito
sobre esquemas variables, se ve como en 3)\kern2pt2 que, si se supone
demostrado el teorema para los $U_i$, \ie para $X_1$, también vale
para~$X$.

3) 4. Caso en que $X$ es liso y afín sobre~$S$.

Como en 3)1, se ve que, tras restringir $S$ a un abierto no vacío y
efectuar una extensión étale sobreyectiva y una extensión radicial
sobreyectiva de~$S$, puede encontrarse un diagrama conmutativo
$$
\xymatrix{ X \ar[r]^{i} \ar[d]_{f} & P \ar[ld]^{q} \\ Y }
$$
donde $P$ es un esquema liso sobre~$S$, $X$ es el complementario en~$P$
de un divisor de cruzamientos normales relativamente a~$S$
\ifthenelse{\boolean{orig}}
{y $g$ es un morfismo propio.}
{y $q$ es un morfismo propio.}
Si $F$ es un haz constante de orden primo con las características
residuales de~$S$, todo torsor bajo~$F$ está moderadamente ramificado
sobre~$P$ relativamente a~$S$ (\Ref{XIII.2.3}~b)). El hecho de que
$(F,f)$ sea cohomológicamente propio relativamente a~$S$ en dimensión
$\leqslant 1$ y que $f_*\Phi$ sea constructible resulta entonces
de~\Ref{XIII.2.7}.


\ifthenelse{\boolean{orig}}
{4) \textbf{Demostración de } 2)\kern2pt b)}
{\Subsection*{4) Demostración de 2)\kern2pt b)}}

4) 1. Caso en que $\Phi$ es un gerbe.

Puede encontrarse un morfismo étale sobreyectivo de tipo finito $p
\colon X_1 \to X$, tal que $p^*\Phi$ sea un gerbe trivial. Por
descenso, como en 3)\kern2pt2, se ve que basta demostrar el teorema
para $X_1$, $X_1\times_X X_1$ y
\marginpar{403}
$X_1\times_X X_1\times_X X_1$. Se reduce, pues, al caso en que
$\Phi$ es el gerbe de los torsores bajo un haz de grupos
constructible~$F$, cuyas fibras tienen orden primo con las
características residuales de~$S$.

Según SGA~4 IX~2.14, puede encontrarse una familia finita de
morfismos finitos $p_i \colon Z_i \to X$ y, para cada $i$, un
haz de grupos constante finito $C_i$, de orden primo con las
características residuales de~$S$, de modo que se tenga un
\ifthenelse{\boolean{orig}}
{morfismo}
{monomorfismo}
$$
j\colon F \to \prod_{i} p_{i*}C_i=G \quoi.
$$
Sea $\Phi_i$ el stack de los torsores bajo~$C_i$ y $\Psi$ el stack de
los torsores bajo~$G$. Resulta de 2)\kern2pt a) que, tras restringir
$S$ a un abierto no vacío, puede suponerse que los $(C_i,fp_i)$ son
cohomológicamente propios relativamente a~$S$ en dimensión
$\leqslant 1$ y que los stacks $f_*p_{i*}\Phi_i$ son
$1$-constructibles. Resulta entonces de~\Ref{XIII.1.9} que los
$(p_{i*}C_i,f)$ son cohomológicamente propios relativamente a~$S$
en dimensión $\leqslant 1$; lo mismo vale, por tanto, para~$(G,f)$.
Además, como $p_{i*}\Phi_i$ es equivalente al stack de los torsores
bajo el grupo $p_{i*}C_i$ (SGA~4 VIII~5.8), se ve que $f_*\Psi$ es
$1$-constructible.

Como $\R^1f_*G$ es constructible, puede encontrarse un haz
representable por un $Y$-esquema étale de tipo finito~$T$ y un
epimorfismo
$$
a \colon T \to \R^1f_*G
$$
(SGA~4 IX~2.7); además, puede suponerse que la imagen de la sección
identidad de~$T(T)$ está definida por un torsor~$Q$ sobre~$X\times_Y
T=X_T$, bajo~$G|X_T$. Sean $f_T \colon X_T \to T$ el morfismo
canónico, $F_T=F|X_T$, etc. Según 1)\kern2pt b), tras restringir
$S$ a un abierto no vacío, puede suponerse que $(Q/F_T,f_T)$ es
cohomológicamente propio relativamente a~$S$ en dimensión
$\leqslant 0$ y que $f_{T*}(Q/F_T)$ es constructible. Resulta entonces
de~\Ref{XIII.3.1.2} que
\ifthenelse{\boolean{orig}}
{$f_*\Psi$ es constructible.}
{$f_*\Phi$ es $1$-constructible.}

Demostremos que $(F,f)$ es cohomológicamente propio relativamente a
\marginpar{404}
~$S$ en dimensión $\leqslant 1$. Según \Ref{XIII.1.13}~2), basta
demostrar que, para todo esquema $Y_1$ étale sobre~$Y$ y para todo
\ifthenelse{\boolean{orig}}
{torsor $Q_1$ sobre~$X_1=X\times_Y Y_1$,}
{torsor $Q_1$ bajo~$G_1$ sobre~$X_1=X\times_Y Y_1$,}
si
$f_1 \colon X_1 \to Y_1$ es el morfismo canónico, entonces
$(Q_1/F_1,f_1)$ es cohomológicamente propio relativamente a~$S$ en
dimensión $\leqslant 0$. Ahora bien, por definición de~$T$, $Q_1$ es,
localmente para la topología étale de~$Y_1$, imagen inversa de~$Q$,
lo que demuestra nuestra reducción.


4) 2. Caso general.\par\nopagebreak
Se ve, utilizando el lema~\Ref{XIII.6.1.1}, 4)\kern2pt1 y
\ifthenelse{\boolean{orig}}
{1)\kern2pt a),}
{1)\kern2pt b),}
que, tras restringir $S$ a un abierto no vacío, puede suponerse que
$S(f_*\Phi)$ es constructible y que $(S\Phi,f)$ es cohomológicamente
propio relativamente a~$S$ en dimensión $\leqslant 0$. Puede entonces
encontrarse un haz representable por un $Y$-esquema étale de tipo
finito~$T$ y un epimorfismo
$$
a \colon T \to S(f_*\Phi)\quoi.
$$
Se retoman las notaciones de 4)\kern2pt1 y se pone además
$Z=T\times_Y T$, $X_Z=X\times_Y Z$; se denota por
$f_Z \colon X_Z \to Z$ el morfismo canónico. Puede suponerse que~$T$
se ha elegido de modo que la imagen~$q$ de la sección identidad de~$T(T)$
por~$a$ esté definida por un objeto~$p$ de
$(f_*\Phi)_T=\Phi_{X_T}$. Sean $p_1$ y $p_2$ (\resp $q_1$ y $q_2$)
las imágenes inversas de~$p$ (\resp $q$) por las dos proyecciones
de~$Z$ en~$T$. Tras restringir $S$ a un abierto no vacío, puede
suponerse que $(\SheafAut_{X_T}(p),f_T)$ es cohomológicamente propio
relativamente a~$S$ en dimensión $\leqslant 1$, que
$(\SheafHom_{X_Z}(p_1,p_2),f_Z)$ es cohomológicamente propio
relativamente a~$S$ en dimensión $\leqslant 0$, y que
$f_{T*}(\SheafAut_{X_T}(p))$ y
$f_{Z*}(\SheafHom_{X_Z}(p_1,p_2))$ son constructibles
\ifthenelse{\boolean{orig}}
{(a)\kern2pt1) y 4)\kern2pt1)).}
{(1)\kern2pt b) y 4)\kern2pt1)).}

Se deduce en primer lugar que, para todo esquema $Y_1$ étale sobre~$Y$
y para todo par de objetos $y$, $y_1$ de~$(f_*\Phi)_{Y_1}$, el haz
$\SheafAut_{Y_1}(y)$ (\resp $\SheafHom_{Y_1}(y,y_1)$) es
constructible, pues tal haz es, localmente para la topología étale
de~$Y_1$, imagen inversa de~$f_{T*}(\SheafAut_{X_T}(p))$
(\resp de~$f_{Z*}(\SheafHom_{X_Z}(p_1,p_2))$).

Queda por demostrar que $(\Phi,f)$ es cohomológicamente propio
relativamente
\marginpar{405}
a~$S$ en dimensión $\leqslant 1$. Para ello basta demostrar que,
\ifthenelse{\boolean{orig}}
{para todo $S'$-esquema~$S$ y para todo punto geométrico
$\overline{y}'$ de~$Y$, si se denota por~$\overline{Y}$ la localización
estricta de~$Y$ en~$y'$, $\overline{X}=X\times_Y\overline{Y}$, etc.,}
{para todo $S$-esquema~$S'$ y para todo punto geométrico
$\overline{y}'$ de~$Y'$, si $\overline{y}$ denota su imagen en~$Y$ y
$\overline{Y}$ (\resp $\overline{Y'}$) denota la localización estricta
de~$Y$ en~$\overline{y}$ (\resp de~$Y'$ en~$\overline{y}'$), y se pone
$\overline{X}=X\times_Y\overline{Y}$,
$\overline{X'}=X'\times_{Y'}\overline{Y'}$, etc.,}
entonces el funtor canónico
$$
\overline{\varphi} \colon \Phi(\overline{X}) \to
\Phi'(\overline{X'})
$$
es una equivalencia de categorías.

Demostremos que~$\overline{\varphi}$ es plenamente fiel. Sean
$x$, $y \in \Phi(\overline{X})$, y sean $x'$, $y'$ sus imágenes en
$\Phi'(\overline{X'})$; demostremos que el morfismo canónico
\begin{equation*}
\label{eq:XIII.4.*} \tag{$*$} {\Hom_{\overline{X}}(x,y)
\to \Hom_{\overline{X'}}(x',y') }
\end{equation*}
es biyectivo. Por definición de~$T$, existen dos morfismos de
$\overline{Y}$ en~$T$ tales que $x$ e~$y$ son las imágenes inversas
de~$p$ por esos dos morfismos. Esto equivale a decir que existe un
morfismo $\overline{Y} \to Z$ y, por tanto, un morfismo
$h\colon \overline{X} \to X_Z$ tal que
$$
h^*(p_1)=x \qquad h^*(p_2)=y.
$$
Por consiguiente, se tiene un isomorfismo canónico
$$
\SheafHom_{\overline{X}}(x,y)=h^*(\SheafHom_{X_Z}(p_1,p_2) )\quoi.
$$
Pero, teniendo en cuenta que $(\SheafHom_{X_Z}(p_1,p_2),f_Z)$ es
cohomológicamente propio relativamente a~$S$ en dimensión
$\leqslant 0$, se ve que lo mismo vale para
$(\SheafHom_{\overline{X}}(x,y),\overline{f})$, lo que demuestra que
el morfismo~\eqref{eq:XIII.4.*} es biyectivo.

Demostremos que~$\overline{\varphi}$ es esencialmente sobreyectivo. Sea
$x'\in \Phi'(\overline{X'})$. Como $(S\Phi,f)$ es cohomológicamente
propio relativamente a~$S$ en dimensión $\leqslant 0$, el morfismo
canónico
$$
\H^0(\overline{X},\overline{S\Phi}) \to
\H^0(\overline{X'},\overline{S\Phi}')
$$
es biyectivo. Sea $G'$ el subgerbe maximal de~$\Phi'$ generado
por~$x'$; existe entonces un subgerbe maximal~$G$ de
$\overline{\Phi}$ cuya imagen inversa sobre~$\overline{X'}$ es~$G'$.
Según 4)\kern2pt1, $(G,\overline{f})$ es cohomológicamente propio
relativamente a~$S$ en dimensión $\leqslant 1$; por consiguiente, el
funtor canónico
$$
G(\overline{X}) \to G'(\overline{X'})
$$
\marginpar{406}%
es una equivalencia de categorías, lo que demuestra la existencia
de un objeto~$x$ de~$\Phi(\overline{X})$ cuya imagen en
\ifthenelse{\boolean{orig}}
{$\Phi'(X')$}
{$\Phi'(\overline{X'})$}
es isomorfa a~$x'$ y concluye la demostración del teorema.

\shorthandoff{<>}% Babel español trata como abreviaturas los delimitadores de curvatura de Xy-pic.

\setcounter{subsection}{1}
\setcounter{subsubsection}{0}
\begin{sublemme}
\label{XIII.3.1.1}
Sean $S$ un esquema localmente noetheriano y
$$
\xymatrix@C=.5cm{ \Phi \ar[r]^p & \Phi_1 \ar@<0.6ex>[rr]^{p_1}
\ar@<-0.6ex>[rr]_{p_2} && \Phi_2 }
$$
un diagrama exacto de stacks sobre~$S$ \eqref{XIII.1.10.1}. Si $\Phi_1$
es constructible, lo mismo ocurre con~$\Phi$. Si, para todo esquema
$S'$ étale sobre~$S$ y todo par
\ifthenelse{\boolean{orig}}
{de objeto}
{de objetos}
$x_1,~y_1$ de~$(\Phi_1)_{S'}$, el haz $\SheafHom_{S'}(x_1,y_1)$ es
constructible, entonces, para todo par de objetos $x,~y$ de
$\Phi_{S'}$, lo mismo ocurre con~$\SheafHom_{S'}(x,y)$. Supongamos que
$\Phi_1$ es $1$-constructible (\Ref{XIII.0}) y que $\Phi_2$ es
constructible; entonces $\Phi$ es $1$-constructible.
\end{sublemme}

Para todo esquema $S'$ étale sobre~$S$ y todo objeto $x$ de
$\Phi_{S'}$, se tiene un monomorfismo
$$
\SheafAut_{S'}(x)\to\SheafAut_{S'}(p(x)).
$$
Resulta, pues, de SGA~4 IX~2.9 que, si $\Phi_1$ es constructible,
lo mismo ocurre con~$\Phi$, y la segunda afirmación del lema se
demuestra del mismo modo.

Supongamos ahora que $\Phi_1$ es $1$-constructible y que $\Phi_2$ es
constructible. El morfismo $p$ induce entre los haces de los subgerbes
maximales un morfismo
$$
\varphi\colon S\Phi\to S\Phi_1.
$$
Sea $G$ la imagen de~$S\Phi$ por~$\varphi$; según SGA~4 IX~2.9,
$G$ es un haz constructible. Puede encontrarse, por tanto, un haz
representable por un $S$-esquema étale de tipo finito~$T$ y un
epimorfismo
$$
a\colon T\to G
$$
(SGA~4 IX~2.7).
\marginpar{407}
Además, puede elegirse~$T$ de modo que la imagen~$y$ de la sección
identidad de~$T(T)$ por~$a$ esté definida por un objeto~$x_1$ de
$(\Phi_1)_T$ de la forma $x_1=p(x)$, donde $x\in\Phi_T$.

Basta demostrar que, para todo punto $s$ de~$S$, existe un abierto no
vacío~$U$ de~$\overline{\{ s\} }$ tal que $S\Phi|U$ sea localmente
constante y constructible. Sean $s\in S$, $\overline{s}$ un punto
geométrico sobre~$s$ y $\overline{y}_1,\dots,\overline{y}_n$ los
elementos de~$G_{\overline{s}}$. Por definición de~$T$, existen
morfismos $h_i\colon\overline{s}\to T$ tales que
$h_i^*(y)=\overline{y}_i$. Sea $S'$ el producto fibrado sobre~$S$ de
$n$ esquemas isomorfos a~$T$, sea $h\colon\overline{s}\to S'$ el
producto fibrado de los~$h_i$, y sean $y_i$ (\resp $x_i$) la imagen
inversa de~$y$ (\resp $x$) por la $i$-ésima proyección de~$S'$ sobre~$T$.
Sea $F_i$ el subhaz de~$S\Phi |S'$ que es la imagen inversa de~$y_i$, y demostremos
que los~$F_i$ son constructibles. El haz~$F_i$ es un cociente del haz
$F'_i$ tal que, para todo esquema~$S''$ étale sobre~$S'$, se tiene
\begin{multline*}
F'_i(S'') = \{\text{clases, salvo isomorfismo, de}
\ifthenelse{\boolean{orig}}
{\text{ objeto $z$ de $\Phi_{S''}$ provisto}}
{\text{ objetos $z$ de $\Phi_{S''}$ provistos}}\\
\ifthenelse{\boolean{orig}}
{\text{ de un isomorfismo $i\colon p(z)\simeq p(x_i|S'')$}.}
{\text{ de un isomorfismo $i\colon p(z)\simeq p(x_i|S'')$}\}.}
\end{multline*}
Basta demostrar que los~$F_i'$ son constructibles. Ahora bien, si se
pone $z_i=p_1p(x_i)$, $z_i'=p_2p(x_i)$, se tiene un monomorfismo
$$
\Psi_i\colon F'_i\to\SheafIsom_{S'}(z_i,z_i'),
$$
obtenido asociando, a todo esquema~$S''$ étale sobre~$S'$ y a todo
objeto~$z$ de~$\Phi_{S''}$ provisto de un isomorfismo
$i\colon p(z)\simeq p(x_i|S'')$, el isomorfismo de~$z_i$ en~$z_i'$
definido por la condición de que el diagrama
$$
\xymatrix@C=1.2cm{
p_1p(z) \ar[r]^-{p_1(i)}\ar[d]_-j& z_i\ar[d]\\ p_2p(z) \ar[r]^-{p_2(i)}& z'_i
}
$$
sea conmutativo ($j$ es el morfismo canónico asociado al diagrama
exacto).
\marginpar{408}

El morfismo~$\Psi_i$ es inyectivo, pues decir que dos objetos $z,~z'$,
provistos de isomorfismos $i\colon p(z)\isomto
p(x_i|S''),~i'\colon p(z')\isomto p(x_i|S'')$, definen el mismo
elemento de~$\SheafIsom_{S''}(z_i,z'_i)$ equivale a decir que
$p_1(i'^{-1}i)=p_2(i'^{-1}i)$, es decir, que $i'^{-1}i$ proviene de
un isomorfismo $z\to z'$. Como el haz~$F'_i$ es un subhaz de
$\SheafIsom_{S'}(z_i,z'_i)$, es constructible.

Como puede encontrarse un abierto no vacío~$U$ de~$\overline{\{s\}}$
tal que $y_1|U,\dots,y_n|U$ engendren~$G$, resulta del
lema~\Ref{XIII.6.1.2} que figura más abajo que $S\Phi |U$ es
constructible; por tanto, tras restringir~$U$, $S\Phi |U$ es localmente
constante y constructible.


\begin{sublemme}
\label{XIII.3.1.2}
Sean $S$ un esquema localmente noetheriano, $f\colon X\to S$ un
morfismo, $F\to G$ un monomorfismo de haces de grupos sobre~$X$,
$\Psi$ (\resp $\Psi_1$) el stack de los torsores bajo~$F$ (\resp bajo~$G$),
$\Phi=f_*\Psi$, $\Phi_1=f_*\Psi_1$. Se supone dado un haz sobre~$S$
representable por un $S$-esquema étale de tipo finito~$T$ y un morfismo
sobreyectivo
$$
a\colon T\to S\Phi_1\simeq\R^1f_*G,
$$
de modo que exista un torsor~$Q$ sobre~$X_T=X\times_ST$, bajo~$G|X_T$,
cuya clase en~$\R^1f_*G(T)$ sea la imagen por~$a$ de la sección identidad
de~$T(T)$. Sea $f_T\colon X_T\to T$ el morfismo canónico y pongamos
$F_T=F|X_T$. Se supone que $\Phi_1$ es $1$-constructible y que
$f_{T*}(Q/F_T)$ es constructible. Entonces $\Phi$ es $1$-constructible.
\end{sublemme}

En efecto, basta repetir la demostración de~\Ref{XIII.3.1.1},
sustituyendo la existencia de los morfismos~$\Psi_i$ por la de los
isomorfismos
$$
\ifthenelse{\boolean{orig}}
{F'_i\isomto f_*(Q/F_T)|S'.}
{F'_i\isomto f_{T*}(Q/F_T)|S'.}
$$


\begin{subremarque}
\label{XIII.3.1.3}
Supongamos
\marginpar{409}
que $k=\kres(s)$ es de característica cero; entonces los esquemas de
tipo finito sobre~$\overline{k}$ de dimensión $\leqslant\dim X$ son
fuertemente desingularizables, y la demostración de~\Ref{XIII.3.1}
permite probar los resultados siguientes:
\begin{enumerate}
\item[a)] Existe un abierto no vacío~$S_1$ de~$S$ tal que, para todo
esquema~$S'$ sobre~$S_1$ cuyos puntos maximales son de característica
cero y para todo haz de conjuntos localmente constante y constructible
$F$ sobre~$X'=X\times_SS'$, $(F,f')$ sea cohomológicamente propio
relativamente a~$S'$ en dimensión $\leqslant 0$.

\item[b)] Si todas las características residuales de~$S$ son nulas,
existe un abierto no vacío~$S_1$ de~$S$ tal que, para todo esquema~$S'$
sobre~$S_1$ y para todo haz de grupos localmente constante y constructible
$F$ sobre~$X'$, $(F,f')$ sea cohomológicamente propio relativamente
a~$S'$ en dimensión $\leqslant 1$.
\end{enumerate}

En efecto, basta repetir la demostración de~\Ref{XIII.3.1}\ifthenelse{\boolean{orig}}
{.}
{~}%
2)\kern2pt a). La proposición~\Ref{XIII.2.7} utilizada en 3)\kern2pt 4
se aplica al caso de un haz localmente constante~$F$, pues, retomando las
notaciones de 3)\kern2pt4, todo torsor bajo~$F$ está moderadamente
ramificado sobre~$P$ relativamente a~$S$, ya que todas las características
residuales de~$S$ son nulas.
\end{subremarque}

\shorthandoff{<>}% Babel español trata como abreviaturas los delimitadores de curvatura de Xy-pic.

\begin{corollaire}
\label{XIII.3.2}
Sean $k$ un cuerpo de característica $p\geq 0$, $p'$ el conjunto de los
números primos distintos de~$p$ y $f\colon X\to k$ un morfismo coherente.
\begin{enumerate}
\item[1)] Para todo haz de conjuntos~$F$, $(F,f)$ es cohomológicamente
propio en dimensión $\leqslant 0$.
\item[2)] Supongamos que se cumple una de las dos condiciones siguientes:
\begin{enumerate}
\item[a)] $f$ es de tipo finito y los esquemas de tipo finito de dimensión
$\leqslant\dim X$ sobre una clausura algebraica de~$k$ son fuertemente
desingularizables.
\item[b)] Los esquemas de tipo finito sobre una clausura algebraica de~$k$
son
\marginpar{410}
fuertemente desingularizables.
\end{enumerate}
Entonces, para todo haz de ind-$p'$-grupos~$F$, $(F,f)$ es
cohomológicamente propio en dimensión $\leqslant 1$.
\end{enumerate}
\end{corollaire}

Sea $F$ un haz de conjuntos (\resp de ind-$p'$-grupos). Según
SGA~4 IX~2.7.2, puede escribirse~$F$ como límite inductivo filtrante
$$
F=\varinjlim F_i,
$$
donde los~$F_i$ son haces de conjuntos (\resp de ind-$p'$-grupos)
constructibles. Como $f$ es coherente, $f_*$ (\resp $\R^1f_*$) conmuta
con los límites inductivos (SGA~4 VII~3.3). Si se sabe que los~$(F_i,f)$
son cohomológicamente propios en dimensión $\leqslant 0$ (\resp que todo
haz de conjuntos es cohomológicamente propio en dimensión $\leqslant 0$
y que los~$(F_i,f)$ son cohomológicamente propios en dimensión
$\leqslant 1$), lo mismo valdrá para~$(F,f)$. Puede suponerse, por tanto,
que~$F$ es constructible.

Si se supone que~$f$ es de tipo finito (\resp que satisface la condición
a)), el resultado se sigue de~\Ref{XIII.3.1}~1)\kern2pt b) (\resp de
\Ref{XIII.3.1}~2)\kern2pt b)). Demostremos ahora~\Ref{XIII.3.2} sin
suponer ya que~$f$ sea de tipo finito. Para todo esquema~$S'$ sobre~$k$
y todo punto geométrico~$\overline{s}$ de~$S'$, se denota por
$\overline{k}$ (\resp $\overline{S'}$) la localización estricta de~$k$
en~$\overline{s}$ (\resp de~$S'$ en~$\overline{s}$), por~$\overline{X}$
la imagen inversa de~$X$ sobre~$\overline{k}$, y se considera el cuadrado
cartesiano
$$
\xymatrix@C=1.2cm{
\overline{X}\ar[d]_-{\overline{f}}&\ar[l]_g \overline{X'}\ar[d]^-{\overline{f}'}\\ \overline{k}&\ar[l]\overline{S'}
}
$$
Basta demostrar que, para todo~$S'$ y todo~$\overline{s}$, el morfismo
canónico
$$
\H^0(\overline{X},\overline{F})\to
\H^0(\overline{X'},\overline{F'})\qquad
\text{(\resp $\H^1(\overline{X},\overline{F})\to
\H^1(\overline{X'},\overline{F'})$)}
$$
es un isomorfismo. Basta demostrar que se tienen las relaciones
\marginpar{411}
\begin{equation*}
\label{eq:XIII.3.*}
\tag{$*$} \overline{F}\simeq g_*g^*\overline{F}\qquad
\ifthenelse{\boolean{orig}}
{\text{(\resp $\R^1g_*(g^*\overline{F}))=0$).}}
{\text{(\resp $\R^1g_*(g^*\overline{F})=0$).}}
\end{equation*}
Ahora bien, en esta forma la cuestión es local sobre~$X$ para la topología
étale. Puede suponerse, por tanto, que~$X$ es afín; por paso al límite,
puede suponerse que~$X$ es de tipo finito sobre~$k$. Se sabe entonces que
$(F,f)$ es cohomológicamente propio en dimensión $\leqslant 0$ (\resp
$\leqslant 1$), y que lo mismo ocurre al reemplazar~$X$ por un esquema
étale de tipo finito sobre~$X$, lo que demuestra~\eqref{eq:XIII.3.*}.

\shorthandoff{<>}% Babel español trata como abreviaturas los delimitadores de curvatura de Xy-pic.

\begin{theoreme}
\label{XIII.3.3}
Sean $S$ un esquema irreducible con punto genérico~$s$ y
$f\colon X\to S$ un morfismo de presentación finita. Supongamos que los
esquemas de tipo finito de dimensión $\leqslant\dim X_s$ sobre una clausura
algebraica~$\bar{k}$ de
\ifthenelse{\boolean{orig}}
{$k$}
{$\kres(s)$}
son desingularizables \textup{(EGA~IV 7.9.1)}. Entonces, si~$\LL$ designa
el conjunto de los números primos distintos de las características
residuales de~$S$, puede encontrarse un abierto no vacío~$S_1$ de~$S$ tal
que el morfismo~$f|S_1$ sea universalmente localmente $1$-asférico
para~$\LL$.
\end{theoreme}

Puede suponerse que~$S$ es íntegro y que~$X$ es reducido (SGA~4 VIII~1.1).
Por paso al límite, puede suponerse que~$S$ es noetheriano. Además, para
demostrar el teorema basta hacerlo tras una extensión finita de la base
$S_1\to S$, donde~$S_1$ es un esquema íntegro y~$S_1\to S$ es composición
de extensiones étale y de extensiones radiciales sobreyectivas.

Demostremos primero que, tras restringir~$S$ a un abierto no vacío, $f$ es
universalmente localmente $0$-acíclico. Tras efectuar una extensión
radicial de~$\kres(s)$, puede suponerse que el morfismo
$(X_s)_\red\to s$ es separable; por tanto, tras restringir~$S$ a un abierto
no vacío y efectuar una extensión radicial sobreyectiva de~$S$, puede
suponerse que el morfismo~$f$ es plano con fibras geométricas separables
(EGA~IV 12.1.1), lo que implica que~$f$ es
\ifthenelse{\boolean{orig}}
{universalmente $0$-acíclico}
{universalmente localmente $0$-acíclico}
(SGA~4 XV~4.1).

Demostremos que, tras restringir~$S$ a un abierto no vacío, $f$ es
universalmente localmente $1$-asférico para~$\LL$. Tras efectuar una
\marginpar{412}
extensión finita de~$\kres(s)$, lo cual es lícito porque dicha extensión
puede factorizarse como una extensión étale seguida de una extensión
radicial, puede
encontrarse un morfismo propio sobreyectivo $p_s\colon Y_s\to X_s$, donde
$Y_s$ es un esquema liso sobre
\ifthenelse{\boolean{orig}}
{$S$}
{$s$},
de la misma dimensión que~$X_s$, y puede suponerse que~$p_s$ proviene de
un morfismo propio sobreyectivo $p\colon Y\to X$, donde~$Y$ es un esquema
liso sobre~$S$ (EGA~IV 9.6.1 y 12.1.6). Basta demostrar que, tras restringir
$S$ a un abierto no vacío, para todo diagrama de cuadrados cartesianos
$$
\xymatrix{ S''\ar[d]_i & X ''\ar[d]^j\ar[l]_{f''} \\ S'\ar[d] &
X'\ar[d]\ar[l]_{f'} \\ S & X, \ar[l]_f }
$$
donde~$i$ es étale de presentación finita, y para todo haz de
ind-$\LL$-grupos~$F$ sobre~$S''$, si~$\Phi$ es el stack de los torsores
bajo~$F$, entonces el morfismo canónico
$$
{f'}^*i_*\Phi \to j_*{f''}^*\Phi
$$
es una equivalencia. Sean $Z=Y\times_XY,
T=Y\times_XY\times_XY$. Se tiene de manera natural un diagrama
conmutativo
$$
\xymatrix{ S'' \ar[d]_i & X''\ar[l]_{f''}\ar[d]_j &
Y''\ar[l]_{p''}\ar[d] & Z''\ar[d]\ar@<-.4ex>[l]\ar@<.4ex>[l] &
T''\ar[d]\ar[l]\ar@<-.8ex>[l]\ar@<.8ex>[l] \\ S' \ar[d] & X'
\ar[l]_{f'}\ar[d] & Y'\ar[d]\ar[l]_{p'}&
Z'\ar[d]\ar@<-.4ex>[l]\ar@<.4ex>[l] &
T'\ar[d]\ar[l]\ar@<-.8ex>[l]\ar@<.8ex>[l] \\ S & X\ar[l]_f & Y\ar[l]_p
& Z\ifthenelse{\boolean{orig}}{'}{}\ar@<-.4ex>[l]\ar@<.4ex>[l]
& T\ifthenelse{\boolean{orig}}{'}{}\ar[l]\ar@<-.8ex>[l]\ar@<.8ex>[l]
}
$$
Sean
\marginpar{413}%
$q\colon Z\to X$ y $r\colon T\to X$ los morfismos canónicos; las
notaciones $q',r',q'',r''$ tienen un sentido evidente. Según
\Ref{XIII.1.11}~2), se tiene el siguiente diagrama esencialmente
conmutativo, cuyas filas son exactas:
$$
\ifthenelse{\boolean{orig}}
{\xymatrix{ f'^* i_*\Phi \ar[r]\ar[d]_a &
p'_*p'^*(f'^*i_x\Phi)\ar[d]_b\ar@<-.4ex>[r]\ar@<.4ex>[r] &
q'_*q'^*(f'^*i_*\Phi)\ar[d]_c\ar[r]\ar@<-.8ex>[r]\ar@<.8ex>[r]&
r'_*r'^*(f'^*i_*\Phi)\ar[d]_d \\ j_*f''^*\Phi \ar[r]&
j_*p''_*p''^*(f''^*\Phi)\ar@<-.4ex>[r]\ar@<.4ex>[r] &
j_*q''_*q''^*(f''^*\Phi)\ar[r]\ar@<-.8ex>[r]\ar@<.8ex>[r]&
j_*r''_*r''^*(f''^*\Phi) }}
{\xymatrix{ f'^* i_*\Phi \ar[r]\ar[d]_a &
p'_*p'^*(f'^*i_*\Phi)\ar[d]_b\ar@<-.4ex>[r]\ar@<.4ex>[r] &
q'_*q'^*(f'^*i_*\Phi)\ar[d]_c\ar[r]\ar@<-.8ex>[r]\ar@<.8ex>[r]&
r'_*r'^*(f'^*i_*\Phi)\ar[d]_d \\ j_*f''^*\Phi \ar[r]&
j_*p''_*p''^*(f''^*\Phi)\ar@<-.4ex>[r]\ar@<.4ex>[r] &
j_*q''_*q''^*(f''^*\Phi)\ar[r]\ar@<-.8ex>[r]\ar@<.8ex>[r]&
j_*r''_*r''^*(f''^*\Phi) }}
$$
Como~$Y$ es liso sobre~$S$, el morfismo~$fp$ es universalmente localmente
$1$-asférico para~$\LL$
\ifthenelse{\boolean{orig}}
{(SGA~XV~2.1),}
{(SGA~4 XV~2.1),}
y de~\cite[VII 2.1.7]{XIII.2} resulta que~$b$ es una equivalencia de
categorías. Por otra parte, tras restringir~$S$ a un abierto no vacío,
puede suponerse que los morfismos $Z\to S$ y $T\to S$ son universalmente
localmente $0$-acíclicos. Resulta de ello que los funtores~$c$ y~$d$ son
plenamente fieles, y el diagrama anterior muestra entonces que~$a$ es una
equivalencia, lo que concluye la demostración.


\begin{corollaire}
\label{XIII.3.4}
Sean $k$ un cuerpo de característica $p\geqslant0$, $p'$ el conjunto de
los números primos distintos de~$p$ y $f\colon X\to k$ un morfismo
coherente. Supongamos que se cumple una de las condiciones siguientes:
\begin{enumerate}
\item[a)] $f$ es de tipo finito y los esquemas de tipo finito de dimensión
$\leqslant\dim X$ sobre una clausura algebraica de~$k$ son
desingularizables.
\item[b)] Los esquemas de tipo finito sobre una clausura algebraica de~$k$
son desingularizables.
\end{enumerate}
Entonces $f$ es universalmente localmente $1$-asférico para~$p'$.
\end{corollaire}

El caso a) resulta de~\Ref{XIII.3.3}. En el caso b), como la cuestión es
local sobre~$X$, puede suponerse que~$X$ es afín; por paso al límite
(SGA~4 XV~1.3), se reduce al caso en que~$X$ es de tipo finito sobre~$k$.

\begin{corollaire}
\label{XIII.3.5}
Sean
\marginpar{414}
$S$ un esquema irreducible con punto genérico~$s$ y $f\colon X\to S$ un
morfismo de presentación finita. Supongamos que los esquemas de tipo finito
de dimensión $\leqslant\dim X_s$ sobre una clausura
algebraica~$\overline{k}$ de~$\kres(s)$ son fuertemente desingularizables
\textup{(SGA~5 I~3.1.5)}. Si~$\LL$ designa el conjunto de los números primos
distintos de las características residuales de~$S$, puede encontrarse un
abierto no vacío~$S_1$ de~$S$ tal que, para toda especialización
$\bar{s}_1\to\bar{s}_2$ de puntos geométricos de~$S_1$, el morfismo de
especialización~\eqref{XIII.2.10}
$$
\pi_1^{\LL}(X_{\bar{s}_1})\to\pi_1^{\LL}(X_{\bar{s}_2})
$$
sea biyectivo.
\end{corollaire}

Según~\Ref{XIII.3.1} y~\Ref{XIII.3.3}, tras restringir~$S$ a un abierto no
vacío puede suponerse que~$f$ es localmente $1$-asférico para~$\LL$, y que,
para todo haz constante finito de~$\LL$-grupos~$F$ sobre~$X$, $(F,f)$ es
cohomológicamente propio en dimensión $\leqslant 1$. Resulta entonces
de~\Ref{XIII.1.14} que, para toda especialización
$\bar{s}_1\to\bar{s}_2$ de puntos geométricos de~$S_1$, el morfismo de
especialización
$$
\left(\R^1 f_* F\right)_{\bar{s}_2} \to \left(\R^1 f_*
F\right)_{\bar{s}_1}
$$
es biyectivo. El corolario no es sino la traducción de lo anterior en
términos de grupos fundamentales.

\section{Sucesiones exactas de homotopía}\label{XIII.4}\setcounter{subsection}{-1}
\subsection{}
\label{XIII.4.0}
Sean $X$ y $S$ dos esquemas conexos, $f\colon X\to S$ un morfismo, $a$ un
punto geométrico de~$X$ y $\LL$ un conjunto de números primos. Sea~$K$ el
núcleo del homomorfismo canónico $\pi_1(X,a)\to\pi_1(S,a)$ y~$N$ el menor
prosubgrupo normal de~$K$ tal que $K/N$ sea un grupo pro-$\LL$,
denotado por~$K^{\LL}$. Entonces $N$ es normal en~$\pi_1(X,a)$, y se denota por
$$
\pi_1'(X,a)
$$
\label{indnot:mj}\oldindexnot{$\pi_1'(X,a)$|hyperpage}%
el
\marginpar{415}
cociente de~$\pi_1(X,a)$ por~$N$. Si~$a$ es un punto geométrico de una
fibra geométrica~$X_{\bar{s}}$, los morfismos canónicos
$$
\pi_1(X_{\bar{s}},a)\to \pi_1(X,a) \to
\pi_1(S,a)
$$
permiten definir morfismos canónicos
$$
\pi_1^{\LL}(X_{\bar{s}},a)\lto{u} \pi_1'(X,a) \lto{v} \pi_1(S,a)
$$
Se tiene $vu=0$.


\begin{proposition}
\label{XIII.4.1}
\setcounter{subsection}{\value{proposition}}% sincroniza la etiqueta visible XIII.4.1.1
Sean $S$ un esquema conexo y $f\colon X\to S$ un morfismo localmente
$0$-acíclico \textup{(SGA~4 XV~1.11)}; supongamos además que~$f$ es
$0$-acíclico (lo que, cuando~$f$ es coherente, equivale a decir que las
fibras geométricas de~$f$ son conexas \textup{(SGA~4 XV~1.16)}). Sea~$\LL$
un conjunto de números primos. Si~$S'$ es un esquema étale sobre~$S$, se
denotan por~$X'$ y~$f'$ las imágenes inversas de~$X$ y~$f$ sobre~$S'$.
Supongamos que, para todo recubrimiento étale~$S'$ de~$S$ y para todo
recubrimiento étale~$E$ de~$X'$, cociente de un recubrimiento galoisiano
cuyo grupo es un~$\LL$-grupo, $(E,f')$ es cohomológicamente propio
relativamente a~$S'$ en dimensión $\leqslant 0$ y~$f'_*E$ es constructible.
Entonces, si~$\bar{s}$ es un punto geométrico de~$S$ y~$a$ un punto
geométrico de la fibra~$X_{\bar{s}}$, la sucesión de homomorfismos de grupos
\begin{equation*}
\label{eq:XIII.4.1.1}
\tag{\thesubsection.1}
\pi_1^{\LL}(X_{\bar{s}},a)\lto{u} \pi_1'(X,a)\lto{v}\pi_1(S,a)\to1
\end{equation*}
es exacta.
\end{proposition}

Este enunciado generaliza X.\Ref{X.1.4}, cuya demostración vamos a repetir.

Demostremos primero que~$v$ es sobreyectivo. Basta demostrar que, para todo
recubrimiento étale conexo~$S'$ de~$S$, $X'$ también es conexo
(\Ref{V.6.9}). Sea~$C$ un conjunto con al menos dos elementos. Del hecho
de que~$f$ es $0$-acíclico resulta que el morfismo canónico
$$
\H^0(S',C_{S'})\to \H^0(X',C_{X'})
$$
es biyectivo; por tanto, $X'$ es conexo, de donde se sigue la sobreyectividad
de~$v$.

Por
\marginpar{416}
definición de~$K^{\LL}$ \eqref{XIII.4.0}, se tiene la sucesión exacta
$$
1\to K^{\LL} \to \pi_1'(X,a)\to \pi_1(S,a)\to 1
$$
Sean $\widetilde{S}$ el recubrimiento universal de~$S$ y
$\widetilde{X}=\widetilde{S}\times_SX$; el grupo~$K^{\LL}$ clasifica los
recubrimientos galoisianos~$P$ cuyo grupo es un~$\LL$-grupo y para los
cuales existen un recubrimiento étale~$S'$ de~$S$ y un recubrimiento
galoisiano~$Q$ de~$X'=X\times_SS'$ tales que se tiene un isomorfismo
$P\simeq Q\times_{X'}\widetilde{X}$. Para que la sucesión
\eqref{eq:XIII.4.1.1} sea exacta, es necesario y suficiente que el morfismo
canónico
$$
\pi_1^{\LL}(X_{\bar{s}},a)\to K^{\LL}
$$
sea sobreyectivo. Según la interpretación de~$K^{\LL}$, esto equivale a
decir que, para todo recubrimiento étale~$S'$ de~$S$ y para todo
recubrimiento galoisiano~$Q$ de~$X'$ cuyo grupo es un~$\LL$-grupo, si
$P=Q\times_{X'}\widetilde{X}$ es conexo, entonces $Q|X_{\bar{s}}$ es
conexo. Demostremos que se cumple esta última condición. Sean, en efecto,
$S'$ un recubrimiento étale de~$S$ y~$Q$ un recubrimiento galoisiano de~$X'$
cuyo grupo es un~$\LL$-grupo~$F$, tal que~$Q|X_{\bar{s}}$ es disconexo, y
demostremos que, tras reemplazar~$S'$ por un recubrimiento étale, $Q$ pasa
a ser disconexo. Existe un subgrupo~$G$ de~$F$, distinto de~$F$, y un
torsor~$R$ bajo~$G|X_{\bar{s}}$ tal que~$Q|X_{\bar{s}}$ se obtiene, por
extensión del grupo estructural $G\to F$, a partir de~$R$. El recubrimiento
étale
$E=Q/G$ de~$X'$ es tal que~$E|X_{\bar{s}}$ tiene una sección. Según
\Ref{XIII.1.16},
\ifthenelse{\boolean{orig}}
{$f_*E$}
{$f'_*E$}
es localmente constante y constructible y, tras reemplazar~$S'$ por un
recubrimiento étale, puede incluso suponerse que
\ifthenelse{\boolean{orig}}
{$f_*E$}
{$f'_*E$}
es constante.
Como~$(E,f')$ es cohomológicamente propio relativamente a~$S'$ en dimensión
$\leqslant0$ y~$\H^0(X_{\bar{s}},E|X_{\bar{s}})$ no es vacío, se ve que~$E$
tiene una sección. Pero esto demuestra que~$Q$ es disconexo, lo que
concluye la demostración.

De~\Ref{XIII.4.1} se deduce el lema siguiente, que se utilizará
en~\Ref{XIII.4.6}.


\begin{lemme}
\label{XIII.4.2}
Sean $S$ un esquema conexo, $f\colon X\to S$ un morfismo localmente
$0$-acíclico y $0$-acíclico, y~$\LL$ un conjunto de números primos. Supongamos
\marginpar{417}
que, para todo haz constante finito de~$\LL$-grupos~$F$ sobre~$X$,
$(F,f)$ es cohomológicamente propio relativamente a~$S$ en dimensión
$\leqslant1$, y que, para todo recubrimiento étale~$S'$ de~$S$ y para todo
recubrimiento étale~$E$ de~$X'$, cociente de un recubrimiento galoisiano
cuyo grupo es un~$\LL$-grupo, $f'_*E$ es constructible. Entonces, si~$\bar{s}$
es un punto geométrico de~$S$ y~$a$ un punto geométrico de la fibra
$X_{\bar{s}}$, la sucesión de homomorfismos de grupos
$$
\pi_1^{\LL}(X_{\bar{s}},a)\to\pi_1'(X,a)\to\pi_1(S,a)\to1
$$
es exacta.
\end{lemme}

Se satisfacen las hipótesis de~\Ref{XIII.4.1}. En efecto, de
\Ref{XIII.1.13}~3) se sigue que, para todo esquema~$S'$ étale sobre~$S$ y
para todo recubrimiento étale~$E$ de~$X'$, cociente de un recubrimiento
galoisiano cuyo grupo es un~$\LL$-grupo, $(E,f')$ es cohomológicamente propio
relativamente a~$S'$ en dimensión $\leqslant 0$.

\begin{proposition}
\label{XIII.4.3}
\setcounter{subsection}{\value{proposition}}% sincroniza la etiqueta visible XIII.4.3.1
Sean $S$ un esquema conexo, $\LL$ un conjunto de números primos,
$f\colon X\to S$ un morfismo $0$-acíclico y localmente $1$-asférico
para~$\LL$ \textup{(SGA~4 XV~1.11)}, y~$g\colon S\to X$ una sección de~$f$.
Sean~$\bar{s}$ un punto geométrico de~$S$ y~$a$ un punto geométrico de la
fibra~$X_{\bar{s}}$. Supongamos que, para todo haz de
\ifthenelse{\boolean{orig}}
{$\LL$-grupo}
{$\LL$-grupos}
constante~$F$, $(F,f)$ es cohomológicamente propio en dimensión
$\leqslant1$, que la imagen directa por~$f$ del stack de los torsores bajo
$F$ es un stack $1$-constructible \eqref{XIII.0} y que,
\ifthenelse{\boolean{orig}}
{para todo recubrimiento étale~$E$ de~$X'$}
{para todo recubrimiento étale~$S'$ de~$S$ y para todo
recubrimiento étale~$E$ de~$X'$}, cociente de un recubrimiento galoisiano
cuyo grupo es un~$\LL$-grupo, $f'_*E$ es constructible. Entonces, la sucesión
de homomorfismos de grupos
\begin{equation*}
\label{eq:XIII.4.3.1}
\tag{\thesubsection.1}
1\to\pi_1^{\LL}(X_{\bar{s}},a)\lto{u} \pi_1'(X,a)\lto{v}\pi_1(S,a)\to1
\end{equation*}
es exacta.
\end{proposition}

Habida cuenta de~\Ref{XIII.4.2}, basta demostrar la inyectividad del
morfismo~$u$, es decir, probar que, para todo recubrimiento principal
$\overline{Z}$ de~$X_{\bar{s}}$
\marginpar{418}
cuyo grupo es un~$\LL$-grupo~$C$, existen un recubrimiento étale~$Z$ de~$X$
y un morfismo de una componente conexa de~$Z|X_{\bar{s}}$ hacia
$\overline{Z}$ (\Ref{V.6.8}). Sean, pues, $\overline{Z}$ un recubrimiento
principal de~$X_{\bar{s}}$ cuyo grupo es un~$\LL$-grupo~$C$ y~$\bar{z}$ su
clase en~$\H^1(X_{\bar{s}},C_{X_{\bar{s}}})$. Según \Ref{XIII.1.5}~d), se
tiene un isomorfismo canónico
$$
\left(\R^1f_*C_X\right)_{\bar{s}}\isomto\H^1(X_{\bar{s}},C_{X_{\bar{s}}}),
$$
y, según~\Ref{XIII.1.16}, $\R^1f_*C_X$ es un haz localmente constante y
constructible. Puede encontrarse, por tanto, un recubrimiento étale~$S'$ de~$S$
tal que~$\R^1f_*C_X|S'$ sea constante. Si~$\bar{s}\to S'$ es un punto
geométrico situado sobre el punto geométrico~$\bar{s}\to S$, existe un
elemento~$z$ de~$\H^0(S', \R^1f_*C_X)$ cuya imagen en
$\H^1(X_{\bar{s}},C_{X_{\bar{s}}})$ es~$\bar{z}$. Según el lema
\Ref{XIII.4.3.1} que sigue, puede encontrarse un recubrimiento étale~$S'_1$
de~$S'$ y un torsor~$P$ sobre~$X'_1$ bajo~$C$ cuya imagen en
$\H^0(S'_1,\R^1f_*C)$ sea igual a la restricción de~$z$. El torsor~$P$ es
representable por un recubrimiento étale~$Z$ de~$X'_1=X\times_SS'_1$ tal que
$Z\times_{X'_1}X_{\bar{s}}$ sea isomorfo a~$\overline{Z}$. Si se considera
$Z$ como un recubrimiento étale de~$X$, se tiene entonces un morfismo
de~$Z\times_XX_{\bar{s}}$ hacia~$\overline{Z}$, lo que concluye la
demostración.


\setcounter{subsection}{3}
\setcounter{subsubsection}{0}
\begin{sublemme}
\label{XIII.4.3.1}
Sean $f\colon X\to S$ un morfismo $0$-acíclico y localmente
$0$-acíclico, y~$g$ una sección de~$f$. Sea~$C$ un grupo constante
finito tal que $(C_X,f)$ sea cohomológicamente propio en dimensión
$\leqslant0$ y que la imagen directa por~$f$ del stack de los torsores bajo
$C_X$ sea constructible. Entonces, para toda sección~$z$ de
$\H^0(S,\R^1f_*C_X)$, puede encontrarse un recubrimiento étale
$S_1$ de~$S$ y,
\ifthenelse{\boolean{orig}}
{si $X_1=X\times_SS_1$, un elemento de}
{si $X_1=X\times_SS_1$ y $f_1\colon X_1\to S_1$ es la proyección, un elemento de}
$\H^1(X_1,C_{X_1})$ cuya imagen por el morfismo canónico
$$
\H^1(X_1,C_{X_1})\to
\ifthenelse{\boolean{orig}}
{\H^0(S_1,\R^1f_*C_{X_1})}
{\H^0(S_1,\R^1f_{1*}C_{X_1})}
$$
sea igual a la restricción de~$z$ a
\ifthenelse{\boolean{orig}}
{$\H^0(S_1,\R^1f_*C_{X_1})$}
{$\H^0(S_1,\R^1f_{1*}C_{X_1})$}.
\end{sublemme}

Para todo esquema~$S'$ étale sobre~$S$, se pone $X'=X\times_SS'$ y
\ifthenelse{\boolean{orig}}
{se denota por~$g'$ (\resp $F'$, etc.) la imagen inversa de~$g$
(\resp $F$, etc.) por el morfismo}
{se denota por~$g'$ la imagen inversa de~$g$ por el morfismo}
\marginpar{419}
$S'\to S$. El prehaz~$G$ sobre~$S$ definido por
\newlength{\tmpRacinet} \setlength{\tmpRacinet}{\textwidth}
\addtolength{\tmpRacinet}{-7em}
$$
G(S') = \left\{\text{\parbox{\tmpRacinet}{clases, salvo isomorfismo, de
		torsores $P$ sobre~$X'$ bajo~$C_{X'}$, provistos de un
		isomorfismo $g'^*P\isomto C_{S'}$ }} \right\}
$$
es entonces un haz. Esto resulta, en efecto, por descenso, del hecho
de que un automorfismo del torsor~$P$ sobre~$X'$ queda bien
determinado por su restricción a~$g'(S')$. Además, se tiene un
morfismo sobreyectivo
$$
\ifthenelse{\boolean{orig}}
{G\to \R^1f_*F.}
{G\to \R^1f_*C_X.}
$$
Sean~$z$ un elemento de~$\H^0(S,\R^1f_*C_X)$ y~$H$ el
subhaz de~$G$ que es la imagen inversa de~$z$. Basta demostrar
que~$H$ es un haz localmente constante y constructible. Como esta
propiedad es local sobre~$S$, puede suponerse que~$z$
proviene de un elemento de~$\H^1(X,C_X)$ representado por
un torsor~$P$ tal que $g^*P$ sea isomorfo a~$C_S$. Dar un
isomorfismo $i\colon g^*P\isomto C_S$ equivale a dar una
sección global de~$\SheafAut_{C_S}(g^*P)$, y dos isomorfismos~$i$
e~$i'$ definen el mismo elemento de
\ifthenelse{\boolean{orig}}
{$G(X)$}
{$G(S)$}
si y solo si
\ifthenelse{\boolean{orig}}
{$ii'^{-1}$}
{$i^{-1}i'$}
es la imagen de un elemento de
$\Aut_{C_X}(P)$. Si se considera la inyección canónica
$$
f_*\SheafAut_{C_X}(P)\to\SheafAut_{C_S}(g^*P)\simeq C,
$$
$H$ se identifica, por tanto, con el cociente de
$\SheafAut_{C_S}(g^*P)$ por~$f_*\SheafAut_{C_X}(P)$. Según
\Ref{XIII.1.16}, $f_*\SheafAut_{C_X}(P)$ es localmente constante;
lo mismo ocurre, pues, con~$H$, lo que concluye la demostración.


\begin{exemples}
\label{XIII.4.4}
Obsérvese que, si~$S$ es un esquema conexo, las hipótesis
de~\Ref{XIII.4.1} se satisfacen cuando el morfismo~$f$ es propio y
plano, de presentación finita, con fibras geométricas separables y
conexas, siendo~$\LL$ arbitrario
(\cf~\Ref{X.1.3}). Las hipótesis de~\Ref{XIII.4.3} se
satisfacen si, además, $f$ es liso y posee una sección, y si se
designa por~$\LL$ el conjunto de los números primos distintos de las
características residuales de~$S$ (SGA~4 XV~2.1 y XVI~5.2).

Las hipótesis de~\Ref{XIII.4.1} también se satisfacen si~$S$ es conexo
y existe un esquema~$Z$ propio, de presentación finita y plano sobre~$S$,
con
\marginpar{420}
fibras geométricas separables y conexas, tal que~$X$ sea
el complementario en~$Z$ de un divisor de cruzamientos normales
relativamente a~$S$, siendo~$\LL$ el conjunto de los números
primos distintos de las características residuales de~$S$
\eqref{XIII.2.9}. Las hipótesis de~\Ref{XIII.4.3} se
satisfacen si, además, $f$ es liso y posee una sección.
\end{exemples}


\setcounter{subsection}{4}
\subsection{}
\label{XIII.4.5}
Retomemos las notaciones y las hipótesis de~\Ref{XIII.4.3}. Si
$\bar{s}$ es un punto geométrico de~$S$ y $a=g(\bar{s})$,
la sección~$g$ permite definir un morfismo
$$
w\colon \pi_1(S,a)\to \pi_1'(X,a),
$$
de modo que $\pi_1'(X,a)$ se identifica con el producto semidirecto de
$\pi_1(S,a)$ por
\ifthenelse{\boolean{orig}}
{$\pi_1(X_{\bar{s}},a)$}
{$\pi_1^\LL(X_{\bar{s}},a)$}.
El grupo profinito $\pi_1(S,a)$ actúa, por tanto, sobre
\ifthenelse{\boolean{orig}}
{$\pi_1(X_{\bar{s}},a)$}
{$\pi_1^\LL(X_{\bar{s}},a)$}.
Como $\pi_1^\LL(X_{\bar{s}},a)$ es límite proyectivo estricto de grupos
invariantes bajo la acción de~$\pi_1(S,a)$, dar
$\pi_1^\LL(X_{\bar{s}},a)$ provisto de esta acción equivale
a dar un sistema proyectivo estricto de esquemas
en grupos finitos y étales sobre~$S$, que se denota por
$$
\pi_1^\LL(X/S,g,\bar{s})\quad\text{o
simplemente}\quad\pi_1^\LL(X/S,g).
$$
\label{indnot:mk}\oldindexnot{$\pi_1^\LL(X/S,g,\bar{s})$ o $\pi_1^\LL(X/S,g)$|hyperpage}%

Se tienen entonces las propiedades siguientes:
\setcounter{subsubsection}{0}

\subsubsection{}
\label{XIII.4.5.1}
Para todo esquema en grupos finito y étale~$G$ sobre~$S$, cuyas
fibras son~$\LL$-grupos, el conjunto~$E$ de las clases de torsores
$P$ bajo la imagen inversa~$G_X$ de~$G$ sobre~$X$, provistos de un
isomorfismo $g^*P\isomto G$, es canónicamente isomorfo al conjunto
$$
\Hom_S(\pi_1^\LL(X/S,g,\bar{s}),G) \quad\text{módulo automorfismos
interiores de~$G$.}
$$

\subsubsection{}
\label{XIII.4.5.2}
Para todo esquema en grupos finito y étale~$G$ sobre~$S$ cuyas
fibras son~$\LL$-grupos, el haz $\R^1f_*G_X$ es
canónicamente isomorfo al haz asociado al prehaz
$$
S'\mto \Hom_{S'}(\pi_1^\LL(X/S,g,\bar{s}),G)\
\text{módulo automorfismos interiores de~$G$.}
$$
($S'$ designa un esquema étale sobre~$S$).

\subsubsection{}
\label{XIII.4.5.3}
Sean
\marginpar{421}
$S'$ un $S$-esquema conexo, $\bar{s}$ un punto geométrico de~$S'$, y~$X'$,
$g'$ las imágenes inversas respectivas de~$X,g$ sobre~$S'$. Entonces,
$\pi_1^\LL(X'/S',g',\bar{s})$ es canónicamente isomorfo a
la imagen inversa de~$\pi_1^\LL(X/S,g,\bar{s})$ sobre~$S'$. Para todo
punto geométrico~$\xi$ de~$S$, la fibra
$\pi_1^\LL(X/S,g,\bar{s})_\xi$ es isomorfa a~$\pi_1^\LL(X_\xi)$.

Dar~$G$ equivale, en efecto, a dar un~$\LL$-grupo abstracto~$\mathbf{G}$
sobre el cual actúa~$\pi_1(S,a)$, de donde resulta una acción de
$\pi_1'(X,a)$ sobre~$\mathbf{G}$. El isomorfismo definido
en~\Ref{XIII.4.5.1} se obtiene entonces, por restricción al subconjunto
$E$, a partir del morfismo canónico
\[
\H^1(\pi_1'(X,a),\mathbf{G})\to
\H^1(\pi_1^\LL(X_{\bar{s}},a),\mathbf{G}) =
\Hom(\pi_1^\LL(X_{\bar{s}},a),\mathbf{G})/
\ifthenelse{\boolean{orig}}
{\text{aut. int. }G}
{\text{aut. int. }\mathbf{G}},
\]
y este morfismo envía biyectivamente~$E$ sobre el subconjunto de los
morfismos de~$\pi_1^\LL(X_{\bar{s}},a)$ a~$\mathbf{G}$ compatibles con
la acción de~$\pi_1(S,a)$. La afirmación~\Ref{XIII.4.5.3} resulta
entonces de la definición de~$\pi_1^\LL(X/S,g,\bar{s})$, habida cuenta
de la sucesión exacta de homotopía~\eqref{eq:XIII.4.3.1}, y
\Ref{XIII.4.5.2} se deduce de~\Ref{XIII.4.5.1} y~\Ref{XIII.4.5.3}.


\setcounter{subsection}{6}% sincronización del rótulo visible XIII.4.6.0 de la fuente
\setcounter{proposition}{5}% sincronización con la numeración visible XIII.4.6 de la fuente
\begin{proposition}[Fórmula de Künneth]
\label{XIII.4.6}
\index{Kunneth (formula de)@K\"unneth (fórmula de)|hyperpage}%
Sean $k$ un cuerpo separablemente cerrado de característica
$p\geq 0$, $X$ e~$Y$ dos $k$-esquemas conexos, $a$ un punto
geométrico de~$X$, $b$ un punto geométrico de~$Y$ y~$c$ un punto
geométrico de~$X\times_k Y$ situado sobre~$a$ y~$b$. Se supone que se
satisface una de las dos condiciones siguientes:
\begin{enumerate}
\item[a)] $X$ es de tipo finito sobre~$k$ y los esquemas de tipo finito
sobre una clausura algebraica~$\bar{k}$ de~$k$, de dimensión
$\leq\dim X$, son fuertemente desingularizables. \textup{(SGA~5 I~3.1.5)}.
\item[b)] $X$ es cuasicompacto y cuasiseparado y todo esquema de tipo
finito sobre~$\bar{k}$ es fuertemente desingularizable.
\end{enumerate}

Entonces, si $p'$ es el conjunto de los números primos distintos de~$p$,
el morfismo
\begin{equation*}
\label{eq:XIII.4.6.0} \tag{\thesubsection.0} {}
\pi_1^{p'}(X\times_k Y,c)\to\pi_1^{p'}(X,a)\times\pi_1^{p'}(Y,b),
\end{equation*}
deducido
\marginpar{422}
de los homomorfismos de grupos fundamentales asociados a las proyecciones
$$
X\times_k Y\to X \quad \text{y} \quad X\times_k Y\to Y,
$$
es un isomorfismo.
\end{proposition}

Puede suponerse que~$k$ es algebraicamente cerrado y que~$X$ es reducido
(SGA~4 VIII~1.1). Sean $Z=X\times_k Y$, $g\colon X\to k$ y $f\colon
Z\to Y$ los morfismos canónicos. El morfismo~$g$, y por tanto también~$f$,
es universalmente localmente $1$-asférico para~$p'$
según~
\ifthenelse{\boolean{orig}}
{\Ref{XIII.3.5}}
{\Ref{XIII.3.4}}.
Como~$X$ es conexo, $f$ es
$0$-acíclico (SGA~4 XV~1.16). Por otra parte, de~\Ref{XIII.3.2} resulta
que, para todo $p'$-grupo finito~$C$, $(C_X,g)$ es cohomológicamente propio
relativamente a~$k$ en dimensión $\leq 1$. De ello resulta que~$(C_Z,f)$
es cohomológicamente propio relativamente a~$Y$ en dimensión $\leq 1$
(\Ref{XIII.1.5}~c)) y que $f_*C_Z$ y $\R^1f_*C_Z$ son haces constantes.
Por consiguiente,
\ifthenelse{\boolean{orig}}
{$g$}
{$f$}
satisface todas las hipótesis de~\Ref{XIII.4.2}. Se tiene, pues, la sucesión
exacta
$$
\pi_1^{p'}(X_b,c)\to\pi_1^{p'}(Z,c)\to\pi_1^{p'}(Y,b)\to 1.
$$
Además, el morfismo compuesto
$$
\pi_1^{p'}(X_b,c)\to\pi_1^{p'}(Z,c)\to
\ifthenelse{\boolean{orig}}
{\pi_1^{p'}(Y,b)}
{\pi_1^{p'}(X,a)}
$$
es un isomorfismo y se tiene, por tanto, la sucesión exacta
$$
1\to\pi_1^{p'}(X,a)\to\pi_1^{p'}(Z,c)\to\pi_1^{p'}(Y,b)\to 1.
$$
Por otra parte, el morfismo~\eqref{eq:XIII.4.6.0} permite definir
un morfismo de esta sucesión exacta en la sucesión exacta
$$
1\to\pi_1^{p'}(X,a)\to
\pi_1^{p'}(X,a)\times\pi_1^{p'}(Y,b)
\ifthenelse{\boolean{orig}}
{\times}
{\to}
\pi_1^{p'}(Y,b)\to 1,
$$
y de ello resulta que el morfismo~\eqref{eq:XIII.4.6.0} es un
isomorfismo.


\subsection{}
\label{XIII.4.7} Sea
$$
\xymatrix@C=1.2cm{
X \ar[d]_-f&\ar[l] X_U \ar[d]_-{f_U}& \ar[l]X_{\bar{s}}\ar[d] \\
S&\ar[l] U &\ar[l]\bar{s}
}
$$
un
\marginpar{423}
diagrama cuyos cuadrados son cartesianos, donde $S$ es un esquema conexo
por arcos (SGA~4 IX~2.12), $U$ un abierto conexo de~$S$ y~$\bar{s}$ un
punto geométrico de~$U$. Sean~$a$ un punto geométrico de~$X_{\bar{s}}$
y~$\LL$ un conjunto de números primos. Sea~$g$ una sección de~$f$ y
supongamos que se satisfacen las condiciones siguientes:
\begin{enumerate}
\item[a)] El morfismo $f$ es $0$-acíclico, localmente $0$-acíclico,
y, para todo recubrimiento étale $S'$ de~$S$ y todo recubrimiento
étale $E$ de~$X\times_S S'$, cociente de un recubrimiento galoisiano
cuyo grupo es un~$\LL$-grupo, $(E,f_{(S')})$ es cohomológicamente propio
relativamente a~$S'$ en dimensión $\leq 0$.
\item[b)] El morfismo $f_U$ es localmente $1$-asférico para~$\LL$ y,
para todo haz constante finito de~$\LL$-grupos~$F$ sobre~$X_U$, $(F,f_U)$
es cohomológicamente propio en dimensión $\leq 1$ y las fibras de
$\R^1 f_{U*}F$ son finitas.
\end{enumerate}

Se deduce entonces de~\Ref{XIII.4.1} y~\Ref{XIII.4.3} el diagrama
conmutativo siguiente, cuyas filas son exactas:
\begin{equation*}
\label{eq:XIII.4.7.0} \tag{\thesubsection.0} {}
\begin{array}{c}
\xymatrix{
1 \ar[r]& \pi_1^{\LL}(X_{\bar{s}},a) \ar[r]\ar@{=}[d]& \pi'_1(X_U,a) \ar[r]\ar[d]& \pi_1(U,a)\ar[r]\ar[d]& 1 \\
&\pi_1^{\LL}(X_{\bar{s}},a) \ar[r]&\pi'_1(X,a) \ar[r]& \pi_1(S,a) \ar[r]& 1
}
\end{array}
\end{equation*}

Gracias a la sección~$g$, se tienen morfismos
\ifthenelse{\boolean{orig}}
{$\pi_1(U,a)\to\pi'_1(X_U,a)$ $\pi_1(S,a)\to\pi'_1(X,a)$;}
{$$\pi_1(U,a)\to\pi'_1(X_U,a),\quad \pi_1(S,a)\to\pi'_1(X,a)\,\;$$}
de ellos se deduce un morfismo de la suma amalgamada de~$\pi_1(S,a)$ y
$\pi'_1(X_U,a)$ sobre~$\pi_1(U,a)$ en~$\pi'_1(X,a)$:
\begin{equation*}
\label{eq:XIII.4.7.1} \tag{\thesubsection.1} {}
\varphi\colon\pi=\pi_1(S,a)\coprod_{\pi_1(U,a)}\pi'_1(X_U,a)\to\pi'_1(X,a).
\end{equation*}
Supongamos que se satisface la condición siguiente:
\begin{enumerate}
\item[c)] Si $T=S-U$, se tiene $\prof \et_T(S)\geq 2$ (SGA~2 XIV~1.1).
\end{enumerate}
Entonces, el funtor que a un recubrimiento étale de~$S$ le hace
corresponder su restricción a~$U$ es plenamente fiel (SGA~2 XVI~1.4).
Se sigue que el morfismo $\pi_1(U,a)\to\pi_1(S,a)$ es sobreyectivo
(\Ref{V.6.9}), y del
\marginpar{424}
diagrama~\eqref{eq:XIII.4.7.0} se deduce que lo mismo ocurre con el
morfismo $\pi'_1(X_U,a)\to\pi'_1(X,a)$; con mayor razón, $\varphi$ es un
epimorfismo. Sea
\begin{equation*}
\label{eq:XIII.4.7.2} \tag{\thesubsection.2} {}
K=\Ker(\pi_1(U,a)\to\pi_1(S,a)).
\end{equation*}
El grupo~$\pi$ de~\eqref{eq:XIII.4.7.1} se identifica con el cociente de
\ifthenelse{\boolean{orig}}
{$\pi'_1(X_U,a)$ por el subgrupo normal cerrado engendrado por la imagen~$L$
de~$K$ en~$\pi'_1(X_U,a)$.}
{$\pi'_1(X_U,a)$ por el subgrupo normal cerrado~$L$ engendrado por la imagen
de~$K$ en~$\pi'_1(X_U,a)$.}
Consideremos $\pi'_1(X_U,a)$ como producto
semidirecto de~$\pi_1(U,a)$ por~$\pi_1^\LL (X_{\bar{s}},a)$. El grupo~$K$
actúa entonces por automorfismos interiores sobre
$\pi_1^\LL (X_{\bar{s}},a)$, y el cociente
$\pi=\pi'_1(X_U,a)/L$ se identifica con el producto semidirecto de
$\pi_1(U,a)/K=\pi_1(S,a)$ por el grupo
$\pi_1^\LL(X_{\bar{s}},a)_K$ de coinvariantes de
$\pi_1^{\LL}(X_{\bar{s}},a)$ bajo la acción de~$K$. Finalmente, se tiene un
epimorfismo
\begin{equation*}
\label{eq:XIII.4.7.3} \tag{\thesubsection.3} {}
\varphi\colon\pi=\pi_1^\LL(X_{\bar{s}},a)_K\cdot\pi_1(S,a)\to\pi'_1(X,a).
\end{equation*}
\label{indnot:ml}\oldindexnot{$\pi_1^\LL(X_{\bar{s}},a)_K$|hyperpage}%

La proposición siguiente da condiciones bajo las cuales el morfismo
$\varphi$ es un isomorfismo.


\setcounter{subsubsection}{3}

\begin{subproposition}
\label{XIII.4.7.4} Las notaciones son las de~\Ref{XIII.4.7}. Se supone
que, además de las condiciones \textup{a), b), c)}, se satisfacen las
condiciones siguientes:
\begin{enumerate}
\item[d)] Para todo punto $t$ de~$T=S-U$, el morfismo $f$ es localmente
$1$-asférico para~$\LL$ en $g(t)$.
\item[e)] Para todo punto $t\in T$, toda componente irreducible de la
fibra $X_t$ contiene a~$g(t)$ y, para todo punto $x$ de
$X_t-\{g(t)\}$ que no sea maximal, se tiene
$$
\prof\hop_x(X)\geq 3 \quad\textup{(SGA~2 XIV~1.2)}
$$
y el anillo $\cal{O}_{X,x}$ es noetheriano.
\end{enumerate}

Entonces el morfismo~\eqref{eq:XIII.4.7.3} es un isomorfismo.
\end{subproposition}

Como se ha dicho anteriormente, el grupo~$\pi$ se identifica con el
cociente de~$\pi'_1(X_U,a)$ por el subgrupo normal cerrado~$L$ engendrado
por la imagen de~$K$ \eqref{eq:XIII.4.7.2} en~$\pi'_1(X_U,a)$. Esto
equivale a decir que~$\pi$ clasifica los
\marginpar{425}
recubrimientos principales~$Z$ de~$X_U$ tales que $g^{-1}_U(Z)$ se
prolonga a un recubrimiento étale de~$S$, y que inducen sobre
$X_{\bar{s}}$ un recubrimiento obtenido por extensión del grupo
estructural a partir de un recubrimiento principal cuyo grupo es un
$\LL$-grupo. Para demostrar que~$\varphi$ es un isomorfismo, basta
demostrar que tal recubrimiento~$Z$ se prolonga a todo~$X$.

Demostremos primero que~$Z$ se prolonga a un recubrimiento étale de un
abierto que contiene a~$X_U$ y~$g(S)$. Sea~$W$ un esquema étale sobre~$X$
cuya imagen contiene a~$X_U$ y~$g(S)$, y pongamos $W_U=W\times_S U$. Del
hecho de que el morfismo $W\to S$ es $0$-acíclico y de la relación
$\mathrm{prof}.\,\et_T S\geq 2$ resulta que se tiene
$$
\mathrm{prof}.\,\et_{(W-W_U)} W\geq 2
$$
(SGA~2 XIV~1.13); por consiguiente, si $Z|W_U$ se prolonga a un
recubrimiento étale de~$W$, dicha prolongación es única salvo isomorfismo
único. De ello resulta que el problema de prolongar~$Z$ a un entorno de
$g(S)\cup X_U$ es local para la topología étale en un entorno de los
puntos de~$g(T)$. Si~$t$ es un punto de~$T$, pongamos $x=g(t)$ y denotemos
por~$\bar{X}$ (\resp por~$\bar{S}$) la localización estricta de~$X$ en
$\bar{x}$ (\resp de~$S$ en~$\bar{t}$), por
$\bar{U}=U\times_S\bar{S}$, por $\bar{X}_U=\bar{X}\times_X X_U$ y por
$\bar{g}\colon\bar{S}\to\bar{X}$ el morfismo deducido de~$g$. Basta
demostrar que, para todo punto~$t$ de~$T$, la imagen inversa~$\bar{Z}$
de~$Z$ sobre~$\bar{X}_U$ se prolonga a~$\bar{X}$ o, lo que equivale a lo
mismo, es trivial. Ahora bien, por definición de~$Z$, la imagen inversa
de~$\bar{Z}$ sobre~$\bar{U}$ es trivial. Para demostrar que~$\bar{Z}$ es
trivial, basta demostrar que es de la forma $\bar{f}^*_U E$, donde~$E$ es
un recubrimiento principal de~$\bar{U}$; en efecto, se tendrá
$E\simeq\bar{g}^*_U\bar{f}^*_U E\simeq\bar{g}^*_U\bar{Z}$, de donde se
sigue el resultado, puesto que $\bar{g}^*_U\bar{Z}$ es trivial. Ahora
bien, como el morfismo~$\bar{f}_U$ es $0$-acíclico y localmente
$0$-acíclico, para demostrar que~$\bar{Z}$ proviene de~$\bar{U}$ basta
demostrar que, para todo punto geométrico algebraico sobre un
punto de~$\bar{U}$, que puede suponerse que es el punto~$\bar{s}$,
\ifthenelse{\boolean{orig}}
{$\bar{Z}|X_{\bar{s}}$}
{$\bar{Z}|\bar{X}_{\bar{s}}$}
es trivial (SGA~4 XV~1.15). Pero, como
\ifthenelse{\boolean{orig}}
{$\bar{Z}|X_{\bar{s}}$}
{$\bar{Z}|\bar{X}_{\bar{s}}$}
se obtiene por extensión del grupo estructural a
partir de un recubrimiento principal cuyo grupo es un~$\LL$-grupo, esto
resulta del hecho de que el morfismo~$\bar{f}$ es $1$-asférico para~$\LL$.

Hemos
\marginpar{426}
demostrado, por tanto, que existe un entorno abierto~$V$ de
$g(S)\cup X_U$ tal que~$Z$ se prolonga a un recubrimiento étale~$Z_V$
de~$V$. Demostremos que~$Z_V$ se prolonga a todo~$X$. Basta ver que, para
todo punto~$x$ de~$X-V$, se tiene
$$
\prof\hop_x X\geq 3.
$$
Ahora bien, esto resulta de la hipótesis~e) y del hecho de que un punto
$x$ de~$X-V$ no puede ser maximal en su fibra~$X_t$, pues, como toda
componente irreducible de~$X_t$ contiene a~$g(t)$, todo punto maximal
de~$X_t$ pertenece a~$V$.


\setcounter{proposition}{7}% adaptación del contador autónomo: el rótulo fuente siguiente es XIII.4.8
\begin{corollaire} \label{XIII.4.8}
Las hipótesis son las de~\Ref{XIII.4.7.4}, pero se supone además que
$\pi_1(S,a)=1$. Entonces se tiene un isomorfismo
\begin{equation*}
\label{eq:XIII.4.8.*} \tag{$*$}
{}\pi_1^\LL(X,a)\isomto\pi_1^\LL(X_{\bar{s}},a)_K.
\end{equation*}
En particular, si $\pi_1^\LL(X_{\bar{s}},a)$ es topológicamente de
presentación finita y si $K$ actúa sobre~$\pi_1^\LL(X_{\bar{s}},a)$ a
través de un grupo finitamente generado, entonces $\pi_1^\LL(X,a)$ es
topológicamente de presentación finita.
\end{corollaire}

El isomorfismo~\eqref{eq:XIII.4.8.*} se demostró
en~\Ref{XIII.4.7.4}. Supongamos que $\pi_1^\LL(X_{\bar{s}},a)$ es el
cociente del grupo pro-$\LL$ libre con $n$ generadores
$L(x_1,\dots,x_n)$ por el subgrupo normal cerrado engendrado por los
elementos $y_1,\dots,y_p$ de~$L(x_1,\dots,x_n)$, y supongamos que $K$
actúa a través de un grupo engendrado por los elementos
$k_1,\dots,k_q$. Si, para cada $i\in[1,n]$ y $j\in[1,q]$, se denota
por~$z_{ij}$ un elemento de~$L(x_1,\dots,x_n)$ que levanta el elemento
$(k_j\cdot x_i)x_i^{-1}$, entonces $\pi_1^\LL(X_{\bar{s}},a)_K$ es el
cociente de~$L(x_1,\dots,x_n)$ por el subgrupo normal cerrado engendrado
por los elementos $(y_i)_{i\in[1,p]}$ y
$(z_{ij})_{i\in[1,n],j\in[1,q]}$.

\ifthenelse{\boolean{orig}}
{\setcounter{proposition}{5}}
{}

\begin{remarques}
\label{rem:XIII.4.6}
a) Las condiciones a) a e) de~\Ref{XIII.4.7} se satisfacen cuando
$S$ es un esquema normal conexo, $U$ un abierto denso de~$S$ con
inmersión cuasicompacta, y $f$ un morfismo propio de
presentación finita cuyas fibras son geométricamente conexas e
irreducibles en todo punto~$t$ de~$T$, siendo además~$f$ separable,
\marginpar{427}
liso en los puntos de~$X_U\cup g(T)$, siendo~$\LL$ el conjunto de los
números primos distintos de las características residuales de~$S$, y
siendo~$X$ regular en todo punto de~$X_t$. La condición a) se deduce,
en efecto, de SGA~4 XV~4.1 y 1.4; las condiciones b) y d) se deducen
de~\Ref{XIII.1.4} y SGA~4 XV~2.1 y XVI~5.2. Finalmente, e) se deduce de
\ifthenelse{\boolean{orig}}
{SGA~XIV~1.11}
{SGA~2 XIV~1.11}.
\par\smallskip
b) El corolario~
\ifthenelse{\boolean{orig}}
{\Ref{XIII.4.5}}
{\Ref{XIII.4.8}}
se aplica para calcular el grupo fundamental $\pi_1^{p'}(X)$ de una
superficie~$X$ propia y lisa sobre
un cuerpo separablemente cerrado~$k$ de característica~$p$
($p'$~designa el conjunto de los números primos distintos de~$p$).
El método nos fue comunicado por J.P\ptbl \textsc{Murre}; consiste en
reducirse, haciendo una explosión de~$X$, al caso en que se tiene una
fibración $X\to\PP^1_k$ y un abierto~$U$ de~$\PP^1_k$ que satisface
las hipótesis de~\Ref{XIII.4.7} (véase SGA~7 para más detalles). El
mismo método puede utilizarse más generalmente (\loccit) para demostrar
que, si~$X$ es un $k$-esquema conexo de tipo finito y si los esquemas
de tipo finito de dimensión $\leq\dim X$ sobre una clausura algebraica
de~$k$ son fuertemente desingularizables (SGA~5 I~3.1.5), entonces
$\pi_1^{p'}(X)$ es topológicamente de presentación finita.
\end{remarques}


\section{Apéndice I: Variaciones sobre el lema de Abhyankar}
\label{XIII.5}
\index{lema de Abhyankar|hyperpage}%
Este apéndice contiene distintas variantes del lema de Abhyankar.

\begin{proposition}
\label{XIII.5.1} Sean $X=\Spec A$ un esquema local regular,
$D=\sum_{1\leq i\leq r} \divisor f_i$ un divisor de cruzamientos
normales, donde los $f_i$ son elementos del ideal maximal de~$A$ que
forman parte de un sistema regular de parámetros. Sean $n_i$,
$1\leq i\leq r$, enteros
\ifthenelse{\boolean{orig}}{$\geq 0$}{$>0$}, y pongamos
$$
X'=X[T_1,\dots,T_r]/(T_1^{n_1}-f_1,\dots,T_r^{n_r}-f_r),
$$
\ifthenelse{\boolean{orig}}
{$U'=U\times_X X'$.}
{Sea $U=X-\Supp D$ y pongamos $U'=U\times_X X'$.}
Entonces $X'$ es regular y $U'$ es el complemento
en~$X'$ del divisor de cruzamientos normales
$\sum_{1\leq i\leq r}\divisor T_i$. Si los enteros $n_i$ son primos
\marginpar{428}
con la característica residual $p$ de~$X$, $U'$ es un recubrimiento
étale conexo de~$U$, moderadamente ramificado relativamente
a~$D$ \textup{(\Ref{XIII.2.3}~c))}.
\end{proposition}

En efecto, $X'$ es el espectro de un anillo local $A'$ cuyo ideal
maximal está engendrado por $T_1,\dots,T_r$.
\ifthenelse{\boolean{orig}}
{\ignorespaces}
{(Puede suponerse que $r=\dim(A)$.)}
Como $A'$ es finito y plano sobre~$A$, y por tanto de dimensión $r$,
$A'$ es regular (EGA~$0_{\textup{IV}}$~17.1.1) y los $T_i$ forman un
sistema regular de parámetros de~$A'$. Supongamos que los $n_i$ son
primos con~$p$. Como todos los $f_i$ son invertibles sobre~$U$, el
hecho de que $U'$ sea étale sobre~$U$ se deduce de
\ifthenelse{\boolean{orig}}
{I~\Ref{I.7.4}}
{\Ref{I.7.4}}.
Además, $U'$ está moderadamente ramificado relativamente a~$D$; sea,
en efecto, $x_i$ el punto genérico de~$V(f_i)$, $\bar{R}$ la
localización estricta de~$R=\cal{O}_{X,x_i}$ y $\bar{K}$ el cuerpo de
fracciones de~$\bar{R}$. Entonces la $\bar{K}$-álgebra que
representa $U'|\bar{K}$ se obtiene del cuerpo
$\bar{K}[T_i]/(T_i^{n_i}-f_i)$ mediante una extensión no ramificada;
por tanto, está moderadamente ramificada sobre~$\bar{R}$.


\begin{proposition}[Lema de Abhyankar absoluto]
\label{XIII.5.2}
\index{Abhyankar absoluto (lema de)|hyperpage}%
Sea $X$ un esquema local regular,
$$
D=\sum_{1\leq i\leq r} \divisor f_i
$$
un divisor de cruzamientos normales como en~\Ref{XIII.5.1},
$Y=\Supp D$, $U=X-Y$. Sea~$V$ un recubrimiento étale de~$U$
moderadamente ramificado relativamente a~$D$. Si $x_i$ es el punto
genérico del cerrado $V(f_i)$, $\cal{O}_{X,x_i}$ es un anillo de
valoración discreta con cuerpo de fracciones $K_i$, y se tiene
$V|K_i=\Spec(\prod_{j\in J_i}L_j)$, donde los $L_j$ son extensiones
finitas separables de~$K_i$; denotemos por $n_j$ el orden del grupo de
inercia de una extensión galoisiana generada por $L_j$, y por $n_i$
el m.c.m. de los $n_j$ cuando $j$ recorre~$J_i$. Si se pone
$$
X'=X[T_1,\dots,T_r]/(T_1^{n_1}-f_1,\dots,T_r^{n_r}-f_r),
$$
$U'=U_{(X')}$, $V'=V_{(X')}$, etc., el recubrimiento étale $V'$
de~$U'$ se prolonga, de manera única salvo isomorfismo único, a un
recubrimiento étale de~$X'$, y los $n_i$ son primos con la
característica residual $p$ de~$X$.
\end{proposition}

La unicidad se deduce de que $X'$ es normal \eqref{XIII.5.1}; en
efecto,
\marginpar{429}
un recubrimiento étale de~$X'$ que prolonga $V'$ es isomorfo a la
normalización de~$X'$ en la fibra de~$V'$ en el punto genérico
de~$X'$ \eqref{I.10.2}. Si $\bar{x}'$ es un punto geométrico de~$Y'$,
se denota por $\bar{X}'$ la localización estricta de~$X'$ en
$\bar{x}'$, y se pone $\bar{V}'=V'_{(\bar{X}')}$, etc. Por descenso,
habida cuenta de la unicidad, basta demostrar que, para todo punto
geométrico $\bar{x}'$ de~$Y'$, el recubrimiento étale $\bar{V}'$
de~$\bar{U}'$ se prolonga a~$\bar{X}'$. Puesto que un recubrimiento
étale de un abierto del esquema regular~$X'$ que contiene todos los
puntos~$x'$ tales que $\dim\cal{O}_{X',x'}\leq 1$ se prolonga
\ifthenelse{\boolean{orig}}
{a todo~$X$}
{a todo~$X'$}
(SGA~2 XIV~1.11), basta incluso limitarse a los puntos
$\bar{x}'$ que se proyectan sobre un punto maximal de~$Y'$. Ahora
bien, en tal punto $\bar{x}'$, el hecho de que $\bar{V}'$ se prolongue
a un recubrimiento étale de~$\bar{X}'$ se deduce de~\Ref{X.3.6}.

Demostremos que los $n_i$ son primos con~$p$. En efecto, si no fuera
así, se tendría, por ejemplo, $p|n_1$. Tras reemplazar $X$ por
\ifthenelse{\boolean{orig}}
{$X[T_1,\dots,T_r]/(T_1^{n_1/p}-f_1,T_2^{n_2}-f_2,\dots,T_r^{n_r}-f_r)$,}
{$$X[T_1,\dots,T_r]/(T_1^{n_1/p}-f_1,T_2^{n_2}-f_2,\dots,T_r^{n_r}-f_r),$$}
se reduce al caso en que $X'=X[T_1]/(T_1^p-f_1)$. Basta demostrar que
$V$ se prolonga a un recubrimiento étale de~$X$, pues entonces se
tendrá $n_1=1$, en contra de la hipótesis. Para ello puede suponerse
que $X$ es estrictamente local. Sea $Z$ el subesquema cerrado de~$X$
de ecuación $p=0$, y sea $Z_1=Z\cap X_{f_1}$; $Z_1$ es un abierto no
vacío de~$Z$. Por lo que precede, el recubrimiento étale $V'$ de~$U'$
se prolonga a un recubrimiento étale $W'$ de~$X'$. Sean $W''_1$ y
$W''_2$ las imágenes inversas de~$W'$ por las dos proyecciones
$X''=X'\times_X X'\rightrightarrows X'$, y demostremos que el
isomorfismo de descenso $u\colon W''_1|U''\to W''_2|U''$ se prolonga
a un $X''$-morfismo $W''_1\to W''_2$, que será necesariamente un dato
de descenso sobre~$W'$ relativamente a~$X'\to X$. Sea $Z''$
(\resp $Z''_1$) la imagen inversa de~$Z$ (\resp $Z_1$) en~$X''$. Como
el morfismo $Z''\to Z$ es radicial, existe un isomorfismo
$v\colon W''_1|Z''\to W''_2|Z''$ que prolonga el isomorfismo
$u|Z''_1$. Pero, como $X$ es henseliano, se tiene una biyección
$$
\Hom_{X''}(W''_1,W''_2)\simeq\Hom_{Z''}(W''_1|Z'',W_2''|Z''),
$$
de donde se obtiene un morfismo $w\colon W''_1\to W''_2$ que levanta
a~$v$. El subesquema a la vez abierto
\marginpar{430}
y cerrado de~$X''$ sobre el cual coinciden $u$ y $w$ contiene
$Z''_1$; por tanto, es igual a~$X''$, y de aquí se sigue que $V$ se
prolonga a~$X$.


\setcounter{subsection}{3} \setcounter{subsubsection}{-1}

\subsubsection{}
\label{XIII.5.3.0} Retomemos las hipótesis y las notaciones
de~\Ref{XIII.5.2}, suponiendo además que
\ifthenelse{\boolean{orig}}
{$S$}
{$X$}
es estrictamente local. Entonces se deduce de \loccit que todo
recubrimiento étale conexo de~$U$, moderadamente ramificado
relativamente a~$D$, es cociente de un recubrimiento moderadamente
ramificado de la forma
$$
U'=U[T_1,\dots,T_r]/(T_1^{n_1}-f_1,\dots,T_r^{n_r}-f_r),
$$
donde los $n_i$ son enteros primos con~$p$. Sea $\bbmu_n$ el grupo de
las raíces $n$-ésimas de la unidad de~$U$. El grupo de los
$U$-automorfismos de~$U'$ no es otro que el grupo
$\bbmu_{n_1}\times\dots\times\bbmu_{n_r}$: una raíz $n_i$-ésima de
la unidad~$\xi_i$ actúa sobre~$U'$ transformando $T_i$ en
$\xi_i T_i$. Se obtiene, pues, el enunciado siguiente:

\setcounter{subsection}{2}

\begin{corollaire}
\label{XIII.5.3} Sea $X$ un esquema estrictamente local regular de
característica residual $p\geq 0$,
$D=\sum_{1\leq i\leq
\ifthenelse{\boolean{orig}}{n}{r}
} \divisor f_i$ un divisor de cruzamientos normales sobre~$X$,
$U=X-\Supp D$. Pongamos
$$
\tilde U=\varprojlim_{(n_i)}
U[T_1,\dots,T_r]/(T_1^{n_1}-f_1,\dots,T_r^{n_r}-f_r),
$$
donde el límite proyectivo se toma sobre el conjunto ordenado filtrante
(para la relación de divisibilidad) de las familias de enteros $n_i>0$
primos con~$p$. Entonces $\tilde U$ es un recubrimiento universal
moderadamente ramificado de~$U$. Por consiguiente, el grupo fundamental
moderadamente ramificado de~$U$ es
$$
\pi_1^{\tame}(U)\simeq \prod_{\ell\neq p} \ZZ_\ell [1]^r
\qquad \text{(isomorfismo canónico)}
$$
donde se ha puesto
$\ZZ_\ell[1]=\varprojlim_{n>0}\bbmu_{\ell^n}$. El grupo
$\pi_1^{\tame}(U)$ es isomorfo no canónicamente a
$\prod_{\ell\neq p}\ZZ_\ell^r$.
\label{indnot:mm}\oldindexnot{$\ZZ_\ell[1]$|hyperpage}%
\end{corollaire}


\begin{proposition}
\label{XIII.5.4} Sean $f\colon X\to S$ un morfismo de esquemas y
$D=\sum_{1\leq i\leq r} \divisor f_i$ un divisor de cruzamientos
normales relativamente a~$S$ \eqref{XIII.2.1}, donde,
\marginpar{431}
para cada punto $x$ de~$Y=\Supp D$, si $I(x)\subset[1,r]$ es el conjunto de los $i$
tales que $f_i(x)=0$, el subesquema $V((f_i)_{i\in I(x)})$ es liso
sobre~$S$ de codimensión
\ifthenelse{\boolean{orig}}
{$\mathrm{card}.\,I(x)$}
{$\card I(x)$}
en~$X$. Sea $U=X-Y$. Sean $x$ un punto de~$Y$,
$X_1=\Spec\cal{O}_{X,x}$, $U_1=U\times_X X_1$, $n_i$, $i\in I(x)$,
\ifthenelse{\boolean{orig}}
{enteros}
{enteros $>0$}, y pongamos
$$
X'=X[T_i]_{i\in I(x)}/(T_i^{n_i}-f_i).
$$
Entonces, si $x'$ es el punto de~$X'$ situado sobre~$x$, $X'$ es liso
sobre~$S$ en~$x'$. Si los enteros~$n_i$ son primos con la
característica~$p$ de~$\kres(x)$, $U'_1=U_1\times_X X'$ es un
recubrimiento étale conexo de~$U_1$, moderadamente ramificado sobre~$X_1$
relativamente a
\ifthenelse{\boolean{orig}}
{$S_1$}
{$S$}
\eqref{XIII.2.1.1}.
\end{proposition}

Si $s=f(x)$, la fibra geométrica $X'_{\bar{s}}$ es regular en el punto
$x'$ \eqref{XIII.5.1}; como $X'$ es plano sobre~$S$ en un entorno
\ifthenelse{\boolean{orig}}
{de~$Y$}
{de la imagen inversa de~$Y$ en~$X'$},
esto demuestra que $X'$ es liso sobre~$S$ en~$x'$
(EGA~IV~12.1.6). Si los enteros $n_i$ son primos con~$p$, $U'_1$ es un
recubrimiento étale de~$U_1$ \eqref{I.7.4}; está moderadamente
ramificado sobre~$X_1$ relativamente a~$S$ porque así ocurre en las
fibras geométricas sobre cada punto de~$S$ \eqref{XIII.5.1}. Por último,
el hecho de que $U'_1$ sea conexo se deduce de (SGA~4 XVI~3.2).


\begin{proposition}[Lema de Abhyankar relativo]
\label{XIII.5.5}
\index{Abhyankar relativo (lema de)|hyperpage}%
Sean $X$ un $S$-esquema y $D$ un divisor de cruzamientos normales
relativamente a~$S$, como en~\eqref{XIII.5.4}. Sean
\ifthenelse{\boolean{orig}}
{$Y=X-\Supp D$}
{$Y=\Supp D$},
$U=X-Y$, $x$ un punto de~$Y$, $X_{1}$ la localización estricta de~$X$
en un punto geométrico situado sobre~$x$, $U_{1}=U\times_{X}X_{1}$ y
$V_{1}$ un recubrimiento étale de~$U_{1}$. Se supone que, para todo
punto maximal $s$ de~$S$, $V_{1\overline{s}}$ está moderadamente
ramificado sobre $X_{1\overline{s}}$ relativamente a~$\overline{s}$.
Entonces pueden encontrarse enteros $n_{i}$, primos con la
característica~$p$ de~$\kres(x)$, con
$i\in I(x)$, tales que, si se pone
$$
X'_{1}=X_{1}[T_{i}]_{i\in I(x)}\big/(T_{i}^{n_{i}}-f_{i})\quoi,
$$
$U'_{1}=U_{1}\times_{X_{1}}X'_{1}$, etc., el recubrimiento étale
$V'_{1}$ de~$U'_{1}$ se prolonga, de manera única salvo isomorfismo
único, a un recubrimiento étale de~$X'_{1}$. En particular,
\marginpar{432}
$V_{1}$ está moderadamente ramificado sobre~$X_{1}$ relativamente a~$S$.
\end{proposition}

Puede suponerse que $S$ es local noetheriano, con punto cerrado~$f(x)$.
Para cada punto maximal $s$ de~$S$ y cada $i\in I(x)$, sea $x_{i}$ el
punto genérico del cerrado $V(f_{i})$ de la fibra $X_{1\overline{s}}$.
El anillo local $(\cal{O}_{X_{1},x_{i}})_{\red}$ es un anillo de
valoración discreta con cuerpo de fracciones $K_{i}$, y se tiene
\ifthenelse{\boolean{orig}}
{$V_{1}|K_{i}=\Spec(\prod_{j\in J(x_{i})}L_{j})$,}
{$V_{1}|K_{i}=\Spec(\prod_{j\in J(x_{i})}L_{j})$,}
donde $L_{j}$ es una extensión finita separable de
\ifthenelse{\boolean{orig}}
{$K$}
{$K_{i}$};
se denota por $n_{j}$ el orden del grupo de inercia de una extensión
galoisiana generada por~$L_{j}$, y por $n_{i}$ el m.c.m. de los $n_{j}$
cuando $s$ recorre los puntos maximales de~$S$ y
\ifthenelse{\boolean{orig}}
{$j\in J(x_{i})$.}
{$j\in J(x_{i})$.}

Elegidos así los $n_{i}$, demostraremos que $V'_{1}$ se prolonga de
manera única a un recubrimiento étale de~$X'_{1}$. La unicidad se
deduce de que, como
\ifthenelse{\boolean{orig}}
{$X'$ es liso sobre~$S$ en los puntos de~$Y$}
{$X'$ es liso sobre~$S$ en los puntos de~$Y'$},
se tiene $\prof\et_{Y'_{1}}(X'_{1})\geqslant 2$ (SGA~4 XVI~3.2 o
SGA~2 XIV~1.19). Sean $x'_{1}$ un punto de~$Y'_{1}$,
$\overline{x}'_{1}$ un punto geométrico situado sobre~$x'_{1}$, y
denotemos por $\overline{X'_{1}}$ la localización estricta de~$X'_{1}$
en~$\overline{x}'_{1}$, $\overline{U'_{1}}=U'_{1(\overline{X'_{1}})}$,
etc. Por descenso, habida cuenta de la unicidad, basta demostrar que
$\overline{V}'_{1}$ se prolonga a~$\overline{X'_{1}}$. Además, puede
uno limitarse a tomar como $x'_{1}$ los puntos maximales de~$Y'_{1}$;
en efecto, se tendrá entonces una prolongación de~$V'_{1}$ sobre un
abierto $W'_{1}$ de~$X'_{1}$ que contiene los puntos maximales
de~$Y'_{1}$. Ahora bien, si $Z'_{1}=X'_{1}-W'_{1}$, se tiene
$\codim(Z'_{1s},X'_{1s})\geqslant 2$ cuando $s$ es un punto maximal
de~$S$, y $\codim(Z'_{1s},X'_{1s})\geqslant 1$,
$\prof\et_{s}(S)\geqslant 1$ cuando $s$ es un punto de~$S$ que no es
maximal; el hecho de que $V'_{1}$ se prolongue a todo~$X'_{1}$ se deduce
entonces de (SGA~2 XIV~1.20). Pero, en un punto geométrico
$\overline{x}'_{1}$ situado sobre un punto maximal de~$Y'_{1}$,
\ifthenelse{\boolean{orig}}
{$X'_{1\red}$}
{$\overline{X'_{1}}_{\red}$}
es el espectro de un anillo de valoración discreta, y el hecho de que
$\overline{V'_{1}}$ se prolongue a~$\overline{X'_{1}}$ se deduce entonces
de~\Ref{X.3.6}.

Demostremos que los $n_{i}$ son primos con~$p$. En efecto, si no fuera
así, habría un índice $i_{0}\in I(x)$ tal que $p$ divide a~$n_{i_{0}}$.
Tras reemplazar, si es preciso, $X$ por
\ifthenelse{\boolean{orig}}
{$X_{1}[T_{i_{0}},T_{i}]_{i\in I(x)}\big/
(T_{i}^{n_{i_{0}}/p}-f_{i_{0}},T_{i}^{n_{i}}-f_{i})$,}
{$X_{1}[T_{i}]_{i\in I(x)}\big/
(T_{i_{0}}^{n_{i_{0}}/p}-f_{i_{0}},
(T_{i}^{n_{i}}-f_{i})_{i\in I(x),\,i\ne i_{0}})$,}
se reduce al caso en que
\ifthenelse{\boolean{orig}}
{$X'_{1}=X_{1}[T]/T^{p}-f_{i_{0}}$}
{$X'_{1}=X_{1}[T]/(T^{p}-f_{i_{0}})$}.
Por lo que
precede, el recubrimiento
\marginpar{433}
étale $V'_{1}$ de~$U'_{1}$ se prolonga a un recubrimiento étale $E'_{1}$
de~$X'_{1}$. Sea $\eta$ el punto cerrado de~$S$; como el morfismo
$X'_{1\eta}\to X_{1\eta}$ es radicial, $V_{1\eta}$ se prolonga a un
recubrimiento étale $E_{1\eta}$ de~$X_{1\eta}$. Se deduce entonces,
como en~\Ref{XIII.5.2}, que $E'_{1}$ está provisto de un dato de descenso
relativamente al morfismo $X'_{1}\to X_{1}$, que prolonga el dato de
descenso natural que se tiene sobre~$E'_{1}|U'_{1}$. De aquí resulta que
$V_{1}$ se prolonga a~$X_{1}$; pero esto implica $n_{i_{0}}=1$, en
contra de la hipótesis $n_{i_{0}}=p$.


\begin{corollaire}
\label{XIII.5.6}
Sean $X$ un $S$-esquema y
$D=\sum_{1\leqslant i\leqslant r}\divisor f_i$ un divisor de
cruzamientos normales relativamente a~$S$, como en~\Ref{XIII.5.4}.
Sean $\overline{x}$ un punto geométrico de~$X$, $\overline{X}$ la
localización estricta de~$X$ en~$\overline{x}$,
$\overline{Y}=Y_{(\overline{X})}$,
$\overline{U}=\overline{X}-\overline{Y}$ y
$$
\widetilde{U}=\varprojlim_{(n_i)}
\overline{U}[T_i]_{i\in I(x)}\big/(T_i^{n_i}-f_i)\quoi,
$$
donde el límite proyectivo se toma sobre el conjunto dirigido de las
familias de enteros $n_i>0$, primos con la característica~$p$
de~$\kres(x)$. Entonces $\widetilde{U}$ es un recubrimiento universal
moderadamente ramificado de~$\overline{U}$ relativamente a~$S$. Por
consiguiente, el grupo fundamental moderadamente ramificado
de~$\overline{U}$ es
$$
\pi_1^{\tame}(\overline{U})\simeq
\prod_{\ell\neq p}\ZZ_\ell[1]^{I(x)}
\qquad \textrm{(isomorfismo canónico)\quoi.}
$$
El grupo $\pi_1^{\tame}(\overline{U})$ es isomorfo, de manera no
canónica, a~$\prod_{\ell\neq p}\ZZ_\ell^{I(x)}$.
\end{corollaire}

\setcounter{subsection}{6}
\setcounter{subsubsection}{0}
\begin{subremarque}
\label{XIII.5.6.1}
Sean $X$ un $S$-esquema y
$D=\sum_{1\leqslant i\leqslant r}\divisor f_i$ un divisor de
cruzamientos normales relativamente a~$S$, como en~\Ref{XIII.5.4}, y
$U=X-\Supp D$. Para todo subconjunto $I\subset[1,r]$, sea
$$
X_I=\Big(\tbigcap_{i\in I}V(f_i)\Big)\cap
\Big(\tbigcap_{i\in\complement I}X_{f_i}\Big)\quoi.
$$
Sea $p$ un número primo o cero, y sea $Z$ un subconjunto de~$X_I$
todos cuyos puntos son de característica~$p$. Sea
$$
\widetilde{U}_I=\varprojlim_{(n_i)}
\ifthenelse{\boolean{orig}}
{U[T_i]}
{U[T_i]_{i\in I}}
\big/(T_i^{n_i}-f_i)\quoi,
$$
donde
\marginpar{434}
el límite proyectivo se toma sobre el conjunto dirigido de las familias
de enteros $n_i>0$, primos con~$p$. Entonces, para todo punto geométrico
$\overline{x}$ de~$Z$, la imagen inversa de~$\widetilde{U}_I$ sobre
$\overline{U}$ se identifica con el recubrimiento universal moderadamente
ramificado de~$\overline{U}$.
\end{subremarque}


\begin{corollaire}
\label{XIII.5.7}
Las notaciones son las de~\Ref{XIII.5.6}. Sean $\overline{S}$ la
localización estricta de~$S$ en~$\overline{x}$,
$$
\overline{g}\colon\overline{U}\to\overline{S}
\quad \text{y} \quad
\tilde{g}\colon\widetilde{U}\to
\ifthenelse{\boolean{orig}}
{\widetilde{S}}
{\overline{S}}
$$
los morfismos canónicos. Entonces los morfismos $\overline{g}$ y
$\tilde g$ son $0$-acíclicos \textup{(SGA~4 XV~1.3)}. Sean $G$ un haz
de grupos constructible sobre~$\overline{S}$,
$F=\overline{g}^*G$ y $P$ un torsor bajo~$F$. Para que $P$ esté
moderadamente ramificado sobre~$\overline{X}$ relativamente
a~$\overline{S}$, es necesario y suficiente que su imagen inversa
$\widetilde{P}$ sobre~$\widetilde{U}$ sea trivial.
\end{corollaire}

En efecto, para todo esquema
$\overline{X'}=\overline{X}[T_i]_{i\in I(x)}\big/
(T_i^{n_i}-f_i)$, donde los $n_i$ son enteros $>0$ primos con~$p$, el
morfismo $\overline{f}'\colon\overline{X'}\to\overline{S}$ es
$0$-acíclico. Por tanto, las fibras geométricas de~$\overline{f}'$ en
los distintos puntos de~$\overline{S}$ son conexas e incluso
irreducibles. Lo mismo sucede con las fibras geométricas de los
morfismos $\overline{g}'\colon\overline{U'}\to\overline{S}$, lo que
demuestra que los $\overline{g}'$, y por tanto también~$\tilde g$, son
$0$-acíclicos (SGA~4 XV~1.16).

Es claro que un torsor $P$ sobre~$\overline{U}$, bajo~$F$, cuya imagen
inversa sobre~$\widetilde{U}$ es trivial, está moderadamente ramificado
sobre~$\overline{X}$ relativamente a~$\overline{S}$. Demostremos el
recíproco: si $P$ está moderadamente ramificado sobre~$\overline{X}$
relativamente a~$\overline{S}$, su imagen inversa sobre~$\widetilde{U}$
es trivial.

Se deduce de (SGA~4 IX~2.14~(ii)) que pueden encontrarse un morfismo
finito $n\colon S_1\to\overline{S}$, un haz de grupos constante~$C$
sobre~$S_1$ y un monomorfismo $G\to n_*C$. Consideremos el siguiente
diagrama conmutativo formado por cuadrados cartesianos:
$$
\xymatrix{ \widetilde{U}_{1} \ar[r] \ar[d]_{r} &U_{1} \ar[r]^{g_{1}}
\ar[d]_{q} &S_{1} \ar[d]_{n} \cr \widetilde{U} \ar[r] &\overline{U}
\ar[r]^{\overline{g}} &\overline{S} \rlap{\quoi.}\cr }
$$
Sean
\marginpar{435}
$C_1$ (\resp $\widetilde{C}_1$) la imagen inversa de~$C$ sobre~$U_1$
(\resp sobre~$\widetilde{U}_1$). Se tiene un diagrama conmutativo en el
que $i$ y $j$ son isomorfismos (SGA~4 VIII~5.8):
\begin{equation*}
\label{eq:XIII.5.7.*}
\tag{$*$}
\begin{array}{c}
\xymatrix{ \H^{1}(\overline{U},q_*C_{1}) \ar[r]^{i}
\ar[d] &\H^{1}(U_{1},C_{1}) \ar[d] \cr
\H^{1}(\widetilde{U},r_*\widetilde{C}_{1}) \ar[r]^{j}
&\H^{1}(\widetilde{U}_{1},\widetilde{C}_{1}) \rlap{\quoi.}\cr }
\end{array}
\end{equation*}
Sea $Q$ el torsor bajo~$q_*C_1$ deducido de~$P$ por extensión del grupo
estructural $F\to q_*C_1$. Según~\Ref{XIII.2.1.4}, $Q$ está
moderadamente ramificado sobre~$\overline{X}$ relativamente
a~$\overline{S}$. Al torsor~$Q$ corresponde, mediante~$i$, un torsor
$Q_1$ bajo~$C_1$, y es claro que $Q_1$ está moderadamente ramificado
sobre $X_1=\overline{X}\times_{\overline{S}}S_1$ relativamente a~$S_1$.
Se deduce entonces de~\Ref{XIII.5.6} que la imagen inversa
$\widetilde{Q}_1$ de~$Q_1$ sobre~$\widetilde{U}_1$ es trivial, y el
diagrama~\eqref{eq:XIII.5.7.*} demuestra que la imagen inversa
$\widetilde{Q}$ de~$Q$ sobre~$\widetilde{U}$ es trivial.

Consideremos el siguiente diagrama conmutativo, cuya segunda fila es
exacta (SGA~4 XII~3.1):
\begin{equation*}
\label{eq:XIII.5.7.**}
\tag{$**$}
\begin{array}{c}
\xymatrix{ \H^{0}(\overline{S},n_*C/G) \ar[r]
\ar[d]_{k} &\H^{1}(\overline{S},G)\rlap{$~=~1$} \ar[d] \cr
\H^{0}(\widetilde{U},r_*\widetilde{C}_{1}/\widetilde{F}) \ar[r]
&\H^{1}(\widetilde{U},\widetilde{F}) \ar[r]
&\H^{1}(\widetilde{U},r_*\widetilde{C}_{1}) \rlap{\quoi.} \cr }
\end{array}
\end{equation*}
Como el morfismo $\widetilde{U}\to\overline{S}$ es $0$-acíclico,
$k$ es un isomorfismo. De~\eqref{eq:XIII.5.7.**} se deduce entonces
que $\widetilde{P}$ es trivial.


\section[Apéndice II: finitud para las imágenes directas de los stacks]
{Apéndice II: teorema de finitud para las imágenes directas de los stacks}
\label{XIII.6}

\begin{proposition}
\label{XIII.6.1}
Sean $S$ un esquema localmente noetheriano y $f\colon X\to S$ un
morfismo. Si $S'$ es un $S$-esquema, se denota por $X'$ (\resp por
$f'$, etc.) la imagen inversa de~$X$ (\resp de~$f$, etc.) por el
morfismo $S'\to S$. Supongamos que, para todo
\marginpar{436}
esquema $S'$ étale sobre~$S$ y todo haz de conjuntos constructible~$F$
sobre~$X'$, $f'_*F$ es constructible, y que, para todo haz de grupos
constructible~$F$ sobre~$X'$, $\R^1f'_*F$ es constructible. Sea $\Phi$
un stack $1$-constructible sobre~$X$~\eqref{XIII.0}. Entonces $f_*\Phi$
es $1$-constructible.
\end{proposition}

Para todo esquema $S'$ étale sobre~$S$ y todo objeto~$x$
de~$(f_*\Phi)_{S'}$, se tiene un isomorfismo
$$
\SheafAut_{S'}(x)\simeq f'_*\SheafAut_{X'}(x)\quoi,
$$
donde, en el segundo miembro de la igualdad, $x$ se considera como
objeto de~$\Phi_{X'}$. Por tanto, las hipótesis implican que $f_*\Phi$
es constructible. Sea $S\Phi$ el haz de los subgerbes maximales
de~$\Phi$~\cite[III~2.1.7]{XIII.1}. Como $f_*(S\Phi)$ es constructible,
puede aplicársele (SGA~4 IX~2.7), y el hecho de que $f_*\Phi$ sea
$1$-constructible se deduce entonces del lema siguiente.

\setcounter{subsection}{1}
\setcounter{subsubsection}{0}
\begin{sublemme}
\label{XIII.6.1.1}
Sean $S$ un esquema localmente noetheriano, $f\colon X\to S$ un
morfismo y $\Phi$ un stack sobre~$X$. Supongamos dados un haz sobre~$S$,
representable por un $S$-esquema~$T$ étale de tipo finito, un morfismo
sobreyectivo
$$
a\colon T\to f_*(S\Phi)
$$
y un objeto~$p$ de la fibra $\Phi_{X_T}$, donde
$X_T=X\times_ST$, que define en
$f_*(S\Phi)(T)=\allowbreak S\Phi(X_T)$ un elemento igual a la imagen~$q$
por~$a$ de la sección identidad de~$T(T)$. Sea
$f_T\colon X_T\to T$ el morfismo canónico, y supongamos que el haz
$\R^1f_{T*}(\SheafAut_{X_T}(p))$ es constructible; entonces también
$S(f_*\Phi)$ es constructible.
\end{sublemme}

El morfismo canónico $f^*f_*\Phi\to\Phi$ da un morfismo
$$
S(f^*f_*\Phi)\simeq f^*(S(f_*\Phi))\to S\Phi,
$$
y por tanto un morfismo canónico
$$
\varphi\colon S(f_*\Phi)\to f_*(S\Phi).
$$
Sean
\marginpar{437}
$F=S(f_*\Phi)$ y $G$ la imagen de~$F$ por~$\varphi$; según
(SGA~4 IX~2.9), $G$ es un haz constructible.

Basta demostrar que, para todo punto~$s$ de~$S$, existe un abierto no
vacío~$U$
\ifthenelse{\boolean{orig}}
{de~$s$}
{de~$\overline{\{s\}}$}
tal que $F|U$ sea localmente constante y constructible.
Sean $s\in S$, $\overline{s}$ un punto geométrico situado sobre~$s$, y
$\overline{q}_1,\dots,\overline{q}_n$ los elementos de~$G_{\overline{s}}$.
Por definición de~$T$, existen $S$-morfismos
\ifthenelse{\boolean{orig}}
{$h_i\colon\overline{s}\to S'$}
{$h_i\colon\overline{s}\to T$}
tales que $\overline{q}_i=h_i^*(q)$. Sean $S'$ el producto fibrado
sobre~$S$ de $n$ esquemas isomorfos a~$T$, $\overline{s}\to S'$ el
producto fibrado de los $h_i$, $X'=X\times_SS'$, y $q_i$ (\resp $p_i$)
la imagen inversa de~$q$ (\resp de~$p$) por la $i$-ésima proyección
de~$S'$ sobre~$T$. Si $\Psi_i$ es el subgerbe maximal de~$\Phi|X'$
engendrado por~$p_i$, el haz
$F_i=\R^1f_*'(\SheafAut_{X'}(p_i))$ no es otro que el haz
$S(f_*'\Psi_i)$ de los subgerbes maximales de~$f_*'\Psi_i$. En
particular, la inyección canónica $\Psi_i\to\Phi|X'$ da un morfismo
$$
\alpha_i\colon F_i\to F|S'.
$$

Vamos a demostrar que $\alpha_i$ es una biyección de~$F_i$ sobre la
imagen inversa de~$q_i$ en~$F|S'$. Para todo esquema $S''$ étale
sobre~$S'$, toda sección~$y$ de~$F_i(S'')$ tiene por imagen
$q_i|S''$ en~$F(S'')$, pues, localmente para la topología étale
sobre~$S''$, $y$ está definida por un objeto~$x$ de~$\Phi_{X''}$
isomorfo a~$p_i|X''$. Recíprocamente, si $y\in F(S'')$ tiene por imagen
$q_i|S''$ en~$F(S'')$, localmente para la topología étale sobre~$S''$,
$y$ está definida por un objeto~$x$ de~$\Phi_{X''}$ isomorfo a~$p_i$;
por consiguiente, $x$ es un objeto de~$\Psi_{iX''}$ y
$y\in F_i(S'')$.

La demostración concluye utilizando~\Ref{XIII.6.1.2}. En efecto, puede
encontrarse un entorno abierto~$U'$ de~$s$ tal que
$q_1|U',\dots,q_n|U'$ sean secciones de~$G(U')$ y engendren este haz.
Como $F_i|U'$ y $G|U'$ son constructibles, también $F|U'$ lo es
según~\Ref{XIII.6.1.2}; tras sustituir, si es preciso, $U$ por un
abierto más pequeño, $F|U$ es localmente constante, lo que concluye la
demostración.


\begin{sublemme}
\label{XIII.6.1.2}
Sean
\marginpar{438}
$S$ un esquema localmente noetheriano y $F\to G$ un morfismo
sobreyectivo de haces
\ifthenelse{\boolean{orig}}
{de grupos}
{de conjuntos}
sobre~$S$. Sean los $q_i$ una familia finita de secciones de~$G$
\ifthenelse{\boolean{orig}}
{sobre~$X$}
{sobre~$S$}
que engendran~$G$, y, para cada~$i$, sea $F_i$ el subhaz de~$F$ que es
la imagen inversa de~$q_i$. Entonces, si $G$ y los $F_i$ son
constructibles, también $F$ es constructible.
\end{sublemme}

Para demostrar que $F$ es constructible, basta probar que, para todo
punto~$s$ de~$S$, existe un entorno abierto~$U$ de~$s$ tal que $F|U$
sea localmente constante y constructible. Sea, pues, $s$ un punto de~$S$.
Como los haces $F_i$ y $G$ son constructibles, puede encontrarse un
entorno abierto~$U$ de~$s$ tal que $F_i|U$ y $G|U$ sean localmente
constantes. Demostremos entonces que $F|U$ es localmente constante.
Según (SGA~4 IX~2.13~(i)), basta ver que, si $\overline{s}$ es un punto
geométrico situado sobre~$s$, $\tilde{s}$ un punto geométrico de~$U$ y
$\overline{s}\to\tilde{s}$ un morfismo de especialización, el morfismo
canónico
$$
F_{\tilde{s}}\to F_{\overline{s}}
$$
es biyectivo.

Consideremos los diagramas conmutativos
$$
\xymatrix{ (F_i)_{\tilde{s}} \ar[r]^-{\tilde{a}_i} \ar[d]^{\wr} &
F_{\tilde{s}} \ar[r]^-{\tilde{a}} \ar[d]_b & G_{\tilde{s}}
\ar[d]^{\wr}\\ (F_i)_{\overline{s}} \ar[r]^-{\overline{a}_i} &
F_{\overline{s}} \ar[r]^-{\overline{a}} & G_{\overline{s}} }
$$
Sea $\overline{q}_i$ (\resp $\tilde{q}_i$) la imagen inversa de~$q_i$
en~$G_{\overline{s}}$ (\resp en~$G_{\tilde{s}}$); los morfismos
$\overline{a}$ y $\tilde{a}$ son sobreyectivos, y el morfismo
$\overline{a}_i$ (\resp $\tilde{a}_i$) induce una biyección de
$(F_i)_{\overline{s}}$ (\resp $(F_i)_{\tilde{s}}$) sobre
\ifthenelse{\boolean{orig}}
{$\overline{a}^{-1}(q_i)$ (\resp sobre~$\tilde{a}^{-1}(q_i)$)}
{$\overline{a}^{-1}(\overline{q}_i)$
(\resp sobre~$\tilde{a}^{-1}(\tilde{q}_i)$)}. Se deduce,
pues, del diagrama anterior que $b$ es un isomorfismo.


\begin{corollaire}
\label{XIII.6.2}
Sean
\marginpar{439}
$S$ un esquema localmente noetheriano y $f\colon X\to S$ un morfismo
propio. Sea $\Phi$ un stack $1$-constructible sobre~$X$; entonces
$f_*\Phi$ es un stack \hbox{$1$-constructible}.
\end{corollaire}

La demostración de~\Ref{XIII.6.1} prueba también el resultado siguiente,
teniendo en cuenta~\Ref{XIII.2.4}~2).

\begin{corollaire}
\label{XIII.6.3}
Sean $S$ un esquema localmente noetheriano, $f\colon X\to S$ un
morfismo, $D$ un divisor sobre~$X$ de cruzamientos normales
relativamente a~$S$~\eqref{XIII.2.1}, $Y=\Supp D$, $U=X-Y$ e
$i\colon U\to X$ la inmersión canónica. Sea $\Phi$ un stack sobre~$U$
dado, localmente para la topología étale sobre~$X$ y~$S$, como imagen
inversa de un stack~$\Psi$ $1$-constructible sobre~$S$. Entonces el
stack $i_*^{\tame}\Phi$ es $1$-constructible.
\end{corollaire}



\begin{sgabibliography}{5}{XIII.7}

\bibitem[1]{XIII.1} J\ptbl \textsc{Giraud}, Cohomologie non abélienne de
degré~$2$, tesis, París (1966).

\bibitem[2]{XIII.2} J\ptbl \textsc{Giraud},
\ifthenelse{\boolean{orig}}
{Cohomologie non abélienne de degré~$2$, en prensa.}
{Cohomologie non abélienne, Springer, Berlín--Nueva York (1971).}

\bibitem[3]{XIII.3}
\ifthenelse{\boolean{orig}}
{\textsc{Seifert}-\textsc{Threlfall}, Lehrbuch der Topologie, Chelsea,
Nueva York (1934).}
{H\ptbl\textsc{Seifert}, W\ptbl\textsc{Threlfall}, Lehrbuch der Topologie,
Chelsea, Nueva York (1934).}

\bibitem[4]{XIII.4}
\ifthenelse{\boolean{orig}}{J.P\ptbl}{J-P\ptbl}%
\textsc{Serre}, Cohomologie galoisienne, Lecture Notes
\no 5, Springer, Berlín (1964).

\bibitem[5]{XIII.5}
\ifthenelse{\boolean{orig}}{J.P\ptbl}{J-P\ptbl}%
\textsc{Serre}, Corps locaux, Hermann, París (1962).
\end{sgabibliography}

\backmatter
\def\indexname{Índice terminológico}
\printindex

\def\indexname{Índice de notaciones}

\begin{theindex}

\medskip\item $\Delta_{X/Y}$ o simplemente $\Delta$\dotfill\pageref{indnot:ab}
\medskip\item $\mathit{\Omega}^1_{X/Y}$\dotfill\pageref{indnot:ac}
\medskip\item $\cal{P}_{X/Y}^n$\dotfill\pageref{indnot:ad}, \pageref{indnot:ad1}
\medskip\item $\Delta_{X/Y}^n$\dotfill\pageref{indnot:ae}
\medskip\item $\goth{d}_{X/Y}$\dotfill\pageref{indnot:af}
\medskip\item $d_{X/S}^n$\dotfill\pageref{indnot:bb}
\medskip\item $\goth{g}_{X/S}$\dotfill\pageref{indnot:cb}
\medskip\item $\cal{C}(\pi)$\dotfill\pageref{indnot:eb}
\medskip\item $\pi$\dotfill\pageref{indnot:ec}
\medskip\item $\Prodash\cal{C}(\pi)$\dotfill\pageref{indnot:ed}
\medskip\item $\Gamma$\dotfill\pageref{indnot:ef}
\medskip\item $\pi_1(S,a)$\dotfill\pageref{indnot:eg}
\medskip\item $\pi_1(S;a,a')$\dotfill\pageref{indnot:eh}
\medskip\item $\cal{C}(S)$\dotfill\pageref{indnot:ei}
\medskip\item $\pi_1(f;a')$\dotfill\pageref{indnot:ej}
\medskip\item $\Ens$\dotfill\pageref{indnot:fb}
\medskip\item $\Cat$\dotfill\pageref{indnot:fc}
\medskip\item $\Ob(\cal{C})$\dotfill\pageref{indnot:fd}
\medskip\item $\Fl(\cal{C})$\dotfill\pageref{indnot:fe}
\medskip\item $\SheafHom(\cal{C},\cal{C}')$\dotfill\pageref{indnot:ff}
\medskip\item $\cal{C}^\circ$\dotfill\pageref{indnot:fg}
\medskip\item $\Cat_{/\cal{E}}$\dotfill\pageref{indnot:fh}
\medskip\item $\SheafHom_{\cal{E}/-}(\cal{F},\cal{G})$\dotfill\pageref{indnot:fi}
\medskip\item $v*u$\dotfill\pageref{indnot:fj}
\medskip\item $\cal{F}\times_{\cal{E}}\cal{G}$\dotfill\pageref{indnot:fk}
\medskip\item $\lambda^*\colon\Cat_{/\cal{E}}\to\Cat_{/\cal{E}'}$\dotfill\pageref{indnot:fl}
\medskip\item $\bf{\Gamma}(\cal{G}/\cal{E})$ y $\Gamma(\cal{G}/\cal{E})$\dotfill\pageref{indnot:fm}
\medskip\item $\cal{F}_S$\dotfill\pageref{indnot:fn}
\medskip\item $f^*_{\cal{F}}(\xi)$ o $f^*(\xi)$\dotfill\pageref{indnot:fo}
\medskip\item $\alpha_f(\xi)$\dotfill\pageref{indnot:fp}
\medskip\item $\SheafHom_\cart(\cal{F},\cal{G})$\dotfill\pageref{indnot:fq}
\medskip\item $\Cat^\cart_{/\cal{E}}$\dotfill\pageref{indnot:fr}
\medskip\item $\varprojLim\cal{F}/\cal{E}$\dotfill\pageref{indnot:fs}
\medskip\item $f_*^{\cal{F}}$ o~$f_*$\dotfill\pageref{indnot:ft}
\medskip\item $\Sch$\dotfill\pageref{indnot:hb}
\medskip\item $\bbmu_{n\,S}$\dotfill\pageref{indnot:kb}
\medskip\item $X^\an$\dotfill\pageref{indnot:lb}
\medskip\item $f^\an$\dotfill\pageref{indnot:lc}
\medskip\item $SF$ o $S(F)$\dotfill\pageref{indnot:mb}
\medskip\item $\H^1_\tame(U,F)$\dotfill\pageref{indnot:mc}
\medskip\item $\R^1_\tame g_*F$\dotfill\pageref{indnot:md}
\medskip\item $C_t((U,X)/S)$ o $C_t$\dotfill\pageref{indnot:me}
\medskip\item $\pi_1^\tame((U,X)/S,a)$ o $\pi_1^\tame(U,a)$ o $\pi_1^\tame(U)$\dotfill\pageref{indnot:mf}
\medskip\item $(g_*^\tame\Phi)_{T'}$\dotfill\pageref{indnot:mg}
\medskip\item $\H^0(V,C_V)^\Pi$\dotfill\pageref{indnot:mh}
\medskip\item $\pi_1^\LL(U,a)$\dotfill\pageref{indnot:mi}
\medskip\item $\pi_1'(X,a)$\dotfill\pageref{indnot:mj}
\medskip\item $\pi_1^\LL(X/S,g,\bar{s})$ o $\pi_1^\LL(X/S,g)$\dotfill\pageref{indnot:mk}
\medskip\item $\pi_1^\LL(X_{\bar{s}},a)_K$\dotfill\pageref{indnot:ml}
\medskip\item $\ZZ_\ell[1]$\dotfill\pageref{indnot:mm}
\par
\end{theindex}

\end{document}
